%% Generated by Sphinx.
\def\sphinxdocclass{report}
\IfFileExists{luatex85.sty}
 {\RequirePackage{luatex85}}
 {\ifdefined\luatexversion\ifnum\luatexversion>84\relax
  \PackageError{sphinx}
  {** With this LuaTeX (\the\luatexversion),Sphinx requires luatex85.sty **}
  {** Add the LaTeX package luatex85 to your TeX installation, and try again **}
  \endinput\fi\fi}
\documentclass[letterpaper,10pt,spanish]{sphinxmanual}
\ifdefined\pdfpxdimen
   \let\sphinxpxdimen\pdfpxdimen\else\newdimen\sphinxpxdimen
\fi \sphinxpxdimen=.75bp\relax
\ifdefined\pdfimageresolution
    \pdfimageresolution= \numexpr \dimexpr1in\relax/\sphinxpxdimen\relax
\fi
%% let collapsible pdf bookmarks panel have high depth per default
\PassOptionsToPackage{bookmarksdepth=5}{hyperref}
%% turn off hyperref patch of \index as sphinx.xdy xindy module takes care of
%% suitable \hyperpage mark-up, working around hyperref-xindy incompatibility
\PassOptionsToPackage{hyperindex=false}{hyperref}
%% memoir class requires extra handling
\makeatletter\@ifclassloaded{memoir}
{\ifdefined\memhyperindexfalse\memhyperindexfalse\fi}{}\makeatother

\PassOptionsToPackage{borderless}{sphinx}
\PassOptionsToPackage{colorrows}{sphinx}

\PassOptionsToPackage{warn}{textcomp}

\catcode`^^^^00a0\active\protected\def^^^^00a0{\leavevmode\nobreak\ }
\usepackage{cmap}
\usepackage{fontspec}
\defaultfontfeatures[\rmfamily,\sffamily,\ttfamily]{}
\usepackage{amsmath,amssymb,amstext}
\usepackage{polyglossia}
\setmainlanguage{spanish}



\setmainfont{FreeSerif}[
  Extension      = .otf,
  UprightFont    = *,
  ItalicFont     = *Italic,
  BoldFont       = *Bold,
  BoldItalicFont = *BoldItalic
]
\setsansfont{FreeSans}[
  Extension      = .otf,
  UprightFont    = *,
  ItalicFont     = *Oblique,
  BoldFont       = *Bold,
  BoldItalicFont = *BoldOblique,
]
\setmonofont{FreeMono}[Scale=0.9,
  Extension      = .otf,
  UprightFont    = *,
  ItalicFont     = *Oblique,
  BoldFont       = *Bold,
  BoldItalicFont = *BoldOblique,
]



\usepackage[Sonny]{fncychap}
\ChNameVar{\Large\normalfont\sffamily}
\ChTitleVar{\Large\normalfont\sffamily}
\usepackage{sphinx}

\fvset{fontsize=auto}
\usepackage{geometry}


% Include hyperref last.
\usepackage{hyperref}
% Fix anchor placement for figures with captions.
\usepackage{hypcap}% it must be loaded after hyperref.
% Set up styles of URL: it should be placed after hyperref.
\urlstyle{same}

\addto\captionsspanish{\renewcommand{\contentsname}{INTRO}}

\usepackage{sphinxmessages}
\setcounter{tocdepth}{1}



\title{Libro didáctico para el estudiante}
\date{18 de junio de 2026}
\release{}
\author{Julian Camilo Fonseca Romero}
\newcommand{\sphinxlogo}{\sphinxincludegraphics{logo_light.png}\par}
\renewcommand{\releasename}{}
\makeindex
\begin{document}

\pagestyle{empty}
\sphinxmaketitle
\pagestyle{plain}
\sphinxtableofcontents
\pagestyle{normal}
\phantomsection\label{\detokenize{index::doc}}


\sphinxAtStartPar
\sphinxincludegraphics{{aula115}.png}

\sphinxAtStartPar
El presente libro didáctico digital, denominado \sphinxstylestrong{aula 115}, es una herramienta utilizada por estudiantes de educación básica y media de la \sphinxstyleemphasis{Institución Educativa John F. Kennedy} en Puerto Boyacá (Colombia), con el propósito de facilitar el \sphinxstyleemphasis{aprendizaje en ofimática} mediante el uso de \sphinxstylestrong{software libre}, promoviendo así la productividad, la autonomía e “independencia tecnológica” en el \sphinxstyleemphasis{manejo de la información} del siglo XXI.


\begin{savenotes}\sphinxattablestart
\sphinxthistablewithglobalstyle
\centering
\begin{tabular}[t]{*{2}{\X{1}{2}}}
\sphinxtoprule
\sphinxtableatstartofbodyhook

\begin{savenotes}\sphinxattablestart
\sphinxthistablewithglobalstyle
\centering
\begin{tabulary}{\linewidth}[t]{T}
\sphinxtoprule
\sphinxtableatstartofbodyhook

\begin{sphinxfigure-in-table}
\centering

\noindent\sphinxincludegraphics[width=180\sphinxpxdimen]{{camilo.fonseca}.png}
\end{sphinxfigure-in-table}\relax
\\
\sphinxbottomrule
\end{tabulary}
\sphinxtableafterendhook\par
\sphinxattableend\end{savenotes}
&

\begin{savenotes}\sphinxattablestart
\sphinxthistablewithglobalstyle
\centering
\begin{tabular}[t]{*{1}{\X{1}{1}}}
\sphinxtoprule
\sphinxtableatstartofbodyhook
\sphinxAtStartPar
\sphinxstylestrong{DOCENTE:} Julian Camilo Fonseca Romero

\sphinxAtStartPar
\sphinxstylestrong{email:} camilo.fonseca@jfkennedy.edu.co
\begin{itemize}
\item {} 
\sphinxAtStartPar
Ingeniero en Mecatrónica (2012)

\item {} 
\sphinxAtStartPar
Especialista en Seguridad de la Información (2015)

\item {} 
\sphinxAtStartPar
Magister en Ingeniería de Software y Sistemas Informáticos (2017)

\end{itemize}
\\
\sphinxbottomrule
\end{tabular}
\sphinxtableafterendhook\par
\sphinxattableend\end{savenotes}
\\
\sphinxbottomrule
\end{tabular}
\sphinxtableafterendhook\par
\sphinxattableend\end{savenotes}

\begin{sphinxadmonition}{note}{PRINCIPIOS DE LA CLASE}
\begin{enumerate}
\sphinxsetlistlabels{\arabic}{enumi}{enumii}{}{.}%
\item {} 
\sphinxAtStartPar
“Se aprende haciendo”.

\item {} 
\sphinxAtStartPar
“Se aprende enseñando”.

\item {} 
\sphinxAtStartPar
“Todo se sustenta”.

\end{enumerate}
\end{sphinxadmonition}

\begin{sphinxadmonition}{note}{SECUENCIA DIDÁCTICA (MODELO ACTIVO)}
\begin{enumerate}
\sphinxsetlistlabels{\arabic}{enumi}{enumii}{}{.}%
\item {} 
\sphinxAtStartPar
\DUrole{sd-sphinx-override}{\DUrole{sd-badge}{\DUrole{sd-bg-light}{\DUrole{sd-bg-text-light}{Exploración}}}}

\item {} 
\sphinxAtStartPar
\DUrole{sd-sphinx-override}{\DUrole{sd-badge}{\DUrole{sd-bg-info}{\DUrole{sd-bg-text-info}{Estructuración}}}}.

\item {} 
\sphinxAtStartPar
\DUrole{sd-sphinx-override}{\DUrole{sd-badge}{\DUrole{sd-bg-warning}{\DUrole{sd-bg-text-warning}{Actividad práctica}}}}.

\item {} 
\sphinxAtStartPar
\DUrole{sd-sphinx-override}{\DUrole{sd-badge}{\DUrole{sd-bg-danger}{\DUrole{sd-bg-text-danger}{Transferencia}}}}.

\end{enumerate}
\end{sphinxadmonition}

\sphinxstepscope


\chapter{Información}
\label{\detokenize{0_intro/informacion:informacion}}\label{\detokenize{0_intro/informacion::doc}}
\sphinxAtStartPar
La siguiente información es útil para el estudiante.


\bigskip\hrule\bigskip



\section{Útiles escolares}
\label{\detokenize{0_intro/informacion:utiles-escolares}}\begin{itemize}
\item {} 
\sphinxAtStartPar
Cuaderno cocido (100 hojas y cuadriculado).

\item {} 
\sphinxAtStartPar
Lapiz, Borrador, Tajalapiz.

\item {} 
\sphinxAtStartPar
Esferos azul, negro.

\item {} 
\sphinxAtStartPar
Corrector.

\end{itemize}


\bigskip\hrule\bigskip



\section{Porcentaje de notas \%}
\label{\detokenize{0_intro/informacion:porcentaje-de-notas}}
\begin{sphinxadmonition}{note}{50\% + 40\% + 10\% = 100\%}
\begin{itemize}
\item {} 
\sphinxAtStartPar
\sphinxstylestrong{50\% COGNITIVA:} Guías didácticas.

\item {} 
\sphinxAtStartPar
\sphinxstylestrong{40\% PROCEDIMENTAL:} Resultado general prueba INSTRUIMOS.

\item {} 
\sphinxAtStartPar
\sphinxstylestrong{10\% ACTITUDINAL} = 5\% PERFIL Kennediano \sphinxstylestrong{+} 5\% Trabajo en equipo.

\end{itemize}
\end{sphinxadmonition}


\bigskip\hrule\bigskip



\section{Equivalencia calificación INSTRUIMOS}
\label{\detokenize{0_intro/informacion:equivalencia-calificacion-instruimos}}
\sphinxAtStartPar
La siguiente tabla muestra el equivalente del puntaje de nota Instruimos (1\sphinxhyphen{}100) respecto al valor de la nota en la Institución Educativa (1\sphinxhyphen{}5).


\begin{savenotes}\sphinxattablestart
\sphinxthistablewithglobalstyle
\centering
\begin{tabulary}{\linewidth}[t]{TT}
\sphinxtoprule
\sphinxstyletheadfamily 
\sphinxAtStartPar
Puntaje INSTRUIMOS
&\sphinxstyletheadfamily 
\sphinxAtStartPar
Nota IE
\\
\sphinxmidrule
\sphinxtableatstartofbodyhook
\sphinxAtStartPar
10\sphinxhyphen{}15
&
\sphinxAtStartPar
10
\\
\sphinxhline
\sphinxAtStartPar
16\sphinxhyphen{}25
&
\sphinxAtStartPar
15
\\
\sphinxhline
\sphinxAtStartPar
26\sphinxhyphen{}35
&
\sphinxAtStartPar
20
\\
\sphinxhline
\sphinxAtStartPar
36\sphinxhyphen{}40
&
\sphinxAtStartPar
25
\\
\sphinxhline
\sphinxAtStartPar
41\sphinxhyphen{}45
&
\sphinxAtStartPar
29
\\
\sphinxhline
\sphinxAtStartPar
46\sphinxhyphen{}55
&
\sphinxAtStartPar
30
\\
\sphinxhline
\sphinxAtStartPar
56\sphinxhyphen{}65
&
\sphinxAtStartPar
35
\\
\sphinxhline
\sphinxAtStartPar
66\sphinxhyphen{}75
&
\sphinxAtStartPar
40
\\
\sphinxhline
\sphinxAtStartPar
76\sphinxhyphen{}85
&
\sphinxAtStartPar
45
\\
\sphinxhline
\sphinxAtStartPar
86\sphinxhyphen{}95
&
\sphinxAtStartPar
48
\\
\sphinxhline
\sphinxAtStartPar
96\sphinxhyphen{}100
&
\sphinxAtStartPar
50
\\
\sphinxbottomrule
\end{tabulary}
\sphinxtableafterendhook\par
\sphinxattableend\end{savenotes}


\bigskip\hrule\bigskip



\section{Sello docente}
\label{\detokenize{0_intro/informacion:sello-docente}}\begin{itemize}
\item {} 
\sphinxAtStartPar
Estudiante coloca la \sphinxstylestrong{fecha}.

\item {} 
\sphinxAtStartPar
La \sphinxstylestrong{nota} puede ser cualitativa o cuantitativa.

\item {} 
\sphinxAtStartPar
Por lo general, cada sello se marca después de terminar una guía.

\end{itemize}


\begin{savenotes}\sphinxattablestart
\sphinxthistablewithglobalstyle
\centering
\begin{tabulary}{\linewidth}[t]{T}
\sphinxtoprule
\sphinxtableatstartofbodyhook

\begin{sphinxfigure-in-table}
\centering
\capstart
\noindent\sphinxincludegraphics[width=400\sphinxpxdimen]{{sello}.png}
\sphinxfigcaption{Sello docente}\label{\detokenize{0_intro/informacion:id1}}\end{sphinxfigure-in-table}\relax
\\
\sphinxbottomrule
\end{tabulary}
\sphinxtableafterendhook\par
\sphinxattableend\end{savenotes}


\bigskip\hrule\bigskip



\section{Acceso dentro y fuera del aula}
\label{\detokenize{0_intro/informacion:acceso-dentro-y-fuera-del-aula}}
\sphinxAtStartPar
Los \sphinxstylestrong{estudiantes} pueden acceder al \sphinxstylestrong{libro digital} desde cualquier lugar con acceso a \sphinxstylestrong{INTERNET}. También pueden acceder por medio de la INTRANET si estan conectados a la red del aula.

\sphinxAtStartPar
Si los estudiantes no tienen acceso a \sphinxstylestrong{INTERNET}, pueden descargar el \sphinxstylestrong{libro digital} previamente en sus equipos o smartphones.
\begin{itemize}
\item {} 
\sphinxAtStartPar
INTRANET: \sphinxurl{http://192.168.115.100}

\item {} 
\sphinxAtStartPar
INTERNET: \sphinxurl{https://aula115.readthedocs.io}

\end{itemize}


\bigskip\hrule\bigskip



\section{Modelo de almacenamiento Google WorkSpace®}
\label{\detokenize{0_intro/informacion:modelo-de-almacenamiento-google-workspace}}
\sphinxAtStartPar
La siguiente tabla presenta la \sphinxstylestrong{capacidad máxima de almacenamiento} disponible para cada miembro de la comunidad educativa, de acuerdo con su rol dentro de la \sphinxhref{https://www.jfkennedy.edu.co}{Institución Educativa John F. Kennedy}.


\begin{savenotes}\sphinxattablestart
\sphinxthistablewithglobalstyle
\centering
\begin{tabulary}{\linewidth}[t]{TTTTT}
\sphinxtoprule
\sphinxstyletheadfamily 
\sphinxAtStartPar
rol
&\sphinxstyletheadfamily 
\sphinxAtStartPar
ACTIVO
&\sphinxstyletheadfamily 
\sphinxAtStartPar
EGRESADO
&\sphinxstyletheadfamily 
\sphinxAtStartPar
RETIRADO
&\sphinxstyletheadfamily 
\sphinxAtStartPar
PENSIONADO
\\
\sphinxmidrule
\sphinxtableatstartofbodyhook
\sphinxAtStartPar
\sphinxstylestrong{ADMINISTRATIVO}
&
\sphinxAtStartPar
1TB
&
\sphinxAtStartPar
\sphinxstylestrong{N/A}
&
\sphinxAtStartPar
16GB
&
\sphinxAtStartPar
16GB
\\
\sphinxhline
\sphinxAtStartPar
\sphinxstylestrong{DOCENTE}
&
\sphinxAtStartPar
256GB
&
\sphinxAtStartPar
\sphinxstylestrong{N/A}
&
\sphinxAtStartPar
16GB
&
\sphinxAtStartPar
16GB
\\
\sphinxhline
\sphinxAtStartPar
\sphinxstylestrong{ESTUDIANTE}
&
\sphinxAtStartPar
16GB
&
\sphinxAtStartPar
16GB
&
\sphinxAtStartPar
\sphinxguilabel{X}
&
\sphinxAtStartPar
\sphinxstylestrong{N/A}
\\
\sphinxbottomrule
\end{tabulary}
\sphinxtableafterendhook\par
\sphinxattableend\end{savenotes}

\sphinxAtStartPar
\sphinxguilabel{X} = Suspensión de la cuenta


\bigskip\hrule\bigskip



\section{Acceso al aula 115 en contrajornada}
\label{\detokenize{0_intro/informacion:acceso-al-aula-115-en-contrajornada}}
\sphinxAtStartPar
Para \sphinxstylestrong{asegurar y mantener la correcta funcionalidad de los equipos}, se debe tener un control y seguimiento en el aula. Es por eso que los estudiantes que deseen \sphinxstyleemphasis{\sphinxstylestrong{hacer uso del aula de informática 115 en contra\sphinxhyphen{}jornada}}, deben solicitar el acceso enviando un correo a:
\begin{quote}\begin{description}
\sphinxlineitem{Cuenta}
\sphinxAtStartPar
camilo.fonseca@jfkennedy.edu.co

\sphinxlineitem{Asunto}
\sphinxAtStartPar
Autorización acceso aula 115 en contra\sphinxhyphen{}jornada

\end{description}\end{quote}

\sphinxAtStartPar
Los horarios en contrajornada son:
\begin{itemize}
\item {} 
\sphinxAtStartPar
\sphinxstylestrong{Mañana:} 10AM\sphinxhyphen{}12PM

\item {} 
\sphinxAtStartPar
\sphinxstylestrong{Tarde:} 01PM\sphinxhyphen{}03PM

\end{itemize}


\bigskip\hrule\bigskip



\section{Acceso de estudiantes al PC}
\label{\detokenize{0_intro/informacion:acceso-de-estudiantes-al-pc}}

\begin{savenotes}\sphinxattablestart
\sphinxthistablewithglobalstyle
\centering
\begin{tabular}[t]{*{1}{\X{1}{1}}}
\sphinxtoprule
\sphinxtableatstartofbodyhook
\begin{sphinxVerbatimintable}[commandchars=\\\{\},numbers=left,firstnumber=1,stepnumber=1]
\PYG{n}{cursoActual} \PYG{o}{=} \PYG{n+nb}{input}\PYG{p}{(}\PYG{l+s+s2}{\PYGZdq{}}\PYG{l+s+s2}{CURSO: }\PYG{l+s+s2}{\PYGZdq{}}\PYG{p}{)}

\PYG{k}{if} \PYG{p}{(}\PYG{n}{cursoActual} \PYG{o}{==} \PYG{l+s+s2}{\PYGZdq{}}\PYG{l+s+s2}{8A}\PYG{l+s+s2}{\PYGZdq{}} \PYG{o+ow}{or} \PYG{n}{cursoActual} \PYG{o}{==} \PYG{l+s+s2}{\PYGZdq{}}\PYG{l+s+s2}{8B}\PYG{l+s+s2}{\PYGZdq{}}\PYG{p}{)}\PYG{p}{:}
    \PYG{n}{usuarioEstudiante} \PYG{o}{=} \PYG{l+s+s2}{\PYGZdq{}}\PYG{l+s+s2}{usuario}\PYG{l+s+s2}{\PYGZdq{}}

\PYG{k}{elif} \PYG{p}{(}\PYG{n}{cursoActual} \PYG{o}{==} \PYG{l+s+s2}{\PYGZdq{}}\PYG{l+s+s2}{10A}\PYG{l+s+s2}{\PYGZdq{}}\PYG{p}{)}\PYG{p}{:}
    \PYG{n}{usuarioEstudiante} \PYG{o}{=} \PYG{l+s+s2}{\PYGZdq{}}\PYG{l+s+s2}{10A}\PYG{l+s+s2}{\PYGZdq{}}
    \PYG{n+nb}{print}\PYG{p}{(}\PYG{l+s+s2}{\PYGZdq{}}\PYG{l+s+s2}{dar contraseña de acceso a 10A}\PYG{l+s+s2}{\PYGZdq{}}\PYG{p}{)}

\PYG{k}{elif} \PYG{p}{(}\PYG{n}{cursoActual} \PYG{o}{==} \PYG{l+s+s2}{\PYGZdq{}}\PYG{l+s+s2}{10B}\PYG{l+s+s2}{\PYGZdq{}}\PYG{p}{)}\PYG{p}{:}
    \PYG{n}{usuarioEstudiante} \PYG{o}{=} \PYG{l+s+s2}{\PYGZdq{}}\PYG{l+s+s2}{10B}\PYG{l+s+s2}{\PYGZdq{}}
    \PYG{n+nb}{print}\PYG{p}{(}\PYG{l+s+s2}{\PYGZdq{}}\PYG{l+s+s2}{dar contraseña de acceso a 10B}\PYG{l+s+s2}{\PYGZdq{}}\PYG{p}{)}

\PYG{k}{elif} \PYG{p}{(}\PYG{n}{cursoActual} \PYG{o}{==} \PYG{l+s+s2}{\PYGZdq{}}\PYG{l+s+s2}{11A}\PYG{l+s+s2}{\PYGZdq{}}\PYG{p}{)}\PYG{p}{:}
    \PYG{n}{usuarioEstudiante} \PYG{o}{=} \PYG{l+s+s2}{\PYGZdq{}}\PYG{l+s+s2}{11A}\PYG{l+s+s2}{\PYGZdq{}}
    \PYG{n+nb}{print}\PYG{p}{(}\PYG{l+s+s2}{\PYGZdq{}}\PYG{l+s+s2}{dar contraseña de acceso a 11A}\PYG{l+s+s2}{\PYGZdq{}}\PYG{p}{)}

\PYG{k}{elif} \PYG{p}{(}\PYG{n}{cursoActual} \PYG{o}{==} \PYG{l+s+s2}{\PYGZdq{}}\PYG{l+s+s2}{11B}\PYG{l+s+s2}{\PYGZdq{}}\PYG{p}{)}\PYG{p}{:}
    \PYG{n}{usuarioEstudiante} \PYG{o}{=} \PYG{l+s+s2}{\PYGZdq{}}\PYG{l+s+s2}{11B}\PYG{l+s+s2}{\PYGZdq{}}
    \PYG{n+nb}{print}\PYG{p}{(}\PYG{l+s+s2}{\PYGZdq{}}\PYG{l+s+s2}{dar contraseña de acceso a 11B}\PYG{l+s+s2}{\PYGZdq{}}\PYG{p}{)}

\PYG{k}{else}\PYG{p}{:}
    \PYG{n}{usuarioEstudiante} \PYG{o}{=} \PYG{l+s+s2}{\PYGZdq{}}\PYG{l+s+s2}{NO INFO}\PYG{l+s+s2}{\PYGZdq{}}
\end{sphinxVerbatimintable}
\\
\sphinxbottomrule
\end{tabular}
\sphinxtableafterendhook\par
\sphinxattableend\end{savenotes}


\bigskip\hrule\bigskip



\section{Estructura carpeta de evidencias}
\label{\detokenize{0_intro/informacion:estructura-carpeta-de-evidencias}}
\sphinxAtStartPar
Los \sphinxstylestrong{estudiantes} tienen la siguiente estructura de \sphinxstylestrong{carpeta de evidencias}:


\begin{savenotes}\sphinxattablestart
\sphinxthistablewithglobalstyle
\centering
\begin{tabulary}{\linewidth}[t]{T}
\sphinxtoprule
\sphinxtableatstartofbodyhook

\begin{sphinxfigure-in-table}
\centering
\capstart
\noindent\sphinxincludegraphics[width=180\sphinxpxdimen]{{image01}.png}
\sphinxfigcaption{Carpeta de evidencias}\label{\detokenize{0_intro/informacion:id2}}\end{sphinxfigure-in-table}\relax
\\
\sphinxbottomrule
\end{tabulary}
\sphinxtableafterendhook\par
\sphinxattableend\end{savenotes}

\sphinxAtStartPar
Los \sphinxstylestrong{archivos} que se trabajen en clase, deben ser clasificados en las carpetas según el periodo y la asignatura.


\bigskip\hrule\bigskip



\section{Actividades inicio de año escolar}
\label{\detokenize{0_intro/informacion:actividades-inicio-de-ano-escolar}}\begin{itemize}
\item {} 
\sphinxAtStartPar
 ASIGNAR equipos de pareja (hombre y mujer).

\item {} 
\sphinxAtStartPar
 DAR hoja para que anoten sus nombres, correo institucional y equipo asignado por pareja.

\item {} 
\sphinxAtStartPar
 PRESENTACIÓN del docente.

\item {} 
\sphinxAtStartPar
 Explicación \sphinxstylestrong{útiles escolares}.

\item {} 
\sphinxAtStartPar
 Explicación de \sphinxstylestrong{porcentaje de notas \%}.

\item {} 
\sphinxAtStartPar
 Explicación \sphinxstylestrong{equivalencia de notas INSTRUIMOS}.

\item {} 
\sphinxAtStartPar
 Explicación \sphinxstylestrong{sello docente}.

\item {} 
\sphinxAtStartPar
 Explicación \sphinxstylestrong{acceso dentro y fuera del aula}.

\item {} 
\sphinxAtStartPar
 Explicación \sphinxstylestrong{modelo de almacenamiento}.

\item {} 
\sphinxAtStartPar
 Explicación \sphinxstylestrong{acceso al aula 115 en contrajornada}.

\item {} 
\sphinxAtStartPar
 \sphinxstylestrong{Reglas del aula} (cada estudiante lee una regla y escriben las reglas en su cuaderno).

\item {} 
\sphinxAtStartPar
 FIRMA DE COMPROMISOS en el cuaderno (Desbloqueo de equipo).

\item {} 
\sphinxAtStartPar
 Acceso de estudiantes al PC.

\item {} 
\sphinxAtStartPar
 Estructura carpeta de evidencias.

\end{itemize}


\bigskip\hrule\bigskip



\section{DESEMPEÑOS}
\label{\detokenize{0_intro/informacion:desempenos}}

\subsection{Lógica computacional 8}
\label{\detokenize{0_intro/informacion:logica-computacional-8}}\begin{quote}\begin{description}
\sphinxlineitem{PERIODO 01}
\sphinxAtStartPar
Explica los conceptos de hardware, software, sistemas computacionales, informática y los relaciona con la definición de \sphinxstylestrong{lógica computacional}.

\sphinxlineitem{PERIODO 02}
\sphinxAtStartPar
Clasifica los elementos de la Jerarquía de datos para facilitar el procesamiento y la comunicación de ideas en contextos específicos.

\sphinxlineitem{PERIODO 03}
\sphinxlineitem{PERIODO 04}
\end{description}\end{quote}


\subsection{Tecnología e Informática 10}
\label{\detokenize{0_intro/informacion:tecnologia-e-informatica-10}}\begin{quote}\begin{description}
\sphinxlineitem{PERIODO 01}
\sphinxAtStartPar
Sistematiza la evolución de las redes, la computación personal y el software, analizando de qué manera estas innovaciones incidieron en los cambios estructurales de la sociedad y la cultura a lo largo de la historia.

\sphinxlineitem{PERIODO 02}
\sphinxAtStartPar
Conoce los conceptos introductorios a las temáticas de redes, mantenimiento y ofimática para evaluar la viabilidad técnica y social de diversas alternativas de solución, asegurando que el tema elegido para su anteproyecto de grado favorezca la productividad y el bienestar en su contexto específico.

\sphinxlineitem{PERIODO 03}
\sphinxlineitem{PERIODO 04}
\end{description}\end{quote}


\subsection{Ofimática 10}
\label{\detokenize{0_intro/informacion:ofimatica-10}}\begin{quote}\begin{description}
\sphinxlineitem{PERIODO 01}
\sphinxAtStartPar
Conoce el concepto de ofimática, argumenta las ventajas y desventajas de su uso con software libre en especial con la suite LibreOffice.

\sphinxlineitem{PERIODO 02}
\sphinxlineitem{PERIODO 03}
\sphinxlineitem{PERIODO 04}
\end{description}\end{quote}


\subsection{Ofimática 11}
\label{\detokenize{0_intro/informacion:ofimatica-11}}\begin{quote}\begin{description}
\sphinxlineitem{PERIODO 01}
\sphinxAtStartPar
Diferencia las partes físicas de la computadora de sus programas y comprende cómo la CPU y la memoria RAM trabajan con los datos para generar información útil. Identifica la organización jerárquica de la información y utiliza correctamente los términos hardware, software y programa en contextos técnicos sencillos.

\sphinxlineitem{PERIODO 02}
\sphinxlineitem{PERIODO 03}
\sphinxlineitem{PERIODO 04}
\end{description}\end{quote}

\sphinxstepscope


\chapter{Reglas del aula}
\label{\detokenize{0_intro/reglas:reglas-del-aula}}\label{\detokenize{0_intro/reglas::doc}}
\sphinxAtStartPar
A continuación se muestran las \sphinxstylestrong{reglas del aula} que los estudiantes se comprometen a cumplir:


\bigskip\hrule\bigskip



\section{Reglas del aula 115}
\label{\detokenize{0_intro/reglas:reglas-del-aula-115}}\begin{quote}\begin{description}
\sphinxlineitem{REGLA 01}
\sphinxAtStartPar
Las personas que hagan uso del aula 115, deben \sphinxstylestrong{atender} y \sphinxstylestrong{obedecer} las instrucciones del profesor (o encargado del aula), así como mantener respeto, buen comportamiento y disciplina.

\sphinxlineitem{REGLA 02}
\sphinxAtStartPar
Los estudiantes deben acomodarse de la siguiente manera:
\begin{itemize}
\item {} 
\sphinxAtStartPar
Bolsos en la parte de atrás de la silla.

\item {} 
\sphinxAtStartPar
Cuadernos y útiles en el soporte frontal.

\item {} 
\sphinxAtStartPar
Si hay evaluación, los bolsos deben estar en el tablero.

\end{itemize}

\sphinxlineitem{REGLA 03}
\sphinxAtStartPar
Para el aseo al final de la jornada escolar:
\begin{itemize}
\item {} 
\sphinxAtStartPar
Primero se dejan las sillas organizadas en la esquina del aula.

\item {} 
\sphinxAtStartPar
Posteriormente se procede a barrer y luego trapear.

\item {} 
\sphinxAtStartPar
Se debe botar la basura y el agua sucia.

\item {} 
\sphinxAtStartPar
También se debe reciclar las botellas plásticas.

\end{itemize}

\sphinxlineitem{REGLA 04}
\sphinxAtStartPar
Cualquier duda sobre el uso del aula, preguntar al docente o al encargado.

\sphinxlineitem{REGLA 05}
\sphinxAtStartPar
Cumplimiento de la hora establecida:
\begin{itemize}
\item {} 
\sphinxAtStartPar
De entrada los estudiantes deben formar dos filas. Una fila de hombres y una fila de mujeres.

\item {} 
\sphinxAtStartPar
Para permitir la entrada los estudiantes deben cumplir:
\begin{itemize}
\item {} 
\sphinxAtStartPar
Portar bien el uniforme.

\item {} 
\sphinxAtStartPar
Celulares guardados en bolsos y en silencio.

\item {} 
\sphinxAtStartPar
No tener ni estar consumiendo comida o bebida alguna.

\end{itemize}

\end{itemize}

\sphinxlineitem{REGLA 06}
\sphinxAtStartPar
Llegada tarde:
\begin{itemize}
\item {} 
\sphinxAtStartPar
Una vez cumplida la hora de entrada, se cierra la puerta y empieza la clase.
\begin{itemize}
\item {} 
\sphinxAtStartPar
Se hace llamado a lista.

\item {} 
\sphinxAtStartPar
Los estudiantes que lleguen tarde permanecen afuera sin distraer, hacer ruido o desorden y se les coloca la falla por tardanza.

\item {} 
\sphinxAtStartPar
Una vez terminado el llamado a lista, los estudiantes faltantes ingresan en silencio al aula. (Profesores dan 5 min para hacer aseo antes de entrar a la siguiente clase).

\end{itemize}

\end{itemize}

\sphinxlineitem{REGLA 07}
\sphinxAtStartPar
El aula de informática es exclusivamente para \sphinxstylestrong{realizar trabajos académicos} (que necesiten de un computador). Los trabajos deben estar relacionados con la Institución Educativa John F. Kennedy. Está prohibido hacer trabajos de cualquier otro tipo o institución ajenos.

\sphinxlineitem{REGLA 08}
\sphinxAtStartPar
Si el estudiante es identificado \sphinxstylestrong{realizando cualquier actividad ajena a lo establecido por el docente} o incumpliendo cualquiera de las presentes reglas, su equipo será bloqueado.

\end{description}\end{quote}


\bigskip\hrule\bigskip



\section{PRØHIBICIONES}
\label{\detokenize{0_intro/reglas:prohibiciones}}\begin{itemize}
\item {} 
\sphinxAtStartPar
Se \sphinxstylestrong{restringe la salida al baño} durante la clase, salvo en casos de emergencias médicas debidamente justificadas mediante certificado, siguiendo las directrices de coordinación.

\item {} 
\sphinxAtStartPar
Consumir o ingresar cualquier tipo de \sphinxstylestrong{comida} o \sphinxstylestrong{bebida} en el aula de informática. Los residuos de comida o bebida caen sobre los equipos.

\item {} 
\sphinxAtStartPar
Dejar bolsos o cualquier otro elemento encima de las mesas de los equipos. (\sphinxstyleemphasis{Los cables de alimentación se pueden desconectar y esto daña los equipos}).

\item {} 
\sphinxAtStartPar
\sphinxstylestrong{Acostarse} en el piso del aula.

\item {} 
\sphinxAtStartPar
\sphinxstylestrong{Sentarse de forma incorrecta en las sillas} del aula o \sphinxstylestrong{sentarse en las mesas}.

\item {} 
\sphinxAtStartPar
\sphinxstylestrong{Gritos extemporáneos} o hablar en \sphinxstylestrong{voz alta}.

\item {} 
\sphinxAtStartPar
\sphinxstylestrong{Dañar o escribir} en los muebles del aula. (lápiz, bolígrafos, pinturas, etc.)

\item {} 
\sphinxAtStartPar
\sphinxstylestrong{Maquillarse} en el aula.

\item {} 
\sphinxAtStartPar
Realizar cualquier tipo de bailes o \sphinxstylestrong{expresiones corporales} en el aula.

\item {} 
\sphinxAtStartPar
\sphinxstylestrong{Consumir} cualquier tipo de droga o estupefaciente.

\item {} 
\sphinxAtStartPar
\sphinxstylestrong{Arrojar basura} en el aula.

\item {} 
\sphinxAtStartPar
\sphinxstylestrong{No hacer el aseo} o hacerlo de forma mediocre.

\item {} 
\sphinxAtStartPar
Esconder y robar cualquier objeto del aula o de los compañeros.

\item {} 
\sphinxAtStartPar
Ocupar el puesto del compañero.

\item {} 
\sphinxAtStartPar
Realizar trabajos diferentes a la asignatura en el aula.

\item {} 
\sphinxAtStartPar
Quitarse los zapatos.

\item {} 
\sphinxAtStartPar
Uso inapropiado del aula:
\begin{itemize}
\item {} 
\sphinxAtStartPar
Invadir el espacio del docente.

\item {} 
\sphinxAtStartPar
Bajar los tacos (breakers).

\item {} 
\sphinxAtStartPar
Escuchar música en el computador o celular sin autorización del docente.

\item {} 
\sphinxAtStartPar
Ver videos en el computador o celular sin autorización del docente.

\item {} 
\sphinxAtStartPar
Jugar videojuegos en el computador o celular sin autorización del docente.

\item {} 
\sphinxAtStartPar
Manipular, portar o \sphinxstylestrong{dejar visible el celular} sin previa autorización del docente (jugar, redes sociales, etc.)

\item {} 
\sphinxAtStartPar
Tomar fotos, selfis o videos sin autorización del docente.

\item {} 
\sphinxAtStartPar
Manipular aires acondicionados, ventiladores, televisor, tablero.

\item {} 
\sphinxAtStartPar
Acceder al gabinete, armario y escritorio del docente.

\item {} 
\sphinxAtStartPar
Desarmar o dar golpes a los equipos.

\item {} 
\sphinxAtStartPar
Instalar software en los equipos sin previa autorización del profesor.

\item {} 
\sphinxAtStartPar
Borrar archivos de los equipos.

\end{itemize}

\end{itemize}


\bigskip\hrule\bigskip



\section{¿Qué pasa si se incumplen las reglas?}
\label{\detokenize{0_intro/reglas:que-pasa-si-se-incumplen-las-reglas}}
\sphinxAtStartPar
Si el estudiante incumple alguna de las reglas, se procede con las siguientes \sphinxstylestrong{acciones formativas}:
\begin{itemize}
\item {} 
\sphinxAtStartPar
Disminución de valor de la nota en \sphinxstylestrong{perfil Kennediano}.

\item {} 
\sphinxAtStartPar
Disminución de valor de la nota \sphinxstylestrong{trabajo en equipo}.

\item {} 
\sphinxAtStartPar
Posible cambio de puesto de trabajo.

\end{itemize}

\sphinxAtStartPar
Mientras que el estudiante NO CUMPLA con las \sphinxstylestrong{acciones formativas}, se procede a:
\begin{itemize}
\item {} 
\sphinxAtStartPar
No calificar o evaluar cualquiera de las actividades académicas en proceso.

\item {} 
\sphinxAtStartPar
No firmar paz y salvo.

\end{itemize}


\bigskip\hrule\bigskip



\section{¿Qué pasa si el incumplimiento es reiterativo?}
\label{\detokenize{0_intro/reglas:que-pasa-si-el-incumplimiento-es-reiterativo}}

\begin{savenotes}\sphinxattablestart
\sphinxthistablewithglobalstyle
\centering
\begin{tabulary}{\linewidth}[t]{T}
\sphinxtoprule
\sphinxtableatstartofbodyhook

\begin{sphinxfigure-in-table}
\centering
\capstart
\noindent\sphinxincludegraphics[width=300\sphinxpxdimen]{{image011}.png}
\sphinxfigcaption{Sin excusas}\label{\detokenize{0_intro/reglas:id1}}\end{sphinxfigure-in-table}\relax
\\
\sphinxbottomrule
\end{tabulary}
\sphinxtableafterendhook\par
\sphinxattableend\end{savenotes}

\sphinxAtStartPar
En caso que el estudiante incumpla alguna de las reglas continuamente, se iniciará \sphinxstylestrong{proceso disciplinario} citando al acudiente, representante de curso o director de curso.

\sphinxstepscope


\chapter{Descargas}
\label{\detokenize{0_intro/descargas:descargas}}\label{\detokenize{0_intro/descargas::doc}}
\sphinxAtStartPar
Identifique el recurso y PULSE el botón \sphinxstylestrong{Descargar}.
\begin{itemize}
\item {} 
\sphinxAtStartPar
{\hyperref[\detokenize{0_intro/descargas::doc}]{\sphinxcrossref{\DUrole{std}{\DUrole{std-doc}{\DUrole{sd-sphinx-override}{\DUrole{sd-badge}{\DUrole{sd-bg-primary}{\DUrole{sd-bg-text-primary}{Descargar}}}}}}}}} Servidor principal

\item {} 
\sphinxAtStartPar
{\hyperref[\detokenize{0_intro/descargas::doc}]{\sphinxcrossref{\DUrole{std}{\DUrole{std-doc}{\DUrole{sd-sphinx-override}{\DUrole{sd-badge}{\DUrole{sd-bg-secondary}{\DUrole{sd-bg-text-secondary}{Descargar}}}}}}}}} Servidor secundario

\end{itemize}


\bigskip\hrule\bigskip



\section{Institucional}
\label{\detokenize{0_intro/descargas:institucional}}\begin{itemize}
\item {} 
\sphinxAtStartPar
MANUAL DE CONVIVENCIA \DUrole{xref}{\DUrole{download}{\DUrole{myst}{\DUrole{sd-sphinx-override}{\DUrole{sd-badge}{\DUrole{sd-bg-primary}{\DUrole{sd-bg-text-primary}{Descargar}}}}}}} \sphinxhref{https://drive.google.com/file/d/101ObkDOMInftIxeh4Q\_yyKYM9gwauDN7/view?usp=drive\_link}{\DUrole{sd-sphinx-override}{\DUrole{sd-badge}{\DUrole{sd-bg-secondary}{\DUrole{sd-bg-text-secondary}{Descargar}}}}}

\item {} 
\sphinxAtStartPar
LIBRO DIGITAL: aula 115 (PDF) \DUrole{xref}{\DUrole{download}{\DUrole{myst}{\DUrole{sd-sphinx-override}{\DUrole{sd-badge}{\DUrole{sd-bg-primary}{\DUrole{sd-bg-text-primary}{Descargar}}}}}}}

\end{itemize}


\bigskip\hrule\bigskip



\section{Proyecto de grado}
\label{\detokenize{0_intro/descargas:proyecto-de-grado}}\begin{itemize}
\item {} 

\begin{savenotes}\sphinxattablestart
\sphinxthistablewithglobalstyle
\raggedright
\begin{tabulary}{\linewidth}[t]{TTTTT}
\sphinxtoprule
\sphinxstyletheadfamily 
\sphinxAtStartPar
Plantilla
&\sphinxstyletheadfamily 
\sphinxAtStartPar
Microsoft Word
&\sphinxstyletheadfamily 
\sphinxAtStartPar
LibreOffice Writer
&\sphinxstyletheadfamily 
\sphinxAtStartPar
Google Docs
&\sphinxstyletheadfamily 
\sphinxAtStartPar
LaTeX
\\
\sphinxmidrule
\sphinxtableatstartofbodyhook
\sphinxAtStartPar
Marco teórico
&&
\sphinxAtStartPar
\DUrole{xref}{\DUrole{download}{\DUrole{myst}{\DUrole{sd-sphinx-override}{\DUrole{sd-badge}{\DUrole{sd-bg-primary}{\DUrole{sd-bg-text-primary}{Descargar}}}}}}} \sphinxhref{https://drive.google.com/file/d/1CgOI-NqmOQrc3K\_oq9lFJrJ5Q9hqbKQH/view?usp=drive\_link}{\DUrole{sd-sphinx-override}{\DUrole{sd-badge}{\DUrole{sd-bg-secondary}{\DUrole{sd-bg-text-secondary}{Descargar}}}}}
&&\\
\sphinxhline
\sphinxAtStartPar
Anteproyecto
&&
\sphinxAtStartPar
\DUrole{xref}{\DUrole{download}{\DUrole{myst}{\DUrole{sd-sphinx-override}{\DUrole{sd-badge}{\DUrole{sd-bg-primary}{\DUrole{sd-bg-text-primary}{Descargar}}}}}}} \sphinxhref{https://drive.google.com/file/d/1klPSl7wJwT1KyA956JvUN11-gPWVMJKR/view?usp=drive\_link}{\DUrole{sd-sphinx-override}{\DUrole{sd-badge}{\DUrole{sd-bg-secondary}{\DUrole{sd-bg-text-secondary}{Descargar}}}}}
&&\\
\sphinxhline
\sphinxAtStartPar
Proyecto
&
\sphinxAtStartPar
\DUrole{xref}{\DUrole{download}{\DUrole{myst}{\DUrole{sd-sphinx-override}{\DUrole{sd-badge}{\DUrole{sd-bg-primary}{\DUrole{sd-bg-text-primary}{Descargar}}}}}}} \sphinxhref{https://docs.google.com/document/d/1Ela1bpCYS60FAh82kLc8D9H6ZBJNU8nZ/edit?usp=drive\_link\&amp;ouid=112729153804362979178\&amp;rtpof=true\&amp;sd=true}{\DUrole{sd-sphinx-override}{\DUrole{sd-badge}{\DUrole{sd-bg-secondary}{\DUrole{sd-bg-text-secondary}{Descargar}}}}}
&&&\\
\sphinxbottomrule
\end{tabulary}
\sphinxtableafterendhook\par
\sphinxattableend\end{savenotes}

\end{itemize}

\sphinxstepscope


\chapter{Lógica computacional 8}
\label{\detokenize{1_guias/8logica/G820:logica-computacional-8}}\label{\detokenize{1_guias/8logica/G820::doc}}
\sphinxstepscope


\section{Guía 821: Concepto Lógica computacional (P1)}
\label{\detokenize{1_guias/8logica/G821:guia-821-concepto-logica-computacional-p1}}\label{\detokenize{1_guias/8logica/G821::doc}}

\begin{savenotes}\sphinxattablestart
\sphinxthistablewithglobalstyle
\centering
\begin{tabular}[t]{\X{35}{100}\X{65}{100}}
\sphinxtoprule
\sphinxtableatstartofbodyhook

\begin{savenotes}\sphinxattablestart
\sphinxthistablewithglobalstyle
\centering
\begin{tabulary}{\linewidth}[t]{T}
\sphinxtoprule
\sphinxtableatstartofbodyhook
\sphinxAtStartPar
\sphinxincludegraphics{{escudo}.png}
\\
\sphinxbottomrule
\end{tabulary}
\sphinxtableafterendhook\par
\sphinxattableend\end{savenotes}
&

\begin{savenotes}\sphinxattablestart
\sphinxthistablewithglobalstyle
\raggedright
\begin{tabular}[t]{*{1}{\X{1}{1}}}
\sphinxtoprule
\sphinxtableatstartofbodyhook
\sphinxAtStartPar
\sphinxstylestrong{INSTITUCIÓN EDUCATIVA}: John F. Kennedy
\begin{itemize}
\item {} 
\sphinxAtStartPar
\sphinxstylestrong{DOCENTE}: Julian Camilo Fonseca Romero

\item {} 
\sphinxAtStartPar
\sphinxstylestrong{ASIGNATURA}: Lógica computacional

\item {} 
\sphinxAtStartPar
\sphinxstylestrong{GRADO}: OCTAVO

\item {} 
\sphinxAtStartPar
\sphinxstylestrong{PERIODO}: 01

\item {} 
\sphinxAtStartPar
\sphinxstylestrong{TIEMPO}: 10 semanas (1 hora / semana)

\item {} 
\sphinxAtStartPar
\sphinxstylestrong{TIEMPO DE TRABAJO}: 10/10 horas

\end{itemize}
\\
\sphinxbottomrule
\end{tabular}
\sphinxtableafterendhook\par
\sphinxattableend\end{savenotes}
\\
\sphinxbottomrule
\end{tabular}
\sphinxtableafterendhook\par
\sphinxattableend\end{savenotes}

\begin{sphinxuseclass}{sd-tab-set}
\begin{sphinxuseclass}{sd-tab-item}\subsubsection*{Objetivo}

\begin{sphinxuseclass}{sd-tab-content}
\sphinxAtStartPar
\DUrole{sd-sphinx-override}{\DUrole{sd-badge}{\DUrole{sd-bg-light}{\DUrole{sd-bg-text-light}{Exploración}}}}
\DUrole{sd-sphinx-override}{\DUrole{sd-badge}{\DUrole{sd-bg-info}{\DUrole{sd-bg-text-info}{Estructuración}}}}
\DUrole{sd-sphinx-override}{\DUrole{sd-badge}{\DUrole{sd-bg-warning}{\DUrole{sd-bg-text-warning}{Actividad práctica}}}}
\DUrole{sd-sphinx-override}{\DUrole{sd-badge}{\DUrole{sd-bg-danger}{\DUrole{sd-bg-text-danger}{Transferencia}}}}
\begin{itemize}
\item {} 
\sphinxAtStartPar
\sphinxstylestrong{Temática}: Concepto de lógica computacional.

\item {} 
\sphinxAtStartPar
\sphinxstylestrong{Objetivo de aprendizaje}: Repasar los conceptos de hardware, software, sistemas computacionales y el razonamiento lógico para comprender el concepto de lógica computacional.

\end{itemize}

\end{sphinxuseclass}
\end{sphinxuseclass}
\begin{sphinxuseclass}{sd-tab-item}\subsubsection*{Orientaciones curriculares}

\begin{sphinxuseclass}{sd-tab-content}\begin{itemize}
\item {} 
\sphinxAtStartPar
\sphinxstylestrong{Componente}: Naturaleza y Evolución de la T\&I.

\item {} 
\sphinxAtStartPar
\sphinxstylestrong{Competencia}: Relaciono saberes, conocimientos tecnológicos e informáticos con los conocimientos de otras disciplinas.

\item {} 
\sphinxAtStartPar
\sphinxstylestrong{Evidencia de aprendizaje}: Comprendo los principios y conocimientos tecnológicos e informáticos que hacen posible el funcionamiento de productos tecnológicos actuales.

\end{itemize}

\end{sphinxuseclass}
\end{sphinxuseclass}
\begin{sphinxuseclass}{sd-tab-item}\subsubsection*{Criterios de evaluación}

\begin{sphinxuseclass}{sd-tab-content}\begin{itemize}
\item {} 
\sphinxAtStartPar
Actividades desarrolladas en su totalidad.

\item {} 
\sphinxAtStartPar
Participación en clase.

\item {} 
\sphinxAtStartPar
Explicación clara de conceptos.

\item {} 
\sphinxAtStartPar
Argumentación de ideas.

\item {} 
\sphinxAtStartPar
Cumplimiento de las reglas del aula.

\end{itemize}

\end{sphinxuseclass}
\end{sphinxuseclass}
\begin{sphinxuseclass}{sd-tab-item}\subsubsection*{Recursos (0)}

\begin{sphinxuseclass}{sd-tab-content}\begin{itemize}
\item {} 
\sphinxAtStartPar
Esta guía no contiene recursos.

\end{itemize}

\end{sphinxuseclass}
\end{sphinxuseclass}
\end{sphinxuseclass}

\subsection{Mensaje del autor}
\label{\detokenize{1_guias/8logica/G821:mensaje-del-autor}}\begin{quote}

\sphinxAtStartPar
“En esta guía, el estudiante comprenderá los conceptos de hardware y software. Asimismo, argumentará la importancia del uso de la CPU y la memoria RAM en el procesamiento de la información. Además, analizará el concepto de sistema computacional, el cual forma parte de la definición de informática. Finalmente, integrará los conceptos de razonamiento y lógica para analizar la definición de \sphinxstylestrong{lógica computacional}.”

\begin{flushright}
---Camilo Fonseca
\end{flushright}
\end{quote}


\bigskip\hrule\bigskip



\subsection{A00: Introducción (+0,1)}
\label{\detokenize{1_guias/8logica/G821:a00-introduccion-0-1}}
\sphinxAtStartPar
\DUrole{sd-sphinx-override}{\DUrole{sd-badge}{\DUrole{sd-bg-light}{\DUrole{sd-bg-text-light}{Exploración}}}}
\begin{enumerate}
\sphinxsetlistlabels{\arabic}{enumi}{enumii}{}{.}%
\item {} 
\sphinxAtStartPar
Transcriba en su cuaderno el objetivo, las orientaciones curriculares y los criterios de evaluación de la presente guía.

\item {} 
\sphinxAtStartPar
Lea el \sphinxstylestrong{mensaje del autor} y en su cuaderno responda:
\begin{itemize}
\item {} 
\sphinxAtStartPar
¿Qué tienen en común un smartphone, un smartTV y un computador?

\end{itemize}

\end{enumerate}


\bigskip\hrule\bigskip



\subsection{A01: Concepto de lógica computacional (+0,1)}
\label{\detokenize{1_guias/8logica/G821:a01-concepto-de-logica-computacional-0-1}}
\sphinxAtStartPar
\DUrole{sd-sphinx-override}{\DUrole{sd-badge}{\DUrole{sd-bg-info}{\DUrole{sd-bg-text-info}{Estructuración}}}}
\begin{enumerate}
\sphinxsetlistlabels{\arabic}{enumi}{enumii}{}{.}%
\item {} 
\sphinxAtStartPar
Lea los siguientes conceptos:

\end{enumerate}


\bigskip\hrule\bigskip



\subsubsection{¿Qué es el HARDWARE?}
\label{\detokenize{1_guias/8logica/G821:que-es-el-hardware}}
\sphinxAtStartPar
%
\begin{footnote}[1]\sphinxAtStartFootnote
Deitel, P. J., \& Deitel, H. (2014). C++: cómo programar (A. V. Romero Elizondo, Trans.). Pearson Educación.
%
\end{footnote} El \sphinxstylestrong{HARDWARE}, es el conjunto de componentes de un \sphinxstylestrong{sistema informático} (por ejemplo un computador) los cuales se pueden ver o tocar.

\sphinxAtStartPar
Estos componentes se clasifican en componentes de entrada, salida, procesamiento, memoria y almacenamiento.

\sphinxAtStartPar
A continuación se muestran algunos ejemplos:
\begin{itemize}
\item {} 
\sphinxAtStartPar
\sphinxstylestrong{Dispositivos de entrada:} Teclados, mouse, micrófonos

\item {} 
\sphinxAtStartPar
\sphinxstylestrong{Dispositivos de salida:} Monitores, impresoras, bafles

\item {} 
\sphinxAtStartPar
\sphinxstylestrong{Procesamiento:} CPU

\item {} 
\sphinxAtStartPar
\sphinxstylestrong{Memoria:} memoria RAM, memoria CACHE

\item {} 
\sphinxAtStartPar
\sphinxstylestrong{Almacenamiento:} Discos duros, memorias USB

\end{itemize}


\paragraph{CPU (Unidad de procesamiento central)}
\label{\detokenize{1_guias/8logica/G821:cpu-unidad-de-procesamiento-central}}

\begin{savenotes}\sphinxattablestart
\sphinxthistablewithglobalstyle
\centering
\begin{tabulary}{\linewidth}[t]{T}
\sphinxtoprule
\sphinxtableatstartofbodyhook

\begin{sphinxfigure-in-table}
\centering
\capstart
\noindent\sphinxincludegraphics[width=300\sphinxpxdimen]{{image0114}.png}
\sphinxfigcaption{CPU}\label{\detokenize{1_guias/8logica/G821:id13}}\end{sphinxfigure-in-table}\relax
\\
\sphinxbottomrule
\end{tabulary}
\sphinxtableafterendhook\par
\sphinxattableend\end{savenotes}
\begin{quote}

\sphinxAtStartPar
\sphinxstyleemphasis{“Muchas personas confunden la CPU con una torre de PC. Realmente son cosas diferentes”}.
\end{quote}

\sphinxAtStartPar
\sphinxfootnotemark[1] La CPU es considerada el “cerebro” del computador. Es el elemento principal que se encarga de:
\begin{itemize}
\item {} 
\sphinxAtStartPar
Procesar todas las instrucciones y cálculos del sistema

\item {} 
\sphinxAtStartPar
Coordinar el funcionamiento de los demás componentes

\item {} 
\sphinxAtStartPar
Ejecutar los programas y aplicaciones

\end{itemize}

\sphinxAtStartPar
La CPU está compuesta principalmente por:
\begin{itemize}
\item {} 
\sphinxAtStartPar
\sphinxstylestrong{Unidad de control}: Coordina y dirige todas las operaciones.

\item {} 
\sphinxAtStartPar
\sphinxstylestrong{ALU}: Unidad Lógica Aritmética. Realiza cálculos matemáticos y operaciones lógicas

\item {} 
\sphinxAtStartPar
\sphinxstylestrong{Registros}: Son unidades pequeñas de memoria de alta velocidad

\end{itemize}

\sphinxAtStartPar
La velocidad de la CPU se mide en Hertz (Hz) y determina qué tan rápido puede procesar las instrucciones. A mayor velocidad, mejor rendimiento del sistema.


\paragraph{MEMORIA RAM}
\label{\detokenize{1_guias/8logica/G821:memoria-ram}}

\begin{savenotes}\sphinxattablestart
\sphinxthistablewithglobalstyle
\centering
\begin{tabulary}{\linewidth}[t]{T}
\sphinxtoprule
\sphinxtableatstartofbodyhook

\begin{sphinxfigure-in-table}
\centering
\capstart
\noindent\sphinxincludegraphics[width=300\sphinxpxdimen]{{image0211}.png}
\sphinxfigcaption{Memoria RAM}\label{\detokenize{1_guias/8logica/G821:id14}}\end{sphinxfigure-in-table}\relax
\\
\sphinxbottomrule
\end{tabulary}
\sphinxtableafterendhook\par
\sphinxattableend\end{savenotes}

\sphinxAtStartPar
\sphinxfootnotemark[1] La memoria RAM (Memoria de acceso aleatorio) tiene las siguientes funciones:
\begin{itemize}
\item {} 
\sphinxAtStartPar
ALMACENA temporalmente los datos y programas que el computador está usando activamente.

\item {} 
\sphinxAtStartPar
ACCEDE rápidamente a la información (SI y SOLO SI el equipo está encendido).

\end{itemize}

\sphinxAtStartPar
Las características de la memoria RAM son:
\begin{itemize}
\item {} 
\sphinxAtStartPar
Se borra cuando se apaga el computador (memoria volátil).

\item {} 
\sphinxAtStartPar
A mayor cantidad de RAM, más programas pueden ejecutarse simultáneamente.

\item {} 
\sphinxAtStartPar
Su velocidad afecta directamente el rendimiento del sistema.

\item {} 
\sphinxAtStartPar
Es diferente al almacenamiento permanente (disco duro) porque la memoria RAM es temporal.

\end{itemize}


\bigskip\hrule\bigskip



\subsubsection{¿Qué es el SOFTWARE?}
\label{\detokenize{1_guias/8logica/G821:que-es-el-software}}

\begin{savenotes}\sphinxattablestart
\sphinxthistablewithglobalstyle
\centering
\begin{tabulary}{\linewidth}[t]{T}
\sphinxtoprule
\sphinxtableatstartofbodyhook

\begin{sphinxfigure-in-table}
\centering
\capstart
\noindent\sphinxincludegraphics[width=300\sphinxpxdimen]{{image0310}.png}
\sphinxfigcaption{Ejemplo de interfaz gráfica}\label{\detokenize{1_guias/8logica/G821:id15}}\end{sphinxfigure-in-table}\relax
\\
\sphinxbottomrule
\end{tabulary}
\sphinxtableafterendhook\par
\sphinxattableend\end{savenotes}

\sphinxAtStartPar
\sphinxfootnotemark[1] El SOFTWARE, es el conjunto de: \sphinxstylestrong{programas}, \sphinxstylestrong{bases de datos}, diferentes tipos de \sphinxstylestrong{archivos} y \sphinxstylestrong{documentación} que le dicen al HARDWARE qué hacer. A diferencia del hardware que es físico y tangible, el software no se puede tocar.

\sphinxAtStartPar
El software solo puede verse su funcionamiento en la pantalla. El software es la parte lógica e intangible de un sistema informático. El software se clasifica de la siguiente manera:
\begin{itemize}
\item {} 
\sphinxAtStartPar
\sphinxstylestrong{Sistema:} Software que controla el funcionamiento básico del computador como por ejemplo: Firmware BIOS, Firmware UEFI, OS Windows, OS Linux, MacOS, Controladores o Drivers.

\item {} 
\sphinxAtStartPar
\sphinxstylestrong{Aplicaciones:} Diferente software que funciona sobre un sistema operativo como por ejemplo: Procesadores de texto, navegadores web, juegos, editores de imagen, modelado en 3D. La mayoría suelen tener interfáz gráfica.

\end{itemize}

\sphinxAtStartPar
Una \sphinxstylestrong{interfaz gráfica} es un sistema que permite a los usuarios interactuar con dispositivos electrónicos a través de elementos visuales en la pantalla, como iconos, botones y menús. Esto facilita la realización de acciones sin necesidad de conocimientos de programación, haciendo la experiencia más intuitiva y accesible.


\bigskip\hrule\bigskip



\subsubsection{¿Qué es una SISTEMA COMPUTACIONAL?}
\label{\detokenize{1_guias/8logica/G821:que-es-una-sistema-computacional}}
\sphinxAtStartPar
%
\begin{footnote}[2]\sphinxAtStartFootnote
Universidad Nacional del Nordeste. (n.d.). Informática. Conceptos fundamentales. https://exa.unne.edu.ar/ingenieria/computacion/Tema1.pdf
%
\end{footnote} Un \sphinxstylestrong{sistema computacional} es el conjunto de hardware que interactúa con un software para almacenar, procesar y mostrar información.


\begin{savenotes}\sphinxattablestart
\sphinxthistablewithglobalstyle
\centering
\begin{tabulary}{\linewidth}[t]{T}
\sphinxtoprule
\sphinxtableatstartofbodyhook

\begin{sphinxfigure-in-table}
\centering
\capstart
\noindent\sphinxincludegraphics[width=300\sphinxpxdimen]{{image048}.png}
\sphinxfigcaption{Computadora}\label{\detokenize{1_guias/8logica/G821:id16}}\end{sphinxfigure-in-table}\relax
\\
\sphinxbottomrule
\end{tabulary}
\sphinxtableafterendhook\par
\sphinxattableend\end{savenotes}


\bigskip\hrule\bigskip



\subsubsection{¿Qué es la INFORMÁTICA?}
\label{\detokenize{1_guias/8logica/G821:que-es-la-informatica}}
\sphinxAtStartPar
\sphinxfootnotemark[2] \sphinxstylestrong{La informática} es la ciencia o disciplina que estudia el diseño, construcción e implementación de \sphinxstylestrong{sistemas computacionales} para garantizar la \sphinxstylestrong{integridad} y \sphinxstylestrong{disponibilidad} de la información.
\begin{itemize}
\item {} 
\sphinxAtStartPar
REGLA 01: (Integridad de la información)
\begin{quote}

\sphinxAtStartPar
\sphinxstyleemphasis{“Toda información debe conservar su integridad para que no pueda ser modifica sin autorización”.}
\end{quote}

\item {} 
\sphinxAtStartPar
REGLA 02: (Disponibilidad de la información)
\begin{quote}

\sphinxAtStartPar
\sphinxstyleemphasis{“Toda información debe estar disponible y llegar oportunamente”.}
\end{quote}

\end{itemize}


\bigskip\hrule\bigskip



\subsubsection{¿Qué es RAZONAR?}
\label{\detokenize{1_guias/8logica/G821:que-es-razonar}}
\sphinxAtStartPar
%
\begin{footnote}[3]\sphinxAtStartFootnote
Trejos Buriticá, O. I. (2017). Lógica de programación. Ediciones de la U.
%
\end{footnote} \sphinxstylestrong{Razonar} es la \sphinxstyleemphasis{actividad mental} mediante la cual se relacionan ideas o premisas para inferir una conclusión de manera coherente o justificada.

\sphinxAtStartPar
Por ejemplo:
\begin{itemize}
\item {} 
\sphinxAtStartPar
Premisa 01: Todos los humanos son mortales

\item {} 
\sphinxAtStartPar
Premisa 02: Sócrates es humano

\item {} 
\sphinxAtStartPar
\sphinxstylestrong{Conclusión:} Sócrates es mortal

\end{itemize}


\bigskip\hrule\bigskip



\subsubsection{¿Qué es la LÓGICA?}
\label{\detokenize{1_guias/8logica/G821:que-es-la-logica}}
\sphinxAtStartPar
\sphinxfootnotemark[3] La \sphinxstylestrong{lógica} es la \sphinxstyleemphasis{disciplina} que permite analizar y verificar la validez de los razonamientos, determinando si una conclusión se deriva correctamente de unas premisas.

\sphinxAtStartPar
Por ejemplo:

\sphinxAtStartPar
Para el siguiente razonamiento, verifique si es correcto o no:
\begin{itemize}
\item {} 
\sphinxAtStartPar
Premisa 01: Todos los humanos son mortales

\item {} 
\sphinxAtStartPar
Premisa 02: Sócrates es humano

\item {} 
\sphinxAtStartPar
\sphinxstylestrong{Conclusión:} Sócrates es mortal

\end{itemize}
\begin{quote}

\sphinxAtStartPar
\sphinxstylestrong{ANÁLISIS:} Si Sócrates es un ser humano, entonces Sócrates es mortal porque todos los humanos son mortales. El razonamiento es correcto.
\end{quote}


\bigskip\hrule\bigskip



\subsubsection{¿Qué es lógica computacional?}
\label{\detokenize{1_guias/8logica/G821:que-es-logica-computacional}}
\sphinxAtStartPar
\sphinxfootnotemark[3] La \sphinxstylestrong{lógica computacional} es el área de la \sphinxstylestrong{informática} que aplica los principios de la \sphinxstylestrong{lógica} para representar y automatizar el \sphinxstylestrong{razonamiento}, con el objetivo de resolver problemas de forma estructurada y eficiente.


\bigskip\hrule\bigskip



\subsection{A02: Taller (+0,8)}
\label{\detokenize{1_guias/8logica/G821:a02-taller-0-8}}
\sphinxAtStartPar
\DUrole{sd-sphinx-override}{\DUrole{sd-badge}{\DUrole{sd-bg-warning}{\DUrole{sd-bg-text-warning}{Actividad práctica}}}}
\begin{enumerate}
\sphinxsetlistlabels{\arabic}{enumi}{enumii}{}{.}%
\item {} 
\sphinxAtStartPar
En su cuaderno responda las siguientes preguntas:
\begin{itemize}
\item {} 
\sphinxAtStartPar
¿Qué es el hardware?

\item {} 
\sphinxAtStartPar
¿De que se compone el hardware de un sistema informático?

\item {} 
\sphinxAtStartPar
¿El concepto de hardware solo puede aplicarse a computadores? ¿Por qué?

\end{itemize}

\item {} 
\sphinxAtStartPar
En su cuaderno cree un \sphinxstylestrong{mapa conceptual} que muestre los diferentes componentes del hardware y dos ejemplos por cada componente.

\item {} 
\sphinxAtStartPar
En su cuaderno responda las siguientes preguntas:
\begin{itemize}
\item {} 
\sphinxAtStartPar
¿Qué es la CPU?

\item {} 
\sphinxAtStartPar
¿Por qué la CPU es considerada el cerebro del computador?

\item {} 
\sphinxAtStartPar
¿En que se mide la velocidad de las CPUs?

\item {} 
\sphinxAtStartPar
¿Qué significa que una velocidad en la CPU sea alta o baja?

\end{itemize}

\item {} 
\sphinxAtStartPar
En su cuaderno cree un \sphinxstylestrong{mapa conceptual} de los tres principales componentes de la CPU y descríbalos.

\item {} 
\sphinxAtStartPar
En INTERNET, consulte 3 marcas de fabricantes de CPU y descríbalas en su cuaderno.

\item {} 
\sphinxAtStartPar
En su cuaderno responda las siguientes preguntas:
\begin{itemize}
\item {} 
\sphinxAtStartPar
¿Qué es la memoria RAM?

\item {} 
\sphinxAtStartPar
¿De que se encarga la memoria RAM?

\item {} 
\sphinxAtStartPar
¿Por qué la memoria RAM es tan importante al momento de procesar la información?

\item {} 
\sphinxAtStartPar
¿Cuales son las características de la memoria RAM?

\end{itemize}

\item {} 
\sphinxAtStartPar
En Internet, consulte las diferentes generaciones de memoria RAM y con esta información en su cuaderno CREE un mapa conceptual.

\item {} 
\sphinxAtStartPar
Con los temas vistos hasta ahora, en \sphinxstylestrong{UNA PÁGINA DE CUADERNO}, con sus palabras (no usar IA) responda la siguiente pregunta:
\begin{itemize}
\item {} 
\sphinxAtStartPar
¿Por qué la memoria RAM y la CPU son tan importantes al momento de procesar información?

\end{itemize}

\item {} 
\sphinxAtStartPar
En su cuaderno responda las siguientes preguntas:
\begin{itemize}
\item {} 
\sphinxAtStartPar
¿Qué es el software?

\item {} 
\sphinxAtStartPar
¿En qué se diferencian el HARDWARE y el SOFTWARE?

\item {} 
\sphinxAtStartPar
¿Por qué los controladores de HARDWARE no necesitan interfaz gráfica?

\end{itemize}

\item {} 
\sphinxAtStartPar
En su cuaderno, cree un mapa conceptual que presente la \sphinxstylestrong{clasificación del software}.

\item {} 
\sphinxAtStartPar
Consulte en Internet y luego en su cuaderno realice una descripción de lo siguiente:
\begin{itemize}
\item {} 
\sphinxAtStartPar
10 sistemas operativos de la empresa MICROSOFT.

\item {} 
\sphinxAtStartPar
1 sistema operativo de la empresa GOOGLE.

\item {} 
\sphinxAtStartPar
3 sistemas operativos de la empresa APPLE.

\end{itemize}

\item {} 
\sphinxAtStartPar
En Internet, consulte dos ejemplos para cada una de las siguientes aplicaciones, luego descríbalas en su cuaderno:
\begin{itemize}
\item {} 
\sphinxAtStartPar
Procesadores de texto

\item {} 
\sphinxAtStartPar
Navegadores web

\item {} 
\sphinxAtStartPar
Videojuegos

\item {} 
\sphinxAtStartPar
Editores de imágenes

\item {} 
\sphinxAtStartPar
Modelado 3D

\end{itemize}

\item {} 
\sphinxAtStartPar
En su cuaderno responda las siguientes preguntas:
\begin{itemize}
\item {} 
\sphinxAtStartPar
¿Qué es un sistema computacional?

\item {} 
\sphinxAtStartPar
¿De qué manera un hardware interactúa con un software? (piense en un smartphone, un smart\sphinxhyphen{}TV, etc.)

\item {} 
\sphinxAtStartPar
¿Cómo un sistema computacional interactúa con la información? (piense en un software como Microsoft EXCEL o WHATSAPP)

\end{itemize}

\item {} 
\sphinxAtStartPar
En su cuaderno, con sus palabras escriba 5 ejemplos de sistemas computacionales (no usar IA). Tenga en cuenta el siguiente ejemplo:
\begin{quote}

\sphinxAtStartPar
\sphinxstylestrong{EJEMPLO 01}: Un \sphinxstyleemphasis{portátil} es un ejemplo de \sphinxstylestrong{sistema computacional} porque su \sphinxstyleemphasis{hardware} (teclado, pantalla, touchpad) y \sphinxstyleemphasis{software} (OS, aplicaciones de escritorio) permite registrar, almacenar y mostrar información para un usuario”.
\end{quote}

\item {} 
\sphinxAtStartPar
responda las siguientes preguntas en su cuaderno:
\begin{itemize}
\item {} 
\sphinxAtStartPar
¿En que se relacionan la informática y la información?

\item {} 
\sphinxAtStartPar
Aparte de la \sphinxstylestrong{integridad} y \sphinxstylestrong{disponibilidad}, ¿Qué otras reglas pueden ser aplicadas en la información?

\end{itemize}

\item {} 
\sphinxAtStartPar
En su cuaderno, escriba los siguientes casos y por cada uno indique si aplica la regla 01 (\sphinxstyleemphasis{Integridad de la información}) o aplica la regla 02 (\sphinxstyleemphasis{disponibilidad de la información}) (escribir el por qué)
\begin{itemize}
\item {} 
\sphinxAtStartPar
\sphinxstylestrong{CASO 01}: \sphinxstyleemphasis{“De nada sirve que un banco entregue el valor de un saldo 10 meses después de ser solicitado. Es por eso que se usan las computadoras, porque permiten que la información este disponible y llegue oportunamente (milésimas de segundo)”}.

\item {} 
\sphinxAtStartPar
\sphinxstylestrong{CASO 02:} \sphinxstyleemphasis{“De nada sirve que una persona vaya al banco a consultar su saldo si este no corresponde al saldo real. Es por eso que se usan las computadoras, porque permiten que la información conserve su estado sin los riesgos que representa tenerla en otros medios como el papel o en la memoria de las personas”}.

\end{itemize}


\begin{savenotes}\sphinxattablestart
\sphinxthistablewithglobalstyle
\centering
\begin{tabulary}{\linewidth}[t]{T}
\sphinxtoprule
\sphinxtableatstartofbodyhook

\begin{sphinxfigure-in-table}
\centering
\capstart
\noindent\sphinxincludegraphics[width=300\sphinxpxdimen]{{image058}.png}
\sphinxfigcaption{\sphinxstylestrong{BANCOS:} Los bancos deben garantizar integridad y disponibilidad}\label{\detokenize{1_guias/8logica/G821:id17}}\end{sphinxfigure-in-table}\relax
\\
\sphinxbottomrule
\end{tabulary}
\sphinxtableafterendhook\par
\sphinxattableend\end{savenotes}

\item {} 
\sphinxAtStartPar
En su cuaderno responda las siguientes preguntas:
\begin{itemize}
\item {} 
\sphinxAtStartPar
¿Qué es razonar?

\item {} 
\sphinxAtStartPar
¿Qué es lógica?

\item {} 
\sphinxAtStartPar
¿Cómo se relacionan los conceptos de razonar y lógica?

\item {} 
\sphinxAtStartPar
¿Se podrían comprobar resultados de razonamiento sin aplicar lógica? (explique su respuesta).

\end{itemize}

\item {} 
\sphinxAtStartPar
Razone la siguiente situación y en su cuaderno DEDUZCA que pudo haber pasado (3 conclusiones con temáticas diferentes):
\begin{itemize}
\item {} 
\sphinxAtStartPar
Premisa 01: Una persona escuchó ruidos extraños en su casa.

\item {} 
\sphinxAtStartPar
Premisa 02: Los ruidos fueron en horas de la noche.

\item {} 
\sphinxAtStartPar
Premisa 03: Habían varios residuos de comida en el piso.

\item {} 
\sphinxAtStartPar
Premisa 04: El computador estaba desconectado.

\item {} 
\sphinxAtStartPar
Premisa 05: No hay ningún objeto perdido.

\end{itemize}
\begin{itemize}
\item {} 
\sphinxAtStartPar
\sphinxstylestrong{Conclusión 01:} ?

\item {} 
\sphinxAtStartPar
\sphinxstylestrong{Conclusión 02:} ?

\item {} 
\sphinxAtStartPar
\sphinxstylestrong{Conclusión 03:} ?

\end{itemize}

\item {} 
\sphinxAtStartPar
En su cuaderno responda las siguientes preguntas:
\begin{itemize}
\item {} 
\sphinxAtStartPar
¿Qué es un algoritmo?

\item {} 
\sphinxAtStartPar
¿Qué es lógica computacional?

\item {} 
\sphinxAtStartPar
¿Cómo se relaciona la lógica computacional y la informática?

\item {} 
\sphinxAtStartPar
¿Cómo se relacionan los algoritmos y la lógica computacional?

\end{itemize}

\item {} 
\sphinxAtStartPar
En su cuaderno escriba las siguientes frases e indique si es FALSO, VERDADERO y explique el por qué:
\begin{itemize}
\item {} 
\sphinxAtStartPar
“\sphinxstyleemphasis{La \sphinxstylestrong{lógica computacional} establece las reglas para programar una persona y modificar su comportamiento}”. \_\_

\item {} 
\sphinxAtStartPar
“\sphinxstyleemphasis{Optimizar \sphinxstylestrong{algoritmos} permite solucionar problemas más rápido pero con los mismos recursos}”. \_\_

\end{itemize}

\end{enumerate}


\bigskip\hrule\bigskip



\subsection{A03: Examen escrito (+4,0)}
\label{\detokenize{1_guias/8logica/G821:a03-examen-escrito-4-0}}
\sphinxAtStartPar
\DUrole{sd-sphinx-override}{\DUrole{sd-badge}{\DUrole{sd-bg-danger}{\DUrole{sd-bg-text-danger}{Transferencia}}}}


\begin{savenotes}\sphinxattablestart
\sphinxthistablewithglobalstyle
\centering
\begin{tabulary}{\linewidth}[t]{TT}
\sphinxtoprule
\sphinxtableatstartofbodyhook
\sphinxAtStartPar
\sphinxstylestrong{TIEMPO}
&
\sphinxAtStartPar
29 minutos
\\
\sphinxhline
\sphinxAtStartPar
\sphinxstylestrong{PREGUNTAS}
&
\sphinxAtStartPar
10
\\
\sphinxhline
\sphinxAtStartPar
\sphinxstylestrong{OBJETIVO}
&
\sphinxAtStartPar
Comprobar conocimientos adquiridos.
\\
\sphinxhline
\sphinxAtStartPar
\sphinxstylestrong{ACTIVIDAD}
&
\sphinxAtStartPar
Responda las siguientes preguntas. Cuadernos y celulares guardados. Solo se permite hoja, lapiz y borrador.
\\
\sphinxbottomrule
\end{tabulary}
\sphinxtableafterendhook\par
\sphinxattableend\end{savenotes}

\begin{sphinxadmonition}{note}{Posibles preguntas}
\begin{enumerate}
\sphinxsetlistlabels{\arabic}{enumi}{enumii}{}{.}%
\item {} 
\sphinxAtStartPar
Explique con sus propias palabras qué es el \sphinxstylestrong{hardware} y mencione al menos tres categorías en las que se clasifican sus componentes, dando un ejemplo de cada una.

\item {} 
\sphinxAtStartPar
¿Por qué la \sphinxstylestrong{CPU} es considerada el “cerebro” del computador? Describa sus tres componentes principales y su función.

\item {} 
\sphinxAtStartPar
Escriba las funciones principales de la \sphinxstylestrong{memoria RAM} y explique por qué es fundamental en el procesamiento de información.

\item {} 
\sphinxAtStartPar
Compare y contraste el \sphinxstylestrong{hardware} y el \sphinxstylestrong{software}, resaltando sus diferencias fundamentales.

\item {} 
\sphinxAtStartPar
¿Qué es el \sphinxstylestrong{software} y cómo se clasifica? Describa brevemente cada una de sus categorías principales.

\item {} 
\sphinxAtStartPar
Describa el concepto de \sphinxstylestrong{“interfaz gráfica”} y explique su importancia para los usuarios.

\item {} 
\sphinxAtStartPar
Defina qué es un \sphinxstylestrong{“sistema computacional”} y cómo interactúan sus componentes para procesar información.

\item {} 
\sphinxAtStartPar
¿Qué es la \sphinxstylestrong{informática} y cuáles son sus objetivos principales en relación con los sistemas computacionales?

\item {} 
\sphinxAtStartPar
Explique las dos \sphinxstylestrong{reglas fundamentales de la informática} (Integridad y Disponibilidad de la información) y por qué son importantes.

\item {} 
\sphinxAtStartPar
Describa el concepto de \sphinxstylestrong{“razonar”} y proporcione un ejemplo propio (diferente al ejemplo de sócrates), que incluya premisas y una conclusión.

\item {} 
\sphinxAtStartPar
¿Qué es la \sphinxstylestrong{lógica} y cómo se relaciona con el razonamiento? Explique su papel en la verificación de la validez de los argumentos.

\item {} 
\sphinxAtStartPar
¿Se podría comprobar la \sphinxstylestrong{validez de un razonamiento} sin aplicar principios de lógica? Justifica tu respuesta.

\item {} 
\sphinxAtStartPar
Defina qué es un \sphinxstylestrong{“algoritmo”} y cuáles son sus características principales.

\item {} 
\sphinxAtStartPar
¿Qué es la \sphinxstylestrong{“lógica computacional”} y cómo aplica los principios de la lógica y la informática?

\item {} 
\sphinxAtStartPar
¿Cómo se relacionan los \sphinxstylestrong{algoritmos} y la \sphinxstylestrong{lógica computacional} en la resolución de problemas?

\item {} 
\sphinxAtStartPar
Basándose en la información sobre \sphinxstylestrong{hardware} y \sphinxstylestrong{software}, ¿cómo cree que un smartphone o una tablet funcionan como sistemas computacionales?

\item {} 
\sphinxAtStartPar
¿Por qué es importante que la memoria \sphinxstylestrong{RAM} y la \sphinxstylestrong{CPU} trabajen de manera coordinada para un buen rendimiento del sistema?

\item {} 
\sphinxAtStartPar
Escriba un ejemplo de su vida diaria donde se aplique el concepto de \sphinxstylestrong{“disponibilidad de la información”}.

\item {} 
\sphinxAtStartPar
Explique un escenario donde la \sphinxstylestrong{“integridad de la información”} sea crucial y qué podría suceder si esta se viera comprometida.

\item {} 
\sphinxAtStartPar
¿Cómo la evolución del \sphinxstylestrong{hardware} y el \sphinxstylestrong{software} ha impactado la forma en que utilizamos los \sphinxstylestrong{sistemas computacionales} hoy en día?

\end{enumerate}
\end{sphinxadmonition}


\bigskip\hrule\bigskip


\sphinxstepscope


\section{Guía 822: Jerarquía de datos (P2)}
\label{\detokenize{1_guias/8logica/G822:guia-822-jerarquia-de-datos-p2}}\label{\detokenize{1_guias/8logica/G822::doc}}

\begin{savenotes}\sphinxattablestart
\sphinxthistablewithglobalstyle
\centering
\begin{tabular}[t]{\X{35}{100}\X{65}{100}}
\sphinxtoprule
\sphinxtableatstartofbodyhook

\begin{savenotes}\sphinxattablestart
\sphinxthistablewithglobalstyle
\centering
\begin{tabulary}{\linewidth}[t]{T}
\sphinxtoprule
\sphinxtableatstartofbodyhook
\sphinxAtStartPar
\sphinxincludegraphics{{escudo}.png}
\\
\sphinxbottomrule
\end{tabulary}
\sphinxtableafterendhook\par
\sphinxattableend\end{savenotes}
&

\begin{savenotes}\sphinxattablestart
\sphinxthistablewithglobalstyle
\raggedright
\begin{tabular}[t]{*{1}{\X{1}{1}}}
\sphinxtoprule
\sphinxtableatstartofbodyhook
\sphinxAtStartPar
\sphinxstylestrong{INSTITUCIÓN EDUCATIVA}: John F. Kennedy
\begin{itemize}
\item {} 
\sphinxAtStartPar
\sphinxstylestrong{DOCENTE}: Julian Camilo Fonseca Romero

\item {} 
\sphinxAtStartPar
\sphinxstylestrong{ASIGNATURA}: Lógica computacional

\item {} 
\sphinxAtStartPar
\sphinxstylestrong{GRADO}: OCTAVO

\item {} 
\sphinxAtStartPar
\sphinxstylestrong{PERIODO}: 02

\item {} 
\sphinxAtStartPar
\sphinxstylestrong{TIEMPO}: 10 semanas (1 hora / semana)

\item {} 
\sphinxAtStartPar
\sphinxstylestrong{TIEMPO DE TRABAJO}: 10/10 horas

\end{itemize}
\\
\sphinxbottomrule
\end{tabular}
\sphinxtableafterendhook\par
\sphinxattableend\end{savenotes}
\\
\sphinxbottomrule
\end{tabular}
\sphinxtableafterendhook\par
\sphinxattableend\end{savenotes}

\begin{sphinxuseclass}{sd-tab-set}
\begin{sphinxuseclass}{sd-tab-item}\subsubsection*{Objetivo}

\begin{sphinxuseclass}{sd-tab-content}
\sphinxAtStartPar
\DUrole{sd-sphinx-override}{\DUrole{sd-badge}{\DUrole{sd-bg-light}{\DUrole{sd-bg-text-light}{Exploración}}}}
\DUrole{sd-sphinx-override}{\DUrole{sd-badge}{\DUrole{sd-bg-info}{\DUrole{sd-bg-text-info}{Estructuración}}}}
\DUrole{sd-sphinx-override}{\DUrole{sd-badge}{\DUrole{sd-bg-warning}{\DUrole{sd-bg-text-warning}{Actividad práctica}}}}
\DUrole{sd-sphinx-override}{\DUrole{sd-badge}{\DUrole{sd-bg-danger}{\DUrole{sd-bg-text-danger}{Transferencia}}}}
\begin{itemize}
\item {} 
\sphinxAtStartPar
\sphinxstylestrong{Temática}: Jerarquía de datos.

\item {} 
\sphinxAtStartPar
\sphinxstylestrong{Objetivo de aprendizaje}: Comprender la diferencia entre \sphinxstylestrong{datos}, \sphinxstylestrong{información} e identificar los elementos de la \sphinxstylestrong{Jerarquía de datos}.

\end{itemize}

\end{sphinxuseclass}
\end{sphinxuseclass}
\begin{sphinxuseclass}{sd-tab-item}\subsubsection*{Orientaciones curriculares}

\begin{sphinxuseclass}{sd-tab-content}\begin{itemize}
\item {} 
\sphinxAtStartPar
\sphinxstylestrong{Componente}: Naturaleza y Evolución de la T\&I.

\item {} 
\sphinxAtStartPar
\sphinxstylestrong{Competencia}: Relaciono saberes, conocimientos tecnológicos e informáticos con los conocimientos de otras disciplinas.

\item {} 
\sphinxAtStartPar
\sphinxstylestrong{Evidencia de aprendizaje}: Comprendo los principios y conocimientos tecnológicos e informáticos que hacen posible el funcionamiento de productos tecnológicos actuales.

\end{itemize}

\end{sphinxuseclass}
\end{sphinxuseclass}
\begin{sphinxuseclass}{sd-tab-item}\subsubsection*{Criterios de evaluación}

\begin{sphinxuseclass}{sd-tab-content}\begin{itemize}
\item {} 
\sphinxAtStartPar
Actividades desarrolladas en su totalidad.

\item {} 
\sphinxAtStartPar
Participación en clase.

\item {} 
\sphinxAtStartPar
Explicación clara de conceptos.

\item {} 
\sphinxAtStartPar
Argumentación de ideas.

\item {} 
\sphinxAtStartPar
Cumplimiento de las reglas del aula.

\end{itemize}

\end{sphinxuseclass}
\end{sphinxuseclass}
\begin{sphinxuseclass}{sd-tab-item}\subsubsection*{Recursos (0)}

\begin{sphinxuseclass}{sd-tab-content}\begin{itemize}
\item {} 
\sphinxAtStartPar
Esta guía no contiene recursos.

\end{itemize}

\end{sphinxuseclass}
\end{sphinxuseclass}
\end{sphinxuseclass}

\subsection{Mensaje del autor}
\label{\detokenize{1_guias/8logica/G822:mensaje-del-autor}}\begin{quote}

\sphinxAtStartPar
“En esta guía, el estudiante comprenderá la \sphinxstylestrong{diferencia entre datos e información}, entendiendo cómo un conjunto de registros estructurados cobra significado para guiar la toma de decisiones. Al identificar los elementos de la jerarquía de datos, el estudiante estará dando el primer paso para dominar el lenguaje que emula el pensamiento humano y da vida a la tecnología actual.”

\begin{flushright}
---Camilo Fonseca
\end{flushright}
\end{quote}


\bigskip\hrule\bigskip



\subsection{A00: Introducción (+0,1)}
\label{\detokenize{1_guias/8logica/G822:a00-introduccion-0-1}}
\sphinxAtStartPar
\DUrole{sd-sphinx-override}{\DUrole{sd-badge}{\DUrole{sd-bg-light}{\DUrole{sd-bg-text-light}{Exploración}}}}
\begin{enumerate}
\sphinxsetlistlabels{\arabic}{enumi}{enumii}{}{.}%
\item {} 
\sphinxAtStartPar
Transcriba en su cuaderno el objetivo, las orientaciones curriculares y los criterios de evaluación de la presente guía.

\item {} 
\sphinxAtStartPar
Lea el \sphinxstylestrong{mensaje del autor} y en su cuaderno con sus palabras responda:
\begin{itemize}
\item {} 
\sphinxAtStartPar
¿En que se diferencian los datos de la información?

\end{itemize}

\end{enumerate}


\bigskip\hrule\bigskip



\subsection{A01: Información y Jerarquía de datos (+0,1)}
\label{\detokenize{1_guias/8logica/G822:a01-informacion-y-jerarquia-de-datos-0-1}}
\sphinxAtStartPar
\DUrole{sd-sphinx-override}{\DUrole{sd-badge}{\DUrole{sd-bg-info}{\DUrole{sd-bg-text-info}{Estructuración}}}}
\begin{enumerate}
\sphinxsetlistlabels{\arabic}{enumi}{enumii}{}{.}%
\item {} 
\sphinxAtStartPar
Lea los siguientes conceptos:

\end{enumerate}


\bigskip\hrule\bigskip



\subsubsection{¿Qué son los datos?}
\label{\detokenize{1_guias/8logica/G822:que-son-los-datos}}
\sphinxAtStartPar
%
\begin{footnote}[1]\sphinxAtStartFootnote
Trejos Buriticá, O. I. (2017). Lógica de programación. Ediciones de la U
%
\end{footnote} Los \sphinxstylestrong{datos} son cualquier cantidad, valor o hecho \sphinxstylestrong{SIN ORGANIZAR, SIN ANALIZAR o SIN PROCESAR}.
\begin{quote}

\sphinxAtStartPar
“Los datos NO ESPECIFICAN significado alguno”.
\end{quote}

\sphinxAtStartPar
Ejemplo de datos:
\begin{itemize}
\item {} 
\sphinxAtStartPar
rojo

\item {} 
\sphinxAtStartPar
llueve

\item {} 
\sphinxAtStartPar
Alexa

\item {} 
\sphinxAtStartPar
33cm

\item {} 
\sphinxAtStartPar
5kg

\item {} 
\sphinxAtStartPar
\$1,000,000.00

\end{itemize}


\bigskip\hrule\bigskip



\subsubsection{¿Qué es la información?}
\label{\detokenize{1_guias/8logica/G822:que-es-la-informacion}}
\sphinxAtStartPar
\sphinxfootnotemark[1] La \sphinxstylestrong{información} es el conjunto de datos organizados, analizados y procesados que tienen significado para las personas.
\begin{quote}

\sphinxAtStartPar
“La información tiene contexto y significado”.
\end{quote}

\sphinxAtStartPar
Ejemplo de información:
\begin{itemize}
\item {} 
\sphinxAtStartPar
El carro del ingeniero es de color \sphinxstylestrong{rojo}.

\item {} 
\sphinxAtStartPar
En el campo \sphinxstylestrong{llueve}, mientras que en la ciudad no lo está.

\item {} 
\sphinxAtStartPar
El nombre de la IA de AMAZON es \sphinxstylestrong{Alexa}.

\item {} 
\sphinxAtStartPar
La hoja tiene una longitud de \sphinxstylestrong{33cm}.

\item {} 
\sphinxAtStartPar
Para le receta se necesitan \sphinxstylestrong{5kg} de carne de res.

\item {} 
\sphinxAtStartPar
Contando los gastos, al usuario le quedó un saldo de \sphinxstylestrong{\$1,000,000.00} COP.

\end{itemize}


\bigskip\hrule\bigskip



\subsubsection{¿Qué es la Jerarquía de datos?}
\label{\detokenize{1_guias/8logica/G822:que-es-la-jerarquia-de-datos}}
\sphinxAtStartPar
%
\begin{footnote}[2]\sphinxAtStartFootnote
Deitel, P. J., \& Deitel, H. (2014). C++: cómo programar (A. V. Romero Elizondo, Trans.). Pearson Educación.
%
\end{footnote} La \sphinxstylestrong{Jerarquía de datos} es una forma de clasificar la información. La jerarquía de datos se compone de los siguientes elementos organizados así:
\begin{equation*}
\begin{split}
Base\_de\_datos(Tabla(Registro(Dato(Carácter(Byte(bit))))))
\end{split}
\end{equation*}
\sphinxAtStartPar
A continuación se muestra la definición de cada uno de los elementos (ordenados de MENOR a MAYOR) que componen la jerarquía de datos:


\bigskip\hrule\bigskip



\paragraph{bit}
\label{\detokenize{1_guias/8logica/G822:bit}}
\sphinxAtStartPar
Un bit es la unidad básica de información en informática. Es como un interruptor que puede estar en dos estados: \sphinxstylestrong{encendido} representado por \sphinxcode{\sphinxupquote{1}} o \sphinxstylestrong{apagado} representado por \sphinxcode{\sphinxupquote{0}}.

\begin{sphinxadmonition}{note}{¿Cómo se representa un bit?}

\sphinxAtStartPar
Imagine una \sphinxstylestrong{bombilla}. Puede estar \sphinxstyleemphasis{encendida} o \sphinxstyleemphasis{apagada}.  Un bit es como esa bombilla, solo que en vez de luz, representa un valor digital.

\begin{center}\sphinxcode{\sphinxupquote{1}} o \sphinxcode{\sphinxupquote{0}}
\end{center}\end{sphinxadmonition}


\bigskip\hrule\bigskip



\paragraph{Byte (8 bits)}
\label{\detokenize{1_guias/8logica/G822:byte-8-bits}}
\sphinxAtStartPar
Un byte es un grupo de \sphinxstylestrong{8 bits} que trabajan juntos.
\begin{quote}

\sphinxAtStartPar
\sphinxstyleemphasis{“Un Byte, es como si tuviera \sphinxstylestrong{8 bombillas} que pueden estar encendidas o apagadas, y cada combinación representa un valor diferente”}.
\end{quote}

\sphinxAtStartPar
Por ejemplo, un byte (8bits) puede representar un \sphinxstylestrong{carácter} como la letra \sphinxstylestrong{“A”}, un número como \sphinxstylestrong{“255”}, o incluso un símbolo como \sphinxstylestrong{“\#”}. Por ejemplo:
\begin{itemize}
\item {} 
\sphinxAtStartPar
\sphinxcode{\sphinxupquote{10010101}} representa \sphinxkeyboard{\sphinxupquote{A}}

\item {} 
\sphinxAtStartPar
\sphinxcode{\sphinxupquote{11111111}} representa \sphinxkeyboard{\sphinxupquote{255}}

\item {} 
\sphinxAtStartPar
\sphinxcode{\sphinxupquote{10000101}} representa \sphinxkeyboard{\sphinxupquote{\#}}

\end{itemize}


\bigskip\hrule\bigskip



\paragraph{Carácter}
\label{\detokenize{1_guias/8logica/G822:caracter}}
\sphinxAtStartPar
Son representaciones de bytes.

\sphinxAtStartPar
Un carácter es como una “letra” o un “símbolo” que se ve en pantalla. Es una unidad lógica que representa un elemento individual de texto, como una letra, un número, un signo de puntuación o un espacio en blanco.

\sphinxAtStartPar
\sphinxkeyboard{\sphinxupquote{A}}, \sphinxkeyboard{\sphinxupquote{c}}, \sphinxkeyboard{\sphinxupquote{\#}}, \sphinxkeyboard{\sphinxupquote{X}}, \sphinxkeyboard{\sphinxupquote{255}}, \sphinxkeyboard{\sphinxupquote{3}}, \sphinxkeyboard{\sphinxupquote{\_}}


\bigskip\hrule\bigskip



\paragraph{Dato}
\label{\detokenize{1_guias/8logica/G822:dato}}
\sphinxAtStartPar
Los datos se componen de caracteres.

\sphinxAtStartPar
Es un atributo codificado entendible por un sistema de información. Se puede manejar y se puede comparar. Por ejemplo:

\sphinxAtStartPar
\sphinxguilabel{CAMILO}, \sphinxguilabel{ROJO}, \sphinxguilabel{33°}, \sphinxguilabel{23} , \sphinxguilabel{\$33.000,999}, \sphinxguilabel{TRUE}


\bigskip\hrule\bigskip



\paragraph{Registro}
\label{\detokenize{1_guias/8logica/G822:registro}}
\sphinxAtStartPar
Un registro es un conjunto de campos en donde cada campo almacena un dato y en donde todos los datos pertenecen a un solo ente informático. Por ejemplo:


\begin{savenotes}\sphinxattablestart
\sphinxthistablewithglobalstyle
\centering
\sphinxcapstartof{table}
\sphinxthecaptionisattop
\sphinxcaption{Registro de un estudiante}\label{\detokenize{1_guias/8logica/G822:id6}}
\sphinxaftertopcaption
\begin{tabulary}{\linewidth}[t]{TTTTT}
\sphinxtoprule
\sphinxtableatstartofbodyhook
\sphinxAtStartPar
\sphinxstylestrong{ID}
&
\sphinxAtStartPar
\sphinxstylestrong{Nombre}
&
\sphinxAtStartPar
\sphinxstylestrong{Apellido}
&
\sphinxAtStartPar
\sphinxstylestrong{Edad}
&
\sphinxAtStartPar
\sphinxstylestrong{Grado}
\\
\sphinxhline
\sphinxAtStartPar
00
&
\sphinxAtStartPar
Camilo
&
\sphinxAtStartPar
Fonseca
&
\sphinxAtStartPar
16
&
\sphinxAtStartPar
11
\\
\sphinxbottomrule
\end{tabulary}
\sphinxtableafterendhook\par
\sphinxattableend\end{savenotes}


\bigskip\hrule\bigskip



\paragraph{Tabla}
\label{\detokenize{1_guias/8logica/G822:tabla}}
\sphinxAtStartPar
Una tabla es un conjunto de registros que tienen la misma estructura y que pueden ser manejados como una sola unidad.
\begin{quote}

\sphinxAtStartPar
\sphinxstyleemphasis{“Todo REGISTRO debe tener un ID para poder identificar el registro como único”}.
\end{quote}

\sphinxAtStartPar
Por ejemplo:


\begin{savenotes}\sphinxattablestart
\sphinxthistablewithglobalstyle
\centering
\sphinxcapstartof{table}
\sphinxthecaptionisattop
\sphinxcaption{Tabla de estudiantes}\label{\detokenize{1_guias/8logica/G822:id7}}
\sphinxaftertopcaption
\begin{tabulary}{\linewidth}[t]{TTTT}
\sphinxtoprule
\sphinxtableatstartofbodyhook
\sphinxAtStartPar
\sphinxstylestrong{ID}
&
\sphinxAtStartPar
\sphinxstylestrong{Nombre}
&
\sphinxAtStartPar
\sphinxstylestrong{Apellido}
&
\sphinxAtStartPar
\sphinxstylestrong{Edad}
\\
\sphinxhline
\sphinxAtStartPar
00
&
\sphinxAtStartPar
Camilo
&
\sphinxAtStartPar
Fonseca
&
\sphinxAtStartPar
36
\\
\sphinxhline
\sphinxAtStartPar
01
&
\sphinxAtStartPar
John
&
\sphinxAtStartPar
Kennedy
&
\sphinxAtStartPar
16
\\
\sphinxhline
\sphinxAtStartPar
02
&
\sphinxAtStartPar
Daniel
&
\sphinxAtStartPar
Lopez
&
\sphinxAtStartPar
19
\\
\sphinxhline
\sphinxAtStartPar
03
&
\sphinxAtStartPar
Luis
&
\sphinxAtStartPar
Ruiz
&
\sphinxAtStartPar
23
\\
\sphinxbottomrule
\end{tabulary}
\sphinxtableafterendhook\par
\sphinxattableend\end{savenotes}


\bigskip\hrule\bigskip



\paragraph{Base de datos}
\label{\detokenize{1_guias/8logica/G822:base-de-datos}}
\sphinxAtStartPar
Es un conjunto de tablas técnicamente organizadas.
\begin{quote}

\sphinxAtStartPar
“Una \sphinxstylestrong{base de datos} esta compuesta por varias \sphinxstylestrong{tablas} relacionadas”.
\end{quote}

\sphinxAtStartPar
Por ejemplo:


\begin{savenotes}\sphinxattablestart
\sphinxthistablewithglobalstyle
\centering
\sphinxcapstartof{table}
\sphinxthecaptionisattop
\sphinxcaption{Tabla estudiantes}\label{\detokenize{1_guias/8logica/G822:id8}}
\sphinxaftertopcaption
\begin{tabular}[t]{*{1}{\X{1}{1}}}
\sphinxtoprule
\sphinxtableatstartofbodyhook

\begin{savenotes}\sphinxattablestart
\sphinxthistablewithglobalstyle
\centering
\begin{tabulary}{\linewidth}[t]{TTTT}
\sphinxtoprule
\sphinxtableatstartofbodyhook
\sphinxAtStartPar
\sphinxstylestrong{ID}
&
\sphinxAtStartPar
\sphinxstylestrong{Nombre}
&
\sphinxAtStartPar
\sphinxstylestrong{Apellido}
&
\sphinxAtStartPar
\sphinxstylestrong{Grado}
\\
\sphinxhline
\sphinxAtStartPar
00
&
\sphinxAtStartPar
Camilo
&
\sphinxAtStartPar
Fonseca
&
\sphinxAtStartPar
9
\\
\sphinxhline
\sphinxAtStartPar
01
&
\sphinxAtStartPar
John
&
\sphinxAtStartPar
Díaz
&
\sphinxAtStartPar
11
\\
\sphinxhline
\sphinxAtStartPar
02
&
\sphinxAtStartPar
Daniel
&
\sphinxAtStartPar
Lopez
&
\sphinxAtStartPar
11
\\
\sphinxhline
\sphinxAtStartPar
03
&
\sphinxAtStartPar
Luis
&
\sphinxAtStartPar
Ruiz
&
\sphinxAtStartPar
8
\\
\sphinxbottomrule
\end{tabulary}
\sphinxtableafterendhook\par
\sphinxattableend\end{savenotes}
\\
\sphinxbottomrule
\end{tabular}
\sphinxtableafterendhook\par
\sphinxattableend\end{savenotes}


\begin{savenotes}\sphinxattablestart
\sphinxthistablewithglobalstyle
\centering
\sphinxcapstartof{table}
\sphinxthecaptionisattop
\sphinxcaption{Tabla docentes}\label{\detokenize{1_guias/8logica/G822:id9}}
\sphinxaftertopcaption
\begin{tabulary}{\linewidth}[t]{TTTT}
\sphinxtoprule
\sphinxtableatstartofbodyhook
\sphinxAtStartPar
\sphinxstylestrong{ID}
&
\sphinxAtStartPar
\sphinxstylestrong{Nombre}
&
\sphinxAtStartPar
\sphinxstylestrong{Apellido}
&
\sphinxAtStartPar
\sphinxstylestrong{Dirección curso}
\\
\sphinxhline
\sphinxAtStartPar
00
&
\sphinxAtStartPar
Saul
&
\sphinxAtStartPar
Duarte
&
\sphinxAtStartPar
10
\\
\sphinxhline
\sphinxAtStartPar
01
&
\sphinxAtStartPar
Dioselina
&
\sphinxAtStartPar
Mena
&
\sphinxAtStartPar
11
\\
\sphinxhline
\sphinxAtStartPar
02
&
\sphinxAtStartPar
Julian
&
\sphinxAtStartPar
Fonseca
&
\sphinxAtStartPar
8
\\
\sphinxbottomrule
\end{tabulary}
\sphinxtableafterendhook\par
\sphinxattableend\end{savenotes}


\begin{savenotes}\sphinxattablestart
\sphinxthistablewithglobalstyle
\centering
\sphinxcapstartof{table}
\sphinxthecaptionisattop
\sphinxcaption{Tabla CURSO}\label{\detokenize{1_guias/8logica/G822:id10}}
\sphinxaftertopcaption
\begin{tabulary}{\linewidth}[t]{TTTTTT}
\sphinxtoprule
\sphinxtableatstartofbodyhook
\sphinxAtStartPar
\sphinxstylestrong{ID}
&
\sphinxAtStartPar
\sphinxstylestrong{ID estudiante}
&
\sphinxAtStartPar
\sphinxstylestrong{Nombre}
&
\sphinxAtStartPar
\sphinxstylestrong{Apellido}
&
\sphinxAtStartPar
\sphinxstylestrong{ID docente}
&
\sphinxAtStartPar
\sphinxstylestrong{Director curso}
\\
\sphinxhline
\sphinxAtStartPar
00
&
\sphinxAtStartPar
01
&
\sphinxAtStartPar
John
&
\sphinxAtStartPar
Díaz
&
\sphinxAtStartPar
01
&
\sphinxAtStartPar
Dioselina
\\
\sphinxhline
\sphinxAtStartPar
01
&
\sphinxAtStartPar
02
&
\sphinxAtStartPar
Daniel
&
\sphinxAtStartPar
Lopez
&
\sphinxAtStartPar
01
&
\sphinxAtStartPar
Dioselina
\\
\sphinxbottomrule
\end{tabulary}
\sphinxtableafterendhook\par
\sphinxattableend\end{savenotes}


\bigskip\hrule\bigskip



\subsection{A02: Taller (+0,8)}
\label{\detokenize{1_guias/8logica/G822:a02-taller-0-8}}
\sphinxAtStartPar
\DUrole{sd-sphinx-override}{\DUrole{sd-badge}{\DUrole{sd-bg-warning}{\DUrole{sd-bg-text-warning}{Actividad práctica}}}}
\begin{enumerate}
\sphinxsetlistlabels{\arabic}{enumi}{enumii}{}{.}%
\item {} 
\sphinxAtStartPar
En su cuaderno responda las siguientes preguntas:
\begin{itemize}
\item {} 
\sphinxAtStartPar
¿Cuál es la diferencia entre los datos e información?

\item {} 
\sphinxAtStartPar
¿Podría existir la información sin datos? ¿Por qué?

\item {} 
\sphinxAtStartPar
¿Podría existir los datos sin información? ¿Por qué?

\end{itemize}

\item {} 
\sphinxAtStartPar
En su cuaderno escriba 10 ejemplos de \sphinxstylestrong{DATOS} y luego transfórmelos en \sphinxstylestrong{INFORMACIÓN}. Tenga en cuenta lo siguiente:
\begin{itemize}
\item {} 
\sphinxAtStartPar
Al momento de transformar los datos a información indicar:
\begin{itemize}
\item {} 
\sphinxAtStartPar
\sphinxstyleemphasis{Contexto}.

\item {} 
\sphinxAtStartPar
\sphinxstyleemphasis{Uso del dato}.

\item {} 
\sphinxAtStartPar
\sphinxstyleemphasis{Consecuencia}.

\end{itemize}

\item {} 
\sphinxAtStartPar
Usar 5 datos cualitativos.

\item {} 
\sphinxAtStartPar
Usar 5 datos cuantitativos.

\end{itemize}

\end{enumerate}


\begin{savenotes}\sphinxattablestart
\sphinxthistablewithglobalstyle
\centering
\begin{tabulary}{\linewidth}[t]{TT}
\sphinxtoprule
\sphinxstyletheadfamily 
\sphinxAtStartPar
Datos
&\sphinxstyletheadfamily 
\sphinxAtStartPar
Información
\\
\sphinxmidrule
\sphinxtableatstartofbodyhook
\sphinxAtStartPar
vacío
&
\sphinxAtStartPar
“El agua fue bebida por el ciclista. El vaso ahora está \sphinxstylestrong{vacío}. Es necesario volver a llenar”.
\\
\sphinxhline
\sphinxAtStartPar
11
&
\sphinxAtStartPar
“En una hora abren las puertas del colegio. Son las \sphinxstylestrong{11} de la mañana. Llegaré tarde”.
\\
\sphinxbottomrule
\end{tabulary}
\sphinxtableafterendhook\par
\sphinxattableend\end{savenotes}
\begin{enumerate}
\sphinxsetlistlabels{\arabic}{enumi}{enumii}{}{.}%
\setcounter{enumi}{2}
\item {} 
\sphinxAtStartPar
En su cuaderno responda las siguientes preguntas:
\begin{itemize}
\item {} 
\sphinxAtStartPar
¿Qué es la Jerarquía de datos?

\item {} 
\sphinxAtStartPar
Con la Jerarquía de datos, ¿Es posible también convertir los datos en INFORMACIÓN? ¿Por qué?

\end{itemize}

\item {} 
\sphinxAtStartPar
En su cuaderno complete la siguiente tabla:


\begin{savenotes}\sphinxattablestart
\sphinxthistablewithglobalstyle
\centering
\begin{tabulary}{\linewidth}[t]{TTT}
\sphinxtoprule
\sphinxstyletheadfamily 
\sphinxAtStartPar
Elemento
&\sphinxstyletheadfamily 
\sphinxAtStartPar
Descripción
&\sphinxstyletheadfamily 
\sphinxAtStartPar
Ejemplo
\\
\sphinxmidrule
\sphinxtableatstartofbodyhook
\sphinxAtStartPar
bit
&&\\
\sphinxhline
\sphinxAtStartPar
Byte
&&\\
\sphinxhline
\sphinxAtStartPar
Carácter
&&\\
\sphinxhline
\sphinxAtStartPar
Dato
&&\\
\sphinxhline
\sphinxAtStartPar
Registro
&&\\
\sphinxhline
\sphinxAtStartPar
Tabla
&&\\
\sphinxhline
\sphinxAtStartPar
Base de datos
&&\\
\sphinxbottomrule
\end{tabulary}
\sphinxtableafterendhook\par
\sphinxattableend\end{savenotes}

\item {} 
\sphinxAtStartPar
DIBUJE la siguiente gráfica en su cuaderno y en cada anillo ubique un elemento de la \sphinxstylestrong{Jerarquía de datos}:

\end{enumerate}


\begin{savenotes}\sphinxattablestart
\sphinxthistablewithglobalstyle
\centering
\begin{tabulary}{\linewidth}[t]{T}
\sphinxtoprule
\sphinxtableatstartofbodyhook

\begin{sphinxfigure-in-table}
\centering
\capstart
\noindent\sphinxincludegraphics[width=400\sphinxpxdimen]{{image0115}.png}
\sphinxfigcaption{\sphinxstylestrong{Anillo} Jerarquía de datos}\label{\detokenize{1_guias/8logica/G822:id11}}\end{sphinxfigure-in-table}\relax
\\
\sphinxbottomrule
\end{tabulary}
\sphinxtableafterendhook\par
\sphinxattableend\end{savenotes}


\begin{enumerate}
\sphinxsetlistlabels{\arabic}{enumi}{enumii}{}{.}%
\setcounter{enumi}{5}
\item {} 
\sphinxAtStartPar
DIBUJE el siguiente gráfico en su cuaderno:


\begin{savenotes}\sphinxattablestart
\sphinxthistablewithglobalstyle
\centering
\begin{tabulary}{\linewidth}[t]{T}
\sphinxtoprule
\sphinxtableatstartofbodyhook

\begin{sphinxfigure-in-table}
\centering
\capstart
\noindent\sphinxincludegraphics[width=550\sphinxpxdimen]{{image0212}.png}
\sphinxfigcaption{EJEMPLO Jerarquía de datos}\label{\detokenize{1_guias/8logica/G822:id12}}\end{sphinxfigure-in-table}\relax
\\
\sphinxbottomrule
\end{tabulary}
\sphinxtableafterendhook\par
\sphinxattableend\end{savenotes}

\sphinxAtStartPar
Acorde al gráfico anterior, en su cuaderno responda:
\begin{itemize}
\item {} 
\sphinxAtStartPar
¿Qué elementos de la jerarquía de datos hacen falta y por qué?

\item {} 
\sphinxAtStartPar
¿Cuantas tablas hay?

\item {} 
\sphinxAtStartPar
¿Cuáles son los registros?

\item {} 
\sphinxAtStartPar
¿Cuáles son los datos?

\item {} 
\sphinxAtStartPar
¿Cuáles son los carácteres?

\item {} 
\sphinxAtStartPar
¿Cuáles son los Bytes?

\item {} 
\sphinxAtStartPar
¿Cuáles son los bits?

\end{itemize}

\item {} 
\sphinxAtStartPar
En su cuaderno (de forma horizontal) realice una \sphinxstylestrong{tabla} de 5 registros de sus compañeros con los siguientes campos (invente los datos):

\end{enumerate}


\begin{savenotes}\sphinxattablestart
\sphinxthistablewithglobalstyle
\centering
\begin{tabular}[t]{*{1}{\X{1}{1}}}
\sphinxtoprule
\sphinxtableatstartofbodyhook

\begin{savenotes}\sphinxattablestart
\sphinxthistablewithglobalstyle
\centering
\begin{tabulary}{\linewidth}[t]{TTTTTTTTT}
\sphinxtoprule
\sphinxstyletheadfamily 
\sphinxAtStartPar
ID
&\sphinxstyletheadfamily 
\sphinxAtStartPar
Nombre
&\sphinxstyletheadfamily 
\sphinxAtStartPar
Apellido
&\sphinxstyletheadfamily 
\sphinxAtStartPar
Edad
&\sphinxstyletheadfamily 
\sphinxAtStartPar
Curso de Ingreso
&\sphinxstyletheadfamily 
\sphinxAtStartPar
Año de Ingreso
&\sphinxstyletheadfamily 
\sphinxAtStartPar
Año de Egreso
&\sphinxstyletheadfamily 
\sphinxAtStartPar
Software más usado
&\sphinxstyletheadfamily 
\sphinxAtStartPar
Hardware más usado
\\
\sphinxmidrule
\sphinxtableatstartofbodyhook
\sphinxAtStartPar
00
&&&&&&&&\\
\sphinxhline
\sphinxAtStartPar
01
&&&&&&&&\\
\sphinxhline
\sphinxAtStartPar
02
&&&&&&&&\\
\sphinxhline
\sphinxAtStartPar
03
&&&&&&&&\\
\sphinxhline
\sphinxAtStartPar
04
&&&&&&&&\\
\sphinxbottomrule
\end{tabulary}
\sphinxtableafterendhook\par
\sphinxattableend\end{savenotes}
\\
\sphinxhline
\begin{center}Tabla Estudiantes
\end{center}\\
\sphinxbottomrule
\end{tabular}
\sphinxtableafterendhook\par
\sphinxattableend\end{savenotes}

\begin{sphinxadmonition}{note}{Advertencia}

\sphinxAtStartPar
Los datos de \sphinxstylestrong{Edad}, \sphinxstylestrong{Curso de Ingreso}, \sphinxstylestrong{Año de Ingreso} y \sphinxstylestrong{Año de Egreso} deben tener lógica.
\end{sphinxadmonition}
\begin{enumerate}
\sphinxsetlistlabels{\arabic}{enumi}{enumii}{}{.}%
\setcounter{enumi}{7}
\item {} 
\sphinxAtStartPar
Construya el \sphinxstylestrong{correo electrónico} del estudiante y agregarlo en la siguiente tabla, teniendo en cuenta la siguiente estructura (ignore los paréntesis):

\end{enumerate}
\begin{equation*}
\begin{split}
(nombre).(apellido)(CursoIngreso)(AñoIngreso)@jfkennedy.edu.co
\end{split}
\end{equation*}

\begin{savenotes}\sphinxattablestart
\sphinxthistablewithglobalstyle
\centering
\sphinxcapstartof{table}
\sphinxthecaptionisattop
\sphinxcaption{Tabla Cuentas Institucionales Estudiantes}\label{\detokenize{1_guias/8logica/G822:id13}}
\sphinxaftertopcaption
\begin{tabular}[t]{*{1}{\X{1}{1}}}
\sphinxtoprule
\sphinxtableatstartofbodyhook

\begin{savenotes}\sphinxattablestart
\sphinxthistablewithglobalstyle
\centering
\begin{tabulary}{\linewidth}[t]{TTTTT}
\sphinxtoprule
\sphinxstyletheadfamily 
\sphinxAtStartPar
ID
&\sphinxstyletheadfamily 
\sphinxAtStartPar
ID Estudiantes
&\sphinxstyletheadfamily 
\sphinxAtStartPar
Nombre
&\sphinxstyletheadfamily 
\sphinxAtStartPar
Apellido
&\sphinxstyletheadfamily 
\sphinxAtStartPar
correo electrónico institucional
\\
\sphinxmidrule
\sphinxtableatstartofbodyhook
\sphinxAtStartPar
00
&&&&\\
\sphinxhline
\sphinxAtStartPar
01
&&&&\\
\sphinxhline
\sphinxAtStartPar
02
&&&&\\
\sphinxhline
\sphinxAtStartPar
03
&&&&\\
\sphinxhline
\sphinxAtStartPar
04
&&&&\\
\sphinxbottomrule
\end{tabulary}
\sphinxtableafterendhook\par
\sphinxattableend\end{savenotes}
\\
\sphinxhline
\begin{center}Tabla Cuentas Institucionales Estudiantes
\end{center}\\
\sphinxbottomrule
\end{tabular}
\sphinxtableafterendhook\par
\sphinxattableend\end{savenotes}
\begin{enumerate}
\sphinxsetlistlabels{\arabic}{enumi}{enumii}{}{.}%
\setcounter{enumi}{8}
\item {} 
\sphinxAtStartPar
En una hoja de su cuaderno (de forma horizontal), dibuje la siguiente base de datos:


\begin{savenotes}\sphinxattablestart
\sphinxthistablewithglobalstyle
\centering
\begin{tabulary}{\linewidth}[t]{T}
\sphinxtoprule
\sphinxtableatstartofbodyhook

\begin{sphinxfigure-in-table}
\centering
\capstart
\noindent\sphinxincludegraphics[width=620\sphinxpxdimen]{{image0311}.png}
\sphinxfigcaption{EJEMPLO Base de datos}\label{\detokenize{1_guias/8logica/G822:id14}}\end{sphinxfigure-in-table}\relax
\\
\sphinxbottomrule
\end{tabulary}
\sphinxtableafterendhook\par
\sphinxattableend\end{savenotes}

\item {} 
\sphinxAtStartPar
Acorde a la \sphinxstylestrong{base de datos}, en su cuaderno responda las siguientes preguntas:
\begin{itemize}
\item {} 
\sphinxAtStartPar
¿Cuál sería el nombre de la base de datos y por qué?

\item {} 
\sphinxAtStartPar
¿Cuantas tablas hay en la base de datos? (colocar el nombre de cada tabla)

\item {} 
\sphinxAtStartPar
¿Cuantos registros hay por cada tabla?

\item {} 
\sphinxAtStartPar
¿Cuantos datos hay por cada registro de cada tabla?

\end{itemize}

\item {} 
\sphinxAtStartPar
En la base de datos (dibujada en su cuaderno), identifique las tablas, registros y los datos.

\item {} 
\sphinxAtStartPar
Con sus palabras, en su cuaderno explique cómo esta compuesta la \sphinxstylestrong{Tabla facturas}.

\item {} 
\sphinxAtStartPar
En su cuaderno, cree un nuevo registro para la \sphinxstylestrong{Tabla clientes}.

\item {} 
\sphinxAtStartPar
En su cuaderno, cree un nuevo registro para la \sphinxstylestrong{Tabla productos}.

\item {} 
\sphinxAtStartPar
Con los nuevos registros \sphinxstylestrong{Tabla clientes} y \sphinxstylestrong{Tabla productos}, ahora cree una nueva \sphinxstyleemphasis{base de datos} compuesta por la \sphinxstylestrong{Tabla facturas} y la \sphinxstylestrong{Tabla elementos en línea}.

\end{enumerate}


\bigskip\hrule\bigskip



\subsection{A03: Examen escrito (+4,0)}
\label{\detokenize{1_guias/8logica/G822:a03-examen-escrito-4-0}}
\sphinxAtStartPar
\DUrole{sd-sphinx-override}{\DUrole{sd-badge}{\DUrole{sd-bg-danger}{\DUrole{sd-bg-text-danger}{Transferencia}}}}


\begin{savenotes}\sphinxattablestart
\sphinxthistablewithglobalstyle
\centering
\begin{tabulary}{\linewidth}[t]{TT}
\sphinxtoprule
\sphinxtableatstartofbodyhook
\sphinxAtStartPar
\sphinxstylestrong{TIEMPO}
&
\sphinxAtStartPar
29 minutos
\\
\sphinxhline
\sphinxAtStartPar
\sphinxstylestrong{PREGUNTAS}
&
\sphinxAtStartPar
10
\\
\sphinxhline
\sphinxAtStartPar
\sphinxstylestrong{OBJETIVO}
&
\sphinxAtStartPar
Comprobar conocimientos adquiridos.
\\
\sphinxhline
\sphinxAtStartPar
\sphinxstylestrong{ACTIVIDAD}
&
\sphinxAtStartPar
Responda las siguientes preguntas a continuación. Cuadernos y celulares guardados. Solo se permite hoja, lapiz y borrador.
\\
\sphinxbottomrule
\end{tabulary}
\sphinxtableafterendhook\par
\sphinxattableend\end{savenotes}

\begin{sphinxadmonition}{note}{Posibles preguntas}
\begin{enumerate}
\sphinxsetlistlabels{\arabic}{enumi}{enumii}{}{.}%
\item {} 
\sphinxAtStartPar
¿Cómo se define un \sphinxstylestrong{bit} y cuáles son sus dos posibles estados?

\item {} 
\sphinxAtStartPar
Explique la analogía de la bombilla utilizada para comprender el funcionamiento de un bit.

\item {} 
\sphinxAtStartPar
¿Cuántos bits conforman un \sphinxstylestrong{BYTE} y por qué es importante esta agrupación?

\item {} 
\sphinxAtStartPar
Si un BYTE puede representar una letra, un número o un símbolo, ¿Cómo se logra esto mediante la combinación de \sphinxstylestrong{8 bits}? De un ejemplo.

\item {} 
\sphinxAtStartPar
¿Cuál es la diferencia fundamental entre un \sphinxstylestrong{BYTE} y un \sphinxstylestrong{carácter}?

\item {} 
\sphinxAtStartPar
Defina qué es un \sphinxstylestrong{carácter} desde la perspectiva de la interfaz de usuario.

\item {} 
\sphinxAtStartPar
¿Qué se entiende por \sphinxstylestrong{Dato} dentro del contexto de los sistemas de información?

\item {} 
\sphinxAtStartPar
¿Por qué se dice que los datos son \sphinxstylestrong{atributos codificados}?

\item {} 
\sphinxAtStartPar
Mencione 10 ejemplos de datos.

\item {} 
\sphinxAtStartPar
¿Cuál es la función principal de un “dato” dentro de la \sphinxstylestrong{Jerarquía de datos}?

\item {} 
\sphinxAtStartPar
Explique la relación que existe entre un \sphinxstylestrong{dato} y un \sphinxstylestrong{registro}.

\item {} 
\sphinxAtStartPar
¿Qué es un \sphinxstylestrong{registro} y por qué es necesario identificarlos?

\item {} 
\sphinxAtStartPar
Describa un escenario (ejemplo propio) de un \sphinxstylestrong{registro} que contenga al menos cuatro tipos de datos distintos.

\item {} 
\sphinxAtStartPar
¿Por qué es necesario que todos los datos dentro de un registro pertenezcan a un mismo \sphinxstylestrong{ente informático}?

\item {} 
\sphinxAtStartPar
¿Qué es una \sphinxstylestrong{tabla} y qué característica estructural debe cumplir para ser considerado como tal?

\item {} 
\sphinxAtStartPar
¿Qué es una \sphinxstylestrong{base de datos} y cómo se diferencia de una tabla individual?

\item {} 
\sphinxAtStartPar
Explique qué significa que una \sphinxstylestrong{base de datos} este compuesta por \sphinxstyleemphasis{\sphinxstylestrong{tablas relacionadas}}.

\item {} 
\sphinxAtStartPar
¿Cuál es la ventaja de utilizar \sphinxstylestrong{tablas cruzadas} en una \sphinxstyleemphasis{base de datos}?

\item {} 
\sphinxAtStartPar
Si tuviera que explicar la \sphinxstylestrong{jerarquía de datos} completa (desde el bit hasta la base de datos) a alguien que no sabe nada de informática, ¿Cómo resumiría la relación entre estos niveles?

\item {} 
\sphinxAtStartPar
Cree un ejemplo de una \sphinxstylestrong{base de datos} con tres tablas de mínimo \sphinxstyleemphasis{3 registros} y \sphinxstyleemphasis{5 datos}.

\end{enumerate}
\end{sphinxadmonition}


\bigskip\hrule\bigskip


\sphinxstepscope


\section{Guía 829: Plan de refuerzo}
\label{\detokenize{1_guias/8logica/G829:guia-829-plan-de-refuerzo}}\label{\detokenize{1_guias/8logica/G829::doc}}\begin{quote}\begin{description}
\sphinxlineitem{ASIGNATURA}
\sphinxAtStartPar
Lógica computacional

\sphinxlineitem{GRADO}
\sphinxAtStartPar
Octavo

\end{description}\end{quote}


\bigskip\hrule\bigskip



\subsection{PRIMER PERIODO}
\label{\detokenize{1_guias/8logica/G829:primer-periodo}}
\begin{sphinxadmonition}{note}{EXAMEN ESCRITO: Plan de refuerzo (10 minutos)}

\sphinxAtStartPar
Estudiar los conceptos a continuación. Se realiza examen de 5 preguntas. Si el estudiante no asiste en la fecha programada, este será evaluado la próxima vez que asista a clase.
\end{sphinxadmonition}


\subsubsection{Hardware}
\label{\detokenize{1_guias/8logica/G829:hardware}}
\sphinxAtStartPar
El hardware es el conjunto de componentes de un sistema informático o computacional que se puede ver o tocar. Por ejemplo:
\begin{itemize}
\item {} 
\sphinxAtStartPar
\sphinxstylestrong{Dispositivos de entrada:} Teclados, mouse, micrófonos.

\item {} 
\sphinxAtStartPar
\sphinxstylestrong{Dispositivos de salida:} Monitores, impresoras, bafles.

\item {} 
\sphinxAtStartPar
\sphinxstylestrong{Procesamiento:} CPU.

\item {} 
\sphinxAtStartPar
\sphinxstylestrong{Memoria:} memoria RAM, memoria CACHE.

\item {} 
\sphinxAtStartPar
\sphinxstylestrong{Almacenamiento:} Discos duros, memorias USB.

\end{itemize}


\subsubsection{Software}
\label{\detokenize{1_guias/8logica/G829:software}}
\sphinxAtStartPar
El SOFTWARE es el conjunto de \sphinxstylestrong{programas}, \sphinxstylestrong{bases de datos}, diferentes tipos de \sphinxstylestrong{archivos} y \sphinxstylestrong{documentación} que le dicen al HARDWARE qué hacer.

\sphinxAtStartPar
A diferencia del hardware que es físico y tangible, el software no se puede tocar. El software solo puede verse su funcionamiento en la pantalla.


\subsubsection{Sistema computacional}
\label{\detokenize{1_guias/8logica/G829:sistema-computacional}}
\sphinxAtStartPar
Es el conjunto de \sphinxstylestrong{hardware} que interactúa con un \sphinxstylestrong{software} para \sphinxstyleemphasis{almacenar}, \sphinxstyleemphasis{procesar} y \sphinxstyleemphasis{mostrar} información. El computador es solo un ejemplo de un sistema computacional.


\subsubsection{Informática}
\label{\detokenize{1_guias/8logica/G829:informatica}}
\sphinxAtStartPar
La informática es la \sphinxstyleemphasis{ciencia} o \sphinxstyleemphasis{disciplina} que estudia el diseño, construcción e implementación de sistemas computacionales (hardware que interactúa con un software para registrar, almacenar, procesar y mostrar la información) para garantizar la integridad y disponibilidad de la información.


\subsubsection{Razonar}
\label{\detokenize{1_guias/8logica/G829:razonar}}
\sphinxAtStartPar
Actividad mental que permite organizar las ideas para llegar a una conclusión (comprobar, deducir).


\subsubsection{Lógica}
\label{\detokenize{1_guias/8logica/G829:logica}}
\sphinxAtStartPar
La lógica es la disciplina que permite analizar y verificar la validez de los razonamientos, determinando si una conclusión se deriva correctamente de unas premisas.


\subsubsection{Lógica computacional}
\label{\detokenize{1_guias/8logica/G829:logica-computacional}}
\sphinxAtStartPar
Aplica los principios de la \sphinxstylestrong{lógica} en los \sphinxstylestrong{sistemas computacionales} para representar la \sphinxstylestrong{información}. La lógica computacional se encarga de:
\begin{itemize}
\item {} 
\sphinxAtStartPar
Establecer las reglas para programar un computador.

\item {} 
\sphinxAtStartPar
Diseñar algoritmos (serie de pasos que se ejecutan de forma ordenada para resolver un problema).

\item {} 
\sphinxAtStartPar
Optimizar los procesos computacionales (hacer más con menos recursos).

\end{itemize}


\bigskip\hrule\bigskip



\subsection{SEGUNDO PERIODO}
\label{\detokenize{1_guias/8logica/G829:segundo-periodo}}
\begin{sphinxadmonition}{note}{EXAMEN ESCRITO: Plan de refuerzo (10 minutos)}

\sphinxAtStartPar
Estudiar los conceptos a continuación. Se realiza examen de 5 preguntas. Si el estudiante no asiste en la fecha programada, este será evaluado la próxima vez que asista a clase.
\end{sphinxadmonition}


\subsubsection{Datos}
\label{\detokenize{1_guias/8logica/G829:datos}}
\sphinxAtStartPar
Los datos son cualquier cantidad, valor o hecho SIN ORGANIZAR, ANALIZAR o PROCESAR. Por ejemplo: rojo, llave, Alexa, 12, 5kg.


\subsubsection{Información}
\label{\detokenize{1_guias/8logica/G829:informacion}}
\sphinxAtStartPar
Es el conjunto de datos organizados, analizados y procesados que tienen significado para las personas. Por ejemplo cuando decimos que: “el carro del ingeniero es de color rojo”, estamos dando una información.


\subsubsection{Jerarquía de datos}
\label{\detokenize{1_guias/8logica/G829:jerarquia-de-datos}}\begin{equation*}
\begin{split}
Base\_de\_datos(Tabla(Registro(Dato(Carácter(Byte(bit))))))
\end{split}
\end{equation*}
\sphinxAtStartPar
La información que procesa una computadora se clasifica en una jerarquía de datos. La información en la computadora se clasifica de MENOR a MAYOR así:
\begin{itemize}
\item {} 
\sphinxAtStartPar
\sphinxstylestrong{bit:} Unidad básica de información en informática. Representado por \sphinxcode{\sphinxupquote{1}} o \sphinxcode{\sphinxupquote{0}}.

\item {} 
\sphinxAtStartPar
\sphinxstylestrong{Byte:} Grupo de ocho bits \sphinxcode{\sphinxupquote{1001 0101}}.

\item {} 
\sphinxAtStartPar
\sphinxstylestrong{Carácter:} Representaciones de Bytes para mostrar en pantalla una letra o un número. \sphinxcode{\sphinxupquote{a}}, \sphinxcode{\sphinxupquote{1}}, \sphinxcode{\sphinxupquote{\$}}.

\item {} 
\sphinxAtStartPar
\sphinxstylestrong{Dato:} Es un valor. Se componen de caracteres. \sphinxcode{\sphinxupquote{175}}, \sphinxcode{\sphinxupquote{camilo}}, \sphinxcode{\sphinxupquote{12}}, \sphinxcode{\sphinxupquote{0.99}}, \sphinxcode{\sphinxupquote{True}}.

\item {} 
\sphinxAtStartPar
\sphinxstylestrong{Registro:} Conjunto de datos que pertenecen a un solo ente informático. En una tabla, el registro correspondería a una fila.

\item {} 
\sphinxAtStartPar
\sphinxstylestrong{Tabla:} Conjunto de registros que tienen la misma estructura.

\item {} 
\sphinxAtStartPar
\sphinxstylestrong{Base de datos:} Es un conjunto de tablas técnicamente organizadas y relacionadas. Una base de datos esta compuesta por varias tablas.

\end{itemize}

\sphinxstepscope


\chapter{Informática 10}
\label{\detokenize{1_guias/10informatica/G1000:informatica-10}}\label{\detokenize{1_guias/10informatica/G1000::doc}}
\sphinxstepscope


\section{Guía 1001: Historia de la Informática (P1)}
\label{\detokenize{1_guias/10informatica/G1001:guia-1001-historia-de-la-informatica-p1}}\label{\detokenize{1_guias/10informatica/G1001::doc}}

\begin{savenotes}\sphinxattablestart
\sphinxthistablewithglobalstyle
\centering
\begin{tabular}[t]{\X{35}{100}\X{65}{100}}
\sphinxtoprule
\sphinxtableatstartofbodyhook

\begin{savenotes}\sphinxattablestart
\sphinxthistablewithglobalstyle
\centering
\begin{tabulary}{\linewidth}[t]{T}
\sphinxtoprule
\sphinxtableatstartofbodyhook
\sphinxAtStartPar
\sphinxincludegraphics{{escudo}.png}
\\
\sphinxbottomrule
\end{tabulary}
\sphinxtableafterendhook\par
\sphinxattableend\end{savenotes}
&

\begin{savenotes}\sphinxattablestart
\sphinxthistablewithglobalstyle
\raggedright
\begin{tabular}[t]{*{1}{\X{1}{1}}}
\sphinxtoprule
\sphinxtableatstartofbodyhook
\sphinxAtStartPar
\sphinxstylestrong{INSTITUCIÓN EDUCATIVA}: John F. Kennedy
\begin{itemize}
\item {} 
\sphinxAtStartPar
\sphinxstylestrong{DOCENTE}: Julian Camilo Fonseca Romero

\item {} 
\sphinxAtStartPar
\sphinxstylestrong{ASIGNATURA}: Tecnología e Informática

\item {} 
\sphinxAtStartPar
\sphinxstylestrong{GRADO}: DÉCIMO

\item {} 
\sphinxAtStartPar
\sphinxstylestrong{PERIODO}: 01

\item {} 
\sphinxAtStartPar
\sphinxstylestrong{TIEMPO}: 10 semanas (2 horas / semana)

\item {} 
\sphinxAtStartPar
\sphinxstylestrong{TIEMPO DE TRABAJO}: 20/20 horas

\end{itemize}
\\
\sphinxbottomrule
\end{tabular}
\sphinxtableafterendhook\par
\sphinxattableend\end{savenotes}
\\
\sphinxbottomrule
\end{tabular}
\sphinxtableafterendhook\par
\sphinxattableend\end{savenotes}

\begin{sphinxuseclass}{sd-tab-set}
\begin{sphinxuseclass}{sd-tab-item}\subsubsection*{Objetivo}

\begin{sphinxuseclass}{sd-tab-content}
\sphinxAtStartPar
\DUrole{sd-sphinx-override}{\DUrole{sd-badge}{\DUrole{sd-bg-light}{\DUrole{sd-bg-text-light}{Exploración}}}}
\DUrole{sd-sphinx-override}{\DUrole{sd-badge}{\DUrole{sd-bg-info}{\DUrole{sd-bg-text-info}{Estructuración}}}}
\DUrole{sd-sphinx-override}{\DUrole{sd-badge}{\DUrole{sd-bg-warning}{\DUrole{sd-bg-text-warning}{Actividad práctica}}}}
\DUrole{sd-sphinx-override}{\DUrole{sd-badge}{\DUrole{sd-bg-danger}{\DUrole{sd-bg-text-danger}{Transferencia}}}}
\begin{itemize}
\item {} 
\sphinxAtStartPar
\sphinxstylestrong{Temática}: Historia de la computación personal.

\item {} 
\sphinxAtStartPar
\sphinxstylestrong{Objetivo de aprendizaje}: Conocer los orígenes de la informática por medio del estudio de la evolución del internet, la computación personal y el software.

\end{itemize}

\end{sphinxuseclass}
\end{sphinxuseclass}
\begin{sphinxuseclass}{sd-tab-item}\subsubsection*{Orientaciones curriculares}

\begin{sphinxuseclass}{sd-tab-content}\begin{itemize}
\item {} 
\sphinxAtStartPar
\sphinxstylestrong{Componente}: Naturaleza y evolución de la TI.

\item {} 
\sphinxAtStartPar
\sphinxstylestrong{Competencia}: Construyo conocimientos y saberes de base tecnológica e informática para la toma de decisiones en el desarrollo de productos tecnológicos.

\item {} 
\sphinxAtStartPar
\sphinxstylestrong{Evidencia de aprendizaje}: Sistematizo la evolución tecnológica e informática y la manera cómo incidieron en los cambios estructurales de la sociedad y la cultura a lo largo de la historia.

\end{itemize}

\end{sphinxuseclass}
\end{sphinxuseclass}
\begin{sphinxuseclass}{sd-tab-item}\subsubsection*{Criterios de evaluación}

\begin{sphinxuseclass}{sd-tab-content}\begin{itemize}
\item {} 
\sphinxAtStartPar
Actividades desarrolladas en su totalidad.

\item {} 
\sphinxAtStartPar
Participación en clase.

\item {} 
\sphinxAtStartPar
Explicación clara de conceptos.

\item {} 
\sphinxAtStartPar
Argumentación de ideas.

\item {} 
\sphinxAtStartPar
Cumplimiento de las reglas del aula.

\end{itemize}

\end{sphinxuseclass}
\end{sphinxuseclass}
\begin{sphinxuseclass}{sd-tab-item}\subsubsection*{Recursos (0)}

\begin{sphinxuseclass}{sd-tab-content}\begin{itemize}
\item {} 
\sphinxAtStartPar
Esta guía no contiene recursos.

\end{itemize}

\end{sphinxuseclass}
\end{sphinxuseclass}
\end{sphinxuseclass}

\subsection{Mensaje del autor}
\label{\detokenize{1_guias/10informatica/G1001:mensaje-del-autor}}\begin{quote}

\sphinxAtStartPar
“En esta guía, el estudiante aprende sobre el origen de la computación personal para entender la evolución de nuestra propia cultura. Se invita a explorar cómo el desarrollo de Internet, junto con la evolución del hardware y el software, transformó la estructura de la sociedad actual. Al finalizar, el estudiante podrá sistematizar estos hitos históricos y reconocer la huella que la informática ha dejado en cada aspecto de nuestra vida cotidiana.”

\begin{flushright}
---Camilo Fonseca
\end{flushright}
\end{quote}


\bigskip\hrule\bigskip



\subsection{A00: Introducción (+0,1)}
\label{\detokenize{1_guias/10informatica/G1001:a00-introduccion-0-1}}
\sphinxAtStartPar
\DUrole{sd-sphinx-override}{\DUrole{sd-badge}{\DUrole{sd-bg-light}{\DUrole{sd-bg-text-light}{Exploración}}}}
\begin{enumerate}
\sphinxsetlistlabels{\arabic}{enumi}{enumii}{}{.}%
\item {} 
\sphinxAtStartPar
Transcriba en su cuaderno el objetivo, las orientaciones curriculares y los criterios de evaluación de la presente guía.

\item {} 
\sphinxAtStartPar
Lea el \sphinxstylestrong{mensaje del autor} y en su cuaderno responda:
\begin{itemize}
\item {} 
\sphinxAtStartPar
¿Por qué es importante estudiar la historia de la informática?

\end{itemize}

\end{enumerate}


\bigskip\hrule\bigskip



\subsection{A01: Historia de la Informática (+0,1)}
\label{\detokenize{1_guias/10informatica/G1001:a01-historia-de-la-informatica-0-1}}
\sphinxAtStartPar
\DUrole{sd-sphinx-override}{\DUrole{sd-badge}{\DUrole{sd-bg-info}{\DUrole{sd-bg-text-info}{Estructuración}}}}
\begin{enumerate}
\sphinxsetlistlabels{\arabic}{enumi}{enumii}{}{.}%
\item {} 
\sphinxAtStartPar
Lea las siguientes historias del INTERNET, la computación personal y el software:

\end{enumerate}


\bigskip\hrule\bigskip



\subsubsection{Historia del INTERNET (ARPANET)}
\label{\detokenize{1_guias/10informatica/G1001:historia-del-internet-arpanet}}
\sphinxAtStartPar
%
\begin{footnote}[1]\sphinxAtStartFootnote
ARPANET. (2005, Feb 10). Sistemas.com. Retrieved April 25, 2026, from https://sistemas.com/arpanet.php
%
\end{footnote} El Internet tiene sus orígenes en \sphinxstylestrong{ARPANET}, un proyecto militar desarrollado en 1969 por el \sphinxstyleemphasis{Departamento de defensa de Estados Unidos}. Esta red inicial conectaba un número reducido de computadoras en distintas universidades.


\begin{savenotes}\sphinxattablestart
\sphinxthistablewithglobalstyle
\centering
\begin{tabulary}{\linewidth}[t]{T}
\sphinxtoprule
\sphinxtableatstartofbodyhook

\begin{sphinxfigure-in-table}
\centering
\capstart
\noindent\sphinxincludegraphics[width=500\sphinxpxdimen]{{image012}.png}
\sphinxfigcaption{\sphinxstylestrong{ARPANET:} Primera red del mundo}\label{\detokenize{1_guias/10informatica/G1001:id26}}\end{sphinxfigure-in-table}\relax
\\
\sphinxbottomrule
\end{tabulary}
\sphinxtableafterendhook\par
\sphinxattableend\end{savenotes}

\sphinxAtStartPar
En los años 70 y 80 (exactamente en el año 1983), se desarrollaron los protocolos TCP/IP, que permitieron la comunicación entre diferentes redes de computadoras.

\sphinxAtStartPar
%
\begin{footnote}[2]\sphinxAtStartFootnote
Rojas, J. (2022, May 17). El primer sitio web de la historia. Un hito que cambió vidas. Telefónica. https://www.telefonica.com/es/sala\sphinxhyphen{}comunicacion/blog/el\sphinxhyphen{}primer\sphinxhyphen{}sitio\sphinxhyphen{}web\sphinxhyphen{}de\sphinxhyphen{}la\sphinxhyphen{}historia\sphinxhyphen{}un\sphinxhyphen{}hito\sphinxhyphen{}que\sphinxhyphen{}cambio\sphinxhyphen{}vidas/
%
\end{footnote} La década de los 90 marcó un punto de inflexión con la creación de la \sphinxstylestrong{World Wide Web} por \sphinxstylestrong{Tim Berners\sphinxhyphen{}Lee}, que introdujo las páginas web y los navegadores en 1989. A continuación se muestra la primera página web de la historia:


\begin{savenotes}\sphinxattablestart
\sphinxthistablewithglobalstyle
\centering
\begin{tabulary}{\linewidth}[t]{T}
\sphinxtoprule
\sphinxtableatstartofbodyhook

\begin{sphinxfigure-in-table}
\centering
\capstart
\noindent\sphinxincludegraphics[width=700\sphinxpxdimen]{{image02}.png}
\sphinxfigcaption{\sphinxstylestrong{Primera página web:} Word Wide Web}\label{\detokenize{1_guias/10informatica/G1001:id27}}\end{sphinxfigure-in-table}\relax
\\
\sphinxbottomrule
\end{tabulary}
\sphinxtableafterendhook\par
\sphinxattableend\end{savenotes}

\sphinxAtStartPar
%
\begin{footnote}[3]\sphinxAtStartFootnote
Pastor, J. (2018, April 28). Cuando Mosaic dominaba el mundo (de los navegadores). Xataka.com; Xataka. https://www.xataka.com/historia\sphinxhyphen{}tecnologica/cuando\sphinxhyphen{}mosaic\sphinxhyphen{}dominaba\sphinxhyphen{}el\sphinxhyphen{}mundo\sphinxhyphen{}de\sphinxhyphen{}los\sphinxhyphen{}navegadores
%
\end{footnote} Para facilitar el acceso a diferentes páginas web, en 1993 se creo el navegador web MOSAIC. Previamente existía otro navegador de nombre Nexus pero MOSAIC fue el primer navegador en usarse de forma masiva.


\begin{savenotes}\sphinxattablestart
\sphinxthistablewithglobalstyle
\centering
\begin{tabulary}{\linewidth}[t]{T}
\sphinxtoprule
\sphinxtableatstartofbodyhook

\begin{sphinxfigure-in-table}
\centering
\capstart
\noindent\sphinxincludegraphics[width=700\sphinxpxdimen]{{image03}.png}
\sphinxfigcaption{\sphinxstylestrong{MOSAIC:} Primer navegador web en usarse de forma masiva.}\label{\detokenize{1_guias/10informatica/G1001:id28}}\end{sphinxfigure-in-table}\relax
\\
\sphinxbottomrule
\end{tabulary}
\sphinxtableafterendhook\par
\sphinxattableend\end{savenotes}

\sphinxAtStartPar
%
\begin{footnote}[4]\sphinxAtStartFootnote
Eadicicco, L. (2025, May 11). La tecnologí­a que definió el internet moderno está cambiando y Silicon Valley finalmente lo admite. CNN en Español. https://cnnespanol.cnn.com/2025/05/11/ciencia/tecnologia\sphinxhyphen{}internet\sphinxhyphen{}moderno\sphinxhyphen{}silicon\sphinxhyphen{}valley\sphinxhyphen{}trax
%
\end{footnote} El \sphinxstylestrong{Internet moderno} ha evolucionado considerablemente desde entonces, pasando por la era de la Web 2.0 (redes sociales e interactividad las cuales tienen su origen en el año 2004), hasta la actual era del Internet móvil cuyo origen fue en el año 2007.  Hoy en día la nube y el Internet de las Cosas (IoT) conectan a miles de millones de dispositivos y usuarios en todo el mundo.


\bigskip\hrule\bigskip



\subsubsection{Historia de la Computación Personal}
\label{\detokenize{1_guias/10informatica/G1001:historia-de-la-computacion-personal}}
\sphinxAtStartPar
La historia de la computadora personal comienza en la década de 1970.


\begin{savenotes}\sphinxattablestart
\sphinxthistablewithglobalstyle
\centering
\begin{tabulary}{\linewidth}[t]{T}
\sphinxtoprule
\sphinxtableatstartofbodyhook

\begin{sphinxfigure-in-table}
\centering
\capstart
\noindent\sphinxincludegraphics[width=300\sphinxpxdimen]{{image04}.png}
\sphinxfigcaption{\sphinxstylestrong{REVISTA:} Popular Electronics}\label{\detokenize{1_guias/10informatica/G1001:id29}}\end{sphinxfigure-in-table}\relax
\\
\sphinxbottomrule
\end{tabulary}
\sphinxtableafterendhook\par
\sphinxattableend\end{savenotes}

\sphinxAtStartPar
%
\begin{footnote}[5]\sphinxAtStartFootnote
Pastor, J. (2015, January 27). Altair 8800: todo comenzó hace ya 40 años. Xataka.com; Xataka. https://www.xataka.com/historia\sphinxhyphen{}tecnologica/altair\sphinxhyphen{}8800\sphinxhyphen{}todo\sphinxhyphen{}comenzo\sphinxhyphen{}hace\sphinxhyphen{}ya\sphinxhyphen{}40\sphinxhyphen{}anos
%
\end{footnote} En 1975, la \sphinxstylestrong{Altair 8800} se convirtió en la \sphinxstylestrong{primera computadora personal} comercialmente exitosa, aunque requería ser ensamblada por el usuario. Esta computadora más adelante sería más fácil de programar por la creación de un lenguaje de programación llamado \sphinxstylestrong{BASIC}. Anteriormente se programaba con el lenguaje \sphinxstylestrong{ENSAMBLADOR}.

\sphinxAtStartPar
%
\begin{footnote}[6]\sphinxAtStartFootnote
Catucci, A. (2020, September 9). La historia de Apple: repaso por este caso de éxito tecnológico. Marketing Insider Review. https://marketinginsiderreview.com/historia\sphinxhyphen{}apple\sphinxhyphen{}exito\sphinxhyphen{}tecnologico/
%
\end{footnote} En 1976, Apple Computer lanzó la \sphinxstylestrong{Apple I}, seguida en 1977 por la revolucionaria \sphinxstylestrong{Apple II}, una de las primeras computadoras con gráficos a color.

\sphinxAtStartPar
%
\begin{footnote}[7]\sphinxAtStartFootnote
The Conversation. (2024, January 22). La Mac cumple 40 años: debutó con este antológico aviso dirigido por Ridley Scott para mostrar cómo había otra forma de usar una computadora. LA NACION. https://www.lanacion.com.ar/tecnologia/la\sphinxhyphen{}mac\sphinxhyphen{}cumple\sphinxhyphen{}40\sphinxhyphen{}anos\sphinxhyphen{}debuto\sphinxhyphen{}con\sphinxhyphen{}este\sphinxhyphen{}antologico\sphinxhyphen{}aviso\sphinxhyphen{}dirigido\sphinxhyphen{}por\sphinxhyphen{}ridley\sphinxhyphen{}scott\sphinxhyphen{}para\sphinxhyphen{}mostrar\sphinxhyphen{}como\sphinxhyphen{}nid22012024/
%
\end{footnote} Los años 80 vieron la llegada de la primera \sphinxstylestrong{Macintosh}, conocida como \sphinxstylestrong{Macintosh 128K}, fue presentada el \sphinxstylestrong{24 de enero de 1984} por \sphinxstylestrong{Steve Jobs}. Este ordenador personal marcó un hito en la historia de la computación por varias razones:
\begin{itemize}
\item {} 
\sphinxAtStartPar
\sphinxstylestrong{Interfaz Gráfica de Usuario (GUI)}: Fue uno de los primeros ordenadores en popularizar el uso de una interfaz gráfica, lo que facilitó su uso para personas no técnicas.

\item {} 
\sphinxAtStartPar
\sphinxstylestrong{Ratón}: Incluía un ratón, lo que era innovador en ese momento y permitía una interacción más intuitiva con el sistema.

\item {} 
\sphinxAtStartPar
\sphinxstylestrong{Diseño Compacto}: Su diseño era más pequeño y accesible en comparación con otros ordenadores de la época, lo que lo hacía ideal para el uso personal.

\item {} 
\sphinxAtStartPar
\sphinxstylestrong{Software}: Venía con aplicaciones básicas como MacPaint y MacWrite, que demostraban las capacidades gráficas y de procesamiento de texto del dispositivo.

\end{itemize}


\begin{savenotes}\sphinxattablestart
\sphinxthistablewithglobalstyle
\centering
\begin{tabulary}{\linewidth}[t]{T}
\sphinxtoprule
\sphinxtableatstartofbodyhook

\begin{sphinxfigure-in-table}
\centering
\capstart
\noindent\sphinxincludegraphics[width=600\sphinxpxdimen]{{image05}.png}
\sphinxfigcaption{\sphinxstylestrong{Artículo:} Descripción de la Macintosh}\label{\detokenize{1_guias/10informatica/G1001:id30}}\end{sphinxfigure-in-table}\relax
\\
\sphinxbottomrule
\end{tabulary}
\sphinxtableafterendhook\par
\sphinxattableend\end{savenotes}

\sphinxAtStartPar
La \sphinxstylestrong{Macintosh 128K} no solo fue un éxito comercial, sino que también sentó las bases para el desarrollo de futuras computadoras de Apple.

\sphinxAtStartPar
%
\begin{footnote}[8]\sphinxAtStartFootnote
Historia de Microsoft. (2012, November 10). Blog de tecnología; Sistemas Ibertrónica. https://ibertronica.es/blog/productos/microsoft\sphinxhyphen{}pasado\sphinxhyphen{}presente\sphinxhyphen{}y\sphinxhyphen{}futuro/
%
\end{footnote} En los años 90, Microsoft democratizaron el uso de las computadoras personales con el desarrollo de su sistema operativo Windows. El siglo XXI trajo las \sphinxstyleemphasis{laptops ultradelgadas}, \sphinxstyleemphasis{tablets} y \sphinxstyleemphasis{mini\sphinxhyphen{}computadoras}.


\begin{savenotes}\sphinxattablestart
\sphinxthistablewithglobalstyle
\centering
\begin{tabulary}{\linewidth}[t]{T}
\sphinxtoprule
\sphinxtableatstartofbodyhook

\begin{sphinxfigure-in-table}
\centering
\capstart
\noindent\sphinxincludegraphics[width=500\sphinxpxdimen]{{image06}.png}
\sphinxfigcaption{\sphinxstylestrong{PC:} Raspberry Pi}\label{\detokenize{1_guias/10informatica/G1001:id31}}\end{sphinxfigure-in-table}\relax
\\
\sphinxbottomrule
\end{tabulary}
\sphinxtableafterendhook\par
\sphinxattableend\end{savenotes}

\sphinxAtStartPar
%
\begin{footnote}[9]\sphinxAtStartFootnote
Andrés, R. (2019, September 1). Raspberry Pi: la historia de un mini\sphinxhyphen{}PC que se ha convertido en un éxito de masas. Computer Hoy. https://computerhoy.20minutos.es/reportajes/tecnologia/historia\sphinxhyphen{}raspberry\sphinxhyphen{}pi\sphinxhyphen{}482477
%
\end{footnote} En 2012, la \sphinxstylestrong{Raspberry Pi} revolucionó el concepto de computadora personal al ofrecer un dispositivo del tamaño de una tarjeta de crédito a bajo costo, enfocado en la educación y proyectos de programación.


\bigskip\hrule\bigskip



\subsubsection{Historia del software}
\label{\detokenize{1_guias/10informatica/G1001:historia-del-software}}
\sphinxAtStartPar
El software tiene su origen en los primeros días de la computación. En la década de 1940, las primeras computadoras ejecutaban instrucciones básicas mediante cables e interruptores físicos.


\begin{savenotes}\sphinxattablestart
\sphinxthistablewithglobalstyle
\centering
\begin{tabulary}{\linewidth}[t]{T}
\sphinxtoprule
\sphinxtableatstartofbodyhook

\begin{sphinxfigure-in-table}
\centering
\capstart
\noindent\sphinxincludegraphics[width=600\sphinxpxdimen]{{image07}.png}
\sphinxfigcaption{Primera computadora programable 1940}\label{\detokenize{1_guias/10informatica/G1001:id32}}\end{sphinxfigure-in-table}\relax
\\
\sphinxbottomrule
\end{tabulary}
\sphinxtableafterendhook\par
\sphinxattableend\end{savenotes}

\sphinxAtStartPar
%
\begin{footnote}[10]\sphinxAtStartFootnote
Huertos, A. A. (2019, May 28). La historia de los lenguajes de programación. Computer Hoy. https://computerhoy.20minutos.es/reportajes/tecnologia/historia\sphinxhyphen{}lenguajes\sphinxhyphen{}programacion\sphinxhyphen{}428041
%
\end{footnote} Los años 50 marcaron el inicio de los lenguajes de programación modernos, con la creación del lenguaje de alto nivel \sphinxstylestrong{FORTRAN} (1957).

\sphinxAtStartPar
%
\begin{footnote}[11]\sphinxAtStartFootnote
Velásquez, W. (2022, April 1). 3 historias del software libre que deberías conocer. Wajari. https://wajari.com/blog/historias\sphinxhyphen{}software\sphinxhyphen{}libre\sphinxhyphen{}open\sphinxhyphen{}source/
%
\end{footnote} En el año 1969 se desarrolló el sistema operativo \sphinxstylestrong{UNIX}. La mayoría de sistemas operativos en la actualidad se basan en \sphinxstylestrong{UNIX}.


\begin{savenotes}\sphinxattablestart
\sphinxthistablewithglobalstyle
\centering
\begin{tabulary}{\linewidth}[t]{T}
\sphinxtoprule
\sphinxtableatstartofbodyhook

\begin{sphinxfigure-in-table}
\centering
\capstart
\noindent\sphinxincludegraphics[width=600\sphinxpxdimen]{{image08}.png}
\sphinxfigcaption{Consola terminal UNIX}\label{\detokenize{1_guias/10informatica/G1001:id33}}\end{sphinxfigure-in-table}\relax
\\
\sphinxbottomrule
\end{tabulary}
\sphinxtableafterendhook\par
\sphinxattableend\end{savenotes}

\sphinxAtStartPar
Con la llegada de \sphinxstylestrong{ALTAIR} en 1975, Microsoft desarrolla el lenguaje de programación BASIC. Este permitió programar en dicha computadora.

\sphinxAtStartPar
%
\begin{footnote}[12]\sphinxAtStartFootnote
Maturana, J. (2023, June 2). 25 aniversario de Computer Hoy: La historia de Windows, cómo ha sido la evolución del servicio más popular de Microsoft. Computer Hoy. https://computerhoy.20minutos.es/windows/historia\sphinxhyphen{}windows\sphinxhyphen{}como\sphinxhyphen{}ha\sphinxhyphen{}sido\sphinxhyphen{}evolucion\sphinxhyphen{}servicio\sphinxhyphen{}popular\sphinxhyphen{}microsoft\sphinxhyphen{}1237508
%
\end{footnote} La década de 1980 fue revolucionaria con la llegada del software comercial para PC, incluyendo procesadores de texto, hojas de cálculo y los primeros sistemas operativos con interfaz gráfica. Esto inicio el año 1981 con la creación del primer sistema operativo de Microsoft: \sphinxstylestrong{MS\sphinxhyphen{}DOS}.


\begin{savenotes}\sphinxattablestart
\sphinxthistablewithglobalstyle
\centering
\begin{tabulary}{\linewidth}[t]{T}
\sphinxtoprule
\sphinxtableatstartofbodyhook

\begin{sphinxfigure-in-table}
\centering
\capstart
\noindent\sphinxincludegraphics[width=200\sphinxpxdimen]{{image09}.png}
\sphinxfigcaption{Logo MSDOS}\label{\detokenize{1_guias/10informatica/G1001:id34}}\end{sphinxfigure-in-table}\relax
\\
\sphinxbottomrule
\end{tabulary}
\sphinxtableafterendhook\par
\sphinxattableend\end{savenotes}

\sphinxAtStartPar
\sphinxfootnotemark[11] En el año 1991 una comunidad de ingenieros y desarrolladores a nivel mundial crean el sistema operativo \sphinxstylestrong{LINUX}. El cual es de gran importancia hasta nuestros días.


\begin{savenotes}\sphinxattablestart
\sphinxthistablewithglobalstyle
\centering
\begin{tabulary}{\linewidth}[t]{T}
\sphinxtoprule
\sphinxtableatstartofbodyhook

\begin{sphinxfigure-in-table}
\centering
\capstart
\noindent\sphinxincludegraphics[width=400\sphinxpxdimen]{{image10}.png}
\sphinxfigcaption{Logo LINUX}\label{\detokenize{1_guias/10informatica/G1001:id35}}\end{sphinxfigure-in-table}\relax
\\
\sphinxbottomrule
\end{tabulary}
\sphinxtableafterendhook\par
\sphinxattableend\end{savenotes}

\sphinxAtStartPar
En la actualidad, el software ha evolucionado hacia aplicaciones en la nube, inteligencia artificial y desarrollo ágil, transformando completamente la manera en la que interactuamos con la tecnología.


\bigskip\hrule\bigskip



\subsection{A02: Taller (+0,8)}
\label{\detokenize{1_guias/10informatica/G1001:a02-taller-0-8}}
\sphinxAtStartPar
\DUrole{sd-sphinx-override}{\DUrole{sd-badge}{\DUrole{sd-bg-warning}{\DUrole{sd-bg-text-warning}{Actividad práctica}}}}
\begin{enumerate}
\sphinxsetlistlabels{\arabic}{enumi}{enumii}{}{.}%
\item {} 
\sphinxAtStartPar
En su cuaderno, responda las siguientes preguntas:
\begin{itemize}
\item {} 
\sphinxAtStartPar
¿Cómo se llamó la primera red?

\item {} 
\sphinxAtStartPar
¿En que año y en que lugar se creó la primera red?

\item {} 
\sphinxAtStartPar
¿Quienes crearon la tecnología para que funcionara el Internet?

\item {} 
\sphinxAtStartPar
¿En que año se creó la tecnología para que funcionara el Internet y que protocolos fueron usados en dicha tecnología?

\end{itemize}

\item {} 
\sphinxAtStartPar
En su cuaderno dibuje el mapa de la primera red y responda las siguientes preguntas:
\begin{itemize}
\item {} 
\sphinxAtStartPar
¿Cuántos nodos en el occidente estaban interconectados?

\item {} 
\sphinxAtStartPar
¿Cuántos nodos en el oriente estaban interconectados?

\end{itemize}

\item {} 
\sphinxAtStartPar
Responda la siguiente pregunta en su cuaderno:
\begin{itemize}
\item {} 
\sphinxAtStartPar
¿En que año se crea el primer navegador y la primera página web?

\end{itemize}

\item {} 
\sphinxAtStartPar
Traduzca los contenidos de la primera página web y haga un breve resumen de dichos contenidos en su cuaderno.

\item {} 
\sphinxAtStartPar
En su cuaderno dibuje un \sphinxstylestrong{boceto} de la \sphinxstyleemphasis{interfaz gráfica} del primer navegador web \sphinxstylestrong{MOSAIC} y realice una breve descripción de las diferencias de este primer navegador con navegadores actuales como GOOGLE CHROME, MOZILLA FIREFOX o MICROSOFT EDGE.

\item {} 
\sphinxAtStartPar
En su cuaderno responda las siguientes preguntas:
\begin{itemize}
\item {} 
\sphinxAtStartPar
¿Cuál fue la primera computadora personal de la historia y en que año se creó?

\item {} 
\sphinxAtStartPar
¿Qué lenguajes de programación se usaron en la primera computadora personal?

\item {} 
\sphinxAtStartPar
¿Cómo APPLE se relaciona con la historia de la computadora personal?

\item {} 
\sphinxAtStartPar
¿Quién fue Steve Jobs?

\item {} 
\sphinxAtStartPar
¿Cómo se democratizó el uso de las computadoras personales?

\end{itemize}

\item {} 
\sphinxAtStartPar
En su cuaderno responda:
\begin{itemize}
\item {} 
\sphinxAtStartPar
¿Qué características tenía la primera Macintosh y en que año fue creada?

\end{itemize}

\item {} 
\sphinxAtStartPar
En su cuaderno dibuje un bosquejo de la primera Macintosh.

\item {} 
\sphinxAtStartPar
Observe el artículo donde se describe la Macintosh y en su cuaderno describa el significado de las siguientes frases:
\begin{itemize}
\item {} 
\sphinxAtStartPar
\sphinxstyleemphasis{“La primera APPLE que se puede llevar en un bolso”}

\item {} 
\sphinxAtStartPar
\sphinxstyleemphasis{“Si puedes señalar, puedes usar Macintosh”}

\item {} 
\sphinxAtStartPar
\sphinxstyleemphasis{“PERSONALIDAD DE MACINTOSH: Lado divertido, Lado serio”}

\item {} 
\sphinxAtStartPar
\sphinxstyleemphasis{“Muy pronto habrán dos tipos de personas: Aquellos que usan computadoras y aquellos que usan APPLES”}

\end{itemize}

\item {} 
\sphinxAtStartPar
En su cuaderno responda las siguientes preguntas:
\begin{itemize}
\item {} 
\sphinxAtStartPar
¿En que año y como operaban las computadoras en sus orígenes?

\item {} 
\sphinxAtStartPar
¿Qué es y en que año se crea \sphinxstylestrong{FORTRAN}?

\item {} 
\sphinxAtStartPar
¿Por qué la década de los 80 fue revolucionaria en el software?

\end{itemize}

\item {} 
\sphinxAtStartPar
Investigue en \sphinxstylestrong{INTERNET} o \sphinxstylestrong{IA} y realice una descripción del sistema operativo \sphinxstylestrong{UNIX} luego en su cuaderno responda la siguiente pregunta:
\begin{itemize}
\item {} 
\sphinxAtStartPar
¿Por qué UNIX es considerado el padre de la mayoría de sistemas operativos en la actualidad?

\end{itemize}

\item {} 
\sphinxAtStartPar
Investigue en \sphinxstylestrong{INTERNET} la diferencia entre un \sphinxstylestrong{equipo cliente} y un \sphinxstylestrong{equipo servidor} y luego realice una descripción en su cuaderno.

\item {} 
\sphinxAtStartPar
Investigue en \sphinxstylestrong{INTERNET} que es el sistema operativo \sphinxstylestrong{LINUX} luego realice una descripción en su cuaderno e indique por que se usa en la mayoría de servidores a nivel mundial.

\item {} 
\sphinxAtStartPar
En su cuaderno describa 10 distribuciones \sphinxstylestrong{LINUX} (con sus fechas de creación).

\item {} 
\sphinxAtStartPar
En su cuaderno responda:
\begin{itemize}
\item {} 
\sphinxAtStartPar
¿En que año se creo el sistema operativo \sphinxstylestrong{Microsoft DOS (MS\sphinxhyphen{}DOS)} y por que fue tan importante?

\end{itemize}

\item {} 
\sphinxAtStartPar
Con ayuda de INTERNET o IA, en su cuaderno realice una descripción (con fechas de creación) de los siguientes sistemas operativos de Microsoft:
\begin{itemize}
\item {} 
\sphinxAtStartPar
Windows 3.1

\item {} 
\sphinxAtStartPar
Windows 95

\item {} 
\sphinxAtStartPar
Windows 98

\item {} 
\sphinxAtStartPar
Windows 2000

\item {} 
\sphinxAtStartPar
Windows Me

\item {} 
\sphinxAtStartPar
Windows XP

\item {} 
\sphinxAtStartPar
Windows Vista

\item {} 
\sphinxAtStartPar
Windows 7

\item {} 
\sphinxAtStartPar
Windows 8

\item {} 
\sphinxAtStartPar
Windows 8.1

\item {} 
\sphinxAtStartPar
Windows 10

\item {} 
\sphinxAtStartPar
Windows 11

\end{itemize}

\end{enumerate}


\bigskip\hrule\bigskip



\subsection{A03: Línea de tiempo (+1,0)}
\label{\detokenize{1_guias/10informatica/G1001:a03-linea-de-tiempo-1-0}}
\sphinxAtStartPar
\DUrole{sd-sphinx-override}{\DUrole{sd-badge}{\DUrole{sd-bg-warning}{\DUrole{sd-bg-text-warning}{Actividad práctica}}}}
\begin{enumerate}
\sphinxsetlistlabels{\arabic}{enumi}{enumii}{}{.}%
\item {} 
\sphinxAtStartPar
En \sphinxstylestrong{medio pliego de cartulina}, cree una línea de tiempo de la \sphinxstylestrong{historia de la computación personal}, desde los orígenes de la primera red con \sphinxstylestrong{ARPANET} hasta la creación de los \sphinxstylestrong{sistemas operativos de Microsoft}.

\end{enumerate}


\bigskip\hrule\bigskip



\subsection{A04: Examen escrito (+3,0)}
\label{\detokenize{1_guias/10informatica/G1001:a04-examen-escrito-3-0}}
\sphinxAtStartPar
\DUrole{sd-sphinx-override}{\DUrole{sd-badge}{\DUrole{sd-bg-danger}{\DUrole{sd-bg-text-danger}{Transferencia}}}}


\begin{savenotes}\sphinxattablestart
\sphinxthistablewithglobalstyle
\centering
\begin{tabulary}{\linewidth}[t]{TT}
\sphinxtoprule
\sphinxtableatstartofbodyhook
\sphinxAtStartPar
\sphinxstylestrong{TIEMPO}
&
\sphinxAtStartPar
29 minutos
\\
\sphinxhline
\sphinxAtStartPar
\sphinxstylestrong{PREGUNTAS}
&
\sphinxAtStartPar
10
\\
\sphinxhline
\sphinxAtStartPar
\sphinxstylestrong{OBJETIVO}
&
\sphinxAtStartPar
Comprobar conocimientos adquiridos.
\\
\sphinxhline
\sphinxAtStartPar
\sphinxstylestrong{ACTIVIDAD}
&
\sphinxAtStartPar
Responda las siguientes preguntas a continuación. Cuadernos y celulares guardados. Solo se permite hoja, lapiz y borrador.
\\
\sphinxbottomrule
\end{tabulary}
\sphinxtableafterendhook\par
\sphinxattableend\end{savenotes}

\begin{sphinxadmonition}{note}{posibles preguntas}
\begin{enumerate}
\sphinxsetlistlabels{\arabic}{enumi}{enumii}{}{.}%
\item {} 
\sphinxAtStartPar
¿Qué es una red de datos y cuál fue la primera red a nivel mundial, (en que año y en que lugar se creó)?

\item {} 
\sphinxAtStartPar
¿Quién fue la persona que creo la tecnología para que funcionara el INTERNET?

\item {} 
\sphinxAtStartPar
¿En que año se creo la tecnología para que funcionara el INTERNET y que protocolos fueron usados en dicha tecnología?

\item {} 
\sphinxAtStartPar
¿En que año se crea el primer navegador y la primera página web?

\item {} 
\sphinxAtStartPar
Realice un descripción del navegador MOSAIC e indique sus diferencias con navegadores actuales

\item {} 
\sphinxAtStartPar
¿Qué es una computadora y que características debe tener para que sea considerada personal?

\item {} 
\sphinxAtStartPar
¿Cuál fue la primera computadora personal de la historia?

\item {} 
\sphinxAtStartPar
¿Qué lenguajes de programación se usaron en la primera computadora personal?

\item {} 
\sphinxAtStartPar
¿Cómo APPLE se relaciona con la historia de la computadora personal?

\item {} 
\sphinxAtStartPar
¿Cómo se democratizó el uso de las computadoras personales?

\item {} 
\sphinxAtStartPar
¿Qué características tenía la primera Macintosh y en que año fue creada?

\item {} 
\sphinxAtStartPar
¿En que año se creo el sistema operativo Microsoft DOS y por que fue tan importante?

\item {} 
\sphinxAtStartPar
Realice una descripción de tres sistemas operativos de Microsoft e indique por que se caracterizan

\item {} 
\sphinxAtStartPar
¿En que año y como operaban las computadoras en sus orígenes?

\item {} 
\sphinxAtStartPar
¿Qué es y en que año se crea FORTRAN?

\item {} 
\sphinxAtStartPar
¿Qué es un sistema operativo?

\item {} 
\sphinxAtStartPar
¿Por qué la década de los 80 fue revolucionaria en el software?

\item {} 
\sphinxAtStartPar
¿Qué es el sistema operativo LINUX y en que se diferencia con el sistema operativo UNIX?

\item {} 
\sphinxAtStartPar
Realice la descripción de tres distribuciones de sistema operativo tipo LINUX

\item {} 
\sphinxAtStartPar
¿En que se diferencian los sistemas operativos LINUX y Microsoft WINDOWS?

\end{enumerate}
\end{sphinxadmonition}


\bigskip\hrule\bigskip


\sphinxstepscope


\section{Guía 1002: OFIMÁTICA y PROGRAMACIÓN (P2)}
\label{\detokenize{1_guias/10informatica/G1002:guia-1002-ofimatica-y-programacion-p2}}\label{\detokenize{1_guias/10informatica/G1002::doc}}

\begin{savenotes}\sphinxattablestart
\sphinxthistablewithglobalstyle
\centering
\begin{tabular}[t]{\X{35}{100}\X{65}{100}}
\sphinxtoprule
\sphinxtableatstartofbodyhook

\begin{savenotes}\sphinxattablestart
\sphinxthistablewithglobalstyle
\centering
\begin{tabulary}{\linewidth}[t]{T}
\sphinxtoprule
\sphinxtableatstartofbodyhook
\sphinxAtStartPar
\sphinxincludegraphics{{escudo}.png}
\\
\sphinxbottomrule
\end{tabulary}
\sphinxtableafterendhook\par
\sphinxattableend\end{savenotes}
&

\begin{savenotes}\sphinxattablestart
\sphinxthistablewithglobalstyle
\raggedright
\begin{tabular}[t]{*{1}{\X{1}{1}}}
\sphinxtoprule
\sphinxtableatstartofbodyhook
\sphinxAtStartPar
\sphinxstylestrong{INSTITUCIÓN EDUCATIVA}: John F. Kennedy
\begin{itemize}
\item {} 
\sphinxAtStartPar
\sphinxstylestrong{DOCENTE}: Julian Camilo Fonseca Romero

\item {} 
\sphinxAtStartPar
\sphinxstylestrong{ASIGNATURA}: Tecnología e Informática

\item {} 
\sphinxAtStartPar
\sphinxstylestrong{GRADO}: DÉCIMO

\item {} 
\sphinxAtStartPar
\sphinxstylestrong{PERIODO}: 02

\item {} 
\sphinxAtStartPar
\sphinxstylestrong{TIEMPO}: 10 semanas (2 horas / semana)

\item {} 
\sphinxAtStartPar
\sphinxstylestrong{TIEMPO DE TRABAJO}: 4/20 horas

\end{itemize}
\\
\sphinxbottomrule
\end{tabular}
\sphinxtableafterendhook\par
\sphinxattableend\end{savenotes}
\\
\sphinxbottomrule
\end{tabular}
\sphinxtableafterendhook\par
\sphinxattableend\end{savenotes}

\begin{sphinxuseclass}{sd-tab-set}
\begin{sphinxuseclass}{sd-tab-item}\subsubsection*{Objetivo}

\begin{sphinxuseclass}{sd-tab-content}
\sphinxAtStartPar
\DUrole{sd-sphinx-override}{\DUrole{sd-badge}{\DUrole{sd-bg-light}{\DUrole{sd-bg-text-light}{Exploración}}}}
\DUrole{sd-sphinx-override}{\DUrole{sd-badge}{\DUrole{sd-bg-info}{\DUrole{sd-bg-text-info}{Estructuración}}}}
\DUrole{sd-sphinx-override}{\DUrole{sd-badge}{\DUrole{sd-bg-warning}{\DUrole{sd-bg-text-warning}{Actividad práctica}}}}
\DUrole{sd-sphinx-override}{\DUrole{sd-badge}{\DUrole{sd-bg-danger}{\DUrole{sd-bg-text-danger}{Transferencia}}}}
\begin{itemize}
\item {} 
\sphinxAtStartPar
\sphinxstylestrong{Temática}: Reflexión sobre la Ofimática y Programación.

\item {} 
\sphinxAtStartPar
\sphinxstylestrong{Objetivo de aprendizaje}: Identificar el nivel de interés del estudiante sobre las temáticas de Ofimática y Programación vistas en el periodo anterior.

\end{itemize}

\end{sphinxuseclass}
\end{sphinxuseclass}
\begin{sphinxuseclass}{sd-tab-item}\subsubsection*{Orientaciones curriculares}

\begin{sphinxuseclass}{sd-tab-content}\begin{itemize}
\item {} 
\sphinxAtStartPar
\sphinxstylestrong{Componente}: Uso y apropiación de la Tecnología y la Informática.

\item {} 
\sphinxAtStartPar
\sphinxstylestrong{Competencia}: Genero propuestas innovadoras para el uso y aprovechamiento de productos tecnológicos.

\item {} 
\sphinxAtStartPar
\sphinxstylestrong{Evidencia de aprendizaje}: Utilizo adecuadamente herramientas informáticas para la búsqueda, organización, procesamiento, sistematización, comunicación y difusión de ideas.

\end{itemize}

\end{sphinxuseclass}
\end{sphinxuseclass}
\begin{sphinxuseclass}{sd-tab-item}\subsubsection*{Criterios de evaluación}

\begin{sphinxuseclass}{sd-tab-content}\begin{itemize}
\item {} 
\sphinxAtStartPar
Actividades desarrolladas en su totalidad.

\item {} 
\sphinxAtStartPar
Participación en clase.

\item {} 
\sphinxAtStartPar
Explicación clara de conceptos.

\item {} 
\sphinxAtStartPar
Argumentación de ideas.

\item {} 
\sphinxAtStartPar
Cumplimiento de las reglas del aula.

\end{itemize}

\end{sphinxuseclass}
\end{sphinxuseclass}
\begin{sphinxuseclass}{sd-tab-item}\subsubsection*{Recursos (0)}

\begin{sphinxuseclass}{sd-tab-content}\begin{itemize}
\item {} 
\sphinxAtStartPar
La presente guía no contiene recursos.

\end{itemize}

\end{sphinxuseclass}
\end{sphinxuseclass}
\end{sphinxuseclass}

\subsection{Mensaje del autor}
\label{\detokenize{1_guias/10informatica/G1002:mensaje-del-autor}}\begin{quote}

\sphinxAtStartPar
“Hasta el momento, se ha explorado las temáticas de: la \sphinxstylestrong{ofimática} (que nos da el poder de organizar, analizar y presentar información de forma profesional) y la \sphinxstylestrong{programación}, que nos permite dar instrucciones a una computadora para crear soluciones desde cero y dar vida a nuestras propias ideas”.

\sphinxAtStartPar
“Esta guía no trata de un examen para calificar sus conocimientos teóricos, sino de un espacio para identificar si siente interés por las temáticas”.

\sphinxAtStartPar
“A continuación encontrará una actividad de exploración y dos actividades con preguntas de reflexión. Desarrolle las actividades con honestidad; sus respuestas le ayudarán a identificar cuál de estas dos rutas (Ofimática o Programación) le despierta más curiosidad y podría ser su temática elegida para su proyecto de grado”.

\begin{flushright}
---Camilo Fonseca
\end{flushright}
\end{quote}


\bigskip\hrule\bigskip



\subsection{A00: Recordando los temas vistos (+1,0)}
\label{\detokenize{1_guias/10informatica/G1002:a00-recordando-los-temas-vistos-1-0}}
\sphinxAtStartPar
\DUrole{sd-sphinx-override}{\DUrole{sd-badge}{\DUrole{sd-bg-light}{\DUrole{sd-bg-text-light}{Exploración}}}}
\DUrole{sd-sphinx-override}{\DUrole{sd-badge}{\DUrole{sd-bg-info}{\DUrole{sd-bg-text-info}{Estructuración}}}}
\DUrole{sd-sphinx-override}{\DUrole{sd-badge}{\DUrole{sd-bg-warning}{\DUrole{sd-bg-text-warning}{Actividad práctica}}}}
\begin{enumerate}
\sphinxsetlistlabels{\arabic}{enumi}{enumii}{}{.}%
\item {} 
\sphinxAtStartPar
Transcriba en su cuaderno el objetivo, las orientaciones curriculares y los criterios de evaluación de la presente guía.

\item {} 
\sphinxAtStartPar
Lea el mensaje del autor y en su cuaderno responda la siguiente pregunta:
\begin{itemize}
\item {} 
\sphinxAtStartPar
¿Que tienen que ver la \sphinxstylestrong{ofimática} y la \sphinxstylestrong{programación} en el estudio de la informática?

\end{itemize}

\item {} 
\sphinxAtStartPar
Acceda a las guías de \sphinxstylestrong{Ofimática} y \sphinxstylestrong{Programación} que trabajó anteriormente para recordar conceptos.
\begin{itemize}
\item {} 
\sphinxAtStartPar
{\hyperref[\detokenize{1_guias/10ofimatica/G1011::doc}]{\sphinxcrossref{\DUrole{std}{\DUrole{std-doc}{Guía Ofimática 01}}}}}

\item {} 
\sphinxAtStartPar
{\hyperref[\detokenize{1_guias/10ofimatica/G1012::doc}]{\sphinxcrossref{\DUrole{std}{\DUrole{std-doc}{Guía Ofimática 02}}}}}

\item {} 
\sphinxAtStartPar
{\hyperref[\detokenize{1_guias/10programacion/G1021::doc}]{\sphinxcrossref{\DUrole{std}{\DUrole{std-doc}{Guía Programación 01}}}}}

\item {} 
\sphinxAtStartPar
{\hyperref[\detokenize{1_guias/10programacion/G1022::doc}]{\sphinxcrossref{\DUrole{std}{\DUrole{std-doc}{Guía Programación 02}}}}}

\end{itemize}

\item {} 
\sphinxAtStartPar
Con sus palabras, realizar un resumen de cada guía en su cuaderno. (Una página de cuaderno por cada guía).

\end{enumerate}


\bigskip\hrule\bigskip



\subsection{A01: Análisis y reflexión (OFIMÁTICA) (+2,0)}
\label{\detokenize{1_guias/10informatica/G1002:a01-analisis-y-reflexion-ofimatica-2-0}}
\sphinxAtStartPar
\DUrole{sd-sphinx-override}{\DUrole{sd-badge}{\DUrole{sd-bg-danger}{\DUrole{sd-bg-text-danger}{Transferencia}}}}


\begin{savenotes}\sphinxattablestart
\sphinxthistablewithglobalstyle
\centering
\begin{tabulary}{\linewidth}[t]{TT}
\sphinxtoprule
\sphinxtableatstartofbodyhook
\sphinxAtStartPar
\sphinxstylestrong{TIEMPO}
&
\sphinxAtStartPar
59 minutos
\\
\sphinxhline
\sphinxAtStartPar
\sphinxstylestrong{PREGUNTAS}
&
\sphinxAtStartPar
10
\\
\sphinxhline
\sphinxAtStartPar
\sphinxstylestrong{OBJETIVO}
&
\sphinxAtStartPar
Identificar qué le llamó la atención de la temática de \sphinxstylestrong{OFIMÁTICA}.
\\
\sphinxhline
\sphinxAtStartPar
\sphinxstylestrong{ACTIVIDAD}
&
\sphinxAtStartPar
En una hoja al final de su cuaderno, desarrolle un \sphinxstylestrong{texto reflexivo} utilizando las siguientes preguntas como guía para estructurar su escrito.
\\
\sphinxbottomrule
\end{tabulary}
\sphinxtableafterendhook\par
\sphinxattableend\end{savenotes}

\begin{sphinxadmonition}{note}{Preguntas de reflexión}
\begin{enumerate}
\sphinxsetlistlabels{\arabic}{enumi}{enumii}{}{.}%
\item {} 
\sphinxAtStartPar
¿Qué aspecto de la ofimática le parece más interesante o útil para su vida académica o personal?

\end{enumerate}


\bigskip\hrule\bigskip

\begin{enumerate}
\sphinxsetlistlabels{\arabic}{enumi}{enumii}{}{.}%
\setcounter{enumi}{1}
\item {} 
\sphinxAtStartPar
¿Le gustaría aprender más sobre alguna herramienta específica de ofimática, como Google Workspace o Microsoft 365? ¿Cuál y por qué?

\end{enumerate}


\bigskip\hrule\bigskip

\begin{enumerate}
\sphinxsetlistlabels{\arabic}{enumi}{enumii}{}{.}%
\setcounter{enumi}{2}
\item {} 
\sphinxAtStartPar
¿Cree que el manejo de herramientas ofimáticas puede facilitar su trabajo diario o sus estudios? ¿Qué le motiva o desmotiva de esta temática?

\end{enumerate}


\bigskip\hrule\bigskip

\begin{enumerate}
\sphinxsetlistlabels{\arabic}{enumi}{enumii}{}{.}%
\setcounter{enumi}{3}
\item {} 
\sphinxAtStartPar
¿Le interesa conocer y usar programas cuyo código fuente está disponible para que cualquiera pueda estudiarlo y modificarlo? (¿por qué?)

\end{enumerate}


\bigskip\hrule\bigskip

\begin{enumerate}
\sphinxsetlistlabels{\arabic}{enumi}{enumii}{}{.}%
\setcounter{enumi}{4}
\item {} 
\sphinxAtStartPar
¿Ha utilizado alguna vez software libre como LibreOffice, GIMP o Linux? ¿Qué le pareció la experiencia?

\end{enumerate}


\bigskip\hrule\bigskip

\begin{enumerate}
\sphinxsetlistlabels{\arabic}{enumi}{enumii}{}{.}%
\setcounter{enumi}{5}
\item {} 
\sphinxAtStartPar
Cree que es importante que el software sea accesible y modificable por todos, sin restricciones de licencia propietaria? (¿por qué?)

\end{enumerate}


\bigskip\hrule\bigskip

\begin{enumerate}
\sphinxsetlistlabels{\arabic}{enumi}{enumii}{}{.}%
\setcounter{enumi}{6}
\item {} 
\sphinxAtStartPar
¿Le gustaría aprender a contribuir con proyectos de software libre, como mejorar programas o crear nuevas herramientas? (¿por qué?)

\end{enumerate}


\bigskip\hrule\bigskip

\begin{enumerate}
\sphinxsetlistlabels{\arabic}{enumi}{enumii}{}{.}%
\setcounter{enumi}{7}
\item {} 
\sphinxAtStartPar
¿Qué opina sobre la idea de compartir y colaborar en el desarrollo de software en lugar de usar programas cerrados y comerciales?

\end{enumerate}


\bigskip\hrule\bigskip

\begin{enumerate}
\sphinxsetlistlabels{\arabic}{enumi}{enumii}{}{.}%
\setcounter{enumi}{8}
\item {} 
\sphinxAtStartPar
¿Le atrae la posibilidad de personalizar el software que usa para adaptarlo a sus necesidades específicas? (¿por qué?)

\end{enumerate}


\bigskip\hrule\bigskip

\begin{enumerate}
\sphinxsetlistlabels{\arabic}{enumi}{enumii}{}{.}%
\setcounter{enumi}{9}
\item {} 
\sphinxAtStartPar
¿Considera que el software libre puede ser una alternativa viable y segura frente al software propietario? (¿por qué?)

\end{enumerate}
\end{sphinxadmonition}


\bigskip\hrule\bigskip



\subsection{A02: Análisis y reflexión (PROGRAMACIÓN) (+2,0)}
\label{\detokenize{1_guias/10informatica/G1002:a02-analisis-y-reflexion-programacion-2-0}}
\sphinxAtStartPar
\DUrole{sd-sphinx-override}{\DUrole{sd-badge}{\DUrole{sd-bg-danger}{\DUrole{sd-bg-text-danger}{Transferencia}}}}


\begin{savenotes}\sphinxattablestart
\sphinxthistablewithglobalstyle
\centering
\begin{tabulary}{\linewidth}[t]{TT}
\sphinxtoprule
\sphinxtableatstartofbodyhook
\sphinxAtStartPar
\sphinxstylestrong{TIEMPO}
&
\sphinxAtStartPar
59 minutos
\\
\sphinxhline
\sphinxAtStartPar
\sphinxstylestrong{PREGUNTAS}
&
\sphinxAtStartPar
10
\\
\sphinxhline
\sphinxAtStartPar
\sphinxstylestrong{OBJETIVO}
&
\sphinxAtStartPar
Identificar qué le llamó la atención de la temática de \sphinxstylestrong{PROGRAMACIÓN}.
\\
\sphinxhline
\sphinxAtStartPar
\sphinxstylestrong{ACTIVIDAD}
&
\sphinxAtStartPar
En una hoja al final de su cuaderno, desarrolle un \sphinxstylestrong{texto reflexivo} utilizando las siguientes preguntas como guía para estructurar su escrito.
\\
\sphinxbottomrule
\end{tabulary}
\sphinxtableafterendhook\par
\sphinxattableend\end{savenotes}

\begin{sphinxadmonition}{note}{Preguntas de reflexión}
\begin{enumerate}
\sphinxsetlistlabels{\arabic}{enumi}{enumii}{}{.}%
\item {} 
\sphinxAtStartPar
¿Se siente motivado al pensar que, mediante la creación de código, se pueden dar instrucciones exactas a una computadora para que resuelva un problema complejo por usted? (¿por qué?)

\end{enumerate}


\bigskip\hrule\bigskip

\begin{enumerate}
\sphinxsetlistlabels{\arabic}{enumi}{enumii}{}{.}%
\setcounter{enumi}{1}
\item {} 
\sphinxAtStartPar
Al conocer la diferencia entre programador y desarrollador, ¿le atrae más la parte técnica de “escribir código” para construir funcionalidades específicas, o le interesa más la visión global de “planificar y gestionar” todo el ciclo de vida de un software? (¿por qué?)

\end{enumerate}


\bigskip\hrule\bigskip

\begin{enumerate}
\sphinxsetlistlabels{\arabic}{enumi}{enumii}{}{.}%
\setcounter{enumi}{2}
\item {} 
\sphinxAtStartPar
Sobre el software y la innovación: Al entender que existen diferentes tipos de software (sistema, aplicación, embebido, programación), ¿hay alguna categoría en particular que le genere curiosidad sobre cómo se construyen o cómo funcionan por dentro? (¿por qué?)

\end{enumerate}


\bigskip\hrule\bigskip

\begin{enumerate}
\sphinxsetlistlabels{\arabic}{enumi}{enumii}{}{.}%
\setcounter{enumi}{3}
\item {} 
\sphinxAtStartPar
De las tecnologías mencionadas (bases de datos, servidores web, entornos como Node.js o Frameworks), ¿hay alguna que le entusiasme aprender a usar para crear tu propio proyecto? (¿por qué?)

\end{enumerate}


\bigskip\hrule\bigskip

\begin{enumerate}
\sphinxsetlistlabels{\arabic}{enumi}{enumii}{}{.}%
\setcounter{enumi}{4}
\item {} 
\sphinxAtStartPar
¿Qué parte del desarrollo web le parece más interesante: el Front\sphinxhyphen{}end (crear la parte visual con la que interactúan los usuarios) o el Back\sphinxhyphen{}end (trabajar con la lógica, servidores y bases de datos que hacen que todo funcione)? (¿por qué?)

\end{enumerate}


\bigskip\hrule\bigskip

\begin{enumerate}
\sphinxsetlistlabels{\arabic}{enumi}{enumii}{}{.}%
\setcounter{enumi}{5}
\item {} 
\sphinxAtStartPar
De los lenguajes mencionados (HTML, CSS, JavaScript, PHP, Java, Python), ¿alguno te llama especialmente la atención para empezar a crear sus primeros programas o sitios web? (¿por qué?)

\end{enumerate}


\bigskip\hrule\bigskip

\begin{enumerate}
\sphinxsetlistlabels{\arabic}{enumi}{enumii}{}{.}%
\setcounter{enumi}{6}
\item {} 
\sphinxAtStartPar
Considerando que el desarrollo de software implica análisis, diseño, pruebas y mantenimiento, ¿le atrae la idea de enfrentar el reto de identificar errores y optimizar código para mejorar continuamente una aplicación? (¿por qué?)

\end{enumerate}


\bigskip\hrule\bigskip

\begin{enumerate}
\sphinxsetlistlabels{\arabic}{enumi}{enumii}{}{.}%
\setcounter{enumi}{7}
\item {} 
\sphinxAtStartPar
Después de analizar cómo los datos simples se transforman en información útil para la toma de decisiones, ¿considera que comprender la jerarquía de datos cambia tu perspectiva sobre cómo la tecnología influye en tu vida cotidiana?. ¿Tendría interés en profundizar cómo se organiza la información en sistemas complejos?

\end{enumerate}


\bigskip\hrule\bigskip

\begin{enumerate}
\sphinxsetlistlabels{\arabic}{enumi}{enumii}{}{.}%
\setcounter{enumi}{8}
\item {} 
\sphinxAtStartPar
Después de conocer que los diagramas de flujo nos permiten visualizar procesos paso a paso para resolver problemas de manera organizada, ¿se siente cómodo(a) utilizando esta herramienta gráfica para planear sus propias soluciones, o prefiere otro método para organizar sus ideas? (¿por qué?)

\end{enumerate}


\bigskip\hrule\bigskip

\begin{enumerate}
\sphinxsetlistlabels{\arabic}{enumi}{enumii}{}{.}%
\setcounter{enumi}{9}
\item {} 
\sphinxAtStartPar
Considerando que los algoritmos nos permiten descomponer cualquier situación de la vida diaria en una secuencia de pasos claros para resolverla, ¿qué tan interesante le resulta el desafío de diseñar instrucciones precisas para que otros (o una computadora) puedan completar una tarea siguiendo su lógica, o prefiere otro tipo de desafíos intelectuales? (¿por qué?)

\end{enumerate}
\end{sphinxadmonition}

\sphinxstepscope


\section{Guía 1003: REDES (P2)}
\label{\detokenize{1_guias/10informatica/G1003:guia-1003-redes-p2}}\label{\detokenize{1_guias/10informatica/G1003::doc}}

\begin{savenotes}\sphinxattablestart
\sphinxthistablewithglobalstyle
\centering
\begin{tabular}[t]{\X{35}{100}\X{65}{100}}
\sphinxtoprule
\sphinxtableatstartofbodyhook

\begin{savenotes}\sphinxattablestart
\sphinxthistablewithglobalstyle
\centering
\begin{tabulary}{\linewidth}[t]{T}
\sphinxtoprule
\sphinxtableatstartofbodyhook
\sphinxAtStartPar
\sphinxincludegraphics{{escudo}.png}
\\
\sphinxbottomrule
\end{tabulary}
\sphinxtableafterendhook\par
\sphinxattableend\end{savenotes}
&

\begin{savenotes}\sphinxattablestart
\sphinxthistablewithglobalstyle
\raggedright
\begin{tabular}[t]{*{1}{\X{1}{1}}}
\sphinxtoprule
\sphinxtableatstartofbodyhook
\sphinxAtStartPar
\sphinxstylestrong{INSTITUCIÓN EDUCATIVA}: John F. Kennedy
\begin{itemize}
\item {} 
\sphinxAtStartPar
\sphinxstylestrong{DOCENTE}: Julian Camilo Fonseca Romero

\item {} 
\sphinxAtStartPar
\sphinxstylestrong{ASIGNATURA}: Tecnología e Informática

\item {} 
\sphinxAtStartPar
\sphinxstylestrong{GRADO}: DÉCIMO

\item {} 
\sphinxAtStartPar
\sphinxstylestrong{PERIODO}: 02

\item {} 
\sphinxAtStartPar
\sphinxstylestrong{TIEMPO}: 10 semanas (2 horas / semana)

\item {} 
\sphinxAtStartPar
\sphinxstylestrong{TIEMPO DE TRABAJO}: 8/20 horas

\end{itemize}
\\
\sphinxbottomrule
\end{tabular}
\sphinxtableafterendhook\par
\sphinxattableend\end{savenotes}
\\
\sphinxbottomrule
\end{tabular}
\sphinxtableafterendhook\par
\sphinxattableend\end{savenotes}

\begin{sphinxuseclass}{sd-tab-set}
\begin{sphinxuseclass}{sd-tab-item}\subsubsection*{Objetivo}

\begin{sphinxuseclass}{sd-tab-content}
\sphinxAtStartPar
\DUrole{sd-sphinx-override}{\DUrole{sd-badge}{\DUrole{sd-bg-light}{\DUrole{sd-bg-text-light}{Exploración}}}}
\DUrole{sd-sphinx-override}{\DUrole{sd-badge}{\DUrole{sd-bg-info}{\DUrole{sd-bg-text-info}{Estructuración}}}}
\DUrole{sd-sphinx-override}{\DUrole{sd-badge}{\DUrole{sd-bg-warning}{\DUrole{sd-bg-text-warning}{Actividad práctica}}}}
\DUrole{sd-sphinx-override}{\DUrole{sd-badge}{\DUrole{sd-bg-danger}{\DUrole{sd-bg-text-danger}{Transferencia}}}}
\begin{itemize}
\item {} 
\sphinxAtStartPar
\sphinxstylestrong{Temática}: Redes de datos.

\item {} 
\sphinxAtStartPar
\sphinxstylestrong{Objetivo de aprendizaje}: Introducir los conceptos relacionados con la temática de redes de datos.

\end{itemize}

\end{sphinxuseclass}
\end{sphinxuseclass}
\begin{sphinxuseclass}{sd-tab-item}\subsubsection*{Orientaciones curriculares}

\begin{sphinxuseclass}{sd-tab-content}\begin{itemize}
\item {} 
\sphinxAtStartPar
\sphinxstylestrong{Componente}: Uso y apropiación de la Tecnología y la Informática.

\item {} 
\sphinxAtStartPar
\sphinxstylestrong{Competencia}: Genero propuestas innovadoras para el uso y aprovechamiento de productos tecnológicos.

\item {} 
\sphinxAtStartPar
\sphinxstylestrong{Evidencia de aprendizaje}: Diseño nuevos escenarios de uso, adaptación o transformación de tecnologías emergentes y componentes de hardware y software, en contextos específicos de mi entorno, que favorezcan la vida, la productividad o el bienestar.

\end{itemize}

\end{sphinxuseclass}
\end{sphinxuseclass}
\begin{sphinxuseclass}{sd-tab-item}\subsubsection*{Criterios de evaluación}

\begin{sphinxuseclass}{sd-tab-content}\begin{itemize}
\item {} 
\sphinxAtStartPar
Actividades desarrolladas en su totalidad.

\item {} 
\sphinxAtStartPar
Participación en clase.

\item {} 
\sphinxAtStartPar
Explicación clara de conceptos.

\item {} 
\sphinxAtStartPar
Argumentación de ideas.

\item {} 
\sphinxAtStartPar
Cumplimiento de las reglas del aula.

\end{itemize}

\end{sphinxuseclass}
\end{sphinxuseclass}
\begin{sphinxuseclass}{sd-tab-item}\subsubsection*{Recursos (0)}

\begin{sphinxuseclass}{sd-tab-content}\begin{itemize}
\item {} 
\sphinxAtStartPar
Esta guía no contiene recursos

\end{itemize}

\end{sphinxuseclass}
\end{sphinxuseclass}
\end{sphinxuseclass}

\subsection{Mensaje del autor}
\label{\detokenize{1_guias/10informatica/G1003:mensaje-del-autor}}\begin{quote}

\sphinxAtStartPar
“Esta guía da a conocer las características de las redes de datos y la ciberseguridad. El objetivo principal es lograr comprender la infraestructura que sostiene el mundo digital, desde el funcionamiento de un simple router en casa hasta la compleja interconexión de servidores a nivel global. Además, es fundamental que el estudiante desarrolle una visión crítica sobre la seguridad de la información, el hacking ético y las implicaciones legales de los delitos informáticos en Colombia”.

\sphinxAtStartPar
“Al final de la guía encontrará unas preguntas de reflexión para identificar si el tema de \sphinxstylestrong{redes} o \sphinxstylestrong{ciberseguridad} le genera interés. Alguna de estas temáticas podría ser su elección para su proyecto de grado”.

\begin{flushright}
---Camilo Fonseca
\end{flushright}
\end{quote}


\bigskip\hrule\bigskip



\subsection{A00: Introducción (+0,1)}
\label{\detokenize{1_guias/10informatica/G1003:a00-introduccion-0-1}}
\sphinxAtStartPar
\DUrole{sd-sphinx-override}{\DUrole{sd-badge}{\DUrole{sd-bg-light}{\DUrole{sd-bg-text-light}{Exploración}}}}
\begin{enumerate}
\sphinxsetlistlabels{\arabic}{enumi}{enumii}{}{.}%
\item {} 
\sphinxAtStartPar
Transcriba en su cuaderno el objetivo, las orientaciones curriculares y los criterios de evaluación de la presente guía.

\item {} 
\sphinxAtStartPar
Lea el mensaje del autor y responda en su cuaderno la siguiente pregunta:
\begin{itemize}
\item {} 
\sphinxAtStartPar
¿Qué tienen que ver las \sphinxstylestrong{redes de datos} y la \sphinxstylestrong{ciberseguridad} con el estudio de la informática?

\end{itemize}

\end{enumerate}


\bigskip\hrule\bigskip



\subsection{A01: Redes de datos, Internet y Ciberseguridad (+0,1)}
\label{\detokenize{1_guias/10informatica/G1003:a01-redes-de-datos-internet-y-ciberseguridad-0-1}}
\sphinxAtStartPar
\DUrole{sd-sphinx-override}{\DUrole{sd-badge}{\DUrole{sd-bg-info}{\DUrole{sd-bg-text-info}{Estructuración}}}}
\begin{enumerate}
\sphinxsetlistlabels{\arabic}{enumi}{enumii}{}{.}%
\item {} 
\sphinxAtStartPar
Lea los siguientes conceptos relacionados con redes de datos e Internet.

\end{enumerate}


\bigskip\hrule\bigskip



\subsubsection{¿Qué son las Redes de datos?}
\label{\detokenize{1_guias/10informatica/G1003:que-son-las-redes-de-datos}}
\begin{sphinxadmonition}{note}{Concepto de redes de datos}

\sphinxAtStartPar
Una \sphinxstylestrong{red de datos} es un sistema compuesto de dispositivos electrónicos y equipos de comunicación que permiten el intercambio de información en diferentes formatos a través de una red.

\sphinxAtStartPar
Las organizaciones y empresas dependen de la infraestructura de red para realizar muchas de sus tareas y procesos de negocio.


\begin{savenotes}\sphinxattablestart
\sphinxthistablewithglobalstyle
\centering
\begin{tabulary}{\linewidth}[t]{T}
\sphinxtoprule
\sphinxtableatstartofbodyhook

\begin{sphinxfigure-in-table}
\centering
\capstart
\noindent\sphinxincludegraphics[width=400\sphinxpxdimen]{{image013}.png}
\sphinxfigcaption{Dispositivo de red}\label{\detokenize{1_guias/10informatica/G1003:id1}}\end{sphinxfigure-in-table}\relax
\\
\sphinxbottomrule
\end{tabulary}
\sphinxtableafterendhook\par
\sphinxattableend\end{savenotes}
\end{sphinxadmonition}


\bigskip\hrule\bigskip



\subsubsection{Tipos de redes}
\label{\detokenize{1_guias/10informatica/G1003:tipos-de-redes}}
\sphinxAtStartPar
Los tipos de redes son:

\begin{sphinxadmonition}{note}{Red de área local (LAN)}
\begin{quote}

\sphinxAtStartPar
Una red de área local es un tipo de red de telecomunicaciones que conecta dispositivos electrónicos en un área geográfica relativamente pequeña, como un edificio, una oficina, una casa o una sala de informática. Estas redes están diseñadas para facilitar la comunicación y el intercambio de datos entre dispositivos que se encuentran en proximidad física.
\end{quote}
\end{sphinxadmonition}

\begin{sphinxadmonition}{note}{Red de área metropolitana (MAN)}
\begin{quote}

\sphinxAtStartPar
Una red de área metropolitana es una red de telecomunicaciones que conecta diferentes ubicaciones dentro de una ciudad o área metropolitana, situándose como una solución intermedia entre las redes LAN y las redes WAN.
\end{quote}
\end{sphinxadmonition}

\begin{sphinxadmonition}{note}{Red de área extensa (WAN)}
\begin{quote}

\sphinxAtStartPar
Las redes de área extensa, son redes de telecomunicaciones que cubren áreas geográficas muy amplias, incluso llegando a interconectar ubicaciones que se encuentran a mucha distancia o a nivel global. Las redes WAN se emplean en entornos empresariales para conectar sucursales, oficinas o puntos de venta en la industria, para supervisar y controlar dispositivos a distancia, y en aplicaciones gubernamentales y de investigación para la transmisión de datos entre ubicaciones distantes.
\end{quote}
\end{sphinxadmonition}


\bigskip\hrule\bigskip



\subsubsection{Clasificación de redes}
\label{\detokenize{1_guias/10informatica/G1003:clasificacion-de-redes}}
\sphinxAtStartPar
Las redes se clasifican en:

\begin{sphinxadmonition}{note}{Redes alámbricas (USO CABLES)}

\sphinxAtStartPar
Las \sphinxstylestrong{redes alámbricas} utilizan tecnologías de comunicación basadas en cables físicos (como cables de par trenzado, coaxiales o fibra óptica) para transmitir datos de un punto a otro, dependiendo de conexiones estructuradas y tangibles para el transporte de la información.

\sphinxAtStartPar
Este tipo de redes de telecomunicaciones ofrecen una mayor estabilidad, velocidad y seguridad en la transmisión de datos, ya que la señal no está sujeta a interferencias ambientales ni a obstáculos físicos. Son redes ideales para conectar dispositivos que requieren un alto rendimiento y están situados en una zona fija, como ordenadores de escritorio, servidores, consolas de videojuegos o sistemas de almacenamiento en red (NAS).

\sphinxAtStartPar
Entre las distintas clases de \sphinxstyleemphasis{redes alámbricas} se tienen:
\begin{itemize}
\item {} 
\sphinxAtStartPar
Basadas en tecnología Ethernet: redes locales domésticas o empresariales (LAN) que utilizan cables de cobre (UTP/STP).

\item {} 
\sphinxAtStartPar
Usando tecnologías de alta velocidad: redes de fibra óptica, ideales para conexiones troncales o internet de banda ancha simétrica.

\item {} 
\sphinxAtStartPar
Redes de infraestructura dedicada: conexiones de redes de área amplia (WAN) o líneas telefónicas tradicionales (como las antiguas conexiones ADSL).

\end{itemize}
\end{sphinxadmonition}

\begin{sphinxadmonition}{note}{Redes inalámbricas (SIN CLABLES) (USO DE ONDAS)}

\sphinxAtStartPar
Las \sphinxstylestrong{redes inalámbricas} utilizan tecnologías de comunicación basadas en señales de \sphinxstylestrong{radiofrecuencia} o \sphinxstylestrong{microondas} para transmitir datos a través del aire, eliminando la necesidad de conexiones físicas por medio de cables.

\sphinxAtStartPar
Este tipo de redes de telecomunicaciones ofrecen una mayor flexibilidad y movilidad, pues no es necesario que los dispositivos estén situados en una zona fija. Son redes ideales para conectar dispositivos móviles como smartphones, tablets, cámaras de video\sphinxhyphen{}vigilancia u ordenadores portátiles, entre otros.

\sphinxAtStartPar
Entre las distintas clases de \sphinxstyleemphasis{redes inalámbricas} se tienen:
\begin{itemize}
\item {} 
\sphinxAtStartPar
Basadas en tecnologías de radiofrecuencia: \sphinxstylestrong{redes wifi}.

\item {} 
\sphinxAtStartPar
Usando tecnologías como 3G, 4G o 5G: redes de \sphinxstylestrong{telefonía móvil}.

\item {} 
\sphinxAtStartPar
Redes ad hoc (se crean automáticamente entre dispositivos cercanos utilizando tecnologías como \sphinxstylestrong{bluetooth} o \sphinxstylestrong{NFC}).

\end{itemize}
\end{sphinxadmonition}


\bigskip\hrule\bigskip



\subsubsection{¿Qué es una Intranet?}
\label{\detokenize{1_guias/10informatica/G1003:que-es-una-intranet}}
\begin{sphinxadmonition}{note}{Concepto de Intranet}

\sphinxAtStartPar
Una \sphinxstylestrong{Intranet} es una \sphinxstylestrong{red privada local} limitada a una organización. Sus características principales son:
\begin{itemize}
\item {} 
\sphinxAtStartPar
Acceso restringido solo a miembros de la organización

\item {} 
\sphinxAtStartPar
Mayor control y seguridad

\item {} 
\sphinxAtStartPar
Optimizada para necesidades específicas

\item {} 
\sphinxAtStartPar
Compartición de recursos internos

\item {} 
\sphinxAtStartPar
Facilita la colaboración entre empleados

\item {} 
\sphinxAtStartPar
Usa protocolos TCP/IP

\end{itemize}


\begin{savenotes}\sphinxattablestart
\sphinxthistablewithglobalstyle
\centering
\begin{tabulary}{\linewidth}[t]{T}
\sphinxtoprule
\sphinxtableatstartofbodyhook

\begin{sphinxfigure-in-table}
\centering
\capstart
\noindent\sphinxincludegraphics[width=360\sphinxpxdimen]{{image021}.png}
\sphinxfigcaption{Dispositivos de una Intranet}\label{\detokenize{1_guias/10informatica/G1003:id2}}\end{sphinxfigure-in-table}\relax
\\
\sphinxbottomrule
\end{tabulary}
\sphinxtableafterendhook\par
\sphinxattableend\end{savenotes}
\end{sphinxadmonition}


\bigskip\hrule\bigskip



\subsubsection{¿Qué es INTERNET?}
\label{\detokenize{1_guias/10informatica/G1003:que-es-internet}}
\begin{sphinxadmonition}{note}{Concepto de Internet}

\sphinxAtStartPar
\sphinxstylestrong{INTERNET} es una \sphinxstylestrong{red pública global} que conecta Intranets en todo el mundo. Sus principales características son:
\begin{itemize}
\item {} 
\sphinxAtStartPar
Acceso público y global

\item {} 
\sphinxAtStartPar
Conecta millones de redes diferentes

\item {} 
\sphinxAtStartPar
Requiere medidas de seguridad robustas

\item {} 
\sphinxAtStartPar
Permite acceder a servicios y recursos públicos

\item {} 
\sphinxAtStartPar
Usa protocolos TCP/IP

\end{itemize}


\begin{savenotes}\sphinxattablestart
\sphinxthistablewithglobalstyle
\centering
\begin{tabulary}{\linewidth}[t]{T}
\sphinxtoprule
\sphinxtableatstartofbodyhook

\begin{sphinxfigure-in-table}
\centering
\capstart
\noindent\sphinxincludegraphics[width=270\sphinxpxdimen]{{image031}.png}
\sphinxfigcaption{INTERNET}\label{\detokenize{1_guias/10informatica/G1003:id3}}\end{sphinxfigure-in-table}\relax
\\
\sphinxbottomrule
\end{tabulary}
\sphinxtableafterendhook\par
\sphinxattableend\end{savenotes}
\end{sphinxadmonition}


\bigskip\hrule\bigskip



\subsubsection{Componentes de una red}
\label{\detokenize{1_guias/10informatica/G1003:componentes-de-una-red}}
\begin{sphinxadmonition}{note}{\sphinxstylestrong{SERVIDORES}}

\sphinxAtStartPar
Un servidor se trata de una computadora o aplicación encargada de proveer a otros equipos de unos determinados servicios, como el almacenamiento de información.

\sphinxAtStartPar
La mayoría de servidores ni siquiera tienen periféricos como el monitor, teclado o mouse. Para acceder y trabajar en estos, se suele trabajar de forma remota por línea de comandos.
\begin{quote}

\sphinxAtStartPar
Los servidores suelen ser \sphinxstylestrong{equipos} que \sphinxstylestrong{muy rara vez se apagan} y una vez configurados quedan trabajando con poca intervención del usuario. Solo se intervienen cuando sea necesario.
\end{quote}

\sphinxAtStartPar
Los servidores se encargan de procesar las solicitudes de esos otros equipos, que en ocasiones se denominan clientes, a los que les entregan datos por medio de una red que puede ser local o por Internet.


\begin{savenotes}\sphinxattablestart
\sphinxthistablewithglobalstyle
\centering
\begin{tabulary}{\linewidth}[t]{T}
\sphinxtoprule
\sphinxtableatstartofbodyhook

\begin{sphinxfigure-in-table}
\centering
\capstart
\noindent\sphinxincludegraphics[width=400\sphinxpxdimen]{{image041}.png}
\sphinxfigcaption{SERVIDORES}\label{\detokenize{1_guias/10informatica/G1003:id4}}\end{sphinxfigure-in-table}\relax
\\
\sphinxbottomrule
\end{tabulary}
\sphinxtableafterendhook\par
\sphinxattableend\end{savenotes}
\end{sphinxadmonition}

\begin{sphinxadmonition}{note}{\sphinxstylestrong{Equipos cliente}}

\sphinxAtStartPar
Un equipo cliente es un dispositivo destinado a realizar una determinada labor profesional, técnica o científica. En resumen, en una oficina, los equipos cliente serían cada uno de los ordenadores con los que se lleva a cabo el trabajo y que generalmente están conectados entre sí a través de un servidor para facilitar los flujos de información, conectadas a periféricos (como impresoras, escáneres, pantallas de proyección, etc.) y a la vez conectadas a Internet.


\begin{savenotes}\sphinxattablestart
\sphinxthistablewithglobalstyle
\centering
\begin{tabulary}{\linewidth}[t]{T}
\sphinxtoprule
\sphinxtableatstartofbodyhook

\begin{sphinxfigure-in-table}
\centering
\capstart
\noindent\sphinxincludegraphics[width=550\sphinxpxdimen]{{image051}.png}
\sphinxfigcaption{Equipos cliente}\label{\detokenize{1_guias/10informatica/G1003:id5}}\end{sphinxfigure-in-table}\relax
\\
\sphinxbottomrule
\end{tabulary}
\sphinxtableafterendhook\par
\sphinxattableend\end{savenotes}
\end{sphinxadmonition}

\begin{sphinxadmonition}{note}{\sphinxstylestrong{Switch}: Conecta diferentes dispositivos de forma alámbrica}

\sphinxAtStartPar
Los switches son dispositivos que se encargan de interconectar dos o más segmentos de red. Permiten formar lo que se conoce como una red de área local (LAN, Local Area Network), cuyas especificaciones técnicas siguen el estándar denominado Ethernet.


\begin{savenotes}\sphinxattablestart
\sphinxthistablewithglobalstyle
\centering
\begin{tabulary}{\linewidth}[t]{T}
\sphinxtoprule
\sphinxtableatstartofbodyhook

\begin{sphinxfigure-in-table}
\centering
\capstart
\noindent\sphinxincludegraphics[width=400\sphinxpxdimen]{{image061}.png}
\sphinxfigcaption{SWITCH}\label{\detokenize{1_guias/10informatica/G1003:id6}}\end{sphinxfigure-in-table}\relax
\\
\sphinxbottomrule
\end{tabulary}
\sphinxtableafterendhook\par
\sphinxattableend\end{savenotes}
\end{sphinxadmonition}

\begin{sphinxadmonition}{note}{\sphinxstylestrong{Access Point}: Conecta diferentes dispositivos de forma inalámbrica}

\sphinxAtStartPar
Los puntos de acceso (ACCESS POINT) son esenciales en las redes inalámbricas, ya que, ayudan a ampliar el alcance y la capacidad de conectividad. Recibe la señal de acceso a la red del router y aumenta su potencia cubriendo un área determinada. A diferencia de un repetidor, crea un nueva señal independiente en lugar de amplificarla.


\begin{savenotes}\sphinxattablestart
\sphinxthistablewithglobalstyle
\centering
\begin{tabulary}{\linewidth}[t]{T}
\sphinxtoprule
\sphinxtableatstartofbodyhook

\begin{sphinxfigure-in-table}
\centering
\capstart
\noindent\sphinxincludegraphics[width=400\sphinxpxdimen]{{image071}.png}
\sphinxfigcaption{Access Point}\label{\detokenize{1_guias/10informatica/G1003:id7}}\end{sphinxfigure-in-table}\relax
\\
\sphinxbottomrule
\end{tabulary}
\sphinxtableafterendhook\par
\sphinxattableend\end{savenotes}
\end{sphinxadmonition}

\begin{sphinxadmonition}{note}{\sphinxstylestrong{ROUTER}: Conecta diferentes redes a INTERNET}

\sphinxAtStartPar
También se le denomina enrutador, ya que se encarga de encontrar el mejor camino para la transmisión de información a través de una red. Sus usos pueden ser más o menos complejos aunque el más común es aquel que permite comunicar a varios equipos, ya sea en casa o en una oficina (aprovechando la misma conexión a Internet).

\sphinxAtStartPar
El router recibe la conexión de red y la distribuye a todos los dispositivos conectados simultáneamente escogiendo la mejor vía para hacerlo de forma que la información completa llegue de manera adecuada y en el menor tiempo posible.


\begin{savenotes}\sphinxattablestart
\sphinxthistablewithglobalstyle
\centering
\begin{tabulary}{\linewidth}[t]{T}
\sphinxtoprule
\sphinxtableatstartofbodyhook

\begin{sphinxfigure-in-table}
\centering
\capstart
\noindent\sphinxincludegraphics[width=400\sphinxpxdimen]{{image081}.png}
\sphinxfigcaption{ROUTER: Dispositivo para salir a INTERNET}\label{\detokenize{1_guias/10informatica/G1003:id8}}\end{sphinxfigure-in-table}\relax
\\
\sphinxbottomrule
\end{tabulary}
\sphinxtableafterendhook\par
\sphinxattableend\end{savenotes}
\end{sphinxadmonition}

\begin{sphinxadmonition}{note}{MODEM: (Switch + Access Point + Router)}

\sphinxAtStartPar
Algunos dispositivos de red son dispositivos todo en uno. Es el caso de los modem que son un Switch, un Access Point y un Router a la vez. Suelen ser facilitados por los prestadores de servicio de Internet.

\sphinxAtStartPar
En empresas grandes donde están conectados varios edificios como las universidades, lo normal es que hayan varios dispositivos de red cumpliendo una única función (Switch, Access Point o Router). Esto se debe a que en dichas empresas suele haber mucha gente conectada y un dispositivo de red todo en uno no resulta suficiente.

\sphinxAtStartPar
En hogares el número de personas para conectarse a Internet suele ser de una a diez, por ahorro de costos es mejor usar los modem.


\begin{savenotes}\sphinxattablestart
\sphinxthistablewithglobalstyle
\centering
\begin{tabulary}{\linewidth}[t]{T}
\sphinxtoprule
\sphinxtableatstartofbodyhook

\begin{sphinxfigure-in-table}
\centering
\capstart
\noindent\sphinxincludegraphics[width=300\sphinxpxdimen]{{image091}.png}
\sphinxfigcaption{MODEM (Switch + Access Point + Router)}\label{\detokenize{1_guias/10informatica/G1003:id9}}\end{sphinxfigure-in-table}\relax
\\
\sphinxbottomrule
\end{tabulary}
\sphinxtableafterendhook\par
\sphinxattableend\end{savenotes}
\end{sphinxadmonition}

\begin{sphinxadmonition}{note}{\sphinxstylestrong{RACK de telecomunicaciones}}

\sphinxAtStartPar
Un rack de telecomunicaciones es una estructura metálica diseñada para alojar, organizar y proteger de forma centralizada equipos informáticos y de redes, como switches, routers, servidores, paneles de parcheo (patch panels) y sistemas de alimentación ininterrumpida (UPS).

\sphinxAtStartPar
Su principal función es optimizar el espacio en los cuartos de cableado o centros de datos, facilitando la ventilación de los dispositivos y el orden del cableado estructurado. Sus dimensiones internas están estandarizadas en anchura (generalmente 19 pulgadas) y la altura de los equipos se mide en Unidades de Rack (U).


\begin{savenotes}\sphinxattablestart
\sphinxthistablewithglobalstyle
\centering
\begin{tabulary}{\linewidth}[t]{T}
\sphinxtoprule
\sphinxtableatstartofbodyhook

\begin{sphinxfigure-in-table}
\centering
\capstart
\noindent\sphinxincludegraphics[width=400\sphinxpxdimen]{{image101}.png}
\sphinxfigcaption{RACK servidores OPENBSD 2009}\label{\detokenize{1_guias/10informatica/G1003:id10}}\end{sphinxfigure-in-table}\relax
\\
\sphinxbottomrule
\end{tabulary}
\sphinxtableafterendhook\par
\sphinxattableend\end{savenotes}
\end{sphinxadmonition}


\bigskip\hrule\bigskip



\subsubsection{Servicios de red}
\label{\detokenize{1_guias/10informatica/G1003:servicios-de-red}}
\sphinxAtStartPar
A continuación se muestran los servicios que suelen haber en una red de datos.

\begin{sphinxadmonition}{note}{\sphinxstylestrong{SSH}: Secure Shell}

\sphinxAtStartPar
Son las siglas de \sphinxstylestrong{Secure Shell} y es un \sphinxstylestrong{protocolo de red destinado principalmente a la conexión con máquinas a las que accedemos por línea de comandos}. En otras palabras, con SSH podemos conectarnos con servidores usando la red Internet como vía para las comunicaciones.

\sphinxAtStartPar
Debido a su seguridad, SSH es el modo preferido para la realización de conexión con servidores que necesitamos administrar. La diferencia con respecto a otros protocolos más antiguos como Telnet es que \sphinxstylestrong{el protocolo} \sphinxstylestrong{SSH siempre es seguro.} Sin embargo, aprovechando la seguridad de las comunicaciones, también se utiliza para otros objetivos como:
\begin{itemize}
\item {} 
\sphinxAtStartPar
\sphinxstylestrong{Transferencia de Archivos Segura}: Permite transferir archivos de forma segura entre sistemas locales y remotos utilizando herramientas como el comando SCP o SFTP.

\item {} 
\sphinxAtStartPar
\sphinxstylestrong{Creación de Túneles de Red}: SSH se utiliza para crear túneles de datos seguros que redirigen el tráfico de red a través de conexiones SSH, lo que puede ayudar a proteger la comunicación en redes no seguras. Se usan en sistemas como Ngrok, un software que permite a los desarrolladores exponer de manera remota los trabajos, tal como los tienen funcionando en su servidor de desarrollo local.

\end{itemize}


\begin{savenotes}\sphinxattablestart
\sphinxthistablewithglobalstyle
\centering
\begin{tabulary}{\linewidth}[t]{T}
\sphinxtoprule
\sphinxtableatstartofbodyhook

\begin{sphinxfigure-in-table}
\centering
\capstart
\noindent\sphinxincludegraphics[width=400\sphinxpxdimen]{{image12}.png}
\sphinxfigcaption{OpenSSH}\label{\detokenize{1_guias/10informatica/G1003:id11}}\end{sphinxfigure-in-table}\relax
\\
\sphinxbottomrule
\end{tabulary}
\sphinxtableafterendhook\par
\sphinxattableend\end{savenotes}
\end{sphinxadmonition}

\begin{sphinxadmonition}{note}{\sphinxstylestrong{DHCP}: Proveedor de direcciones IP (ejemplo: 192.168.115.100)}

\sphinxAtStartPar
El Protocolo de configuración dinámica de host (DHCP) es un protocolo cliente\sphinxhyphen{}servidor que proporciona automáticamente un host de protocolo de Internet (IP) con su dirección IP y otra información de configuración relacionada, como la máscara de subred y la puerta de enlace de predeterminada.
\end{sphinxadmonition}

\begin{sphinxadmonition}{note}{\sphinxstylestrong{DNS}: Proveedor de nombres de dominio (ejemplo: www.jfkennedy.edu.co)}

\sphinxAtStartPar
Sistema de nombres de dominio. Traduce los nombres de dominios aptos para lectura humana (por ejemplo, \sphinxhref{http://www.amazon.com/}{www.amazon.com}) a direcciones IP aptas para lectura por parte de máquinas (por ejemplo, \sphinxguilabel{18.155.246.193}).
\begin{quote}

\sphinxAtStartPar
El sistema DNS de Internet funciona como una agenda telefónica donde se administra el mapa entre los nombres y los números. Los servidores DNS convierten las solicitudes de nombres en direcciones IP, controlando a qué servidor se dirigirá un usuario final cuando escriba un nombre de dominio en su navegador web. Estas solicitudes se denominan \sphinxstylestrong{consultas}.
\end{quote}
\end{sphinxadmonition}

\begin{sphinxadmonition}{note}{\sphinxstylestrong{FTP}: Proveedor de archivos}

\sphinxAtStartPar
Son las siglas en inglés de File Transfer Protocol (protocolo de transferencia de archivos). FTP es el conjunto de reglas que los dispositivos de una red TCP/IP (Internet) utilizan para transferir archivos. Cuando usted usa Internet, en realidad se utiliza una variedad de diferentes protocolos. El FTP es simplemente el protocolo usado para transferir archivos.
\end{sphinxadmonition}

\begin{sphinxadmonition}{note}{\sphinxstylestrong{HTTP://} Proveedor de páginas web}

\sphinxAtStartPar
El http (del inglés HyperText Transfer Protocol o Protocolo de Transferencia de Hiper Textos) es el protocolo de transmisión de información de la World Wide Web. Se trata de un protocolo “sin estado”, es decir, que no lleva registro de visitas anteriores sino que siempre empieza de nuevo. La información relativa a visitas previas se almacena en estos sistemas en las llamadas “cookies”, almacenadas en el sistema cliente.
\end{sphinxadmonition}

\begin{sphinxadmonition}{note}{\sphinxstylestrong{HTTPS://} Proveedor de páginas web seguras}

\sphinxAtStartPar
Por \sphinxstylestrong{https} se entiende HyperText Transfer Procotol Secure o Protocolo Seguro de Transferencia de Hipertexto, que no es más que la versión segura del http, es decir, una variante del mismo protocolo que se basa en la creación de un canal cifrado para la transmisión de la información, lo cual lo hace más apropiado para ciertos datos de tipo sensible (como claves y usuarios personales).
\end{sphinxadmonition}

\begin{sphinxadmonition}{note}{\sphinxstylestrong{SMTP\sphinxhyphen{}POP\sphinxhyphen{}IMAP}: Proveedor de correo electrónico}

\sphinxAtStartPar
Los servidores de correo electrónico son indispensables para proveer de este servicio. Funcionan de forma similar a un buzón de correos central, donde se almacena y se gestionan mensajes de forma segura, en espera a ser leído por un grupo de clientes (usuarios).
\begin{itemize}
\item {} 
\sphinxAtStartPar
\sphinxstylestrong{SMTP (Servidores de correo de salida)}: Es el protocolo más usado por la flexibilidad en su operación al transferir información de un servidor de salida a uno de entrada, guardando la dirección de correo electrónico de quién lo envió.

\item {} 
\sphinxAtStartPar
\sphinxstylestrong{POP (Post Office Protocol Version 3)}: Es una de las versiones más antiguas de protocolo y confiables. Cuando el dispositivo se conecta al servidor de correo, éste detecta el ordenador que realiza la acción y los aloja ahí, una vez descargados se eliminan de inmediato del servidor de entrada.

\item {} 
\sphinxAtStartPar
\sphinxstylestrong{IMAP (Internet Message Access Protocol)}: Es un protocolo en los servidores de correo electrónico entrante que funciona para sincronizar el email a diferentes dispositivos. Se alojan siempre en el servidor y se visualiza a través de una copia de descarga.

\end{itemize}
\end{sphinxadmonition}

\begin{sphinxadmonition}{note}{\sphinxstylestrong{VOIP}: Voz IP}

\sphinxAtStartPar
En inglés, las siglas VoIP significan \sphinxstylestrong{Voice over Internet Protocol}, es decir, voz sobre protocolo de Internet. La telefonía VoIP es un método que permite realizar llamadas de voz a través de la red, utilizando la banda ancha. Estas llamadas capturan el sonido procedente de un dispositivo móvil o fijo con conexión a Internet a través de un micrófono y lo transforman en paquetes de datos digitales. Los datos se transfieren a través de la red IP hasta el dispositivo receptor, al igual que el resto de contenidos transmitidos a través de Internet.
\end{sphinxadmonition}

\begin{sphinxadmonition}{note}{\sphinxstylestrong{RDP}: Escritorio remoto}

\sphinxAtStartPar
Las siglas RDP responden a \sphinxstylestrong{Remote Desktop Protocol}, o, en español, Protocolo de Escritorio Remoto. El protocolo RDP, entonces, permite que el escritorio de un equipo informático sea controlado a distancia por un usuario remoto. Existen diferentes programas de RDP como el RDP Remote Desktop de Microsoft –también conocido como RDP Windows–, o los más populares TeamViewer o PC Anywhere.
\end{sphinxadmonition}

\begin{sphinxadmonition}{note}{\sphinxstylestrong{FIREWALL}: Protección y control de tráfico en la red}

\sphinxAtStartPar
Un firewall es un sistema de seguridad que supervisa y controla el tráfico de la red en base a un conjunto de reglas de seguridad. Los firewalls suelen situarse entre una red de confianza y una red no fiable; con frecuencia, la red no fiable es Internet. Por ejemplo, las redes de oficinas suelen utilizar un firewall para proteger su red de las amenazas en línea.
\end{sphinxadmonition}

\begin{sphinxadmonition}{note}{\sphinxstylestrong{PROXY}: Intermediario en la red}

\sphinxAtStartPar
Un servidor proxy proporciona una puerta de enlace entre los usuarios e Internet. Es un servidor denominado “intermediario”, porque está entre los usuarios finales y las páginas web que visitan en línea. Cuando una computadora se conecta a Internet, utiliza una dirección IP.  Algunas personas usan los proxies para fines personales, como ocultar su ubicación mientras miran películas en línea. Sin embargo, para una compañía, pueden utilizarse para realizar varias tareas clave como las siguientes:
\begin{itemize}
\item {} 
\sphinxAtStartPar
Mejorar la seguridad.

\item {} 
\sphinxAtStartPar
Proteger la actividad de los empleados en Internet de las personas que intentan espiarlas.

\item {} 
\sphinxAtStartPar
Equilibrar el tráfico de Internet para evitar choques.

\item {} 
\sphinxAtStartPar
Controlar el acceso de los empleados a los sitios web.

\item {} 
\sphinxAtStartPar
Guardar el ancho de banda almacenando archivos en caché o comprimiendo el tráfico entrante.

\end{itemize}
\end{sphinxadmonition}

\begin{sphinxadmonition}{note}{\sphinxstylestrong{VPN}: Red Privada Virtual}

\sphinxAtStartPar
Una VPN o red privada virtual crea una conexión de red privada entre dispositivos a través de Internet. Las VPN se utilizan para transmitir datos de forma segura y anónima a través de redes públicas. Su funcionamiento consiste en ocultar las direcciones IP de los usuarios y cifrar los datos para que nadie que no esté autorizado a recibirlos pueda leerlos.
\end{sphinxadmonition}


\bigskip\hrule\bigskip



\subsubsection{¿Qué es la ciberseguridad?}
\label{\detokenize{1_guias/10informatica/G1003:que-es-la-ciberseguridad}}
\begin{sphinxadmonition}{note}{Protección de sistemas y redes}

\sphinxAtStartPar
La ciberseguridad es la implementación de tecnologías, prácticas y políticas diseñadas para \sphinxstylestrong{proteger} sistemas, redes y programas de ataques digitales.


\begin{savenotes}\sphinxattablestart
\sphinxthistablewithglobalstyle
\centering
\begin{tabulary}{\linewidth}[t]{T}
\sphinxtoprule
\sphinxtableatstartofbodyhook

\begin{sphinxfigure-in-table}
\centering
\capstart
\noindent\sphinxincludegraphics[width=300\sphinxpxdimen]{{image11}.png}
\sphinxfigcaption{Acceso remoto al sistema}\label{\detokenize{1_guias/10informatica/G1003:id12}}\end{sphinxfigure-in-table}\relax
\\
\sphinxbottomrule
\end{tabulary}
\sphinxtableafterendhook\par
\sphinxattableend\end{savenotes}

\sphinxAtStartPar
La ciberseguridad tiene los siguiente enfoques:
\begin{enumerate}
\sphinxsetlistlabels{\arabic}{enumi}{enumii}{}{.}%
\item {} 
\sphinxAtStartPar
\sphinxstylestrong{Protección de datos}: Se enfoca en salvaguardar la información sensible de accesos no autorizados y ataques maliciosos.

\item {} 
\sphinxAtStartPar
\sphinxstylestrong{Defensa de sistemas}: Incluye la protección de computadoras, servidores, dispositivos móviles y redes contra amenazas cibernéticas.

\item {} 
\sphinxAtStartPar
\sphinxstylestrong{Prevención de ciberataques}: Busca mitigar el impacto de ataques como malware, phishing y ransomware.

\item {} 
\sphinxAtStartPar
\sphinxstylestrong{Convergencia de personas y tecnología}: La ciberseguridad combina la colaboración de personas, procesos y tecnología para crear un entorno seguro.

\end{enumerate}

\sphinxAtStartPar
La ciberseguridad se puede clasificar en \sphinxstylestrong{seguridad informática} y \sphinxstylestrong{seguridad de la información}.
\begin{quote}

\sphinxAtStartPar
Una empresa debe implementar \sphinxstyleemphasis{Antivirus} y \sphinxstyleemphasis{Firewalls} (\sphinxstylestrong{seguridad informática}) pero también debe establecer políticas sobre por ejemplo, quién puede \sphinxstyleemphasis{acceder a documentos confidenciales} o cómo \sphinxstyleemphasis{deben manejarse las copias físicas} (\sphinxstylestrong{seguridad de la información}).
\end{quote}
\end{sphinxadmonition}


\bigskip\hrule\bigskip



\subsubsection{Seguridad informática}
\label{\detokenize{1_guias/10informatica/G1003:seguridad-informatica}}
\sphinxAtStartPar
La \sphinxstylestrong{seguridad informática} es el estudio e implementación de tecnologías para proteger sistemas informáticos, redes y datos en una organización.
\begin{quote}

\sphinxAtStartPar
\sphinxstyleemphasis{“Someter un sistema informático a constantes intentos de intrusión controlados, puede ser la mejor vía para garantizar su seguridad”}.
\end{quote}

\sphinxAtStartPar
La seguridad informática estudia los siguientes campos:

\begin{sphinxadmonition}{note}{\sphinxstylestrong{Ingeniería social}: Hacking de seres humanos}

\sphinxAtStartPar
La ingeniería social es una técnica de manipulación que aprovecha el error humano para obtener información privada, acceso a sistemas u objetos de valor. En el caso del delito cibernético, estas estafas de “hackeo de humanos” tienden a hacer que los usuarios desprevenidos expongan datos, propaguen infecciones de malware o den acceso a sistemas restringidos.

\sphinxAtStartPar
Los ataques de ingeniería social pueden ocurrir en línea, en persona y a través de otras interacciones.

\sphinxAtStartPar
Las estafas basadas en la ingeniería social se basan en cómo las personas piensan y actúan. Como tal, los ataques de ingeniería social son muy útiles para manipular el comportamiento de un usuario. Una vez que un atacante comprende qué motiva las acciones de un usuario, puede engañar y manipular al usuario de manera efectiva. Además, los piratas informáticos intentan explotar la falta de conocimiento de un usuario.

\sphinxAtStartPar
Gracias a la velocidad de la tecnología, muchos consumidores y empleados no son conscientes de ciertas amenazas, como las descargas no autorizadas (drive\sphinxhyphen{}in). Es posible que los usuarios no se den cuenta del valor real de sus datos personales, como el número de teléfono. Como resultado, muchos usuarios no están seguros de cómo protegerse mejor a sí mismos y a su información.
\end{sphinxadmonition}

\begin{sphinxadmonition}{note}{\sphinxstylestrong{Hacking en conexiones inalámbricas}: (WIFI, infrarojo, microondas, Bluetooth)}

\sphinxAtStartPar
Las redes WIFI se encuentran en muchos lugares desde el hogar hasta en la empresa. Una gran cantidad de la lugares públicos también las adoptan y usan. Esto resulta en un canal de comunicación muy usado y a su vez vulnerable.

\sphinxAtStartPar
En redes WIFI se debe conocer de:
\begin{itemize}
\item {} 
\sphinxAtStartPar
Protocolos.

\item {} 
\sphinxAtStartPar
Estándares.

\item {} 
\sphinxAtStartPar
Algoritmos de cifrado.

\item {} 
\sphinxAtStartPar
Sentencias legales.

\end{itemize}

\sphinxAtStartPar
Con estos conceptos en mente es posible tratar temas de seguridad en diferentes redes inalámbricas, el equipo de laboratorio necesario y las herramientas que se pueden usar. Con este equipo ya sería posible analizar señales, monitorizar tráfico y simular diferentes ataques sobre redes abiertas como WEP, WPA y WPA2.
\end{sphinxadmonition}

\begin{sphinxadmonition}{note}{\sphinxstylestrong{Ataques controlados a la red}: (PENTEST)}

\sphinxAtStartPar
El \sphinxstylestrong{pentesting} es la técnica de ciberseguridad consistente en atacar entornos informáticos con la intención de descubrir vulnerabilidades en los mismos, con el objetivo de reunir la información necesaria para poder prevenir en el futuro ataques externos hacia esos mismos equipos.

\sphinxAtStartPar
El \sphinxstylestrong{pentesting} es una técnica 100\% legal… siempre y cuando los ataques se realicen en sistemas de nuestra propiedad o con permiso del propietario.

\sphinxAtStartPar
La realización de un \sphinxcode{\sphinxupquote{pentest}} requiere de un amplio conocimiento de múltiples tecnologías y sistemas operativos. Las herramientas necesarias (escáner de puertos, software de sniffing de redes, etc) suelen venir agrupadas en ‘toolkits’, en algún caso como parte integrante de versiones ‘live’ de GNU/Linux, como el popular y reputado Kali Linux (el antiguo Backtrack).
\end{sphinxadmonition}

\begin{sphinxadmonition}{note}{\sphinxstylestrong{Defensa de la red:} Seguridad perimetral y monitorización}

\sphinxAtStartPar
Las \sphinxstylestrong{redes de telecomunicaciones} envían y reciben grandes cantidades de información. Este campo se centra en proteger los datos que circulan por ellas. Para conocer el nivel de seguridad es necesario interceptar las comunicaciones desde un punto de vista ético y legal. De hacerlo de forma diferente conllevaría a problemas de millonarias multas o cárcel.

\sphinxAtStartPar
La monitorización y defensa de redes tiene como objetivo mantener seguras las redes por medio de la detección de intrusos y pruebas de pentesting. La confidencialidad de la información debe asegurarse con herramientas como:
\begin{itemize}
\item {} 
\sphinxAtStartPar
Certificados digitales.

\item {} 
\sphinxAtStartPar
Cifrado de datos.

\item {} 
\sphinxAtStartPar
Comunicaciones seguras con SSH.

\item {} 
\sphinxAtStartPar
IPsec.

\item {} 
\sphinxAtStartPar
VNP\sphinxhyphen{}SSL.

\item {} 
\sphinxAtStartPar
Wireshark.

\end{itemize}
\end{sphinxadmonition}

\begin{sphinxadmonition}{note}{\sphinxstylestrong{Hacking controlado en aplicaciones web}}

\sphinxAtStartPar
El pentesting web o la seguridad en páginas o aplicaciones web consiste en diferentes técnicas de ataque de manera guiada contra aplicativos web, así como las contramedidas necesarias para proteger este tipo de sistemas.

\sphinxAtStartPar
Las aplicaciones web son unas de las tecnologías de mayor uso hoy en día. Esto se debe a que los ingenieros quieren llevar todos los servicios a la web. Para mejorar la seguridad de los sitios web es necesario saber como se realizan los ataques y como se explotan vulnerabilidades. A continuación se listan técnicas usadas en el pentesting de páginas web:
\begin{itemize}
\item {} 
\sphinxAtStartPar
SQL Injection.

\item {} 
\sphinxAtStartPar
Blind SQL injection.

\item {} 
\sphinxAtStartPar
XSS Reflected.

\item {} 
\sphinxAtStartPar
XSS stored.

\item {} 
\sphinxAtStartPar
CSRF.

\item {} 
\sphinxAtStartPar
XML Injection.

\item {} 
\sphinxAtStartPar
Session fixation.

\item {} 
\sphinxAtStartPar
Botnet DDOS html5.

\end{itemize}
\end{sphinxadmonition}


\bigskip\hrule\bigskip



\subsubsection{Seguridad de la información}
\label{\detokenize{1_guias/10informatica/G1003:seguridad-de-la-informacion}}
\sphinxAtStartPar
La \sphinxstylestrong{seguridad de la información} es algo más que la seguridad informática. Se centra en toda la empresa. Al fin y al cabo, la seguridad de la información no sólo se dirige a los datos procesados por los sistemas electrónicos, también abarca todos los activos de la empresa que deben protegerse, incluidos los que se encuentran en soportes de datos analógicos, como el papel.
\begin{quote}

\sphinxAtStartPar
La \sphinxstylestrong{seguridad de la información} tiene un alcance más amplio que la seguridad informática. \sphinxstylestrong{La seguridad informática}, se refiere “sólo” al uso e implementación de tecnología para proteger la información. No tiene en cuenta normatividades. Por tanto la seguridad informática hace parte de la seguridad de la información.
\end{quote}

\begin{sphinxadmonition}{note}{3 principios de seguridad de la información}

\sphinxAtStartPar
Los tres principios de la seguridad de la información son:
\begin{itemize}
\item {} 
\sphinxAtStartPar
\sphinxstylestrong{Confidencialidad:} Garantizar que solo personas autorizadas accedan a la información.

\item {} 
\sphinxAtStartPar
\sphinxstylestrong{Integridad:} Asegurar que la información no sea alterada indebidamente.

\item {} 
\sphinxAtStartPar
\sphinxstylestrong{Disponibilidad:} Garantizar el acceso a la información cuando se necesite.

\end{itemize}

\sphinxAtStartPar
Se aplican, por tanto, también a documentos importantes, que debe llegar a la puerta de su destinatario a tiempo, de forma fiable e intacta, transportada por un mensajero, pero totalmente analógico. Y estos principios se aplican igualmente a una hoja de papel que contenga información confidencial, pero que esté sobre un escritorio desatendido para que cualquiera pueda verla o que esté esperando en la fotocopiadora, libremente accesible, para un acceso no autorizado.

\sphinxAtStartPar
A continuación se muestran unos temas relacionados con la seguridad de la información:
\begin{itemize}
\item {} 
\sphinxAtStartPar
Asegurar alimentación energética del hardware.

\item {} 
\sphinxAtStartPar
Medidas contra el sobre\sphinxhyphen{}calentamiento del hardware.

\item {} 
\sphinxAtStartPar
Análisis de virus y programas seguros.

\item {} 
\sphinxAtStartPar
Organización de las estructuras de las carpetas.

\item {} 
\sphinxAtStartPar
Configuración y actualización de los cortafuegos.

\item {} 
\sphinxAtStartPar
Formación de los empleados en temas de de ingeniería social.

\end{itemize}

\sphinxAtStartPar
En la seguridad de la información, lo importante es identificar los \sphinxstylestrong{activos de información} y lograr una clasificación certera de los \sphinxstylestrong{activos de información críticos}. De esta forma encaminar todos los recursos disponibles para poder proteger la información de estos activos críticos. Para ello entonces se usa el análisis de riesgos.
\end{sphinxadmonition}

\begin{sphinxadmonition}{note}{\sphinxstylestrong{ISO 27001}: Sistema de gestión de seguridad de la Información}

\sphinxAtStartPar
La \sphinxstylestrong{norma ISO 27001} es un estándar internacional que establece los requisitos para la implementación, mantenimiento y mejora continua de un Sistema de Gestión de la Seguridad de la Información (SGSI). Este sistema se utiliza para proteger la confidencialidad, integridad y disponibilidad de la información. La norma proporciona un marco para la seguridad de la información que ayuda a las organizaciones a identificar y gestionar sus riesgos de seguridad de la información de manera efectiva.

\sphinxAtStartPar
La \sphinxstylestrong{norma ISO 27001} se aplica a cualquier tipo de organización, incluyendo pequeñas y medianas empresas, grandes corporaciones, instituciones gubernamentales y sin fines de lucro. También se puede aplicar en cualquier sector, incluyendo tecnología de la información, finanzas, salud y servicios públicos.

\sphinxAtStartPar
Para la implementación de \sphinxstylestrong{ISO 27001} en una empresa se deben cumplir las siguientes fases:
\begin{itemize}
\item {} 
\sphinxAtStartPar
\sphinxstylestrong{Fase de planificación:} Durante la fase de planificación, la organización identifica sus requisitos de seguridad de la información y establece un plan para implementar el SGSI.

\item {} 
\sphinxAtStartPar
\sphinxstylestrong{Fase de implementación:} La fase de implementación incluye la creación de políticas, procedimientos y controles para proteger la información.

\item {} 
\sphinxAtStartPar
\sphinxstylestrong{Fase de evaluación:} Durante la fase de evaluación, la organización evalúa la eficacia de su SGSI e identifica áreas de mejora.

\item {} 
\sphinxAtStartPar
\sphinxstylestrong{Fase de mejora continua:} La fase de mejora continua implica la identificación y aplicación de mejoras a los procesos y controles del SGSI.

\end{itemize}
\end{sphinxadmonition}

\begin{sphinxadmonition}{note}{Seguridad Informática Vs Seguridad de la Información}


\begin{savenotes}\sphinxattablestart
\sphinxthistablewithglobalstyle
\centering
\begin{tabulary}{\linewidth}[t]{TT}
\sphinxtoprule
\sphinxstyletheadfamily 
\sphinxAtStartPar
Seguridad Informática
&\sphinxstyletheadfamily 
\sphinxAtStartPar
Seguridad de la Información
\\
\sphinxmidrule
\sphinxtableatstartofbodyhook
\sphinxAtStartPar
Se centra en la protección de la infraestructura tecnológica y sistemas computacionales
&
\sphinxAtStartPar
Abarca la protección de toda la información independiente de su formato o medio.
\\
\sphinxhline
\sphinxAtStartPar
Protege contra amenazas técnicas como Virus, Malware y ciber\sphinxhyphen{}delincuentes
&
\sphinxAtStartPar
Protege contra amenazas físicas, digitales y humanas.
\\
\sphinxhline
\sphinxAtStartPar
Implementa soluciones como Firewalls, antivirus y cifrado
&
\sphinxAtStartPar
Implementa políticas, procedimientos y controles organizacionales.
\\
\sphinxbottomrule
\end{tabulary}
\sphinxtableafterendhook\par
\sphinxattableend\end{savenotes}
\end{sphinxadmonition}


\bigskip\hrule\bigskip



\subsubsection{Delitos informáticos en Colombia}
\label{\detokenize{1_guias/10informatica/G1003:delitos-informaticos-en-colombia}}
\begin{sphinxadmonition}{note}{Regulaciones en la legislación Colombiana}

\sphinxAtStartPar
La legislación en el ámbito de la tecnología e informática es fundamental por los siguientes aspectos:
\begin{itemize}
\item {} 
\sphinxAtStartPar
\sphinxstylestrong{Protección de datos personales:} Establece normas para el manejo seguro y ético de la información personal de los usuarios.

\item {} 
\sphinxAtStartPar
\sphinxstylestrong{Derechos digitales:} Define y protege los derechos de los ciudadanos en el entorno digital.

\item {} 
\sphinxAtStartPar
\sphinxstylestrong{Seguridad cibernética:} Establece marcos legales para combatir delitos informáticos y proteger la infraestructura digital.

\item {} 
\sphinxAtStartPar
\sphinxstylestrong{Comercio electrónico:} Regula las transacciones en línea y protege tanto a consumidores como a empresas.

\item {} 
\sphinxAtStartPar
\sphinxstylestrong{Propiedad intelectual:} Protege los derechos de autor en contenidos digitales y desarrollos tecnológicos.

\item {} 
\sphinxAtStartPar
\sphinxstylestrong{Responsabilidad digital:} Define las obligaciones y responsabilidades de usuarios y proveedores de servicios tecnológicos.

\end{itemize}

\sphinxAtStartPar
En Colombia, estas regulaciones son especialmente importantes para garantizar un desarrollo tecnológico seguro y equitativo en la sociedad digital.
\end{sphinxadmonition}

\begin{sphinxadmonition}{note}{\sphinxstylestrong{Ley 1273 de 2009}: Penas y multas por delitos informáticos en Colombia}

\sphinxAtStartPar
De acuerdo con la \sphinxstylestrong{Ley 1273 de 2009}, las penas por delitos informáticos en Colombia varían según el delito:
\begin{itemize}
\item {} 
\sphinxAtStartPar
\sphinxstylestrong{Acceso abusivo a sistemas informáticos:} Prisión de 48 a 96 meses y multas de 100 a 1000 salarios mínimos legales mensuales vigentes.

\item {} 
\sphinxAtStartPar
\sphinxstylestrong{Interceptación de datos:} Prisión de 36 a 72 meses y multas de 100 a 1000 salarios mínimos.

\item {} 
\sphinxAtStartPar
\sphinxstylestrong{Daño informático:} Prisión de 48 a 96 meses y multas de 100 a 1000 salarios mínimos.

\item {} 
\sphinxAtStartPar
\sphinxstylestrong{Suplantación de sitios web:} Prisión de 48 a 96 meses y multas de 100 a 1000 salarios mínimos.

\end{itemize}

\sphinxAtStartPar
Las penas pueden aumentar de la mitad a las tres cuartas partes si el delito se comete sobre redes o sistemas informáticos gubernamentales, del sector financiero, o contra sistemas de seguridad nacional.
\end{sphinxadmonition}


\subsubsection{Hacking ético}
\label{\detokenize{1_guias/10informatica/G1003:hacking-etico}}
\begin{sphinxadmonition}{note}{Ataques autorizados}

\sphinxAtStartPar
El \sphinxstylestrong{hacking ético} es la práctica de utilizar las mismas técnicas y herramientas que los hackers maliciosos, pero con el permiso y el objetivo de mejorar la seguridad de un sistema o red. En esencia, los hackers éticos actúan como “inspectores de seguridad” para identificar vulnerabilidades antes de que los atacantes puedan explotarlas.

\sphinxAtStartPar
Aquí hay algunos puntos clave sobre el hacking ético:
\begin{itemize}
\item {} 
\sphinxAtStartPar
\sphinxstylestrong{Autorización:} Los hackers éticos siempre deben tener el permiso explícito del propietario del sistema o red que están probando.

\item {} 
\sphinxAtStartPar
\sphinxstylestrong{Objetivo:} El objetivo principal es identificar y solucionar vulnerabilidades de seguridad.

\item {} 
\sphinxAtStartPar
\sphinxstylestrong{Metodología:} Siguen un proceso estructurado que incluye:
\begin{itemize}
\item {} 
\sphinxAtStartPar
\sphinxstylestrong{Reconocimiento:} Recopilación de información sobre el objetivo.

\item {} 
\sphinxAtStartPar
\sphinxstylestrong{Escaneo:} Identificación de puertos abiertos, servicios y sistemas operativos.

\item {} 
\sphinxAtStartPar
\sphinxstylestrong{Obtención de acceso:} Explotación de vulnerabilidades para acceder al sistema (con permiso).

\item {} 
\sphinxAtStartPar
\sphinxstylestrong{Mantenimiento del acceso:} Asegurar el acceso continuo para evaluar la seguridad.

\item {} 
\sphinxAtStartPar
\sphinxstylestrong{Análisis y reporte:} Documentación de las vulnerabilidades encontradas y recomendaciones para su corrección.

\end{itemize}

\item {} 
\sphinxAtStartPar
\sphinxstylestrong{Herramientas:} Utilizan una amplia gama de herramientas, muchas de las cuales son las mismas que usan los hackers maliciosos, como escáneres de vulnerabilidades, analizadores de red, herramientas de cracking de contraseñas, etc.

\item {} 
\sphinxAtStartPar
\sphinxstylestrong{Certificaciones:} Existen certificaciones profesionales como Certified Ethical Hacker (CEH) que validan los conocimientos y habilidades de los hackers éticos.

\end{itemize}

\sphinxAtStartPar
El hacking ético es una parte crucial de la \sphinxstylestrong{ciberseguridad moderna}. Ayuda a las organizaciones a proteger sus sistemas y datos de ataques maliciosos al identificar y solucionar las debilidades de seguridad.
\end{sphinxadmonition}


\bigskip\hrule\bigskip



\subsection{A02: Taller (+0,8)}
\label{\detokenize{1_guias/10informatica/G1003:a02-taller-0-8}}
\sphinxAtStartPar
\DUrole{sd-sphinx-override}{\DUrole{sd-badge}{\DUrole{sd-bg-warning}{\DUrole{sd-bg-text-warning}{Actividad práctica}}}}
\begin{enumerate}
\sphinxsetlistlabels{\arabic}{enumi}{enumii}{}{.}%
\item {} 
\sphinxAtStartPar
En una hoja de cuaderno, cree un mapa conceptual con los temas de \sphinxstylestrong{A01: Redes de datos, Internet y Ciberseguridad}.

\item {} 
\sphinxAtStartPar
En su cuaderno, responda las siguientes preguntas:
\begin{itemize}
\item {} 
\sphinxAtStartPar
¿Qué es una red de datos?

\item {} 
\sphinxAtStartPar
¿En que se diferencian el INTERNET de una Intranet?

\item {} 
\sphinxAtStartPar
¿En que se diferencian los dispositivos de \sphinxstylestrong{tipo cliente} y los dispositivos de \sphinxstylestrong{tipo servidor}?

\item {} 
\sphinxAtStartPar
¿Por que existen los dispositivos todo en uno?

\item {} 
\sphinxAtStartPar
¿Que son los servicios de red?

\end{itemize}

\item {} 
\sphinxAtStartPar
Escriba las siguientes frases en su cuaderno y responda si es \sphinxguilabel{falso} o \sphinxguilabel{verdadero} y diga el por qué:
\begin{itemize}
\item {} 
\sphinxAtStartPar
INTERNET permite conectar personas al rededor el mundo así como acceder a diferentes servicios en la nube. \_\_\_

\item {} 
\sphinxAtStartPar
En un INTRANET se deben tener accesos restringidos. No todo el mundo puede ingresar. \_\_\_

\item {} 
\sphinxAtStartPar
INTERNET interconecta varias INTRANETS de forma controlada. \_\_\_

\item {} 
\sphinxAtStartPar
A la página web de la institución educativa John F. Kennedy solo se puede acceder desde la INTRANET. \_\_\_

\item {} 
\sphinxAtStartPar
La INTRANET del portal del aula 115 se puede acceder desde cualquier lugar del mundo. \_\_\_

\end{itemize}

\item {} 
\sphinxAtStartPar
Indique qué es y en su cuaderno haga un dibujo de los siguientes dispositivos:
\begin{itemize}
\item {} 
\sphinxAtStartPar
Servidores.

\item {} 
\sphinxAtStartPar
Equipos cliente.

\item {} 
\sphinxAtStartPar
Switch.

\item {} 
\sphinxAtStartPar
Access Point.

\item {} 
\sphinxAtStartPar
Router.

\item {} 
\sphinxAtStartPar
Modem.

\end{itemize}

\item {} 
\sphinxAtStartPar
El siguiente es un plano de la ubicación de diferentes \sphinxstyleemphasis{dispositivos}* del aula 115 de la \sphinxstylestrong{Institución Educativa John F. Kennedy}. En una hoja de cuaderno, identifique en el plano los siguientes dispositivos:
\begin{itemize}
\item {} 
\sphinxAtStartPar
11 portátiles (numerarlos)

\item {} 
\sphinxAtStartPar
11 PCs de escritorio (numerarlos)

\item {} 
\sphinxAtStartPar
1 impresora

\item {} 
\sphinxAtStartPar
1 smart\sphinxhyphen{}TV

\item {} 
\sphinxAtStartPar
1 servidor

\item {} 
\sphinxAtStartPar
1 switch

\item {} 
\sphinxAtStartPar
1 modem

\end{itemize}
\begin{quote}

\sphinxAtStartPar
\sphinxstylestrong{ATENCIÓN:} El plano puede estar desactualizado. En ese caso ajustar el plano al estado actual del aula.
\end{quote}


\begin{savenotes}\sphinxattablestart
\sphinxthistablewithglobalstyle
\centering
\begin{tabulary}{\linewidth}[t]{T}
\sphinxtoprule
\sphinxtableatstartofbodyhook

\begin{sphinxfigure-in-table}
\centering
\capstart
\noindent\sphinxincludegraphics[width=400\sphinxpxdimen]{{image13}.png}
\sphinxfigcaption{Plano de red aula 115}\label{\detokenize{1_guias/10informatica/G1003:id13}}\end{sphinxfigure-in-table}\relax
\\
\sphinxbottomrule
\end{tabulary}
\sphinxtableafterendhook\par
\sphinxattableend\end{savenotes}

\item {} 
\sphinxAtStartPar
Con el plano anterior terminado, conecte los dispositivos de la siguiente forma:
\begin{itemize}
\item {} 
\sphinxAtStartPar
Dibuje con líneas, las conexiones de los dispositivos que están conectados con \sphinxstylestrong{cable de red} a la INTRANET.

\item {} 
\sphinxAtStartPar
Dibuje con puntos, las conexiones de los dispositivos que están conectados de \sphinxstylestrong{forma inalámbrica} a la INTRANET.

\end{itemize}

\item {} 
\sphinxAtStartPar
Copie la siguiente sopa de letras en su cuaderno e IDENTIFIQUE 10 ejemplos de servicios de red:

\end{enumerate}


\begin{savenotes}\sphinxattablestart
\sphinxthistablewithglobalstyle
\centering
\begin{tabulary}{\linewidth}[t]{T}
\sphinxtoprule
\sphinxtableatstartofbodyhook

\begin{sphinxfigure-in-table}
\centering
\capstart
\noindent\sphinxincludegraphics[width=550\sphinxpxdimen]{{image14}.png}
\sphinxfigcaption{Servicios de red}\label{\detokenize{1_guias/10informatica/G1003:id14}}\end{sphinxfigure-in-table}\relax
\\
\sphinxbottomrule
\end{tabulary}
\sphinxtableafterendhook\par
\sphinxattableend\end{savenotes}

\sphinxAtStartPar
Ejemplo:
\begin{quote}

\sphinxAtStartPar
\sphinxstylestrong{FIREWALL:} Supervisa y controla un tráfico de red en base a unas reglas de seguridad.
\end{quote}
\begin{enumerate}
\sphinxsetlistlabels{\arabic}{enumi}{enumii}{}{.}%
\setcounter{enumi}{7}
\item {} 
\sphinxAtStartPar
Complete la siguiente tabla en su cuaderno:

\end{enumerate}


\begin{savenotes}\sphinxattablestart
\sphinxthistablewithglobalstyle
\centering
\begin{tabulary}{\linewidth}[t]{TTT}
\sphinxtoprule
\sphinxstyletheadfamily 
\sphinxAtStartPar
Aspecto
&\sphinxstyletheadfamily 
\sphinxAtStartPar
Seguridad de la Información
&\sphinxstyletheadfamily 
\sphinxAtStartPar
Seguridad Informática
\\
\sphinxmidrule
\sphinxtableatstartofbodyhook
\sphinxAtStartPar
Enfoque
&
\sphinxAtStartPar

&
\sphinxAtStartPar

\\
\sphinxhline
\sphinxAtStartPar
Amenazas
&
\sphinxAtStartPar

&
\sphinxAtStartPar

\\
\sphinxhline
\sphinxAtStartPar
Elementos
&
\sphinxAtStartPar

&
\sphinxAtStartPar

\\
\sphinxbottomrule
\end{tabulary}
\sphinxtableafterendhook\par
\sphinxattableend\end{savenotes}
\begin{enumerate}
\sphinxsetlistlabels{\arabic}{enumi}{enumii}{}{.}%
\setcounter{enumi}{8}
\item {} 
\sphinxAtStartPar
En su cuaderno escriba y analice los siguientes casos e indique cual de los 3 principios de la seguridad de la información se ajusta a dicho caso. Justifique su respuesta.

\end{enumerate}
\begin{quote}

\sphinxAtStartPar
\sphinxstylestrong{CASO 01:} Por medio de un ataque cibernético el día de ayer a la empresa X le robaron toda la información personal de sus empleados. Ahora sus datos como su localización, teléfono y lugar de nacimiento están expuestos en internet. Esto dio a que suplantaran muchas identidades.
\end{quote}
\begin{quote}

\sphinxAtStartPar
\sphinxstylestrong{CASO 02:} Debido a un software malicioso en los servidores principales del banco X los servicios de cajero automático no van a estar disponibles por 2 días. Esto hará que muchas personas pierdan su confianza en la nueva planta administrativa del banco elegida recientemente.
\end{quote}
\begin{quote}

\sphinxAtStartPar
\sphinxstylestrong{CASO 03:} Un error humano de un técnico practicante modificó la información en las bases de datos de la nómina de este mes. Esto genera incumplimiento en el pago de salarios de los trabajadores de la empresa. Los funcionarios tendrán que esperar una semana más para que reciban su sueldo.
\end{quote}
\begin{enumerate}
\sphinxsetlistlabels{\arabic}{enumi}{enumii}{}{.}%
\setcounter{enumi}{9}
\item {} 
\sphinxAtStartPar
Lea la sección de \sphinxstylestrong{Delitos informáticos en Colombia} y complete las siguientes frases en su cuaderno:
\begin{itemize}
\item {} 
\sphinxAtStartPar
Un artista está involucrado en un problema de plagio de una canción. La \_\_\_\_\_\_\_\_\_\_\_\_\_\_\_ corresponde al abuelo de una actriz que debutó en los primeros años en YouTube.

\item {} 
\sphinxAtStartPar
Nuevamente atentaron contra la \_\_\_\_\_\_\_\_\_\_\_\_\_\_\_\_\_. Esta vez ingresaron sin autorización al servidor y robaron los datos personales de varias personas.

\item {} 
\sphinxAtStartPar
El prestador del servicio de INTERNET últimamente ha invertido bastante en calidad. Tanto es así que ya gano varios cliente nuevos. Definitivamente se tomaron en serio su \_\_\_\_\_\_\_\_\_\_\_\_\_\_\_\_.

\item {} 
\sphinxAtStartPar
Estoy cansado de que cada vez que intento ingresar al portal me manden a leer sobre los riesgos de exponer mi información personal. Ya se que tengo derecho a la \_\_\_\_\_\_\_\_\_\_\_\_\_\_\_\_\_\_\_\_\_ pero están muy intensos.

\item {} 
\sphinxAtStartPar
Ya está listo el programa con la nueva funcionalidad de \_\_\_\_\_\_\_\_\_\_\_\_\_\_\_\_\_\_\_\_. Ahora si tenemos pasarela de pagos en línea.

\end{itemize}

\item {} 
\sphinxAtStartPar
Suponga que usted es un juez de delitos informáticos. Copie y analice los siguientes casos en su cuaderno, e imponga las penas necesarias según la \sphinxstylestrong{ley 1273 del 2009}:
\begin{quote}

\sphinxAtStartPar
\sphinxstylestrong{CASO 01:} El acusado se declara culpable de haber clonado el sitio web y así engañar a varias amas de casa vendiéndoles productos falsos de aseo.
\end{quote}
\begin{quote}

\sphinxAtStartPar
\sphinxstylestrong{CASO 02:} La víctima testificó con evidencia probatoria que sus comunicaciones fueron chuzadas por su ex\sphinxhyphen{}esposa mientras conducía la noche anterior.
\end{quote}
\begin{quote}

\sphinxAtStartPar
\sphinxstylestrong{CASO 03:} El ingeniero accedió sin permiso al sistema empresarial. Esto fue 24 horas después de haber sido despedido. Por tanto todos sus accesos ya habían sido retirados. Sin embargo previamente había dejado una backdoor.
\end{quote}

\end{enumerate}


\bigskip\hrule\bigskip



\subsection{A03: Análisis y reflexión (+4,0)}
\label{\detokenize{1_guias/10informatica/G1003:a03-analisis-y-reflexion-4-0}}
\sphinxAtStartPar
\DUrole{sd-sphinx-override}{\DUrole{sd-badge}{\DUrole{sd-bg-danger}{\DUrole{sd-bg-text-danger}{Transferencia}}}}


\begin{savenotes}\sphinxattablestart
\sphinxthistablewithglobalstyle
\centering
\begin{tabulary}{\linewidth}[t]{TT}
\sphinxtoprule
\sphinxtableatstartofbodyhook
\sphinxAtStartPar
\sphinxstylestrong{TIEMPO}
&
\sphinxAtStartPar
59 minutos
\\
\sphinxhline
\sphinxAtStartPar
\sphinxstylestrong{PREGUNTAS}
&
\sphinxAtStartPar
10
\\
\sphinxhline
\sphinxAtStartPar
\sphinxstylestrong{OBJETIVO}
&
\sphinxAtStartPar
Identificar qué le llamó la atención de la temática de \sphinxstylestrong{REDES y Ciberseguridad}.
\\
\sphinxhline
\sphinxAtStartPar
\sphinxstylestrong{ACTIVIDAD}
&
\sphinxAtStartPar
En una hoja al final de su cuaderno, desarrolle un \sphinxstylestrong{texto reflexivo} utilizando las siguientes preguntas como guía para estructurar su escrito.
\\
\sphinxbottomrule
\end{tabulary}
\sphinxtableafterendhook\par
\sphinxattableend\end{savenotes}

\begin{sphinxadmonition}{note}{Preguntas de reflexión}
\begin{enumerate}
\sphinxsetlistlabels{\arabic}{enumi}{enumii}{}{.}%
\item {} 
\sphinxAtStartPar
¿Le genera curiosidad saber cómo viaja la información desde un dispositivo hasta un servidor en otro continente? ¿Por qué?

\end{enumerate}


\bigskip\hrule\bigskip

\begin{enumerate}
\sphinxsetlistlabels{\arabic}{enumi}{enumii}{}{.}%
\setcounter{enumi}{1}
\item {} 
\sphinxAtStartPar
Cuando el internet falla en casa, ¿Se siente motivado por intentar diagnosticar si es un problema del módem, del cableado o de la señal inalámbrica? ¿Por qué?

\end{enumerate}


\bigskip\hrule\bigskip

\begin{enumerate}
\sphinxsetlistlabels{\arabic}{enumi}{enumii}{}{.}%
\setcounter{enumi}{2}
\item {} 
\sphinxAtStartPar
¿Le llama la atención la idea de configurar dispositivos para que otras personas puedan comunicarse de forma segura dentro de una organización (Intranet)? ¿Por qué?

\end{enumerate}


\bigskip\hrule\bigskip

\begin{enumerate}
\sphinxsetlistlabels{\arabic}{enumi}{enumii}{}{.}%
\setcounter{enumi}{3}
\item {} 
\sphinxAtStartPar
¿Prefiere entender cómo funciona un “Switch” o un “Router” por dentro en lugar de simplemente usarlos como cajas negras que dan internet? ¿Por qué?

\end{enumerate}


\bigskip\hrule\bigskip

\begin{enumerate}
\sphinxsetlistlabels{\arabic}{enumi}{enumii}{}{.}%
\setcounter{enumi}{4}
\item {} 
\sphinxAtStartPar
¿Le parece interesante el concepto de “Hacking Ético” como una forma de ayudar a las empresas a protegerse, en lugar de ser solo un hacker malicioso? ¿Por qué?

\end{enumerate}


\bigskip\hrule\bigskip

\begin{enumerate}
\sphinxsetlistlabels{\arabic}{enumi}{enumii}{}{.}%
\setcounter{enumi}{5}
\item {} 
\sphinxAtStartPar
¿Siente curiosidad por aprender a detectar vulnerabilidades en páginas web (como inyecciones SQL o XSS) para poder corregirlas? ¿Por qué?

\end{enumerate}


\bigskip\hrule\bigskip

\begin{enumerate}
\sphinxsetlistlabels{\arabic}{enumi}{enumii}{}{.}%
\setcounter{enumi}{6}
\item {} 
\sphinxAtStartPar
¿Le interesaría aprender a administrar servidores de forma remota, trabajando mediante líneas de comandos y protocolos seguros como SSH? ¿Por qué?

\end{enumerate}


\bigskip\hrule\bigskip

\begin{enumerate}
\sphinxsetlistlabels{\arabic}{enumi}{enumii}{}{.}%
\setcounter{enumi}{7}
\item {} 
\sphinxAtStartPar
¿Le motiva aprender sobre leyes y normas (como la Ley 1273 de 2009 en Colombia) para entender qué es legal y qué no en el entorno digital? ¿Por qué?

\end{enumerate}


\bigskip\hrule\bigskip

\begin{enumerate}
\sphinxsetlistlabels{\arabic}{enumi}{enumii}{}{.}%
\setcounter{enumi}{8}
\item {} 
\sphinxAtStartPar
¿Le gustaría trabajar en el diseño de planes de seguridad para proteger activos críticos, asegurando que la información sea confidencial, íntegra y disponible? ¿Por qué?

\end{enumerate}


\bigskip\hrule\bigskip

\begin{enumerate}
\sphinxsetlistlabels{\arabic}{enumi}{enumii}{}{.}%
\setcounter{enumi}{9}
\item {} 
\sphinxAtStartPar
¿Le gustaría la idea de estar constantemente actualizado sobre nuevas amenazas informáticas para implementar medidas de defensa como Firewalls o cifrados avanzados? ¿Por qué?

\end{enumerate}
\end{sphinxadmonition}

\sphinxstepscope


\section{Guía 1004: MANTENIMIENTO (P2)}
\label{\detokenize{1_guias/10informatica/G1004:guia-1004-mantenimiento-p2}}\label{\detokenize{1_guias/10informatica/G1004::doc}}

\begin{savenotes}\sphinxattablestart
\sphinxthistablewithglobalstyle
\centering
\begin{tabular}[t]{\X{35}{100}\X{65}{100}}
\sphinxtoprule
\sphinxtableatstartofbodyhook

\begin{savenotes}\sphinxattablestart
\sphinxthistablewithglobalstyle
\centering
\begin{tabulary}{\linewidth}[t]{T}
\sphinxtoprule
\sphinxtableatstartofbodyhook
\sphinxAtStartPar
\sphinxincludegraphics{{escudo}.png}
\\
\sphinxbottomrule
\end{tabulary}
\sphinxtableafterendhook\par
\sphinxattableend\end{savenotes}
&

\begin{savenotes}\sphinxattablestart
\sphinxthistablewithglobalstyle
\raggedright
\begin{tabular}[t]{*{1}{\X{1}{1}}}
\sphinxtoprule
\sphinxtableatstartofbodyhook
\sphinxAtStartPar
\sphinxstylestrong{INSTITUCIÓN EDUCATIVA}: John F. Kennedy
\begin{itemize}
\item {} 
\sphinxAtStartPar
\sphinxstylestrong{DOCENTE}: Julian Camilo Fonseca Romero

\item {} 
\sphinxAtStartPar
\sphinxstylestrong{ASIGNATURA}: Tecnología e Informática

\item {} 
\sphinxAtStartPar
\sphinxstylestrong{GRADO}: DÉCIMO

\item {} 
\sphinxAtStartPar
\sphinxstylestrong{PERIODO}: 02

\item {} 
\sphinxAtStartPar
\sphinxstylestrong{TIEMPO}: 10 semanas (2 horas / semana)

\item {} 
\sphinxAtStartPar
\sphinxstylestrong{TIEMPO DE TRABAJO}: 8/20 horas

\end{itemize}
\\
\sphinxbottomrule
\end{tabular}
\sphinxtableafterendhook\par
\sphinxattableend\end{savenotes}
\\
\sphinxbottomrule
\end{tabular}
\sphinxtableafterendhook\par
\sphinxattableend\end{savenotes}

\begin{sphinxuseclass}{sd-tab-set}
\begin{sphinxuseclass}{sd-tab-item}\subsubsection*{Objetivo}

\begin{sphinxuseclass}{sd-tab-content}
\sphinxAtStartPar
\DUrole{sd-sphinx-override}{\DUrole{sd-badge}{\DUrole{sd-bg-light}{\DUrole{sd-bg-text-light}{Exploración}}}}
\DUrole{sd-sphinx-override}{\DUrole{sd-badge}{\DUrole{sd-bg-info}{\DUrole{sd-bg-text-info}{Estructuración}}}}
\DUrole{sd-sphinx-override}{\DUrole{sd-badge}{\DUrole{sd-bg-warning}{\DUrole{sd-bg-text-warning}{Actividad práctica}}}}
\DUrole{sd-sphinx-override}{\DUrole{sd-badge}{\DUrole{sd-bg-danger}{\DUrole{sd-bg-text-danger}{Transferencia}}}}
\begin{itemize}
\item {} 
\sphinxAtStartPar
\sphinxstylestrong{Temática}: Mantenimiento.

\item {} 
\sphinxAtStartPar
\sphinxstylestrong{Objetivo de aprendizaje}: Introducir los conceptos relacionados con la temática de mantenimiento.

\end{itemize}

\end{sphinxuseclass}
\end{sphinxuseclass}
\begin{sphinxuseclass}{sd-tab-item}\subsubsection*{Orientaciones curriculares}

\begin{sphinxuseclass}{sd-tab-content}\begin{itemize}
\item {} 
\sphinxAtStartPar
\sphinxstylestrong{Componente}: Uso y apropiación de la Tecnología y la Informática.

\item {} 
\sphinxAtStartPar
\sphinxstylestrong{Competencia}: Genero propuestas innovadoras para el uso y aprovechamiento de productos tecnológicos.

\item {} 
\sphinxAtStartPar
\sphinxstylestrong{Evidencia de aprendizaje}: Diseño y aplico planes sistemáticos de mantenimiento correctivo en productos tecnológicos analógicos y digitales utilizados en la vida diaria.

\end{itemize}

\end{sphinxuseclass}
\end{sphinxuseclass}
\begin{sphinxuseclass}{sd-tab-item}\subsubsection*{Criterios de evaluación}

\begin{sphinxuseclass}{sd-tab-content}\begin{itemize}
\item {} 
\sphinxAtStartPar
Actividades desarrolladas en su totalidad.

\item {} 
\sphinxAtStartPar
Participación en clase.

\item {} 
\sphinxAtStartPar
Explicación clara de conceptos.

\item {} 
\sphinxAtStartPar
Argumentación de ideas.

\item {} 
\sphinxAtStartPar
Cumplimiento de las reglas del aula.

\end{itemize}

\end{sphinxuseclass}
\end{sphinxuseclass}
\begin{sphinxuseclass}{sd-tab-item}\subsubsection*{Recursos (0)}

\begin{sphinxuseclass}{sd-tab-content}\begin{itemize}
\item {} 
\sphinxAtStartPar
Esta guía no contiene recursos.

\end{itemize}

\end{sphinxuseclass}
\end{sphinxuseclass}
\end{sphinxuseclass}

\subsection{Mensaje del autor}
\label{\detokenize{1_guias/10informatica/G1004:mensaje-del-autor}}\begin{quote}

\sphinxAtStartPar
“En esta guía el estudiante descubre que el mantenimiento es mucho más que ‘arreglar cosas’; es una habilidad fundamental para preservar la tecnología que nos rodea, optimizar recursos y garantizar la seguridad digital. En esta guía se aprende a diferenciar entre un mantenimiento preventivo y uno correctivo, así como también las características básica de la electricidad y la electrónica. También encuentra una explicación de hardware libre con enfoque en ARUDINO y Raspberry Pi”.

\sphinxAtStartPar
“Al final de la guía encontrará unas preguntas de reflexión para identificar si el tema de \sphinxstylestrong{mantenimiento} y \sphinxstylestrong{electrónica} le genera interés. Una de estas temáticas podría ser su elección para su proyecto de grado”.

\begin{flushright}
---Camilo Fonseca
\end{flushright}
\end{quote}


\bigskip\hrule\bigskip



\subsection{A00: Introducción (+0,1)}
\label{\detokenize{1_guias/10informatica/G1004:a00-introduccion-0-1}}
\sphinxAtStartPar
\DUrole{sd-sphinx-override}{\DUrole{sd-badge}{\DUrole{sd-bg-light}{\DUrole{sd-bg-text-light}{Exploración}}}}
\begin{enumerate}
\sphinxsetlistlabels{\arabic}{enumi}{enumii}{}{.}%
\item {} 
\sphinxAtStartPar
Transcriba en su cuaderno el objetivo, las orientaciones curriculares y los criterios de evaluación de la presente guía.

\end{enumerate}


\bigskip\hrule\bigskip



\subsection{A01: Mantenimiento (+0,1)}
\label{\detokenize{1_guias/10informatica/G1004:a01-mantenimiento-0-1}}
\sphinxAtStartPar
\DUrole{sd-sphinx-override}{\DUrole{sd-badge}{\DUrole{sd-bg-info}{\DUrole{sd-bg-text-info}{Estructuración}}}}
\begin{enumerate}
\sphinxsetlistlabels{\arabic}{enumi}{enumii}{}{.}%
\item {} 
\sphinxAtStartPar
Lea los siguientes conceptos:

\end{enumerate}


\bigskip\hrule\bigskip



\subsubsection{¿Qué es el mantenimiento?}
\label{\detokenize{1_guias/10informatica/G1004:que-es-el-mantenimiento}}
\begin{sphinxadmonition}{note}{Concepto de mantenimiento}

\sphinxAtStartPar
El \sphinxstylestrong{mantenimiento} se refiere al conjunto de actividades y procesos destinados a conservar, reparar o mejorar un objeto, sistema o infraestructura. Su objetivo principal es asegurar que un bien funcione adecuadamente y prolongar su vida útil.


\begin{savenotes}\sphinxattablestart
\sphinxthistablewithglobalstyle
\centering
\begin{tabulary}{\linewidth}[t]{T}
\sphinxtoprule
\sphinxtableatstartofbodyhook

\begin{sphinxfigure-in-table}
\centering
\capstart
\noindent\sphinxincludegraphics[width=400\sphinxpxdimen]{{image014}.png}
\sphinxfigcaption{Torre de PC}\label{\detokenize{1_guias/10informatica/G1004:id1}}\end{sphinxfigure-in-table}\relax
\\
\sphinxbottomrule
\end{tabulary}
\sphinxtableafterendhook\par
\sphinxattableend\end{savenotes}

\sphinxAtStartPar
Los mantenimiento son importantes debido a las siguientes razones:
\begin{itemize}
\item {} 
\sphinxAtStartPar
\sphinxstylestrong{Prolongar la vida útil:} Permite mantener los equipos y dispositivos funcionando correctamente durante más tiempo, maximizando la inversión realizada.

\item {} 
\sphinxAtStartPar
\sphinxstylestrong{Prevención de problemas:} Ayuda a identificar y resolver problemas antes de que causen daños mayores como el daño en una persona o pérdida de información sensible.

\item {} 
\sphinxAtStartPar
\sphinxstylestrong{Ahorro económico:} Un mantenimiento adecuado evita costosas reparaciones y reemplazos prematuros de equipos.

\item {} 
\sphinxAtStartPar
\sphinxstylestrong{Seguridad digital:} El mantenimiento incluye aspectos de seguridad informática que protegen contra Malware o pérdidas de información.

\item {} 
\sphinxAtStartPar
\sphinxstylestrong{Eficiencia operativa:} Los equipos bien mantenidos funcionan de manera más eficiente, mejorando la productividad.

\end{itemize}

\sphinxAtStartPar
Los tipos de mantenimiento son:
\begin{itemize}
\item {} 
\sphinxAtStartPar
Mantenimiento preventivo.

\item {} 
\sphinxAtStartPar
Mantenimiento correctivo.

\item {} 
\sphinxAtStartPar
Mantenimiento de actualización.

\end{itemize}
\end{sphinxadmonition}


\bigskip\hrule\bigskip



\subsubsection{Mantenimiento preventivo (evitar el daño)}
\label{\detokenize{1_guias/10informatica/G1004:mantenimiento-preventivo-evitar-el-dano}}
\sphinxAtStartPar
Es ideal para:
\begin{itemize}
\item {} 
\sphinxAtStartPar
Equipos críticos cuya falla causaría grandes pérdidas.

\item {} 
\sphinxAtStartPar
Sistemas que requieren alta confiabilidad.

\item {} 
\sphinxAtStartPar
Cuando el costo de una falla es mayor que el costo del mantenimiento.

\item {} 
\sphinxAtStartPar
Para extender la vida útil de los equipos.

\end{itemize}

\begin{sphinxadmonition}{note}{Clasificación del \sphinxstylestrong{mantenimiento preventivo}}

\sphinxAtStartPar
El mantenimiento preventivo tiene la siguiente clasificación:
\begin{itemize}
\item {} 
\sphinxAtStartPar
\sphinxstylestrong{Mantenimiento preventivo programado:} Cuando el mantenimiento se efectúa automáticamente, en función del tiempo de vida transcurrido.

\item {} 
\sphinxAtStartPar
\sphinxstylestrong{Mantenimiento preventivo predictivo:} Es aquel que se realiza cuando se ha ido revisando periódicamente el equipo, de manera que se puede anticipar cuando va a ocurrir un fallo, haciendo en ese momento la respectiva reparación.

\item {} 
\sphinxAtStartPar
\sphinxstylestrong{Mantenimiento preventivo de oportunidad:} Es el mantenimiento que se desarrolla aprovechando que el equipo no está siendo utilizado, por ejemplo, cuando se para la actividad en una temporada de baja demanda. De ese modo, se evita que se tenga que detener la producción en momentos donde sería inoportuno y más costoso. Si el equipo dejara de funcionar en una coyuntura de alta demanda, la empresa tendría que alquilar otra maquinaria o perdería ventas.

\end{itemize}
\end{sphinxadmonition}


\bigskip\hrule\bigskip



\subsubsection{Mantenimiento correctivo (el daño ocurre)}
\label{\detokenize{1_guias/10informatica/G1004:mantenimiento-correctivo-el-dano-ocurre}}
\sphinxAtStartPar
Es el mantenimiento que se realiza después de que ocurre una falla para reparar el equipo y devolverlo a su estado operativo.

\sphinxAtStartPar
Es apropiado para:
\begin{itemize}
\item {} 
\sphinxAtStartPar
Equipos no críticos cuya falla no afecta significativamente la operación.

\item {} 
\sphinxAtStartPar
Cuando el costo de la falla es menor que el costo del mantenimiento preventivo.

\item {} 
\sphinxAtStartPar
Componentes con patrones de falla aleatorios.

\item {} 
\sphinxAtStartPar
Situaciones donde la reparación es simple y rápida.

\end{itemize}

\begin{sphinxadmonition}{note}{Clasifición del \sphinxstylestrong{mantenimiento correctivo}}

\sphinxAtStartPar
El mantenimiento correctivo tiene la siguiente clasificación:
\begin{itemize}
\item {} 
\sphinxAtStartPar
\sphinxstylestrong{Mantenimiento correctivo inmediato:} Es aquel que se realiza en el mismo momento en el que se identifica el daño.

\item {} 
\sphinxAtStartPar
\sphinxstylestrong{Mantenimiento correctivo diferido:} Cuando se detiene la actividad del elemento afectado, pudiendo luego efectuarse la reparación correspondiente.

\end{itemize}
\end{sphinxadmonition}


\bigskip\hrule\bigskip



\subsubsection{Mantenimiento de actualización}
\label{\detokenize{1_guias/10informatica/G1004:mantenimiento-de-actualizacion}}
\sphinxAtStartPar
Se refiere a las inversiones necesarias frente a la obsolescencia tecnológica. Por ejemplo, puede tratarse de la instalación de un software o hardware que potencia el rendimiento de los ordenadores.


\bigskip\hrule\bigskip



\subsubsection{¿Qué tipo de mantenimiento elegir?}
\label{\detokenize{1_guias/10informatica/G1004:que-tipo-de-mantenimiento-elegir}}
\begin{sphinxadmonition}{note}{El mantenimiento a elegir depende de ciertos factores}

\sphinxAtStartPar
La elección entre un \sphinxstylestrong{mantenimiento preventivo}, \sphinxstylestrong{correctivo} o de \sphinxstylestrong{actualización} depende de factores como:
\begin{itemize}
\item {} 
\sphinxAtStartPar
Criticidad del equipo.

\item {} 
\sphinxAtStartPar
Costos asociados.

\item {} 
\sphinxAtStartPar
Impacto de las fallas.

\item {} 
\sphinxAtStartPar
Recursos disponibles.

\end{itemize}

\sphinxAtStartPar
En la práctica, muchas organizaciones implementan una combinación de ambos tipos de mantenimiento según sus necesidades específicas.
\end{sphinxadmonition}


\bigskip\hrule\bigskip



\subsubsection{¿Qué es la electricidad?}
\label{\detokenize{1_guias/10informatica/G1004:que-es-la-electricidad}}
\begin{sphinxadmonition}{note}{Concepto de electricidad}

\sphinxAtStartPar
La \sphinxstylestrong{electricidad} es una forma de energía basada en el movimiento de cargas eléctricas (electrones) a través de materiales conductores. Es un fenómeno físico fundamental que está presente tanto en la naturaleza como en la tecnología moderna. Maneja voltajes y corrientes altos.


\begin{savenotes}\sphinxattablestart
\sphinxthistablewithglobalstyle
\centering
\begin{tabulary}{\linewidth}[t]{T}
\sphinxtoprule
\sphinxtableatstartofbodyhook

\begin{sphinxfigure-in-table}
\centering
\capstart
\noindent\sphinxincludegraphics[width=400\sphinxpxdimen]{{image022}.png}
\sphinxfigcaption{Flujo de electrones a través de una PILA}\label{\detokenize{1_guias/10informatica/G1004:id2}}\end{sphinxfigure-in-table}\relax
\\
\sphinxbottomrule
\end{tabulary}
\sphinxtableafterendhook\par
\sphinxattableend\end{savenotes}
\end{sphinxadmonition}

\sphinxAtStartPar
Las aplicaciones de la electricidad son:
\begin{itemize}
\item {} 
\sphinxAtStartPar
\sphinxstylestrong{Funcionamiento de dispositivos:} Todos los dispositivos electrónicos requieren electricidad para funcionar.

\item {} 
\sphinxAtStartPar
\sphinxstylestrong{Procesamiento de datos:} La electricidad permite la transmisión y procesamiento de información digital.

\item {} 
\sphinxAtStartPar
\sphinxstylestrong{Infraestructura tecnológica:} Las redes de comunicación y centros de datos dependen del suministro eléctrico.

\item {} 
\sphinxAtStartPar
\sphinxstylestrong{Innovación:} El desarrollo de nuevas tecnologías está directamente relacionado con avances en el campo eléctrico.

\item {} 
\sphinxAtStartPar
\sphinxstylestrong{Eficiencia energética:} La gestión eficiente de la electricidad es crucial para la sostenibilidad tecnológica.

\end{itemize}
\begin{quote}

\sphinxAtStartPar
La \sphinxstylestrong{electricidad} se enfoca en la generación, transmisión y distribución de energía y trabaja principalmente con corriente alterna (AC).
\end{quote}

\sphinxAtStartPar
Algunos ejemplos de la electricidad son las \sphinxstyleemphasis{instalaciones eléctricas}, \sphinxstyleemphasis{motores eléctricos} y los \sphinxstyleemphasis{transformadores}.


\begin{savenotes}\sphinxattablestart
\sphinxthistablewithglobalstyle
\centering
\begin{tabulary}{\linewidth}[t]{T}
\sphinxtoprule
\sphinxtableatstartofbodyhook

\begin{sphinxfigure-in-table}
\centering
\capstart
\noindent\sphinxincludegraphics[width=300\sphinxpxdimen]{{image032}.png}
\sphinxfigcaption{Motor eléctrico}\label{\detokenize{1_guias/10informatica/G1004:id3}}\end{sphinxfigure-in-table}\relax
\\
\sphinxbottomrule
\end{tabulary}
\sphinxtableafterendhook\par
\sphinxattableend\end{savenotes}

\sphinxAtStartPar
Los componentes de la electricidad son:
\begin{itemize}
\item {} 
\sphinxAtStartPar
\sphinxstylestrong{Conductores:} Materiales que permiten el flujo de corriente eléctrica

\item {} 
\sphinxAtStartPar
\sphinxstylestrong{Aislantes:} Materiales que impiden el paso de la corriente eléctrica

\item {} 
\sphinxAtStartPar
\sphinxstylestrong{Interruptores:} Dispositivos para abrir o cerrar circuitos eléctricos

\item {} 
\sphinxAtStartPar
\sphinxstylestrong{Fusibles:} Dispositivos de protección contra sobre\sphinxhyphen{}corrientes

\item {} 
\sphinxAtStartPar
\sphinxstylestrong{Transformadores:} Modifican los niveles de voltaje

\item {} 
\sphinxAtStartPar
\sphinxstylestrong{Generadores:} Producen energía eléctrica

\item {} 
\sphinxAtStartPar
\sphinxstylestrong{Motores:} Convierten energía eléctrica en energía mecánica

\item {} 
\sphinxAtStartPar
\sphinxstylestrong{Tableros eléctricos:} Dispositivos de distribución y control

\end{itemize}


\bigskip\hrule\bigskip



\subsubsection{¿Qué es la electrónica?}
\label{\detokenize{1_guias/10informatica/G1004:que-es-la-electronica}}
\begin{sphinxadmonition}{note}{Concepto de electrónica}

\sphinxAtStartPar
La electrónica se ocupa del control del flujo de electrones a menor escala. Maneja voltajes y corrientes más pequeños.
\end{sphinxadmonition}
\begin{quote}

\sphinxAtStartPar
La \sphinxstylestrong{electrónica} se enfoca en el procesamiento de señales e información y trabaja principalmente con corriente continua (DC).
\end{quote}

\sphinxAtStartPar
Algunos ejemplos de la electrónica son las \sphinxstyleemphasis{computadoras}, \sphinxstyleemphasis{celulares} y \sphinxstyleemphasis{circuitos integrados}.


\begin{savenotes}\sphinxattablestart
\sphinxthistablewithglobalstyle
\centering
\begin{tabulary}{\linewidth}[t]{T}
\sphinxtoprule
\sphinxtableatstartofbodyhook

\begin{sphinxfigure-in-table}
\centering
\capstart
\noindent\sphinxincludegraphics[width=300\sphinxpxdimen]{{image042}.png}
\sphinxfigcaption{Microprocesador}\label{\detokenize{1_guias/10informatica/G1004:id4}}\end{sphinxfigure-in-table}\relax
\\
\sphinxbottomrule
\end{tabulary}
\sphinxtableafterendhook\par
\sphinxattableend\end{savenotes}

\sphinxAtStartPar
Los componentes de la electrónica son:
\begin{itemize}
\item {} 
\sphinxAtStartPar
\sphinxstylestrong{Resistencias:} Controlan el flujo de corriente

\item {} 
\sphinxAtStartPar
\sphinxstylestrong{Capacitores:} Almacenan carga eléctrica

\item {} 
\sphinxAtStartPar
\sphinxstylestrong{Inductores:} Almacenan energía en campo magnético

\item {} 
\sphinxAtStartPar
\sphinxstylestrong{Diodos:} Permiten el flujo de corriente en una dirección

\item {} 
\sphinxAtStartPar
\sphinxstylestrong{Transistores:} Amplifican o conmutan señales

\item {} 
\sphinxAtStartPar
\sphinxstylestrong{Circuitos integrados:} Combinan múltiples componentes

\item {} 
\sphinxAtStartPar
\sphinxstylestrong{Sensores:} Detectan cambios en el entorno

\item {} 
\sphinxAtStartPar
\sphinxstylestrong{Micro\sphinxhyphen{}controladores:} Procesan señales y ejecutan programas

\end{itemize}


\bigskip\hrule\bigskip



\subsubsection{Electrónica digital}
\label{\detokenize{1_guias/10informatica/G1004:electronica-digital}}

\begin{savenotes}\sphinxattablestart
\sphinxthistablewithglobalstyle
\centering
\begin{tabulary}{\linewidth}[t]{T}
\sphinxtoprule
\sphinxtableatstartofbodyhook

\begin{sphinxfigure-in-table}
\centering
\capstart
\noindent\sphinxincludegraphics[width=500\sphinxpxdimen]{{image052}.png}
\sphinxfigcaption{Elementos electrónica digital}\label{\detokenize{1_guias/10informatica/G1004:id5}}\end{sphinxfigure-in-table}\relax
\\
\sphinxbottomrule
\end{tabulary}
\sphinxtableafterendhook\par
\sphinxattableend\end{savenotes}

\sphinxAtStartPar
La \sphinxstyleemphasis{electrónica digital} es una rama de la electrónica que se enfoca en el procesamiento de señales y datos utilizando señales discretas, generalmente representadas como combinaciones de ceros y unos (bits).

\sphinxAtStartPar
A diferencia de la \sphinxstyleemphasis{electrónica analógica}, que trabaja con señales continuas, la \sphinxstyleemphasis{electrónica digital} opera con información en formato digital, lo que la hace especialmente adecuada para el procesamiento, almacenamiento y transmisión de información de manera precisa y confiable.

\sphinxAtStartPar
Los elementos de la electrónica digital son:
\begin{itemize}
\item {} 
\sphinxAtStartPar
\sphinxstylestrong{BITS:} El bit es la unidad más básica de información en electrónica digital y puede tener dos valores: 0 o 1. Los bits se utilizan para representar información binaria, como números, caracteres, imágenes, etc.

\item {} 
\sphinxAtStartPar
\sphinxstylestrong{Compuertas lógicas:} Las compuertas lógicas son circuitos electrónicos que realizan operaciones lógicas en señales binarias. Las compuertas lógicas más comunes incluyen: AND, OR, NOT y XOR

\item {} 
\sphinxAtStartPar
\sphinxstylestrong{Circuitos de combinación:} Estos circuitos están formados por compuertas lógicas interconectadas y no tienen memoria. Su salida depende únicamente de las entradas en ese momento.

\item {} 
\sphinxAtStartPar
\sphinxstylestrong{Circuitos secuenciales:}   A diferencia de los circuitos de combinación, los circuitos secuenciales tienen memoria. Utilizan elementos como flip\sphinxhyphen{}flops y registros para almacenar información y mantener un estado interno. Ejemplos de circuitos secuenciales incluyen contadores y máquinas de estados finitos.

\item {} 
\sphinxAtStartPar
\sphinxstylestrong{Flip\sphinxhyphen{}flops:} Los flip\sphinxhyphen{}flops son dispositivos de almacenamiento que contienen un solo bit y se utilizan para mantener el estado en circuitos secuenciales. Los tipos más comunes de flip\sphinxhyphen{}flops incluyen el flip\sphinxhyphen{}flop D, el flip\sphinxhyphen{}flop T, el flip\sphinxhyphen{}flop JK y el flip\sphinxhyphen{}flop SR.

\item {} 
\sphinxAtStartPar
\sphinxstylestrong{Multiplexores y Demultiplexores:} Los (multi\sphinxhyphen{}plexores) son utilizados para seleccionar una de varias entradas y enrutarla a una sola salida. Los demultiplexores hacen lo contrario, tomando una entrada y seleccionando una de varias salidas posibles.

\item {} 
\sphinxAtStartPar
\sphinxstylestrong{Decodificadores y Codificadores:} Los decodificadores se utilizan para convertir una entrada binaria en una selección de una de varias salidas posibles. Los codificadores hacen lo contrario, convierten una selección en una representación binaria.

\item {} 
\sphinxAtStartPar
\sphinxstylestrong{Contadores:}  Los contadores son circuitos secuenciales que generan una secuencia de números binarios en respuesta a un reloj de entrada. Se utilizan en una variedad de aplicaciones, como la medición del tiempo y la creación de divisores de frecuencia.

\item {} 
\sphinxAtStartPar
\sphinxstylestrong{Memorias:} Las memorias digitales se utilizan para almacenar datos en sistemas digitales. Esto puede incluir RAM (memoria de acceso aleatorio) para almacenamiento temporal y ROM (memoria de solo lectura) para almacenamiento permanente. También existen tipos de memoria intermedios, como las memorias Flash.

\item {} 
\sphinxAtStartPar
\sphinxstylestrong{Convertidores Analógico\sphinxhyphen{}Digitales (ADC) y Digitales\sphinxhyphen{}Analógicos (DAC):} Estos dispositivos permiten la conversión de señales entre dominios analógicos y digitales. Los ADC convierten señales analógicas en digitales, mientras que los DAC hacen lo contrario.

\item {} 
\sphinxAtStartPar
\sphinxstylestrong{Puertos de Entrada/Salida (E/S):} Estos son interfaces que permiten que un sistema digital se comunique con el mundo exterior, ya sea para recibir datos (E/S de entrada) o enviar datos (E/S de salida).

\item {} 
\sphinxAtStartPar
\sphinxstylestrong{Micro\sphinxhyphen{}controladores y Micro\sphinxhyphen{}procesadores:} Estos son dispositivos complejos que incorporan CPU (unidad central de procesamiento), memoria y periféricos en un solo chip. Se utilizan en una gran variedad de aplicaciones, desde electrodomésticos hasta sistemas embebidos.

\end{itemize}


\bigskip\hrule\bigskip



\subsubsection{Diferencia entre microprocesador y microcontrolador}
\label{\detokenize{1_guias/10informatica/G1004:diferencia-entre-microprocesador-y-microcontrolador}}
\begin{sphinxadmonition}{note}{Un micro\sphinxhyphen{}procesador requiere periféricos para funcionar. Un micro\sphinxhyphen{}controlador es un elemento completo}

\sphinxAtStartPar
Un \sphinxstylestrong{micro\sphinxhyphen{}procesador} es un elemento que realiza operaciones lógico aritméticas. No dispone de entradas y salidas como un \sphinxstylestrong{micro\sphinxhyphen{}controlador}. Requiere de más periféricos adicionales para funcionar, como memorias o controladores de bus. Sin embargo, son más veloces al realizar estas operaciones que un \sphinxstylestrong{micro\sphinxhyphen{}controlador}.


\begin{savenotes}\sphinxattablestart
\sphinxthistablewithglobalstyle
\centering
\begin{tabulary}{\linewidth}[t]{T}
\sphinxtoprule
\sphinxtableatstartofbodyhook

\begin{sphinxfigure-in-table}
\centering
\capstart
\noindent\sphinxincludegraphics[width=320\sphinxpxdimen]{{image042}.png}
\sphinxfigcaption{Microprocesador}\label{\detokenize{1_guias/10informatica/G1004:id6}}\end{sphinxfigure-in-table}\relax
\\
\sphinxbottomrule
\end{tabulary}
\sphinxtableafterendhook\par
\sphinxattableend\end{savenotes}

\sphinxAtStartPar
Un \sphinxstylestrong{micro\sphinxhyphen{}controlador} es un circuito integrado compuesto de entradas salidas, memoria y unidades lógico aritméticas. Son en sí, un elemento completo y funcional para realizar operaciones digitales. Comparados con un micro\sphinxhyphen{}procesador, los micro\sphinxhyphen{}controladores son más “lentos” dado que realizan menos instrucciones por segundo.


\begin{savenotes}\sphinxattablestart
\sphinxthistablewithglobalstyle
\centering
\begin{tabulary}{\linewidth}[t]{T}
\sphinxtoprule
\sphinxtableatstartofbodyhook

\begin{sphinxfigure-in-table}
\centering
\capstart
\noindent\sphinxincludegraphics[width=300\sphinxpxdimen]{{image062}.png}
\sphinxfigcaption{Microcontrolador}\label{\detokenize{1_guias/10informatica/G1004:id7}}\end{sphinxfigure-in-table}\relax
\\
\sphinxbottomrule
\end{tabulary}
\sphinxtableafterendhook\par
\sphinxattableend\end{savenotes}
\end{sphinxadmonition}


\bigskip\hrule\bigskip



\subsubsection{Hardware libre}
\label{\detokenize{1_guias/10informatica/G1004:hardware-libre}}
\sphinxAtStartPar
El \sphinxstylestrong{hardware libre} son los dispositivos cuyas especificaciones y diagramas son de acceso público, de manera que cualquiera puede replicarlos. Esto quiere decir que Arduino ofrece las bases \sphinxstylestrong{para que cualquier otra persona o empresa pueda crear sus propias placas}, pudiendo ser diferentes entre ellas pero igualmente funcionales al partir de la misma base.


\paragraph{ARDUINO}
\label{\detokenize{1_guias/10informatica/G1004:arduino}}
\begin{sphinxadmonition}{note}{Orientado a micro\sphinxhyphen{}controladores}


\begin{savenotes}\sphinxattablestart
\sphinxthistablewithglobalstyle
\centering
\begin{tabulary}{\linewidth}[t]{T}
\sphinxtoprule
\sphinxtableatstartofbodyhook

\begin{sphinxfigure-in-table}
\centering
\capstart
\noindent\sphinxincludegraphics[width=500\sphinxpxdimen]{{image072}.png}
\sphinxfigcaption{Placa ARDUINO}\label{\detokenize{1_guias/10informatica/G1004:id8}}\end{sphinxfigure-in-table}\relax
\\
\sphinxbottomrule
\end{tabulary}
\sphinxtableafterendhook\par
\sphinxattableend\end{savenotes}

\sphinxAtStartPar
\sphinxstylestrong{Arduino} es una plataforma de creación de electrónica de código abierto, la cual está basada en hardware y software libre, flexible y fácil de utilizar para los creadores y desarrolladores.

\sphinxAtStartPar
Esta plataforma permite crear diferentes tipos de micro\sphinxhyphen{}controladores de una sola placa a los que la comunidad de creadores puede darles diferentes tipos de uso.

\sphinxAtStartPar
Los \sphinxstylestrong{micro\sphinxhyphen{}controladores} son circuitos integrados en los que se pueden grabar instrucciones, las cuales usted escribe con el lenguaje de programación que se utiliza en el entorno \sphinxstylestrong{Arduino IDE}. Estas instrucciones permiten crear programas que interactúan con los circuitos de la placa.

\sphinxAtStartPar
\sphinxstylestrong{Arduino} posee lo que se llama una interfaz de entrada, que es una conexión en la que se puede conectar en la placa diferentes tipos de periféricos. La información de estos periféricos que usted conecte se trasladará al micro\sphinxhyphen{}controlador, el cual se encargará de procesar los datos que le lleguen a través de ellos.

\sphinxAtStartPar
El tipo de periféricos que usted pueda utilizar para enviar datos al micro\sphinxhyphen{}controlador depende en gran medida de qué uso le esté pensando dar. Pueden ser:
\begin{itemize}
\item {} 
\sphinxAtStartPar
Cámaras para obtener imágenes

\item {} 
\sphinxAtStartPar
Teclados para introducir datos

\item {} 
\sphinxAtStartPar
Diferentes tipos de sensores.

\end{itemize}

\sphinxAtStartPar
También cuenta con una interfaz de salida, que es la que se encarga de llevar la información que se ha procesado en el Arduino a otros periféricos. Estos periféricos pueden ser:
\begin{itemize}
\item {} 
\sphinxAtStartPar
Pantallas

\item {} 
\sphinxAtStartPar
Altavoces en los que reproducir los datos procesados

\item {} 
\sphinxAtStartPar
Otras placas o controladores.

\end{itemize}
\end{sphinxadmonition}


\paragraph{Raspberry PI}
\label{\detokenize{1_guias/10informatica/G1004:raspberry-pi}}
\begin{sphinxadmonition}{note}{Orientado a microprocesadores}


\begin{savenotes}\sphinxattablestart
\sphinxthistablewithglobalstyle
\centering
\begin{tabulary}{\linewidth}[t]{T}
\sphinxtoprule
\sphinxtableatstartofbodyhook

\begin{sphinxfigure-in-table}
\centering
\capstart
\noindent\sphinxincludegraphics[width=400\sphinxpxdimen]{{image082}.png}
\sphinxfigcaption{Placa RASPBERRY PI 5}\label{\detokenize{1_guias/10informatica/G1004:id9}}\end{sphinxfigure-in-table}\relax
\\
\sphinxbottomrule
\end{tabulary}
\sphinxtableafterendhook\par
\sphinxattableend\end{savenotes}

\sphinxAtStartPar
Una \sphinxstylestrong{Raspberry PI} es un \sphinxstylestrong{ordenador} del tamaño de una tarjeta de crédito. Consiste en una placa base que soporta distintos componentes de un ordenador como un micro\sphinxhyphen{}procesador ARM de hasta 1500 MHz, un chip gráfico y una memoria RAM de hasta 8 GB.

\sphinxAtStartPar
Gracias a sus puertos y entradas, permite conectar dispositivos periféricos. Por ejemplo:
\begin{itemize}
\item {} 
\sphinxAtStartPar
Una pantalla táctil.

\item {} 
\sphinxAtStartPar
Un teclado e incluso un televisor.

\end{itemize}

\sphinxAtStartPar
Contiene un procesador gráfico VideoCoreIV, con lo que permite la reproducción de vídeo “incluso en alta definición”. Permite la conexión a la red a través del puerto de Ethernet, y algunos modelos permiten conexión Wifi y Bluetooth. Consta de una ranura SD que permite instalar (a través de una tarjeta microSD) sistemas operativos libres basados en LINUX.
\end{sphinxadmonition}


\bigskip\hrule\bigskip



\subsection{A02: Taller (+0,8)}
\label{\detokenize{1_guias/10informatica/G1004:a02-taller-0-8}}
\sphinxAtStartPar
\DUrole{sd-sphinx-override}{\DUrole{sd-badge}{\DUrole{sd-bg-warning}{\DUrole{sd-bg-text-warning}{Actividad práctica}}}}
\begin{enumerate}
\sphinxsetlistlabels{\arabic}{enumi}{enumii}{}{.}%
\item {} 
\sphinxAtStartPar
Dibuje la siguiente pirámide en su cuaderno. Divídala por secciones y de menor a mayor coloque (según su criterio) 5 razones del por que son importantes los mantenimientos. La base de la pirámide es la razón menos importante.

\end{enumerate}


\begin{savenotes}\sphinxattablestart
\sphinxthistablewithglobalstyle
\centering
\begin{tabulary}{\linewidth}[t]{T}
\sphinxtoprule
\sphinxtableatstartofbodyhook

\begin{sphinxfigure-in-table}
\centering
\capstart
\noindent\sphinxincludegraphics[width=500\sphinxpxdimen]{{image092}.png}
\sphinxfigcaption{Nivel de importancia}\label{\detokenize{1_guias/10informatica/G1004:id10}}\end{sphinxfigure-in-table}\relax
\\
\sphinxbottomrule
\end{tabulary}
\sphinxtableafterendhook\par
\sphinxattableend\end{savenotes}
\begin{enumerate}
\sphinxsetlistlabels{\arabic}{enumi}{enumii}{}{.}%
\setcounter{enumi}{1}
\item {} 
\sphinxAtStartPar
En su cuaderno escriba los siguientes casos, luego especifique que tipo de mantenimiento haría por cada caso y explique por qué:
\begin{quote}

\sphinxAtStartPar
\sphinxstylestrong{CASO 01:} \sphinxstyleemphasis{“El lote de equipos de cómputo que recién llegaron a la empresa es de la misma tecnología del lote de equipos anteriores. Serán asignados a las oficinas de secretaría y tesorería. La última vez en que fallaron los equipos, los funcionarios trabajaron con su celular. Ese día al final entregamos el pedido al cliente a tiempo”}.
\end{quote}
\begin{quote}

\sphinxAtStartPar
\sphinxstylestrong{CASO 02:} \sphinxstyleemphasis{“La máquina de \sphinxstylestrong{torno} llegó hace dos días y nada que se ha puesto en \sphinxstylestrong{producción}. Estamos perdiendo \sphinxstylestrong{\$6,000,000} al día. Este rublo es considerado \sphinxstylestrong{crítico}. Se espera que la máquina esté en producción a más tardar en tres días y garantizar que funcione 24/7”}.
\end{quote}
\begin{quote}

\sphinxAtStartPar
\sphinxstylestrong{CASO 03:} \sphinxstyleemphasis{“Este año tampoco llegaron computadoras nuevas al colegio. Así que tenemos que volver a usar los equipos obsoletos de hace 8 años. Necesitamos la ayuda de los estudiantes para que cuiden estos equipos. Asimismo, para extender la vida útil de los equipos, será necesario instalar el sistema operativo \sphinxstylestrong{LINUX}”}.
\end{quote}

\item {} 
\sphinxAtStartPar
Responda las siguientes preguntas en su cuaderno:
\begin{itemize}
\item {} 
\sphinxAtStartPar
¿Cómo funciona la electricidad?

\item {} 
\sphinxAtStartPar
¿De que forma se transmite la electricidad?

\end{itemize}

\item {} 
\sphinxAtStartPar
Describa y dibuje en su cuaderno 10 dispositivos electrónicos, luego por cada uno responda:
\begin{itemize}
\item {} 
\sphinxAtStartPar
¿Que tocaría hacer si no existiera dicho dispositivo electrónico?

\end{itemize}

\item {} 
\sphinxAtStartPar
Consulte en INTERNET qué es un \sphinxstylestrong{PANEL SOLAR} y luego escríba la información en su cuaderno.

\item {} 
\sphinxAtStartPar
Consulte en INTERNET los \sphinxstylestrong{precios de compra e instalación de paneles solares} en el municipio. Luego escriba la información consultada en el cuaderno. (También escribir los enlaces a la página web donde encontró la información).

\item {} 
\sphinxAtStartPar
Complete la siguiente tabla en su cuaderno:

\end{enumerate}


\begin{savenotes}\sphinxattablestart
\sphinxthistablewithglobalstyle
\centering
\begin{tabulary}{\linewidth}[t]{TTT}
\sphinxtoprule
\sphinxstyletheadfamily 
\sphinxAtStartPar
Característica
&\sphinxstyletheadfamily 
\sphinxAtStartPar
Electricidad
&\sphinxstyletheadfamily 
\sphinxAtStartPar
Electrónica
\\
\sphinxmidrule
\sphinxtableatstartofbodyhook
\sphinxAtStartPar
\sphinxstylestrong{Definición}
&
\sphinxAtStartPar

&
\sphinxAtStartPar

\\
\sphinxhline
\sphinxAtStartPar
\sphinxstylestrong{Aplicaciones}
&
\sphinxAtStartPar

&
\sphinxAtStartPar

\\
\sphinxhline
\sphinxAtStartPar
\sphinxstylestrong{Componentes}
&
\sphinxAtStartPar

&
\sphinxAtStartPar

\\
\sphinxbottomrule
\end{tabulary}
\sphinxtableafterendhook\par
\sphinxattableend\end{savenotes}
\begin{enumerate}
\sphinxsetlistlabels{\arabic}{enumi}{enumii}{}{.}%
\setcounter{enumi}{7}
\item {} 
\sphinxAtStartPar
En INTERNET, consulte 5 ejemplos de proyectos que se pueden hacer usando \sphinxstylestrong{ARDUINO} y descríbalos en su cuaderno (incluya dibujos).

\item {} 
\sphinxAtStartPar
En INTERNET, consulte 5 ejemplos de proyectos que se pueden hacer usando una \sphinxstylestrong{Rapsberry Pi} y descríbalos en su cuaderno (incluya dibujos).

\end{enumerate}


\bigskip\hrule\bigskip



\subsection{A03: Análisis y reflexión (+4,0)}
\label{\detokenize{1_guias/10informatica/G1004:a03-analisis-y-reflexion-4-0}}
\sphinxAtStartPar
\DUrole{sd-sphinx-override}{\DUrole{sd-badge}{\DUrole{sd-bg-danger}{\DUrole{sd-bg-text-danger}{Transferencia}}}}


\begin{savenotes}\sphinxattablestart
\sphinxthistablewithglobalstyle
\centering
\begin{tabulary}{\linewidth}[t]{TT}
\sphinxtoprule
\sphinxtableatstartofbodyhook
\sphinxAtStartPar
\sphinxstylestrong{TIEMPO}
&
\sphinxAtStartPar
59 minutos
\\
\sphinxhline
\sphinxAtStartPar
\sphinxstylestrong{PREGUNTAS}
&
\sphinxAtStartPar
10
\\
\sphinxhline
\sphinxAtStartPar
\sphinxstylestrong{OBJETIVO}
&
\sphinxAtStartPar
Identificar qué le llamó la atención de la temática de \sphinxstylestrong{Mantenimiento}.
\\
\sphinxhline
\sphinxAtStartPar
\sphinxstylestrong{ACTIVIDAD}
&
\sphinxAtStartPar
En una hoja al final de su cuaderno, desarrolle un \sphinxstylestrong{texto reflexivo} utilizando las siguientes preguntas como guía para estructurar su escrito.
\\
\sphinxbottomrule
\end{tabulary}
\sphinxtableafterendhook\par
\sphinxattableend\end{savenotes}

\begin{sphinxadmonition}{note}{Preguntas de reflexión}
\begin{enumerate}
\sphinxsetlistlabels{\arabic}{enumi}{enumii}{}{.}%
\item {} 
\sphinxAtStartPar
En el mantenimiento, ¿se ve más como alguien que planifica todo para que nada falle (preventivo) o como la persona que llega a salvar el día cuando algo ya se rompió (correctivo)? (¿Por qué?)

\end{enumerate}


\bigskip\hrule\bigskip

\begin{enumerate}
\sphinxsetlistlabels{\arabic}{enumi}{enumii}{}{.}%
\setcounter{enumi}{1}
\item {} 
\sphinxAtStartPar
¿Le parece valioso aprender a revisar un equipo periódicamente para predecir cuándo va a fallar antes de que ocurra un desastre? (¿Por qué?)

\end{enumerate}


\bigskip\hrule\bigskip

\begin{enumerate}
\sphinxsetlistlabels{\arabic}{enumi}{enumii}{}{.}%
\setcounter{enumi}{2}
\item {} 
\sphinxAtStartPar
¿Le entusiasma más la idea de trabajar con grandes energías, motores y voltajes altos (electricidad), o prefiere el mundo detallado de procesar información con voltajes pequeños (electrónica)? (¿Por qué?)

\end{enumerate}


\bigskip\hrule\bigskip

\begin{enumerate}
\sphinxsetlistlabels{\arabic}{enumi}{enumii}{}{.}%
\setcounter{enumi}{3}
\item {} 
\sphinxAtStartPar
¿Le parece interesante que toda la tecnología actual funcione mediante la combinación de ceros y unos para que sea precisa y confiable? (¿Por qué?)

\end{enumerate}


\bigskip\hrule\bigskip

\begin{enumerate}
\sphinxsetlistlabels{\arabic}{enumi}{enumii}{}{.}%
\setcounter{enumi}{4}
\item {} 
\sphinxAtStartPar
Al ver un dispositivo electrónico o eléctrico por dentro, ¿siente curiosidad de entender qué función cumplen componentes como las resistencias, los transistores o los diodos? (¿Por qué?)

\end{enumerate}


\bigskip\hrule\bigskip

\begin{enumerate}
\sphinxsetlistlabels{\arabic}{enumi}{enumii}{}{.}%
\setcounter{enumi}{5}
\item {} 
\sphinxAtStartPar
¿Le agrada la idea del hardware libre, donde usted puede ver los planos de un invento para aprender de él o incluso mejorarlo? (¿Por qué?)

\end{enumerate}


\bigskip\hrule\bigskip

\begin{enumerate}
\sphinxsetlistlabels{\arabic}{enumi}{enumii}{}{.}%
\setcounter{enumi}{6}
\item {} 
\sphinxAtStartPar
¿Le motiva aprender a escribir instrucciones para que un cerebro electrónico (micro\sphinxhyphen{}controlador) interactúe con el mundo físico y no se quede solo dentro de una pantalla? (¿Por qué?)

\end{enumerate}


\bigskip\hrule\bigskip

\begin{enumerate}
\sphinxsetlistlabels{\arabic}{enumi}{enumii}{}{.}%
\setcounter{enumi}{7}
\item {} 
\sphinxAtStartPar
¿Le genera curiosidad saber por qué un micro\sphinxhyphen{}procesador es tan veloz pero un micro\sphinxhyphen{}controlador es más práctico para realizar tareas por sí solo? (¿Por qué?)

\end{enumerate}


\bigskip\hrule\bigskip

\begin{enumerate}
\sphinxsetlistlabels{\arabic}{enumi}{enumii}{}{.}%
\setcounter{enumi}{8}
\item {} 
\sphinxAtStartPar
¿Aceptaría el reto de configurar su propio ordenador del tamaño de una tarjeta de crédito e instalar sistemas operativos como Linux? (¿Por qué?)

\end{enumerate}


\bigskip\hrule\bigskip

\begin{enumerate}
\sphinxsetlistlabels{\arabic}{enumi}{enumii}{}{.}%
\setcounter{enumi}{9}
\item {} 
\sphinxAtStartPar
¿Le llama la atención aprender sobre el hardware libre (ARDUINO o Raspberry PI) (¿Por qué?)

\end{enumerate}
\end{sphinxadmonition}

\sphinxstepscope


\section{Guía 1009: Plan de refuerzo}
\label{\detokenize{1_guias/10informatica/G1009:guia-1009-plan-de-refuerzo}}\label{\detokenize{1_guias/10informatica/G1009::doc}}\begin{quote}\begin{description}
\sphinxlineitem{ASIGNATURA}
\sphinxAtStartPar
Informática

\sphinxlineitem{GRADO}
\sphinxAtStartPar
Décimo

\end{description}\end{quote}


\bigskip\hrule\bigskip



\subsection{PRIMER PERIODO}
\label{\detokenize{1_guias/10informatica/G1009:primer-periodo}}
\begin{sphinxadmonition}{note}{EXAMEN ESCRITO: Plan de refuerzo (10 minutos)}

\sphinxAtStartPar
Estudiar los conceptos a continuación. Se realiza examen de 5 preguntas. Si el estudiante no asiste en la fecha programada, este será evaluado la próxima vez que asista a clase.
\end{sphinxadmonition}


\subsubsection{Historia del Internet}
\label{\detokenize{1_guias/10informatica/G1009:historia-del-internet}}
\sphinxAtStartPar
El Internet tiene sus orígenes en ARPANET, un proyecto militar desarrollado en 1969 por el Departamento de Defensa de Estados Unidos. Esta red inicial conectaba solo cuatro computadoras en diferentes universidades.

\sphinxAtStartPar
En los años 70 y 80 (exactamente en el año 1983), se desarrollaron los protocolos TCP/IP, que permitieron la comunicación entre diferentes redes de computadoras. La década de los 90 marcó un punto de inflexión con la creación de la World Wide Web por Tim Berners\sphinxhyphen{}Lee, que introdujo las páginas web y los navegadores en 1989.

\sphinxAtStartPar
Para facilitar el acceso a diferentes páginas web, en 1993 se creó el primer navegador web con interfaz gráfica de la historia. Este tenía por nombre \sphinxstylestrong{MOSAIC}.

\sphinxAtStartPar
El Internet moderno ha evolucionado, pasando por la era de la Web 2.0 (redes sociales e interactividad con su origen en el año 2004), hasta la actual era del Internet móvil desde el año 2007. Hoy en día la nube y el Internet de las Cosas (IoT) conectan a miles de millones de dispositivos y usuarios en todo el mundo.


\subsubsection{Historia de la computación personal}
\label{\detokenize{1_guias/10informatica/G1009:historia-de-la-computacion-personal}}
\sphinxAtStartPar
La historia de la computadora personal comienza en la década de 1970. En 1975, la Altair 8800 se convirtió en la primera computadora personal comercialmente exitosa, aunque requería ser ensamblada por el usuario.

\sphinxAtStartPar
En 1976, Apple Computer lanzó la \sphinxstylestrong{Apple I}, seguida en 1977 por la revolucionaria \sphinxstylestrong{Apple II}, una de las primeras computadoras con gráficos a color.

\sphinxAtStartPar
IBM entró al mercado en 1981 con la \sphinxstylestrong{IBM PC}, estableciendo un estándar que dominaría la industria. Los años 80 vieron la llegada de la \sphinxstylestrong{Commodore 64} (1982) y la \sphinxstylestrong{Apple Macintosh} (1984), que introdujo la interfaz gráfica de usuario (GUI) al mercado masivo.

\sphinxAtStartPar
En los 90, \sphinxstylestrong{Windows 3.0} y posteriores versiones de Microsoft democratizaron el uso de las computadoras personales.

\sphinxAtStartPar
El siglo XXI trajo las laptops ultra\sphinxhyphen{}delgadas, tablets y mini\sphinxhyphen{}computadoras. En 2012, la \sphinxstylestrong{Raspberry\sphinxhyphen{}Pi} revolucionó el concepto de computadora personal al ofrecer un dispositivo del tamaño de una tarjeta de crédito a bajo costo, enfocado en la educación y proyectos de programación.


\subsubsection{Historia del Software}
\label{\detokenize{1_guias/10informatica/G1009:historia-del-software}}
\sphinxAtStartPar
El software tiene su origen en los primeros días de la computación. \sphinxstylestrong{En la década de 1940}, las primeras computadoras ejecutaban instrucciones básicas mediante cables e interruptores físicos.

\sphinxAtStartPar
\sphinxstylestrong{Los años 50} marcaron el inicio de los lenguajes de programación modernos, con la creación del lenguaje de alto nivel \sphinxstylestrong{FORTRAN (1957)}. En el \sphinxstylestrong{año 1969} se desarrolló el sistema operativo \sphinxstylestrong{UNIX}. La mayoría de sistemas operativos en la actualidad se basan en UNIX.

\sphinxAtStartPar
Con la llegada de \sphinxstylestrong{ALTAIR en 1975,} Microsoft desarrolla el lenguaje de programación \sphinxstylestrong{BASIC}. Este permitió programar en dicha computadora.

\sphinxAtStartPar
\sphinxstylestrong{La década de 1980} fue revolucionaria con la llegada del software comercial para PC, incluyendo procesadores de texto, hojas de cálculo y los primeros sistemas operativos con interfaz gráfica. Esto inicio el \sphinxstylestrong{año 1981} con la creación del primer sistema operativo de Microsoft: \sphinxstylestrong{MS\sphinxhyphen{}DOS.}

\sphinxAtStartPar
Luego en el \sphinxstylestrong{año 1991} la comunidad de ingenieros y desarrolladores a nivel mundial crean el sistema operativo \sphinxstylestrong{LINUX}. El cual es de gran importancia hasta nuestros días.

\sphinxAtStartPar
En la actualidad, el software ha evolucionado con aplicaciones en la nube, inteligencia artificial y desarrollo ágil, transformando completamente la manera en la que interactuamos con la tecnología.


\bigskip\hrule\bigskip



\subsection{SEGUNDO PERIODO}
\label{\detokenize{1_guias/10informatica/G1009:segundo-periodo}}
\begin{sphinxadmonition}{note}{EXAMEN ESCRITO: Plan de refuerzo (10 minutos)}

\sphinxAtStartPar
Estudiar los conceptos a continuación. Se realiza examen de 5 preguntas. Si el estudiante no asiste en la fecha programada, este será evaluado la próxima vez que asista a clase.
\end{sphinxadmonition}


\subsubsection{Ofimática y Programación}
\label{\detokenize{1_guias/10informatica/G1009:ofimatica-y-programacion}}
\sphinxAtStartPar
La \sphinxstylestrong{ofimática} es el conjunto de técnicas, aplicaciones y herramientas informáticas que se utilizan para optimizar, automatizar y mejorar los procedimientos o tareas relacionadas con el trabajo de oficina. El término “ofimática” es una combinación de las palabras “oficina” e “informática”.

\sphinxAtStartPar
Las principales áreas que abarca la ofimática incluyen:
\begin{itemize}
\item {} 
\sphinxAtStartPar
Procesamiento de textos (creación y edición de documentos).

\item {} 
\sphinxAtStartPar
Hojas de cálculo (manejo de datos numéricos y fórmulas).

\item {} 
\sphinxAtStartPar
Presentaciones (creación de diapositivas).

\item {} 
\sphinxAtStartPar
Gestión de bases de datos.

\item {} 
\sphinxAtStartPar
Correo electrónico y comunicaciones.

\item {} 
\sphinxAtStartPar
Gestión de calendarios y agendas.

\item {} 
\sphinxAtStartPar
Gestión documental.

\end{itemize}

\sphinxAtStartPar
La \sphinxstylestrong{programación} se refiere específicamente a la escritura de código para crear funcionalidades dentro del software. Es una parte esencial del desarrollo, pero no abarca todo el proceso.

\sphinxAtStartPar
En la programación, las actividades a desarrollar están orientadas a la creación de código fuente. Es decir, una serie de instrucciones que se le da al PC para realizar alguna acción. También se realiza depuración y optimización de código lo cual permite que solucionen muchos errores. Es muy importante que el programador entienda la responsabilidad de salvaguardar el código fuente construido, por tanto debe sacar tiempo para tareas de copias de seguridad y control de versiones.

\sphinxAtStartPar
El \sphinxstylestrong{desarrollo de software} es un proceso que abarca la programación, pero también incluye planificación, diseño, pruebas y mantenimiento. Es un enfoque más holístico que busca crear software funcional y útil.

\sphinxAtStartPar
En el desarrollo de software se realizan tareas como el análisis de requisitos que busca identificar que funciones debe tener dicho software. También se hacen tareas de diseño donde se trabaja con bocetos iniciales para tener una idea general de como estará compuesto y de que forma funcionará el software. Otra tarea muy común es la programación la cual consiste en crear código fuente para la construcción del programa. Finalmente se realizan pruebas, implementación, mantenimiento y documentación.


\subsubsection{Redes}
\label{\detokenize{1_guias/10informatica/G1009:redes}}
\sphinxAtStartPar
Una red de datos es un sistema compuesto de dispositivos electrónicos y equipos de comunicación que permiten el intercambio de información en diferentes formatos a través de una red.

\sphinxAtStartPar
Las organizaciones y empresas dependen de la infraestructura de red para realizar muchas de sus tareas y procesos de negocio.

\sphinxAtStartPar
Los tipos de redes son:
\begin{itemize}
\item {} 
\sphinxAtStartPar
\sphinxstylestrong{RED LAN}: Es una red de área local.

\item {} 
\sphinxAtStartPar
\sphinxstylestrong{RED MAN}: Es una red de área metropolitana.

\item {} 
\sphinxAtStartPar
\sphinxstylestrong{RED WAN}: Es una red de área global (INTERNET).

\end{itemize}


\begin{savenotes}\sphinxattablestart
\sphinxthistablewithglobalstyle
\centering
\begin{tabulary}{\linewidth}[t]{T}
\sphinxtoprule
\sphinxtableatstartofbodyhook
\sphinxAtStartPar
Las \sphinxstylestrong{redes alámbricas} utilizan tecnologías de comunicación basadas en cables físicos (como cables de par trenzado, coaxiales o fibra óptica) para transmitir datos de un punto a otro, dependiendo de conexiones estructuradas y tangibles para el transporte de la información.  Este tipo de redes de telecomunicaciones ofrecen una mayor estabilidad, velocidad y seguridad en la transmisión de datos, ya que la señal no está sujeta a interferencias ambientales ni a obstáculos físicos. Son redes ideales para conectar dispositivos que requieren un alto rendimiento y están situados en una zona fija, como ordenadores de escritorio, servidores, consolas de videojuegos o sistemas de almacenamiento en red (NAS).
\\
\sphinxbottomrule
\end{tabulary}
\sphinxtableafterendhook\par
\sphinxattableend\end{savenotes}


\begin{savenotes}\sphinxattablestart
\sphinxthistablewithglobalstyle
\centering
\begin{tabulary}{\linewidth}[t]{T}
\sphinxtoprule
\sphinxtableatstartofbodyhook
\sphinxAtStartPar
Las \sphinxstylestrong{redes inalámbricas} utilizan tecnologías de comunicación basadas en señales de radiofrecuencia o microondas para transmitir datos a través del aire, eliminando la necesidad de conexiones físicas por medio de cables.  Este tipo de redes de telecomunicaciones ofrecen una mayor flexibilidad y movilidad, pues no es necesario que los dispositivos estén situados en una zona fija. Son redes ideales para conectar dispositivos móviles como smartphones, tablets, cámaras de video\sphinxhyphen{}vigilancia u ordenadores portátiles, entre otros.
\\
\sphinxbottomrule
\end{tabulary}
\sphinxtableafterendhook\par
\sphinxattableend\end{savenotes}


\subsubsection{Mantenimiento}
\label{\detokenize{1_guias/10informatica/G1009:mantenimiento}}
\sphinxAtStartPar
El mantenimiento se refiere al conjunto de actividades y procesos destinados a conservar, reparar o mejorar un objeto, sistema o infraestructura. Su objetivo principal es asegurar que un bien funcione adecuadamente y prolongar su vida útil.

\sphinxAtStartPar
Los mantenimiento son importantes debido a las siguientes razones:
\begin{itemize}
\item {} 
\sphinxAtStartPar
\sphinxstylestrong{Prolongar la vida útil}: Permite mantener los equipos y dispositivos funcionando correctamente durante más tiempo, maximizando la inversión realizada.

\item {} 
\sphinxAtStartPar
\sphinxstylestrong{Prevención de problemas}: Ayuda a identificar y resolver problemas antes de que causen daños mayores como el daño en una persona o pérdida de información sensible.

\item {} 
\sphinxAtStartPar
\sphinxstylestrong{Ahorro económico}: Un mantenimiento adecuado evita costosas reparaciones y reemplazos prematuros de equipos.

\item {} 
\sphinxAtStartPar
\sphinxstylestrong{Seguridad digital}: El mantenimiento incluye aspectos de seguridad informática que protegen contra Malware o pérdidas de información.

\item {} 
\sphinxAtStartPar
\sphinxstylestrong{Eficiencia operativa}: Los equipos bien mantenidos funcionan de manera más eficiente, mejorando la productividad.

\end{itemize}

\sphinxAtStartPar
Los tipos de mantenimiento son:
\begin{itemize}
\item {} 
\sphinxAtStartPar
Mantenimiento preventivo.

\item {} 
\sphinxAtStartPar
Mantenimiento correctivo.

\item {} 
\sphinxAtStartPar
Mantenimiento de actualización.

\end{itemize}

\sphinxstepscope


\chapter{Ofimática 10}
\label{\detokenize{1_guias/10ofimatica/G1010:ofimatica-10}}\label{\detokenize{1_guias/10ofimatica/G1010::doc}}
\sphinxstepscope


\section{Guía 1011: OFIMÁTICA (P1)}
\label{\detokenize{1_guias/10ofimatica/G1011:guia-1011-ofimatica-p1}}\label{\detokenize{1_guias/10ofimatica/G1011::doc}}

\begin{savenotes}\sphinxattablestart
\sphinxthistablewithglobalstyle
\centering
\begin{tabular}[t]{\X{35}{100}\X{65}{100}}
\sphinxtoprule
\sphinxtableatstartofbodyhook

\begin{savenotes}\sphinxattablestart
\sphinxthistablewithglobalstyle
\centering
\begin{tabulary}{\linewidth}[t]{T}
\sphinxtoprule
\sphinxtableatstartofbodyhook
\sphinxAtStartPar
\sphinxincludegraphics{{escudo}.png}
\\
\sphinxbottomrule
\end{tabulary}
\sphinxtableafterendhook\par
\sphinxattableend\end{savenotes}
&

\begin{savenotes}\sphinxattablestart
\sphinxthistablewithglobalstyle
\raggedright
\begin{tabular}[t]{*{1}{\X{1}{1}}}
\sphinxtoprule
\sphinxtableatstartofbodyhook
\sphinxAtStartPar
\sphinxstylestrong{INSTITUCIÓN EDUCATIVA}: John F. Kennedy
\begin{itemize}
\item {} 
\sphinxAtStartPar
\sphinxstylestrong{DOCENTE}: Julian Camilo Fonseca Romero

\item {} 
\sphinxAtStartPar
\sphinxstylestrong{ASIGNATURA}: Ofimática

\item {} 
\sphinxAtStartPar
\sphinxstylestrong{GRADO}: DÉCIMO

\item {} 
\sphinxAtStartPar
\sphinxstylestrong{PERIODO:} 01

\item {} 
\sphinxAtStartPar
\sphinxstylestrong{TIEMPO}: 10 semanas (2 horas / semana)

\item {} 
\sphinxAtStartPar
\sphinxstylestrong{TIEMPO DE TRABAJO}: 8/20 horas

\end{itemize}
\\
\sphinxbottomrule
\end{tabular}
\sphinxtableafterendhook\par
\sphinxattableend\end{savenotes}
\\
\sphinxbottomrule
\end{tabular}
\sphinxtableafterendhook\par
\sphinxattableend\end{savenotes}

\begin{sphinxuseclass}{sd-tab-set}
\begin{sphinxuseclass}{sd-tab-item}\subsubsection*{Objetivo}

\begin{sphinxuseclass}{sd-tab-content}
\sphinxAtStartPar
\DUrole{sd-sphinx-override}{\DUrole{sd-badge}{\DUrole{sd-bg-light}{\DUrole{sd-bg-text-light}{Exploración}}}}
\DUrole{sd-sphinx-override}{\DUrole{sd-badge}{\DUrole{sd-bg-info}{\DUrole{sd-bg-text-info}{Estructuración}}}}
\DUrole{sd-sphinx-override}{\DUrole{sd-badge}{\DUrole{sd-bg-warning}{\DUrole{sd-bg-text-warning}{Actividad práctica}}}}
\DUrole{sd-sphinx-override}{\DUrole{sd-badge}{\DUrole{sd-bg-danger}{\DUrole{sd-bg-text-danger}{Transferencia}}}}
\begin{itemize}
\item {} 
\sphinxAtStartPar
\sphinxstylestrong{Temática}: Concepto de ofimática.

\item {} 
\sphinxAtStartPar
\sphinxstylestrong{Objetivo de aprendizaje}: Entender que es la ofimática y conocer los principales proveedores de servicios ofimáticos, productividad y colaboración en la nube.

\end{itemize}

\end{sphinxuseclass}
\end{sphinxuseclass}
\begin{sphinxuseclass}{sd-tab-item}\subsubsection*{Orientaciones curriculares}

\begin{sphinxuseclass}{sd-tab-content}\begin{itemize}
\item {} 
\sphinxAtStartPar
\sphinxstylestrong{Componente}: Uso y apropiación de la Tecnología y la Informática.

\item {} 
\sphinxAtStartPar
\sphinxstylestrong{Competencia}: Genero propuestas innovadoras para el uso y aprovechamiento de productos tecnológicos.

\item {} 
\sphinxAtStartPar
\sphinxstylestrong{Evidencia de aprendizaje}: Utilizo adecuadamente herramientas informáticas para la búsqueda, organización, procesamiento, sistematización, comunicación y difusión de ideas.

\end{itemize}

\end{sphinxuseclass}
\end{sphinxuseclass}
\begin{sphinxuseclass}{sd-tab-item}\subsubsection*{Criterios de evaluación}

\begin{sphinxuseclass}{sd-tab-content}\begin{itemize}
\item {} 
\sphinxAtStartPar
Actividades desarrolladas en su totalidad.

\item {} 
\sphinxAtStartPar
Participación en clase.

\item {} 
\sphinxAtStartPar
Explicación clara de conceptos.

\item {} 
\sphinxAtStartPar
Argumentación de ideas.

\item {} 
\sphinxAtStartPar
Cumplimiento de las reglas del aula.

\end{itemize}

\end{sphinxuseclass}
\end{sphinxuseclass}
\begin{sphinxuseclass}{sd-tab-item}\subsubsection*{Recursos (0)}

\begin{sphinxuseclass}{sd-tab-content}\begin{itemize}
\item {} 
\sphinxAtStartPar
La presente guía no contiene recursos.

\end{itemize}

\end{sphinxuseclass}
\end{sphinxuseclass}
\end{sphinxuseclass}

\subsection{Mensaje del autor}
\label{\detokenize{1_guias/10ofimatica/G1011:mensaje-del-autor}}\begin{quote}

\sphinxAtStartPar
“En esta guía, el estudiante conoce las herramientas informáticas para organizar, procesar y difundir información de manera efectiva. Al entender la utilidad de suites como Google Workspace o Microsoft 365, se deduce cómo el trabajo colaborativo en la nube es esencial para las tareas de oficina o empresariales del siglo XXI.”

\begin{flushright}
---Camilo Fonseca
\end{flushright}
\end{quote}


\bigskip\hrule\bigskip



\subsection{A00: Introducción (+0,1)}
\label{\detokenize{1_guias/10ofimatica/G1011:a00-introduccion-0-1}}
\sphinxAtStartPar
\DUrole{sd-sphinx-override}{\DUrole{sd-badge}{\DUrole{sd-bg-light}{\DUrole{sd-bg-text-light}{Exploración}}}}
\begin{enumerate}
\sphinxsetlistlabels{\arabic}{enumi}{enumii}{}{.}%
\item {} 
\sphinxAtStartPar
Transcriba en su cuaderno el objetivo, las orientaciones curriculares y los criterios de evaluación de la presente guía.

\item {} 
\sphinxAtStartPar
Lea el \sphinxstylestrong{mensaje del autor} y en su cuaderno responda:
\begin{itemize}
\item {} 
\sphinxAtStartPar
Imagine que tiene que coordinar un proyecto con un equipo de personas ubicadas en diferentes lugares ¿Qué herramientas, dispositivos o software podría usar para desarrollar este proyecto?

\end{itemize}

\end{enumerate}


\bigskip\hrule\bigskip



\subsection{A01: OFIMÁTICA (+0,1)}
\label{\detokenize{1_guias/10ofimatica/G1011:a01-ofimatica-0-1}}
\sphinxAtStartPar
\DUrole{sd-sphinx-override}{\DUrole{sd-badge}{\DUrole{sd-bg-info}{\DUrole{sd-bg-text-info}{Estructuración}}}}
\begin{enumerate}
\sphinxsetlistlabels{\arabic}{enumi}{enumii}{}{.}%
\item {} 
\sphinxAtStartPar
Lea los siguientes conceptos:

\end{enumerate}


\bigskip\hrule\bigskip



\subsubsection{¿Qué es la ofimática?}
\label{\detokenize{1_guias/10ofimatica/G1011:que-es-la-ofimatica}}
\begin{sphinxadmonition}{note}{Concepto de ofimática}

\sphinxAtStartPar
%
\begin{footnote}[1]\sphinxAtStartFootnote
Superior, E. F. P. (2023, August 31). ¿Qué es la ofimática y para qué sirve?: herramientas. Esic.edu; ESIC. https://www.esic.edu/business/que\sphinxhyphen{}es\sphinxhyphen{}la\sphinxhyphen{}ofimatica\sphinxhyphen{}para\sphinxhyphen{}que\sphinxhyphen{}sirve\sphinxhyphen{}c
%
\end{footnote} La \sphinxstylestrong{ofimática} es el conjunto de técnicas, aplicaciones y herramientas informáticas que se utilizan para optimizar, automatizar y mejorar los procedimientos o tareas relacionadas con el trabajo de oficina. El término “ofimática” es una combinación de las palabras “oficina” e “informática”.

\sphinxAtStartPar
Las principales áreas que abarca la ofimática incluyen:
\begin{itemize}
\item {} 
\sphinxAtStartPar
Procesamiento de textos (creación y edición de documentos).

\item {} 
\sphinxAtStartPar
Hojas de cálculo (manejo de datos numéricos y fórmulas).

\item {} 
\sphinxAtStartPar
Presentaciones (creación de diapositivas).

\item {} 
\sphinxAtStartPar
Gestión de bases de datos.

\item {} 
\sphinxAtStartPar
Correo electrónico y comunicaciones.

\item {} 
\sphinxAtStartPar
Gestión de calendarios y agendas.

\item {} 
\sphinxAtStartPar
Gestión documental.

\end{itemize}

\sphinxAtStartPar
La ofimática es fundamental en el entorno laboral moderno, ya que permite aumentar la productividad y eficiencia en las tareas administrativas y de gestión de información.
\end{sphinxadmonition}


\bigskip\hrule\bigskip



\subsubsection{Google Workspace}
\label{\detokenize{1_guias/10ofimatica/G1011:google-workspace}}

\begin{savenotes}\sphinxattablestart
\sphinxthistablewithglobalstyle
\centering
\begin{tabulary}{\linewidth}[t]{T}
\sphinxtoprule
\sphinxtableatstartofbodyhook

\begin{sphinxfigure-in-table}
\centering
\capstart
\noindent\sphinxincludegraphics[width=700\sphinxpxdimen]{{image015}.png}
\sphinxfigcaption{Suite de aplicaciones \sphinxstylestrong{Google WorkSpace}}\label{\detokenize{1_guias/10ofimatica/G1011:id9}}\end{sphinxfigure-in-table}\relax
\\
\sphinxbottomrule
\end{tabulary}
\sphinxtableafterendhook\par
\sphinxattableend\end{savenotes}

\sphinxAtStartPar
%
\begin{footnote}[2]\sphinxAtStartFootnote
Hill, S. (2024, March 24). Qué es Google One y cómo saber si lo necesitas. WIRED. https://es.wired.com/articulos/que\sphinxhyphen{}es\sphinxhyphen{}google\sphinxhyphen{}one\sphinxhyphen{}y\sphinxhyphen{}como\sphinxhyphen{}saber\sphinxhyphen{}si\sphinxhyphen{}lo\sphinxhyphen{}necesitas
%
\end{footnote} \sphinxstylestrong{Google Workspace} es una solución de productividad en línea desarrollada por \sphinxstylestrong{Google} que permite a los usuarios conectarse, crear y colaborar de manera segura. Incluye herramientas esenciales como:
\begin{itemize}
\item {} 
\sphinxAtStartPar
\sphinxstylestrong{Gmail}: Servicio de correo electrónico.

\item {} 
\sphinxAtStartPar
\sphinxstylestrong{Documentos}: Para crear y editar documentos de texto.

\item {} 
\sphinxAtStartPar
\sphinxstylestrong{Hojas de cálculo}: Para trabajar con datos y realizar cálculos.

\item {} 
\sphinxAtStartPar
\sphinxstylestrong{Presentaciones}: Para crear presentaciones visuales.

\item {} 
\sphinxAtStartPar
\sphinxstylestrong{Drive}: Almacenamiento en la nube para guardar y compartir archivos.

\item {} 
\sphinxAtStartPar
\sphinxstylestrong{Calendario}: Para gestionar eventos y citas.

\end{itemize}


\paragraph{Google ONE \sphinxhyphen{} AI PLUS (\$19,900/mes) (2026)}
\label{\detokenize{1_guias/10ofimatica/G1011:google-one-ai-plus-19-900-mes-2026}}
\sphinxAtStartPar
\sphinxfootnotemark[2] Plan de almacenamiento y la IA de GOOGLE en una sola suscripción.

\sphinxAtStartPar
Teniendo activa esta suscripción, el usuario recibe los siguientes servicios:
\begin{itemize}
\item {} 
\sphinxAtStartPar
Acceso a App de GEMINI con mejores funciones.

\item {} 
\sphinxAtStartPar
200 GB de almacenamiento en la nube.

\item {} 
\sphinxAtStartPar
Copias de seguridad automáticas de dispositivos.

\item {} 
\sphinxAtStartPar
Llamadas de mayor calidad con funciones de video mejoradas con \sphinxstylestrong{Google Meet}.

\item {} 
\sphinxAtStartPar
Herramientas de agenda avanzadas para ahorrar tiempo.

\end{itemize}


\bigskip\hrule\bigskip



\subsubsection{Microsoft Office}
\label{\detokenize{1_guias/10ofimatica/G1011:microsoft-office}}

\begin{savenotes}\sphinxattablestart
\sphinxthistablewithglobalstyle
\centering
\begin{tabulary}{\linewidth}[t]{T}
\sphinxtoprule
\sphinxtableatstartofbodyhook

\begin{sphinxfigure-in-table}
\centering
\capstart
\noindent\sphinxincludegraphics[width=450\sphinxpxdimen]{{image023}.png}
\sphinxfigcaption{Suite de aplicaciones \sphinxstylestrong{Microsoft Office}}\label{\detokenize{1_guias/10ofimatica/G1011:id10}}\end{sphinxfigure-in-table}\relax
\\
\sphinxbottomrule
\end{tabulary}
\sphinxtableafterendhook\par
\sphinxattableend\end{savenotes}

\sphinxAtStartPar
%
\begin{footnote}[3]\sphinxAtStartFootnote
(Visnuh), M. G. (2024, January 24). Con Microsoft Office 365 dominarás la productividad y con el precio de esta licencia en oferta no tendrás que empeñarte para ello. Xataka.com; Xataka. https://www.xataka.com/seleccion/microsoft\sphinxhyphen{}office\sphinxhyphen{}365\sphinxhyphen{}dominaras\sphinxhyphen{}productividad\sphinxhyphen{}precio\sphinxhyphen{}esta\sphinxhyphen{}licencia\sphinxhyphen{}oferta\sphinxhyphen{}no\sphinxhyphen{}tendras\sphinxhyphen{}que\sphinxhyphen{}empenarte\sphinxhyphen{}para\sphinxhyphen{}ello
%
\end{footnote} \sphinxstylestrong{Microsoft Office} es un conjunto de aplicaciones de software de productividad creado por \sphinxstylestrong{Microsoft}®. Es una herramienta muy popular y ampliamente utilizada en todo el mundo para una variedad de tareas relacionadas con la creación, edición y gestión de documentos, hojas de cálculo, presentaciones y más.

\sphinxAtStartPar
Algunas de las aplicaciones más conocidas que forman parte de Microsoft Office son:
\begin{itemize}
\item {} 
\sphinxAtStartPar
\sphinxstylestrong{Word:} Para crear y editar documentos de texto, como cartas, informes y trabajos académicos.

\item {} 
\sphinxAtStartPar
\sphinxstylestrong{Excel:} Para crear hojas de cálculo, analizar datos, realizar cálculos y crear gráficos.

\item {} 
\sphinxAtStartPar
\sphinxstylestrong{PowerPoint:} Para crear presentaciones visuales con diapositivas, texto, imágenes y animaciones.

\item {} 
\sphinxAtStartPar
\sphinxstylestrong{Outlook:} Para gestionar correos electrónicos, calendarios, contactos y tareas.

\item {} 
\sphinxAtStartPar
\sphinxstylestrong{Access:} Para crear y gestionar bases de datos (menos común en el uso diario).

\end{itemize}
\begin{quote}

\sphinxAtStartPar
\sphinxstylestrong{Microsoft Office} está disponible en diferentes versiones y planes de suscripción, incluyendo versiones de escritorio para \sphinxstylestrong{Windows} y \sphinxstylestrong{macOS}, así como versiones en línea accesibles a través de un navegador web. También existen aplicaciones móviles para dispositivos iOS y Android.
\end{quote}


\paragraph{MICROSOFT 365 Family (\$459,999/año) (2026)}
\label{\detokenize{1_guias/10ofimatica/G1011:microsoft-365-family-459-999-ano-2026}}
\sphinxAtStartPar
\sphinxfootnotemark[3] Con una suscripción activa, el usuario recibe los siguientes servicios:
\begin{itemize}
\item {} 
\sphinxAtStartPar
Compartir hasta con 5 personas.

\item {} 
\sphinxAtStartPar
Iniciar sesión hasta en cinco dispositivos diferentes simultáneamente por persona.

\item {} 
\sphinxAtStartPar
Usar el software tanto en PC, Móviles, tabletas y MAC.

\item {} 
\sphinxAtStartPar
Cada persona tiene hasta 6TB de almacenamiento en la nube.

\item {} 
\sphinxAtStartPar
Las aplicaciones tienen características exclusivas y acceso sin conexión.

\item {} 
\sphinxAtStartPar
Seguridad integrada en datos y dispositivos.

\item {} 
\sphinxAtStartPar
Acceso a correo electrónico sin anuncios.

\end{itemize}


\bigskip\hrule\bigskip



\subsection{A02: Taller (+0,8)}
\label{\detokenize{1_guias/10ofimatica/G1011:a02-taller-0-8}}
\sphinxAtStartPar
\DUrole{sd-sphinx-override}{\DUrole{sd-badge}{\DUrole{sd-bg-warning}{\DUrole{sd-bg-text-warning}{Actividad práctica}}}}
\begin{enumerate}
\sphinxsetlistlabels{\arabic}{enumi}{enumii}{}{.}%
\item {} 
\sphinxAtStartPar
En su cuaderno responda las siguientes preguntas:
\begin{itemize}
\item {} 
\sphinxAtStartPar
¿Qué es la ofimática?

\item {} 
\sphinxAtStartPar
¿Cuáles son las principales \sphinxstylestrong{áreas de trabajo} que cubre la ofimática?

\item {} 
\sphinxAtStartPar
¿Por qué la ofimática es clave para el entorno laboral moderno?

\end{itemize}

\item {} 
\sphinxAtStartPar
Escriba en su cuaderno los siguientes casos y al final de cada uno, indique a que \sphinxstylestrong{área de trabajo} de la ofimática corresponde y explique por qué.
\begin{quote}

\sphinxAtStartPar
\sphinxstylestrong{CASO 01}: “Un grupo de estudiantes crea una carpeta compartida en la nube para subir evidencias fotográficas de una salida de campo, permitiendo que todos editen archivos de forma colaborativa y segura desde cualquier dispositivo.”
\end{quote}
\begin{quote}

\sphinxAtStartPar
\sphinxstylestrong{CASO 02}: “Un profesional envía una propuesta comercial a un cliente potencial, adjuntando documentos técnicos y utilizando firmas digitales personalizadas para mantener una comunicación formal y segura.”
\end{quote}
\begin{quote}

\sphinxAtStartPar
\sphinxstylestrong{CASO 03}: “Un líder de proyecto programa una reunión de equipo sincronizada, enviando invitaciones automáticas con recordatorios y verificando la disponibilidad de horario de todos los colaboradores en tiempo real.”
\end{quote}
\begin{quote}

\sphinxAtStartPar
\sphinxstylestrong{CASO 04}: “El encargado de la biblioteca escolar crea un sistema estructurado para registrar el inventario de libros, los datos de los estudiantes y las fechas de devolución, facilitando la búsqueda rápida de ejemplares disponibles.”
\end{quote}
\begin{quote}

\sphinxAtStartPar
\sphinxstylestrong{CASO 05}: “Un estudiante de grado undécimo redacta su proyecto de grado utilizando herramientas de edición para estructurar el índice, insertar citas bibliográficas y exportar el informe final en formato PDF para su entrega formal.”
\end{quote}
\begin{quote}

\sphinxAtStartPar
\sphinxstylestrong{CASO 06}: “Un equipo de trabajo diseña una exposición visual sobre la cuarta revolución industrial, integrando imágenes, animaciones y transiciones para comunicar sus hallazgos de forma atractiva ante un comité evaluador.”
\end{quote}
\begin{quote}

\sphinxAtStartPar
\sphinxstylestrong{CASO 07}: “El tesorero de una pequeña empresa organiza los gastos mensuales mediante el uso de fórmulas automáticas para calcular el IVA, generar gráficos de barras que muestren las tendencias de consumo y filtrar datos por categorías de proveedores.”
\end{quote}

\item {} 
\sphinxAtStartPar
En su cuaderno responda las siguientes preguntas:
\begin{itemize}
\item {} 
\sphinxAtStartPar
¿Qué es Google Workspace?

\item {} 
\sphinxAtStartPar
¿Qué herramientas incluye Google Workspace?

\item {} 
\sphinxAtStartPar
¿En qué consiste la suscripción de Google ONE?

\item {} 
\sphinxAtStartPar
¿Qué servicios recibe el usuario por suscribirse a \sphinxstylestrong{Google ONE \sphinxhyphen{} IA PLUS} y cuanto cuesta?

\end{itemize}

\item {} 
\sphinxAtStartPar
En Internet, consulte que otros planes de \sphinxstylestrong{Google ONE} existen y por cada plan en su cuaderno responda:
\begin{itemize}
\item {} 
\sphinxAtStartPar
¿Cuanto cuesta el plan?

\item {} 
\sphinxAtStartPar
¿Que servicios carece dicho plan respecto a \sphinxstylestrong{Google ONE \sphinxhyphen{} IA PLUS}?

\item {} 
\sphinxAtStartPar
¿Que servicios extra ofrece dicho plan respecto a \sphinxstylestrong{Google ONE \sphinxhyphen{} IA PLUS}?

\end{itemize}

\item {} 
\sphinxAtStartPar
En su cuaderno responda las siguientes preguntas:
\begin{itemize}
\item {} 
\sphinxAtStartPar
¿Qué es Microsoft Office?

\item {} 
\sphinxAtStartPar
¿Cúales son las aplicaciones más conocidas de Microsoft Office?

\item {} 
\sphinxAtStartPar
¿Qué sistemas operativos son compatibles con Microsoft Office?

\item {} 
\sphinxAtStartPar
¿En qué consiste la suscripción \sphinxstylestrong{Microsoft 365 Family}?

\item {} 
\sphinxAtStartPar
¿Qué servicios recibe el usuario por suscribirse a \sphinxstylestrong{Microsoft 365 Family} y cuanto cuesta?

\end{itemize}

\item {} 
\sphinxAtStartPar
En Internet, consulte que otros planes de Microsoft 365 existen y por cada plan en su cuaderno responda:
\begin{itemize}
\item {} 
\sphinxAtStartPar
¿Cuanto cuesta dicho plan?

\item {} 
\sphinxAtStartPar
¿Que servicios carece dicho plan respecto a \sphinxstylestrong{Microsoft Family 365}?

\item {} 
\sphinxAtStartPar
¿Que servicios extra ofrece dicho plan respecto a \sphinxstylestrong{Microsoft Family 365}?

\end{itemize}

\item {} 
\sphinxAtStartPar
Para cada uno de los casos, seleccione el servicio que más convenga. No olvide explicar el por qué. Los servicios son:
\begin{itemize}
\item {} 
\sphinxAtStartPar
Microsoft 365 Personal.

\item {} 
\sphinxAtStartPar
Microsoft 365 Familia.

\item {} 
\sphinxAtStartPar
Microsoft 365 Premium.

\item {} 
\sphinxAtStartPar
Office Hogar 2024.

\item {} 
\sphinxAtStartPar
Plan gratuito de Google One.

\item {} 
\sphinxAtStartPar
Plan gratuito de Google.

\item {} 
\sphinxAtStartPar
Planes de suscripción de Google One.

\item {} 
\sphinxAtStartPar
Google One con Google AI.

\item {} 
\sphinxAtStartPar
Uso de la aplicación Google One.

\item {} 
\sphinxAtStartPar
Planes familiares de Google One.

\end{itemize}

\sphinxAtStartPar
Los casos son:
\begin{quote}

\sphinxAtStartPar
\sphinxstylestrong{CASO 01:} Sofía es una diseñadora digital que busca optimizar su tiempo usando herramientas avanzadas de Inteligencia Artificial. Necesita acceso exclusivo a las funciones más recientes de Copilot y un uso extensivo para crear y editar imágenes con IA.
\end{quote}
\begin{quote}

\sphinxAtStartPar
\sphinxstylestrong{CASO 02:} Un usuario prefiere gestionar todas sus suscripciones y copias de seguridad directamente desde su smartphone. Necesita una forma sencilla de monitorear cuánto espacio ocupan sus fotos y archivos en tiempo real. El servicio debe suministrar una herramienta de administración móvil centralizada.
\end{quote}
\begin{quote}

\sphinxAtStartPar
\sphinxstylestrong{CASO 03:} Elena es una estudiante interesada en las últimas tendencias tecnológicas. Quiere integrar funciones avanzadas de IA directamente en sus herramientas de trabajo de Google. El servicio debe permitir acceder a los beneficios exclusivos de Google AI.
\end{quote}
\begin{quote}

\sphinxAtStartPar
\sphinxstylestrong{CASO 04:} Carlos estudia ingeniería y necesita usar Word, Excel y PowerPoint tanto en su PC como en su celular. Vive solo y tiene un presupuesto ajustado, pero requiere al menos 1 TB de espacio para sus proyectos pesados.
\end{quote}
\begin{quote}

\sphinxAtStartPar
\sphinxstylestrong{CASO 05:} Lucía utiliza su cuenta de Google principalmente para enviar correos electrónicos académicos y guardar copias de seguridad de las fotos de su celular. No suele manejar archivos muy pesados. Se necesita un almacenamiento básico sin coste adicional.
\end{quote}
\begin{quote}

\sphinxAtStartPar
\sphinxstylestrong{CASO 06:} Jorge utiliza su correo electrónico para comunicaciones básicas y ocasionalmente guarda documentos PDF y fotos personales. No utiliza aplicaciones de escritorio complejas y sus archivos totales no superan los 12 GB. Un almacenamiento gratuito sin coste adicional.
\end{quote}
\begin{quote}

\sphinxAtStartPar
\sphinxstylestrong{CASO 07:} Un pequeño grupo de estudio quiere una solución donde puedan compartir un fondo común de almacenamiento para sus recursos compartidos, pero manteniendo cada uno su privacidad. El servicio debe permitir compartir beneficios con otros miembros.
\end{quote}
\begin{quote}

\sphinxAtStartPar
\sphinxstylestrong{CASO 08:} Martín ha alcanzado el límite de sus 15 GB gratuitos debido a la gran cantidad de videos y trabajos de investigación guardados en su unidad. Necesita un poco más de espacio para no dejar de recibir correos importantes. Se debe adquirir una ampliación de almacenamiento económica para seguir usando su cuenta personal
\end{quote}
\begin{quote}

\sphinxAtStartPar
\sphinxstylestrong{CASO 09:} Una pequeña oficina de contabilidad prefiere no tener gastos recurrentes mensuales o anuales. Solo necesitan instalar las versiones clásicas de Word y Excel en una computadora de la oficina y no requieren almacenamiento en la nube ni actualizaciones constantes de funciones. No se requiere una suscripción.
\end{quote}
\begin{quote}

\sphinxAtStartPar
\sphinxstylestrong{CASO 10:} La familia Martínez tiene 5 integrantes (padres y 3 hijos en el colegio). Todos necesitan las aplicaciones de Office para sus tareas y quieren que cada uno tenga su propio espacio privado de almacenamiento de 1 TB sin pagar suscripciones por separado.
\end{quote}

\end{enumerate}


\bigskip\hrule\bigskip



\subsection{A03: Examen escrito (+4,0)}
\label{\detokenize{1_guias/10ofimatica/G1011:a03-examen-escrito-4-0}}
\sphinxAtStartPar
\DUrole{sd-sphinx-override}{\DUrole{sd-badge}{\DUrole{sd-bg-danger}{\DUrole{sd-bg-text-danger}{Transferencia}}}}


\begin{savenotes}\sphinxattablestart
\sphinxthistablewithglobalstyle
\centering
\begin{tabulary}{\linewidth}[t]{TT}
\sphinxtoprule
\sphinxtableatstartofbodyhook
\sphinxAtStartPar
\sphinxstylestrong{TIEMPO}
&
\sphinxAtStartPar
29 minutos
\\
\sphinxhline
\sphinxAtStartPar
\sphinxstylestrong{PREGUNTAS}
&
\sphinxAtStartPar
10
\\
\sphinxhline
\sphinxAtStartPar
\sphinxstylestrong{OBJETIVO}
&
\sphinxAtStartPar
Comprobar conocimientos adquiridos.
\\
\sphinxhline
\sphinxAtStartPar
\sphinxstylestrong{ACTIVIDAD}
&
\sphinxAtStartPar
Responda las siguientes preguntas. Cuadernos y celulares guardados. Solo se permite hoja, lapiz y borrador.
\\
\sphinxbottomrule
\end{tabulary}
\sphinxtableafterendhook\par
\sphinxattableend\end{savenotes}

\begin{sphinxadmonition}{note}{Posibles preguntas}
\begin{enumerate}
\sphinxsetlistlabels{\arabic}{enumi}{enumii}{}{.}%
\item {} 
\sphinxAtStartPar
Explique con sus propias palabras qué es la \sphinxstylestrong{ofimática} y de qué combinación de términos surge su nombre.

\item {} 
\sphinxAtStartPar
¿Cuál es el propósito fundamental de utilizar \sphinxstylestrong{herramientas ofimáticas} en un entorno de trabajo o estudio?

\item {} 
\sphinxAtStartPar
Mencione al menos cuatro áreas principales que abarca la ofimática.

\item {} 
\sphinxAtStartPar
Describa la función principal de las herramientas de \sphinxstylestrong{procesamiento de textos} e indique un ejemplo de software tanto de Google como de Microsoft.

\item {} 
\sphinxAtStartPar
¿Para qué tipo de tareas es ideal el uso de una \sphinxstylestrong{hoja de cálculo} y cuáles son sus componentes básicos?

\item {} 
\sphinxAtStartPar
Explique la utilidad de aplicaciones como \sphinxstylestrong{PowerPoint} o \sphinxstylestrong{Google Slides} en el ámbito académico o empresarial.

\item {} 
\sphinxAtStartPar
¿En qué consiste la gestión de \sphinxstylestrong{bases de datos} y qué aplicación de \sphinxstylestrong{Microsoft} se menciona para esta tarea?

\item {} 
\sphinxAtStartPar
¿En qué consiste \sphinxstylestrong{Google Workspace} y qué beneficio principal tienen los usuarios cuando necesitan trabajar en equipo?

\item {} 
\sphinxAtStartPar
Describa la función de \sphinxstylestrong{Google Drive} dentro del ecosistema de Google Workspace.

\item {} 
\sphinxAtStartPar
Según el plan \sphinxstylestrong{Google One AI PLUS}, ¿Qué servicios adicionales recibe un usuario al activar esta suscripción específica?

\item {} 
\sphinxAtStartPar
¿Qué es \sphinxstylestrong{Gemini} y bajo qué plan de suscripción de Google se obtienen sus mejores funciones?

\item {} 
\sphinxAtStartPar
¿Qué es \sphinxstylestrong{Microsoft Office} y para qué tipos de dispositivos está disponible?

\item {} 
\sphinxAtStartPar
¿Cuales son las diferencias entre las funciones principales de un \sphinxstylestrong{procesador de texto} y una \sphinxstylestrong{hoja de cálculo}?

\item {} 
\sphinxAtStartPar
¿Qué funciones integra \sphinxstylestrong{Outlook} además de la gestión de \sphinxstylestrong{correos electrónicos}?

\item {} 
\sphinxAtStartPar
Explique tres beneficios clave de contar con la \sphinxstylestrong{suscripción familiar de Microsoft 365}.

\item {} 
\sphinxAtStartPar
Compare la capacidad de almacenamiento y los precios que ofrece el plan \sphinxstylestrong{Google One AI PLUS} frente a lo que recibe cada persona en el plan \sphinxstylestrong{Microsoft 365 Family}.

\item {} 
\sphinxAtStartPar
¿Por qué se considera que las suites en línea permiten “conectar, crear y colaborar de manera segura”?

\item {} 
\sphinxAtStartPar
Según las características de \sphinxstylestrong{Microsoft 365}, ¿qué ventajas ofrece en cuanto a seguridad y acceso a funciones en comparación con versiones estáticas?

\item {} 
\sphinxAtStartPar
¿Por qué la \sphinxstylestrong{ofimática} es considerada fundamental en el \sphinxstylestrong{entorno laboral moderno}?

\item {} 
\sphinxAtStartPar
¿Qué herramientas específicas se mencionan para la gestión de calendarios, agendas y documentos en ambas suites (Google y Microsoft)?

\end{enumerate}
\end{sphinxadmonition}


\bigskip\hrule\bigskip


\sphinxstepscope


\section{Guía 1012: Software Libre (P1)}
\label{\detokenize{1_guias/10ofimatica/G1012:guia-1012-software-libre-p1}}\label{\detokenize{1_guias/10ofimatica/G1012::doc}}

\begin{savenotes}\sphinxattablestart
\sphinxthistablewithglobalstyle
\centering
\begin{tabular}[t]{\X{35}{100}\X{65}{100}}
\sphinxtoprule
\sphinxtableatstartofbodyhook

\begin{savenotes}\sphinxattablestart
\sphinxthistablewithglobalstyle
\centering
\begin{tabulary}{\linewidth}[t]{T}
\sphinxtoprule
\sphinxtableatstartofbodyhook
\sphinxAtStartPar
\sphinxincludegraphics{{escudo}.png}
\\
\sphinxbottomrule
\end{tabulary}
\sphinxtableafterendhook\par
\sphinxattableend\end{savenotes}
&

\begin{savenotes}\sphinxattablestart
\sphinxthistablewithglobalstyle
\raggedright
\begin{tabular}[t]{*{1}{\X{1}{1}}}
\sphinxtoprule
\sphinxtableatstartofbodyhook
\sphinxAtStartPar
\sphinxstylestrong{INSTITUCIÓN EDUCATIVA}: John F. Kennedy
\begin{itemize}
\item {} 
\sphinxAtStartPar
\sphinxstylestrong{DOCENTE}: Julian Camilo Fonseca Romero

\item {} 
\sphinxAtStartPar
\sphinxstylestrong{ASIGNATURA}: Ofimática

\item {} 
\sphinxAtStartPar
\sphinxstylestrong{GRADO}: DÉCIMO

\item {} 
\sphinxAtStartPar
\sphinxstylestrong{PERIODO:} 01

\item {} 
\sphinxAtStartPar
\sphinxstylestrong{TIEMPO}: 10 semanas (2 horas / semana)

\item {} 
\sphinxAtStartPar
\sphinxstylestrong{TIEMPO DE TRABAJO}: 12/20 horas

\end{itemize}
\\
\sphinxbottomrule
\end{tabular}
\sphinxtableafterendhook\par
\sphinxattableend\end{savenotes}
\\
\sphinxbottomrule
\end{tabular}
\sphinxtableafterendhook\par
\sphinxattableend\end{savenotes}

\begin{sphinxuseclass}{sd-tab-set}
\begin{sphinxuseclass}{sd-tab-item}\subsubsection*{Objetivo}

\begin{sphinxuseclass}{sd-tab-content}
\sphinxAtStartPar
\DUrole{sd-sphinx-override}{\DUrole{sd-badge}{\DUrole{sd-bg-light}{\DUrole{sd-bg-text-light}{Exploración}}}}
\DUrole{sd-sphinx-override}{\DUrole{sd-badge}{\DUrole{sd-bg-info}{\DUrole{sd-bg-text-info}{Estructuración}}}}
\DUrole{sd-sphinx-override}{\DUrole{sd-badge}{\DUrole{sd-bg-warning}{\DUrole{sd-bg-text-warning}{Actividad práctica}}}}
\DUrole{sd-sphinx-override}{\DUrole{sd-badge}{\DUrole{sd-bg-danger}{\DUrole{sd-bg-text-danger}{Transferencia}}}}
\begin{itemize}
\item {} 
\sphinxAtStartPar
\sphinxstylestrong{Temática}: Software libre.

\item {} 
\sphinxAtStartPar
\sphinxstylestrong{Objetivo de aprendizaje}: Entender la importancia del desarrollo y uso del software libre.

\end{itemize}

\end{sphinxuseclass}
\end{sphinxuseclass}
\begin{sphinxuseclass}{sd-tab-item}\subsubsection*{Orientaciones curriculares}

\begin{sphinxuseclass}{sd-tab-content}\begin{itemize}
\item {} 
\sphinxAtStartPar
\sphinxstylestrong{Componente}: Tecnología, Informática y Sociedad.

\item {} 
\sphinxAtStartPar
\sphinxstylestrong{Competencia}: Actúo críticamente y de forma argumentada frente a las implicaciones éticas, sociales y ambientales del desarrollo, implementación, uso y disposición final de los productos tecnológicos.

\item {} 
\sphinxAtStartPar
\sphinxstylestrong{Evidencia de aprendizaje}: Juzgo las implicaciones que la protección a la propiedad intelectual tiene sobre el desarrollo y uso de diversas manifestaciones tecnológicas en el mundo.

\end{itemize}

\end{sphinxuseclass}
\end{sphinxuseclass}
\begin{sphinxuseclass}{sd-tab-item}\subsubsection*{Criterios de evaluación}

\begin{sphinxuseclass}{sd-tab-content}\begin{itemize}
\item {} 
\sphinxAtStartPar
Actividades desarrolladas en su totalidad.

\item {} 
\sphinxAtStartPar
Participación en clase.

\item {} 
\sphinxAtStartPar
Explicación clara de conceptos.

\item {} 
\sphinxAtStartPar
Argumentación de ideas.

\item {} 
\sphinxAtStartPar
Cumplimiento de las reglas del aula.

\end{itemize}

\end{sphinxuseclass}
\end{sphinxuseclass}
\begin{sphinxuseclass}{sd-tab-item}\subsubsection*{Recursos (0)}

\begin{sphinxuseclass}{sd-tab-content}\begin{itemize}
\item {} 
\sphinxAtStartPar
La presente guía no contiene recursos.

\end{itemize}

\end{sphinxuseclass}
\end{sphinxuseclass}
\end{sphinxuseclass}

\subsection{Mensaje del autor}
\label{\detokenize{1_guias/10ofimatica/G1012:mensaje-del-autor}}\begin{quote}

\sphinxAtStartPar
“El estudio del software libre es esencial ya que, como manifestación tecnológica, surge como una alternativa ética y política a los modelos de propiedad intelectual. Esta filosofía no solo garantiza la equidad en el acceso a recursos digitales, sino que permite trascender del saber tecnológico al conocimiento tecnológico profundo. Al permitir el análisis de su funcionamiento interno (posibilidad negada por el software propietario), se faculta al estudiante para adoptar, adaptar y generar soluciones que respondan a necesidades sociales.”

\begin{flushright}
---Camilo Fonseca
\end{flushright}
\end{quote}


\bigskip\hrule\bigskip



\subsection{A00: Introducción (+0,1)}
\label{\detokenize{1_guias/10ofimatica/G1012:a00-introduccion-0-1}}
\sphinxAtStartPar
\DUrole{sd-sphinx-override}{\DUrole{sd-badge}{\DUrole{sd-bg-light}{\DUrole{sd-bg-text-light}{Exploración}}}}
\begin{enumerate}
\sphinxsetlistlabels{\arabic}{enumi}{enumii}{}{.}%
\item {} 
\sphinxAtStartPar
Transcriba en su cuaderno el objetivo, las orientaciones curriculares y los criterios de evaluación de la presente guía.

\item {} 
\sphinxAtStartPar
Lea el mensaje del autor luego identifique y describa en su cuaderno \sphinxstylestrong{5 ejemplos de software} gratuitos que utilice, reflexionando sobre cómo estos facilitan sus actividades cotidianas.

\end{enumerate}


\bigskip\hrule\bigskip



\subsection{A01: Software libre (+0,1)}
\label{\detokenize{1_guias/10ofimatica/G1012:a01-software-libre-0-1}}
\sphinxAtStartPar
\DUrole{sd-sphinx-override}{\DUrole{sd-badge}{\DUrole{sd-bg-info}{\DUrole{sd-bg-text-info}{Estructuración}}}}
\begin{enumerate}
\sphinxsetlistlabels{\arabic}{enumi}{enumii}{}{.}%
\item {} 
\sphinxAtStartPar
Lea los siguientes conceptos relacionados con software libre.

\end{enumerate}


\bigskip\hrule\bigskip



\subsubsection{¿Qué es el software?}
\label{\detokenize{1_guias/10ofimatica/G1012:que-es-el-software}}
\sphinxAtStartPar
%
\begin{footnote}[1]\sphinxAtStartFootnote
Deitel, P. J., \& Deitel, H. (2014). C++: cómo programar (A. V. Romero Elizondo, Trans.). Pearson Educación.
%
\end{footnote} “El software es un conjunto de programas, instrucciones y reglas informáticas que permiten ejecutar distintas tareas en una computadora. Se trata de las aplicaciones y sistemas que hacen posible que los usuarios puedan trabajar con un \sphinxstylestrong{ordenador} o \sphinxstylestrong{dispositivo electrónico}”.

\begin{sphinxadmonition}{note}{Tipos de software}

\sphinxAtStartPar
\sphinxfootnotemark[1] El software tiene los siguientes tipos:
\begin{itemize}
\item {} 
\sphinxAtStartPar
\sphinxstylestrong{Software de sistema:} Programas básicos que permiten que el hardware funcione (sistemas operativos, controladores).

\item {} 
\sphinxAtStartPar
\sphinxstylestrong{Software de aplicación:} Programas que realizan tareas específicas para los usuarios (procesadores de texto, navegadores web).

\end{itemize}
\end{sphinxadmonition}

\begin{sphinxadmonition}{note}{Clasificación de software}

\sphinxAtStartPar
%
\begin{footnote}[2]\sphinxAtStartFootnote
Tipos de licencias de software: software libre, propietario y demás. (2022, July 19). Velneo.com. https://velneo.com/blog/licencias\sphinxhyphen{}software\sphinxhyphen{}libre\sphinxhyphen{}propietario\sphinxhyphen{}otros/
%
\end{footnote} El software se puede clasificar de la siguiente manera:
\begin{itemize}
\item {} 
\sphinxAtStartPar
\sphinxstylestrong{Software libre:} Software que puede ser usado con \sphinxstyleemphasis{libertad} y de \sphinxstyleemphasis{forma legal}.

\item {} 
\sphinxAtStartPar
\sphinxstylestrong{Software propietario:} Software que puede ser usado con \sphinxstyleemphasis{restricciones} y de \sphinxstyleemphasis{forma legal}.

\item {} 
\sphinxAtStartPar
\sphinxstylestrong{Software pirata:} Software que es usado de \sphinxstylestrong{forma ilegal} (independiente si es usado con libertad o restricciones).

\end{itemize}
\end{sphinxadmonition}


\bigskip\hrule\bigskip



\subsubsection{Software libre}
\label{\detokenize{1_guias/10ofimatica/G1012:software-libre}}
\sphinxAtStartPar
\sphinxfootnotemark[2] El \sphinxstylestrong{software libre} es un tipo de software que respeta la libertad de los usuarios y la comunidad. Los usuarios tienen la libertad de ejecutar, copiar, distribuir, estudiar, modificar y mejorar el software.

\begin{sphinxadmonition}{note}{COLABORACIÓN}

\sphinxAtStartPar
El software libre promueve la colaboración, la transparencia y el conocimiento compartido, permitiendo que las comunidades de desarrolladores y usuarios trabajen juntos para mejorar el software continuamente.
\end{sphinxadmonition}

\sphinxAtStartPar
El software libre se caracteriza por:
\begin{itemize}
\item {} 
\sphinxAtStartPar
Se centra en la \sphinxstylestrong{libertad} como un principio ético y filosófico.

\item {} 
\sphinxAtStartPar
Enfatiza las cuatro libertades fundamentales mencionadas anteriormente.

\item {} 
\sphinxAtStartPar
Impulsado por la Free Software Foundation (FSF).

\item {} 
\sphinxAtStartPar
Usa licencias que garantizan que las versiones modificadas también sean libres (Copyleft).

\end{itemize}

\begin{sphinxadmonition}{note}{4 Libertades del software libre}

\sphinxAtStartPar
\sphinxfootnotemark[2] Las cuatro libertades esenciales del software libre son:
\begin{itemize}
\item {} 
\sphinxAtStartPar
\sphinxstylestrong{(Libertad 0):} La libertad de ejecutar el programa como se desee, con cualquier propósito.

\item {} 
\sphinxAtStartPar
\sphinxstylestrong{(Libertad 1):} La libertad de estudiar cómo funciona el programa y cambiarlo para que haga lo que uno quiera. (El acceso al código fuente es una condición necesaria para ello).

\item {} 
\sphinxAtStartPar
\sphinxstylestrong{(Libertad 2):} La libertad de redistribuir copias para ayudar a otros.

\item {} 
\sphinxAtStartPar
\sphinxstylestrong{(Libertad 3):} La libertad de distribuir copias de sus versiones modificadas a terceros.

\end{itemize}

\sphinxAtStartPar
Estas \sphinxstylestrong{libertades} son esenciales, no opcionales. Un \sphinxstylestrong{software} es \sphinxstylestrong{libre} si otorga a los usuarios todas estas libertades de manera adecuada.
\end{sphinxadmonition}


\bigskip\hrule\bigskip



\subsubsection{Ejemplos de software libre}
\label{\detokenize{1_guias/10ofimatica/G1012:ejemplos-de-software-libre}}
\sphinxAtStartPar
A continuación se muestran algunos ejemplos del software libre:

\begin{sphinxadmonition}{note}{Sistema Operativo \sphinxstylestrong{GNU\sphinxhyphen{}LINUX} (\$0)}


\begin{savenotes}\sphinxattablestart
\sphinxthistablewithglobalstyle
\centering
\begin{tabulary}{\linewidth}[t]{T}
\sphinxtoprule
\sphinxtableatstartofbodyhook

\begin{sphinxfigure-in-table}
\centering
\capstart
\noindent\sphinxincludegraphics[width=300\sphinxpxdimen]{{image016}.png}
\sphinxfigcaption{Logo GNU/Linux}\label{\detokenize{1_guias/10ofimatica/G1012:id41}}\end{sphinxfigure-in-table}\relax
\\
\sphinxbottomrule
\end{tabulary}
\sphinxtableafterendhook\par
\sphinxattableend\end{savenotes}

\sphinxAtStartPar
%
\begin{footnote}[3]\sphinxAtStartFootnote
Rohaut, S. Linux: preparación a la certificación LPIC 1 (exámenes LPI 101 y LPI 102), 3ª Edición. Barcelona: ENI ediciones, 2015.
%
\end{footnote} GNU/Linux es un sistema operativo libre y de código abierto que combina el núcleo Linux con las herramientas del proyecto GNU.

\sphinxAtStartPar
GNU/Linux incluye las siguientes características:
\begin{itemize}
\item {} 
\sphinxAtStartPar
Gratuito y libre para usar, modificar y distribuir.

\item {} 
\sphinxAtStartPar
Alta seguridad y estabilidad.

\item {} 
\sphinxAtStartPar
Gran variedad de distribuciones para diferentes necesidades (Ubuntu, Fedora, Debian, etc.).

\item {} 
\sphinxAtStartPar
Funciona en diversos dispositivos, desde computadores personales hasta servidores.

\end{itemize}
\begin{quote}

\sphinxAtStartPar
El sistema fue creado inicialmente por \sphinxstylestrong{Linus Torvalds} (kernel Linux) en conjunto con el proyecto GNU de \sphinxstylestrong{Richard Stallman}, y hoy es mantenido por una gran comunidad de desarrolladores en todo el mundo.
\end{quote}
\end{sphinxadmonition}

\begin{sphinxadmonition}{note}{Navegador web \sphinxstylestrong{Mozilla FIREFOX} (\$0)}


\begin{savenotes}\sphinxattablestart
\sphinxthistablewithglobalstyle
\centering
\begin{tabulary}{\linewidth}[t]{T}
\sphinxtoprule
\sphinxtableatstartofbodyhook

\begin{sphinxfigure-in-table}
\centering
\capstart
\noindent\sphinxincludegraphics[width=360\sphinxpxdimen]{{image024}.png}
\sphinxfigcaption{Logo Navegador Mozilla Firefox}\label{\detokenize{1_guias/10ofimatica/G1012:id42}}\end{sphinxfigure-in-table}\relax
\\
\sphinxbottomrule
\end{tabulary}
\sphinxtableafterendhook\par
\sphinxattableend\end{savenotes}

\sphinxAtStartPar
%
\begin{footnote}[4]\sphinxAtStartFootnote
Andreu, A. (2024, September 17). Qué es Mozilla Firefox, diferencias con otros navegadores y extensiones imprescindibles. Computer Hoy. https://computerhoy.20minutos.es/tecnologia/mozilla\sphinxhyphen{}firefox\sphinxhyphen{}diferencias\sphinxhyphen{}otros\sphinxhyphen{}navegadores\sphinxhyphen{}extensiones\sphinxhyphen{}imprescindibles\sphinxhyphen{}1404587
%
\end{footnote} \sphinxstylestrong{Mozilla Firefox} es un navegador web de código abierto y gratuito que se caracteriza por:
\begin{itemize}
\item {} 
\sphinxAtStartPar
Alta velocidad y rendimiento en la navegación.

\item {} 
\sphinxAtStartPar
Fuerte énfasis en la privacidad y seguridad del usuario.

\item {} 
\sphinxAtStartPar
Personalizable a través de extensiones y temas.

\item {} 
\sphinxAtStartPar
Compatible con los estándares web modernos.

\end{itemize}

\sphinxAtStartPar
Firefox es desarrollado por Mozilla Foundation, una organización sin fines de lucro que promueve una internet abierta y accesible para todos. El navegador representa una alternativa libre a opciones privativas como \sphinxstylestrong{Google Chrome}.
\end{sphinxadmonition}

\begin{sphinxadmonition}{note}{Editor de imágenes \sphinxstylestrong{GIMP} (\$0)}


\begin{savenotes}\sphinxattablestart
\sphinxthistablewithglobalstyle
\centering
\begin{tabulary}{\linewidth}[t]{T}
\sphinxtoprule
\sphinxtableatstartofbodyhook

\begin{sphinxfigure-in-table}
\centering
\capstart
\noindent\sphinxincludegraphics[width=300\sphinxpxdimen]{{image033}.png}
\sphinxfigcaption{Logo \sphinxstylestrong{Gimp}}\label{\detokenize{1_guias/10ofimatica/G1012:id43}}\end{sphinxfigure-in-table}\relax
\\
\sphinxbottomrule
\end{tabulary}
\sphinxtableafterendhook\par
\sphinxattableend\end{savenotes}

\sphinxAtStartPar
%
\begin{footnote}[5]\sphinxAtStartFootnote
Delgado, J. M. (2024, September 15). Razones por las que GIMP es la mejor alternativa a Photoshop para editar tus fotos. Computer Hoy. https://computerhoy.20minutos.es/apps/razones\sphinxhyphen{}gimp\sphinxhyphen{}mejor\sphinxhyphen{}alternativa\sphinxhyphen{}photoshop\sphinxhyphen{}editar\sphinxhyphen{}fotos\sphinxhyphen{}1400576
%
\end{footnote} GIMP (\sphinxstylestrong{GNU Image Manipulation Program}) es un editor de imágenes libre y de código abierto que ofrece:
\begin{itemize}
\item {} 
\sphinxAtStartPar
Herramientas profesionales para edición y retoque de imágenes.

\item {} 
\sphinxAtStartPar
Soporte para múltiples formatos de archivo.

\item {} 
\sphinxAtStartPar
Capacidad de trabajar con capas y filtros.

\item {} 
\sphinxAtStartPar
Sistema de plugins para extender funcionalidades.

\end{itemize}

\sphinxAtStartPar
Es una alternativa gratuita a programas comerciales como \sphinxstylestrong{Photoshop}, siendo especialmente popular en la comunidad Linux aunque también está disponible para Windows y MacOS.
\end{sphinxadmonition}

\begin{sphinxadmonition}{note}{Editor de gráficos vectoriales \sphinxstylestrong{INKSCAPE} (\$0)}


\begin{savenotes}\sphinxattablestart
\sphinxthistablewithglobalstyle
\centering
\begin{tabulary}{\linewidth}[t]{T}
\sphinxtoprule
\sphinxtableatstartofbodyhook

\begin{sphinxfigure-in-table}
\centering
\capstart
\noindent\sphinxincludegraphics[width=360\sphinxpxdimen]{{image043}.png}
\sphinxfigcaption{Logo \sphinxstylestrong{Inkscape}}\label{\detokenize{1_guias/10ofimatica/G1012:id44}}\end{sphinxfigure-in-table}\relax
\\
\sphinxbottomrule
\end{tabulary}
\sphinxtableafterendhook\par
\sphinxattableend\end{savenotes}

\sphinxAtStartPar
%
\begin{footnote}[6]\sphinxAtStartFootnote
Inkscape, tu ilustrador de cabecera. (2021, July 19). INTEF. https://intef.es/observatorio\_tecno/inkscape\sphinxhyphen{}tu\sphinxhyphen{}ilustrador\sphinxhyphen{}de\sphinxhyphen{}cabecera/
%
\end{footnote} Inkscape es un editor de gráficos vectoriales libre y de código abierto que ofrece:
\begin{itemize}
\item {} 
\sphinxAtStartPar
Creación y edición de gráficos vectoriales profesionales.

\item {} 
\sphinxAtStartPar
Compatible con estándares como SVG (Scalable Vector Graphics).

\item {} 
\sphinxAtStartPar
Herramientas avanzadas de dibujo y manipulación de formas.

\item {} 
\sphinxAtStartPar
Ideal para crear logotipos, ilustraciones y diseños escalables.

\end{itemize}

\sphinxAtStartPar
Es una alternativa gratuita a programas comerciales como \sphinxstylestrong{Adobe Illustrator}, permitiendo crear gráficos que mantienen su calidad sin importar el tamaño al que se escalen.
\end{sphinxadmonition}

\begin{sphinxadmonition}{note}{Reproductor multimedia \sphinxstylestrong{VLC} (\$0)}


\begin{savenotes}\sphinxattablestart
\sphinxthistablewithglobalstyle
\centering
\begin{tabulary}{\linewidth}[t]{T}
\sphinxtoprule
\sphinxtableatstartofbodyhook

\begin{sphinxfigure-in-table}
\centering
\capstart
\noindent\sphinxincludegraphics[width=400\sphinxpxdimen]{{image053}.png}
\sphinxfigcaption{Logo \sphinxstylestrong{VLC}}\label{\detokenize{1_guias/10ofimatica/G1012:id45}}\end{sphinxfigure-in-table}\relax
\\
\sphinxbottomrule
\end{tabulary}
\sphinxtableafterendhook\par
\sphinxattableend\end{savenotes}

\sphinxAtStartPar
%
\begin{footnote}[7]\sphinxAtStartFootnote
Delgado, J. M. (2025, June 14). VLC sigue muy vivo: cosas sorprendentes que puedes hacer con el mejor reproductor de vídeo para PC. Computer Hoy. https://computerhoy.20minutos.es/tecnologia/vlc\sphinxhyphen{}sigue\sphinxhyphen{}muy\sphinxhyphen{}vivo\sphinxhyphen{}cosas\sphinxhyphen{}sorprendentes\sphinxhyphen{}puedes\sphinxhyphen{}hacer\sphinxhyphen{}mejor\sphinxhyphen{}reproductor\sphinxhyphen{}video\sphinxhyphen{}pc\sphinxhyphen{}1463759
%
\end{footnote} VLC Media Player es un reproductor multimedia libre y de código abierto que se caracteriza por:
\begin{itemize}
\item {} 
\sphinxAtStartPar
Capacidad para reproducir casi cualquier formato de audio y video.

\item {} 
\sphinxAtStartPar
No requiere códecs adicionales para funcionar.

\item {} 
\sphinxAtStartPar
Permite streaming y conversión de archivos multimedia.

\item {} 
\sphinxAtStartPar
Funciona en múltiples sistemas operativos (Windows, Linux, Mac).

\end{itemize}

\sphinxAtStartPar
Desarrollado por VideoLAN Organization, VLC es uno de los reproductores más populares debido a su versatilidad y facilidad de uso. Es completamente gratuito y se mantiene actualizado por una comunidad activa de desarrolladores.
\end{sphinxadmonition}

\begin{sphinxadmonition}{note}{Suite de creación 3D BLENDER (\$0)}


\begin{savenotes}\sphinxattablestart
\sphinxthistablewithglobalstyle
\centering
\begin{tabulary}{\linewidth}[t]{T}
\sphinxtoprule
\sphinxtableatstartofbodyhook

\begin{sphinxfigure-in-table}
\centering
\capstart
\noindent\sphinxincludegraphics[width=360\sphinxpxdimen]{{image063}.png}
\sphinxfigcaption{Logo \sphinxstylestrong{BLENDER}}\label{\detokenize{1_guias/10ofimatica/G1012:id46}}\end{sphinxfigure-in-table}\relax
\\
\sphinxbottomrule
\end{tabulary}
\sphinxtableafterendhook\par
\sphinxattableend\end{savenotes}

\sphinxAtStartPar
%
\begin{footnote}[8]\sphinxAtStartFootnote
Nava, J. (2026, February 19). ¿Qué es blender? sus usos y beneficios. IMSED Business \& Tech School. https://imsed.com/postgrados/tecnologia/que\sphinxhyphen{}es\sphinxhyphen{}blender\sphinxhyphen{}usos\sphinxhyphen{}beneficios/
%
\end{footnote} Blender es una suite de \sphinxstylestrong{creación 3D} de código abierto y gratuita que abarca todo el ciclo de producción digital. Permite realizar modelado, renderizado, animación, edición de video, composición y creación de efectos visuales (VFX).

\sphinxAtStartPar
Características principales:
\begin{itemize}
\item {} 
\sphinxAtStartPar
\sphinxstylestrong{Modelado y Esculpido:} Herramientas avanzadas para crear formas complejas y personajes.

\item {} 
\sphinxAtStartPar
\sphinxstylestrong{Animación y Rigging:} Sistema profesional para dar movimiento a objetos y modelos orgánicos.

\item {} 
\sphinxAtStartPar
\sphinxstylestrong{Motores de Renderizado:} Incluye Cycles (trazado de rayos de alta fidelidad) y Eevee (renderizado en tiempo real).

\item {} 
\sphinxAtStartPar
\sphinxstylestrong{Versatilidad:} Al ser multiplataforma y ligero, es utilizado tanto por artistas independientes como por estudios de videojuegos y cine.

\end{itemize}
\end{sphinxadmonition}

\begin{sphinxadmonition}{note}{Edición de sonido Audacity (\$0)}


\begin{savenotes}\sphinxattablestart
\sphinxthistablewithglobalstyle
\centering
\begin{tabulary}{\linewidth}[t]{T}
\sphinxtoprule
\sphinxtableatstartofbodyhook

\begin{sphinxfigure-in-table}
\centering
\capstart
\noindent\sphinxincludegraphics[width=360\sphinxpxdimen]{{image073}.png}
\sphinxfigcaption{Logo \sphinxstylestrong{Audacity}}\label{\detokenize{1_guias/10ofimatica/G1012:id47}}\end{sphinxfigure-in-table}\relax
\\
\sphinxbottomrule
\end{tabulary}
\sphinxtableafterendhook\par
\sphinxattableend\end{savenotes}

\sphinxAtStartPar
%
\begin{footnote}[9]\sphinxAtStartFootnote
Audacity: una forma diferente de estar en la onda. (2023, November 20). INTEF. https://intef.es/observatorio\_tecno/audacity\sphinxhyphen{}una\sphinxhyphen{}forma\sphinxhyphen{}diferente\sphinxhyphen{}de\sphinxhyphen{}estar\sphinxhyphen{}en\sphinxhyphen{}la\sphinxhyphen{}onda/
%
\end{footnote} Audacity es un \sphinxstylestrong{editor y grabador de audio digital} de código abierto y gratuito, ampliamente utilizado para la producción de podcasts, edición de música y procesamiento de sonido. Audacity tiene las siguientes características:
\begin{itemize}
\item {} 
\sphinxAtStartPar
\sphinxstylestrong{Grabación multipista:} Permite capturar audio en vivo a través de micrófonos o mezcladores.

\item {} 
\sphinxAtStartPar
\sphinxstylestrong{Edición sencilla:} Ofrece herramientas para cortar, copiar, pegar y mezclar sonidos con una interfaz intuitiva.

\item {} 
\sphinxAtStartPar
\sphinxstylestrong{Procesamiento de efectos:} Incluye funciones para reducción de ruido, ecualización, normalización y cambio de tono o velocidad.

\item {} 
\sphinxAtStartPar
\sphinxstylestrong{Compatibilidad:} Soporta diversos formatos de archivo como MP3, WAV, AIFF y Ogg Vorbis.

\end{itemize}
\end{sphinxadmonition}

\begin{sphinxadmonition}{note}{Suite ofimática Libreoffice (\$0)}


\begin{savenotes}\sphinxattablestart
\sphinxthistablewithglobalstyle
\centering
\begin{tabulary}{\linewidth}[t]{T}
\sphinxtoprule
\sphinxtableatstartofbodyhook

\begin{sphinxfigure-in-table}
\centering
\capstart
\noindent\sphinxincludegraphics[width=360\sphinxpxdimen]{{image083}.png}
\sphinxfigcaption{\sphinxstylestrong{LibreOffice}}\label{\detokenize{1_guias/10ofimatica/G1012:id48}}\end{sphinxfigure-in-table}\relax
\\
\sphinxbottomrule
\end{tabulary}
\sphinxtableafterendhook\par
\sphinxattableend\end{savenotes}

\sphinxAtStartPar
%
\begin{footnote}[10]\sphinxAtStartFootnote
Onieva, D. (2020, February 17). LibreOffice, la suite gratuita que compite cara a cara con Office. SoftZone. https://www.softzone.es/programas/oficina/libreoffice/
%
\end{footnote} LibreOffice es una suite ofimática libre y de código abierto que incluye:
\begin{itemize}
\item {} 
\sphinxAtStartPar
\sphinxstylestrong{Writer:} Procesador de textos similar a Microsoft Word.

\item {} 
\sphinxAtStartPar
\sphinxstylestrong{Calc:} Hoja de cálculo comparable a Microsoft Excel.

\item {} 
\sphinxAtStartPar
\sphinxstylestrong{Impress:} Programa de presentaciones equivalente a PowerPoint.

\item {} 
\sphinxAtStartPar
\sphinxstylestrong{Base:} Gestor de bases de datos.

\item {} 
\sphinxAtStartPar
\sphinxstylestrong{Draw:} Editor de dibujos y diagramas.

\item {} 
\sphinxAtStartPar
\sphinxstylestrong{Math:} Editor de fórmulas matemáticas.

\end{itemize}

\sphinxAtStartPar
LibreOffice es completamente gratuito, compatible con los formatos de \sphinxstylestrong{Microsoft Office} y está disponible en múltiples idiomas. Es desarrollado por \sphinxstylestrong{The Document Foundation} y cuenta con una gran comunidad de usuarios y desarrolladores que lo mantienen actualizado constantemente.
\end{sphinxadmonition}


\bigskip\hrule\bigskip



\subsubsection{Software propietariØ}
\label{\detokenize{1_guias/10ofimatica/G1012:software-propietario}}
\sphinxAtStartPar
\sphinxfootnotemark[2] El software propietario, es aquel que no permite a los usuarios acceder, modificar o distribuir su código fuente. Los usuarios tienen restricciones legales y técnicas que les impiden ejercer las libertades que caracterizan al software libre.

\begin{sphinxadmonition}{note}{Control total}

\sphinxAtStartPar
El \sphinxstylestrong{software propietario} mantiene el control total sobre el programa en manos de su propietario, lo que puede limitar la libertad de elección y personalización por parte de los usuarios.
\end{sphinxadmonition}

\sphinxAtStartPar
El software propietario se caracteriza por:
\begin{itemize}
\item {} 
\sphinxAtStartPar
El código fuente no está disponible para los usuarios.

\item {} 
\sphinxAtStartPar
Prohibición de copiar, modificar o redistribuir el software.

\item {} 
\sphinxAtStartPar
Generalmente requiere el pago de licencias para su uso.

\item {} 
\sphinxAtStartPar
Las actualizaciones y mejoras dependen exclusivamente del propietario.

\end{itemize}


\bigskip\hrule\bigskip



\subsubsection{Ejemplos de software propietario}
\label{\detokenize{1_guias/10ofimatica/G1012:ejemplos-de-software-propietario}}
\sphinxAtStartPar
A continuación se muestran algunos ejemplos del software propietario:

\begin{sphinxadmonition}{note}{Sistema operativo Microsoft Windows (\$769,999)}


\begin{savenotes}\sphinxattablestart
\sphinxthistablewithglobalstyle
\centering
\begin{tabulary}{\linewidth}[t]{T}
\sphinxtoprule
\sphinxtableatstartofbodyhook

\begin{sphinxfigure-in-table}
\centering
\capstart
\noindent\sphinxincludegraphics[width=300\sphinxpxdimen]{{image093}.png}
\sphinxfigcaption{Logo Microsoft \sphinxstylestrong{Windows}}\label{\detokenize{1_guias/10ofimatica/G1012:id49}}\end{sphinxfigure-in-table}\relax
\\
\sphinxbottomrule
\end{tabulary}
\sphinxtableafterendhook\par
\sphinxattableend\end{savenotes}

\sphinxAtStartPar
%
\begin{footnote}[11]\sphinxAtStartFootnote
Corrales, R. (2024, August 19). Microsoft, el gigante de la informática: historia, Windows, todos sus servicios y la IA de Copilot. Computer Hoy. https://computerhoy.20minutos.es/industria/microsoft\sphinxhyphen{}gigante\sphinxhyphen{}informatica\sphinxhyphen{}historia\sphinxhyphen{}windows\sphinxhyphen{}todos\sphinxhyphen{}servicios\sphinxhyphen{}ia\sphinxhyphen{}copilot\sphinxhyphen{}1399819
%
\end{footnote} Microsoft Windows es el sistema operativo más utilizado en computadoras personales. Sus características principales incluyen:
\begin{itemize}
\item {} 
\sphinxAtStartPar
Interfaz gráfica intuitiva y fácil de usar.

\item {} 
\sphinxAtStartPar
Amplia compatibilidad con hardware y software.

\item {} 
\sphinxAtStartPar
Actualizaciones regulares de seguridad y funcionalidad.

\item {} 
\sphinxAtStartPar
Gran variedad de aplicaciones disponibles.

\end{itemize}

\sphinxAtStartPar
Desarrollado por \sphinxstylestrong{Microsoft Corporation}, Windows es un software privativo que requiere la compra de una licencia para su uso legal. Desde su primera versión en 1985, ha evolucionado constantemente.
\end{sphinxadmonition}

\begin{sphinxadmonition}{note}{Microsoft Office (suite ofimática) (\$459,999/año)}


\begin{savenotes}\sphinxattablestart
\sphinxthistablewithglobalstyle
\centering
\begin{tabulary}{\linewidth}[t]{T}
\sphinxtoprule
\sphinxtableatstartofbodyhook

\begin{sphinxfigure-in-table}
\centering
\capstart
\noindent\sphinxincludegraphics[width=300\sphinxpxdimen]{{image102}.png}
\sphinxfigcaption{Microsoft \sphinxstylestrong{Office}}\label{\detokenize{1_guias/10ofimatica/G1012:id50}}\end{sphinxfigure-in-table}\relax
\\
\sphinxbottomrule
\end{tabulary}
\sphinxtableafterendhook\par
\sphinxattableend\end{savenotes}

\sphinxAtStartPar
%
\begin{footnote}[12]\sphinxAtStartFootnote
Delgado, J. M. (2024, December 15). Funciones ocultas de Microsoft Office que deberías usar ahora mismo para mejorar tu productividad. Computer Hoy. https://computerhoy.20minutos.es/apps/funciones\sphinxhyphen{}ocultas\sphinxhyphen{}microsoft\sphinxhyphen{}office\sphinxhyphen{}deberias\sphinxhyphen{}usar\sphinxhyphen{}ahora\sphinxhyphen{}mismo\sphinxhyphen{}mejorar\sphinxhyphen{}productividad\sphinxhyphen{}1428291
%
\end{footnote} Microsoft Office es una suite de programas de oficina desarrollada por Microsoft que incluye:
\begin{itemize}
\item {} 
\sphinxAtStartPar
\sphinxstylestrong{Word:} Procesador de textos para crear documentos.

\item {} 
\sphinxAtStartPar
\sphinxstylestrong{Excel:} Hoja de cálculo para análisis de datos y cálculos.

\item {} 
\sphinxAtStartPar
\sphinxstylestrong{PowerPoint:} Programa para crear presentaciones.

\item {} 
\sphinxAtStartPar
\sphinxstylestrong{Outlook:} Cliente de correo electrónico y calendario.

\item {} 
\sphinxAtStartPar
\sphinxstylestrong{Access:} Sistema de gestión de bases de datos.

\end{itemize}

\sphinxAtStartPar
Es una de las suites ofimáticas más utilizadas en el mundo empresarial y educativo. Requiere la compra de una licencia o suscripción para su uso legal, y está disponible tanto en versión de escritorio como en la nube (Microsoft 365).
\end{sphinxadmonition}

\begin{sphinxadmonition}{note}{Adobe Photoshop (editor de imágenes) (\$48,971/mes)}


\begin{savenotes}\sphinxattablestart
\sphinxthistablewithglobalstyle
\centering
\begin{tabulary}{\linewidth}[t]{T}
\sphinxtoprule
\sphinxtableatstartofbodyhook

\begin{sphinxfigure-in-table}
\centering
\capstart
\noindent\sphinxincludegraphics[width=360\sphinxpxdimen]{{image111}.png}
\sphinxfigcaption{Logo Adobe \sphinxstylestrong{Photoshop}}\label{\detokenize{1_guias/10ofimatica/G1012:id51}}\end{sphinxfigure-in-table}\relax
\\
\sphinxbottomrule
\end{tabulary}
\sphinxtableafterendhook\par
\sphinxattableend\end{savenotes}

\sphinxAtStartPar
%
\begin{footnote}[13]\sphinxAtStartFootnote
de Luis, E. R. (2026, March 15). En 1990 un estudiante tuvo un problema mostrando imágenes en su Mac. La solución fue crear uno de los softwares más exitosos de la historia. Xataka.com; Xataka. https://www.xataka.com/aplicaciones/1987\sphinxhyphen{}tuvo\sphinxhyphen{}problema\sphinxhyphen{}mostrando\sphinxhyphen{}imagenes\sphinxhyphen{}su\sphinxhyphen{}mac\sphinxhyphen{}asi\sphinxhyphen{}que\sphinxhyphen{}creo\sphinxhyphen{}app\sphinxhyphen{}hoy\sphinxhyphen{}editor\sphinxhyphen{}imagenes\sphinxhyphen{}usado\sphinxhyphen{}historia
%
\end{footnote} Adobe Photoshop es el software de edición de imágenes más popular y potente del mercado. Sus características principales incluyen:
\begin{itemize}
\item {} 
\sphinxAtStartPar
Herramientas avanzadas para edición y manipulación de imágenes.

\item {} 
\sphinxAtStartPar
Trabajo con capas, máscaras y efectos profesionales.

\item {} 
\sphinxAtStartPar
Soporte para múltiples formatos de archivo.

\item {} 
\sphinxAtStartPar
Integración con otros productos de Adobe Creative Suite.

\item {} 
\sphinxAtStartPar
Herramientas de retoque fotográfico profesional.

\end{itemize}

\sphinxAtStartPar
Desarrollado por Adobe Inc., Photoshop es el estándar de la industria para el diseño gráfico, fotografía y arte digital. Requiere una suscripción mensual a través de Adobe Creative Cloud para su uso legal.
\end{sphinxadmonition}

\begin{sphinxadmonition}{note}{AutoCAD (software de diseño) (\$2,201,999/año)}


\begin{savenotes}\sphinxattablestart
\sphinxthistablewithglobalstyle
\centering
\begin{tabulary}{\linewidth}[t]{T}
\sphinxtoprule
\sphinxtableatstartofbodyhook

\begin{sphinxfigure-in-table}
\centering
\capstart
\noindent\sphinxincludegraphics[width=360\sphinxpxdimen]{{image121}.png}
\sphinxfigcaption{Logo Autodesk \sphinxstylestrong{AutoCAD}}\label{\detokenize{1_guias/10ofimatica/G1012:id52}}\end{sphinxfigure-in-table}\relax
\\
\sphinxbottomrule
\end{tabulary}
\sphinxtableafterendhook\par
\sphinxattableend\end{savenotes}

\sphinxAtStartPar
%
\begin{footnote}[14]\sphinxAtStartFootnote
Lautrec, T. (n.d.). ¿Qué es AutoCAD y para qué sirve? Toulouse Lautrec. Retrieved May 7, 2026, from https://www.toulouselautrec.edu.pe/blogs/que\sphinxhyphen{}es\sphinxhyphen{}autocad
%
\end{footnote} AutoCAD es un software de diseño asistido por computadora (CAD) desarrollado por Autodesk. Sus características principales incluyen:
\begin{itemize}
\item {} 
\sphinxAtStartPar
Diseño y dibujo técnico en 2D y 3D.

\item {} 
\sphinxAtStartPar
Herramientas precisas para arquitectura e ingeniería.

\item {} 
\sphinxAtStartPar
Creación de planos y modelos profesionales.

\item {} 
\sphinxAtStartPar
Amplia biblioteca de componentes y símbolos.

\item {} 
\sphinxAtStartPar
Capacidad de colaboración y compartir proyectos.

\end{itemize}

\sphinxAtStartPar
Es ampliamente utilizado en campos como la arquitectura, ingeniería, diseño industrial y construcción. AutoCAD es un software privativo que requiere una suscripción para su uso legal, siendo uno de los estándares de la industria para el diseño técnico profesional.
\end{sphinxadmonition}


\bigskip\hrule\bigskip



\subsubsection{Software pirata}
\label{\detokenize{1_guias/10ofimatica/G1012:software-pirata}}\begin{quote}

\sphinxAtStartPar
\sphinxstyleemphasis{“El software pirata se refiere a copias ilegales de \sphinxstylestrong{programas informáticos} que se distribuyen o utilizan sin la debida autorización de sus propietarios legales. Esta práctica viola los derechos de autor y la propiedad intelectual”.}
\end{quote}

\sphinxAtStartPar
%
\begin{footnote}[15]\sphinxAtStartFootnote
Textos: Jorge Eliécer Lozano / Ilustraciones: Gissell Andreína García Rodríguez. (n.d.). El software pirata representa un riesgo inminente. Edu.co. Retrieved May 7, 2026, from https://notired.univalle.edu.co/el\sphinxhyphen{}software\sphinxhyphen{}pirata\sphinxhyphen{}representa\sphinxhyphen{}un\sphinxhyphen{}riesgo\sphinxhyphen{}inminente/
%
\end{footnote} El \sphinxstylestrong{software pirata} se caracteriza por:
\begin{itemize}
\item {} 
\sphinxAtStartPar
Se obtiene y distribuye sin pagar las licencias correspondientes.

\item {} 
\sphinxAtStartPar
Viola los derechos de autor y la propiedad intelectual.

\item {} 
\sphinxAtStartPar
Puede contener Malware por ejemplo:
\begin{itemize}
\item {} 
\sphinxAtStartPar
\sphinxstylestrong{Virus}: impacta la integridad del sistema. Su objetivo es dañar.

\item {} 
\sphinxAtStartPar
\sphinxstylestrong{Troyanos:} impacta la confidencialidad del sistema. Su objetivo es hacerse pasar por un software confiable pero en realidad esta ejecutando acciones maliciosas.

\end{itemize}

\item {} 
\sphinxAtStartPar
No recibe actualizaciones oficiales ni soporte técnico.

\item {} 
\sphinxAtStartPar
Su uso puede resultar en problemas legales y sanciones.

\end{itemize}

\begin{sphinxadmonition}{note}{ADVERTENCIA}

\sphinxAtStartPar
El uso de \sphinxstylestrong{software pirata} es ilegal y puede tener consecuencias graves, incluyendo:
\begin{itemize}
\item {} 
\sphinxAtStartPar
Multas y sanciones legales (CARCEL Y MULTAS MUY ALTAS).

\item {} 
\sphinxAtStartPar
Riesgos de seguridad informática (VIRUS, TROYANOS).

\item {} 
\sphinxAtStartPar
Pérdida de funcionalidades y actualizaciones (SOFTWARE A MEDIAS, POCO RENDIMIENTO).

\item {} 
\sphinxAtStartPar
Daños a la industria del software y su desarrollo (PERDIDAS ECONÓMICAS, DESPIDOS).

\end{itemize}
\end{sphinxadmonition}


\bigskip\hrule\bigskip



\subsubsection{Suite Ofimática LibreOffice}
\label{\detokenize{1_guias/10ofimatica/G1012:suite-ofimatica-libreoffice}}

\begin{savenotes}\sphinxattablestart
\sphinxthistablewithglobalstyle
\centering
\begin{tabulary}{\linewidth}[t]{T}
\sphinxtoprule
\sphinxtableatstartofbodyhook

\begin{sphinxfigure-in-table}
\centering
\capstart
\noindent\sphinxincludegraphics[width=300\sphinxpxdimen]{{image131}.png}
\sphinxfigcaption{Suite LibreOffice}\label{\detokenize{1_guias/10ofimatica/G1012:id53}}\end{sphinxfigure-in-table}\relax
\\
\sphinxbottomrule
\end{tabulary}
\sphinxtableafterendhook\par
\sphinxattableend\end{savenotes}

\sphinxAtStartPar
%
\begin{footnote}[16]\sphinxAtStartFootnote
Antonio Fernández, B. (2025, May 9). Guía de iniciación 24.8 \sphinxhyphen{} Prefacio. Libreoffice.org. https://books.libreoffice.org/es/GS248/GS24800\sphinxhyphen{}Prefacio.html
%
\end{footnote} \sphinxstylestrong{LibreOffice} es una suite de productividad de oficina de código abierto, disponible de forma gratuita, que es compatible con otras suites de oficina y está disponible en una variedad de plataformas. Su formato de archivo nativo es \sphinxstylestrong{Open Document Format (ODF)} y puede abrir y guardar documentos en muchos formatos, incluidos los utilizados por varias versiones de Microsoft Office.

\sphinxAtStartPar
La suite de \sphinxstylestrong{LibreOffice} esta compuesta por las siguientes herramientas:
\begin{itemize}
\item {} 
\sphinxAtStartPar
El procesador de textos \sphinxcode{\sphinxupquote{Writer}}.

\item {} 
\sphinxAtStartPar
La hoja de cálculo \sphinxcode{\sphinxupquote{Calc}}.

\item {} 
\sphinxAtStartPar
El programa de presentaciones \sphinxcode{\sphinxupquote{Impress}}.

\item {} 
\sphinxAtStartPar
El programa de dibujo vectorial \sphinxcode{\sphinxupquote{Draw}}.

\item {} 
\sphinxAtStartPar
La base de datos \sphinxcode{\sphinxupquote{Base}}.

\item {} 
\sphinxAtStartPar
El editor de ecuaciones \sphinxcode{\sphinxupquote{Math}}.

\end{itemize}

\sphinxAtStartPar
LibreOffice es una suite ofimática multi\sphinxhyphen{}plataforma, lo que quiere decir que se puede ejecutar en diferentes plataformas y sistemas operativos. Por consiguiente usted puede disfrutar del mismo programa en \sphinxstylestrong{Windows}, \sphinxstylestrong{GNU/Linux} o \sphinxstylestrong{MacOS}, facilitando la portabilidad entre diferentes sistemas de los documentos generados.


\bigskip\hrule\bigskip



\subsubsection{Entidades que usan LibreOffice}
\label{\detokenize{1_guias/10ofimatica/G1012:entidades-que-usan-libreoffice}}
\sphinxAtStartPar
%
\begin{footnote}[17]\sphinxAtStartFootnote
(N.d.). Libreoffice.org. Retrieved May 17, 2026, from https://es.libreoffice.org/who\sphinxhyphen{}uses\sphinxhyphen{}libreoffice/
%
\end{footnote} Millones de personas utilizan LibreOffice todos los días en sus hogares, negocios, organizaciones benéficas y sectores gubernamentales. De entre estos grupos de usuarios pueden destacarse:


\begin{savenotes}\sphinxattablestart
\sphinxthistablewithglobalstyle
\centering
\begin{tabulary}{\linewidth}[t]{TT}
\sphinxtoprule
\sphinxstyletheadfamily 
\sphinxAtStartPar
País
&\sphinxstyletheadfamily 
\sphinxAtStartPar
Descripción
\\
\sphinxmidrule
\sphinxtableatstartofbodyhook
\sphinxAtStartPar
\sphinxstylestrong{FRANCIA}
&
\sphinxAtStartPar
El equipo de trabajo interministerial para el software libre del Gobierno francés, ejecuta LibreOffice en aproximadamente 500,000 equipos. El software se emplea en varios ministerios, incluidos los de Energía, Defensa, Agricultura y Educación.
\\
\sphinxhline
\sphinxAtStartPar
\sphinxstylestrong{ESPAÑA}
&
\sphinxAtStartPar
La administración de la comunidad autónoma de Valencia, en España, ha instalado LibreOffice en 120,000 equipos. Con esto, la administración ha ganado independencia frente a proveedores de TI y ha disminuido los gastos en licencias de software privativo.
\\
\sphinxhline
\sphinxAtStartPar
\sphinxstylestrong{ITALIA}
&
\sphinxAtStartPar
El Ministerio de Defensa de Italia se encuentra en un proceso de transición hacia LibreOffice y el formato OpenDocument (ODF) en más de 100,000 equipos. El Ministerio, asimismo, ha desarrollado cursos en línea para ayudar con el cambio a LibreOffice.
\\
\sphinxhline
\sphinxAtStartPar
\sphinxstylestrong{ALEMANIA}
&
\sphinxAtStartPar
En Alemania, la ciudad de Múnich ha migrado más de 15,000 equipos a LibreOffice. Esto forma parte de una transición más ambiciosa hacia el software libre y de código abierto (FOSS), que incluye el sistema operativo GNU/Linux.
\\
\sphinxhline
\sphinxAtStartPar
\sphinxstylestrong{TAIWAN}
&
\sphinxAtStartPar
El Ministerio de Finanzas de Taiwán ha instalado LibreOffice en más de 24,000 PC. Con esta iniciativa, el formato estandarizado OpenDocument está convirtiéndose en la norma para el intercambio de datos entre los diversos departamentos.
\\
\sphinxhline
\sphinxAtStartPar
\sphinxstylestrong{BRASIL}
&
\sphinxAtStartPar
En Brasil, la UNESP ―Universidade Estadual Paulista― ha realizado más de 10,000 instalaciones de LibreOffice como parte de una migración general hacia el software libre y de código abierto, la cual incluye el sistema operativo GNU/Linux.
\\
\sphinxhline
\sphinxAtStartPar
\sphinxstylestrong{COLOMBIA}
&
\sphinxAtStartPar
???
\\
\sphinxbottomrule
\end{tabulary}
\sphinxtableafterendhook\par
\sphinxattableend\end{savenotes}


\bigskip\hrule\bigskip



\subsubsection{Herramientas de LibreOffice}
\label{\detokenize{1_guias/10ofimatica/G1012:herramientas-de-libreoffice}}
\sphinxAtStartPar
\sphinxfootnotemark[16] A continuación se explican las diferentes herramientas de LibreOffice:

\begin{sphinxadmonition}{note}{LibreOffice Writer (procesador de textos)}


\begin{savenotes}\sphinxattablestart
\sphinxthistablewithglobalstyle
\centering
\begin{tabulary}{\linewidth}[t]{T}
\sphinxtoprule
\sphinxtableatstartofbodyhook

\begin{sphinxfigure-in-table}
\centering
\capstart
\noindent\sphinxincludegraphics[width=270\sphinxpxdimen]{{image141}.png}
\sphinxfigcaption{LibreOffice Writer}\label{\detokenize{1_guias/10ofimatica/G1012:id54}}\end{sphinxfigure-in-table}\relax
\\
\sphinxbottomrule
\end{tabulary}
\sphinxtableafterendhook\par
\sphinxattableend\end{savenotes}

\sphinxAtStartPar
\sphinxstylestrong{Writer} es una herramienta rica en funciones para crear cartas, libros, informes, boletines, folletos y otros documentos. Puede insertar gráficos y objetos de otros componentes de LibreOffice en documentos de Writer.

\sphinxAtStartPar
\sphinxstylestrong{Writer} puede exportar archivos a HTML, XHTML, XML, PDF y EPUB, puede guardar archivos en muchos formatos incluidas varias versiones de archivo de Microsoft Word y conectarse con su cliente de correo electrónico.
\end{sphinxadmonition}

\begin{sphinxadmonition}{note}{LibreOffice Calc (hoja de cálculo)}


\begin{savenotes}\sphinxattablestart
\sphinxthistablewithglobalstyle
\centering
\begin{tabulary}{\linewidth}[t]{T}
\sphinxtoprule
\sphinxtableatstartofbodyhook

\begin{sphinxfigure-in-table}
\centering
\capstart
\noindent\sphinxincludegraphics[width=190\sphinxpxdimen]{{image15}.png}
\sphinxfigcaption{LibreOffice Calc}\label{\detokenize{1_guias/10ofimatica/G1012:id55}}\end{sphinxfigure-in-table}\relax
\\
\sphinxbottomrule
\end{tabulary}
\sphinxtableafterendhook\par
\sphinxattableend\end{savenotes}

\sphinxAtStartPar
\sphinxstylestrong{Calc} tiene todas las funciones avanzadas de análisis, gráficos y toma de decisiones que se esperan de una hoja de cálculo de alta gama. Incluye más de 500 funciones para operaciones financieras, estadísticas y matemáticas, entre otras. El Gestor de escenarios proporciona un análisis de «qué pasaría si».

\sphinxAtStartPar
\sphinxstylestrong{Calc} genera gráficos en 2D y 3D, que se pueden integrar en otros documentos de LibreOffice. También puede abrir y trabajar y guardar hojas de cálculo de Microsoft Excel y otros formatos de hoja de cálculo, así como exportar hojas de cálculo a varios formatos, como archivos de valores separados por comas (CSV), PDF y HTML.
\end{sphinxadmonition}

\begin{sphinxadmonition}{note}{LibreOffice Impress (presentaciones)}


\begin{savenotes}\sphinxattablestart
\sphinxthistablewithglobalstyle
\centering
\begin{tabulary}{\linewidth}[t]{T}
\sphinxtoprule
\sphinxtableatstartofbodyhook

\begin{sphinxfigure-in-table}
\centering
\capstart
\noindent\sphinxincludegraphics[width=270\sphinxpxdimen]{{image16}.png}
\sphinxfigcaption{LibreOffice Impress}\label{\detokenize{1_guias/10ofimatica/G1012:id56}}\end{sphinxfigure-in-table}\relax
\\
\sphinxbottomrule
\end{tabulary}
\sphinxtableafterendhook\par
\sphinxattableend\end{savenotes}

\sphinxAtStartPar
\sphinxstylestrong{Impress} proporciona todas las herramientas de presentación multimedia comunes, como efectos especiales, animaciones y herramientas de dibujo. Está integrado con las capacidades gráficas avanzadas de los componentes de \sphinxstylestrong{LibreOffice Draw} y \sphinxstylestrong{Math}.

\sphinxAtStartPar
Las presentaciones de diapositivas se pueden mejorar utilizando texto de efectos especiales de Fontwork, así como clips de sonido y vídeo. Impress puede abrir, editar y guardar presentaciones de Microsoft PowerPoint y también puede guardar su trabajo en numerosos formatos gráficos.
\end{sphinxadmonition}

\begin{sphinxadmonition}{note}{LibreOffice Draw (Dibujo vectorial)}


\begin{savenotes}\sphinxattablestart
\sphinxthistablewithglobalstyle
\centering
\begin{tabulary}{\linewidth}[t]{T}
\sphinxtoprule
\sphinxtableatstartofbodyhook

\begin{sphinxfigure-in-table}
\centering
\capstart
\noindent\sphinxincludegraphics[width=230\sphinxpxdimen]{{image17}.png}
\sphinxfigcaption{LibreOffice Draw}\label{\detokenize{1_guias/10ofimatica/G1012:id57}}\end{sphinxfigure-in-table}\relax
\\
\sphinxbottomrule
\end{tabulary}
\sphinxtableafterendhook\par
\sphinxattableend\end{savenotes}

\sphinxAtStartPar
\sphinxstylestrong{Draw} es una herramienta de dibujo vectorial que puede producir desde simples diagramas o diagramas de flujo hasta ilustraciones en 3D. Su característica Conectores Inteligentes le permite definir sus propios puntos de conexión. Puede usar \sphinxstylestrong{Draw} para crear dibujos y emplearlos en cualquiera de los componentes de LibreOffice, y puede crear su propia imagen prediseñada para luego agregarla a la Galería.

\sphinxAtStartPar
\sphinxstylestrong{Draw} puede importar archivos gráficos en muchos formatos y guardar su trabajo en muchos formatos, incluidos PNG, GIF, JPEG, BMP, TIFF, SVG, HTML, PDF y WebP.
\end{sphinxadmonition}

\begin{sphinxadmonition}{note}{LibreOffice Math (Editor de fórmulas)}


\begin{savenotes}\sphinxattablestart
\sphinxthistablewithglobalstyle
\centering
\begin{tabulary}{\linewidth}[t]{T}
\sphinxtoprule
\sphinxtableatstartofbodyhook

\begin{sphinxfigure-in-table}
\centering
\capstart
\noindent\sphinxincludegraphics[width=240\sphinxpxdimen]{{image18}.png}
\sphinxfigcaption{LibreOffice Math}\label{\detokenize{1_guias/10ofimatica/G1012:id58}}\end{sphinxfigure-in-table}\relax
\\
\sphinxbottomrule
\end{tabulary}
\sphinxtableafterendhook\par
\sphinxattableend\end{savenotes}

\sphinxAtStartPar
Math es un editor de fórmulas o ecuaciones. Puede usarlo para crear ecuaciones complejas que incluyen símbolos o caracteres que no están disponibles en conjuntos de fuentes estándar. Si bien se usa más comúnmente para crear fórmulas en otros documentos, como archivos Writer e Impress, Math también puede funcionar como una herramienta independiente. Puede guardar fórmulas en el formato estándar de \sphinxstylestrong{Lenguaje de Marcado Matemático (MathML)} para incluirlas en páginas web u otros documentos creados con otros programas.
\end{sphinxadmonition}


\bigskip\hrule\bigskip



\subsubsection{Ventajas de LibreOffice}
\label{\detokenize{1_guias/10ofimatica/G1012:ventajas-de-libreoffice}}
\sphinxAtStartPar
\sphinxfootnotemark[16] A continuación se presentan las ventajas que tiene el uso de LibreOffice frente a otras suites ofimáticas:
\begin{itemize}
\item {} 
\sphinxAtStartPar
\sphinxstylestrong{Sin tarifas de licencia:} LibreOffice es gratuito para que cualquiera lo use y distribuya sin coste alguno. LibreOffice incluye muchas funciones (como exportación a PDF) de manera gratuita y que en otras suites de oficina son de pago. Con LibreOffice no hay cargos ocultos ni ahora, ni en el futuro.

\item {} 
\sphinxAtStartPar
\sphinxstylestrong{Código abierto:} Puede distribuir, copiar y modificar el software tanto como desee, de acuerdo con las licencias de código abierto de LibreOffice.

\item {} 
\sphinxAtStartPar
\sphinxstylestrong{Multiplataforma:} LibreOffice se ejecuta en varias arquitecturas de hardware y en varios sistemas operativos, como Microsoft Windows, macOS y Linux.

\item {} 
\sphinxAtStartPar
\sphinxstylestrong{Amplio soporte de idiomas:} La interfaz de usuario de LibreOffice, diccionarios ortográficos, separación de sílabas y sinónimos, están disponibles en más de 100 idiomas y dialectos. LibreOffice también ofrece compatibilidad con los idiomas de diseño de texto complejo (CTL) y de escritura de derecha a izquierda (RTL) (como urdu, hebreo y árabe).

\item {} 
\sphinxAtStartPar
\sphinxstylestrong{Interfaz de usuario consistente:} Todos los componentes tienen una apariencia similar, lo que los hace fáciles de usar y aprender.

\item {} 
\sphinxAtStartPar
\sphinxstylestrong{Integración:} Los componentes de LibreOffice están integrados entre sí, todos comparten un corrector ortográfico común y otras herramientas, que se utilizan de forma coherente en toda la suite. Por ejemplo, las herramientas de dibujo disponibles en Writer también se encuentran en Calc, y con versiones similares pero mejoradas en Impress y Draw.
\begin{itemize}
\item {} 
\sphinxAtStartPar
No necesita saber qué aplicación se utilizó para crear un archivo en particular. Por ejemplo, puede abrir un archivo Draw desde Writer; se abrirá automáticamente en Draw.

\end{itemize}

\item {} 
\sphinxAtStartPar
\sphinxstylestrong{Granularidad:} Por lo general, si cambia una opción, afecta a todos los componentes. Sin embargo, las opciones de LibreOffice se pueden configurar en el ámbito de componente o incluso en el ámbito de documento.

\item {} 
\sphinxAtStartPar
\sphinxstylestrong{Compatibilidad de archivos:} Además de sus formatos nativos (OpenDocument), LibreOffice puede y guardar archivos en muchos formatos comunes, entre otros: Microsoft Office, HTML, XML, WordPerfect, Lotus 1\sphinxhyphen{}2\sphinxhyphen{}3 y PDF.

\item {} 
\sphinxAtStartPar
\sphinxstylestrong{Sin bloqueo de proveedores:} LibreOffice utiliza OpenDocument, un formato de archivo XML (eXtensible Markup Language) desarrollado como estándar industrial por OASIS (Organization for the Advancement of Structured Information Standards). Estos archivos se pueden descomprimir y leer fácilmente con cualquier editor de texto, y su estructura es abierta y está documentada.

\item {} 
\sphinxAtStartPar
\sphinxstylestrong{Usted tiene voz:} Las mejoras, las correcciones de software y las fechas de lanzamiento están impulsadas por la comunidad. Puede unirse a la comunidad e influir sobre el curso de LibreOffice.

\end{itemize}


\bigskip\hrule\bigskip



\subsection{A02: Taller (+0,8)}
\label{\detokenize{1_guias/10ofimatica/G1012:a02-taller-0-8}}
\sphinxAtStartPar
\DUrole{sd-sphinx-override}{\DUrole{sd-badge}{\DUrole{sd-bg-warning}{\DUrole{sd-bg-text-warning}{Actividad práctica}}}}
\begin{enumerate}
\sphinxsetlistlabels{\arabic}{enumi}{enumii}{}{.}%
\item {} 
\sphinxAtStartPar
En su cuaderno responda:
\begin{itemize}
\item {} 
\sphinxAtStartPar
¿Qué es el software?

\item {} 
\sphinxAtStartPar
¿Cuales son los tipos de software?

\item {} 
\sphinxAtStartPar
¿Cómo se puede clasificar el software?

\item {} 
\sphinxAtStartPar
¿Qué es el software libre?

\item {} 
\sphinxAtStartPar
¿Cuáles son las \sphinxstylestrong{4 libertades esenciales} que debe cumplir un software, (para que sea considerado software libre)?

\end{itemize}

\item {} 
\sphinxAtStartPar
En su cuaderno, escriba una lista de 15 ejemplos de software que usted conozca (de smartphone, de computadores, de internet, etc.) NO OLVIDE HACER UNA BREVE DESCRIPCIÓN.

\item {} 
\sphinxAtStartPar
Para cada uno de los ejemplos escritos anteriormente, identifique si es \sphinxstylestrong{software de sistema}, \sphinxstylestrong{software de aplicación}, \sphinxstylestrong{software libre}, \sphinxstylestrong{software propietario} o \sphinxstylestrong{software pirata} y explique el por qué.

\item {} 
\sphinxAtStartPar
Analice los siguientes casos y en su cuaderno explique si el software es libre o no. Identifique si cumple las \sphinxstylestrong{4 libertades esenciales} y explique el por qué:
\begin{quote}

\sphinxAtStartPar
\sphinxstylestrong{CASO 01}: “La institución educativa utiliza el \sphinxstylestrong{software X} para gestionar sus datos. Al adquirirlo, la comunidad recibió no solo el instalador, sino también el código fuente del programa. Debido a una necesidad específica, el docente de tecnología pudo estudiar cómo funciona el software y modificarlo para que acepte nuevos parámetros de entrada que antes no permitía. Posteriormente, el docente distribuyó copias de esta versión mejorada a otros colegios de la región para que ellos también pudieran beneficiarse de los cambios sin restricciones legales.”
\end{quote}
\begin{quote}

\sphinxAtStartPar
\sphinxstylestrong{CASO 02}: “El \sphinxstylestrong{Software Y} se distribuye bajo un modelo de suscripción de pago obligatorio que requiere una cuenta personal para su activación. Unos usuarios en Rusia lo utilizan para capacitación, pero solo pueden usarlo en un máximo de cinco dispositivos simultáneamente según su plan. Cuando usuarios en España intentaron obtener la versión modificada de Rusia, descubrieron que la licencia de uso prohíbe explícitamente cualquier redistribución o copia del programa a terceros. Además, aunque los usuarios son dueños de los datos que generan, no tienen acceso al código fuente ni el control sobre el software, por lo que no pueden estudiar cómo funciona ni realizar mejoras legales para adaptarlo a sus necesidades.”
\end{quote}
\begin{quote}

\sphinxAtStartPar
\sphinxstylestrong{CASO 03}: “Al \sphinxstylestrong{Software Z} se le implementaron diversas adaptaciones gracias a que la comunidad tuvo acceso total a su código fuente, permitiendo que todas las modificaciones funcionaran sin bloqueos legales. Aunque inicialmente los clientes trabajaban gastando más tiempo del habitual, la libertad de examinar el software permitió a los programadores desentrañar su funcionamiento para adecuar y optimizar las funcionalidades críticas. Al final, estas mejoras se compartieron legalmente, permitiendo que usuarios en España lo eligieran para reemplazar un software privativo que, a diferencia del Software Z, prohibía por contrato su libre distribución y estudio.”
\end{quote}

\item {} 
\sphinxAtStartPar
A partir de los \sphinxstylestrong{ejemplos de software libre}, en su cuaderno diligencie el siguiente crucigrama:


\begin{savenotes}\sphinxattablestart
\sphinxthistablewithglobalstyle
\centering
\begin{tabulary}{\linewidth}[t]{T}
\sphinxtoprule
\sphinxtableatstartofbodyhook

\begin{sphinxfigure-in-table}
\centering
\capstart
\noindent\sphinxincludegraphics[width=500\sphinxpxdimen]{{image20}.png}
\sphinxfigcaption{\sphinxstylestrong{Crucigrama} ejemplos software libre}\label{\detokenize{1_guias/10ofimatica/G1012:id59}}\end{sphinxfigure-in-table}\relax
\\
\sphinxhline
\sphinxAtStartPar
\sphinxstylestrong{1.} Mismas funcionalidades de Microsoft Office pero es software libre.
\\
\sphinxhline
\sphinxAtStartPar
\sphinxstylestrong{2.} No requiere Codecs para funcionar.
\\
\sphinxhline
\sphinxAtStartPar
\sphinxstylestrong{3.} Este sistema fue creado inicialmente por Linus Torvalds.
\\
\sphinxhline
\sphinxAtStartPar
\sphinxstylestrong{4.} Es una alternativa libre al navegador Google CHROME.
\\
\sphinxhline
\sphinxAtStartPar
\sphinxstylestrong{5.} Mismas funcionalidades de Photoshop pero es software libre.
\\
\sphinxhline
\sphinxAtStartPar
\sphinxstylestrong{6.} Se usa para crear logotipos.
\\
\sphinxbottomrule
\end{tabulary}
\sphinxtableafterendhook\par
\sphinxattableend\end{savenotes}

\item {} 
\sphinxAtStartPar
En su cuaderno, REALICE una \sphinxstylestrong{descripción} de cada uno de los ejemplos de software libre, encontrados en el anterior crucigrama.

\item {} 
\sphinxAtStartPar
Desarrolle los siguientes cálculos en su cuaderno:
\begin{quote}

\sphinxAtStartPar
Una empresa necesita adquirir 11 licencias con \sphinxstylestrong{Windows 11} para poder instalar en 11 equipos. El presupuesto asignado fue \$10,000,000. Al final se dieron cuenta que solo 5 equipos cumplían con los requisitos para poder instalar Windows 11. Después de hacer cuentas, ¿Cuánta plata le sobro a la empresa?
\begin{itemize}
\item {} 
\sphinxAtStartPar
\sphinxstylestrong{RESPUESTA:} \$6,150,000

\end{itemize}
\end{quote}
\begin{quote}

\sphinxAtStartPar
Una familia quiere hacer un cálculo para identificar cuanto tienen que invertir en 5 años por tener activa una suscripción a \sphinxstylestrong{Microsoft Office}. Cada año la suscripción aumenta un 3\%.
\begin{itemize}
\item {} 
\sphinxAtStartPar
\sphinxstylestrong{RESPUESTA:} \$2,427,122.84

\end{itemize}
\end{quote}
\begin{quote}

\sphinxAtStartPar
Un diseñador adquirió una suscripción de Adobe Photoshop por un año en diferentes intervalos. Uso Photoshop de ENERO a ABRIL. De MAYO a AGOSTO no uso photoshop porque estaba en vacaciones. De SEPTIEMBRE a NOVIEMBRE Pago nuevamente. En Diciembre la renovó por un año más. ¿Cuánto gasto el diseñador por usar la licencia en el primer año?
\begin{itemize}
\item {} 
\sphinxAtStartPar
\sphinxstylestrong{RESPUESTA:} \$391,768

\end{itemize}
\end{quote}
\begin{quote}

\sphinxAtStartPar
Un profesional independiente decide contratar la suscripción \sphinxstylestrong{Google ONE \sphinxhyphen{} AI PLUS} para potenciar su trabajo con inteligencia artificial. Si planea utilizar el servicio de manera ininterrumpida durante un año y 3/4, ¿cuál será el valor total de su inversión?
\begin{itemize}
\item {} 
\sphinxAtStartPar
\sphinxstylestrong{RESPUESTA:} \$358,200

\end{itemize}
\end{quote}
\begin{quote}

\sphinxAtStartPar
Un grupo de 4 amigos desea adquirir una suscripción de \sphinxstylestrong{Microsoft 365 Family} para compartir los gastos de forma equitativa durante un año. Si cada uno aporta la misma cantidad, ¿cuánto dinero debe pagar cada integrante para cubrir el costo total de la licencia anual?
\begin{itemize}
\item {} 
\sphinxAtStartPar
\sphinxstylestrong{RESPUESTA:} \$114,999

\end{itemize}
\end{quote}

\item {} 
\sphinxAtStartPar
Lea las \sphinxstylestrong{características} del software libre, software propietario y el software pirata. Luego complete el siguiente cuadro en su cuaderno:


\begin{savenotes}\sphinxattablestart
\sphinxthistablewithglobalstyle
\centering
\begin{tabulary}{\linewidth}[t]{TTT}
\sphinxtoprule
\sphinxstyletheadfamily 
\sphinxAtStartPar
SOFTWARE LIBRE
&\sphinxstyletheadfamily 
\sphinxAtStartPar
SOFTWARE PROPIETARIO
&\sphinxstyletheadfamily 
\sphinxAtStartPar
SOFTWARE PIRATA
\\
\sphinxmidrule
\sphinxtableatstartofbodyhook
\sphinxAtStartPar
Impulsado por la \sphinxstylestrong{free Software Foundation}
&
\sphinxAtStartPar
Impulsado por la \sphinxstylestrong{propia empresa} con fines lucrativos
&
\sphinxAtStartPar
Impulsado por \sphinxstylestrong{personas} u \sphinxstylestrong{organizaciones} con fines maliciosos
\\
\sphinxhline
\sphinxAtStartPar
.
&
\sphinxAtStartPar

&
\sphinxAtStartPar

\\
\sphinxhline
\sphinxAtStartPar
.
&
\sphinxAtStartPar

&
\sphinxAtStartPar

\\
\sphinxhline
\sphinxAtStartPar
.
&
\sphinxAtStartPar

&
\sphinxAtStartPar

\\
\sphinxhline
\sphinxAtStartPar
.
&
\sphinxAtStartPar

&
\sphinxAtStartPar

\\
\sphinxbottomrule
\end{tabulary}
\sphinxtableafterendhook\par
\sphinxattableend\end{savenotes}

\item {} 
\sphinxAtStartPar
Responda en su cuaderno:
\begin{itemize}
\item {} 
\sphinxAtStartPar
¿Cuáles son las características de un software pirata?

\item {} 
\sphinxAtStartPar
¿Qué consecuencias tiene un sistema que tenga VIRUS y TROYANOS?

\end{itemize}

\item {} 
\sphinxAtStartPar
LEA el siguiente caso:
\begin{quote}

\sphinxAtStartPar
Una empresa está afectada por el uso de \sphinxstylestrong{software pirata}.
\begin{itemize}
\item {} 
\sphinxAtStartPar
Una auditoria del gobierno multó a la empresa por \$100,000,000.00 COP.

\item {} 
\sphinxAtStartPar
La información sensible como fórmulas secretas o datos personales de los empleados está en riesgo de ser expuesta.

\item {} 
\sphinxAtStartPar
Los PCs de la empresa están lentos y se reinician solos. Esto altera el flujo normal de trabajo.

\end{itemize}
\end{quote}

\sphinxAtStartPar
En una hoja de cuaderno responda:
\begin{itemize}
\item {} 
\sphinxAtStartPar
¿Qué medidas tomaría usted para evitar que la empresa quiebre?

\end{itemize}

\item {} 
\sphinxAtStartPar
En una hoja de cuaderno, complete la siguiente tabla:


\begin{savenotes}\sphinxattablestart
\sphinxthistablewithglobalstyle
\centering
\begin{tabulary}{\linewidth}[t]{TTTT}
\sphinxtoprule
\sphinxstyletheadfamily 
\sphinxAtStartPar
.
&\sphinxstyletheadfamily 
\sphinxAtStartPar
SOFTWARE LIBRE
&\sphinxstyletheadfamily 
\sphinxAtStartPar
SOFTWARE PROPIETARIO
&\sphinxstyletheadfamily 
\sphinxAtStartPar
Descripción
\\
\sphinxmidrule
\sphinxtableatstartofbodyhook
\sphinxAtStartPar
\sphinxstylestrong{Edición de fotos}
&
\sphinxAtStartPar
GIMP
&
\sphinxAtStartPar
PHOTOSHOP
&
\sphinxAtStartPar

\\
\sphinxhline
\sphinxAtStartPar
\sphinxstylestrong{Procesamiento de textos}
&
\sphinxAtStartPar

&
\sphinxAtStartPar

&
\sphinxAtStartPar

\\
\sphinxhline
\sphinxAtStartPar
\sphinxstylestrong{Sistema Operativo}
&
\sphinxAtStartPar

&
\sphinxAtStartPar

&
\sphinxAtStartPar

\\
\sphinxhline
\sphinxAtStartPar
\sphinxstylestrong{Hojas de cálculo}
&
\sphinxAtStartPar

&
\sphinxAtStartPar

&
\sphinxAtStartPar

\\
\sphinxhline
\sphinxAtStartPar
\sphinxstylestrong{Navegadores Web}
&
\sphinxAtStartPar

&
\sphinxAtStartPar

&
\sphinxAtStartPar

\\
\sphinxhline
\sphinxAtStartPar
\sphinxstylestrong{Presentaciones}
&
\sphinxAtStartPar

&
\sphinxAtStartPar

&
\sphinxAtStartPar

\\
\sphinxhline
\sphinxAtStartPar
\sphinxstylestrong{Reproductor Multimedia}
&
\sphinxAtStartPar

&
\sphinxAtStartPar

&
\sphinxAtStartPar

\\
\sphinxhline
\sphinxAtStartPar
\sphinxstylestrong{Gestión de calendario y agendas}
&
\sphinxAtStartPar

&
\sphinxAtStartPar

&
\sphinxAtStartPar

\\
\sphinxbottomrule
\end{tabulary}
\sphinxtableafterendhook\par
\sphinxattableend\end{savenotes}

\item {} 
\sphinxAtStartPar
RESPONDA en su cuaderno las siguientes preguntas:
\begin{itemize}
\item {} 
\sphinxAtStartPar
¿Que es LibreOffice?

\item {} 
\sphinxAtStartPar
¿Qué significa \sphinxstylestrong{ODF}?

\end{itemize}

\item {} 
\sphinxAtStartPar
COPIE las siguientes frases en su cuaderno e INDIQUE si es \sphinxstylestrong{falso}, \sphinxstylestrong{verdadero} y por qué:
\begin{itemize}
\item {} 
\sphinxAtStartPar
\sphinxstylestrong{Calc} es el componente de LibreOffice que se utiliza para la creación y edición de hojas de cálculo. \_\_\_

\item {} 
\sphinxAtStartPar
\sphinxstylestrong{LibreOffice} es una suite ofimática exclusiva para usuarios del sistema operativo Windows. \_\_\_

\item {} 
\sphinxAtStartPar
El programa de la suite destinado al diseño y edición de dibujo vectorial recibe el nombre de \sphinxstylestrong{Math}. \_\_\_

\item {} 
\sphinxAtStartPar
Para la gestión y creación de bases de datos, la suite de LibreOffice incorpora la herramienta llamada \sphinxstylestrong{Impress}. \_\_\_

\item {} 
\sphinxAtStartPar
\sphinxstylestrong{Writer} es el procesador de textos que se incluye dentro del conjunto de programas de LibreOffice. \_\_\_

\item {} 
\sphinxAtStartPar
Al ser una suite \sphinxstylestrong{multi\sphinxhyphen{}plataforma}, los documentos generados en LibreOffice no se pueden abrir en sistemas operativos
diferentes al de origen. \_\_\_

\item {} 
\sphinxAtStartPar
El editor que permite trabajar con ecuaciones dentro de esta suite ofimática se llama \sphinxstylestrong{Calc}. \_\_\_

\item {} 
\sphinxAtStartPar
El programa de presentaciones multimedia de LibreOffice se denomina \sphinxstylestrong{Base}. \_\_\_

\item {} 
\sphinxAtStartPar
\sphinxstylestrong{GNU/Linux} y \sphinxstylestrong{MacOS} son dos de los sistemas operativos en los que se puede ejecutar LibreOffice. \_\_\_

\item {} 
\sphinxAtStartPar
El término \sphinxstylestrong{“multi\sphinxhyphen{}plataforma”} significa que el programa solo funciona de manera óptima si se está conectado a una red de servidores local. \_\_\_

\end{itemize}

\item {} 
\sphinxAtStartPar
En su cuaderno complete la siguiente tabla:


\begin{savenotes}\sphinxattablestart
\sphinxthistablewithglobalstyle
\centering
\begin{tabulary}{\linewidth}[t]{TTTT}
\sphinxtoprule
\sphinxstyletheadfamily 
\sphinxAtStartPar
PAÍS
&\sphinxstyletheadfamily 
\sphinxAtStartPar
¿En donde se emplea LibreOffice?
&\sphinxstyletheadfamily 
\sphinxAtStartPar
¿Número de Instalaciones de LibreOffice?
&\sphinxstyletheadfamily 
\sphinxAtStartPar
¿Para que se emplea LibreOffice?
\\
\sphinxmidrule
\sphinxtableatstartofbodyhook
\sphinxAtStartPar
Francia
&
\sphinxAtStartPar

&
\sphinxAtStartPar

&
\sphinxAtStartPar

\\
\sphinxhline
\sphinxAtStartPar
España
&
\sphinxAtStartPar

&
\sphinxAtStartPar

&
\sphinxAtStartPar

\\
\sphinxhline
\sphinxAtStartPar
Italia
&
\sphinxAtStartPar

&
\sphinxAtStartPar

&
\sphinxAtStartPar

\\
\sphinxhline
\sphinxAtStartPar
Alemania
&
\sphinxAtStartPar

&
\sphinxAtStartPar

&
\sphinxAtStartPar

\\
\sphinxhline
\sphinxAtStartPar
Taiwan
&
\sphinxAtStartPar

&
\sphinxAtStartPar

&
\sphinxAtStartPar

\\
\sphinxhline
\sphinxAtStartPar
Brasil
&
\sphinxAtStartPar

&
\sphinxAtStartPar

&
\sphinxAtStartPar

\\
\sphinxhline
\sphinxAtStartPar
Colombia
&
\sphinxAtStartPar

&
\sphinxAtStartPar

&
\sphinxAtStartPar

\\
\sphinxbottomrule
\end{tabulary}
\sphinxtableafterendhook\par
\sphinxattableend\end{savenotes}

\item {} 
\sphinxAtStartPar
En su cuaderno, diligencie la siguiente tabla colocando el nombre de la herramienta contraparte de Microsoft Office frente a LibreOffice. Explique también la función que tiene dichas herramientas. (\sphinxstylestrong{Si no existe contraparte, escribir N/A}).
\begin{quote}

\sphinxAtStartPar
\sphinxstylestrong{Ejemplo:} La contraparte de \sphinxstylestrong{Microsoft Word} es \sphinxstylestrong{LibreOffice Writer.} Por tanto Microsoft Word como LibreOffice Writer funcionan para edición y formato de textos.
\end{quote}


\begin{savenotes}\sphinxattablestart
\sphinxthistablewithglobalstyle
\centering
\begin{tabulary}{\linewidth}[t]{TTT}
\sphinxtoprule
\sphinxstyletheadfamily 
\sphinxAtStartPar
Microsoft Office
&\sphinxstyletheadfamily 
\sphinxAtStartPar
LibreOffice
&\sphinxstyletheadfamily 
\sphinxAtStartPar
Función
\\
\sphinxmidrule
\sphinxtableatstartofbodyhook
\sphinxAtStartPar
Microsoft Word
&
\sphinxAtStartPar
LibreOffice Writer
&
\sphinxAtStartPar
Tanto Microsoft Word como LibreOffice Writer funcionan para edición y formato de textos
\\
\sphinxhline
\sphinxAtStartPar
Microsoft Excel
&
\sphinxAtStartPar

&
\sphinxAtStartPar

\\
\sphinxhline
\sphinxAtStartPar
Microsoft PowerPoint
&
\sphinxAtStartPar

&
\sphinxAtStartPar

\\
\sphinxhline
\sphinxAtStartPar
Windows Defender
&
\sphinxAtStartPar

&
\sphinxAtStartPar

\\
\sphinxhline
\sphinxAtStartPar
OneDrive
&
\sphinxAtStartPar

&
\sphinxAtStartPar

\\
\sphinxhline
\sphinxAtStartPar
Microsoft Outlook
&
\sphinxAtStartPar

&
\sphinxAtStartPar

\\
\sphinxhline
\sphinxAtStartPar
Microsoft Designer
&
\sphinxAtStartPar

&
\sphinxAtStartPar

\\
\sphinxhline
\sphinxAtStartPar
Clipchamp
&
\sphinxAtStartPar

&
\sphinxAtStartPar

\\
\sphinxhline
\sphinxAtStartPar
OneNote
&
\sphinxAtStartPar

&
\sphinxAtStartPar

\\
\sphinxbottomrule
\end{tabulary}
\sphinxtableafterendhook\par
\sphinxattableend\end{savenotes}

\item {} 
\sphinxAtStartPar
ANALICE los siguientes casos, eliga la mejor herramienta de LibreOffice para cada caso y explique el por qué:
\begin{quote}

\sphinxAtStartPar
\sphinxstylestrong{CASO 01}: Un investigador necesita crear fórmulas matemáticas complejas con símbolos avanzados (integrales, matrices) que deben ser interpretadas correctamente por navegadores web. ¿Qué herramienta específica debe usar para generar el código MathML necesario, y por qué es más eficiente usarla como herramienta independiente en lugar de solo dentro de un procesador de textos?
\end{quote}
\begin{quote}

\sphinxAtStartPar
\sphinxstylestrong{CASO 02}: Una empresa necesita crear diagramas de flujo de procesos industriales donde los nodos deben mantener su posición lógica incluso si se mueven otros elementos, utilizando “Conectores Inteligentes”. Además, requieren exportar estos diagramas en formato vectorial (como SVG o WebP) para una página web. ¿Qué herramienta de LibreOffice usaría y por qué?
\end{quote}
\begin{quote}

\sphinxAtStartPar
\sphinxstylestrong{CASO 03}: Se requiere crear una presentación que incluya no solo diapositivas, sino también efectos de Fontwork, animaciones personalizadas y gráficos técnicos importados directamente de otros componentes de dibujo vectorial. ¿Qué herramienta utilizaría para orquestar este ecosistema de elementos gráficos y multimedia? ¿Por qué?
\end{quote}
\begin{quote}

\sphinxAtStartPar
\sphinxstylestrong{CASO 04}: Un administrador de sistemas necesita una solución que permita gestionar relaciones entre tablas, realizar consultas SQL (específicamente soporte para ANSI\sphinxhyphen{}92) y conectarse a motores externos como PostgreSQL o MySQL. ¿Qué herramienta de la suite debe emplear y cómo se diferencia su interfaz de una simple hoja de cálculo para este propósito?
\end{quote}
\begin{quote}

\sphinxAtStartPar
\sphinxstylestrong{CASO 05}: Un usuario debe redactar un manual que requiere un estilo de formato coherente, una estructura de capítulos compleja, la inserción de gráficos vectoriales y la capacidad de exportar el resultado final a formatos de lectura digital como EPUB y PDF. ¿Qué herramienta de LibreOffice debería elegir y por qué su capacidad de integración con otros componentes es crítica en este caso?
\end{quote}
\begin{quote}

\sphinxAtStartPar
\sphinxstylestrong{CASO 06}: Un equipo de contabilidad necesita evaluar la viabilidad de un proyecto financiero. Deben procesar grandes volúmenes de datos usando funciones estadísticas y aplicar el “Gestor de escenarios” para analizar diferentes variables de tipo: (“qué pasaría si”). ¿Por qué esta herramienta debe ser superior a una simple hoja de cálculo básica?
\end{quote}

\item {} 
\sphinxAtStartPar
Para cada una de las ventajas de LibreOffice, en su cuaderno responda la siguiente pregunta:
\begin{itemize}
\item {} 
\sphinxAtStartPar
¿Cómo se beneficiaría la comunidad académica si se implementara LibreOffice en la Institución Educativa?

\end{itemize}

\end{enumerate}


\bigskip\hrule\bigskip



\subsection{A03: Carta al propietario (+4,0)}
\label{\detokenize{1_guias/10ofimatica/G1012:a03-carta-al-propietario-4-0}}
\sphinxAtStartPar
\DUrole{sd-sphinx-override}{\DUrole{sd-badge}{\DUrole{sd-bg-danger}{\DUrole{sd-bg-text-danger}{Transferencia}}}}


\begin{savenotes}\sphinxattablestart
\sphinxthistablewithglobalstyle
\centering
\begin{tabulary}{\linewidth}[t]{TT}
\sphinxtoprule
\sphinxtableatstartofbodyhook
\sphinxAtStartPar
\sphinxstylestrong{TIEMPO}
&
\sphinxAtStartPar
59 minutos
\\
\sphinxhline
\sphinxAtStartPar
\sphinxstylestrong{PREGUNTAS}
&
\sphinxAtStartPar
0
\\
\sphinxhline
\sphinxAtStartPar
\sphinxstylestrong{OBJETIVO}
&
\sphinxAtStartPar
Aplicar sus conocimientos en software libre para orientar al propietario de un negocio.
\\
\sphinxhline
\sphinxAtStartPar
\sphinxstylestrong{ACTIVIDAD}
&
\sphinxAtStartPar
Situación carta al propietario.
\\
\sphinxbottomrule
\end{tabulary}
\sphinxtableafterendhook\par
\sphinxattableend\end{savenotes}

\begin{sphinxadmonition}{note}{Descripción de la situación}
\begin{enumerate}
\sphinxsetlistlabels{\arabic}{enumi}{enumii}{}{.}%
\item {} 
\sphinxAtStartPar
Imagine la siguiente situación:

\end{enumerate}
\begin{quote}

\sphinxAtStartPar
“Usted trabaja en un negocio de \sphinxstylestrong{venta de frutas y verduras} con muy \sphinxstylestrong{poco presupuesto}. \sphinxstyleemphasis{Las cuentas de las ganancias no coinciden con los productos vendidos}. Su jefe no tiene el dinero suficiente para adquirir un \sphinxstylestrong{software contable}”. El negocio está ganando mala fama y cada día pierde clientes.
\end{quote}
\begin{enumerate}
\sphinxsetlistlabels{\arabic}{enumi}{enumii}{}{.}%
\setcounter{enumi}{1}
\item {} 
\sphinxAtStartPar
En una \sphinxstylestrong{hoja de cuaderno}, REDACTE una \sphinxstylestrong{carta} aconsejando a su jefe sobre el uso del \sphinxstylestrong{software libre} en el negocio. Tenga en cuenta lo siguiente:
\begin{itemize}
\item {} 
\sphinxAtStartPar
Tenga en cuenta que es un modelo de carta (no solo un escrito).

\item {} 
\sphinxAtStartPar
La carta debe tener:
\begin{itemize}
\item {} 
\sphinxAtStartPar
Buena presentación.

\item {} 
\sphinxAtStartPar
Letra entendible.

\item {} 
\sphinxAtStartPar
Redacción coherente.

\item {} 
\sphinxAtStartPar
Ortografía correcta.

\end{itemize}

\item {} 
\sphinxAtStartPar
Invente un nombre para su jefe.

\item {} 
\sphinxAtStartPar
Cree un nombre para el negocio.

\item {} 
\sphinxAtStartPar
Use una dirección ficticia en Puerto Boyacá. (Justificar la ubicación de la dirección).

\item {} 
\sphinxAtStartPar
Indique las características del software libre.

\item {} 
\sphinxAtStartPar
Indique las ventajas y desventajas de usar el software libre.

\item {} 
\sphinxAtStartPar
Adicione un análisis cuantitativo de ahorro.

\item {} 
\sphinxAtStartPar
Añadir un análisis de seguridad de la información (Malware, Virus, Hacking).

\item {} 
\sphinxAtStartPar
Adicione un análisis frente a la competencia.

\item {} 
\sphinxAtStartPar
Nombrar alguna fruta, verdura, producto o servicio que se ofrezca en el negocio.

\end{itemize}


\begin{savenotes}\sphinxattablestart
\sphinxthistablewithglobalstyle
\centering
\begin{tabulary}{\linewidth}[t]{T}
\sphinxtoprule
\sphinxtableatstartofbodyhook

\begin{sphinxfigure-in-table}
\centering
\capstart
\noindent\sphinxincludegraphics[width=400\sphinxpxdimen]{{image21}.png}
\sphinxfigcaption{\sphinxstylestrong{Situación negocio:} frutas y verduras}\label{\detokenize{1_guias/10ofimatica/G1012:id60}}\end{sphinxfigure-in-table}\relax
\\
\sphinxbottomrule
\end{tabulary}
\sphinxtableafterendhook\par
\sphinxattableend\end{savenotes}

\end{enumerate}
\end{sphinxadmonition}


\bigskip\hrule\bigskip


\sphinxstepscope


\section{Guía 1013: LibreOffice Writer (P2)}
\label{\detokenize{1_guias/10ofimatica/G1013:guia-1013-libreoffice-writer-p2}}\label{\detokenize{1_guias/10ofimatica/G1013::doc}}

\begin{savenotes}\sphinxattablestart
\sphinxthistablewithglobalstyle
\centering
\begin{tabular}[t]{\X{35}{100}\X{65}{100}}
\sphinxtoprule
\sphinxtableatstartofbodyhook

\begin{savenotes}\sphinxattablestart
\sphinxthistablewithglobalstyle
\centering
\begin{tabulary}{\linewidth}[t]{T}
\sphinxtoprule
\sphinxtableatstartofbodyhook
\sphinxAtStartPar
\sphinxincludegraphics{{escudo}.png}
\\
\sphinxbottomrule
\end{tabulary}
\sphinxtableafterendhook\par
\sphinxattableend\end{savenotes}
&

\begin{savenotes}\sphinxattablestart
\sphinxthistablewithglobalstyle
\raggedright
\begin{tabular}[t]{*{1}{\X{1}{1}}}
\sphinxtoprule
\sphinxtableatstartofbodyhook
\sphinxAtStartPar
\sphinxstylestrong{INSTITUCIÓN EDUCATIVA}: John F. Kennedy
\begin{itemize}
\item {} 
\sphinxAtStartPar
\sphinxstylestrong{DOCENTE}: Julian Camilo Fonseca Romero

\item {} 
\sphinxAtStartPar
\sphinxstylestrong{ASIGNATURA}: Ofimática

\item {} 
\sphinxAtStartPar
\sphinxstylestrong{GRADO}: DÉCIMO

\item {} 
\sphinxAtStartPar
\sphinxstylestrong{PERIODO:} 02

\item {} 
\sphinxAtStartPar
\sphinxstylestrong{TIEMPO}: 10 semanas (2 horas / semana)

\item {} 
\sphinxAtStartPar
\sphinxstylestrong{TIEMPO DE TRABAJO}: 20/20 horas

\end{itemize}
\\
\sphinxbottomrule
\end{tabular}
\sphinxtableafterendhook\par
\sphinxattableend\end{savenotes}
\\
\sphinxbottomrule
\end{tabular}
\sphinxtableafterendhook\par
\sphinxattableend\end{savenotes}

\begin{sphinxuseclass}{sd-tab-set}
\begin{sphinxuseclass}{sd-tab-item}\subsubsection*{Objetivo}

\begin{sphinxuseclass}{sd-tab-content}
\sphinxAtStartPar
\DUrole{sd-sphinx-override}{\DUrole{sd-badge}{\DUrole{sd-bg-light}{\DUrole{sd-bg-text-light}{Exploración}}}}
\DUrole{sd-sphinx-override}{\DUrole{sd-badge}{\DUrole{sd-bg-info}{\DUrole{sd-bg-text-info}{Estructuración}}}}
\DUrole{sd-sphinx-override}{\DUrole{sd-badge}{\DUrole{sd-bg-warning}{\DUrole{sd-bg-text-warning}{Actividad práctica}}}}
\DUrole{sd-sphinx-override}{\DUrole{sd-badge}{\DUrole{sd-bg-danger}{\DUrole{sd-bg-text-danger}{Transferencia}}}}
\begin{itemize}
\item {} 
\sphinxAtStartPar
\sphinxstylestrong{Temática}: Uso básico de LibreOffice Writer.

\item {} 
\sphinxAtStartPar
\sphinxstylestrong{Objetivo de aprendizaje}: Usar de forma básica el software LibreOffice Writer.

\end{itemize}

\end{sphinxuseclass}
\end{sphinxuseclass}
\begin{sphinxuseclass}{sd-tab-item}\subsubsection*{Orientaciones curriculares}

\begin{sphinxuseclass}{sd-tab-content}\begin{itemize}
\item {} 
\sphinxAtStartPar
\sphinxstylestrong{Componente}: Uso y apropiación de la Tecnología y la Informática.

\item {} 
\sphinxAtStartPar
\sphinxstylestrong{Competencia}: Genero propuestas innovadoras para el uso y aprovechamiento de productos tecnológicos.

\item {} 
\sphinxAtStartPar
\sphinxstylestrong{Evidencia de aprendizaje}: Utilizo adecuadamente herramientas informáticas para la búsqueda, organización, procesamiento, sistematización, comunicación y difusión de ideas.

\end{itemize}

\end{sphinxuseclass}
\end{sphinxuseclass}
\begin{sphinxuseclass}{sd-tab-item}\subsubsection*{Criterios de evaluación}

\begin{sphinxuseclass}{sd-tab-content}\begin{itemize}
\item {} 
\sphinxAtStartPar
Actividades desarrolladas en su totalidad.

\item {} 
\sphinxAtStartPar
Participación en clase.

\item {} 
\sphinxAtStartPar
Explicación clara de conceptos.

\item {} 
\sphinxAtStartPar
Argumentación de ideas.

\item {} 
\sphinxAtStartPar
Cumplimiento de las reglas del aula.

\end{itemize}

\end{sphinxuseclass}
\end{sphinxuseclass}
\begin{sphinxuseclass}{sd-tab-item}\subsubsection*{Recursos (3)}

\begin{sphinxuseclass}{sd-tab-content}\begin{itemize}
\item {} 
\sphinxAtStartPar
\sphinxhref{https://www.libreoffice.org/download/}{Descarga de \sphinxstylestrong{LibreOffice}}.

\item {} 
\sphinxAtStartPar
G1013A07\sphinxhyphen{}mitos\_desordenados.odt \DUrole{xref}{\DUrole{download}{\DUrole{myst}{\DUrole{sd-sphinx-override}{\DUrole{sd-badge}{\DUrole{sd-bg-primary}{\DUrole{sd-bg-text-primary}{Descargar}}}}}}}

\item {} 
\sphinxAtStartPar
G1013A08\sphinxhyphen{}proyecto\_escritorio\_libre.odt \DUrole{xref}{\DUrole{download}{\DUrole{myst}{\DUrole{sd-sphinx-override}{\DUrole{sd-badge}{\DUrole{sd-bg-primary}{\DUrole{sd-bg-text-primary}{Descargar}}}}}}}

\end{itemize}

\end{sphinxuseclass}
\end{sphinxuseclass}
\end{sphinxuseclass}

\subsection{Mensaje del autor}
\label{\detokenize{1_guias/10ofimatica/G1013:mensaje-del-autor}}\begin{quote}

\sphinxAtStartPar
“El propósito de la presente guía es que el estudiante descubra que un procesador de textos es mucho más que una herramienta para escribir; es un aliado fundamental para organizar sus ideas, estructurar informes y elevar la calidad visual de sus proyectos académicos y personales. El objetivo de esta guía es que el estudiante desarrolle las competencias técnicas necesarias para manejar herramientas de edición de texto con fluidez. No solo se aprende a ‘hacer que el texto se vea bien’, sino a dominar la lógica detrás de la documentación profesional, una habilidad indispensable en el mundo laboral actual. El desarrollo de cada actividad fortalece la autonomía al hacer uso de un computador”.

\begin{flushright}
---Camilo Fonseca
\end{flushright}
\end{quote}


\bigskip\hrule\bigskip



\subsection{A00: Introducción (+0,1)}
\label{\detokenize{1_guias/10ofimatica/G1013:a00-introduccion-0-1}}
\sphinxAtStartPar
\DUrole{sd-sphinx-override}{\DUrole{sd-badge}{\DUrole{sd-bg-light}{\DUrole{sd-bg-text-light}{Exploración}}}}
\begin{enumerate}
\sphinxsetlistlabels{\arabic}{enumi}{enumii}{}{.}%
\item {} 
\sphinxAtStartPar
Transcriba en su cuaderno el objetivo, las orientaciones curriculares y los criterios de evaluación de la presente guía.

\item {} 
\sphinxAtStartPar
Lea el mensaje del autor y responda la siguiente pregunta en su cuaderno:
\begin{itemize}
\item {} 
\sphinxAtStartPar
¿De qué forma el aprender a usar una herramienta de edición de texto como LibreOffice Writer, permite desarrollar la autonomía en el estudiante?

\end{itemize}

\end{enumerate}


\bigskip\hrule\bigskip



\subsection{A01: LibreOffice Writer (+0,1)}
\label{\detokenize{1_guias/10ofimatica/G1013:a01-libreoffice-writer-0-1}}
\sphinxAtStartPar
\DUrole{sd-sphinx-override}{\DUrole{sd-badge}{\DUrole{sd-bg-info}{\DUrole{sd-bg-text-info}{Estructuración}}}}
\begin{enumerate}
\sphinxsetlistlabels{\arabic}{enumi}{enumii}{}{.}%
\item {} 
\sphinxAtStartPar
Lea los siguientes conceptos:

\end{enumerate}


\bigskip\hrule\bigskip



\subsubsection{¿Qué es LibreOffice Writer?}
\label{\detokenize{1_guias/10ofimatica/G1013:que-es-libreoffice-writer}}

\begin{savenotes}\sphinxattablestart
\sphinxthistablewithglobalstyle
\centering
\begin{tabulary}{\linewidth}[t]{T}
\sphinxtoprule
\sphinxtableatstartofbodyhook

\begin{sphinxfigure-in-table}
\centering
\capstart
\noindent\sphinxincludegraphics[width=120\sphinxpxdimen]{{image017}.png}
\sphinxfigcaption{LibreOffice Writer}\label{\detokenize{1_guias/10ofimatica/G1013:id1}}\end{sphinxfigure-in-table}\relax
\\
\sphinxbottomrule
\end{tabulary}
\sphinxtableafterendhook\par
\sphinxattableend\end{savenotes}

\sphinxAtStartPar
\sphinxstylestrong{LibreOffice Writer} es un procesador de textos, compatible con todas las versiones de Microsoft Word y además nos provee de nuevas funcionalidades. Writer le permitirá crear documentos con características como:
\begin{itemize}
\item {} 
\sphinxAtStartPar
Un correcto formateado de caracteres y párrafos.

\item {} 
\sphinxAtStartPar
Un uso eficiente de los estilos.

\item {} 
\sphinxAtStartPar
Auto\sphinxhyphen{}texto que le permite reutilizar contenidos con toda comodidad.

\item {} 
\sphinxAtStartPar
Numeraciones y viñetas.

\item {} 
\sphinxAtStartPar
Inserción y manejo de imágenes.

\item {} 
\sphinxAtStartPar
Trabajo con documentos largos esquematizados, con creación de índices.

\end{itemize}

\sphinxAtStartPar
A continuación se muestra un ejemplo de un documento realizado en LibreOffice Writer:


\begin{savenotes}\sphinxattablestart
\sphinxthistablewithglobalstyle
\centering
\begin{tabulary}{\linewidth}[t]{T}
\sphinxtoprule
\sphinxtableatstartofbodyhook

\begin{sphinxfigure-in-table}
\centering
\capstart
\noindent\sphinxincludegraphics[width=600\sphinxpxdimen]{{image025}.png}
\sphinxfigcaption{Hoja de documento LibreOffice Writer}\label{\detokenize{1_guias/10ofimatica/G1013:id2}}\end{sphinxfigure-in-table}\relax
\\
\sphinxbottomrule
\end{tabulary}
\sphinxtableafterendhook\par
\sphinxattableend\end{savenotes}

\sphinxAtStartPar
LibreOffice también puede usarse de forma profesional. Para ello es necesario conocer varias herramientas que permiten hacer que los documentos sean más productivos. Estas herramientas son:
\begin{itemize}
\item {} 
\sphinxAtStartPar
Uso de tablas para crear formularios con valores que se auto\sphinxhyphen{}calculan.

\item {} 
\sphinxAtStartPar
Creación de secciones en columnas en el documento.

\item {} 
\sphinxAtStartPar
Creación de diagramas de datos a partir de datos contenidos en tablas.

\item {} 
\sphinxAtStartPar
Trabajo con cuadros de texto y marcos.

\item {} 
\sphinxAtStartPar
Establecer referencias, vínculos y enlaces en los documentos.

\item {} 
\sphinxAtStartPar
Creación de cartas en serie y etiquetas personalizadas para múltiples destinatarios extraídos de una base de datos.

\item {} 
\sphinxAtStartPar
Asignación de contraseñas para asegurar la integridad de los documentos.

\item {} 
\sphinxAtStartPar
Personalización de las barra de herramientas, menús y atajos de teclado para mayor productividad.

\end{itemize}


\begin{savenotes}\sphinxattablestart
\sphinxthistablewithglobalstyle
\centering
\begin{tabulary}{\linewidth}[t]{T}
\sphinxtoprule
\sphinxtableatstartofbodyhook

\begin{sphinxfigure-in-table}
\centering
\capstart
\noindent\sphinxincludegraphics[width=600\sphinxpxdimen]{{image034}.png}
\sphinxfigcaption{Herramientas LibreOffice Writer}\label{\detokenize{1_guias/10ofimatica/G1013:id3}}\end{sphinxfigure-in-table}\relax
\\
\sphinxbottomrule
\end{tabulary}
\sphinxtableafterendhook\par
\sphinxattableend\end{savenotes}


\bigskip\hrule\bigskip



\subsubsection{LibreOffice Writer Vs Microsoft Word}
\label{\detokenize{1_guias/10ofimatica/G1013:libreoffice-writer-vs-microsoft-word}}

\begin{savenotes}\sphinxattablestart
\sphinxthistablewithglobalstyle
\centering
\begin{tabulary}{\linewidth}[t]{T}
\sphinxtoprule
\sphinxtableatstartofbodyhook

\begin{sphinxfigure-in-table}
\centering
\capstart
\noindent\sphinxincludegraphics[width=200\sphinxpxdimen]{{image044}.png}
\sphinxfigcaption{LibreOffice Writer Vs Microsoft Word}\label{\detokenize{1_guias/10ofimatica/G1013:id4}}\end{sphinxfigure-in-table}\relax
\\
\sphinxbottomrule
\end{tabulary}
\sphinxtableafterendhook\par
\sphinxattableend\end{savenotes}

\sphinxAtStartPar
LibreOffice Writer y Microsoft Word son dos suites de oficina ampliamente utilizadas que ofrecen muchas funcionalidades similares para la creación y edición de documentos. Sin embargo, hay algunas características distintivas de LibreOffice Writer que pueden diferenciarlo de Microsoft Word.

\sphinxAtStartPar
A continuación se muestran algunas funcionalidades de LibreOffice Writer que no están disponibles en Microsoft Word:
\begin{itemize}
\item {} 
\sphinxAtStartPar
\sphinxstylestrong{Licencia de código abierto:} LibreOffice es un software de código abierto bajo la licencia LGPL, lo que significa que es gratuito para descargar, usar, modificar y distribuir libremente. Microsoft Word, en cambio, es parte de Microsoft Office, que es un software comercial con licencia propietaria.

\item {} 
\sphinxAtStartPar
\sphinxstylestrong{Soporte nativo para formatos abiertos:} LibreOffice Writer tiene un soporte nativo robusto para formatos abiertos como ODT (OpenDocument Text), mientras que Microsoft Word generalmente utiliza formatos propietarios como DOCX. Aunque Word puede abrir y guardar archivos ODT, no es su formato nativo.

\item {} 
\sphinxAtStartPar
\sphinxstylestrong{Mayor flexibilidad en la interfaz:} LibreOffice Writer ofrece una interfaz altamente personalizable que puede ser ajustada para adaptarse mejor a las preferencias del usuario. Esto puede incluir la capacidad de personalizar barras de herramientas, atajos de teclado y más, de una manera más flexible que en Microsoft Word.

\item {} 
\sphinxAtStartPar
\sphinxstylestrong{Compatibilidad multi\sphinxhyphen{}plataforma:} LibreOffice está disponible para una amplia gama de sistemas operativos, incluidos Windows, macOS y varias distribuciones de Linux. Microsoft Word, aunque está disponible para Windows y macOS, no tiene una versión nativa para Linux.

\item {} 
\sphinxAtStartPar
\sphinxstylestrong{Extensiones y complementos libres:} LibreOffice tiene una comunidad activa que desarrolla extensiones y complementos que pueden ampliar sus funcionalidades básicas. Estos pueden ser descargados e instalados fácilmente para añadir nuevas características a LibreOffice Writer.

\item {} 
\sphinxAtStartPar
\sphinxstylestrong{No requiere suscripción:} LibreOffice Writer y toda la suite LibreOffice no requieren una suscripción para su uso. Esto puede ser una ventaja para aquellos que prefieren no comprometerse con un modelo de suscripción.

\end{itemize}


\bigskip\hrule\bigskip



\subsubsection{Características de LibreOffice Writer}
\label{\detokenize{1_guias/10ofimatica/G1013:caracteristicas-de-libreoffice-writer}}
\begin{sphinxadmonition}{note}{Paginación e impresión de documentos}

\sphinxAtStartPar
Cuando se desea distribuir documentos de forma impresa o en formato PDF, tan sólo necesitamos disponer de un documento bien estructurado y organizado. Para ello tendremos en cuenta otras configuraciones como el tamaño y formato del papel donde se imprimirá, los márgenes del documento, la orientación, los encabezados y pies que se deben repetir en todas las páginas o la numeración de las páginas.
\end{sphinxadmonition}

\begin{sphinxadmonition}{note}{Imágenes y objetos de dibujo}

\sphinxAtStartPar
Según las necesidades que se tengan, será preciso insertar imágenes en los documentos para ilustrar algún informe o proyecto.

\sphinxAtStartPar
También podemos tener necesidad de dibujar formas geométricas para diagramar algún proceso, o por simple placer creativo. Resolver estas tareas requerirá el conocimiento de las diferentes propiedades que se pueden establecer para obtener los mejores resultados.
\end{sphinxadmonition}

\begin{sphinxadmonition}{note}{Herramientas de revisión y autocorrección}

\sphinxAtStartPar
Ocasionalmente se tienen dudas sobre el uso ortográfico correcto en la redacción de nuestros escritos. O quizás estamos escribiendo en un idioma con el cual no estamos demasiado familiarizados. En esos momentos, el uso de correctores ortográficos es de gran ayuda.

\sphinxAtStartPar
También pueden resultar muy útiles otro tipo de ayudas de escritura como los diccionarios de sinónimos, de corrección gramatical, de separación silábica, así como opciones de formateado y corrección automática.

\sphinxAtStartPar
En ocasiones puede que un documento maneje varios idiomas. En estos casos lo mejor es establecer estilos con el idioma extra y aplicarlos en el documento según se necesite.
\end{sphinxadmonition}

\begin{sphinxadmonition}{note}{Trabajo con documentos largos}

\sphinxAtStartPar
Cuando se elaboran largos informes o memorias, nuestros documentos deben dotarse de una estructura que haga cómodo su seguimiento. Ya sabemos que se puede y se “debe” utilizar estilos de Título para ello.

\sphinxAtStartPar
Pero además, Writer nos proporciona un conjunto de herramientas que nos facilitarán esta tarea. Gracias a estas herramientas podremos alterar con mucha facilidad el esquema de nuestros documentos, los títulos podrán ser numerados y crearemos sumarios y tablas de toda clase de contenido.
\end{sphinxadmonition}

\begin{sphinxadmonition}{note}{Navegador de elementos}

\sphinxAtStartPar
El navegador contiene una lista donde aparecen los elementos que conforman un documento. Estos elementos son:


\begin{savenotes}\sphinxattablestart
\sphinxthistablewithglobalstyle
\centering
\begin{tabulary}{\linewidth}[t]{TT}
\sphinxtoprule
\sphinxstyletheadfamily 
\sphinxAtStartPar
ELEMENTOS
&\sphinxstyletheadfamily 
\sphinxAtStartPar
DESCRIPCIÓN
\\
\sphinxmidrule
\sphinxtableatstartofbodyhook
\sphinxAtStartPar
\sphinxstylestrong{Títulos}
&
\sphinxAtStartPar
Los diferentes títulos del documento, agrupados y sangrados según su nivel, dicho de otra manera, el esquema del documento.
\\
\sphinxhline
\sphinxAtStartPar
\sphinxstylestrong{Tablas}
&
\sphinxAtStartPar
Las tablas son estructuras de texto organizadas en filas y columnas.
\\
\sphinxhline
\sphinxAtStartPar
\sphinxstylestrong{Marcos de texto}
&
\sphinxAtStartPar
Cuadros con texto en su interior que pueden ser tratados como un objeto gráfico.
\\
\sphinxhline
\sphinxAtStartPar
\sphinxstylestrong{Imágenes}
&
\sphinxAtStartPar
Imágenes insertadas en el documento.
\\
\sphinxhline
\sphinxAtStartPar
\sphinxstylestrong{Objetos OLE}
&
\sphinxAtStartPar
Elementos como hojas de cálculo, fórmulas, y, en general, objetos de cualquier programa que admita la tecnología OLE (Object Linking and Embedding \sphinxhyphen{} Objetos vinculados e incrustados)
\\
\sphinxhline
\sphinxAtStartPar
\sphinxstylestrong{Marcadores}
&
\sphinxAtStartPar
Puntos del texto con un nombre que se crean para ser referenciados desde otro lugar del documento.
\\
\sphinxhline
\sphinxAtStartPar
\sphinxstylestrong{Secciones}
&
\sphinxAtStartPar
Divisiones estructurales del documento que permiten aplicar determinados atributos a partes del documento, como la organización en columnas periodísticas o protección, entre otros.
\\
\sphinxhline
\sphinxAtStartPar
\sphinxstylestrong{Hiperenlaces}
&
\sphinxAtStartPar
Vínculos a sitios web, otros documentos, títulos o marcadores.
\\
\sphinxhline
\sphinxAtStartPar
\sphinxstylestrong{Referencias}
&
\sphinxAtStartPar
Textos del documento que apuntan a algún Marcador.
\\
\sphinxhline
\sphinxAtStartPar
\sphinxstylestrong{Índices}
&
\sphinxAtStartPar
Tablas de contenido que pueden ser de tipos muy variados: sumarios, alfabéticos o de ilustraciones entre otros.
\\
\sphinxhline
\sphinxAtStartPar
\sphinxstylestrong{Comentarios}
&
\sphinxAtStartPar
Anotaciones creadas para revisiones del documento o establecer recordatorios.
\\
\sphinxhline
\sphinxAtStartPar
\sphinxstylestrong{Objetivos}
&
\sphinxAtStartPar
Elementos de dibujo insertados en el documento.
\\
\sphinxbottomrule
\end{tabulary}
\sphinxtableafterendhook\par
\sphinxattableend\end{savenotes}
\end{sphinxadmonition}

\begin{sphinxadmonition}{note}{Uso de tablas}

\sphinxAtStartPar
LibreOffice Writer dispone de una herramienta para mostrar información organizada en los  documentos: las tablas. Su uso sustituye de manera ventajosa al uso de los tabuladores. Pero las tablas van mucho más allá, puesto que además de texto, pueden incluir valores numéricos con los que efectuar cálculos con toda comodidad.
\end{sphinxadmonition}


\bigskip\hrule\bigskip



\subsection{A02: Taller (+0,8)}
\label{\detokenize{1_guias/10ofimatica/G1013:a02-taller-0-8}}
\sphinxAtStartPar
\DUrole{sd-sphinx-override}{\DUrole{sd-badge}{\DUrole{sd-bg-warning}{\DUrole{sd-bg-text-warning}{Actividad práctica}}}}
\begin{enumerate}
\sphinxsetlistlabels{\arabic}{enumi}{enumii}{}{.}%
\item {} 
\sphinxAtStartPar
En su cuaderno responda las siguientes preguntas:
\begin{itemize}
\item {} 
\sphinxAtStartPar
¿Para que sirve LibreOffice Writer?

\item {} 
\sphinxAtStartPar
¿Qué características tiene el software LibreOffice Writer?

\item {} 
\sphinxAtStartPar
¿Cuáles son las 6 ventajas de usar \sphinxstylestrong{LibreOffice Writer} frente \sphinxstylestrong{Microsoft Word}?

\item {} 
\sphinxAtStartPar
¿Qué posibles motivos tendría LibreOffice para no ofrecer un servicio de edición de documentos en la nube?

\end{itemize}

\item {} 
\sphinxAtStartPar
Observe las siguientes imágenes y en \sphinxstylestrong{una hoja de cuaderno} realice una descripción de las cosas que se pueden hacer en LibreOffice Writer:


\begin{savenotes}\sphinxattablestart
\sphinxthistablewithglobalstyle
\centering
\begin{tabulary}{\linewidth}[t]{TT}
\sphinxtoprule
\sphinxtableatstartofbodyhook

\begin{sphinxfigure-in-table}
\centering

\noindent\sphinxincludegraphics[width=360\sphinxpxdimen]{{image025}.png}
\end{sphinxfigure-in-table}\relax
&

\begin{sphinxfigure-in-table}
\centering

\noindent\sphinxincludegraphics[width=360\sphinxpxdimen]{{image034}.png}
\end{sphinxfigure-in-table}\relax
\\
\sphinxbottomrule
\end{tabulary}
\sphinxtableafterendhook\par
\sphinxattableend\end{savenotes}

\item {} 
\sphinxAtStartPar
En una hoja de cuaderno, redacte un escrito que responda a las siguientes preguntas:
\begin{itemize}
\item {} 
\sphinxAtStartPar
¿Por qué Microsoft Office para PC solo está disponible (de forma legal) para ser usado exclusivamente en Windows y MACOS? (media hoja de cuaderno).

\item {} 
\sphinxAtStartPar
¿Qué posibles motivos tendría Microsoft para que Microsoft Word no sea de código abierto? (media hoja de cuaderno).

\end{itemize}

\end{enumerate}


\bigskip\hrule\bigskip



\subsection{A03: Información de identificación \sphinxstyleliteralintitle{\sphinxupquote{011}} (+0,1)}
\label{\detokenize{1_guias/10ofimatica/G1013:a03-informacion-de-identificacion-011-0-1}}
\sphinxAtStartPar
\DUrole{sd-sphinx-override}{\DUrole{sd-badge}{\DUrole{sd-bg-warning}{\DUrole{sd-bg-text-warning}{Actividad práctica}}}}

\sphinxAtStartPar
En esta actividad, el estudiante aprende a crear la \sphinxstylestrong{información de identificación}.

\begin{sphinxadmonition}{note}{¿Qué es la \sphinxstylestrong{información de identificación}?}

\sphinxAtStartPar
Todo documento que se cree en esta guía, debe tener la \sphinxstylestrong{información de identificación} en la cabecera del documento con la siguiente información y en el siguiente orden:

\begin{center}TÍTULO DEL DOCUMENTO
\end{center}
\begin{center}Nombres y apellidos del estudiante
\end{center}
\begin{center}Curso \sphinxhyphen{} Jornada
\end{center}
\begin{center}HORA \sphinxhyphen{} FECHA
\end{center}\end{sphinxadmonition}

\sphinxAtStartPar
Para esto siga los siguientes pasos:
\begin{enumerate}
\sphinxsetlistlabels{\arabic}{enumi}{enumii}{}{.}%
\item {} 
\sphinxAtStartPar
Ejecute el software LibreOffice Writer.

\item {} 
\sphinxAtStartPar
Realice doble clic en la \sphinxstyleemphasis{cabecera del documento} hasta que salga una opción que diga: \sphinxstylestrong{Cabecera (Estilo de página predeterminado)}.


\begin{savenotes}\sphinxattablestart
\sphinxthistablewithglobalstyle
\centering
\begin{tabulary}{\linewidth}[t]{T}
\sphinxtoprule
\sphinxtableatstartofbodyhook

\begin{sphinxfigure-in-table}
\centering

\noindent\sphinxincludegraphics[width=300\sphinxpxdimen]{{image054}.png}
\end{sphinxfigure-in-table}\relax
\\
\sphinxbottomrule
\end{tabulary}
\sphinxtableafterendhook\par
\sphinxattableend\end{savenotes}

\item {} 
\sphinxAtStartPar
En la cabecera del documento, agregue la \sphinxstylestrong{información de identificación} con el título: \sphinxstylestrong{IDENTIFICACIÓN}. Debería quedar así:


\begin{savenotes}\sphinxattablestart
\sphinxthistablewithglobalstyle
\centering
\begin{tabulary}{\linewidth}[t]{T}
\sphinxtoprule
\sphinxtableatstartofbodyhook

\begin{sphinxfigure-in-table}
\centering

\noindent\sphinxincludegraphics[width=400\sphinxpxdimen]{{image064}.png}
\end{sphinxfigure-in-table}\relax
\\
\sphinxbottomrule
\end{tabulary}
\sphinxtableafterendhook\par
\sphinxattableend\end{savenotes}

\item {} 
\sphinxAtStartPar
Realice una captura de pantalla del documento y guarde la imagen en la \sphinxstylestrong{carpeta de evidencias} del periodo actual con el nombre: \sphinxstylestrong{G1013A03\sphinxhyphen{}Identificacion.png}

\item {} 
\sphinxAtStartPar
Cierre el software LibreOffice Writer.

\item {} 
\sphinxAtStartPar
En su cuaderno:

\begin{sphinxadmonition}{note}{En su cuaderno responda la siguiente pregunta:}
\begin{itemize}
\item {} 
\sphinxAtStartPar
¿Por qué es importante identificar los documentos que creamos?

\end{itemize}
\end{sphinxadmonition}

\end{enumerate}


\bigskip\hrule\bigskip



\subsection{A04: Exploración de la interfaz \sphinxstyleliteralintitle{\sphinxupquote{114}} (+0,2)}
\label{\detokenize{1_guias/10ofimatica/G1013:a04-exploracion-de-la-interfaz-114-0-2}}
\sphinxAtStartPar
\DUrole{sd-sphinx-override}{\DUrole{sd-badge}{\DUrole{sd-bg-warning}{\DUrole{sd-bg-text-warning}{Actividad práctica}}}}
\begin{enumerate}
\sphinxsetlistlabels{\arabic}{enumi}{enumii}{}{.}%
\item {} 
\sphinxAtStartPar
Ejecute el software LibreOffice Writer.

\item {} 
\sphinxAtStartPar
En la cabecera del documento, agregue la \sphinxstylestrong{información de identificación} con el título: \sphinxstylestrong{EXPLORACIÓN DE LA INTERFAZ}. Debería quedar así:


\begin{savenotes}\sphinxattablestart
\sphinxthistablewithglobalstyle
\centering
\begin{tabulary}{\linewidth}[t]{T}
\sphinxtoprule
\sphinxtableatstartofbodyhook

\begin{sphinxfigure-in-table}
\centering

\noindent\sphinxincludegraphics[width=400\sphinxpxdimen]{{image074}.png}
\end{sphinxfigure-in-table}\relax
\\
\sphinxbottomrule
\end{tabulary}
\sphinxtableafterendhook\par
\sphinxattableend\end{savenotes}

\item {} 
\sphinxAtStartPar
Identifique el icono de guardar:


\begin{savenotes}\sphinxattablestart
\sphinxthistablewithglobalstyle
\centering
\begin{tabulary}{\linewidth}[t]{T}
\sphinxtoprule
\sphinxtableatstartofbodyhook

\begin{sphinxfigure-in-table}
\centering

\noindent\sphinxincludegraphics[width=89\sphinxpxdimen]{{image084}.png}
\end{sphinxfigure-in-table}\relax
\\
\sphinxbottomrule
\end{tabulary}
\sphinxtableafterendhook\par
\sphinxattableend\end{savenotes}

\item {} 
\sphinxAtStartPar
Si el \sphinxstylestrong{icono de guardar} tiene un punto rojo, significa que el documento no ha sido guardado. PULSE el \sphinxstylestrong{icono de guardar} y se abrirá una ventana para guardar el documento.

\item {} 
\sphinxAtStartPar
En la ventana de guardado, LOCALICE su carpeta de evidencias del periodo actual y guarde el documento con el nombre: \sphinxstylestrong{G1013A04\sphinxhyphen{}documento.odt}

\item {} 
\sphinxAtStartPar
En la interfaz, IDENTIFIQUE el \sphinxstylestrong{menú superior}.


\begin{savenotes}\sphinxattablestart
\sphinxthistablewithglobalstyle
\centering
\begin{tabulary}{\linewidth}[t]{T}
\sphinxtoprule
\sphinxtableatstartofbodyhook

\begin{sphinxfigure-in-table}
\centering
\capstart
\noindent\sphinxincludegraphics[width=650\sphinxpxdimen]{{image094}.png}
\sphinxfigcaption{Menú superior}\label{\detokenize{1_guias/10ofimatica/G1013:id5}}\end{sphinxfigure-in-table}\relax
\\
\sphinxbottomrule
\end{tabulary}
\sphinxtableafterendhook\par
\sphinxattableend\end{savenotes}

\item {} 
\sphinxAtStartPar
En su cuaderno:

\begin{sphinxadmonition}{note}{En su cuaderno…}
\begin{itemize}
\item {} 
\sphinxAtStartPar
Dibuje el \sphinxstylestrong{menú superior}.

\end{itemize}
\end{sphinxadmonition}

\item {} 
\sphinxAtStartPar
En la interfaz, IDENTIFIQUE la \sphinxstylestrong{barra de herramientas estándar}:


\begin{savenotes}\sphinxattablestart
\sphinxthistablewithglobalstyle
\centering
\begin{tabulary}{\linewidth}[t]{T}
\sphinxtoprule
\sphinxtableatstartofbodyhook

\begin{sphinxfigure-in-table}
\centering
\capstart
\noindent\sphinxincludegraphics[width=700\sphinxpxdimen]{{image103}.png}
\sphinxfigcaption{Barra de herramientas estándar}\label{\detokenize{1_guias/10ofimatica/G1013:id6}}\end{sphinxfigure-in-table}\relax
\\
\sphinxbottomrule
\end{tabulary}
\sphinxtableafterendhook\par
\sphinxattableend\end{savenotes}

\item {} 
\sphinxAtStartPar
SITUAR el mouse en cada uno de los iconos de la \sphinxstylestrong{barra de herramientas estándar}. Al hacerlo el nombre de la herramienta aparecerá resaltado.


\begin{savenotes}\sphinxattablestart
\sphinxthistablewithglobalstyle
\centering
\begin{tabulary}{\linewidth}[t]{T}
\sphinxtoprule
\sphinxtableatstartofbodyhook

\begin{sphinxfigure-in-table}
\centering

\noindent\sphinxincludegraphics[width=100\sphinxpxdimen]{{image112}.png}
\end{sphinxfigure-in-table}\relax
\\
\sphinxbottomrule
\end{tabulary}
\sphinxtableafterendhook\par
\sphinxattableend\end{savenotes}

\item {} 
\sphinxAtStartPar
En su cuaderno:

\begin{sphinxadmonition}{note}{En su cuaderno…}
\begin{itemize}
\item {} 
\sphinxAtStartPar
Dibuje la \sphinxstylestrong{barra de herramientas estándar} y agregue el nombre de cada uno de sus iconos.

\end{itemize}
\end{sphinxadmonition}

\item {} 
\sphinxAtStartPar
En la interfaz, IDENTIFIQUE la barra de herramientas formato:


\begin{savenotes}\sphinxattablestart
\sphinxthistablewithglobalstyle
\centering
\begin{tabulary}{\linewidth}[t]{T}
\sphinxtoprule
\sphinxtableatstartofbodyhook

\begin{sphinxfigure-in-table}
\centering
\capstart
\noindent\sphinxincludegraphics[width=700\sphinxpxdimen]{{image122}.png}
\sphinxfigcaption{Barra de herramientas formato}\label{\detokenize{1_guias/10ofimatica/G1013:id7}}\end{sphinxfigure-in-table}\relax
\\
\sphinxbottomrule
\end{tabulary}
\sphinxtableafterendhook\par
\sphinxattableend\end{savenotes}

\item {} 
\sphinxAtStartPar
SITUAR el mouse en cada uno de los iconos de la \sphinxstylestrong{barra de herramientas formato}. Al hacerlo el nombre de la herramienta aparecerá resaltado.

\item {} 
\sphinxAtStartPar
En su cuaderno:

\begin{sphinxadmonition}{note}{En su cuaderno…}
\begin{itemize}
\item {} 
\sphinxAtStartPar
Dibuje la \sphinxstylestrong{barra de herramientas formato} y agregue el nombre de cada uno de sus iconos.

\end{itemize}
\end{sphinxadmonition}

\item {} 
\sphinxAtStartPar
En el \sphinxstylestrong{menú superior}, PULSE: \sphinxguilabel{Ver} > \sphinxguilabel{Barras de herramientas}. Se despliega una ventana con varias barras de herramientas desactivadas. ACTIVE las siguientes herramientas:
\begin{itemize}
\item {} 
\sphinxAtStartPar
Buscar

\item {} 
\sphinxAtStartPar
Controles de formulario

\item {} 
\sphinxAtStartPar
Dibujo

\item {} 
\sphinxAtStartPar
Tabla

\end{itemize}

\item {} 
\sphinxAtStartPar
REALICE una captura de pantalla de la interfaz de LibreOffice Writer y guarde la imagen en la carpeta de evidencias del periodo actual con el nombre: \sphinxstylestrong{G1013A04\sphinxhyphen{}interfaz.png}

\item {} 
\sphinxAtStartPar
En el \sphinxstylestrong{menú superior}, PULSE: \sphinxguilabel{Ver} > \sphinxguilabel{Barras de herramientas} y \sphinxstylestrong{DESACTIVE todas las herramientas.} SOLO DEJAR ACTIVAS las siguientes herramientas:
\begin{itemize}
\item {} 
\sphinxAtStartPar
Estándar

\item {} 
\sphinxAtStartPar
Formato

\end{itemize}

\item {} 
\sphinxAtStartPar
En la interfaz, IDENTIFIQUE la \sphinxstylestrong{barra de herramientas lateral}:


\begin{savenotes}\sphinxattablestart
\sphinxthistablewithglobalstyle
\centering
\begin{tabulary}{\linewidth}[t]{T}
\sphinxtoprule
\sphinxtableatstartofbodyhook

\begin{sphinxfigure-in-table}
\centering
\capstart
\noindent\sphinxincludegraphics[width=45\sphinxpxdimen]{{image132}.png}
\sphinxfigcaption{Barra de herramientas lateral}\label{\detokenize{1_guias/10ofimatica/G1013:id8}}\end{sphinxfigure-in-table}\relax
\\
\sphinxbottomrule
\end{tabulary}
\sphinxtableafterendhook\par
\sphinxattableend\end{savenotes}

\item {} 
\sphinxAtStartPar
PULSE cada uno de los iconos de la \sphinxstylestrong{barra de herramientas lateral}.


\begin{savenotes}\sphinxattablestart
\sphinxthistablewithglobalstyle
\centering
\begin{tabulary}{\linewidth}[t]{T}
\sphinxtoprule
\sphinxtableatstartofbodyhook

\begin{sphinxfigure-in-table}
\centering
\capstart
\noindent\sphinxincludegraphics[width=400\sphinxpxdimen]{{image142}.png}
\sphinxfigcaption{Propiedades barra de herramientas lateral}\label{\detokenize{1_guias/10ofimatica/G1013:id9}}\end{sphinxfigure-in-table}\relax
\\
\sphinxbottomrule
\end{tabulary}
\sphinxtableafterendhook\par
\sphinxattableend\end{savenotes}

\item {} 
\sphinxAtStartPar
En su cuaderno:

\begin{sphinxadmonition}{note}{En su cuaderno…}

\sphinxAtStartPar
Dibuje y escriba el (título y los subtítulos) de cada una de las secciones de cada icono, de la \sphinxstylestrong{barra de herramientas lateral}.
\end{sphinxadmonition}

\item {} 
\sphinxAtStartPar
Guarde los cambios y cierre el software LibreOffice Writer.

\end{enumerate}


\bigskip\hrule\bigskip



\subsection{A05: Creación y edición de documentos \sphinxstyleliteralintitle{\sphinxupquote{111}} (+0,2)}
\label{\detokenize{1_guias/10ofimatica/G1013:a05-creacion-y-edicion-de-documentos-111-0-2}}
\sphinxAtStartPar
\DUrole{sd-sphinx-override}{\DUrole{sd-badge}{\DUrole{sd-bg-warning}{\DUrole{sd-bg-text-warning}{Actividad práctica}}}}
\begin{enumerate}
\sphinxsetlistlabels{\arabic}{enumi}{enumii}{}{.}%
\item {} 
\sphinxAtStartPar
Ejecute el software LibreOffice Writer.

\item {} 
\sphinxAtStartPar
En la cabecera del documento, agregue la \sphinxstylestrong{información de identificación} con el título: \sphinxstylestrong{CREACIÓN DE DOCUMENTOS}

\item {} 
\sphinxAtStartPar
GUARDE el documento en su \sphinxstylestrong{carpeta de evidencias} del periodo actual con el nombre: \sphinxstylestrong{G1013A05\sphinxhyphen{}creacion\_documento.odt}

\item {} 
\sphinxAtStartPar
En la \sphinxstylestrong{barra de herramientas estándar}, PULSE el botón \sphinxguilabel{Nuevo}.


\begin{savenotes}\sphinxattablestart
\sphinxthistablewithglobalstyle
\centering
\begin{tabulary}{\linewidth}[t]{T}
\sphinxtoprule
\sphinxtableatstartofbodyhook

\begin{sphinxfigure-in-table}
\centering

\noindent\sphinxincludegraphics[width=50\sphinxpxdimen]{{image151}.png}
\end{sphinxfigure-in-table}\relax
\\
\sphinxbottomrule
\end{tabulary}
\sphinxtableafterendhook\par
\sphinxattableend\end{savenotes}

\item {} 
\sphinxAtStartPar
Se abrirá un nuevo documento de LibreOffice Writer con el nombre: \sphinxstylestrong{Sin titulo 2}

\item {} 
\sphinxAtStartPar
VUELVA al documento \sphinxstylestrong{G1013A05\sphinxhyphen{}creacion\_documento.odt} y en el menú superior PULSE: \sphinxguilabel{Archivo} > \sphinxguilabel{Nuevo} > \sphinxguilabel{Documento} de texto. Se abrirá un nuevo documento de LibreOffice Writer con el nombre: \sphinxstylestrong{Sin titulo 3}.

\item {} 
\sphinxAtStartPar
VUELVA al documento \sphinxstylestrong{G1013A05\sphinxhyphen{}creacion\_documento.odt} y PULSE la combinación de teclas: \sphinxkeyboard{\sphinxupquote{CTRL}} + \sphinxkeyboard{\sphinxupquote{U}}. Se abrirá un nuevo documento de LibreOffice Writer con el nombre: \sphinxstylestrong{Sin titulo 4}.

\item {} 
\sphinxAtStartPar
Con los archivos de LibreOffice Writer abiertos hasta el momento, ORGANICE cuatro ventanas (en su pantalla) de tal forma que se vean todas al mismo tiempo así:


\begin{savenotes}\sphinxattablestart
\sphinxthistablewithglobalstyle
\centering
\begin{tabulary}{\linewidth}[t]{T}
\sphinxtoprule
\sphinxtableatstartofbodyhook

\begin{sphinxfigure-in-table}
\centering
\capstart
\noindent\sphinxincludegraphics[width=700\sphinxpxdimen]{{image161}.png}
\sphinxfigcaption{Ventanas de LibreOffice Writer}\label{\detokenize{1_guias/10ofimatica/G1013:id10}}\end{sphinxfigure-in-table}\relax
\\
\sphinxbottomrule
\end{tabulary}
\sphinxtableafterendhook\par
\sphinxattableend\end{savenotes}

\item {} 
\sphinxAtStartPar
REALICE una \sphinxstylestrong{captura de pantalla} de su pantalla y guarde la imagen en la \sphinxstylestrong{carpeta de evidencias} del periodo actual con el nombre: \sphinxstylestrong{G1013A05\sphinxhyphen{}creacion\_documento.png}

\item {} 
\sphinxAtStartPar
En su cuaderno:

\begin{sphinxadmonition}{note}{En su cuaderno…}
\begin{itemize}
\item {} 
\sphinxAtStartPar
Describa las cuatro formas de crear documentos en \sphinxstylestrong{LibreOffice Writer}.

\end{itemize}
\end{sphinxadmonition}

\item {} 
\sphinxAtStartPar
Guarde los cambios y cierre el software LibreOffice Writer.

\end{enumerate}


\bigskip\hrule\bigskip



\subsection{A06: Guardado de documentos \sphinxstyleliteralintitle{\sphinxupquote{301}} (+0,2)}
\label{\detokenize{1_guias/10ofimatica/G1013:a06-guardado-de-documentos-301-0-2}}
\sphinxAtStartPar
\DUrole{sd-sphinx-override}{\DUrole{sd-badge}{\DUrole{sd-bg-warning}{\DUrole{sd-bg-text-warning}{Actividad práctica}}}}
\begin{enumerate}
\sphinxsetlistlabels{\arabic}{enumi}{enumii}{}{.}%
\item {} 
\sphinxAtStartPar
Ejecute el software LibreOffice Writer.

\item {} 
\sphinxAtStartPar
En la cabecera del documento, agregue la \sphinxstylestrong{información de identificación} con el título: \sphinxstylestrong{GUARDADO DE DOCUMENTOS}

\item {} 
\sphinxAtStartPar
GUARDE el documento en su \sphinxstylestrong{carpeta de evidencias} del periodo actual con el nombre: \sphinxstylestrong{G1013A06\sphinxhyphen{}guardado\_documento\_1.odt}

\item {} 
\sphinxAtStartPar
Copie y pegue el siguiente fragmento de texto a el documento:
\begin{quote}

\sphinxAtStartPar
“Lorem ipsum dolor sit amet, consectetur adipiscing elit, sed do eiusmod tempor incididunt ut labore et dolore magna aliqua. Ut enim ad minim veniam, quis nostrud exercitation ullamco laboris nisi ut aliquip ex ea commodo consequat. Duis aute irure dolor in reprehenderit in voluptate velit esse cillum dolore eu fugiat nulla pariatur. Excepteur sint occaecat cupidatat non proident, sunt in culpa qui officia deserunt mollit anim id est laborum”
\end{quote}

\item {} 
\sphinxAtStartPar
BORRE 30 palabras (las que usted quiera) del texto del fragmento anterior y en el \sphinxstylestrong{menú superior} PULSE: \sphinxguilabel{Archivo} > \sphinxguilabel{Guardar una copia…}

\item {} 
\sphinxAtStartPar
En la ventana de guardado, LOCALICE su carpeta de evidencias y guarde el documento con el nombre: \sphinxstylestrong{G1013A06\sphinxhyphen{}guardado\_documento\_2.odt}

\item {} 
\sphinxAtStartPar
Vuelva a abrir el primer documento: \sphinxstylestrong{G1013A06\sphinxhyphen{}guardado\_documento\_1.odt} y AGREGUE más texto al fragmento (el texto que usted quiera). Luego en el \sphinxstylestrong{menú superior} PULSE: \sphinxguilabel{Archivo} > \sphinxguilabel{Guardar como…}

\item {} 
\sphinxAtStartPar
En la ventana de guardado, IDENTIFIQUE las opciones de \sphinxstylestrong{tipo de documento} y seleccione: \sphinxstylestrong{Documento de word 2010\sphinxhyphen{}365 (.docx)} y guarde el documento con el nombre de: \sphinxstylestrong{G1013A06\sphinxhyphen{}guardado\_documento\_3.docx}


\begin{savenotes}\sphinxattablestart
\sphinxthistablewithglobalstyle
\centering
\begin{tabulary}{\linewidth}[t]{T}
\sphinxtoprule
\sphinxtableatstartofbodyhook

\begin{sphinxfigure-in-table}
\centering

\noindent\sphinxincludegraphics[width=500\sphinxpxdimen]{{image171}.png}
\end{sphinxfigure-in-table}\relax
\\
\sphinxbottomrule
\end{tabulary}
\sphinxtableafterendhook\par
\sphinxattableend\end{savenotes}

\begin{sphinxadmonition}{note}{ADVERTENCIA}

\sphinxAtStartPar
Si sale una ventana con una advertencia sobre \sphinxstylestrong{Confirmar el formato del archivo}, PULSE el botón:

\sphinxAtStartPar
\sphinxguilabel{Usar formato Documento de Word 2010\sphinxhyphen{}365}
\end{sphinxadmonition}

\item {} 
\sphinxAtStartPar
En su cuaderno:

\begin{sphinxadmonition}{note}{En su cuaderno…}
\begin{itemize}
\item {} 
\sphinxAtStartPar
Escriba las diferencias entre guardar un documento con la opción \sphinxstylestrong{Guardar una copia…} y la opción \sphinxstylestrong{Guardar como…} en LibreOffice Writer.

\end{itemize}
\end{sphinxadmonition}

\item {} 
\sphinxAtStartPar
Guarde los cambios y cierre el software LibreOffice Writer.

\end{enumerate}


\bigskip\hrule\bigskip



\subsection{A07: Cortar, copiar y pegar \sphinxstyleliteralintitle{\sphinxupquote{101}} (+0,1)}
\label{\detokenize{1_guias/10ofimatica/G1013:a07-cortar-copiar-y-pegar-101-0-1}}
\sphinxAtStartPar
\DUrole{sd-sphinx-override}{\DUrole{sd-badge}{\DUrole{sd-bg-warning}{\DUrole{sd-bg-text-warning}{Actividad práctica}}}}
\begin{enumerate}
\sphinxsetlistlabels{\arabic}{enumi}{enumii}{}{.}%
\item {} 
\sphinxAtStartPar
DESCARGUE el siguiente archivo y guárdelo en su carpeta de evidencias:
\begin{itemize}
\item {} 
\sphinxAtStartPar
G1013A07\sphinxhyphen{}mitos\_desordenados.odt \DUrole{xref}{\DUrole{download}{\DUrole{myst}{\DUrole{sd-sphinx-override}{\DUrole{sd-badge}{\DUrole{sd-bg-primary}{\DUrole{sd-bg-text-primary}{Descargar}}}}}}}

\end{itemize}

\item {} 
\sphinxAtStartPar
Abra el archivo: \sphinxstylestrong{G1013A07\sphinxhyphen{}mitos\_desordenados.odt} y en la cabecera del documento, agregue la \sphinxstylestrong{información de identificación} con el título: \sphinxstylestrong{CORTAR, COPIAR y PEGAR}

\begin{sphinxadmonition}{note}{¿Qué contiene el documento?}

\sphinxAtStartPar
Los párrafos del documento están desordenados.

\sphinxAtStartPar
Cada párrafo tienen una numeración para que sepa el orden en el que deben aparecer.
\end{sphinxadmonition}

\item {} 
\sphinxAtStartPar
Por medio de cortar, copiar y pegar: MUEVA los párrafos de manera que queden ordenados.

\begin{sphinxadmonition}{note}{Atajos de teclado para cortar, copiar y pegar}

\sphinxAtStartPar
El párrafo que desee (cortar, copiar o pegar) debe estar previamente seleccionado.

\sphinxAtStartPar
Para cortar, copiar y pegar se usan las siguientes combinaciones de teclado:


\begin{savenotes}\sphinxattablestart
\sphinxthistablewithglobalstyle
\centering
\begin{tabulary}{\linewidth}[t]{TT}
\sphinxtoprule
\sphinxstyletheadfamily 
\sphinxAtStartPar
Teclas
&\sphinxstyletheadfamily 
\sphinxAtStartPar
Acción
\\
\sphinxmidrule
\sphinxtableatstartofbodyhook
\sphinxAtStartPar
\sphinxkeyboard{\sphinxupquote{CTRL}} + \sphinxkeyboard{\sphinxupquote{X}}
&
\sphinxAtStartPar
Cortar
\\
\sphinxhline
\sphinxAtStartPar
\sphinxkeyboard{\sphinxupquote{CTRL}} + \sphinxkeyboard{\sphinxupquote{C}}
&
\sphinxAtStartPar
Copiar
\\
\sphinxhline
\sphinxAtStartPar
\sphinxkeyboard{\sphinxupquote{CTRL}} + \sphinxkeyboard{\sphinxupquote{V}}
&
\sphinxAtStartPar
Pegar
\\
\sphinxbottomrule
\end{tabulary}
\sphinxtableafterendhook\par
\sphinxattableend\end{savenotes}
\end{sphinxadmonition}

\item {} 
\sphinxAtStartPar
En su cuaderno:

\begin{sphinxadmonition}{note}{En su cuaderno…}
\begin{itemize}
\item {} 
\sphinxAtStartPar
Escriba los títulos ordenados del documento: \sphinxstylestrong{G1013A07\sphinxhyphen{}mitos\_desordenados.odt}

\end{itemize}
\end{sphinxadmonition}

\item {} 
\sphinxAtStartPar
Guarde los cambios y cierre el software LibreOffice Writer.

\end{enumerate}


\bigskip\hrule\bigskip



\subsection{A08: Buscar y reemplazar \sphinxstyleliteralintitle{\sphinxupquote{124}} (+0,1)}
\label{\detokenize{1_guias/10ofimatica/G1013:a08-buscar-y-reemplazar-124-0-1}}
\sphinxAtStartPar
\DUrole{sd-sphinx-override}{\DUrole{sd-badge}{\DUrole{sd-bg-warning}{\DUrole{sd-bg-text-warning}{Actividad práctica}}}}
\begin{enumerate}
\sphinxsetlistlabels{\arabic}{enumi}{enumii}{}{.}%
\item {} 
\sphinxAtStartPar
DESCARGUE el siguiente archivo y guárdelo en su carpeta de evidencias:
\begin{itemize}
\item {} 
\sphinxAtStartPar
G1013A08\sphinxhyphen{}proyecto\_escritorio\_libre.odt \DUrole{xref}{\DUrole{download}{\DUrole{myst}{\DUrole{sd-sphinx-override}{\DUrole{sd-badge}{\DUrole{sd-bg-primary}{\DUrole{sd-bg-text-primary}{Descargar}}}}}}}

\end{itemize}

\item {} 
\sphinxAtStartPar
ABRA el archivo: \sphinxstylestrong{G1013A08\sphinxhyphen{}proyecto\_escritorio\_libre.odt} y en la cabecera del documento, agregue la \sphinxstylestrong{información de identificación} con el título: \sphinxstylestrong{BUSCAR Y REEMPLAZAR}

\item {} 
\sphinxAtStartPar
Se debe realizar un búsqueda de la palabra aplicaciones en el texto del documento. Siga los siguientes pasos:
\begin{itemize}
\item {} 
\sphinxAtStartPar
PULSE las teclas \sphinxkeyboard{\sphinxupquote{CTRL}} + \sphinxkeyboard{\sphinxupquote{B}} para abrir la \sphinxstylestrong{barra de herramientas de Buscar} en la parte inferior de la interfaz.

\item {} 
\sphinxAtStartPar
MARQUE la opción \sphinxstylestrong{Distinguir mayúsculas y minúsculas}.

\item {} 
\sphinxAtStartPar
En el campo de búsqueda ESCRIBA la palabra: \sphinxstylestrong{aplicaciones}.

\item {} 
\sphinxAtStartPar
PULSE la opción que dice: \sphinxstylestrong{Buscar todo}

\end{itemize}


\begin{savenotes}\sphinxattablestart
\sphinxthistablewithglobalstyle
\centering
\begin{tabulary}{\linewidth}[t]{T}
\sphinxtoprule
\sphinxtableatstartofbodyhook

\begin{sphinxfigure-in-table}
\centering
\capstart
\noindent\sphinxincludegraphics[width=670\sphinxpxdimen]{{image181}.png}
\sphinxfigcaption{Barra de herramientas Buscar}\label{\detokenize{1_guias/10ofimatica/G1013:id11}}\end{sphinxfigure-in-table}\relax
\\
\sphinxbottomrule
\end{tabulary}
\sphinxtableafterendhook\par
\sphinxattableend\end{savenotes}
\begin{itemize}
\item {} 
\sphinxAtStartPar
IDENTIFIQUE un mensaje que aparecerá en la parte derecha de la \sphinxstylestrong{barra de herramientas de Buscar} el cual indica el número de veces que se encontró la expresión buscada.

\item {} 
\sphinxAtStartPar
En su cuaderno:

\begin{sphinxadmonition}{note}{En su cuaderno…}
\begin{itemize}
\item {} 
\sphinxAtStartPar
Escriba el número de veces que \sphinxstylestrong{LibreOffice Writer} encontró la palabra \sphinxstylestrong{aplicaciones.}

\end{itemize}
\end{sphinxadmonition}

\end{itemize}

\item {} 
\sphinxAtStartPar
Realice una captura de pantalla que muestre \sphinxstylestrong{el número de veces que se encontró la expresión buscada} y guarde la imagen en la \sphinxstylestrong{carpeta de evidencias} del periodo actual con el nombre: \sphinxstylestrong{G1013A08\sphinxhyphen{}proyecto\_escritorio\_libre\sphinxhyphen{}01.png}

\item {} 
\sphinxAtStartPar
Ahora se debe \sphinxstylestrong{reemplazar} todas las apariciones de la palabra \sphinxstylestrong{aplicaciones}. Siga los siguientes pasos:
\begin{itemize}
\item {} 
\sphinxAtStartPar
PULSE las teclas \sphinxkeyboard{\sphinxupquote{CTRL}} + \sphinxkeyboard{\sphinxupquote{ALT}} + \sphinxkeyboard{\sphinxupquote{B}} para abrir la ventana \sphinxguilabel{Buscar y reemplazar}.

\item {} 
\sphinxAtStartPar
En la ventana de \sphinxguilabel{Buscar y reemplazar} MARQUE las opciones:
\begin{itemize}
\item {} 
\sphinxAtStartPar
\sphinxstylestrong{Distinguir mayúsculas y minúsculas}

\item {} 
\sphinxAtStartPar
\sphinxstylestrong{Solo palabras completas}

\end{itemize}

\item {} 
\sphinxAtStartPar
En el campo \sphinxstylestrong{Buscar,} ESCRIBA la palabra: \sphinxstylestrong{aplicaciones}

\item {} 
\sphinxAtStartPar
PULSE el botón: \sphinxstylestrong{Buscar todo}.

\end{itemize}


\begin{savenotes}\sphinxattablestart
\sphinxthistablewithglobalstyle
\centering
\begin{tabulary}{\linewidth}[t]{T}
\sphinxtoprule
\sphinxtableatstartofbodyhook

\begin{sphinxfigure-in-table}
\centering
\capstart
\noindent\sphinxincludegraphics[width=670\sphinxpxdimen]{{image19}.png}
\sphinxfigcaption{Barra de herramientas Buscar y Reemplazar}\label{\detokenize{1_guias/10ofimatica/G1013:id12}}\end{sphinxfigure-in-table}\relax
\\
\sphinxbottomrule
\end{tabulary}
\sphinxtableafterendhook\par
\sphinxattableend\end{savenotes}
\begin{itemize}
\item {} 
\sphinxAtStartPar
En su cuaderno:

\begin{sphinxadmonition}{note}{En su cuaderno…}
\begin{itemize}
\item {} 
\sphinxAtStartPar
Escriba el número de veces que \sphinxstylestrong{LibreOffice Writer} encontró la palabra \sphinxstylestrong{aplicaciones.} (Debe coincidir con el número de veces mostrado anteriormente).

\end{itemize}
\end{sphinxadmonition}

\item {} 
\sphinxAtStartPar
En el campo \sphinxstylestrong{Reemplazar}, ESCRIBA la palabra: \sphinxstylestrong{programas}

\item {} 
\sphinxAtStartPar
PULSE el botón: \sphinxguilabel{Reemplazar todo}

\item {} 
\sphinxAtStartPar
Realice una \sphinxstylestrong{captura de pantalla} que muestre \sphinxstylestrong{el número de veces que se reemplazó la expresión buscada} y guarde la imagen en la \sphinxstylestrong{carpeta de evidencias} del periodo actual con el nombre: \sphinxstylestrong{G1013A08\sphinxhyphen{}proyecto\_escritorio\_libre\sphinxhyphen{}02.png}

\item {} 
\sphinxAtStartPar
CIERRE la ventana de \sphinxguilabel{Buscar y reemplazar}

\end{itemize}

\item {} 
\sphinxAtStartPar
Con la \sphinxstylestrong{barra de herramientas de Buscar}, ahora realice la búsqueda de la palabra: \sphinxstylestrong{programas}

\item {} 
\sphinxAtStartPar
En su cuaderno:

\begin{sphinxadmonition}{note}{En su cuaderno…}
\begin{itemize}
\item {} 
\sphinxAtStartPar
Escriba el número de veces que \sphinxstylestrong{LibreOffice Writer} encontró la palabra \sphinxstylestrong{programas.} (el número de veces debe ser mayor que 2 pero menor que 20).

\end{itemize}
\end{sphinxadmonition}

\item {} 
\sphinxAtStartPar
En su cuaderno:

\begin{sphinxadmonition}{note}{En su cuaderno…}
\begin{itemize}
\item {} 
\sphinxAtStartPar
Escriba las diferencias encontradas al usar la \sphinxstylestrong{barra de herramientas de Buscar} y la ventana \sphinxguilabel{Buscar y reemplazar}.

\end{itemize}
\end{sphinxadmonition}

\item {} 
\sphinxAtStartPar
Guarde los cambios y cierre el software LibreOffice Writer.

\end{enumerate}


\bigskip\hrule\bigskip



\subsection{A09: Corrector ortográfico \sphinxstyleliteralintitle{\sphinxupquote{102}} (+0,1)}
\label{\detokenize{1_guias/10ofimatica/G1013:a09-corrector-ortografico-102-0-1}}\begin{enumerate}
\sphinxsetlistlabels{\arabic}{enumi}{enumii}{}{.}%
\item {} 
\sphinxAtStartPar
EJECUTE el software \sphinxstylestrong{LibreOffice Writer}.

\item {} 
\sphinxAtStartPar
En la cabecera del documento, agregue la \sphinxstylestrong{información de identificación} con el título: \sphinxstylestrong{CORRECTOR ORTOGRÁFICO}.

\item {} 
\sphinxAtStartPar
GUARDE el documento en su \sphinxstylestrong{carpeta de evidencias} con el nombre: \sphinxstylestrong{G1013A09\sphinxhyphen{}corrector\_ortografico.odt}

\item {} 
\sphinxAtStartPar
Copie y pegue el siguiente texto en el documento:
\begin{quote}

\sphinxAtStartPar
En el año 2023, una extaña situacion ocurrio en la ciudad. 5 gatos desaparecieron sin dejar rastro. La policia investigo durante 3 semanas, pero no encontro ninguna pista. Luego, 2 tesigos afirmaron haber visto algo inusual cerca del parque. Uno de ellos dijo que vio una nave espacial, mientras que el otro juro que eran 7 duendes. El misterio se profundizo cuando se descubrieron 10 huellas giganes en el cesped. Los cientificos analisrron las huellas durante 4 dias y determinaron que eran de un animal desconocido. La prenso publico 12 articulos sobre el tema, generando gran expectacion. Al final, despues de 6 meses, el misterio se resolvio de manera inesperada: los gatos havian sido secuestrados por un grupo de niñoz traviesos que querian jugar con ellos. La polisia encontro a los gatos en una caza abandonada, sanos y salbos. La historia fue un gran exito, pero nadie pudo exxplicar las 15 huellas ni los avistamientos.
\end{quote}

\item {} 
\sphinxAtStartPar
PULSAR el botón: \sphinxstylestrong{Barra de herramientas estándar} > \sphinxguilabel{Revisar ortografía}


\begin{savenotes}\sphinxattablestart
\sphinxthistablewithglobalstyle
\centering
\begin{tabulary}{\linewidth}[t]{T}
\sphinxtoprule
\sphinxtableatstartofbodyhook

\begin{sphinxfigure-in-table}
\centering
\capstart
\noindent\sphinxincludegraphics[width=50\sphinxpxdimen]{{image201}.png}
\sphinxfigcaption{Corrector ortográfico}\label{\detokenize{1_guias/10ofimatica/G1013:id13}}\end{sphinxfigure-in-table}\relax
\\
\sphinxbottomrule
\end{tabulary}
\sphinxtableafterendhook\par
\sphinxattableend\end{savenotes}

\item {} 
\sphinxAtStartPar
En la ventana: \sphinxguilabel{Ortografía: Español (Colombia)}, corregir con las sugerencias que da el asistente de ortografía de LibreOffice Writer.
\begin{itemize}
\item {} 
\sphinxAtStartPar
En su cuaderno:

\begin{sphinxadmonition}{note}{En su cuaderno…}
\begin{itemize}
\item {} 
\sphinxAtStartPar
Escriba los errores ortográficos que \sphinxstylestrong{fueron detectados} por el asistente ortográfico de LibreOffice Writer.

\end{itemize}
\end{sphinxadmonition}

\end{itemize}


\begin{savenotes}\sphinxattablestart
\sphinxthistablewithglobalstyle
\centering
\begin{tabulary}{\linewidth}[t]{T}
\sphinxtoprule
\sphinxtableatstartofbodyhook

\begin{sphinxfigure-in-table}
\centering
\capstart
\noindent\sphinxincludegraphics[width=600\sphinxpxdimen]{{image211}.png}
\sphinxfigcaption{Ventana corrector ortográfico}\label{\detokenize{1_guias/10ofimatica/G1013:id14}}\end{sphinxfigure-in-table}\relax
\\
\sphinxbottomrule
\end{tabulary}
\sphinxtableafterendhook\par
\sphinxattableend\end{savenotes}

\item {} 
\sphinxAtStartPar
Una vez finalice con el asistente de ortografía, lea el texto y solucione manualmente los errores que quedan.

\item {} 
\sphinxAtStartPar
En su cuaderno:

\begin{sphinxadmonition}{note}{En su cuaderno…}
\begin{itemize}
\item {} 
\sphinxAtStartPar
Escriba los errores ortográficos que \sphinxstylestrong{NO fueron detectados} por el asistente ortográfico de LibreOffice Writer.

\end{itemize}
\end{sphinxadmonition}

\item {} 
\sphinxAtStartPar
Guarde los cambios y cierre el software LibreOffice Writer.

\end{enumerate}


\bigskip\hrule\bigskip



\subsection{A10: Presentación de evidencias (+3,0)}
\label{\detokenize{1_guias/10ofimatica/G1013:a10-presentacion-de-evidencias-3-0}}
\sphinxAtStartPar
\DUrole{sd-sphinx-override}{\DUrole{sd-badge}{\DUrole{sd-bg-danger}{\DUrole{sd-bg-text-danger}{Transferencia}}}}

\begin{sphinxadmonition}{note}{¿Cómo se evalúa?}
\begin{enumerate}
\sphinxsetlistlabels{\arabic}{enumi}{enumii}{}{.}%
\item {} 
\sphinxAtStartPar
Los estudiantes deben presentar lo siguiente:
\begin{itemize}
\item {} 
\sphinxAtStartPar
Carpeta de evidencias con todos los archivos editados en la presente guía. (Los archivos deben estar abiertos).

\item {} 
\sphinxAtStartPar
Cuaderno con las actividades desarrolladas.

\end{itemize}

\item {} 
\sphinxAtStartPar
El docente realiza unas preguntas para verificación de conocimientos.

\item {} 
\sphinxAtStartPar
El estudiante debe demostrar de forma práctica sus conocimientos en el uso de \sphinxstylestrong{LibreOffice Writer}.

\end{enumerate}
\end{sphinxadmonition}

\sphinxstepscope


\section{Guía 1014: Formato en LibreOffice Writer (P3)}
\label{\detokenize{1_guias/10ofimatica/G1014:guia-1014-formato-en-libreoffice-writer-p3}}\label{\detokenize{1_guias/10ofimatica/G1014::doc}}

\begin{savenotes}\sphinxattablestart
\sphinxthistablewithglobalstyle
\centering
\begin{tabular}[t]{\X{35}{100}\X{65}{100}}
\sphinxtoprule
\sphinxtableatstartofbodyhook

\begin{savenotes}\sphinxattablestart
\sphinxthistablewithglobalstyle
\centering
\begin{tabulary}{\linewidth}[t]{T}
\sphinxtoprule
\sphinxtableatstartofbodyhook
\sphinxAtStartPar
\sphinxincludegraphics{{escudo}.png}
\\
\sphinxbottomrule
\end{tabulary}
\sphinxtableafterendhook\par
\sphinxattableend\end{savenotes}
&

\begin{savenotes}\sphinxattablestart
\sphinxthistablewithglobalstyle
\raggedright
\begin{tabular}[t]{*{1}{\X{1}{1}}}
\sphinxtoprule
\sphinxtableatstartofbodyhook
\sphinxAtStartPar
\sphinxstylestrong{INSTITUCIÓN EDUCATIVA}: John F. Kennedy
\begin{itemize}
\item {} 
\sphinxAtStartPar
\sphinxstylestrong{DOCENTE}: Julian Camilo Fonseca Romero

\item {} 
\sphinxAtStartPar
\sphinxstylestrong{ASIGNATURA}: Ofimática

\item {} 
\sphinxAtStartPar
\sphinxstylestrong{GRADO}: DÉCIMO

\item {} 
\sphinxAtStartPar
\sphinxstylestrong{PERIODO:} 03

\item {} 
\sphinxAtStartPar
\sphinxstylestrong{TIEMPO}: 10 semanas (2 horas / semana)

\item {} 
\sphinxAtStartPar
\sphinxstylestrong{TIEMPO DE TRABAJO}: 20/20 horas

\end{itemize}
\\
\sphinxbottomrule
\end{tabular}
\sphinxtableafterendhook\par
\sphinxattableend\end{savenotes}
\\
\sphinxbottomrule
\end{tabular}
\sphinxtableafterendhook\par
\sphinxattableend\end{savenotes}

\begin{sphinxuseclass}{sd-tab-set}
\begin{sphinxuseclass}{sd-tab-item}\subsubsection*{Objetivo}

\begin{sphinxuseclass}{sd-tab-content}
\sphinxAtStartPar
\DUrole{sd-sphinx-override}{\DUrole{sd-badge}{\DUrole{sd-bg-light}{\DUrole{sd-bg-text-light}{Exploración}}}}
\DUrole{sd-sphinx-override}{\DUrole{sd-badge}{\DUrole{sd-bg-info}{\DUrole{sd-bg-text-info}{Estructuración}}}}
\DUrole{sd-sphinx-override}{\DUrole{sd-badge}{\DUrole{sd-bg-warning}{\DUrole{sd-bg-text-warning}{Actividad práctica}}}}
\DUrole{sd-sphinx-override}{\DUrole{sd-badge}{\DUrole{sd-bg-danger}{\DUrole{sd-bg-text-danger}{Transferencia}}}}
\begin{itemize}
\item {} 
\sphinxAtStartPar
\sphinxstylestrong{Temática}: Uso de formato de carácter, párrafo y estilo en LibreOffice Writer.

\item {} 
\sphinxAtStartPar
\sphinxstylestrong{Objetivo de aprendizaje}: Aplicar adecuadamente los formatos de carácter, párrafo y estilo en el software LibreOffice Writer.

\end{itemize}

\end{sphinxuseclass}
\end{sphinxuseclass}
\begin{sphinxuseclass}{sd-tab-item}\subsubsection*{Orientaciones curriculares}

\begin{sphinxuseclass}{sd-tab-content}\begin{itemize}
\item {} 
\sphinxAtStartPar
\sphinxstylestrong{Componente}: Uso y apropiación de la Tecnología y la Informática.

\item {} 
\sphinxAtStartPar
\sphinxstylestrong{Competencia}: Genero propuestas innovadoras para el uso y aprovechamiento de productos tecnológicos.

\item {} 
\sphinxAtStartPar
\sphinxstylestrong{Evidencia de aprendizaje}: Utilizo adecuadamente herramientas informáticas para la búsqueda, organización, procesamiento, sistematización, comunicación y difusión de ideas.

\end{itemize}

\end{sphinxuseclass}
\end{sphinxuseclass}
\begin{sphinxuseclass}{sd-tab-item}\subsubsection*{Criterios de evaluación}

\begin{sphinxuseclass}{sd-tab-content}\begin{itemize}
\item {} 
\sphinxAtStartPar
Actividades desarrolladas en su totalidad.

\item {} 
\sphinxAtStartPar
Participación en clase.

\item {} 
\sphinxAtStartPar
Explicación clara de conceptos.

\item {} 
\sphinxAtStartPar
Argumentación de ideas.

\item {} 
\sphinxAtStartPar
Cumplimiento de las reglas del aula.

\end{itemize}

\end{sphinxuseclass}
\end{sphinxuseclass}
\begin{sphinxuseclass}{sd-tab-item}\subsubsection*{Recursos (9)}

\begin{sphinxuseclass}{sd-tab-content}\begin{itemize}
\item {} 
\sphinxAtStartPar
\sphinxhref{https://www.libreoffice.org/download/}{Descarga de \sphinxstylestrong{LibreOffice}}.

\item {} 
\sphinxAtStartPar
G1014A04\sphinxhyphen{}practica\_fuentes.odt \DUrole{xref}{\DUrole{download}{\DUrole{myst}{\DUrole{sd-sphinx-override}{\DUrole{sd-badge}{\DUrole{sd-bg-primary}{\DUrole{sd-bg-text-primary}{Descargar}}}}}}}

\item {} 
\sphinxAtStartPar
G1014A05\sphinxhyphen{}mitos.odt \DUrole{xref}{\DUrole{download}{\DUrole{myst}{\DUrole{sd-sphinx-override}{\DUrole{sd-badge}{\DUrole{sd-bg-primary}{\DUrole{sd-bg-text-primary}{Descargar}}}}}}}

\item {} 
\sphinxAtStartPar
G1014A06\sphinxhyphen{}enfermedades.odt \DUrole{xref}{\DUrole{download}{\DUrole{myst}{\DUrole{sd-sphinx-override}{\DUrole{sd-badge}{\DUrole{sd-bg-primary}{\DUrole{sd-bg-text-primary}{Descargar}}}}}}}

\item {} 
\sphinxAtStartPar
G1014A07\sphinxhyphen{}nueces.odt \DUrole{xref}{\DUrole{download}{\DUrole{myst}{\DUrole{sd-sphinx-override}{\DUrole{sd-badge}{\DUrole{sd-bg-primary}{\DUrole{sd-bg-text-primary}{Descargar}}}}}}}

\item {} 
\sphinxAtStartPar
G1014A08\sphinxhyphen{}certificado.odt \DUrole{xref}{\DUrole{download}{\DUrole{myst}{\DUrole{sd-sphinx-override}{\DUrole{sd-badge}{\DUrole{sd-bg-primary}{\DUrole{sd-bg-text-primary}{Descargar}}}}}}}

\item {} 
\sphinxAtStartPar
G1014A09\sphinxhyphen{}temario.odt \DUrole{xref}{\DUrole{download}{\DUrole{myst}{\DUrole{sd-sphinx-override}{\DUrole{sd-badge}{\DUrole{sd-bg-primary}{\DUrole{sd-bg-text-primary}{Descargar}}}}}}}

\item {} 
\sphinxAtStartPar
G1014A10\sphinxhyphen{}curriculum.odt \DUrole{xref}{\DUrole{download}{\DUrole{myst}{\DUrole{sd-sphinx-override}{\DUrole{sd-badge}{\DUrole{sd-bg-primary}{\DUrole{sd-bg-text-primary}{Descargar}}}}}}}

\item {} 
\sphinxAtStartPar
G1014A11\sphinxhyphen{}estilo\_personalizado.odt \DUrole{xref}{\DUrole{download}{\DUrole{myst}{\DUrole{sd-sphinx-override}{\DUrole{sd-badge}{\DUrole{sd-bg-primary}{\DUrole{sd-bg-text-primary}{Descargar}}}}}}}

\end{itemize}

\end{sphinxuseclass}
\end{sphinxuseclass}
\end{sphinxuseclass}

\subsection{Mensaje del autor}
\label{\detokenize{1_guias/10ofimatica/G1014:mensaje-del-autor}}\begin{quote}

\sphinxAtStartPar
“Esta guía didáctica ha sido diseñada para facilitar el aprendizaje práctico y efectivo del formato de documentos en LibreOffice Writer. El objetivo es que, a través de actividades claras y orientadas, los estudiantes desarrollen competencias en el manejo de herramientas digitales, promoviendo la autonomía y la creatividad en la elaboración de textos”.

\begin{flushright}
---Camilo Fonseca
\end{flushright}
\end{quote}


\bigskip\hrule\bigskip



\subsection{A00: Introducción (+0,1)}
\label{\detokenize{1_guias/10ofimatica/G1014:a00-introduccion-0-1}}
\sphinxAtStartPar
\DUrole{sd-sphinx-override}{\DUrole{sd-badge}{\DUrole{sd-bg-light}{\DUrole{sd-bg-text-light}{Exploración}}}}
\begin{enumerate}
\sphinxsetlistlabels{\arabic}{enumi}{enumii}{}{.}%
\item {} 
\sphinxAtStartPar
Transcriba en su cuaderno el objetivo, las orientaciones curriculares y los criterios de evaluación de la presente guía.

\item {} 
\sphinxAtStartPar
Lea el mensaje del autor.

\end{enumerate}


\bigskip\hrule\bigskip



\subsection{A01: LibreOffice Writer (+0,1)}
\label{\detokenize{1_guias/10ofimatica/G1014:a01-libreoffice-writer-0-1}}
\sphinxAtStartPar
\DUrole{sd-sphinx-override}{\DUrole{sd-badge}{\DUrole{sd-bg-info}{\DUrole{sd-bg-text-info}{Estructuración}}}}
\begin{enumerate}
\sphinxsetlistlabels{\arabic}{enumi}{enumii}{}{.}%
\item {} 
\sphinxAtStartPar
Lea los siguientes conceptos:

\end{enumerate}


\bigskip\hrule\bigskip



\subsubsection{Formato de texto}
\label{\detokenize{1_guias/10ofimatica/G1014:formato-de-texto}}

\begin{savenotes}\sphinxattablestart
\sphinxthistablewithglobalstyle
\centering
\begin{tabulary}{\linewidth}[t]{T}
\sphinxtoprule
\sphinxtableatstartofbodyhook

\begin{sphinxfigure-in-table}
\centering
\capstart
\noindent\sphinxincludegraphics[width=700\sphinxpxdimen]{{image018}.png}
\sphinxfigcaption{Barra de herramientas formato}\label{\detokenize{1_guias/10ofimatica/G1014:id1}}\end{sphinxfigure-in-table}\relax
\\
\sphinxbottomrule
\end{tabulary}
\sphinxtableafterendhook\par
\sphinxattableend\end{savenotes}

\sphinxAtStartPar
La forma más productiva, eficiente y elegante para dar formato al texto del documento es mediante los estilos, entendidos como conjuntos de formatos.

\sphinxAtStartPar
Pero para poder aplicar, modificar o crear estilos es necesario conocer los diferentes formatos que podemos aplicar al tipo de letra, párrafos o listas de nuestro documento.

\sphinxAtStartPar
Por este motivo, es recomendable familiarizarnos con los diferentes formatos que podemos utilizar con el procesador de textos LibreOffice Writer y las diferentes maneras de aplicarlos.


\bigskip\hrule\bigskip



\subsubsection{Trabajando con estilos}
\label{\detokenize{1_guias/10ofimatica/G1014:trabajando-con-estilos}}

\begin{savenotes}\sphinxattablestart
\sphinxthistablewithglobalstyle
\centering
\begin{tabulary}{\linewidth}[t]{T}
\sphinxtoprule
\sphinxtableatstartofbodyhook

\begin{sphinxfigure-in-table}
\centering
\capstart
\noindent\sphinxincludegraphics[width=200\sphinxpxdimen]{{image026}.png}
\sphinxfigcaption{Jerarquía de estilos}\label{\detokenize{1_guias/10ofimatica/G1014:id2}}\end{sphinxfigure-in-table}\relax
\\
\sphinxbottomrule
\end{tabulary}
\sphinxtableafterendhook\par
\sphinxattableend\end{savenotes}

\sphinxAtStartPar
Con \sphinxstylestrong{LibreOffice Writer} podemos aplicar formatos directamente al texto del documento. Pero si trabajamos con documentos largos y estructurados como memorias, proyectos o informes, la aplicación directa de formatos es muy improductiva.

\sphinxAtStartPar
La solución está en trabajar con estilos. Los estilos son conjuntos de formatos agrupados bajo un nombre y gracias a ellos, podemos dotar a nuestros documentos de la estructura adecuada de títulos y subtítulos, crear índices automáticamente o numerar capítulos. Modificando un estilo podemos cambiar en cuatro clics la apariencia de un documento de decenas o cientos de páginas.

\begin{sphinxadmonition}{note}{RECOMENDACIONES USO DE ESTILOS}

\sphinxAtStartPar
Para el uso de estilos se recomienda lo siguiente:
\begin{itemize}
\item {} 
\sphinxAtStartPar
Los colores de texto y fondo deben estar suficientemente contrastados.

\item {} 
\sphinxAtStartPar
El tamaño mínimo recomendable para la tipografía debe ser de 12 puntos. Nunca usar
tipografías de tamaño inferior a 10 puntos.

\item {} 
\sphinxAtStartPar
Uso de tipografías \sphinxstyleemphasis{sans serif} (sin remate) habituales. Arial o Verdana son muy buenas elecciones.

\item {} 
\sphinxAtStartPar
Evitar el uso de variantes \sphinxstyleemphasis{thin}, \sphinxstyleemphasis{light} o \sphinxstyleemphasis{narrow} de las tipografías (variantes más estrechas de los tipos de letra).

\item {} 
\sphinxAtStartPar
Precaución con el uso de las negritas, cursivas o subrayados.

\item {} 
\sphinxAtStartPar
Evitar el uso de efectos de contorno o intermitencias.

\item {} 
\sphinxAtStartPar
Evitar el uso de características incompatibles con otros formatos de documento (sup\_rarayados, por ejemplo)

\item {} 
\sphinxAtStartPar
El interlineado simple no facilita la lectura; casi siempre será más apropiado un interlineado
proporcional del 120\% como mínimo.

\item {} 
\sphinxAtStartPar
Establecer un espaciado anterior y posterior al párrafo.

\item {} 
\sphinxAtStartPar
Evitar el uso de características incompatibles con otros formatos de documento (rellenos
degradados o de imagen, por ejemplo)

\end{itemize}
\end{sphinxadmonition}

\sphinxAtStartPar
Cuando se desea distribuir documentos de forma impresa o en formato PDF, tan sólo se necesita disponer de un documento bien estructurado y organizado. Para ello se tiene en cuenta otras configuraciones como el tamaño y formato del papel donde se imprimirá, los márgenes del documento, la orientación, los encabezados y pies que se deben repetir en todas las páginas o la numeración de las páginas.


\bigskip\hrule\bigskip



\subsection{A02: Taller (+0,8)}
\label{\detokenize{1_guias/10ofimatica/G1014:a02-taller-0-8}}
\sphinxAtStartPar
\DUrole{sd-sphinx-override}{\DUrole{sd-badge}{\DUrole{sd-bg-warning}{\DUrole{sd-bg-text-warning}{Actividad práctica}}}}
\begin{enumerate}
\sphinxsetlistlabels{\arabic}{enumi}{enumii}{}{.}%
\item {} 
\sphinxAtStartPar
Responda las siguientes preguntas en su cuaderno:
\begin{itemize}
\item {} 
\sphinxAtStartPar
¿Qué es dar formato a un documento?

\item {} 
\sphinxAtStartPar
¿Cuál es la forma más eficiente y productiva de dar formato a un texto?

\item {} 
\sphinxAtStartPar
¿Qué son los estilos?

\end{itemize}

\item {} 
\sphinxAtStartPar
Con 10 terminos de las \sphinxstylestrong{recomendaciones en el uso de estilos}, en una hoja de cuaderno cree una \sphinxstyleemphasis{sopa de letras}.

\end{enumerate}


\bigskip\hrule\bigskip



\subsection{A03: Formato de carácter 1.0 \sphinxstyleliteralintitle{\sphinxupquote{101}} (+0,1)}
\label{\detokenize{1_guias/10ofimatica/G1014:a03-formato-de-caracter-1-0-101-0-1}}
\sphinxAtStartPar
\DUrole{sd-sphinx-override}{\DUrole{sd-badge}{\DUrole{sd-bg-warning}{\DUrole{sd-bg-text-warning}{Actividad práctica}}}}
\begin{enumerate}
\sphinxsetlistlabels{\arabic}{enumi}{enumii}{}{.}%
\item {} 
\sphinxAtStartPar
EJECUTE el software LibreOffice Writer.

\item {} 
\sphinxAtStartPar
En la cabecera del documento, agregue la \sphinxstylestrong{información de identificación} con el título: \sphinxstylestrong{FORMATO DE CARÁCTER 1.0}

\item {} 
\sphinxAtStartPar
GUARDE el documento en su \sphinxstylestrong{carpeta de evidencias} del periodo actual con el nombre: \sphinxstylestrong{G1014A03\sphinxhyphen{}formato\_caracter\_1.odt}

\item {} 
\sphinxAtStartPar
AGREGUE el siguiente fragmento de texto al documento:
\begin{quote}

\sphinxAtStartPar
“La ofimática es el conjunto de aplicaciones y herramientas informáticas diseñadas para facilitar y mejorar las tareas relacionadas con la creación, edición, almacenamiento y distribución de información en el entorno de una oficina o empresa”.
\end{quote}

\item {} 
\sphinxAtStartPar
En el documento, SELECCIONE el texto recién copiado y en el Menú superior PULSE: \sphinxguilabel{Formato} > \sphinxguilabel{Carácter…} para abrir la ventana \sphinxguilabel{Carácter}.


\begin{savenotes}\sphinxattablestart
\sphinxthistablewithglobalstyle
\centering
\begin{tabulary}{\linewidth}[t]{T}
\sphinxtoprule
\sphinxtableatstartofbodyhook

\begin{sphinxfigure-in-table}
\centering
\capstart
\noindent\sphinxincludegraphics[width=260\sphinxpxdimen]{{image035}.png}
\sphinxfigcaption{\sphinxstylestrong{Menú} formato de carácter}\label{\detokenize{1_guias/10ofimatica/G1014:id3}}\end{sphinxfigure-in-table}\relax
\\
\sphinxbottomrule
\end{tabulary}
\sphinxtableafterendhook\par
\sphinxattableend\end{savenotes}

\item {} 
\sphinxAtStartPar
En la ventana \sphinxguilabel{Carácter} configure las siguientes opciones:


\begin{savenotes}\sphinxattablestart
\sphinxthistablewithglobalstyle
\centering
\begin{tabulary}{\linewidth}[t]{T}
\sphinxtoprule
\sphinxtableatstartofbodyhook

\begin{sphinxfigure-in-table}
\centering
\capstart
\noindent\sphinxincludegraphics[width=550\sphinxpxdimen]{{image045}.png}
\sphinxfigcaption{\sphinxstylestrong{Menú} Ventana formato de carácter}\label{\detokenize{1_guias/10ofimatica/G1014:id4}}\end{sphinxfigure-in-table}\relax
\\
\sphinxbottomrule
\end{tabulary}
\sphinxtableafterendhook\par
\sphinxattableend\end{savenotes}
\begin{itemize}
\item {} 
\sphinxAtStartPar
En la pestaña \sphinxguilabel{Tipo de letra} configurar:
\begin{itemize}
\item {} 
\sphinxAtStartPar
\sphinxstylestrong{Tipo de letra:} Liberation Serif

\item {} 
\sphinxAtStartPar
\sphinxstylestrong{Estilo:} Negrita Itálica

\item {} 
\sphinxAtStartPar
\sphinxstylestrong{Tamaño:} 12 pt

\end{itemize}

\item {} 
\sphinxAtStartPar
En la pestaña \sphinxguilabel{Efectos tipográficos} configurar:
\begin{itemize}
\item {} 
\sphinxAtStartPar
\sphinxstylestrong{Color de letra:} Verde claro 1

\item {} 
\sphinxAtStartPar
\sphinxstylestrong{Suprarrayado:} Punteado, Color: Azul

\item {} 
\sphinxAtStartPar
\sphinxstylestrong{Tachado:} Con /

\item {} 
\sphinxAtStartPar
\sphinxstylestrong{Subrayado:} Trazo, Color: Magenta

\item {} 
\sphinxAtStartPar
\sphinxstylestrong{Mayusculación:} Mayúsculas iniciales

\item {} 
\sphinxAtStartPar
\sphinxstylestrong{Relieve:} Sin

\item {} 
\sphinxAtStartPar
Marcar opción de \sphinxstylestrong{Contorno}

\item {} 
\sphinxAtStartPar
Marcar opción de \sphinxstylestrong{Sombra}

\item {} 
\sphinxAtStartPar
Desmarcar opción de \sphinxstylestrong{Oculto}

\end{itemize}

\item {} 
\sphinxAtStartPar
En la pestaña \sphinxguilabel{Bordes} configurar:
\begin{itemize}
\item {} 
\sphinxAtStartPar
En la sección \sphinxstylestrong{Disposición de líneas} configurar:
\begin{itemize}
\item {} 
\sphinxAtStartPar
\sphinxstylestrong{Definido por el usuario:} \sphinxstyleemphasis{Seleccionar bordes superior y derecho así:}

\noindent{\hspace*{\fill}\sphinxincludegraphics[width=250\sphinxpxdimen]{{image055}.png}\hspace*{\fill}}

\end{itemize}

\item {} 
\sphinxAtStartPar
En la sección \sphinxstylestrong{Línea} configurar:
\begin{itemize}
\item {} 
\sphinxAtStartPar
\sphinxstylestrong{Estilo:} \sphinxstyleemphasis{Traza\sphinxhyphen{}punto\sphinxhyphen{}punto}

\item {} 
\sphinxAtStartPar
\sphinxstylestrong{Color:} Amarillo

\item {} 
\sphinxAtStartPar
\sphinxstylestrong{Grosor:} Muy grueso (4,5 pt)

\noindent{\hspace*{\fill}\sphinxincludegraphics[width=250\sphinxpxdimen]{{image065}.png}\hspace*{\fill}}

\end{itemize}

\item {} 
\sphinxAtStartPar
En la sección \sphinxstylestrong{Estilo de sombra} configurar:
\begin{itemize}
\item {} 
\sphinxAtStartPar
\sphinxstylestrong{Posición:} Proyectar sombra a la parte inferior derecha

\item {} 
\sphinxAtStartPar
\sphinxstylestrong{Color:} Lima

\end{itemize}

\end{itemize}

\item {} 
\sphinxAtStartPar
Pulsar el botón \sphinxguilabel{Aceptar}

\end{itemize}

\item {} 
\sphinxAtStartPar
En su cuaderno:

\begin{sphinxadmonition}{note}{En su cuaderno…}

\sphinxAtStartPar
Responda las siguientes preguntas:
\begin{itemize}
\item {} 
\sphinxAtStartPar
¿Que pasó con el texto que había pegado inicialmente?

\item {} 
\sphinxAtStartPar
Al momento de crear informes, memorias, cartas o proyectos ¿Qué se necesita para que el contenido tenga un buen formato?

\end{itemize}
\end{sphinxadmonition}

\item {} 
\sphinxAtStartPar
Guarde los cambios y cierre el software \sphinxstylestrong{LibreOffice Writer}.

\end{enumerate}


\bigskip\hrule\bigskip



\subsection{A04: Formato de carácter 2.0 \sphinxstyleliteralintitle{\sphinxupquote{102}} (+0,1)}
\label{\detokenize{1_guias/10ofimatica/G1014:a04-formato-de-caracter-2-0-102-0-1}}
\sphinxAtStartPar
\DUrole{sd-sphinx-override}{\DUrole{sd-badge}{\DUrole{sd-bg-warning}{\DUrole{sd-bg-text-warning}{Actividad práctica}}}}
\begin{enumerate}
\sphinxsetlistlabels{\arabic}{enumi}{enumii}{}{.}%
\item {} 
\sphinxAtStartPar
DESCARGUE el siguiente archivo y guárdelo en su carpeta de evidencias:
\begin{itemize}
\item {} 
\sphinxAtStartPar
G1014A04\sphinxhyphen{}practica\_fuentes.odt \DUrole{xref}{\DUrole{download}{\DUrole{myst}{\DUrole{sd-sphinx-override}{\DUrole{sd-badge}{\DUrole{sd-bg-primary}{\DUrole{sd-bg-text-primary}{Descargar}}}}}}}

\end{itemize}

\item {} 
\sphinxAtStartPar
Abra el archivo: \sphinxstylestrong{G1014A04\sphinxhyphen{}practica\_fuentes.odt} y en la cabecera del documento, agregue la \sphinxstylestrong{información de identificación} con el título: \sphinxstylestrong{FORMATO DE CARÁCTER 2.0}

\begin{sphinxadmonition}{note}{¿Qué contiene el documento?}

\sphinxAtStartPar
El documento tiene varios párrafos que especifican la configuración de formato de carácter que debe tener cada uno.
\end{sphinxadmonition}

\item {} 
\sphinxAtStartPar
SELECCIONE cada párrafo y aplique la \sphinxstylestrong{configuración indicada} PULSANDO en el menú superior: \sphinxguilabel{Formato} > \sphinxguilabel{carácter…}

\begin{sphinxadmonition}{note}{Advertencia}
\begin{itemize}
\item {} 
\sphinxAtStartPar
NO USE la \sphinxstylestrong{barra de herramientas de formato}!!!

\item {} 
\sphinxAtStartPar
Puede que \sphinxstylestrong{algunas opciones no coincidan}…

\end{itemize}

\begin{sphinxadmonition}{note}{En su cuaderno…}

\sphinxAtStartPar
En caso de que \sphinxstylestrong{algunas opciones no coincidan} (por ejemplo los colores), elegir la opción más cercana y en su cuaderno escribir las opciones que no coincidieron y cómo las reemplazó.
\end{sphinxadmonition}
\end{sphinxadmonition}

\item {} 
\sphinxAtStartPar
Después de aplicar todos los formatos a los párrafos, el documento debería tener el siguiente aspecto:


\begin{savenotes}\sphinxattablestart
\sphinxthistablewithglobalstyle
\centering
\begin{tabulary}{\linewidth}[t]{T}
\sphinxtoprule
\sphinxtableatstartofbodyhook

\begin{sphinxfigure-in-table}
\centering

\noindent\sphinxincludegraphics[width=500\sphinxpxdimen]{{image075}.png}
\end{sphinxfigure-in-table}\relax
\\
\sphinxbottomrule
\end{tabulary}
\sphinxtableafterendhook\par
\sphinxattableend\end{savenotes}

\item {} 
\sphinxAtStartPar
En su cuaderno:

\begin{sphinxadmonition}{note}{En su cuaderno…}

\sphinxAtStartPar
Responda la siguiente pregunta:
\begin{itemize}
\item {} 
\sphinxAtStartPar
¿Cuál es la ventaja de aplicar \sphinxstylestrong{formato de carácter} por medio del \sphinxstylestrong{menú superior} frente a aplicar \sphinxstylestrong{formato de carácter} por medio de la \sphinxstylestrong{barra de herramientas formato}?

\end{itemize}
\end{sphinxadmonition}

\item {} 
\sphinxAtStartPar
Guarde los cambios y luego cierre el software \sphinxstylestrong{LibreOffice Writer}.

\end{enumerate}


\bigskip\hrule\bigskip



\subsection{A05: Formato de carácter 3.0 \sphinxstyleliteralintitle{\sphinxupquote{101}} (+0,1)}
\label{\detokenize{1_guias/10ofimatica/G1014:a05-formato-de-caracter-3-0-101-0-1}}
\sphinxAtStartPar
\DUrole{sd-sphinx-override}{\DUrole{sd-badge}{\DUrole{sd-bg-warning}{\DUrole{sd-bg-text-warning}{Actividad práctica}}}}
\begin{enumerate}
\sphinxsetlistlabels{\arabic}{enumi}{enumii}{}{.}%
\item {} 
\sphinxAtStartPar
DESCARGUE el siguiente archivo y guárdelo en su carpeta de evidencias:
\begin{itemize}
\item {} 
\sphinxAtStartPar
G1014A05\sphinxhyphen{}mitos.odt \DUrole{xref}{\DUrole{download}{\DUrole{myst}{\DUrole{sd-sphinx-override}{\DUrole{sd-badge}{\DUrole{sd-bg-primary}{\DUrole{sd-bg-text-primary}{Descargar}}}}}}}

\end{itemize}

\item {} 
\sphinxAtStartPar
Abra el archivo: \sphinxstylestrong{G1014A05\sphinxhyphen{}mitos.odt} y en la cabecera del documento, agregue la \sphinxstylestrong{información de identificación} con el título: \sphinxstylestrong{FORMATO DE CARÁCTER 3.0}

\begin{sphinxadmonition}{note}{¿Que contiene el documento?}

\sphinxAtStartPar
El documento contienen las descripciones de Dioses de la mitología griega. A los títulos: ADONIS, AFRODITA y AGON se les debe aplicar una configuración de \sphinxstylestrong{formato de carácter}.
\end{sphinxadmonition}

\item {} 
\sphinxAtStartPar
Seleccione el título \sphinxstylestrong{ADONIS} y PULSE en el \sphinxstylestrong{Menú superior} > \sphinxguilabel{Formato} > \sphinxguilabel{Carácter…}
\begin{itemize}
\item {} 
\sphinxAtStartPar
En la pestaña \sphinxguilabel{Resalte} configurar:
\begin{itemize}
\item {} 
\sphinxAtStartPar
\sphinxstylestrong{Color:} Verde claro 3

\end{itemize}

\item {} 
\sphinxAtStartPar
En la pestaña \sphinxguilabel{Bordes} configurar:
\begin{itemize}
\item {} 
\sphinxAtStartPar
En la sección \sphinxstylestrong{Disposición de líneas} configurar:
\begin{itemize}
\item {} 
\sphinxAtStartPar
\sphinxstylestrong{Preconfigs:} Los cuatro bordes

\end{itemize}

\item {} 
\sphinxAtStartPar
En la sección \sphinxstylestrong{Línea} configurar:
\begin{itemize}
\item {} 
\sphinxAtStartPar
\sphinxstylestrong{Estilo:} Sólido

\item {} 
\sphinxAtStartPar
\sphinxstylestrong{Color:} Azul

\item {} 
\sphinxAtStartPar
\sphinxstylestrong{Grosor:} Grueso 2.25 pt

\end{itemize}

\item {} 
\sphinxAtStartPar
En la sección \sphinxstylestrong{Separación} configurar:
\begin{itemize}
\item {} 
\sphinxAtStartPar
0.50cm por cada lado (izquierda, derecha, superior e inferior)

\end{itemize}

\item {} 
\sphinxAtStartPar
En la sección \sphinxstylestrong{Estilo de sombra} configurar:
\begin{itemize}
\item {} 
\sphinxAtStartPar
\sphinxstylestrong{Posición:} Abajo y a la derecha

\item {} 
\sphinxAtStartPar
\sphinxstylestrong{Color:} Azul claro 3

\end{itemize}

\end{itemize}

\end{itemize}

\item {} 
\sphinxAtStartPar
Aplique la misma configuración de \sphinxstylestrong{formato de carácter} a los títulos \sphinxstylestrong{AFRODITA} y \sphinxstylestrong{AGÓN}.

\item {} 
\sphinxAtStartPar
En cada párrafo, RESALTE la primera palabra con color amarillo.

\begin{sphinxadmonition}{warning}{Advertencia:}
\sphinxAtStartPar
Use el \sphinxstylestrong{formato de carácter}. No use la \sphinxstylestrong{barra de herramientas formato}.
\end{sphinxadmonition}

\item {} 
\sphinxAtStartPar
Después de aplicar \sphinxstylestrong{todos los formatos} al texto, el documento debería tener el siguiente aspecto:


\begin{savenotes}\sphinxattablestart
\sphinxthistablewithglobalstyle
\centering
\begin{tabulary}{\linewidth}[t]{T}
\sphinxtoprule
\sphinxtableatstartofbodyhook

\begin{sphinxfigure-in-table}
\centering

\noindent\sphinxincludegraphics[width=500\sphinxpxdimen]{{image085}.png}
\end{sphinxfigure-in-table}\relax
\\
\sphinxbottomrule
\end{tabulary}
\sphinxtableafterendhook\par
\sphinxattableend\end{savenotes}

\item {} 
\sphinxAtStartPar
En su cuaderno:

\begin{sphinxadmonition}{note}{En su cuaderno…}

\sphinxAtStartPar
Responda las siguientes preguntas:
\begin{itemize}
\item {} 
\sphinxAtStartPar
¿Que desventajas ha encontrado hasta el momento usando la herramienta de \sphinxstylestrong{formato de carácter (Menú superior > Formato > Carácter…)}?

\item {} 
\sphinxAtStartPar
¿Que pasaría si en vez de 3, hubiesen sido 300 títulos a los que era necesario configurar formato de carácter?

\end{itemize}
\end{sphinxadmonition}

\item {} 
\sphinxAtStartPar
Guarde los cambios y cierre el software LibreOffice Writer.

\end{enumerate}


\bigskip\hrule\bigskip



\subsection{A06: Clonando formato \sphinxstyleliteralintitle{\sphinxupquote{101}} (+0,1)}
\label{\detokenize{1_guias/10ofimatica/G1014:a06-clonando-formato-101-0-1}}
\sphinxAtStartPar
\DUrole{sd-sphinx-override}{\DUrole{sd-badge}{\DUrole{sd-bg-warning}{\DUrole{sd-bg-text-warning}{Actividad práctica}}}}
\begin{enumerate}
\sphinxsetlistlabels{\arabic}{enumi}{enumii}{}{.}%
\item {} 
\sphinxAtStartPar
DESCARGUE el siguiente archivo y guárdelo en su carpeta de evidencias:
\begin{itemize}
\item {} 
\sphinxAtStartPar
G1014A06\sphinxhyphen{}enfermedades.odt \DUrole{xref}{\DUrole{download}{\DUrole{myst}{\DUrole{sd-sphinx-override}{\DUrole{sd-badge}{\DUrole{sd-bg-primary}{\DUrole{sd-bg-text-primary}{Descargar}}}}}}}

\end{itemize}

\item {} 
\sphinxAtStartPar
Abra el archivo: \sphinxstylestrong{G1014A06\sphinxhyphen{}enfermedades.odt} y en la cabecera del documento, agregue la \sphinxstylestrong{información de identificación} con el título: \sphinxstylestrong{CLONANDO FORMATO}

\begin{sphinxadmonition}{note}{¿Qué contiene el documento?}

\sphinxAtStartPar
En este documento se encuentra la descripción de unas enfermedades.

\sphinxAtStartPar
Es necesario dar formato a los títulos de las enfermedades.

\sphinxAtStartPar
Se debe aplicar un formato de carácter en el título Gota. Luego este formato será clonado en los demás títulos (Artrosis y Aftas).
\end{sphinxadmonition}

\item {} 
\sphinxAtStartPar
SELECCIONE el título \sphinxstylestrong{Gota} y aplique las siguientes configuraciones de \sphinxstylestrong{formato de carácter:}
\begin{itemize}
\item {} 
\sphinxAtStartPar
En la pestaña \sphinxguilabel{Tipo de letra} configurar:
\begin{itemize}
\item {} 
\sphinxAtStartPar
\sphinxstylestrong{Familia:} Verdana

\item {} 
\sphinxAtStartPar
\sphinxstylestrong{Tamaño:} 28 pt

\item {} 
\sphinxAtStartPar
\sphinxstylestrong{Estilo:} Negrita

\end{itemize}

\item {} 
\sphinxAtStartPar
En la pestaña \sphinxguilabel{Efectos tipográficos} configurar:
\begin{itemize}
\item {} 
\sphinxAtStartPar
En la sección \sphinxstylestrong{Efectos} configurar:
\begin{itemize}
\item {} 
\sphinxAtStartPar
MARCAR la opción de \sphinxstylestrong{contorno}.

\end{itemize}

\end{itemize}

\item {} 
\sphinxAtStartPar
En la pestaña \sphinxguilabel{Resalte} configurar:
\begin{itemize}
\item {} 
\sphinxAtStartPar
\sphinxstylestrong{Color:} verde claro 4

\end{itemize}

\item {} 
\sphinxAtStartPar
En la pestaña \sphinxguilabel{Bordes} configurar:
\begin{itemize}
\item {} 
\sphinxAtStartPar
En la sección \sphinxstylestrong{Disposición de líneas} configurar:
\begin{itemize}
\item {} 
\sphinxAtStartPar
\sphinxstylestrong{Preconfigs:} los cuatro bordes

\end{itemize}

\item {} 
\sphinxAtStartPar
En la sección \sphinxstylestrong{Separación} configurar:
\begin{itemize}
\item {} 
\sphinxAtStartPar
0,20 cm por cada lado (izquierda, derecha, superior e inferior)

\end{itemize}

\end{itemize}

\item {} 
\sphinxAtStartPar
PULSAR el botón \sphinxstylestrong{Aceptar.}

\end{itemize}

\item {} 
\sphinxAtStartPar
El \sphinxstylestrong{formato de carácter} del título Gota debería quedar así:


\begin{savenotes}\sphinxattablestart
\sphinxthistablewithglobalstyle
\centering
\begin{tabulary}{\linewidth}[t]{T}
\sphinxtoprule
\sphinxtableatstartofbodyhook

\begin{sphinxfigure-in-table}
\centering

\noindent\sphinxincludegraphics[width=200\sphinxpxdimen]{{image095}.png}
\end{sphinxfigure-in-table}\relax
\\
\sphinxbottomrule
\end{tabulary}
\sphinxtableafterendhook\par
\sphinxattableend\end{savenotes}

\item {} 
\sphinxAtStartPar
Al párrafo de la descripción de la enfermedad \sphinxstylestrong{Gota}, APLICAR: \sphinxstylestrong{tipo de letra:} Tahoma y el \sphinxstylestrong{tamaño de letra:} 14pt (En este caso puede usar el formato de la \sphinxstylestrong{barra de herramientas formato})

\item {} 
\sphinxAtStartPar
SELECCIONE el título \sphinxstylestrong{Gota} y en la \sphinxstylestrong{barra de herramientas estándar}, PULSE el botón \sphinxstylestrong{Clonar formato} dos veces para dejar la herramienta activa.


\begin{savenotes}\sphinxattablestart
\sphinxthistablewithglobalstyle
\centering
\begin{tabulary}{\linewidth}[t]{T}
\sphinxtoprule
\sphinxtableatstartofbodyhook

\begin{sphinxfigure-in-table}
\centering

\noindent\sphinxincludegraphics[width=30\sphinxpxdimen]{{image104}.png}
\end{sphinxfigure-in-table}\relax
\\
\sphinxbottomrule
\end{tabulary}
\sphinxtableafterendhook\par
\sphinxattableend\end{savenotes}

\begin{sphinxadmonition}{note}{¿Qué esta pasando en el mouse?}

\sphinxAtStartPar
El \sphinxstylestrong{mouse} ahora tendrá un \sphinxstylestrong{ícono} de un \sphinxstylestrong{bote de pintura derramándose}, esto quiere decir que cualquier palabra que seleccione, recibirá el formato previamente seleccionado.

\sphinxAtStartPar
Cuando termine de clonar formato, presione la tecla \sphinxkeyboard{\sphinxupquote{ESC}}.
\end{sphinxadmonition}

\item {} 
\sphinxAtStartPar
COPIE el formato a los títulos de las enfermedades: \sphinxstylestrong{Artrosis} y \sphinxstylestrong{Aftas}. Cuando termine presione la tecla \sphinxkeyboard{\sphinxupquote{ESC}}.

\item {} 
\sphinxAtStartPar
En su cuaderno:

\begin{sphinxadmonition}{note}{En su cuaderno…}

\sphinxAtStartPar
Responda la siguiente pregunta:
\begin{itemize}
\item {} 
\sphinxAtStartPar
¿Cómo la herramienta \sphinxstylestrong{clonar formato} puede ayudar a superar las desventajas al usar la herramienta: \sphinxstylestrong{formato de carácter?}

\end{itemize}
\end{sphinxadmonition}

\item {} 
\sphinxAtStartPar
Guarde los cambios y cierre el software \sphinxstylestrong{LibreOffice Writer}.

\end{enumerate}


\bigskip\hrule\bigskip



\subsection{A07: Formato de párrafo \sphinxstyleliteralintitle{\sphinxupquote{103}} (+0,1)}
\label{\detokenize{1_guias/10ofimatica/G1014:a07-formato-de-parrafo-103-0-1}}
\sphinxAtStartPar
\DUrole{sd-sphinx-override}{\DUrole{sd-badge}{\DUrole{sd-bg-warning}{\DUrole{sd-bg-text-warning}{Actividad práctica}}}}
\begin{enumerate}
\sphinxsetlistlabels{\arabic}{enumi}{enumii}{}{.}%
\item {} 
\sphinxAtStartPar
DESCARGUE el siguiente archivo y guárdelo en su carpeta de evidencias:
\begin{itemize}
\item {} 
\sphinxAtStartPar
G1014A07\sphinxhyphen{}nueces.odt \DUrole{xref}{\DUrole{download}{\DUrole{myst}{\DUrole{sd-sphinx-override}{\DUrole{sd-badge}{\DUrole{sd-bg-primary}{\DUrole{sd-bg-text-primary}{Descargar}}}}}}}

\end{itemize}

\item {} 
\sphinxAtStartPar
Abra el archivo: \sphinxstylestrong{G1014A07\sphinxhyphen{}nueces.odt} y en la cabecera del documento, agregue la \sphinxstylestrong{información de identificación} con el título: \sphinxstylestrong{FORMATO DE PÁRRAFO}

\begin{sphinxadmonition}{note}{¿Qué contiene el documento?}

\sphinxAtStartPar
En este documento encontrará la descripción de algunas vitaminas presentes en las nueces. Por medio de la herramienta formato de párrafo el documento tendrá una mejor apariencia para el lector.
\end{sphinxadmonition}

\item {} 
\sphinxAtStartPar
En la \sphinxstylestrong{barra de herramientas estándar}, PULSE el botón \sphinxstylestrong{Alternar marcas de formato}.


\begin{savenotes}\sphinxattablestart
\sphinxthistablewithglobalstyle
\centering
\begin{tabulary}{\linewidth}[t]{T}
\sphinxtoprule
\sphinxtableatstartofbodyhook

\begin{sphinxfigure-in-table}
\centering

\noindent\sphinxincludegraphics[width=50\sphinxpxdimen]{{image113}.png}
\end{sphinxfigure-in-table}\relax
\\
\sphinxbottomrule
\end{tabulary}
\sphinxtableafterendhook\par
\sphinxattableend\end{savenotes}

\item {} 
\sphinxAtStartPar
Una vez activó las marcas de formato, en su cuaderno describa que acabó de pasar en el documento.

\begin{sphinxadmonition}{note}{¿Qué significan las marcas de formato?}
\begin{itemize}
\item {} 
\sphinxAtStartPar
Una \sphinxstylestrong{marca} (¶) significa un párrafo.

\item {} 
\sphinxAtStartPar
Un \sphinxstylestrong{punto} (.) significa un espacio.

\item {} 
\sphinxAtStartPar
Una \sphinxstylestrong{flecha} (→) significa una tabulación.

\end{itemize}
\end{sphinxadmonition}

\item {} 
\sphinxAtStartPar
En su cuaderno:

\begin{sphinxadmonition}{note}{En su cuaderno…}

\sphinxAtStartPar
Responda la siguiente pregunta:
\begin{itemize}
\item {} 
\sphinxAtStartPar
¿Cuántas marcas, puntos y flechas salieron en el documento?

\end{itemize}
\end{sphinxadmonition}

\item {} 
\sphinxAtStartPar
SITUE el cursor al final de la segunda marca (segundo párrafo) y en el \sphinxstylestrong{menú superior}, PULSE \sphinxguilabel{Formato} > \sphinxguilabel{Párrafo…} para abrir la ventana \sphinxstylestrong{Párrafo}

\item {} 
\sphinxAtStartPar
En la ventana \sphinxstylestrong{Párrafo} configure las siguientes opciones:


\begin{savenotes}\sphinxattablestart
\sphinxthistablewithglobalstyle
\centering
\begin{tabulary}{\linewidth}[t]{T}
\sphinxtoprule
\sphinxtableatstartofbodyhook

\begin{sphinxfigure-in-table}
\centering

\noindent\sphinxincludegraphics[width=650\sphinxpxdimen]{{image123}.png}
\end{sphinxfigure-in-table}\relax
\\
\sphinxbottomrule
\end{tabulary}
\sphinxtableafterendhook\par
\sphinxattableend\end{savenotes}
\begin{itemize}
\item {} 
\sphinxAtStartPar
En la pestaña \sphinxguilabel{Sangrías y espaciado} configurar:
\begin{itemize}
\item {} 
\sphinxAtStartPar
En la sección \sphinxstylestrong{Espaciado} configurar:
\begin{itemize}
\item {} 
\sphinxAtStartPar
Bajo el párrafo: 0,30 cm

\end{itemize}

\item {} 
\sphinxAtStartPar
En la sección \sphinxstylestrong{Interlineado} configurar:
\begin{itemize}
\item {} 
\sphinxAtStartPar
\sphinxstylestrong{Proporcional} de \sphinxstylestrong{120\%}

\end{itemize}

\end{itemize}

\item {} 
\sphinxAtStartPar
PULSE el botón \sphinxstylestrong{Aceptar}

\end{itemize}

\item {} 
\sphinxAtStartPar
En su cuaderno:

\begin{sphinxadmonition}{note}{En su cuaderno…}

\sphinxAtStartPar
Describa que acabó de pasar en el documento.
\end{sphinxadmonition}

\item {} 
\sphinxAtStartPar
SELECCIONE desde el principio del segundo párrafo hasta el final y aplique la siguiente configuración de \sphinxstylestrong{formato de párrafo}:
\begin{itemize}
\item {} 
\sphinxAtStartPar
En la pestaña \sphinxguilabel{Sangrías y espaciado} configurar:
\begin{itemize}
\item {} 
\sphinxAtStartPar
En la sección \sphinxstylestrong{Espaciado} configurar:
\begin{itemize}
\item {} 
\sphinxAtStartPar
Bajo el párrafo: 0,30 cm

\end{itemize}

\item {} 
\sphinxAtStartPar
En la sección \sphinxstylestrong{Interlineado} configurar:
\begin{itemize}
\item {} 
\sphinxAtStartPar
\sphinxstylestrong{Proporcional} de \sphinxstylestrong{120\%}

\end{itemize}

\end{itemize}

\item {} 
\sphinxAtStartPar
PULSE el botón \sphinxstylestrong{Aceptar}

\end{itemize}

\item {} 
\sphinxAtStartPar
SELECCIONE el título “\sphinxstylestrong{Las nueces contienen”} y aplique la siguiente configuración de \sphinxstylestrong{formato de párrafo}:
\begin{itemize}
\item {} 
\sphinxAtStartPar
En la pestaña \sphinxguilabel{Alineación} configurar:
\begin{itemize}
\item {} 
\sphinxAtStartPar
En la sección \sphinxstylestrong{Opciones} seleccionar \sphinxstylestrong{Centro.}

\end{itemize}

\item {} 
\sphinxAtStartPar
En la pestaña \sphinxguilabel{Área} configurar:
\begin{itemize}
\item {} 
\sphinxAtStartPar
\sphinxstylestrong{Color:} Amarillo

\end{itemize}

\item {} 
\sphinxAtStartPar
PULSE el botón \sphinxstylestrong{Aceptar}.

\end{itemize}

\item {} 
\sphinxAtStartPar
SELECCIONE el título “\sphinxstylestrong{Las nueces contienen”} y aplique la siguiente configuración de \sphinxstylestrong{formato de carácter}:
\begin{itemize}
\item {} 
\sphinxAtStartPar
En la pestaña \sphinxguilabel{Tipo de letra} configurar:
\begin{itemize}
\item {} 
\sphinxAtStartPar
\sphinxstylestrong{Familia:} Liberation Sans

\item {} 
\sphinxAtStartPar
\sphinxstylestrong{Tamaño:} 16 pt

\item {} 
\sphinxAtStartPar
\sphinxstylestrong{Estilo:} Negrita

\end{itemize}

\item {} 
\sphinxAtStartPar
En la pestaña \sphinxguilabel{Efectos tipográficos} configurar:
\begin{itemize}
\item {} 
\sphinxAtStartPar
En la sección \sphinxstylestrong{Color de letra} configurar:
\begin{itemize}
\item {} 
\sphinxAtStartPar
\sphinxstylestrong{Color de letra:} Púrpura

\end{itemize}

\item {} 
\sphinxAtStartPar
En la sección \sphinxstylestrong{Decoración de texto} configurar:
\begin{itemize}
\item {} 
\sphinxAtStartPar
\sphinxstylestrong{Subrayado:} Sencillo

\item {} 
\sphinxAtStartPar
\sphinxstylestrong{Color:} Púrpura

\end{itemize}

\item {} 
\sphinxAtStartPar
En la sección \sphinxstylestrong{Efectos} configurar:
\begin{itemize}
\item {} 
\sphinxAtStartPar
Marcar la opción Sombra.

\end{itemize}

\end{itemize}

\item {} 
\sphinxAtStartPar
PULSE el botón \sphinxstylestrong{Aceptar}.

\end{itemize}

\item {} 
\sphinxAtStartPar
SELECCIONE desde el principio del segundo párrafo hasta el final y aplique la siguiente configuración de \sphinxstylestrong{formato de carácter}:
\begin{itemize}
\item {} 
\sphinxAtStartPar
En la pestaña \sphinxguilabel{Efectos tipográficos} configurar:
\begin{itemize}
\item {} 
\sphinxAtStartPar
En la sección \sphinxstylestrong{Color de letra} configurar:
\begin{itemize}
\item {} 
\sphinxAtStartPar
\sphinxstylestrong{Color de letra:} Gris oscuro 2

\end{itemize}

\end{itemize}

\item {} 
\sphinxAtStartPar
PULSAR el botón \sphinxstylestrong{Aceptar}.

\end{itemize}

\item {} 
\sphinxAtStartPar
SELECCIONE desde el principio del segundo párrafo al penúltimo párrafo (no seleccione el último párrafo) y en el \sphinxstylestrong{menú superior} PULSE \sphinxguilabel{Formato} > \sphinxguilabel{Numeración y bolos} para abrir la ventana de \sphinxstylestrong{Numeración y bolos} y configure las siguientes opciones:


\begin{savenotes}\sphinxattablestart
\sphinxthistablewithglobalstyle
\centering
\begin{tabulary}{\linewidth}[t]{T}
\sphinxtoprule
\sphinxtableatstartofbodyhook

\begin{sphinxfigure-in-table}
\centering

\noindent\sphinxincludegraphics[width=650\sphinxpxdimen]{{image133}.png}
\end{sphinxfigure-in-table}\relax
\\
\sphinxbottomrule
\end{tabulary}
\sphinxtableafterendhook\par
\sphinxattableend\end{savenotes}
\begin{itemize}
\item {} 
\sphinxAtStartPar
En la \sphinxguilabel{pestaña No ordenada} configurar:
\begin{itemize}
\item {} 
\sphinxAtStartPar
SELECCIONAR cualquier modelo.

\end{itemize}

\item {} 
\sphinxAtStartPar
PULSE el botón \sphinxstylestrong{Aceptar}.

\end{itemize}

\item {} 
\sphinxAtStartPar
En el documento SELECCIONE el último párrafo y aplique la siguiente configuración de \sphinxstylestrong{Formato de párrafo}:
\begin{itemize}
\item {} 
\sphinxAtStartPar
En la pestaña de \sphinxguilabel{Sangría y espaciado} configure:
\begin{itemize}
\item {} 
\sphinxAtStartPar
En la sección \sphinxstylestrong{Sangría} configure:
\begin{itemize}
\item {} 
\sphinxAtStartPar
Antes del texto: 4cm

\item {} 
\sphinxAtStartPar
Después del texto: 4cm

\end{itemize}

\item {} 
\sphinxAtStartPar
En la sección \sphinxstylestrong{Espaciado} configure:
\begin{itemize}
\item {} 
\sphinxAtStartPar
Sobre el párrafo: 2cm

\end{itemize}

\end{itemize}

\item {} 
\sphinxAtStartPar
En la pestaña \sphinxguilabel{Alineación} configure:
\begin{itemize}
\item {} 
\sphinxAtStartPar
En la sección \sphinxstylestrong{Opciones} configure \sphinxstylestrong{Centro}.

\end{itemize}

\item {} 
\sphinxAtStartPar
En la pestaña \sphinxguilabel{Bordes} configure:
\begin{itemize}
\item {} 
\sphinxAtStartPar
En las sección \sphinxstylestrong{Disposición de líneas} configure:
\begin{itemize}
\item {} 
\sphinxAtStartPar
\sphinxstylestrong{Preconfigs:} Los cuatro bordes

\end{itemize}

\item {} 
\sphinxAtStartPar
En la sección \sphinxstylestrong{Separación} configure:
\begin{itemize}
\item {} 
\sphinxAtStartPar
0,20 cm en todos los lados

\end{itemize}

\end{itemize}

\item {} 
\sphinxAtStartPar
En la pestaña \sphinxguilabel{Área} configurar:
\begin{itemize}
\item {} 
\sphinxAtStartPar
\sphinxstylestrong{Color:} Gris claro 4.

\end{itemize}

\item {} 
\sphinxAtStartPar
PULSE el botón \sphinxstylestrong{Aceptar}.

\end{itemize}

\item {} 
\sphinxAtStartPar
Después de aplicar todas las configuraciones, el documento debería tener la siguiente apariencia:


\begin{savenotes}\sphinxattablestart
\sphinxthistablewithglobalstyle
\centering
\begin{tabulary}{\linewidth}[t]{T}
\sphinxtoprule
\sphinxtableatstartofbodyhook

\begin{sphinxfigure-in-table}
\centering

\noindent\sphinxincludegraphics[width=600\sphinxpxdimen]{{image143}.png}
\end{sphinxfigure-in-table}\relax
\\
\sphinxbottomrule
\end{tabulary}
\sphinxtableafterendhook\par
\sphinxattableend\end{savenotes}

\item {} 
\sphinxAtStartPar
En su cuaderno:

\begin{sphinxadmonition}{note}{En su cuaderno…}

\sphinxAtStartPar
Responda la siguiente pregunta:
\begin{itemize}
\item {} 
\sphinxAtStartPar
¿Cuantas veces y en que pasos se usaron el \sphinxstylestrong{formato de carácter}, \sphinxstylestrong{formato de párrafo} y \sphinxstylestrong{la numeración y bolos}?

\end{itemize}
\end{sphinxadmonition}

\item {} 
\sphinxAtStartPar
Guarde los cambios y cierre el software LibreOffice Writer.

\end{enumerate}


\bigskip\hrule\bigskip



\subsection{A08: Sangría Francesa \sphinxstyleliteralintitle{\sphinxupquote{101}} (+0,1)}
\label{\detokenize{1_guias/10ofimatica/G1014:a08-sangria-francesa-101-0-1}}
\sphinxAtStartPar
\DUrole{sd-sphinx-override}{\DUrole{sd-badge}{\DUrole{sd-bg-warning}{\DUrole{sd-bg-text-warning}{Actividad práctica}}}}
\begin{enumerate}
\sphinxsetlistlabels{\arabic}{enumi}{enumii}{}{.}%
\item {} 
\sphinxAtStartPar
DESCARGUE el siguiente archivo y guárdelo en su carpeta de evidencias:
\begin{itemize}
\item {} 
\sphinxAtStartPar
G1014A08\sphinxhyphen{}certificado.odt \DUrole{xref}{\DUrole{download}{\DUrole{myst}{\DUrole{sd-sphinx-override}{\DUrole{sd-badge}{\DUrole{sd-bg-primary}{\DUrole{sd-bg-text-primary}{Descargar}}}}}}}

\end{itemize}

\item {} 
\sphinxAtStartPar
Abra el archivo: \sphinxstylestrong{G1014A08\sphinxhyphen{}certificado.odt} y en la cabecera del documento, agregue la \sphinxstylestrong{información de identificación} con el título: \sphinxstylestrong{SANGRÍA FRANCESA}

\begin{sphinxadmonition}{note}{¿Qué contiene el documento?}

\sphinxAtStartPar
El documento muestra información relacionada con un certificado laboral.

\sphinxAtStartPar
Se debe configurar el documento para que el certificado tenga una mejor presentación.
\end{sphinxadmonition}

\item {} 
\sphinxAtStartPar
La palabra “Certifica” debe estar en \sphinxstylestrong{negrita} (usar la \sphinxstylestrong{barra de herramientas formato}).

\item {} 
\sphinxAtStartPar
SELECCIONE todo el documento y aplique la siguiente configuración de \sphinxstylestrong{formato de párrafo:}
\begin{itemize}
\item {} 
\sphinxAtStartPar
En la pestaña \sphinxguilabel{Sangrías y espaciado}:
\begin{itemize}
\item {} 
\sphinxAtStartPar
En la sección \sphinxstylestrong{Espaciado} configurar:
\begin{itemize}
\item {} 
\sphinxAtStartPar
Bajo el párrafo: 0,40 cm

\end{itemize}

\end{itemize}

\item {} 
\sphinxAtStartPar
PULSE el botón \sphinxstylestrong{Aceptar}.

\end{itemize}

\item {} 
\sphinxAtStartPar
Al tercer párrafo APLIQUE el siguiente \sphinxstylestrong{formato de párrafo}:
\begin{itemize}
\item {} 
\sphinxAtStartPar
En la pestaña \sphinxguilabel{Sangría y espaciado} configurar:
\begin{itemize}
\item {} 
\sphinxAtStartPar
En la sección \sphinxstylestrong{Sangría} configurar:
\begin{itemize}
\item {} 
\sphinxAtStartPar
\sphinxstylestrong{Antes del texto:} 3,00 cm

\item {} 
\sphinxAtStartPar
\sphinxstylestrong{Primer renglón:} \sphinxhyphen{}3,00 cm (Esto es la sangría francesa)

\end{itemize}

\end{itemize}

\item {} 
\sphinxAtStartPar
PULSE el botón \sphinxstylestrong{Aceptar}.

\end{itemize}

\item {} 
\sphinxAtStartPar
Después de aplicar todas las configuraciones, el documento debería tener la siguiente apariencia:


\begin{savenotes}\sphinxattablestart
\sphinxthistablewithglobalstyle
\centering
\begin{tabulary}{\linewidth}[t]{T}
\sphinxtoprule
\sphinxtableatstartofbodyhook

\begin{sphinxfigure-in-table}
\centering

\noindent\sphinxincludegraphics[width=720\sphinxpxdimen]{{image152}.png}
\end{sphinxfigure-in-table}\relax
\\
\sphinxbottomrule
\end{tabulary}
\sphinxtableafterendhook\par
\sphinxattableend\end{savenotes}

\item {} 
\sphinxAtStartPar
En su cuaderno:

\begin{sphinxadmonition}{note}{En su cuaderno…}

\sphinxAtStartPar
Responda la siguiente pregunta:
\begin{itemize}
\item {} 
\sphinxAtStartPar
¿Qué es una \sphinxstylestrong{sangría francesa} y explique en que otros tipos de documentos podría usarse?

\end{itemize}
\end{sphinxadmonition}

\item {} 
\sphinxAtStartPar
GUARDE los cambios y cierre el software Libreoffice Writer.

\end{enumerate}


\bigskip\hrule\bigskip



\subsection{A09: Esquema numerado \sphinxstyleliteralintitle{\sphinxupquote{102}} (+0,1)}
\label{\detokenize{1_guias/10ofimatica/G1014:a09-esquema-numerado-102-0-1}}
\sphinxAtStartPar
\DUrole{sd-sphinx-override}{\DUrole{sd-badge}{\DUrole{sd-bg-warning}{\DUrole{sd-bg-text-warning}{Actividad práctica}}}}
\begin{enumerate}
\sphinxsetlistlabels{\arabic}{enumi}{enumii}{}{.}%
\item {} 
\sphinxAtStartPar
DESCARGUE el siguiente archivo y guárdelo en su carpeta de evidencias:
\begin{itemize}
\item {} 
\sphinxAtStartPar
G1014A09\sphinxhyphen{}temario.odt \DUrole{xref}{\DUrole{download}{\DUrole{myst}{\DUrole{sd-sphinx-override}{\DUrole{sd-badge}{\DUrole{sd-bg-primary}{\DUrole{sd-bg-text-primary}{Descargar}}}}}}}

\end{itemize}

\item {} 
\sphinxAtStartPar
Abra el archivo: \sphinxstylestrong{G1014A09\sphinxhyphen{}temario.odt} y en la cabecera del documento, agregue la \sphinxstylestrong{información de identificación} con el título: \sphinxstylestrong{ESQUEMA NUMERADO}

\begin{sphinxadmonition}{note}{¿Qué contiene el documento?}

\sphinxAtStartPar
El documento muestra un contenido el cual no esta numerado. Es necesario aplicar esquemas numerados.
\end{sphinxadmonition}

\item {} 
\sphinxAtStartPar
En la \sphinxstylestrong{barra de herramientas estándar}, PULSE: \sphinxguilabel{Alternar marcas de formato}.

\item {} 
\sphinxAtStartPar
SELECCIONE todo el documento y aplique la siguiente configuración de \sphinxstylestrong{formato de párrafo}:
\begin{itemize}
\item {} 
\sphinxAtStartPar
En la pestaña de \sphinxguilabel{Sangrías y espaciado} configure:
\begin{itemize}
\item {} 
\sphinxAtStartPar
En la sección de \sphinxstylestrong{Espaciado} configure:
\begin{itemize}
\item {} 
\sphinxAtStartPar
Bajo el párrafo: 0.30 cm

\end{itemize}

\end{itemize}

\item {} 
\sphinxAtStartPar
PULSE el botón \sphinxstylestrong{Aceptar}.

\end{itemize}

\item {} 
\sphinxAtStartPar
SELECCIONE todo el documento y en el \sphinxstylestrong{menú superior,} PULSE \sphinxguilabel{Formato} > \sphinxguilabel{Numeración y bolos}.
\begin{itemize}
\item {} 
\sphinxAtStartPar
En la pestaña de \sphinxguilabel{Esquema} seleccione el siguiente:

\end{itemize}


\begin{savenotes}\sphinxattablestart
\sphinxthistablewithglobalstyle
\centering
\begin{tabulary}{\linewidth}[t]{T}
\sphinxtoprule
\sphinxtableatstartofbodyhook

\begin{sphinxfigure-in-table}
\centering

\noindent\sphinxincludegraphics[width=180\sphinxpxdimen]{{image162}.png}
\end{sphinxfigure-in-table}\relax
\\
\sphinxbottomrule
\end{tabulary}
\sphinxtableafterendhook\par
\sphinxattableend\end{savenotes}
\begin{itemize}
\item {} 
\sphinxAtStartPar
PULSE el botón \sphinxstylestrong{Aceptar}.

\end{itemize}

\item {} 
\sphinxAtStartPar
SITUE el cursor al comienzo del párrafo \sphinxstylestrong{2. Generalidades} y PULSE la tecla \sphinxkeyboard{\sphinxupquote{TAB}}.

\item {} 
\sphinxAtStartPar
En su cuaderno:

\begin{sphinxadmonition}{note}{En su cuaderno…}

\sphinxAtStartPar
Responda la siguiente pregunta:
\begin{itemize}
\item {} 
\sphinxAtStartPar
¿Que pasó en el documento cuanto presionó la tecla \sphinxkeyboard{\sphinxupquote{TAB}}?

\end{itemize}
\end{sphinxadmonition}

\item {} 
\sphinxAtStartPar
El título \sphinxstylestrong{Generalidades} aplique \sphinxstylestrong{Tipo de letra:} negrita

\item {} 
\sphinxAtStartPar
APLIQUE \sphinxkeyboard{\sphinxupquote{TAB}} y \sphinxstylestrong{negrita} a los siguientes títulos:
\begin{itemize}
\item {} 
\sphinxAtStartPar
Control de comunicaciones de la red.

\item {} 
\sphinxAtStartPar
Polling

\item {} 
\sphinxAtStartPar
Paso de testigo

\item {} 
\sphinxAtStartPar
Normas estándares para redes locales

\end{itemize}

\item {} 
\sphinxAtStartPar
APLIQUE 2 veces \sphinxkeyboard{\sphinxupquote{TAB}} y \sphinxstylestrong{negrita} a los siguientes títulos:
\begin{itemize}
\item {} 
\sphinxAtStartPar
Concepto de red local

\item {} 
\sphinxAtStartPar
Ventajas de redes locales

\item {} 
\sphinxAtStartPar
Gestores de una red local

\item {} 
\sphinxAtStartPar
Introducción

\item {} 
\sphinxAtStartPar
Protocolos

\item {} 
\sphinxAtStartPar
Protocolos de contienda

\item {} 
\sphinxAtStartPar
Contienda simple

\item {} 
\sphinxAtStartPar
Factores de evaluación del protocolo de paso de testigo

\end{itemize}

\item {} 
\sphinxAtStartPar
APLIQUE 3 veces \sphinxkeyboard{\sphinxupquote{TAB}} a los títulos que no estén en \sphinxstylestrong{negrita}.

\item {} 
\sphinxAtStartPar
Después de aplicar todas las configuraciones, el documento debería tener la siguiente apariencia:


\begin{savenotes}\sphinxattablestart
\sphinxthistablewithglobalstyle
\centering
\begin{tabulary}{\linewidth}[t]{T}
\sphinxtoprule
\sphinxtableatstartofbodyhook

\begin{sphinxfigure-in-table}
\centering

\noindent\sphinxincludegraphics[width=400\sphinxpxdimen]{{image172}.png}
\end{sphinxfigure-in-table}\relax
\\
\sphinxbottomrule
\end{tabulary}
\sphinxtableafterendhook\par
\sphinxattableend\end{savenotes}

\item {} 
\sphinxAtStartPar
En su cuaderno:

\begin{sphinxadmonition}{note}{En su cuaderno…}

\sphinxAtStartPar
Responda la siguiente pregunta:
\begin{itemize}
\item {} 
\sphinxAtStartPar
¿Cuales son los títulos de primer, segundo, tercer y cuarto nivel?

\end{itemize}
\end{sphinxadmonition}

\item {} 
\sphinxAtStartPar
GUARDE los cambios y cierre el software LibreOffice Writer.

\end{enumerate}


\bigskip\hrule\bigskip



\subsection{A10: Currículum de estilos \sphinxstyleliteralintitle{\sphinxupquote{101}} (+0,1)}
\label{\detokenize{1_guias/10ofimatica/G1014:a10-curriculum-de-estilos-101-0-1}}\begin{enumerate}
\sphinxsetlistlabels{\arabic}{enumi}{enumii}{}{.}%
\item {} 
\sphinxAtStartPar
DESCARGUE el siguiente archivo y guárdelo en su carpeta de evidencias:
\begin{itemize}
\item {} 
\sphinxAtStartPar
G1014A10\sphinxhyphen{}curriculum.odt \DUrole{xref}{\DUrole{download}{\DUrole{myst}{\DUrole{sd-sphinx-override}{\DUrole{sd-badge}{\DUrole{sd-bg-primary}{\DUrole{sd-bg-text-primary}{Descargar}}}}}}}

\end{itemize}

\item {} 
\sphinxAtStartPar
Abra el archivo: \sphinxstylestrong{G1014A10\sphinxhyphen{}curriculum.odt} y en la cabecera del documento, agregue la \sphinxstylestrong{información de identificación} con el título: \sphinxstylestrong{FORMATO DE ESTILOS}

\begin{sphinxadmonition}{note}{¿Qué contiene el documento?}

\sphinxAtStartPar
El documento muestra la hoja de vida de una persona. Sin embargo no es entendible.

\sphinxAtStartPar
Es necesario aplicar estilos para mejorar la apariencia del documento y así sea más entendible.
\end{sphinxadmonition}

\item {} 
\sphinxAtStartPar
SELECCIONE los títulos: \sphinxstylestrong{Formación Reglada, Formación No Reglada y Experiencia Laboral}.

\begin{sphinxadmonition}{note}{¿Cómo seleccionar varios títulos al mismo tiempo?}

\sphinxAtStartPar
Para seleccionar varios títulos al mismo tiempo, mantenga PULSADA la tecla \sphinxkeyboard{\sphinxupquote{CTRL}} y selecciona los títulos uno por uno.
\end{sphinxadmonition}


\begin{savenotes}\sphinxattablestart
\sphinxthistablewithglobalstyle
\centering
\begin{tabulary}{\linewidth}[t]{T}
\sphinxtoprule
\sphinxtableatstartofbodyhook

\begin{sphinxfigure-in-table}
\centering

\noindent\sphinxincludegraphics[width=620\sphinxpxdimen]{{image182}.png}
\end{sphinxfigure-in-table}\relax
\\
\sphinxbottomrule
\end{tabulary}
\sphinxtableafterendhook\par
\sphinxattableend\end{savenotes}

\item {} 
\sphinxAtStartPar
En la \sphinxstylestrong{barra de herramientas lateral}, PULSE el ícono Estilos.


\begin{savenotes}\sphinxattablestart
\sphinxthistablewithglobalstyle
\centering
\begin{tabulary}{\linewidth}[t]{T}
\sphinxtoprule
\sphinxtableatstartofbodyhook

\begin{sphinxfigure-in-table}
\centering

\noindent\sphinxincludegraphics[width=89\sphinxpxdimen]{{image191}.png}
\end{sphinxfigure-in-table}\relax
\\
\sphinxbottomrule
\end{tabulary}
\sphinxtableafterendhook\par
\sphinxattableend\end{savenotes}

\item {} 
\sphinxAtStartPar
En la \sphinxstyleemphasis{herramienta de estilos} encontrará 6 herramientas para aplicar formato de estilo al documento. PULSE el \sphinxstylestrong{icono de Estilo de párrafos}.


\begin{savenotes}\sphinxattablestart
\sphinxthistablewithglobalstyle
\centering
\begin{tabulary}{\linewidth}[t]{T}
\sphinxtoprule
\sphinxtableatstartofbodyhook

\begin{sphinxfigure-in-table}
\centering

\noindent\sphinxincludegraphics[width=180\sphinxpxdimen]{{image202}.png}
\end{sphinxfigure-in-table}\relax
\\
\sphinxbottomrule
\end{tabulary}
\sphinxtableafterendhook\par
\sphinxattableend\end{savenotes}

\item {} 
\sphinxAtStartPar
Observe la \sphinxstyleemphasis{Jerarquía de estilos de párrafo} y localice el estilo de nombre: \sphinxstylestrong{Título}.

\begin{sphinxadmonition}{note}{¿Qué es el estilo de título?}

\sphinxAtStartPar
El estilo \sphinxstylestrong{Título}, a su vez tiene mas \sphinxstylestrong{sub\sphinxhyphen{}estilos}. Si no aparecen, PULSAR clic sobre la flecha apuntando hacia abajo a la izquierda del nombre del estilo.
\end{sphinxadmonition}


\begin{savenotes}\sphinxattablestart
\sphinxthistablewithglobalstyle
\centering
\begin{tabulary}{\linewidth}[t]{T}
\sphinxtoprule
\sphinxtableatstartofbodyhook

\begin{sphinxfigure-in-table}
\centering

\noindent\sphinxincludegraphics[width=300\sphinxpxdimen]{{image212}.png}
\end{sphinxfigure-in-table}\relax
\\
\sphinxbottomrule
\end{tabulary}
\sphinxtableafterendhook\par
\sphinxattableend\end{savenotes}

\item {} 
\sphinxAtStartPar
ANTES DE CONTINUAR, ASEGURESE que los títulos del documento previamente seleccionados siguen estando seleccionados!!!

\item {} 
\sphinxAtStartPar
PULSE \sphinxstylestrong{doble clic} sobre el sub\sphinxhyphen{}estilo \sphinxstylestrong{Título 1.} (\sphinxstyleemphasis{Usted acaba de aplicar un estilo en el documento}). \sphinxstyleemphasis{\sphinxstylestrong{Responda las siguientes preguntas en su cuaderno:}}
\begin{itemize}
\item {} 
\sphinxAtStartPar
¿Qué acabo de pasar en el documento cuando aplicó el estilo?

\item {} 
\sphinxAtStartPar
¿Qué es mejor: aplicar \sphinxstylestrong{formato de párrafo} o aplicar \sphinxstylestrong{formato de estilo} y por qué?

\end{itemize}

\item {} 
\sphinxAtStartPar
SELECCIONE todos los párrafos del documento (NO SELECCIONE los títulos) y en la \sphinxstylestrong{Jerarquía de estilos de párrafo}, APLIQUE el estilo: \sphinxstylestrong{Cuerpo de texto}.

\item {} 
\sphinxAtStartPar
En la herramienta: \sphinxstylestrong{Jerarquía de estilos de párrafo}, SELECCIONE nuevamente el estilo \sphinxstylestrong{Cuerpo de texto} y PULSE el botón derecho del mouse para abrir un \sphinxstylestrong{menú contextual}. En el \sphinxstylestrong{menú contextual} de \sphinxstylestrong{Cuerpo de texto}, seleccione la opción: \sphinxstylestrong{Editar estilo…}

\begin{sphinxadmonition}{note}{¿Qué es el menú contextual?}

\sphinxAtStartPar
Un \sphinxstylestrong{menú contextual} es una lista de opciones que aparece cuando PULSA clic derecho (o un gesto similar) en un elemento en la pantalla de una computadora o dispositivo.

\sphinxAtStartPar
Esto permite realizar acciones específicas relacionadas con ese elemento, como copiar, pegar, eliminar, o acceder a propiedades.
\end{sphinxadmonition}


\begin{savenotes}\sphinxattablestart
\sphinxthistablewithglobalstyle
\centering
\begin{tabulary}{\linewidth}[t]{T}
\sphinxtoprule
\sphinxtableatstartofbodyhook

\begin{sphinxfigure-in-table}
\centering

\noindent\sphinxincludegraphics[width=250\sphinxpxdimen]{{image22}.png}
\end{sphinxfigure-in-table}\relax
\\
\sphinxbottomrule
\end{tabulary}
\sphinxtableafterendhook\par
\sphinxattableend\end{savenotes}

\item {} 
\sphinxAtStartPar
En la ventana \sphinxguilabel{Estilo de párrafo: Cuerpo de texto}, configurar las siguientes opciones:


\begin{savenotes}\sphinxattablestart
\sphinxthistablewithglobalstyle
\centering
\begin{tabulary}{\linewidth}[t]{T}
\sphinxtoprule
\sphinxtableatstartofbodyhook

\begin{sphinxfigure-in-table}
\centering

\noindent\sphinxincludegraphics[width=630\sphinxpxdimen]{{image23}.png}
\end{sphinxfigure-in-table}\relax
\\
\sphinxbottomrule
\end{tabulary}
\sphinxtableafterendhook\par
\sphinxattableend\end{savenotes}
\begin{itemize}
\item {} 
\sphinxAtStartPar
En la pestaña \sphinxguilabel{Tipo de letra} configurar:
\begin{itemize}
\item {} 
\sphinxAtStartPar
\sphinxstylestrong{Familia:} Times New Roman

\item {} 
\sphinxAtStartPar
\sphinxstylestrong{Tamaño:} 14pt

\end{itemize}

\item {} 
\sphinxAtStartPar
En la pestaña \sphinxguilabel{Esquema y lista} configurar:
\begin{itemize}
\item {} 
\sphinxAtStartPar
\sphinxstylestrong{Nivel de esquema:} Nivel 1

\item {} 
\sphinxAtStartPar
\sphinxstylestrong{Estilo de esquema:} Bolo \sphinxhyphen{}

\end{itemize}

\item {} 
\sphinxAtStartPar
PULSAR el botón \sphinxstylestrong{Aceptar}.

\end{itemize}

\item {} 
\sphinxAtStartPar
En la \sphinxstylestrong{Jerarquía de estilos de párrafo}, IDENTIFIQUE el estilo: \sphinxstylestrong{Título 1} y en su menú contextual PULSE \sphinxstylestrong{Editar estilo…}

\item {} 
\sphinxAtStartPar
En la ventana: \sphinxstylestrong{Estilo de párrafo: Título 1} configure lo siguiente:
\begin{itemize}
\item {} 
\sphinxAtStartPar
En la pestaña \sphinxguilabel{Tipo de letra} configurar:
\begin{itemize}
\item {} 
\sphinxAtStartPar
\sphinxstylestrong{Familia:} Arial

\item {} 
\sphinxAtStartPar
\sphinxstylestrong{Estilo:} Negrita itálica

\end{itemize}

\item {} 
\sphinxAtStartPar
En la pestaña \sphinxguilabel{Efectos tipográficos} configurar:
\begin{itemize}
\item {} 
\sphinxAtStartPar
\sphinxstylestrong{Color de letra:} Azul

\end{itemize}

\item {} 
\sphinxAtStartPar
En la pestaña \sphinxguilabel{Bordes} configurar:
\begin{itemize}
\item {} 
\sphinxAtStartPar
En la sección \sphinxstylestrong{Disposición de líneas:}
\begin{itemize}
\item {} 
\sphinxAtStartPar
Definido por el usuario: Agregar línea inferior

\end{itemize}

\noindent{\hspace*{\fill}\sphinxincludegraphics[width=180\sphinxpxdimen]{{image24}.png}\hspace*{\fill}}

\item {} 
\sphinxAtStartPar
En la sección \sphinxstylestrong{Línea}:
\begin{itemize}
\item {} 
\sphinxAtStartPar
\sphinxstylestrong{Color:} Azul

\item {} 
\sphinxAtStartPar
\sphinxstylestrong{Grosor:} Grueso (2,25pt)

\end{itemize}

\item {} 
\sphinxAtStartPar
En la sección \sphinxstylestrong{Estilo de sombra}:
\begin{itemize}
\item {} 
\sphinxAtStartPar
\sphinxstylestrong{Posición:} Proyectar sombra a la parte superior izquierda

\end{itemize}

\end{itemize}

\item {} 
\sphinxAtStartPar
PULSE el botón \sphinxstylestrong{Aceptar}.

\end{itemize}

\item {} 
\sphinxAtStartPar
En su cuaderno:

\begin{sphinxadmonition}{note}{En su cuaderno…}

\sphinxAtStartPar
Responda la siguiente pregunta:
\begin{itemize}
\item {} 
\sphinxAtStartPar
¿Cuál es la diferencia entre editar el estilo de párrafo de \sphinxstylestrong{Cuerpo de texto} y editar el estilo de párrafo de \sphinxstylestrong{Título 1?}

\end{itemize}
\end{sphinxadmonition}

\item {} 
\sphinxAtStartPar
Con las configuraciones aplicadas hasta el momento, el documento debería tener la siguiente apariencia:


\begin{savenotes}\sphinxattablestart
\sphinxthistablewithglobalstyle
\centering
\begin{tabulary}{\linewidth}[t]{T}
\sphinxtoprule
\sphinxtableatstartofbodyhook

\begin{sphinxfigure-in-table}
\centering

\noindent\sphinxincludegraphics[width=500\sphinxpxdimen]{{image25}.png}
\end{sphinxfigure-in-table}\relax
\\
\sphinxbottomrule
\end{tabulary}
\sphinxtableafterendhook\par
\sphinxattableend\end{savenotes}

\item {} 
\sphinxAtStartPar
Guarde los cambios y cierre el software LibreOffice Writer.

\end{enumerate}


\bigskip\hrule\bigskip



\subsection{A11: Estilo personalizado \sphinxstyleliteralintitle{\sphinxupquote{101}} (+0,2)}
\label{\detokenize{1_guias/10ofimatica/G1014:a11-estilo-personalizado-101-0-2}}\begin{enumerate}
\sphinxsetlistlabels{\arabic}{enumi}{enumii}{}{.}%
\item {} 
\sphinxAtStartPar
DESCARGUE el siguiente archivo y guárdelo en su carpeta de evidencias:
\begin{itemize}
\item {} 
\sphinxAtStartPar
G1014A11\sphinxhyphen{}estilo\_personalizado.odt \DUrole{xref}{\DUrole{download}{\DUrole{myst}{\DUrole{sd-sphinx-override}{\DUrole{sd-badge}{\DUrole{sd-bg-primary}{\DUrole{sd-bg-text-primary}{Descargar}}}}}}}

\end{itemize}

\item {} 
\sphinxAtStartPar
Abra el archivo: \sphinxstylestrong{G1014A11\sphinxhyphen{}estilo\_personalizado.odt} y en la cabecera del documento, agregue la \sphinxstylestrong{información de identificación} con el título: \sphinxstylestrong{ESTILO PERSONALIZADO}

\item {} 
\sphinxAtStartPar
En la \sphinxstylestrong{Jerarquía de estilos de párrafo}, IDENTIFIQUE el estilo: \sphinxstylestrong{Estilo de párrafo predeterminado} y en su \sphinxstylestrong{menú contextual}, PULSE la opción: \sphinxstylestrong{Nuevo…}

\item {} 
\sphinxAtStartPar
En la ventana: \sphinxstylestrong{Estilo de párrafo: Sin Título1}, configure las siguiente opciones:
\begin{itemize}
\item {} 
\sphinxAtStartPar
En la pestaña \sphinxguilabel{Generales} configurar:
\begin{itemize}
\item {} 
\sphinxAtStartPar
En la sección \sphinxstylestrong{Estilo:}
\begin{itemize}
\item {} 
\sphinxAtStartPar
\sphinxstylestrong{Nombre:} Mi estilo

\end{itemize}

\end{itemize}

\item {} 
\sphinxAtStartPar
En la pestaña \sphinxguilabel{Sangrías y espaciado} configurar:
\begin{itemize}
\item {} 
\sphinxAtStartPar
En la sección \sphinxstylestrong{Sangría}:
\begin{itemize}
\item {} 
\sphinxAtStartPar
\sphinxstylestrong{Antes de texto:} 5,00cm

\item {} 
\sphinxAtStartPar
\sphinxstylestrong{Después del texto:} 5,00cm

\end{itemize}

\item {} 
\sphinxAtStartPar
En la sección \sphinxstylestrong{Espaciado:}
\begin{itemize}
\item {} 
\sphinxAtStartPar
\sphinxstylestrong{Sobre el párrafo:} 1,00cm

\item {} 
\sphinxAtStartPar
\sphinxstylestrong{Bajo el párrafo:} 1,00cm

\end{itemize}

\item {} 
\sphinxAtStartPar
En la sección \sphinxstylestrong{Interlineado}:
\begin{itemize}
\item {} 
\sphinxAtStartPar
1,5 renglones

\end{itemize}

\end{itemize}

\item {} 
\sphinxAtStartPar
En la pestaña \sphinxguilabel{Alineación} configurar:
\begin{itemize}
\item {} 
\sphinxAtStartPar
En la sección \sphinxstylestrong{Opciones}:
\begin{itemize}
\item {} 
\sphinxAtStartPar
Centro

\end{itemize}

\end{itemize}

\item {} 
\sphinxAtStartPar
En la pestaña \sphinxguilabel{Tipo de letra} configurar:
\begin{itemize}
\item {} 
\sphinxAtStartPar
\sphinxstylestrong{Familia:} Arial

\item {} 
\sphinxAtStartPar
\sphinxstylestrong{Estilo:} Itálica

\item {} 
\sphinxAtStartPar
\sphinxstylestrong{Tamaño:} 11pt

\end{itemize}

\item {} 
\sphinxAtStartPar
En la pestaña \sphinxguilabel{Bordes} configurar:
\begin{itemize}
\item {} 
\sphinxAtStartPar
En la sección \sphinxstylestrong{Disposición de líneas}:
\begin{itemize}
\item {} 
\sphinxAtStartPar
Preconfigs: Los cuatro bordes

\end{itemize}

\item {} 
\sphinxAtStartPar
En la sección \sphinxstylestrong{Separación:}
\begin{itemize}
\item {} 
\sphinxAtStartPar
0,20cm en Izquierda, Derecha, Superior e Inferior.

\end{itemize}

\end{itemize}

\item {} 
\sphinxAtStartPar
En la pestaña \sphinxguilabel{Área} configurar:
\begin{itemize}
\item {} 
\sphinxAtStartPar
PULSAR el botón \sphinxstylestrong{Degradado.}
\begin{itemize}
\item {} 
\sphinxAtStartPar
En la sección \sphinxstylestrong{Degradado} seleccionar: Luz solar

\end{itemize}

\noindent{\hspace*{\fill}\sphinxincludegraphics[width=200\sphinxpxdimen]{{image26}.png}\hspace*{\fill}}
\begin{itemize}
\item {} 
\sphinxAtStartPar
En la sección \sphinxstylestrong{Opciones}:
\begin{itemize}
\item {} 
\sphinxAtStartPar
Tipo: Axial

\end{itemize}

\end{itemize}

\end{itemize}

\item {} 
\sphinxAtStartPar
PULSE el botón \sphinxstylestrong{Aceptar}.

\end{itemize}

\item {} 
\sphinxAtStartPar
Al final del texto, AGREGUE el siguiente párrafo:
\begin{quote}

\sphinxAtStartPar
“De necesitar referencias, solicítelas directamente al candidato. Gracias”.
\end{quote}

\item {} 
\sphinxAtStartPar
SELECCIONE el texto que acabó de agregar y APLIQUE el estilo: \sphinxstylestrong{mi estilo.}

\item {} 
\sphinxAtStartPar
Con las configuración aplicada hasta el momento, el texto debería tener la siguiente apariencia:


\begin{savenotes}\sphinxattablestart
\sphinxthistablewithglobalstyle
\centering
\begin{tabulary}{\linewidth}[t]{T}
\sphinxtoprule
\sphinxtableatstartofbodyhook

\begin{sphinxfigure-in-table}
\centering

\noindent\sphinxincludegraphics[width=300\sphinxpxdimen]{{image27}.png}
\end{sphinxfigure-in-table}\relax
\\
\sphinxbottomrule
\end{tabulary}
\sphinxtableafterendhook\par
\sphinxattableend\end{savenotes}

\item {} 
\sphinxAtStartPar
En su cuaderno:

\begin{sphinxadmonition}{note}{En su cuaderno…}

\sphinxAtStartPar
Responda la siguiente pregunta:
\begin{itemize}
\item {} 
\sphinxAtStartPar
¿Para que casos es útil configurar un estilo personalizado en \sphinxstylestrong{LibreOffice Writer}?

\end{itemize}
\end{sphinxadmonition}

\item {} 
\sphinxAtStartPar
GUARDE los cambios y cierre el software \sphinxstylestrong{LibreOffice Writer}.

\end{enumerate}


\bigskip\hrule\bigskip



\subsection{A12: Presentación de evidencias (+3,0)}
\label{\detokenize{1_guias/10ofimatica/G1014:a12-presentacion-de-evidencias-3-0}}
\sphinxAtStartPar
\DUrole{sd-sphinx-override}{\DUrole{sd-badge}{\DUrole{sd-bg-danger}{\DUrole{sd-bg-text-danger}{Transferencia}}}}

\begin{sphinxadmonition}{note}{¿Cómo se evalúa?}
\begin{enumerate}
\sphinxsetlistlabels{\arabic}{enumi}{enumii}{}{.}%
\item {} 
\sphinxAtStartPar
Los estudiantes deben presentar lo siguiente:
\begin{itemize}
\item {} 
\sphinxAtStartPar
Carpeta de evidencias con todos los archivos editados en la presente guía. (Los archivos deben estar abiertos).

\item {} 
\sphinxAtStartPar
Cuaderno con las actividades desarrolladas.

\end{itemize}

\item {} 
\sphinxAtStartPar
El docente realiza unas preguntas para verificación de conocimientos.

\item {} 
\sphinxAtStartPar
El estudiante debe demostrar de forma práctica sus conocimientos en el uso de \sphinxstylestrong{LibreOffice Writer}.

\end{enumerate}
\end{sphinxadmonition}

\sphinxstepscope


\section{Guía 1019: Plan de refuerzo}
\label{\detokenize{1_guias/10ofimatica/G1019:guia-1019-plan-de-refuerzo}}\label{\detokenize{1_guias/10ofimatica/G1019::doc}}\begin{quote}\begin{description}
\sphinxlineitem{ASIGNATURA}
\sphinxAtStartPar
Ofimática

\sphinxlineitem{GRADO}
\sphinxAtStartPar
Décimo

\end{description}\end{quote}


\bigskip\hrule\bigskip



\subsection{PRIMER PERIODO}
\label{\detokenize{1_guias/10ofimatica/G1019:primer-periodo}}
\begin{sphinxadmonition}{note}{EXAMEN ESCRITO: Plan de refuerzo (10 minutos)}

\sphinxAtStartPar
Estudiar los conceptos a continuación. Se realiza examen de 5 preguntas. Si el estudiante no asiste en la fecha programada, este será evaluado la próxima vez que asista a clase.
\end{sphinxadmonition}


\subsubsection{¿Qué es la ofimática?}
\label{\detokenize{1_guias/10ofimatica/G1019:que-es-la-ofimatica}}
\sphinxAtStartPar
La \sphinxstylestrong{ofimática} es el conjunto de técnicas, aplicaciones y herramientas informáticas que se utilizan para optimizar, automatizar y mejorar los procedimientos o tareas relacionadas con el trabajo de oficina. El término “ofimática” es una combinación de las palabras “oficina” e “informática”.


\subsubsection{¿Qué es Google Workspace?}
\label{\detokenize{1_guias/10ofimatica/G1019:que-es-google-workspace}}
\sphinxAtStartPar
\sphinxstylestrong{Google Workspace} es una solución de productividad en línea desarrollada por \sphinxstylestrong{Google} que permite a los usuarios conectarse, crear y colaborar de manera segura. Incluye herramientas esenciales como:
\begin{itemize}
\item {} 
\sphinxAtStartPar
\sphinxstylestrong{Gmail}: Servicio de correo electrónico.

\item {} 
\sphinxAtStartPar
\sphinxstylestrong{Documentos}: Para crear y editar documentos de texto.

\item {} 
\sphinxAtStartPar
\sphinxstylestrong{Hojas de cálculo}: Para trabajar con datos y realizar cálculos.

\item {} 
\sphinxAtStartPar
\sphinxstylestrong{Presentaciones}: Para crear presentaciones visuales.

\item {} 
\sphinxAtStartPar
\sphinxstylestrong{Drive}: Almacenamiento en la nube para guardar y compartir archivos.

\item {} 
\sphinxAtStartPar
\sphinxstylestrong{Calendario}: Para gestionar eventos y citas.

\end{itemize}


\subsubsection{¿Qué es Microsoft Office?}
\label{\detokenize{1_guias/10ofimatica/G1019:que-es-microsoft-office}}
\sphinxAtStartPar
\sphinxstylestrong{Microsoft Office} es un conjunto de aplicaciones de software de productividad creado por \sphinxstylestrong{Microsoft}®. Es una herramienta muy popular y ampliamente utilizada en todo el mundo para una variedad de tareas relacionadas con la creación, edición y gestión de documentos, hojas de cálculo, presentaciones y más.

\sphinxAtStartPar
Algunas de las aplicaciones más conocidas que forman parte de Microsoft Office son:
\begin{itemize}
\item {} 
\sphinxAtStartPar
\sphinxstylestrong{Word:} Para crear y editar documentos de texto, como cartas, informes y trabajos académicos.

\item {} 
\sphinxAtStartPar
\sphinxstylestrong{Excel:} Para crear hojas de cálculo, analizar datos, realizar cálculos y crear gráficos.

\item {} 
\sphinxAtStartPar
\sphinxstylestrong{PowerPoint:} Para crear presentaciones visuales con diapositivas, texto, imágenes y animaciones.

\item {} 
\sphinxAtStartPar
\sphinxstylestrong{Outlook:} Para gestionar correos electrónicos, calendarios, contactos y tareas.

\item {} 
\sphinxAtStartPar
\sphinxstylestrong{Access:} Para crear y gestionar bases de datos (menos común en el uso diario).

\end{itemize}


\subsubsection{¿Qué es el software libre?}
\label{\detokenize{1_guias/10ofimatica/G1019:que-es-el-software-libre}}
\sphinxAtStartPar
El \sphinxstylestrong{software libre} es un tipo de software que respeta la libertad de los usuarios y la comunidad. Los usuarios tienen la libertad de ejecutar, copiar, distribuir, estudiar, modificar y mejorar el software.

\sphinxAtStartPar
Las cuatro libertades esenciales del software libre son:
\begin{itemize}
\item {} 
\sphinxAtStartPar
\sphinxstylestrong{(Libertad 0):} La libertad de ejecutar el programa como se desee, con cualquier propósito.

\item {} 
\sphinxAtStartPar
\sphinxstylestrong{(Libertad 1):} La libertad de estudiar cómo funciona el programa y cambiarlo para que haga lo que uno quiera. (El acceso al código fuente es una condición necesaria para ello).

\item {} 
\sphinxAtStartPar
\sphinxstylestrong{(Libertad 2):} La libertad de redistribuir copias para ayudar a otros.

\item {} 
\sphinxAtStartPar
\sphinxstylestrong{(Libertad 3):} La libertad de distribuir copias de sus versiones modificadas a terceros.

\end{itemize}


\paragraph{Ejemplos de software libre}
\label{\detokenize{1_guias/10ofimatica/G1019:ejemplos-de-software-libre}}\begin{itemize}
\item {} 
\sphinxAtStartPar
Sistema Operativo GNU\sphinxhyphen{}LINUX (\$0)

\item {} 
\sphinxAtStartPar
Navegador web Mozilla FIREFOX (\$0)

\item {} 
\sphinxAtStartPar
Editor de imágenes GIMP (\$0)

\item {} 
\sphinxAtStartPar
Editor de gráficos vectoriales INKSCAPE (\$0)

\item {} 
\sphinxAtStartPar
Reproductor multimedia VLC (\$0)

\item {} 
\sphinxAtStartPar
Suite de creación 3D BLENDER (\$0)

\item {} 
\sphinxAtStartPar
Edición de sonido Audacity (\$0)

\item {} 
\sphinxAtStartPar
Suite ofimática Libreoffice (\$0)

\end{itemize}


\subsubsection{¿Qué es el software propietario?}
\label{\detokenize{1_guias/10ofimatica/G1019:que-es-el-software-propietario}}
\sphinxAtStartPar
El software propietario, es aquel que no permite a los usuarios acceder, modificar o distribuir su código fuente. Los usuarios tienen restricciones legales y técnicas que les impiden ejercer las libertades que caracterizan al software libre.
\begin{itemize}
\item {} 
\sphinxAtStartPar
El código fuente no está disponible para los usuarios.

\item {} 
\sphinxAtStartPar
Prohibición de copiar, modificar o redistribuir el software.

\item {} 
\sphinxAtStartPar
Generalmente requiere el pago de licencias para su uso.

\item {} 
\sphinxAtStartPar
Las actualizaciones y mejoras dependen exclusivamente del propietario.

\end{itemize}


\paragraph{Ejemplos de software propietario}
\label{\detokenize{1_guias/10ofimatica/G1019:ejemplos-de-software-propietario}}\begin{itemize}
\item {} 
\sphinxAtStartPar
Sistema operativo Microsoft Windows (\$769,999)

\item {} 
\sphinxAtStartPar
Microsoft Office (suite ofimática) (\$459,999/año)

\item {} 
\sphinxAtStartPar
Adobe Photoshop (editor de imágenes) (\$48,971/mes)

\item {} 
\sphinxAtStartPar
AutoCAD (software de diseño) (\$2,201,999/año)

\end{itemize}


\subsubsection{¿Qué es el software pirata?}
\label{\detokenize{1_guias/10ofimatica/G1019:que-es-el-software-pirata}}
\sphinxAtStartPar
El software pirata se refiere a copias ilegales de \sphinxstylestrong{programas informáticos} que se distribuyen o utilizan sin la debida autorización de sus propietarios legales. Esta práctica viola los derechos de autor y la propiedad intelectual


\subsubsection{¿Qué es LibreOffice?}
\label{\detokenize{1_guias/10ofimatica/G1019:que-es-libreoffice}}
\sphinxAtStartPar
\sphinxstylestrong{LibreOffice} es una suite de productividad de oficina de código abierto, disponible de forma gratuita, que es compatible con otras suites de oficina y está disponible en una variedad de plataformas. Su formato de archivo nativo es \sphinxstylestrong{Open Document Format (ODF)} y puede abrir y guardar documentos en muchos formatos, incluidos los utilizados por varias versiones de Microsoft Office.

\sphinxAtStartPar
La suite de \sphinxstylestrong{LibreOffice} esta compuesta por las siguientes herramientas:
\begin{itemize}
\item {} 
\sphinxAtStartPar
El procesador de textos \sphinxcode{\sphinxupquote{Writer}}.

\item {} 
\sphinxAtStartPar
La hoja de cálculo \sphinxcode{\sphinxupquote{Calc}}.

\item {} 
\sphinxAtStartPar
El programa de presentaciones \sphinxcode{\sphinxupquote{Impress}}.

\item {} 
\sphinxAtStartPar
El programa de dibujo vectorial \sphinxcode{\sphinxupquote{Draw}}.

\item {} 
\sphinxAtStartPar
La base de datos \sphinxcode{\sphinxupquote{Base}}.

\item {} 
\sphinxAtStartPar
El editor de ecuaciones \sphinxcode{\sphinxupquote{Math}}.

\end{itemize}


\bigskip\hrule\bigskip



\subsection{SEGUNDO PERIODO}
\label{\detokenize{1_guias/10ofimatica/G1019:segundo-periodo}}
\begin{sphinxadmonition}{note}{EXAMEN ESCRITO: Plan de refuerzo (10 minutos)}

\sphinxAtStartPar
Estudiar los conceptos a continuación. Se realiza examen de 5 preguntas. Si el estudiante no asiste en la fecha programada, este será evaluado la próxima vez que asista a clase.
\end{sphinxadmonition}

\sphinxstepscope


\chapter{Programación 10}
\label{\detokenize{1_guias/10programacion/G1020:programacion-10}}\label{\detokenize{1_guias/10programacion/G1020::doc}}
\sphinxstepscope


\section{Guía 1021: PROGRAMACIÓN (P1)}
\label{\detokenize{1_guias/10programacion/G1021:guia-1021-programacion-p1}}\label{\detokenize{1_guias/10programacion/G1021::doc}}

\begin{savenotes}\sphinxattablestart
\sphinxthistablewithglobalstyle
\centering
\begin{tabular}[t]{\X{35}{100}\X{65}{100}}
\sphinxtoprule
\sphinxtableatstartofbodyhook

\begin{savenotes}\sphinxattablestart
\sphinxthistablewithglobalstyle
\centering
\begin{tabulary}{\linewidth}[t]{T}
\sphinxtoprule
\sphinxtableatstartofbodyhook
\sphinxAtStartPar
\sphinxincludegraphics{{escudo}.png}
\\
\sphinxbottomrule
\end{tabulary}
\sphinxtableafterendhook\par
\sphinxattableend\end{savenotes}
&

\begin{savenotes}\sphinxattablestart
\sphinxthistablewithglobalstyle
\raggedright
\begin{tabular}[t]{*{1}{\X{1}{1}}}
\sphinxtoprule
\sphinxtableatstartofbodyhook
\sphinxAtStartPar
\sphinxstylestrong{INSTITUCIÓN EDUCATIVA}: John F. Kennedy
\begin{itemize}
\item {} 
\sphinxAtStartPar
\sphinxstylestrong{DOCENTE}: Julian Camilo Fonseca Romero

\item {} 
\sphinxAtStartPar
\sphinxstylestrong{ASIGNATURA}: Programación

\item {} 
\sphinxAtStartPar
\sphinxstylestrong{GRADO}: ONCE

\item {} 
\sphinxAtStartPar
\sphinxstylestrong{PERIODO}: 01

\item {} 
\sphinxAtStartPar
\sphinxstylestrong{TIEMPO}: 10 semanas (2 horas / semana)

\item {} 
\sphinxAtStartPar
\sphinxstylestrong{TIEMPO DE TRABAJO}: 8/20 horas

\end{itemize}
\\
\sphinxbottomrule
\end{tabular}
\sphinxtableafterendhook\par
\sphinxattableend\end{savenotes}
\\
\sphinxbottomrule
\end{tabular}
\sphinxtableafterendhook\par
\sphinxattableend\end{savenotes}

\begin{sphinxuseclass}{sd-tab-set}
\begin{sphinxuseclass}{sd-tab-item}\subsubsection*{Objetivo}

\begin{sphinxuseclass}{sd-tab-content}
\sphinxAtStartPar
\DUrole{sd-sphinx-override}{\DUrole{sd-badge}{\DUrole{sd-bg-light}{\DUrole{sd-bg-text-light}{Exploración}}}}
\DUrole{sd-sphinx-override}{\DUrole{sd-badge}{\DUrole{sd-bg-info}{\DUrole{sd-bg-text-info}{Estructuración}}}}
\DUrole{sd-sphinx-override}{\DUrole{sd-badge}{\DUrole{sd-bg-warning}{\DUrole{sd-bg-text-warning}{Actividad práctica}}}}
\DUrole{sd-sphinx-override}{\DUrole{sd-badge}{\DUrole{sd-bg-danger}{\DUrole{sd-bg-text-danger}{Transferencia}}}}
\begin{itemize}
\item {} 
\sphinxAtStartPar
\sphinxstylestrong{Temática}: Introducción a la programación.

\item {} 
\sphinxAtStartPar
\sphinxstylestrong{Objetivo de aprendizaje}: Dar a conocer los conceptos introductorios a la programación.

\end{itemize}

\end{sphinxuseclass}
\end{sphinxuseclass}
\begin{sphinxuseclass}{sd-tab-item}\subsubsection*{Orientaciones curriculares}

\begin{sphinxuseclass}{sd-tab-content}\begin{itemize}
\item {} 
\sphinxAtStartPar
\sphinxstylestrong{Componente}: Solución de problemas con Tecnología e Informática.

\item {} 
\sphinxAtStartPar
\sphinxstylestrong{Competencia}: Propongo innovaciones tecnológicas e informáticas para la solución de problemas dando cumplimiento a restricciones, condiciones y especificaciones técnicas y contextuales.

\item {} 
\sphinxAtStartPar
\sphinxstylestrong{Evidencia de aprendizaje}: Represento ideas sobre diseños, innovaciones tecnológicas o informáticas mediante el uso de registros, textos, diagramas, figuras… empleando para ello herramientas informáticas.

\end{itemize}

\end{sphinxuseclass}
\end{sphinxuseclass}
\begin{sphinxuseclass}{sd-tab-item}\subsubsection*{Criterios de evaluación}

\begin{sphinxuseclass}{sd-tab-content}\begin{itemize}
\item {} 
\sphinxAtStartPar
Actividades desarrolladas en su totalidad.

\item {} 
\sphinxAtStartPar
Participación en clase.

\item {} 
\sphinxAtStartPar
Explicación clara de conceptos.

\item {} 
\sphinxAtStartPar
Argumentación de ideas.

\item {} 
\sphinxAtStartPar
Cumplimiento de las reglas del aula.

\end{itemize}

\end{sphinxuseclass}
\end{sphinxuseclass}
\begin{sphinxuseclass}{sd-tab-item}\subsubsection*{Recursos (0)}

\begin{sphinxuseclass}{sd-tab-content}\begin{itemize}
\item {} 
\sphinxAtStartPar
La presente guía no contiene recursos.

\end{itemize}

\end{sphinxuseclass}
\end{sphinxuseclass}
\end{sphinxuseclass}

\subsection{Mensaje del autor}
\label{\detokenize{1_guias/10programacion/G1021:mensaje-del-autor}}\begin{quote}

\sphinxAtStartPar
“En esta guía didáctica, se exploraran los fundamentos esenciales en el desarrollo de software. El objetivo es obtener conocimientos básicos, para entender el cómo funciona el mundo digital y que a su vez impulsen la motivación hacia el aprendizaje de la programación”.

\begin{flushright}
---Camilo Fonseca
\end{flushright}
\end{quote}


\bigskip\hrule\bigskip



\subsection{A00: Introducción (+0,1)}
\label{\detokenize{1_guias/10programacion/G1021:a00-introduccion-0-1}}
\sphinxAtStartPar
\DUrole{sd-sphinx-override}{\DUrole{sd-badge}{\DUrole{sd-bg-light}{\DUrole{sd-bg-text-light}{Exploración}}}}
\begin{enumerate}
\sphinxsetlistlabels{\arabic}{enumi}{enumii}{}{.}%
\item {} 
\sphinxAtStartPar
TRANSCRIBA en su cuaderno el objetivo, las orientaciones curriculares y los criterios de evaluación de la presente guía.

\item {} 
\sphinxAtStartPar
Lea el mensaje del autor y en su cuaderno responda la siguiente pregunta:
\begin{itemize}
\item {} 
\sphinxAtStartPar
¿En qué se relacionan el mundo digital y el software?

\end{itemize}

\end{enumerate}


\bigskip\hrule\bigskip



\subsection{A01: PROGRAMACIÓN (+0,1)}
\label{\detokenize{1_guias/10programacion/G1021:a01-programacion-0-1}}
\sphinxAtStartPar
\DUrole{sd-sphinx-override}{\DUrole{sd-badge}{\DUrole{sd-bg-info}{\DUrole{sd-bg-text-info}{Estructuración}}}}
\begin{enumerate}
\sphinxsetlistlabels{\arabic}{enumi}{enumii}{}{.}%
\item {} 
\sphinxAtStartPar
Lea los siguientes conceptos relacionados con programación:

\end{enumerate}


\bigskip\hrule\bigskip



\subsubsection{¿Qué es la programación?}
\label{\detokenize{1_guias/10programacion/G1021:que-es-la-programacion}}
\sphinxAtStartPar
La programación se refiere específicamente a la escritura de código para crear funcionalidades dentro del software. Es una parte esencial del desarrollo, pero no abarca todo el proceso.

\sphinxAtStartPar
En la programación, las actividades a desarrollar están orientadas a la creación de código fuente. Es decir, una serie de instrucciones que se le da al PC para realizar alguna acción. También se realiza depuración y optimización de código lo cual permite que solucionen muchos errores. Es muy importante que el programador entienda la responsabilidad de salvaguardar el código fuente construido, por tanto debe sacar tiempo para tareas de \sphinxstylestrong{copias de seguridad} y \sphinxstylestrong{control de versiones}.


\begin{savenotes}\sphinxattablestart
\sphinxthistablewithglobalstyle
\centering
\begin{tabulary}{\linewidth}[t]{T}
\sphinxtoprule
\sphinxtableatstartofbodyhook

\begin{sphinxfigure-in-table}
\centering
\capstart
\noindent\sphinxincludegraphics[width=360\sphinxpxdimen]{{image019}.png}
\sphinxfigcaption{Código fuente}\label{\detokenize{1_guias/10programacion/G1021:id1}}\end{sphinxfigure-in-table}\relax
\\
\sphinxbottomrule
\end{tabulary}
\sphinxtableafterendhook\par
\sphinxattableend\end{savenotes}


\bigskip\hrule\bigskip



\subsubsection{¿Qué es el software?}
\label{\detokenize{1_guias/10programacion/G1021:que-es-el-software}}
\sphinxAtStartPar
\sphinxstylestrong{El software} es el conjunto de programas, bases de datos y diferentes tipos de archivos que controlan el hardware a través de instrucciones. El software debe tener requisitos y funcionalidades. Es la parte intangible de un sistema informático, en contraste con el hardware, que es la parte física. El software se puede clasificar en dos tipos:
\begin{itemize}
\item {} 
\sphinxAtStartPar
\sphinxstylestrong{Software de sistema:} Proporciona funciones básicas, como sistemas operativos, gestión de discos, servicios, gestión de hardware y otras necesidades operativas.

\item {} 
\sphinxAtStartPar
\sphinxstylestrong{Software de aplicación:} Ayuda a los usuarios a realizar tareas. Por ejemplo, conjuntos de productividad ofimática, software de gestión de datos, reproductores multimedia y programas de seguridad. Las aplicaciones también se refieren a aplicaciones web y móviles como las que se utilizan para comprar en \sphinxstylestrong{Amazon}, socializar con \sphinxstylestrong{Facebook} o publicar fotos en \sphinxstylestrong{Instagram}.

\item {} 
\sphinxAtStartPar
\sphinxstylestrong{Software de programación:} Proporciona herramientas a los programadores, como editores de texto, compiladores, enlazadores, depuradores y otras herramientas para crear código.

\item {} 
\sphinxAtStartPar
\sphinxstylestrong{Software embebido:} Se utiliza para controlar máquinas y dispositivos que normalmente no se consideran ordenadores: redes de telecomunicaciones, automóviles, robots industriales, etc. Estos dispositivos, y su software pueden conectarse como parte del Internet de las cosas (IoT).

\end{itemize}


\bigskip\hrule\bigskip



\subsubsection{¿Qué es el desarrollo de software?}
\label{\detokenize{1_guias/10programacion/G1021:que-es-el-desarrollo-de-software}}
\sphinxAtStartPar
El \sphinxstylestrong{desarrollo de software} es un proceso que abarca la programación, pero también incluye planificación, diseño, pruebas y mantenimiento. Es un enfoque más holístico que busca crear software funcional y útil.

\sphinxAtStartPar
En el \sphinxstylestrong{desarrollo de software} se realizan tareas como el análisis de requisitos que busca identificar que funciones debe tener dicho software. También se hacen tareas de diseño donde se trabaja con bocetos iniciales para tener una idea general de como estará compuesto y de que forma funcionará el software. Otra tarea muy común es la programación la cual consiste en crear código fuente para la construcción del programa. Finalmente se realizan pruebas, implementación, mantenimiento y documentación.


\begin{savenotes}\sphinxattablestart
\sphinxthistablewithglobalstyle
\centering
\begin{tabulary}{\linewidth}[t]{T}
\sphinxtoprule
\sphinxtableatstartofbodyhook

\begin{sphinxfigure-in-table}
\centering
\capstart
\noindent\sphinxincludegraphics[width=360\sphinxpxdimen]{{image027}.png}
\sphinxfigcaption{Equipo de desarrollo de software}\label{\detokenize{1_guias/10programacion/G1021:id2}}\end{sphinxfigure-in-table}\relax
\\
\sphinxbottomrule
\end{tabulary}
\sphinxtableafterendhook\par
\sphinxattableend\end{savenotes}


\bigskip\hrule\bigskip



\subsubsection{Programadores Vs Desarrolladores}
\label{\detokenize{1_guias/10programacion/G1021:programadores-vs-desarrolladores}}
\sphinxAtStartPar
El desarrollo de software lo llevan a cabo principalmente programadores y desarrolladores de software. Estos roles interactúan y se superponen, y la dinámica entre ellos varía mucho entre los departamentos de desarrollo y las comunidades.
\begin{itemize}
\item {} 
\sphinxAtStartPar
\sphinxstylestrong{Programadores:} Escriben código fuente para programar ordenadores que realicen tareas específicas, como fusionar bases de datos, procesar pedidos en línea, enrutar comunicaciones, realizar búsquedas o mostrar textos y gráficos. Los programadores suelen interpretar las instrucciones de los desarrolladores e ingenieros de software y utilizan lenguajes de programación como C++ o Java para llevarlas a cabo.

\item {} 
\sphinxAtStartPar
\sphinxstylestrong{Desarrollador de software:} Estan estrechamente involucrados en áreas específicas del proyecto, incluida la escritura de código. Al mismo tiempo, dirigen todo el ciclo de vida del desarrollo de software, lo que implica colaborar con los equipos funcionales para transformar los requisitos en características, gestionar los equipos y procesos de desarrollo, y realizar pruebas y mantenimiento del software.

\end{itemize}

\begin{sphinxadmonition}{note}{¿Cuál es la diferencia entre un programador y un desarrollador?}

\sphinxAtStartPar
Construir un software es como construir una misma casa varias veces. El lograr que cada persona aporte significativamente en su campo de acción ya sea como desarrollador o como programador, hará que la casa se construya de forma óptima.
\begin{itemize}
\item {} 
\sphinxAtStartPar
Un \sphinxstylestrong{desarrollador} diseña, construye y mantiene software mediante un proceso compuesto por etapas o fases(análisis de requisitos, diseño, construcción, pruebas, implementación y mantenimiento).

\item {} 
\sphinxAtStartPar
Un \sphinxstylestrong{programador} codifica programas que hacen parte del software.

\end{itemize}
\end{sphinxadmonition}


\bigskip\hrule\bigskip



\subsubsection{Proceso de desarrollo de software}
\label{\detokenize{1_guias/10programacion/G1021:proceso-de-desarrollo-de-software}}
\sphinxAtStartPar
Los pasos del proceso de desarrollo de software se pueden agrupar en las fases del ciclo de vida, pero la importancia del ciclo de vida es que se recicla para permitir una mejora continua. Por ejemplo, los problemas de los usuarios que surgen en la fase de mantenimiento y soporte pueden convertirse en requisitos al inicio del siguiente ciclo.
\begin{itemize}
\item {} 
\sphinxAtStartPar
\sphinxstylestrong{FASE 01 (Análisis de requisitos)}

\sphinxAtStartPar
En esta fase se comprende y documenta lo que necesitan los usuarios y otras partes interesadas. Se analiza la estructura tecnológica donde funcionará el software.

\item {} 
\sphinxAtStartPar
\sphinxstylestrong{FASE 02 (Diseño)}

\sphinxAtStartPar
Se diseña en torno a soluciones a los problemas presentados por los requisitos, lo que a menudo implica modelos de procesos y guiones gráficos.
Se crean los modelos con una herramienta de modelado que utilice un lenguaje de modelado como UML para llevar a cabo la validación temprana, la creación de prototipos y la simulación del diseño.

\item {} 
\sphinxAtStartPar
\sphinxstylestrong{FASE 03 (Construcción)}

\sphinxAtStartPar
En la construcción se usa un lenguaje de programación adecuado al aplicativo que se quiere construir. Implica la revisión por pares y en equipo para eliminar problemas de forma temprana y producir software de calidad más rápido.

\item {} 
\sphinxAtStartPar
\sphinxstylestrong{FASE 04 (Pruebas)}

\sphinxAtStartPar
Se realizan dos tipos de pruebas:
\begin{itemize}
\item {} 
\sphinxAtStartPar
Pruebas con escenarios planificados de antemano como parte del diseño y la codificación del software. Por ejemplo identificación de errores de código, de lógica de negocio o de seguridad.

\item {} 
\sphinxAtStartPar
Pruebas de rendimiento para simular niveles de carga de uso en la aplicación. Por ejemplo Identificar en que momento el aplicativo puede presentar fallos cuando varios usuarios lo usen al mismo tiempo.

\end{itemize}

\item {} 
\sphinxAtStartPar
\sphinxstylestrong{FASE 05 (Implementación)}

\sphinxAtStartPar
El software se instala o pone en producción para capacitar a los usuarios y lograr que se familiaricen con las funciones que este contenga. Se resuelven dudas de los usuarios. Si hay una versión anterior del software, la información alojada en el software anterior debe ser migrada el software nuevo desde aplicaciones o fuentes de datos existentes, si es necesario.

\item {} 
\sphinxAtStartPar
\sphinxstylestrong{FASE 06 (Mantenimiento y soporte)}

\sphinxAtStartPar
Para mantener la calidad y la entrega a lo largo del ciclo de vida de la aplicación, se deben tener fechas de seguimiento y soporte. Los usuarios nuevos deben ser capacitados. El software en lo posible debe ser escalable es decir, en la medida que requiera nuevas funcionalidades, estas se deben ir implementando para que el software no quede desactualizado en la medida que pase el tiempo.

\end{itemize}


\bigskip\hrule\bigskip



\subsubsection{Tecnologías usadas en el desarrollo de software}
\label{\detokenize{1_guias/10programacion/G1021:tecnologias-usadas-en-el-desarrollo-de-software}}
\sphinxAtStartPar
Hoy en día son bastantes las tecnologías que  como herramientas que nos ayudan a construir software. Estas se pueden clasificar de la siguiente forma:

\begin{sphinxadmonition}{note}{Sistemas operativos}

\sphinxAtStartPar
Un sistema operativo (OS) es un software esencial que actúa como intermediario entre el usuario y el hardware de la computadora. Su función principal es gestionar los recursos del sistema y proporcionar una interfaz para que los usuarios y las aplicaciones interactúen con el hardware. Por ejemplo:
\begin{itemize}
\item {} 
\sphinxAtStartPar
\sphinxstylestrong{LINUX:} Sistema operativo libre.

\item {} 
\sphinxAtStartPar
\sphinxstylestrong{Windows:} Sistema operativo privativo.

\item {} 
\sphinxAtStartPar
\sphinxstylestrong{MACOS:} Sistema operativo privativo para equipos de computo de APPLE.

\end{itemize}
\end{sphinxadmonition}

\begin{sphinxadmonition}{note}{Bases de datos}

\sphinxAtStartPar
Una base de datos (DB) es un sistema organizado que permite almacenar, gestionar y recuperar información de manera eficiente. Una base de datos es como un archivo digital muy bien estructurado donde se guarda información que se puede buscar y manipular fácilmente. Por ejemplo:
\begin{itemize}
\item {} 
\sphinxAtStartPar
\sphinxstylestrong{MongoDB:} Almacena los datos en documentos.

\item {} 
\sphinxAtStartPar
\sphinxstylestrong{MySQL:} Utiliza tablas para almacenar datos. La empresa propietaria es ORACLE.

\item {} 
\sphinxAtStartPar
\sphinxstylestrong{SQL\_Server:} Utiliza tablas para almacenar datos. La empresa propietaria es MICROSOFT.

\end{itemize}
\end{sphinxadmonition}

\begin{sphinxadmonition}{note}{Servidor web}

\sphinxAtStartPar
Un servidor web permite alojar páginas web y ejecutar aplicaciones web. Ejemplo:
\begin{itemize}
\item {} 
\sphinxAtStartPar
\sphinxstylestrong{APACHE:} Es un servidor web tradicional que se basa en un modelo de procesamiento de solicitudes por separado. Cada solicitud genera un hilo diferente.

\end{itemize}
\end{sphinxadmonition}

\begin{sphinxadmonition}{note}{Entornos de ejecución}

\sphinxAtStartPar
Un entorno de ejecución es un software que contiene todo lo necesario para que otro software se ejecute correctamente. Esto incluye desde un intermediario con el sistema operativo hasta librerías y demás recursos necesarios como la documentación. Por ejemplo:
\begin{itemize}
\item {} 
\sphinxAtStartPar
\sphinxstylestrong{NODE:} También conocido como NODEJS, es un \sphinxstylestrong{entorno de ejecución} que permite ejecutar código JavaScript en el lado del servidor. A diferencia de JavaScript tradicional, que se ejecuta en el navegador, Node.js permite que los desarrolladores utilicen JavaScript para crear aplicaciones del lado del servidor, lo que lo hace muy versátil y poderoso.

\end{itemize}
\end{sphinxadmonition}

\begin{sphinxadmonition}{note}{Framework}

\sphinxAtStartPar
Un framework es un conjunto de librerías, documentación que ayudan a desarrollar o ejecutar un software en específico. Por ejemplo:
\begin{itemize}
\item {} 
\sphinxAtStartPar
\sphinxstylestrong{.NET:} Creado por Microsoft, permite crear aplicaciones de software más fácil y eficiente.

\item {} 
\sphinxAtStartPar
\sphinxstylestrong{EXPRESS:} Es un \sphinxstylestrong{framework de aplicación web} para Node.js que facilita la creación de aplicaciones que se ejecutan en el servidor web.

\item {} 
\sphinxAtStartPar
\sphinxstylestrong{ANGULAR:} Es un \sphinxstylestrong{framework de aplicación web} para Node.js que facilita la creación de aplicaciones que se ejecutan en el navegador web.

\end{itemize}
\end{sphinxadmonition}

\begin{sphinxadmonition}{note}{IDE}

\sphinxAtStartPar
Un IDE son las siglas para (\sphinxstylestrong{entorno de desarrollo de software)}. Es una aplicación que proporciona un conjunto completo de herramientas para facilitar el desarrollo de software. Un IDE combina diferentes herramientas y funciones en un solo lugar para ayudar a los programadores a escribir, depurar y mantener su código de manera más eficiente. Por ejemplo:
\begin{itemize}
\item {} 
\sphinxAtStartPar
\sphinxstylestrong{VISUAL STUDIO CODE:} IDE de Microsoft para programar con cualquier lenguaje de programación.

\item {} 
\sphinxAtStartPar
\sphinxstylestrong{THONNY:} IDE para programar con el lenguaje de programación PYTHON.

\end{itemize}
\end{sphinxadmonition}


\bigskip\hrule\bigskip



\subsubsection{Desarollo web}
\label{\detokenize{1_guias/10programacion/G1021:desarollo-web}}
\sphinxAtStartPar
Se conoce como desarrollo web al proceso de crear y mantener un sitio web que sea funcional en internet, a través de diferentes lenguajes de programación, según el modelo y la parte de la página que corresponda. Cada sitio tiene una URL única que lo distingue de los demás en la red informática mundial.

\sphinxAtStartPar
Un sitio web puede clasificarse de diferentes formas. Para cuestiones de desarrollo web principalmente se divide en dos partes:
\begin{itemize}
\item {} 
\sphinxAtStartPar
\sphinxstylestrong{Front\sphinxhyphen{}end:} Es la parte que interactúa con el usuario, tanto en imagen como en función. Por ello está íntimamente relacionada con la experiencia del usuario (UX) y la interfaz de usuario (IU).

\item {} 
\sphinxAtStartPar
\sphinxstylestrong{Back\sphinxhyphen{}end:} Se refiere a la parte que está en contacto directo con el servidor; es donde se aplica el código de programación para crear la estructura. Permanece en un segundo plano a cargo de la accesibilidad, actualización, bases de datos y cambios del sitio.

\end{itemize}

\sphinxAtStartPar
Los desarrolladores web se enfocan en la codificación y programación del sitio web utilizando lenguajes como HTML, CSS, JavaScript y otros. Su objetivo es que el sitio responda correctamente a las interacciones que realice el usuario en él y concretar un stack tecnológico adecuado. Algunos ejemplos de STACK son:
\begin{itemize}
\item {} 
\sphinxAtStartPar
LAMP: Linux, Apache, MySQL, PHP/Python/Perl

\item {} 
\sphinxAtStartPar
MEAN: MongoDB, Express.js, Angular, Node.js

\item {} 
\sphinxAtStartPar
.NET stack: .NET, C\#, SQL Server y Visual Studio

\item {} 
\sphinxAtStartPar
DJANGO: Python

\end{itemize}


\bigskip\hrule\bigskip



\subsubsection{Lenguajes de programación}
\label{\detokenize{1_guias/10programacion/G1021:lenguajes-de-programacion}}
\sphinxAtStartPar
Un \sphinxstylestrong{lenguaje de programación} es un conjunto de instrucciones que le damos a una computadora para que realice tareas específicas. Es como un idioma que usamos para comunicarnos con las máquinas. Estas instrucciones se escriben siguiendo reglas gramaticales y de sintaxis definidas, y el objetivo es crear programas que resuelvan problemas o hagan cosas útiles.

\sphinxAtStartPar
A continuación se muestran los 6 lenguajes de programación más usados en el desarrollo web. HTML, CSS y JavaScript suelen pertenecer al Front\sphinxhyphen{}end mientras que PHP, JAVA y Python suelen pertenecer al Back\sphinxhyphen{}end.
\begin{itemize}
\item {} 
\sphinxAtStartPar
\sphinxstylestrong{\sphinxhref{https://developer.mozilla.org/es/docs/Web/HTML}{HTML}}: Es uno de los lenguajes de marcado más importantes que se usa en el Front\sphinxhyphen{}end de un sitio. Su escritura ayuda a dar estructura y organización al contenido de una página web, a través de una acomodación tipo árbol. Se configura por medio de etiquetas o hipertextos que permiten que los sitios web se encuentren en los motores de búsqueda.

\item {} 
\sphinxAtStartPar
\sphinxstylestrong{\sphinxhref{https://developer.mozilla.org/es/docs/Web/CSS}{CSS}}: El nombre extendido de CSS es Cascading Style Sheets, en español significa hojas de estilo en cascada. Este es un lenguaje que trabaja en perfecta armonía con el HTML en el Front\sphinxhyphen{}end. Para los programadores web es una herramienta muy útil para especificar el aspecto y la posición de los elementos en el sitio.

\item {} 
\sphinxAtStartPar
\sphinxstylestrong{\sphinxhref{https://developer.mozilla.org/es/docs/Web/JavaScript}{Javascript}}: Uno de los lenguajes más apreciados es JavaScript, ya que con él es muy fácil crear sitios interactivos y dinámicos (como animaciones, formularios, juegos, galerías, botones, etc.), los cuales son muy demandados hoy en día. Su código también se refleja en el Front\sphinxhyphen{}end. Se basa en objetos que se pueden acomodar y reutilizar de forma sencilla.

\item {} 
\sphinxAtStartPar
\sphinxstylestrong{\sphinxhref{https://www.php.net}{PHP}}: En el lado del Back\sphinxhyphen{}end se tiene el lenguaje PHP, uno de los pioneros de la transición de sitios fijos a sitios interactivos. Es uno de los lenguajes más utilizados por ser de código abierto, con casi 30 años de trayectoria; también es considerado uno de los más extensos que hay. Puede ser incrustado en el HTML sin ningún problema.

\item {} 
\sphinxAtStartPar
\sphinxstylestrong{\sphinxhref{https://www.java.com/es}{Java}}: Una de las grandes ventajas es que se escribe una sola vez y se puede ejecutar en distintos dispositivos y sistemas operativos, gracias a su máquina virtual (JVM). Corresponde al Back\sphinxhyphen{}end de una aplicación web. Tiene un código eficiente, memoria automática y detección de errores oportuna.

\item {} 
\sphinxAtStartPar
\sphinxstylestrong{\sphinxhref{https://www.python.org}{Python}}: Es uno de los lenguajes de desarrollo web más innovadores que hay hasta el momento, debido a su característica multiparadigma, que es capaz de adaptarse a varios estilos de programación y crear aplicaciones de cualquier tipo. Es de código abierto y su escritura es muy parecida al lenguaje humano. Forma parte del Back\sphinxhyphen{}end de un sitio.

\end{itemize}


\bigskip\hrule\bigskip



\subsubsection{Software de aplicación}
\label{\detokenize{1_guias/10programacion/G1021:software-de-aplicacion}}
\sphinxAtStartPar
Dentro del software de aplicación se encuentran varios ejemplos relacionados con áreas específicas como diseño, ofimática, gráfico, financiero, musical, etc.
\begin{itemize}
\item {} 
\sphinxAtStartPar
\sphinxstylestrong{\sphinxhref{https://inkscape.org/}{Inkscape}}: Es un editor de gráficos vectoriales de código abierto, gratuito y disponible para múltiples sistemas operativos. Permite crear y editar imágenes vectoriales de alta calidad, ideal para ilustraciones, logotipos, diagramas y más.

\item {} 
\sphinxAtStartPar
\sphinxstylestrong{\sphinxhref{https://www.gimp.org/}{Gimp}}: Es un editor de imágenes de código abierto, gratuito y multiplataforma. Ofrece una amplia gama de herramientas para manipular, retocar y crear imágenes rasterizadas, similar a Photoshop pero sin costo.

\item {} 
\sphinxAtStartPar
\sphinxstylestrong{\sphinxhref{https://www.blender.org/}{Blender}}: Es un software de modelado 3D, animación y renderizado gratuito y de código abierto. Es una herramienta potente y versátil utilizada para crear gráficos 3D, animaciones, juegos y más.

\item {} 
\sphinxAtStartPar
\sphinxstylestrong{\sphinxhref{https://www.audacityteam.org}{Audacity}}: Audacity es el software gratuito más popular del mundo para grabar y editar audio. Así que, si esta produciendo música, un podcast o simplemente experimentando con audio, Audacity es para usted.

\item {} 
\sphinxAtStartPar
\sphinxstylestrong{\sphinxhref{https://godotengine.org/}{Godot}}: Es un motor de juego de código abierto, gratuito y multiplataforma. Proporciona un entorno completo para desarrollar videojuegos 2D y 3D, con un editor visual intuitivo y una gran comunidad de apoyo.

\end{itemize}


\bigskip\hrule\bigskip



\subsection{A02: Taller (+0,8)}
\label{\detokenize{1_guias/10programacion/G1021:a02-taller-0-8}}
\sphinxAtStartPar
\DUrole{sd-sphinx-override}{\DUrole{sd-badge}{\DUrole{sd-bg-warning}{\DUrole{sd-bg-text-warning}{Actividad práctica}}}}
\begin{enumerate}
\sphinxsetlistlabels{\arabic}{enumi}{enumii}{}{.}%
\item {} 
\sphinxAtStartPar
Responda las siguientes preguntas en su cuaderno:
\begin{itemize}
\item {} 
\sphinxAtStartPar
¿Que tiene que ver la \sphinxstylestrong{programación} con el \sphinxstylestrong{desarrollo de software}?

\item {} 
\sphinxAtStartPar
¿Qué es \sphinxstylestrong{código fuente}?

\end{itemize}

\item {} 
\sphinxAtStartPar
Indique los diferentes \sphinxstylestrong{tipos de software} y escriba 5 ejemplos por cada tipo.

\item {} 
\sphinxAtStartPar
Realice una descripción en su cuaderno sobre el \sphinxstylestrong{desarrollo de software}.

\item {} 
\sphinxAtStartPar
COPIE la siguiente tabla en su cuaderno y escriba las tareas o actividades que realiza un programador y un desarrollador de software:

\end{enumerate}


\begin{savenotes}\sphinxattablestart
\sphinxthistablewithglobalstyle
\centering
\begin{tabulary}{\linewidth}[t]{TTT}
\sphinxtoprule

\sphinxAtStartPar

&\sphinxstyletheadfamily 
\sphinxAtStartPar
Desarrollador de software
&\sphinxstyletheadfamily 
\sphinxAtStartPar
Programador
\\
\sphinxmidrule
\sphinxtableatstartofbodyhook
\sphinxAtStartPar
1.
&
\sphinxAtStartPar

&
\sphinxAtStartPar

\\
\sphinxhline
\sphinxAtStartPar
2.
&
\sphinxAtStartPar

&
\sphinxAtStartPar

\\
\sphinxhline
\sphinxAtStartPar
3.
&
\sphinxAtStartPar

&
\sphinxAtStartPar

\\
\sphinxhline
\sphinxAtStartPar
4.
&
\sphinxAtStartPar

&
\sphinxAtStartPar

\\
\sphinxhline
\sphinxAtStartPar
5.
&
\sphinxAtStartPar

&
\sphinxAtStartPar

\\
\sphinxhline
\sphinxAtStartPar
6.
&
\sphinxAtStartPar

&
\sphinxAtStartPar

\\
\sphinxbottomrule
\end{tabulary}
\sphinxtableafterendhook\par
\sphinxattableend\end{savenotes}
\begin{enumerate}
\sphinxsetlistlabels{\arabic}{enumi}{enumii}{}{.}%
\setcounter{enumi}{4}
\item {} 
\sphinxAtStartPar
COPIE la siguiente tabla en su cuaderno y escriba el nombre de las fases y las actividades que se realizan en cada fase durante el proceso de \sphinxstylestrong{desarrollo de software}:

\end{enumerate}


\begin{savenotes}\sphinxattablestart
\sphinxthistablewithglobalstyle
\centering
\begin{tabulary}{\linewidth}[t]{TTT}
\sphinxtoprule

\sphinxAtStartPar

&\sphinxstyletheadfamily 
\sphinxAtStartPar
Nombre
&\sphinxstyletheadfamily 
\sphinxAtStartPar
Actividades
\\
\sphinxmidrule
\sphinxtableatstartofbodyhook
\sphinxAtStartPar
FASE 1.
&
\sphinxAtStartPar

&
\sphinxAtStartPar

\\
\sphinxhline
\sphinxAtStartPar
FASE 2.
&
\sphinxAtStartPar

&
\sphinxAtStartPar

\\
\sphinxhline
\sphinxAtStartPar
FASE 3.
&
\sphinxAtStartPar

&
\sphinxAtStartPar

\\
\sphinxhline
\sphinxAtStartPar
FASE 4.
&
\sphinxAtStartPar

&
\sphinxAtStartPar

\\
\sphinxhline
\sphinxAtStartPar
FASE 5.
&
\sphinxAtStartPar

&
\sphinxAtStartPar

\\
\sphinxhline
\sphinxAtStartPar
FASE 6.
&
\sphinxAtStartPar

&
\sphinxAtStartPar

\\
\sphinxbottomrule
\end{tabulary}
\sphinxtableafterendhook\par
\sphinxattableend\end{savenotes}
\begin{enumerate}
\sphinxsetlistlabels{\arabic}{enumi}{enumii}{}{.}%
\setcounter{enumi}{5}
\item {} 
\sphinxAtStartPar
COPIE la siguiente \sphinxstylestrong{sopa de letras} en su cuaderno y encuentre 14 términos de tecnologías que se usan para desarrollar software. Posteriormente escriba su significado.

\end{enumerate}


\begin{savenotes}\sphinxattablestart
\sphinxthistablewithglobalstyle
\centering
\begin{tabulary}{\linewidth}[t]{T}
\sphinxtoprule
\sphinxtableatstartofbodyhook

\begin{sphinxfigure-in-table}
\centering

\noindent\sphinxincludegraphics[width=550\sphinxpxdimen]{{image036}.png}
\end{sphinxfigure-in-table}\relax
\\
\sphinxbottomrule
\end{tabulary}
\sphinxtableafterendhook\par
\sphinxattableend\end{savenotes}

\sphinxAtStartPar
Ejemplo:
\begin{quote}

\sphinxAtStartPar
\sphinxstylestrong{SQL\_SERVER:} Sistema de gestión de base de datos relacional de Microsoft.
\end{quote}
\begin{enumerate}
\sphinxsetlistlabels{\arabic}{enumi}{enumii}{}{.}%
\setcounter{enumi}{6}
\item {} 
\sphinxAtStartPar
En su cuaderno dibuje la siguiente \sphinxstyleemphasis{nube de términos} de los \sphinxstylestrong{STACK tecnológicos del desarrollo web}, una el STACK con sus componentes por medio de líneas.

\end{enumerate}


\begin{savenotes}\sphinxattablestart
\sphinxthistablewithglobalstyle
\centering
\begin{tabulary}{\linewidth}[t]{T}
\sphinxtoprule
\sphinxtableatstartofbodyhook

\begin{sphinxfigure-in-table}
\centering

\noindent\sphinxincludegraphics[width=500\sphinxpxdimen]{{image046}.png}
\end{sphinxfigure-in-table}\relax
\\
\sphinxbottomrule
\end{tabulary}
\sphinxtableafterendhook\par
\sphinxattableend\end{savenotes}
\begin{enumerate}
\sphinxsetlistlabels{\arabic}{enumi}{enumii}{}{.}%
\setcounter{enumi}{7}
\item {} 
\sphinxAtStartPar
Para los siguientes casos, en su cuaderno elegir el lenguaje de programación más apropiado. Especificar el por qué.
\begin{quote}

\sphinxAtStartPar
\sphinxstylestrong{CASO 01}: Un programador está trabajando en el Back\sphinxhyphen{}end de una plataforma educativa muy grande. Necesita un lenguaje pionero, de código abierto, con casi tres décadas de trayectoria y que sea sumamente extenso para gestionar los datos. Además, requiere que este código se pueda incrustar directamente dentro del código HTML sin generar conflictos.
\begin{itemize}
\item {} 
\sphinxAtStartPar
¿Qué lenguaje del Back\sphinxhyphen{}end cumple con tener esta larga trayectoria y la capacidad de incrustarse en HTML?

\end{itemize}
\end{quote}
\begin{quote}

\sphinxAtStartPar
\sphinxstylestrong{CASO 02}: Diego ha creado un sitio web para su club de videojuegos. La información ya está en la pantalla, pero se ve muy simple: el fondo es completamente blanco, los textos están desalineados y las imágenes se ven desordenadas. Él quiere cambiar el aspecto del sitio, definir los colores de los fondos, elegir los tipos de letra y especificar la posición exacta de cada elemento en la pantalla para que se vea genial.
\begin{itemize}
\item {} 
\sphinxAtStartPar
¿Qué lenguaje trabaja en armonía con la estructura existente y es la herramienta perfecta para cambiar el aspecto visual y la posición de los elementos?

\end{itemize}
\end{quote}
\begin{quote}

\sphinxAtStartPar
\sphinxstylestrong{CASO 03}: Valentina está diseñando una página web y quiere que deje de ser un sitio fijo y aburrido. Su objetivo es que, cuando los usuarios hagan clic en un botón, aparezca una animación, las galerías de fotos se muevan solas y los formularios validen los datos del usuario de forma inmediata con ventanas interactivas. Busque un lenguaje del lado del Front\sphinxhyphen{}end que se base en objetos reutilizables.
\begin{itemize}
\item {} 
\sphinxAtStartPar
¿Qué lenguaje es el más apreciado para crear este tipo de sitios interactivos, dinámicos y muy demandados hoy en día?

\end{itemize}
\end{quote}
\begin{quote}

\sphinxAtStartPar
\sphinxstylestrong{CASO 04}: El equipo de robótica del instituto quiere crear una aplicación web innovadora en el Back\sphinxhyphen{}end. Como el equipo está aprendiendo, buscan un lenguaje de código abierto cuya escritura sea muy parecida al lenguaje humano (fácil de leer y escribir). Además, como el proyecto cambiará constantemente, necesitan que sea “multiparadigma”, es decir, que se adapte a varios estilos de programación.
\begin{itemize}
\item {} 
\sphinxAtStartPar
¿Cuál es este innovador lenguaje multiparadigma que facilita el desarrollo gracias a su similitud con el lenguaje humano?

\end{itemize}
\end{quote}
\begin{quote}

\sphinxAtStartPar
\sphinxstylestrong{CASO 05}: Sofía quiere crear una página web para la cafetería del colegio. Ella ya tiene todo el texto, los precios de los productos y las imágenes que quiere mostrar, pero necesita organizar esta información de forma ordenada (por ejemplo, los títulos arriba, el menú en el centro y los datos de contacto al final) para que los motores de búsqueda como Google puedan encontrar la página fácilmente.
\begin{itemize}
\item {} 
\sphinxAtStartPar
¿Qué lenguaje debe utilizar Sofía para dar esta estructura tipo árbol y organizar el contenido de su sitio web mediante etiquetas?

\end{itemize}
\end{quote}
\begin{quote}

\sphinxAtStartPar
\sphinxstylestrong{CASO 06}: Carlos está desarrollando el sistema de inscripciones del colegio (una aplicación web en el Back\sphinxhyphen{}end). El director le ha pedido que el sistema funcione de manera eficiente tanto en las computadoras de la sala de informática como en servidores con otros sistemas operativos. Carlos necesita un lenguaje que se escriba una sola vez y corra en cualquier lado gracias a una máquina virtual, además de que detecte errores a tiempo y maneje la memoria automáticamente.
\begin{itemize}
\item {} 
\sphinxAtStartPar
¿Qué lenguaje del Back\sphinxhyphen{}end es famoso por usar una máquina virtual para ejecutarse en distintos dispositivos?

\end{itemize}
\end{quote}

\item {} 
\sphinxAtStartPar
Para los siguientes \sphinxstylestrong{softwares de aplicación}, en su cuaderno realice una descripción de cada uno y dibuje su logotipo (buscar en INTERNET). Al mismo tiempo responda la siguiente pregunta: \sphinxstyleemphasis{¿el logotipo tiene que ver con el software?}
\begin{itemize}
\item {} 
\sphinxAtStartPar
Godot

\item {} 
\sphinxAtStartPar
Audacity

\item {} 
\sphinxAtStartPar
Blender

\item {} 
\sphinxAtStartPar
Gimp

\item {} 
\sphinxAtStartPar
Inkscape

\end{itemize}

\end{enumerate}


\bigskip\hrule\bigskip



\subsection{A03: Examen escrito (+4,0)}
\label{\detokenize{1_guias/10programacion/G1021:a03-examen-escrito-4-0}}
\sphinxAtStartPar
\DUrole{sd-sphinx-override}{\DUrole{sd-badge}{\DUrole{sd-bg-danger}{\DUrole{sd-bg-text-danger}{Transferencia}}}}


\begin{savenotes}\sphinxattablestart
\sphinxthistablewithglobalstyle
\centering
\begin{tabulary}{\linewidth}[t]{TT}
\sphinxtoprule
\sphinxtableatstartofbodyhook
\sphinxAtStartPar
\sphinxstylestrong{TIEMPO}
&
\sphinxAtStartPar
29 minutos
\\
\sphinxhline
\sphinxAtStartPar
\sphinxstylestrong{PREGUNTAS}
&
\sphinxAtStartPar
10
\\
\sphinxhline
\sphinxAtStartPar
\sphinxstylestrong{OBJETIVO}
&
\sphinxAtStartPar
Comprobar conocimientos adquiridos.
\\
\sphinxhline
\sphinxAtStartPar
\sphinxstylestrong{ACTIVIDAD}
&
\sphinxAtStartPar
Responda las siguientes preguntas. Cuadernos y celulares guardados. Solo se permite hoja, lapiz y borrador.
\\
\sphinxbottomrule
\end{tabulary}
\sphinxtableafterendhook\par
\sphinxattableend\end{savenotes}

\begin{sphinxadmonition}{note}{Posibles preguntas}
\begin{enumerate}
\sphinxsetlistlabels{\arabic}{enumi}{enumii}{}{.}%
\item {} 
\sphinxAtStartPar
¿Qué es la \sphinxstylestrong{programación} y cuál es el objetivo principal de escribir código fuente?

\item {} 
\sphinxAtStartPar
Además de escribir código, ¿qué otras dos tareas importantes debe realizar un \sphinxstylestrong{programador} para mantener la calidad del software?

\item {} 
\sphinxAtStartPar
¿Cómo se define el \sphinxstylestrong{software} en contraste con el \sphinxstylestrong{hardware} de un sistema informático?

\item {} 
\sphinxAtStartPar
Explique la diferencia entre el “\sphinxstylestrong{software de sistema}” y el “\sphinxstylestrong{software de aplicación}”, dando un ejemplo de cada uno.

\item {} 
\sphinxAtStartPar
¿Qué es el “\sphinxstylestrong{software embebido}” y en qué tipo de dispositivos suele encontrarse?

\item {} 
\sphinxAtStartPar
¿Qué se entiende por “\sphinxstylestrong{desarrollo de software}” y por qué se considera un proceso más holístico que la simple programación? (holístico es un enfoque donde todos los componentes se ven integrados y hacen parte de un mismo ente).

\item {} 
\sphinxAtStartPar
¿Cuáles son las diferencias fundamentales entre el rol de un \sphinxstylestrong{programador} y el de un \sphinxstylestrong{desarrollador de software}?

\item {} 
\sphinxAtStartPar
Describa brevemente la “\sphinxstylestrong{FASE 01: Análisis de requisitos}” y por qué es fundamental para el éxito del proyecto.

\item {} 
\sphinxAtStartPar
En la \sphinxstylestrong{fase de diseño}, ¿qué herramientas se suelen utilizar para la validación temprana y creación de prototipos?

\item {} 
\sphinxAtStartPar
¿Por qué es importante realizar pruebas de rendimiento durante el proceso de \sphinxstylestrong{desarrollo de software}?

\item {} 
\sphinxAtStartPar
¿En qué consiste la “\sphinxstylestrong{FASE 05: Implementación}” y qué ocurre si ya existe una versión previa del software?

\item {} 
\sphinxAtStartPar
Explique por qué el \sphinxstylestrong{mantenimiento} y \sphinxstylestrong{soporte} son cruciales para evitar que el software quede desactualizado.

\item {} 
\sphinxAtStartPar
¿Qué es un \sphinxstylestrong{sistema operativo} y cómo actúa como intermediario entre el usuario y el hardware?

\item {} 
\sphinxAtStartPar
Explique la función de una \sphinxstylestrong{base de datos} y mencione dos ejemplos de sistemas que utilizan \sphinxstylestrong{tablas} para almacenar información.

\item {} 
\sphinxAtStartPar
¿Qué diferencia principal existe entre un entorno de ejecución (como \sphinxstylestrong{Node.js}) y un \sphinxstylestrong{framework de aplicación web}?

\item {} 
\sphinxAtStartPar
¿Qué es un \sphinxstylestrong{IDE} (Entorno de Desarrollo Integrado) y qué ventajas ofrece para el trabajo de un programador?

\item {} 
\sphinxAtStartPar
Explique la diferencia entre las responsabilidades del “\sphinxstylestrong{Front\sphinxhyphen{}end}” y el “\sphinxstylestrong{Back\sphinxhyphen{}end}” en el desarrollo web.

\item {} 
\sphinxAtStartPar
¿Por qué se dice que el lenguaje \sphinxstylestrong{HTML} ayuda a dar una “\sphinxstylestrong{acomodación tipo árbol}” al contenido de un sitio web?

\item {} 
\sphinxAtStartPar
¿Cuáles son las características distintivas que hacen que el lenguaje \sphinxstylestrong{Python} sea considerado innovador para el \sphinxstylestrong{desarrollo web}?

\item {} 
\sphinxAtStartPar
¿Qué aporta el lenguaje \sphinxstylestrong{JavaScript} a un sitio web en comparación con el uso exclusivo de \sphinxstylestrong{HTML} y \sphinxstylestrong{CSS}?

\end{enumerate}
\end{sphinxadmonition}

\sphinxstepscope


\section{Guía 1022: Datos, Información y Software (P1)}
\label{\detokenize{1_guias/10programacion/G1022:guia-1022-datos-informacion-y-software-p1}}\label{\detokenize{1_guias/10programacion/G1022::doc}}

\begin{savenotes}\sphinxattablestart
\sphinxthistablewithglobalstyle
\centering
\begin{tabular}[t]{\X{35}{100}\X{65}{100}}
\sphinxtoprule
\sphinxtableatstartofbodyhook

\begin{savenotes}\sphinxattablestart
\sphinxthistablewithglobalstyle
\centering
\begin{tabulary}{\linewidth}[t]{T}
\sphinxtoprule
\sphinxtableatstartofbodyhook
\sphinxAtStartPar
\sphinxincludegraphics{{escudo}.png}
\\
\sphinxbottomrule
\end{tabulary}
\sphinxtableafterendhook\par
\sphinxattableend\end{savenotes}
&

\begin{savenotes}\sphinxattablestart
\sphinxthistablewithglobalstyle
\raggedright
\begin{tabular}[t]{*{1}{\X{1}{1}}}
\sphinxtoprule
\sphinxtableatstartofbodyhook
\sphinxAtStartPar
\sphinxstylestrong{INSTITUCIÓN EDUCATIVA}: John F. Kennedy
\begin{itemize}
\item {} 
\sphinxAtStartPar
\sphinxstylestrong{DOCENTE}: Julian Camilo Fonseca Romero

\item {} 
\sphinxAtStartPar
\sphinxstylestrong{ASIGNATURA}: Programación

\item {} 
\sphinxAtStartPar
\sphinxstylestrong{GRADO}: ONCE

\item {} 
\sphinxAtStartPar
\sphinxstylestrong{PERIODO}: 01

\item {} 
\sphinxAtStartPar
\sphinxstylestrong{TIEMPO}: 10 semanas (2 horas / semana)

\item {} 
\sphinxAtStartPar
\sphinxstylestrong{TIEMPO DE TRABAJO}: 12/20 horas

\end{itemize}
\\
\sphinxbottomrule
\end{tabular}
\sphinxtableafterendhook\par
\sphinxattableend\end{savenotes}
\\
\sphinxbottomrule
\end{tabular}
\sphinxtableafterendhook\par
\sphinxattableend\end{savenotes}

\begin{sphinxuseclass}{sd-tab-set}
\begin{sphinxuseclass}{sd-tab-item}\subsubsection*{Objetivo}

\begin{sphinxuseclass}{sd-tab-content}
\sphinxAtStartPar
\DUrole{sd-sphinx-override}{\DUrole{sd-badge}{\DUrole{sd-bg-light}{\DUrole{sd-bg-text-light}{Exploración}}}}
\DUrole{sd-sphinx-override}{\DUrole{sd-badge}{\DUrole{sd-bg-info}{\DUrole{sd-bg-text-info}{Estructuración}}}}
\DUrole{sd-sphinx-override}{\DUrole{sd-badge}{\DUrole{sd-bg-warning}{\DUrole{sd-bg-text-warning}{Actividad práctica}}}}
\DUrole{sd-sphinx-override}{\DUrole{sd-badge}{\DUrole{sd-bg-danger}{\DUrole{sd-bg-text-danger}{Transferencia}}}}
\begin{itemize}
\item {} 
\sphinxAtStartPar
\sphinxstylestrong{Temática}: Concepto de hardware, software, diferencias entre software y programa, información, datos y jerarquía de datos.

\item {} 
\sphinxAtStartPar
\sphinxstylestrong{Objetivo de aprendizaje}: Entender los fundamentos de arquitectura y lógica computacional.

\end{itemize}

\end{sphinxuseclass}
\end{sphinxuseclass}
\begin{sphinxuseclass}{sd-tab-item}\subsubsection*{Orientaciones curriculares}

\begin{sphinxuseclass}{sd-tab-content}\begin{itemize}
\item {} 
\sphinxAtStartPar
\sphinxstylestrong{Componente}: Naturaleza y Evolución de la T\&I.

\item {} 
\sphinxAtStartPar
\sphinxstylestrong{Competencia}: Relaciono saberes, conocimientos tecnológicos e informáticos con los conocimientos de otras disciplinas.

\item {} 
\sphinxAtStartPar
\sphinxstylestrong{Evidencia de aprendizaje}: Comprendo los principios y conocimientos tecnológicos e informáticos que hacen posible el funcionamiento de productos tecnológicos actuales.

\end{itemize}

\end{sphinxuseclass}
\end{sphinxuseclass}
\begin{sphinxuseclass}{sd-tab-item}\subsubsection*{Criterios de evaluación}

\begin{sphinxuseclass}{sd-tab-content}\begin{itemize}
\item {} 
\sphinxAtStartPar
Actividades desarrolladas en su totalidad.

\item {} 
\sphinxAtStartPar
Participación en clase.

\item {} 
\sphinxAtStartPar
Explicación clara de conceptos.

\item {} 
\sphinxAtStartPar
Argumentación de ideas.

\item {} 
\sphinxAtStartPar
Cumplimiento de las reglas del aula.

\end{itemize}

\end{sphinxuseclass}
\end{sphinxuseclass}
\begin{sphinxuseclass}{sd-tab-item}\subsubsection*{Recursos (0)}

\begin{sphinxuseclass}{sd-tab-content}\begin{itemize}
\item {} 
\sphinxAtStartPar
La presente guía no contiene recursos.

\end{itemize}

\end{sphinxuseclass}
\end{sphinxuseclass}
\end{sphinxuseclass}

\subsection{Mensaje del autor}
\label{\detokenize{1_guias/10programacion/G1022:mensaje-del-autor}}\begin{quote}

\sphinxAtStartPar
“En esta guía, el estudiante recuerda conceptos básicos de informática para aplicar más adelante en el aprendizaje de la programación.”

\begin{flushright}
---Camilo Fonseca
\end{flushright}
\end{quote}


\bigskip\hrule\bigskip



\subsection{A00: Introducción (+0,1)}
\label{\detokenize{1_guias/10programacion/G1022:a00-introduccion-0-1}}
\sphinxAtStartPar
\DUrole{sd-sphinx-override}{\DUrole{sd-badge}{\DUrole{sd-bg-light}{\DUrole{sd-bg-text-light}{Exploración}}}}
\begin{enumerate}
\sphinxsetlistlabels{\arabic}{enumi}{enumii}{}{.}%
\item {} 
\sphinxAtStartPar
TRANSCRIBA en su cuaderno el objetivo, las orientaciones curriculares y los criterios de evaluación de la presente guía.

\item {} 
\sphinxAtStartPar
LEA el \sphinxstylestrong{mensaje del autor} y en su cuaderno responda:
\begin{itemize}
\item {} 
\sphinxAtStartPar
¿Es posible entender la tecnología actual sin conocer como funciona el hardware y como se organizan los datos?

\item {} 
\sphinxAtStartPar
¿Qué es para usted un hardware?

\item {} 
\sphinxAtStartPar
¿Qué es para usted un software?

\item {} 
\sphinxAtStartPar
¿Qué es para usted información y en que se diferencia con los datos?

\item {} 
\sphinxAtStartPar
¿Cómo se relacionan los datos con el hardware y el software?

\end{itemize}

\end{enumerate}


\bigskip\hrule\bigskip



\subsection{A01: Información y datos (+0,1)}
\label{\detokenize{1_guias/10programacion/G1022:a01-informacion-y-datos-0-1}}
\sphinxAtStartPar
\DUrole{sd-sphinx-override}{\DUrole{sd-badge}{\DUrole{sd-bg-info}{\DUrole{sd-bg-text-info}{Estructuración}}}}
\begin{enumerate}
\sphinxsetlistlabels{\arabic}{enumi}{enumii}{}{.}%
\item {} 
\sphinxAtStartPar
Lea los siguientes conceptos relacionados con el tema de información y datos.

\end{enumerate}


\bigskip\hrule\bigskip



\subsubsection{¿Qué son los datos?}
\label{\detokenize{1_guias/10programacion/G1022:que-son-los-datos}}
\sphinxAtStartPar
Los datos son valores en bruto, hechos o símbolos que por sí solos no tienen significado. Por ejemplo:
\begin{itemize}
\item {} 
\sphinxAtStartPar
Números sueltos.

\item {} 
\sphinxAtStartPar
Caracteres sin contexto.

\item {} 
\sphinxAtStartPar
Valores sin procesar.

\end{itemize}


\bigskip\hrule\bigskip



\subsubsection{¿Qué es la información?}
\label{\detokenize{1_guias/10programacion/G1022:que-es-la-informacion}}
\sphinxAtStartPar
La información es el resultado de procesar, organizar e interpretar los datos de manera que tengan sentido y sean útiles. La información:
\begin{itemize}
\item {} 
\sphinxAtStartPar
Tiene significado y contexto.

\item {} 
\sphinxAtStartPar
Es comprensible para el usuario.

\item {} 
\sphinxAtStartPar
Ayuda en la toma de decisiones.

\end{itemize}

\sphinxAtStartPar
Por ejemplo:


\begin{savenotes}\sphinxattablestart
\sphinxthistablewithglobalstyle
\centering
\begin{tabulary}{\linewidth}[t]{TT}
\sphinxtoprule
\sphinxstyletheadfamily 
\sphinxAtStartPar
Datos
&\sphinxstyletheadfamily 
\sphinxAtStartPar
Información
\\
\sphinxmidrule
\sphinxtableatstartofbodyhook
\sphinxAtStartPar
37,5
&
\sphinxAtStartPar
La temperatura del paciente es 37.5°C, lo cual indica fiebre
\\
\sphinxhline
\sphinxAtStartPar
4,5
&
\sphinxAtStartPar
La nota final del estudiante es 4.5
\\
\sphinxbottomrule
\end{tabulary}
\sphinxtableafterendhook\par
\sphinxattableend\end{savenotes}

\sphinxAtStartPar
La transformación de datos en información es una parte fundamental del procesamiento computacional, ya que permite convertir valores sin sentido en conocimiento útil y aplicable.


\bigskip\hrule\bigskip



\subsubsection{Valores cuantitativos y cualitativos}
\label{\detokenize{1_guias/10programacion/G1022:valores-cuantitativos-y-cualitativos}}
\sphinxAtStartPar
%
\begin{footnote}[1]\sphinxAtStartFootnote
Trejos Buriticá, O. I. (2017). Lógica de programación. Ediciones de la U.
%
\end{footnote} Los valores cualitativos y cuantitativos son dos tipos diferentes de información:
\begin{itemize}
\item {} 
\sphinxAtStartPar
\sphinxstylestrong{Valores cualitativos:} Describen cualidades o características que no se pueden medir numéricamente. Por ejemplo:
\begin{itemize}
\item {} 
\sphinxAtStartPar
Color de ojos

\item {} 
\sphinxAtStartPar
Estado de ánimo

\item {} 
\sphinxAtStartPar
Nivel de satisfacción (excelente, bueno, regular)

\end{itemize}

\item {} 
\sphinxAtStartPar
\sphinxstylestrong{Valores cuantitativos:} Son aquellos que se pueden medir y expresar en números. Por ejemplo:
\begin{itemize}
\item {} 
\sphinxAtStartPar
Edad: 12 años

\item {} 
\sphinxAtStartPar
Peso: 99 Kg

\item {} 
\sphinxAtStartPar
Temperatura: 40°C

\item {} 
\sphinxAtStartPar
Calificaciones numéricas: 3,0

\end{itemize}

\end{itemize}

\sphinxAtStartPar
La principal diferencia es que los valores cualitativos describen atributos o cualidades, mientras que los cuantitativos expresan cantidades medibles.


\bigskip\hrule\bigskip



\subsubsection{Jerarquía de datos}
\label{\detokenize{1_guias/10programacion/G1022:jerarquia-de-datos}}
\sphinxAtStartPar
%
\begin{footnote}[2]\sphinxAtStartFootnote
Deitel, P. J., \& Deitel, H. (2014). C++: cómo programar (A. V. Romero Elizondo, Trans.). Pearson Educación.
%
\end{footnote} La jerarquía de datos es la organización lógica de los datos en niveles sucesivos de complejidad, desde la unidad más pequeña hasta la más completa. En informática, se utiliza para entender cómo se estructura la información dentro de un sistema o base de datos.

\sphinxAtStartPar
Los datos en la computadora se clasifican de MENOR a MAYOR en el siguiente ranking:
\begin{enumerate}
\sphinxsetlistlabels{\arabic}{enumi}{enumii}{}{.}%
\item {} 
\sphinxAtStartPar
\sphinxstylestrong{bit:} unidad lógica puede ser \sphinxcode{\sphinxupquote{1}} o \sphinxcode{\sphinxupquote{0}}.

\item {} 
\sphinxAtStartPar
\sphinxstylestrong{BYTE:} Son 8 bits. Ejemplo: \sphinxcode{\sphinxupquote{10010010}}

\item {} 
\sphinxAtStartPar
\sphinxstylestrong{Caracteres:} Representaciones de Bytes. Ejemplo: \sphinxstylestrong{\sphinxcode{\sphinxupquote{A}}}: \sphinxcode{\sphinxupquote{10001001}}, \sphinxstylestrong{\sphinxcode{\sphinxupquote{B}}:} \sphinxcode{\sphinxupquote{10001001}}

\item {} 
\sphinxAtStartPar
\sphinxstylestrong{Dato:} Se componen de representaciones de Caracteres. Ejemplo: \sphinxcode{\sphinxupquote{'CAMILO'}}, \sphinxcode{\sphinxupquote{'ROJO'}}, \sphinxcode{\sphinxupquote{'FRIO'}}, \sphinxcode{\sphinxupquote{23}} , \sphinxcode{\sphinxupquote{0.0999}}, \sphinxguilabel{True},

\item {} 
\sphinxAtStartPar
\sphinxstylestrong{Campo:} Nombre que se le coloca a un dato para identificarlo. Ejemplo: \sphinxcode{\sphinxupquote{nombre}}, \sphinxcode{\sphinxupquote{color}}, \sphinxcode{\sphinxupquote{clima}}, \sphinxcode{\sphinxupquote{edad}}, \sphinxcode{\sphinxupquote{temperatura}}, \sphinxcode{\sphinxupquote{estado}}

\item {} 
\sphinxAtStartPar
\sphinxstylestrong{Registro:} Conjunto de campos que conforman un solo objeto. Ejemplo: El observador de un estudiante, la hoja técnica de una moto.

\item {} 
\sphinxAtStartPar
\sphinxstylestrong{Tablas:} Conjunto de registros que tienen la misma estructura y que pertenecen a una organización. Ejemplo:
\begin{quote}

\sphinxAtStartPar
“La \sphinxstylestrong{tabla} de los \sphinxstyleemphasis{estudiantes} de un colegio”.
\end{quote}

\item {} 
\sphinxAtStartPar
\sphinxstylestrong{Bases de datos:} Conjunto de tablas relacionadas técnicamente y organizadas. Ejemplo:
\begin{quote}

\sphinxAtStartPar
“La \sphinxstylestrong{base de datos} donde están relacionadas las tablas de los estudiantes junto con las tablas de las notas”.
\end{quote}

\end{enumerate}


\bigskip\hrule\bigskip



\subsubsection{Hardware}
\label{\detokenize{1_guias/10programacion/G1022:hardware}}
\sphinxAtStartPar
\sphinxfootnotemark[2] El hardware es el conjunto de componentes de un sistema informático o computacional que se puede ver o tocar. Por ejemplo:
\begin{itemize}
\item {} 
\sphinxAtStartPar
\sphinxstylestrong{Dispositivos de entrada:} Teclados, mouse, micrófonos.

\item {} 
\sphinxAtStartPar
\sphinxstylestrong{Dispositivos de salida:} Monitores, impresoras, bafles.

\item {} 
\sphinxAtStartPar
\sphinxstylestrong{Procesamiento:} CPU.

\item {} 
\sphinxAtStartPar
\sphinxstylestrong{Memoria:} memoria RAM, memoria CACHE.

\item {} 
\sphinxAtStartPar
\sphinxstylestrong{Almacenamiento:} Discos duros, memorias USB.

\end{itemize}


\paragraph{La memoria RAM (acceso rápido a los datos)}
\label{\detokenize{1_guias/10programacion/G1022:la-memoria-ram-acceso-rapido-a-los-datos}}
\sphinxAtStartPar
La \sphinxstylestrong{memoria RAM} (\sphinxstyleemphasis{Random Access Memory o Memoria de Acceso Aleatorio}) es un componente fundamental del hardware que:
\begin{itemize}
\item {} 
\sphinxAtStartPar
Almacena temporalmente los datos y programas que el computador está usando activamente.

\item {} 
\sphinxAtStartPar
Permite acceder rápidamente a la información mientras el equipo está encendido.

\item {} 
\sphinxAtStartPar
Se borra cuando se apaga el computador (memoria volátil).

\end{itemize}

\sphinxAtStartPar
Características importantes de la RAM:
\begin{itemize}
\item {} 
\sphinxAtStartPar
A mayor cantidad de RAM, más programas pueden ejecutarse simultáneamente.

\item {} 
\sphinxAtStartPar
Su velocidad afecta directamente el rendimiento del sistema.

\item {} 
\sphinxAtStartPar
Es diferente al almacenamiento permanente (disco duro) porque es temporal.

\end{itemize}


\paragraph{La CPU (procesa los datos)}
\label{\detokenize{1_guias/10programacion/G1022:la-cpu-procesa-los-datos}}
\sphinxAtStartPar
La \sphinxstylestrong{CPU (Unidad Central de Procesamiento)} es el “cerebro” del computador. Es el componente principal que se encarga de:
\begin{itemize}
\item {} 
\sphinxAtStartPar
Procesar todas las instrucciones y cálculos del sistema

\item {} 
\sphinxAtStartPar
Coordinar el funcionamiento de los demás componentes

\item {} 
\sphinxAtStartPar
Ejecutar los programas y aplicaciones

\end{itemize}

\sphinxAtStartPar
La \sphinxstylestrong{CPU} está compuesta principalmente por:
\begin{itemize}
\item {} 
\sphinxAtStartPar
\sphinxstylestrong{Unidad de Control:} Coordina y dirige todas las operaciones

\item {} 
\sphinxAtStartPar
\sphinxstylestrong{Unidad Aritmético\sphinxhyphen{}Lógica (ALU):} Realiza cálculos matemáticos y operaciones lógicas

\item {} 
\sphinxAtStartPar
\sphinxstylestrong{Registros:} Pequeñas unidades de memoria de alta velocidad

\end{itemize}

\sphinxAtStartPar
La velocidad de la \sphinxstylestrong{CPU} se mide en Hertz (Hz) y determina qué tan rápido puede procesar las instrucciones. A mayor velocidad, mejor rendimiento del sistema.


\bigskip\hrule\bigskip



\subsubsection{Software}
\label{\detokenize{1_guias/10programacion/G1022:software}}
\sphinxAtStartPar
\sphinxfootnotemark[2] Un \sphinxstylestrong{Software} es el conjunto de \sphinxstylestrong{programas}, bases de datos, archivos y documentación que permiten realizar diferentes tareas en un sistema computacional.

\sphinxAtStartPar
Existen dos tipos de software: Software de sistema y software de aplicación.
\begin{itemize}
\item {} 
\sphinxAtStartPar
\sphinxstylestrong{Software de sistema:} Ejecutan todas las funciones básicas de una computadora (por ejemplo encender la máquina, mostrar imagen en pantalla, monitorear temperatura, comunicarse con los diferentes periféricos, etc.). El \sphinxstylestrong{software de sistema} controla e intercambia datos con el hardware. Un ejemplo de software de sistema es el sistema operativo Windows. Inclusive los controladores o drivers también hacen parte del software de sistema.

\item {} 
\sphinxAtStartPar
\sphinxstylestrong{Software de aplicación:} Permite al usuario trabajar en una tarea específica (procesar texto, edición de video, edición de imágenes, navegación en INTERNET, desarrollo de software, etc.). Por lo general, el software de aplicación trabaja en función del software de sistema.

\end{itemize}

\sphinxAtStartPar
A diferencia del hardware que es físico y tangible, el software no se puede tocar \sphinxhyphen{} solo puede verse su funcionamiento en la pantalla.


\paragraph{Software Vs Programa}
\label{\detokenize{1_guias/10programacion/G1022:software-vs-programa}}
\sphinxAtStartPar
Un programa es un \sphinxstylestrong{conjunto de instrucciones} que le dicen a una \sphinxstylestrong{computadora} COMO HACER una tarea específica. Las instrucciones están escritas en código y son diseñadas por programadores con conocimiento en un lenguaje de programación.

\sphinxAtStartPar
A diferencia de un \sphinxstylestrong{algoritmo}, las instrucciones de un programa \sphinxstylestrong{NO pueden ser escritas en lenguaje natural}.
\begin{quote}

\sphinxAtStartPar
“Aunque se usa el término \sphinxstylestrong{software} como sinónimo, la realidad es que un \sphinxstylestrong{programa} es solo una parte del \sphinxstylestrong{software}”.
\end{quote}


\bigskip\hrule\bigskip



\subsection{A02: Taller (+0,8)}
\label{\detokenize{1_guias/10programacion/G1022:a02-taller-0-8}}
\sphinxAtStartPar
\DUrole{sd-sphinx-override}{\DUrole{sd-badge}{\DUrole{sd-bg-warning}{\DUrole{sd-bg-text-warning}{Actividad práctica}}}}
\begin{enumerate}
\sphinxsetlistlabels{\arabic}{enumi}{enumii}{}{.}%
\item {} 
\sphinxAtStartPar
En su cuaderno responda las siguientes preguntas:
\begin{itemize}
\item {} 
\sphinxAtStartPar
¿Qué son los datos?

\item {} 
\sphinxAtStartPar
¿Qué es información?

\item {} 
\sphinxAtStartPar
¿Cuál es la diferencia entre un valor cualitativo y un valor cuantitativo?

\item {} 
\sphinxAtStartPar
¿Qué es una jerarquía de datos?

\end{itemize}

\item {} 
\sphinxAtStartPar
En su cuaderno describa cada uno de los componentes de la jerarquía de datos y de un ejemplo de cada uno.

\item {} 
\sphinxAtStartPar
Responda las siguientes preguntas en su cuaderno:
\begin{itemize}
\item {} 
\sphinxAtStartPar
¿Qué es tener sentido?

\item {} 
\sphinxAtStartPar
¿Por qué es importante tener “sentido” para diferenciar entre datos e información?

\end{itemize}

\item {} 
\sphinxAtStartPar
Dibuje una tabla en su cuaderno y en la primera columna escriba 15 ejemplos de datos. En la segunda columna convierta estos datos en información. En la tercera columna indique si la información es un valor cualitativo o cuantitativo.


\begin{savenotes}\sphinxattablestart
\sphinxthistablewithglobalstyle
\centering
\begin{tabulary}{\linewidth}[t]{TTT}
\sphinxtoprule
\sphinxstyletheadfamily 
\sphinxAtStartPar
Datos
&\sphinxstyletheadfamily 
\sphinxAtStartPar
Información
&\sphinxstyletheadfamily 
\sphinxAtStartPar
Valor
\\
\sphinxmidrule
\sphinxtableatstartofbodyhook
\sphinxAtStartPar
Rojo
&
\sphinxAtStartPar
“Semáforo en estado rojo, carros no pueden pasar”
&
\sphinxAtStartPar
cualitativo
\\
\sphinxhline
\sphinxAtStartPar
16
&
\sphinxAtStartPar
“La edad del estudiante es 16 años, por tanto es menor de edad”
&
\sphinxAtStartPar
cuantitativo
\\
\sphinxbottomrule
\end{tabulary}
\sphinxtableafterendhook\par
\sphinxattableend\end{savenotes}

\item {} 
\sphinxAtStartPar
Dibuje los siguientes círculos, ingrese ejemplos de datos según su jerarquía (ver la sección de \sphinxstylestrong{jerarquía de datos}).

\end{enumerate}


\begin{savenotes}\sphinxattablestart
\sphinxthistablewithglobalstyle
\centering
\begin{tabulary}{\linewidth}[t]{T}
\sphinxtoprule
\sphinxtableatstartofbodyhook

\begin{sphinxfigure-in-table}
\centering
\capstart
\noindent\sphinxincludegraphics[width=400\sphinxpxdimen]{{image0110}.png}
\sphinxfigcaption{\sphinxstylestrong{Diagrama:} Jerarquía de datos}\label{\detokenize{1_guias/10programacion/G1022:id7}}\end{sphinxfigure-in-table}\relax
\\
\sphinxbottomrule
\end{tabulary}
\sphinxtableafterendhook\par
\sphinxattableend\end{savenotes}
\begin{enumerate}
\sphinxsetlistlabels{\arabic}{enumi}{enumii}{}{.}%
\setcounter{enumi}{5}
\item {} 
\sphinxAtStartPar
En su cuaderno responda:
\begin{itemize}
\item {} 
\sphinxAtStartPar
¿Qué es hardware?

\item {} 
\sphinxAtStartPar
¿Cómo se puede clasificar el hardware de un sistema computacional?

\end{itemize}

\item {} 
\sphinxAtStartPar
En su cuaderno mencione diferentes ejemplos de hardware y clasifíquelos.

\item {} 
\sphinxAtStartPar
En su cuaderno responda:
\begin{itemize}
\item {} 
\sphinxAtStartPar
¿Qué es la memoria RAM?

\item {} 
\sphinxAtStartPar
¿Cuales son las características de la memoria RAM?

\end{itemize}

\item {} 
\sphinxAtStartPar
En INTERNET consulte 3 fabricantes de memorias RAM y escríbalos en su cuaderno. Tenga en cuenta lo siguiente:
\begin{itemize}
\item {} 
\sphinxAtStartPar
Mencionar la nacionalidad del fabricante.

\item {} 
\sphinxAtStartPar
Mencionar los clientes del fabricante.

\end{itemize}

\item {} 
\sphinxAtStartPar
En su cuaderno responda:
\begin{itemize}
\item {} 
\sphinxAtStartPar
¿Qué es la CPU?

\item {} 
\sphinxAtStartPar
¿Qué son los Hertz(Hz)?

\item {} 
\sphinxAtStartPar
La CPU tiene 3 componentes. En su cuaderno complete el siguiente gráfico:


\begin{savenotes}\sphinxattablestart
\sphinxthistablewithglobalstyle
\centering
\begin{tabulary}{\linewidth}[t]{T}
\sphinxtoprule
\sphinxtableatstartofbodyhook

\begin{sphinxfigure-in-table}
\centering
\capstart
\noindent\sphinxincludegraphics[width=400\sphinxpxdimen]{{image028}.png}
\sphinxfigcaption{\sphinxstylestrong{Diagrama:} Componentes CPU}\label{\detokenize{1_guias/10programacion/G1022:id8}}\end{sphinxfigure-in-table}\relax
\\
\sphinxbottomrule
\end{tabulary}
\sphinxtableafterendhook\par
\sphinxattableend\end{savenotes}

\end{itemize}

\item {} 
\sphinxAtStartPar
Teniendo en cuenta la siguiente imagen, responda en su cuaderno la siguiente pregunta:
\begin{itemize}
\item {} 
\sphinxAtStartPar
¿En que se diferencian los dispositivos de entrada con los dispositivos de salida?

\end{itemize}


\begin{savenotes}\sphinxattablestart
\sphinxthistablewithglobalstyle
\centering
\begin{tabulary}{\linewidth}[t]{T}
\sphinxtoprule
\sphinxtableatstartofbodyhook

\begin{sphinxfigure-in-table}
\centering
\capstart
\noindent\sphinxincludegraphics[width=360\sphinxpxdimen]{{image037}.png}
\sphinxfigcaption{\sphinxstylestrong{Diagrama:} CPU}\label{\detokenize{1_guias/10programacion/G1022:id9}}\end{sphinxfigure-in-table}\relax
\\
\sphinxbottomrule
\end{tabulary}
\sphinxtableafterendhook\par
\sphinxattableend\end{savenotes}

\item {} 
\sphinxAtStartPar
En media hoja de cuaderno, explique por que la CPU y la memoria RAM son los dispositivos más importantes para dar instrucciones.

\item {} 
\sphinxAtStartPar
Respecto a la memoria \sphinxstylestrong{RAM}, en las siguientes frases diga si es falso o verdadero y explique el por qué:
\begin{itemize}
\item {} 
\sphinxAtStartPar
La memoria \sphinxstylestrong{RAM} almacena infinitamente los datos y programas \_\_\_

\item {} 
\sphinxAtStartPar
Para acceder velozmente a la información en el equipo usamos la memoria \sphinxstylestrong{RAM} \_\_\_

\item {} 
\sphinxAtStartPar
La memoriam \sphinxstylestrong{RAM} no es volátil \_\_\_

\item {} 
\sphinxAtStartPar
La velocidad de trabajo de un computador, depende exclusivamente de la velocidad de la memoria \sphinxstylestrong{RAM} \_\_\_

\end{itemize}

\item {} 
\sphinxAtStartPar
En su cuaderno responda la siguientes preguntas:
\begin{itemize}
\item {} 
\sphinxAtStartPar
¿Qué es el software?

\item {} 
\sphinxAtStartPar
¿Cómo se clasifica el software?

\item {} 
\sphinxAtStartPar
¿Cuál es la diferencia entre un software y un programa?

\end{itemize}

\item {} 
\sphinxAtStartPar
Clasifique el siguiente software en la siguiente tabla:
\begin{itemize}
\item {} 
\sphinxAtStartPar
Windows 11, WORD, cpa000.exe, Linux UBUNTU, LibreOffice Calc, UNIX, hii.bin, ANDROID, FIREFOX, LibreOffice Writer, windowSys.exe, IOS, CHROME, FIFA, Photoshop, BIOS

\end{itemize}


\begin{savenotes}\sphinxattablestart
\sphinxthistablewithglobalstyle
\centering
\begin{tabulary}{\linewidth}[t]{TT}
\sphinxtoprule
\sphinxstyletheadfamily 
\sphinxAtStartPar
Software de sistema
&\sphinxstyletheadfamily 
\sphinxAtStartPar
Software de aplicación
\\
\sphinxmidrule
\sphinxtableatstartofbodyhook
\sphinxAtStartPar
\_
&
\sphinxAtStartPar

\\
\sphinxhline
\sphinxAtStartPar
\_
&
\sphinxAtStartPar

\\
\sphinxhline
\sphinxAtStartPar
\_
&
\sphinxAtStartPar

\\
\sphinxhline
\sphinxAtStartPar
\_
&
\sphinxAtStartPar

\\
\sphinxhline
\sphinxAtStartPar
\_
&
\sphinxAtStartPar

\\
\sphinxbottomrule
\end{tabulary}
\sphinxtableafterendhook\par
\sphinxattableend\end{savenotes}

\end{enumerate}


\bigskip\hrule\bigskip



\subsection{A03: Examen final (+4,0)}
\label{\detokenize{1_guias/10programacion/G1022:a03-examen-final-4-0}}
\sphinxAtStartPar
\DUrole{sd-sphinx-override}{\DUrole{sd-badge}{\DUrole{sd-bg-danger}{\DUrole{sd-bg-text-danger}{Transferencia}}}}


\begin{savenotes}\sphinxattablestart
\sphinxthistablewithglobalstyle
\centering
\begin{tabulary}{\linewidth}[t]{TT}
\sphinxtoprule
\sphinxtableatstartofbodyhook
\sphinxAtStartPar
\sphinxstylestrong{TIEMPO}
&
\sphinxAtStartPar
29 minutos
\\
\sphinxhline
\sphinxAtStartPar
\sphinxstylestrong{PREGUNTAS}
&
\sphinxAtStartPar
10
\\
\sphinxhline
\sphinxAtStartPar
\sphinxstylestrong{OBJETIVO}
&
\sphinxAtStartPar
Comprobar conocimientos adquiridos.
\\
\sphinxhline
\sphinxAtStartPar
\sphinxstylestrong{ACTIVIDAD}
&
\sphinxAtStartPar
Responda las siguientes preguntas. Cuadernos y celulares guardados. Solo se permite hoja, lapiz y borrador.
\\
\sphinxbottomrule
\end{tabulary}
\sphinxtableafterendhook\par
\sphinxattableend\end{savenotes}

\begin{sphinxadmonition}{note}{Posibles preguntas}
\begin{enumerate}
\sphinxsetlistlabels{\arabic}{enumi}{enumii}{}{.}%
\item {} 
\sphinxAtStartPar
Explique la diferencia entre \sphinxstylestrong{datos} e \sphinxstylestrong{información} y de un ejemplo de cada uno.

\item {} 
\sphinxAtStartPar
¿Por qué es importante que los \sphinxstylestrong{datos} tengan un contexto para convertirse en \sphinxstylestrong{información útil}?

\item {} 
\sphinxAtStartPar
Describa qué son los \sphinxstylestrong{valores cualitativos} y \sphinxstylestrong{cuantitativos} y proporcione un ejemplo de cada uno.

\item {} 
\sphinxAtStartPar
Explique la \sphinxstylestrong{jerarquía de datos} desde el \sphinxstylestrong{bit} hasta la \sphinxstylestrong{base de datos}, indicando qué representa cada nivel.

\item {} 
\sphinxAtStartPar
¿Por qué un \sphinxstylestrong{byte} está compuesto por 8 bits y qué importancia tiene esta unidad en la informática?

\item {} 
\sphinxAtStartPar
Define qué es un \sphinxstylestrong{registro} y cómo se relaciona con los \sphinxstylestrong{campos} y \sphinxstylestrong{archivos} en la organización de datos.

\item {} 
\sphinxAtStartPar
Explique qué es el \sphinxstylestrong{hardware} y mencione al menos \sphinxstylestrong{cinco componentes físicos} que lo conforman.

\item {} 
\sphinxAtStartPar
Describa la función principal de la \sphinxstylestrong{memoria RAM} y por qué se considera una memoria volátil.

\item {} 
\sphinxAtStartPar
¿Cuál es el papel de la CPU en un sistema computacional y cuáles son sus componentes principales?

\item {} 
\sphinxAtStartPar
Explique la diferencia entre \sphinxstylestrong{software de sistema} y \sphinxstylestrong{software de aplicación} con ejemplos.

\item {} 
\sphinxAtStartPar
¿Qué es un \sphinxstylestrong{programa} y cómo se relaciona con el \sphinxstylestrong{software en general}?

\item {} 
\sphinxAtStartPar
Describa cómo la \sphinxstylestrong{información procesada} por un computador puede influir en la toma de decisiones.

\item {} 
\sphinxAtStartPar
¿Por qué es fundamental entender la \sphinxstylestrong{arquitectura de un computador} para programar eficientemente?

\item {} 
\sphinxAtStartPar
Explique cómo los datos se almacenan y organizan en un \sphinxstylestrong{sistema informático} para facilitar su acceso.

\item {} 
\sphinxAtStartPar
Describa el proceso que sigue un computador para transformar \sphinxstylestrong{datos} en \sphinxstylestrong{información}.

\item {} 
\sphinxAtStartPar
¿Qué papel juegan los \sphinxstylestrong{Hertz} dentro de la \sphinxstylestrong{CPU} y cómo afecta el rendimiento del sistema?

\item {} 
\sphinxAtStartPar
Explique la importancia de la \sphinxstylestrong{memoria RAM} en la ejecución de programas y qué sucede cuando se apaga el computador.

\item {} 
\sphinxAtStartPar
Describa cómo el \sphinxstylestrong{hardware} y el \sphinxstylestrong{software} trabajan juntos para realizar tareas computacionales.

\item {} 
\sphinxAtStartPar
¿Qué diferencias existen entre los distintos tipos de \sphinxstylestrong{memoria} en un computador y para qué se usa cada una?

\item {} 
\sphinxAtStartPar
Reflexione sobre cómo la evolución de la \sphinxstylestrong{arquitectura computacional} ha impactado en el desarrollo de tecnologías actuales.

\end{enumerate}
\end{sphinxadmonition}


\bigskip\hrule\bigskip


\sphinxstepscope


\section{Guía 1023: Algoritmos (P2)}
\label{\detokenize{1_guias/10programacion/G1023:guia-1023-algoritmos-p2}}\label{\detokenize{1_guias/10programacion/G1023::doc}}

\begin{savenotes}\sphinxattablestart
\sphinxthistablewithglobalstyle
\centering
\begin{tabular}[t]{\X{35}{100}\X{65}{100}}
\sphinxtoprule
\sphinxtableatstartofbodyhook

\begin{savenotes}\sphinxattablestart
\sphinxthistablewithglobalstyle
\centering
\begin{tabulary}{\linewidth}[t]{T}
\sphinxtoprule
\sphinxtableatstartofbodyhook
\sphinxAtStartPar
\sphinxincludegraphics{{escudo}.png}
\\
\sphinxbottomrule
\end{tabulary}
\sphinxtableafterendhook\par
\sphinxattableend\end{savenotes}
&

\begin{savenotes}\sphinxattablestart
\sphinxthistablewithglobalstyle
\raggedright
\begin{tabular}[t]{*{1}{\X{1}{1}}}
\sphinxtoprule
\sphinxtableatstartofbodyhook
\sphinxAtStartPar
\sphinxstylestrong{INSTITUCIÓN EDUCATIVA}: John F. Kennedy
\begin{itemize}
\item {} 
\sphinxAtStartPar
\sphinxstylestrong{DOCENTE}: Julian Camilo Fonseca Romero

\item {} 
\sphinxAtStartPar
\sphinxstylestrong{ASIGNATURA}: Programación

\item {} 
\sphinxAtStartPar
\sphinxstylestrong{GRADO}: ONCE

\item {} 
\sphinxAtStartPar
\sphinxstylestrong{PERIODO}: 02

\item {} 
\sphinxAtStartPar
\sphinxstylestrong{TIEMPO}: 10 semanas (2 horas / semana)

\item {} 
\sphinxAtStartPar
\sphinxstylestrong{TIEMPO DE TRABAJO}: 20/20 horas

\end{itemize}
\\
\sphinxbottomrule
\end{tabular}
\sphinxtableafterendhook\par
\sphinxattableend\end{savenotes}
\\
\sphinxbottomrule
\end{tabular}
\sphinxtableafterendhook\par
\sphinxattableend\end{savenotes}

\begin{sphinxuseclass}{sd-tab-set}
\begin{sphinxuseclass}{sd-tab-item}\subsubsection*{Objetivo}

\begin{sphinxuseclass}{sd-tab-content}
\sphinxAtStartPar
\DUrole{sd-sphinx-override}{\DUrole{sd-badge}{\DUrole{sd-bg-light}{\DUrole{sd-bg-text-light}{Exploración}}}}
\DUrole{sd-sphinx-override}{\DUrole{sd-badge}{\DUrole{sd-bg-info}{\DUrole{sd-bg-text-info}{Estructuración}}}}
\DUrole{sd-sphinx-override}{\DUrole{sd-badge}{\DUrole{sd-bg-warning}{\DUrole{sd-bg-text-warning}{Actividad práctica}}}}
\DUrole{sd-sphinx-override}{\DUrole{sd-badge}{\DUrole{sd-bg-danger}{\DUrole{sd-bg-text-danger}{Transferencia}}}}
\begin{itemize}
\item {} 
\sphinxAtStartPar
\sphinxstylestrong{Temática}: Desarrollo de algoritmos.

\item {} 
\sphinxAtStartPar
\sphinxstylestrong{Objetivo de aprendizaje}: Desarrollar la capacidad de abstracción y lógica algorítmica mediante la representación gráfica de procesos, permitiéndole modelar soluciones secuenciales y ordenadas a problemas de la vida cotidiana.

\end{itemize}

\end{sphinxuseclass}
\end{sphinxuseclass}
\begin{sphinxuseclass}{sd-tab-item}\subsubsection*{Orientaciones curriculares}

\begin{sphinxuseclass}{sd-tab-content}\begin{itemize}
\item {} 
\sphinxAtStartPar
\sphinxstylestrong{Componente}: Solución de problemas con T\&I.

\item {} 
\sphinxAtStartPar
\sphinxstylestrong{Competencia}: Propongo innovaciones tecnológicas e informáticas para la solución de problemas dando cumplimiento a restricciones, condiciones y especificaciones técnicas y contextuales.

\item {} 
\sphinxAtStartPar
\sphinxstylestrong{Evidencia de aprendizaje}: Propongo, analizo y comparo diferentes soluciones a un mismo problema de la tecnología o la informática, explicando su origen, ventajas y dificultades.

\end{itemize}

\end{sphinxuseclass}
\end{sphinxuseclass}
\begin{sphinxuseclass}{sd-tab-item}\subsubsection*{Criterios de evaluación}

\begin{sphinxuseclass}{sd-tab-content}\begin{itemize}
\item {} 
\sphinxAtStartPar
Actividades desarrolladas en su totalidad.

\item {} 
\sphinxAtStartPar
Participación en clase.

\item {} 
\sphinxAtStartPar
Explicación clara de conceptos.

\item {} 
\sphinxAtStartPar
Argumentación de ideas.

\item {} 
\sphinxAtStartPar
Cumplimiento de las reglas del aula.

\end{itemize}

\end{sphinxuseclass}
\end{sphinxuseclass}
\begin{sphinxuseclass}{sd-tab-item}\subsubsection*{Recursos (0)}

\begin{sphinxuseclass}{sd-tab-content}\begin{itemize}
\item {} 
\sphinxAtStartPar
La presente guía no contiene recursos.

\end{itemize}

\end{sphinxuseclass}
\end{sphinxuseclass}
\end{sphinxuseclass}

\subsection{Mensaje del autor}
\label{\detokenize{1_guias/10programacion/G1023:mensaje-del-autor}}\begin{quote}

\sphinxAtStartPar
“Entender cómo funciona el mundo actual requiere aprender a pensar de forma organizada. Esta guía invita a usar el desarrollo de algoritmos como una herramienta para abstraer y modelar soluciones. Al visualizar sus ideas, usted podrá comparar distintas rutas lógicas para un mismo problema, evaluando sus beneficios y desafíos”.

\begin{flushright}
---Camilo Fonseca
\end{flushright}
\end{quote}


\bigskip\hrule\bigskip



\subsection{A00: Introducción (+0,1)}
\label{\detokenize{1_guias/10programacion/G1023:a00-introduccion-0-1}}
\sphinxAtStartPar
\DUrole{sd-sphinx-override}{\DUrole{sd-badge}{\DUrole{sd-bg-light}{\DUrole{sd-bg-text-light}{Exploración}}}}
\begin{enumerate}
\sphinxsetlistlabels{\arabic}{enumi}{enumii}{}{.}%
\item {} 
\sphinxAtStartPar
TRANSCRIBA en su cuaderno el objetivo, las orientaciones curriculares y los criterios de evaluación de la presente guía.

\item {} 
\sphinxAtStartPar
LEA el \sphinxstylestrong{mensaje del autor} y en su cuaderno con sus palabras responda:
\begin{itemize}
\item {} 
\sphinxAtStartPar
¿Qué es abstracción?

\item {} 
\sphinxAtStartPar
¿Qué es modelar una solución?

\end{itemize}

\end{enumerate}


\bigskip\hrule\bigskip



\subsection{A01: Algoritmos (+0,1)}
\label{\detokenize{1_guias/10programacion/G1023:a01-algoritmos-0-1}}
\sphinxAtStartPar
\DUrole{sd-sphinx-override}{\DUrole{sd-badge}{\DUrole{sd-bg-info}{\DUrole{sd-bg-text-info}{Estructuración}}}}
\begin{enumerate}
\sphinxsetlistlabels{\arabic}{enumi}{enumii}{}{.}%
\item {} 
\sphinxAtStartPar
Lea los siguientes conceptos:

\end{enumerate}


\bigskip\hrule\bigskip



\subsubsection{¿Qué es un algoritmo?}
\label{\detokenize{1_guias/10programacion/G1023:que-es-un-algoritmo}}
\sphinxAtStartPar
%
\begin{footnote}[2]\sphinxAtStartFootnote
Trejos Buriticá, O. I. (2017). Lógica de programación. Ediciones de la U.
%
\end{footnote} Un algoritmo es un conjunto de \sphinxstylestrong{instrucciones} que le dicen a un ser humano COMO REALIZAR una tarea específica, dichas instrucciones pueden estar escritas en código o lenguaje natural.

\sphinxAtStartPar
Un algoritmo puede ser expresado con palabras o números.


\bigskip\hrule\bigskip



\subsubsection{Tipos de algoritmos}
\label{\detokenize{1_guias/10programacion/G1023:tipos-de-algoritmos}}
\sphinxAtStartPar
Un algoritmo es cualquier actividad o evento con inicio y final, que tenga una serie de pasos lógicos para lograr su ejecución. Y estos suelen ser representados mediante diagramas de flujo.

\sphinxAtStartPar
Existen dos tipos de algoritmos:
\begin{itemize}
\item {} 
\sphinxAtStartPar
Algoritmos cualitativos (informales)

\item {} 
\sphinxAtStartPar
Algoritmos cuantitativos (formales)

\end{itemize}


\paragraph{Algoritmos cualitativos (informales)}
\label{\detokenize{1_guias/10programacion/G1023:algoritmos-cualitativos-informales}}
\sphinxAtStartPar
No requieren de cálculos numéricos para su ejecución. En su lugar, se deben ejecutar secuencias lógicas. Por ejemplo: las instrucciones para ensamblar un artefacto.


\paragraph{Algoritmos cuantitativos (formales)}
\label{\detokenize{1_guias/10programacion/G1023:algoritmos-cuantitativos-formales}}
\sphinxAtStartPar
Requieren de cálculos numéricos (por ejemplo la solución de una ecuación). Si el cálculo numérico puede ser procesado por una persona, entonces es un \sphinxstylestrong{algoritmo cuantitativo no computacional}.

\sphinxAtStartPar
Si los cálculos numéricos son muy complejos entonces requieren poder computacional. A esta clase de algoritmos se les dice que son \sphinxstylestrong{algoritmos cuantitativos computacionales}.
\begin{itemize}
\item {} 
\sphinxAtStartPar
\sphinxstylestrong{Algoritmo cuantitativo no computacional:} Estos algoritmos requieren de operaciones numéricas simples, que un ser humano puede resolver sin necesidad de usar un sistema computacional. Por ejemplo una receta de cocina, un cálculo numérico sencillo o las instrucciones para llegar a un lugar.

\item {} 
\sphinxAtStartPar
\sphinxstylestrong{Algoritmo cuantitativo computacional:} Estos algoritmos requieren de operaciones numéricas complejas que deben resolverse con dispositivos de cálculo, como una computadora o calculadora. Ecuaciones de mucha complejidad o códigos que solo pueden ser interpretados por una máquina, son ejemplos de este tipo de algoritmo.

\end{itemize}


\bigskip\hrule\bigskip



\subsubsection{¿Qué es un diagrama de flujo?}
\label{\detokenize{1_guias/10programacion/G1023:que-es-un-diagrama-de-flujo}}
\sphinxAtStartPar
%
\begin{footnote}[1]\sphinxAtStartFootnote
Deitel, P. J., \& Deitel, H. (2014). C++: cómo programar (A. V. Romero Elizondo, Trans.). Pearson Educación.
%
\end{footnote} Un diagrama de flujo es un \sphinxstylestrong{esquema gráfico} donde se muestra el paso a paso para cumplir un objetivo.

\sphinxAtStartPar
Los diagramas de flujo usados en informática, son útiles porque ayudan a los programadores a organizar y visualizar el código haciendo más fácil encontrar y corregir errores.


\bigskip\hrule\bigskip



\subsubsection{Símbolos de diagramas de flujo}
\label{\detokenize{1_guias/10programacion/G1023:simbolos-de-diagramas-de-flujo}}
\sphinxAtStartPar
A continuación se explican los símbolos que componen los diagramas de flujo del estándar (ANSI/ISO):

\begin{sphinxadmonition}{note}{Operación}


\begin{savenotes}\sphinxattablestart
\sphinxthistablewithglobalstyle
\centering
\begin{tabulary}{\linewidth}[t]{T}
\sphinxtoprule
\sphinxtableatstartofbodyhook

\begin{sphinxfigure-in-table}
\centering
\capstart
\noindent\sphinxincludegraphics[width=250\sphinxpxdimen]{{image0111}.png}
\sphinxfigcaption{Acción o Instrucción}\label{\detokenize{1_guias/10programacion/G1023:id5}}\end{sphinxfigure-in-table}\relax
\\
\sphinxbottomrule
\end{tabulary}
\sphinxtableafterendhook\par
\sphinxattableend\end{savenotes}

\sphinxAtStartPar
Se usa para describir cualquier acción o instrucción. En el interior se escribe una descripción de alguna acción o instrucción a ejecutar.

\begin{sphinxadmonition}{note}{¿Qué es una instrucción?}

\sphinxAtStartPar
Una \sphinxstylestrong{instrucción} es una orden o indicación que se da para realizar una acción específica. Puede estar presente en distintos ámbitos. Por ejemplo:
\begin{itemize}
\item {} 
\sphinxAtStartPar
\sphinxstylestrong{Una instrucción de educación:} Indicaciones que un profesor da a los estudiantes sobre cómo realizar una tarea o actividad.

\item {} 
\sphinxAtStartPar
\sphinxstylestrong{Instrucciones de programación:} Son comandos que le indican a una computadora qué hacer. Las instrucciones son parte del código fuente de un programa.

\item {} 
\sphinxAtStartPar
\sphinxstylestrong{Instrucciones de cocina:} Son pasos a seguir para una receta.

\item {} 
\sphinxAtStartPar
\sphinxstylestrong{Instrucciones técnicas:} Guías que explican cómo usar o ensamblar un producto. También aplica en el software.

\end{itemize}
\begin{quote}

\sphinxAtStartPar
Reglas para dar buenas instrucciones
\begin{itemize}
\item {} 
\sphinxAtStartPar
\sphinxstylestrong{CLARIDAD:} Ayudan a entender qué se espera hacer.

\item {} 
\sphinxAtStartPar
\sphinxstylestrong{EFICIENCIA:} Permiten realizar tareas de manera correcta y rápida.

\item {} 
\sphinxAtStartPar
\sphinxstylestrong{SEGURIDAD:} En ciertos contextos, como el uso de maquinaria, las instrucciones son vitales para evitar accidentes.

\end{itemize}
\end{quote}
\end{sphinxadmonition}
\end{sphinxadmonition}

\begin{sphinxadmonition}{note}{Límites de proceso}


\begin{savenotes}\sphinxattablestart
\sphinxthistablewithglobalstyle
\centering
\begin{tabulary}{\linewidth}[t]{T}
\sphinxtoprule
\sphinxtableatstartofbodyhook

\begin{sphinxfigure-in-table}
\centering
\capstart
\noindent\sphinxincludegraphics[width=250\sphinxpxdimen]{{image029}.png}
\sphinxfigcaption{Inicio \sphinxhyphen{} Fin}\label{\detokenize{1_guias/10programacion/G1023:id6}}\end{sphinxfigure-in-table}\relax
\\
\sphinxbottomrule
\end{tabulary}
\sphinxtableafterendhook\par
\sphinxattableend\end{savenotes}

\sphinxAtStartPar
Indica el inicio (círculo blanco) y final (círculo negro) del algoritmo.
\end{sphinxadmonition}

\begin{sphinxadmonition}{note}{Punto de decisión}


\begin{savenotes}\sphinxattablestart
\sphinxthistablewithglobalstyle
\centering
\begin{tabulary}{\linewidth}[t]{T}
\sphinxtoprule
\sphinxtableatstartofbodyhook

\begin{sphinxfigure-in-table}
\centering
\capstart
\noindent\sphinxincludegraphics[width=150\sphinxpxdimen]{{image038}.png}
\sphinxfigcaption{Condicional}\label{\detokenize{1_guias/10programacion/G1023:id7}}\end{sphinxfigure-in-table}\relax
\\
\sphinxbottomrule
\end{tabulary}
\sphinxtableafterendhook\par
\sphinxattableend\end{savenotes}

\sphinxAtStartPar
Denota que en este punto se toma una decisión. Las únicas posibles opciones son SI/NO.
\end{sphinxadmonition}

\begin{sphinxadmonition}{note}{Conector}


\begin{savenotes}\sphinxattablestart
\sphinxthistablewithglobalstyle
\centering
\begin{tabulary}{\linewidth}[t]{T}
\sphinxtoprule
\sphinxtableatstartofbodyhook

\begin{sphinxfigure-in-table}
\centering
\capstart
\noindent\sphinxincludegraphics[width=150\sphinxpxdimen]{{image047}.png}
\sphinxfigcaption{Conexión}\label{\detokenize{1_guias/10programacion/G1023:id8}}\end{sphinxfigure-in-table}\relax
\\
\sphinxbottomrule
\end{tabulary}
\sphinxtableafterendhook\par
\sphinxattableend\end{savenotes}

\sphinxAtStartPar
Señala que la salida de un proceso/algoritmo, puede ser la entrada de otro proceso.
\end{sphinxadmonition}

\begin{sphinxadmonition}{note}{Dirección de flujo}


\begin{savenotes}\sphinxattablestart
\sphinxthistablewithglobalstyle
\centering
\begin{tabulary}{\linewidth}[t]{T}
\sphinxtoprule
\sphinxtableatstartofbodyhook

\begin{sphinxfigure-in-table}
\centering
\capstart
\noindent\sphinxincludegraphics[width=150\sphinxpxdimen]{{image056}.png}
\sphinxfigcaption{Dirección}\label{\detokenize{1_guias/10programacion/G1023:id9}}\end{sphinxfigure-in-table}\relax
\\
\sphinxbottomrule
\end{tabulary}
\sphinxtableafterendhook\par
\sphinxattableend\end{savenotes}

\sphinxAtStartPar
Denota la dirección y el orden de los pasos del proceso/algoritmo.
\end{sphinxadmonition}

\begin{sphinxadmonition}{note}{Entrada manual}


\begin{savenotes}\sphinxattablestart
\sphinxthistablewithglobalstyle
\centering
\begin{tabulary}{\linewidth}[t]{T}
\sphinxtoprule
\sphinxtableatstartofbodyhook

\begin{sphinxfigure-in-table}
\centering
\capstart
\noindent\sphinxincludegraphics[width=150\sphinxpxdimen]{{image066}.png}
\sphinxfigcaption{Entrada}\label{\detokenize{1_guias/10programacion/G1023:id10}}\end{sphinxfigure-in-table}\relax
\\
\sphinxbottomrule
\end{tabulary}
\sphinxtableafterendhook\par
\sphinxattableend\end{savenotes}

\sphinxAtStartPar
Permite una entrada manual del usuario para ser procesada en el algoritmo.
\end{sphinxadmonition}

\begin{sphinxadmonition}{note}{Documento}


\begin{savenotes}\sphinxattablestart
\sphinxthistablewithglobalstyle
\centering
\begin{tabulary}{\linewidth}[t]{T}
\sphinxtoprule
\sphinxtableatstartofbodyhook

\begin{sphinxfigure-in-table}
\centering
\capstart
\noindent\sphinxincludegraphics[width=150\sphinxpxdimen]{{image076}.png}
\sphinxfigcaption{Documento}\label{\detokenize{1_guias/10programacion/G1023:id11}}\end{sphinxfigure-in-table}\relax
\\
\sphinxbottomrule
\end{tabulary}
\sphinxtableafterendhook\par
\sphinxattableend\end{savenotes}

\sphinxAtStartPar
Documento o registro generado en el proceso/algoritmo.
\end{sphinxadmonition}

\begin{sphinxadmonition}{note}{Base de datos}


\begin{savenotes}\sphinxattablestart
\sphinxthistablewithglobalstyle
\centering
\begin{tabulary}{\linewidth}[t]{T}
\sphinxtoprule
\sphinxtableatstartofbodyhook

\begin{sphinxfigure-in-table}
\centering
\capstart
\noindent\sphinxincludegraphics[width=120\sphinxpxdimen]{{image086}.png}
\sphinxfigcaption{Base de datos}\label{\detokenize{1_guias/10programacion/G1023:id12}}\end{sphinxfigure-in-table}\relax
\\
\sphinxbottomrule
\end{tabulary}
\sphinxtableafterendhook\par
\sphinxattableend\end{savenotes}

\sphinxAtStartPar
Conjunto de tablas donde se retiene la información. En espera que se cumplan otras condiciones para continuar con el proceso/algoritmo.
\end{sphinxadmonition}


\bigskip\hrule\bigskip



\subsubsection{Metodología para desarrollar un algoritmo}
\label{\detokenize{1_guias/10programacion/G1023:metodologia-para-desarrollar-un-algoritmo}}
\sphinxAtStartPar
Cuando se desarrolla un algoritmo, cualquier usuario podría ejecutarlo siempre y cuando cumpla unos requisitos iniciales.

\sphinxAtStartPar
Para desarrollar un algoritmo, se tienen las siguientes fases:

\begin{sphinxadmonition}{note}{FASE 01:  Análisis de requisitos (objetivo y condiciones)}

\sphinxAtStartPar
En esta fase se identifica lo siguiente:
\begin{itemize}
\item {} 
\sphinxAtStartPar
\sphinxstylestrong{Objetivo:} Identificar el objetivo claro del algoritmo.

\item {} 
\sphinxAtStartPar
\sphinxstylestrong{Condiciones iniciales:} Identificar los requisitos que debe cumplir el usuario para poder desarrollar el algoritmo.

\item {} 
\sphinxAtStartPar
\sphinxstylestrong{Condiciones de incumplimiento:} Identificar todos los condicionales del algoritmo.

\item {} 
\sphinxAtStartPar
\sphinxstylestrong{Condiciones de cumplimiento:} Identificar el estado del usuario cuando se cumple el algoritmo.

\end{itemize}
\end{sphinxadmonition}

\begin{sphinxadmonition}{note}{FASE 02: Diseño del algoritmo (diagrama de flujo)}

\sphinxAtStartPar
Realizar un diagrama de flujo donde muestre el paso a paso lógico por medio de las siguientes figuras:
\begin{itemize}
\item {} 
\sphinxAtStartPar
\sphinxstylestrong{Círculos:} INICIO y FIN

\item {} 
\sphinxAtStartPar
\sphinxstylestrong{Rectángulos:} Acciones

\item {} 
\sphinxAtStartPar
\sphinxstylestrong{Rombos:} Condicionales

\item {} 
\sphinxAtStartPar
\sphinxstylestrong{Rombos anidados:} Bucles

\item {} 
\sphinxAtStartPar
\sphinxstylestrong{Flechas:} Orden de secuencia.

\end{itemize}

\sphinxAtStartPar
Después de tener el diagrama de flujo, realizar un análisis.
\end{sphinxadmonition}

\begin{sphinxadmonition}{note}{FASE 03: Construcción del algoritmo (código fuente)}

\sphinxAtStartPar
Conversión del \sphinxstylestrong{diagrama de flujo} a \sphinxstylestrong{CÓDIGO FUENTE}.
El código fuente es cualquier código que usa un lenguaje de programación que entienda una PC.
En este caso sería PYTHON pero de momento se usará PSEUDOCÓDIGO.
Se debe identificar el algoritmo. La identificación debe tener:
\begin{itemize}
\item {} 
\sphinxAtStartPar
\sphinxstylestrong{TÍTULO:} “Nombre del algoritmo”

\item {} 
\sphinxAtStartPar
\sphinxstylestrong{AUTOR:} “Creador del algoritmo”

\item {} 
\sphinxAtStartPar
\sphinxstylestrong{FECHA:} “Fecha de creación o última modificación del algoritmo”

\end{itemize}
\end{sphinxadmonition}


\bigskip\hrule\bigskip



\subsubsection{EJEMPLO: Algoritmo para ir al colegio}
\label{\detokenize{1_guias/10programacion/G1023:ejemplo-algoritmo-para-ir-al-colegio}}
\sphinxAtStartPar
Desarrolle el siguiente ejercicio:

\begin{sphinxadmonition}{note}{EJERCICIO: Algoritmo para ir al colegio}

\sphinxAtStartPar
En su cuaderno desarrolle un \sphinxstylestrong{algoritmo que permita a un estudiante ir al colegio}. Tenga en cuenta lo siguiente:
\begin{quote}

\sphinxAtStartPar
El estudiante inicialmente está en su casa. Luego debe ponerse el uniforme e ir al garaje para intentar prender la moto.

\sphinxAtStartPar
Si la moto NO prende, entonces salir del garaje, luego salir de la casa para llegar al colegio a pie.

\sphinxAtStartPar
Si la moto SI prende, entonces abrir el garaje, luego salir del garaje para llegar al colegio en moto.
\end{quote}
\end{sphinxadmonition}

\sphinxAtStartPar
A continuación se muestra la solución del desarrollo del algoritmo.


\paragraph{FASE 01: Análisis de requisitos}
\label{\detokenize{1_guias/10programacion/G1023:fase-01-analisis-de-requisitos}}
\begin{sphinxadmonition}{note}{(Objetivo y condiciones)}
\begin{itemize}
\item {} 
\sphinxAtStartPar
OBJETIVO:
\begin{itemize}
\item {} 
\sphinxAtStartPar
Llegar al colegio.

\end{itemize}

\item {} 
\sphinxAtStartPar
CONDICIONES INICIALES:
\begin{itemize}
\item {} 
\sphinxAtStartPar
El estudiante debe estar en su casa

\item {} 
\sphinxAtStartPar
La casa del estudiante debe tener un garaje.

\item {} 
\sphinxAtStartPar
El estudiante no debe tener puesto el uniforme.

\item {} 
\sphinxAtStartPar
El estudiante debe tener una moto en el garaje.

\end{itemize}

\item {} 
\sphinxAtStartPar
CONDICIONES DE INCUMPLIMIENTO:
\begin{itemize}
\item {} 
\sphinxAtStartPar
La moto no prende.

\end{itemize}

\item {} 
\sphinxAtStartPar
CONDICIONES DE CUMPLIMIENTO:
\begin{itemize}
\item {} 
\sphinxAtStartPar
El estudiante está en el colegio.

\item {} 
\sphinxAtStartPar
El estudiante tiene el uniforme puesto.

\item {} 
\sphinxAtStartPar
La moto puede estar en el colegio o en el garaje.

\end{itemize}

\end{itemize}
\end{sphinxadmonition}


\paragraph{FASE 02: Diseño del algoritmo}
\label{\detokenize{1_guias/10programacion/G1023:fase-02-diseno-del-algoritmo}}

\begin{savenotes}\sphinxattablestart
\sphinxthistablewithglobalstyle
\centering
\begin{tabulary}{\linewidth}[t]{T}
\sphinxtoprule
\sphinxtableatstartofbodyhook

\begin{sphinxfigure-in-table}
\centering

\noindent\sphinxincludegraphics[width=400\sphinxpxdimen]{{image096}.png}
\end{sphinxfigure-in-table}\relax
\\
\sphinxbottomrule
\end{tabulary}
\sphinxtableafterendhook\par
\sphinxattableend\end{savenotes}


\paragraph{FASE 03: Construcción del algoritmo}
\label{\detokenize{1_guias/10programacion/G1023:fase-03-construccion-del-algoritmo}}
\begin{sphinxadmonition}{note}{(Código fuente)}


\begin{savenotes}\sphinxattablestart
\sphinxthistablewithglobalstyle
\centering
\begin{tabulary}{\linewidth}[t]{TT}
\sphinxtoprule
\sphinxtableatstartofbodyhook
\sphinxAtStartPar
TÍTULO
&
\sphinxAtStartPar
Algoritmo para que un estudiante pueda ir al colegio
\\
\sphinxhline
\sphinxAtStartPar
AUTOR
&
\sphinxAtStartPar
Camilo Fonseca
\\
\sphinxhline
\sphinxAtStartPar
FECHA
&
\sphinxAtStartPar
Domingo, 12 de Abril de 2026 \sphinxhyphen{} 03:38pm
\\
\sphinxbottomrule
\end{tabulary}
\sphinxtableafterendhook\par
\sphinxattableend\end{savenotes}

\begin{sphinxVerbatim}[commandchars=\\\{\},numbers=left,firstnumber=1,stepnumber=1]
 \PYG{n}{INICIO}

 \PYG{l+m+mf}{1.} \PYG{n}{Ponerse} \PYG{n}{el} \PYG{n}{uniforme}
 \PYG{l+m+mf}{3.} \PYG{n}{Ir} \PYG{n}{al} \PYG{n}{garaje}
 \PYG{l+m+mf}{4.} \PYG{n}{Intentar} \PYG{n}{prender} \PYG{n}{la} \PYG{n}{moto}
 \PYG{l+m+mf}{5.} \PYG{n}{Si} \PYG{p}{(}\PYG{n}{la} \PYG{n}{moto} \PYG{n}{SI} \PYG{n}{prende}\PYG{p}{)}
     \PYG{l+m+mf}{5.1} \PYG{n}{Abrir} \PYG{n}{el} \PYG{n}{garaje}
     \PYG{l+m+mf}{5.2} \PYG{n}{Sacar} \PYG{n}{la} \PYG{n}{moto}
     \PYG{l+m+mf}{5.3} \PYG{n}{Cerrar} \PYG{n}{el} \PYG{n}{garaje}
     \PYG{l+m+mf}{5.4} \PYG{n}{Conducir} \PYG{n}{hasta} \PYG{n}{el} \PYG{n}{colegio}
 \PYG{l+m+mf}{6.} \PYG{n}{Si} \PYG{p}{(}\PYG{n}{la} \PYG{n}{moto} \PYG{n}{NO} \PYG{n}{prende}\PYG{p}{)}
     \PYG{l+m+mf}{6.1} \PYG{n}{Salir} \PYG{k}{del} \PYG{n}{garage}
     \PYG{l+m+mf}{6.2} \PYG{n}{Salir} \PYG{n}{de} \PYG{n}{la} \PYG{n}{casa}
     \PYG{l+m+mf}{6.3} \PYG{n}{Cerrar} \PYG{n}{la} \PYG{n}{puerta}
     \PYG{l+m+mf}{6.4} \PYG{n}{Caminar} \PYG{n}{hasta} \PYG{n}{el} \PYG{n}{colegio}
 \PYG{l+m+mf}{7.} \PYG{n}{Llegar} \PYG{n}{al} \PYG{n}{colegio}
 
 \PYG{n}{FIN}
\end{sphinxVerbatim}
\end{sphinxadmonition}


\bigskip\hrule\bigskip



\subsection{A02: Taller (+0,8)}
\label{\detokenize{1_guias/10programacion/G1023:a02-taller-0-8}}
\sphinxAtStartPar
\DUrole{sd-sphinx-override}{\DUrole{sd-badge}{\DUrole{sd-bg-info}{\DUrole{sd-bg-text-info}{Estructuración}}}}
\begin{enumerate}
\sphinxsetlistlabels{\arabic}{enumi}{enumii}{}{.}%
\item {} 
\sphinxAtStartPar
Escriba las siguientes frases en su cuaderno e indique si es \sphinxstylestrong{falso} o \sphinxstylestrong{verdadero} e indique el por qué.
\begin{itemize}
\item {} 
\sphinxAtStartPar
Un algoritmo está compuesto por elementos informáticos que componen una secuencia organizacional únicamente. \_\_\_ (¿por qué?)

\item {} 
\sphinxAtStartPar
Un algoritmo son un conjunto de instrucciones que le dicen a una computadora o un ser humano como realizar una tarea específica, dichas instrucciones pueden estar escritas en código o lenguaje natural. \_\_\_ (¿por qué?)

\item {} 
\sphinxAtStartPar
Un algoritmo se define como cuadros de funcionamiento infinitos que permiten y tienen como propósito reducir la información. \_\_\_ (¿por qué?)

\end{itemize}

\item {} 
\sphinxAtStartPar
En su cuaderno relacione con líneas los siguientes cuadros según corresponda su significado:


\begin{savenotes}\sphinxattablestart
\sphinxthistablewithglobalstyle
\centering
\begin{tabulary}{\linewidth}[t]{T}
\sphinxtoprule
\sphinxtableatstartofbodyhook

\begin{sphinxfigure-in-table}
\centering
\capstart
\noindent\sphinxincludegraphics[width=500\sphinxpxdimen]{{image105}.png}
\sphinxfigcaption{Cuadro comparativo (Tipos de algoritmos)}\label{\detokenize{1_guias/10programacion/G1023:id13}}\end{sphinxfigure-in-table}\relax
\\
\sphinxbottomrule
\end{tabulary}
\sphinxtableafterendhook\par
\sphinxattableend\end{savenotes}

\item {} 
\sphinxAtStartPar
Escriba la siguiente tabla en su cuaderno y clasifique los algoritmos según corresponda:


\begin{savenotes}\sphinxattablestart
\sphinxthistablewithglobalstyle
\centering
\begin{tabulary}{\linewidth}[t]{TTT}
\sphinxtoprule
\sphinxstyletheadfamily 
\sphinxAtStartPar
EJEMPLO
&\sphinxstyletheadfamily 
\sphinxAtStartPar
Clasificación
&\sphinxstyletheadfamily 
\sphinxAtStartPar
¿Por qué?
\\
\sphinxmidrule
\sphinxtableatstartofbodyhook
\sphinxAtStartPar
Instrucciones para llegar a un lugar
&
\sphinxAtStartPar

&
\sphinxAtStartPar

\\
\sphinxhline
\sphinxAtStartPar
Instrucciones con medición de tiempo para llegar a un lugar
&
\sphinxAtStartPar

&
\sphinxAtStartPar

\\
\sphinxhline
\sphinxAtStartPar
A=10,B=12C=A+B
&
\sphinxAtStartPar

&
\sphinxAtStartPar

\\
\sphinxhline
\sphinxAtStartPar
Microsoft Word
&
\sphinxAtStartPar

&
\sphinxAtStartPar

\\
\sphinxhline
\sphinxAtStartPar
LINUX
&
\sphinxAtStartPar

&
\sphinxAtStartPar

\\
\sphinxhline
\sphinxAtStartPar
\sphinxcode{\sphinxupquote{print("Digite su nombre")}}
&
\sphinxAtStartPar

&
\sphinxAtStartPar

\\
\sphinxhline
\sphinxAtStartPar
PASO 01: Abra la bandeja de entrada
&
\sphinxAtStartPar

&
\sphinxAtStartPar

\\
\sphinxhline
\sphinxAtStartPar
PASO 02: Inserte hojas tamaño carta
&
\sphinxAtStartPar

&
\sphinxAtStartPar

\\
\sphinxhline
\sphinxAtStartPar
PASO 03: Cierre la bandeja de entrada
&
\sphinxAtStartPar

&
\sphinxAtStartPar

\\
\sphinxhline
\sphinxAtStartPar
BIOS
&
\sphinxAtStartPar

&
\sphinxAtStartPar

\\
\sphinxhline
\sphinxAtStartPar
\sphinxstylestrong{CANVA}
&
\sphinxAtStartPar

&
\sphinxAtStartPar

\\
\sphinxhline
\sphinxAtStartPar
LibreOffice Writer
&
\sphinxAtStartPar

&
\sphinxAtStartPar

\\
\sphinxhline
\sphinxAtStartPar
Control de vuelo con API de GOOGLE
&
\sphinxAtStartPar

&
\sphinxAtStartPar

\\
\sphinxhline
\sphinxAtStartPar
Sensor Temperatura con ARDUINO
&
\sphinxAtStartPar

&
\sphinxAtStartPar

\\
\sphinxhline
\sphinxAtStartPar
Instrucciones corte cabello
&
\sphinxAtStartPar

&
\sphinxAtStartPar

\\
\sphinxbottomrule
\end{tabulary}
\sphinxtableafterendhook\par
\sphinxattableend\end{savenotes}

\item {} 
\sphinxAtStartPar
En su cuaderno describa las 3 fases para desarrollar un algoritmo.

\item {} 
\sphinxAtStartPar
Responda la siguiente pregunta en su cuaderno:
\begin{itemize}
\item {} 
\sphinxAtStartPar
¿Que diferencia hay entre desarrollar un algoritmo, diseñar un algoritmo y construir un algoritmo?

\end{itemize}

\item {} 
\sphinxAtStartPar
Consulte en INTERNET que hace un arquitecto de software y luego en su cuaderno responda:
\begin{itemize}
\item {} 
\sphinxAtStartPar
¿Cual es la diferencia entre un arquitecto de software y un programador?

\end{itemize}

\item {} 
\sphinxAtStartPar
En su cuaderno responda las siguientes preguntas:
\begin{itemize}
\item {} 
\sphinxAtStartPar
¿Qué es una \sphinxstylestrong{instrucción}?

\item {} 
\sphinxAtStartPar
¿Por que la \sphinxstylestrong{claridad}, la \sphinxstylestrong{eficiencia} y la \sphinxstylestrong{seguridad} son importantes al momento de dar instrucciones?

\end{itemize}

\item {} 
\sphinxAtStartPar
En su cuaderno escriba una serie de \sphinxstylestrong{instrucciones} (seguras, claras y eficientes) para que un estudiante vaya al colegio desde que se despierta hasta que llega al aula de clase.

\item {} 
\sphinxAtStartPar
En su cuaderno describa cada una de las figuras de un \sphinxstylestrong{diagrama de flujo} y dibuje sus símbolos.

\end{enumerate}


\bigskip\hrule\bigskip



\subsection{A03: Orientaciones al extranjero (+1,0)}
\label{\detokenize{1_guias/10programacion/G1023:a03-orientaciones-al-extranjero-1-0}}
\sphinxAtStartPar
\DUrole{sd-sphinx-override}{\DUrole{sd-badge}{\DUrole{sd-bg-warning}{\DUrole{sd-bg-text-warning}{Actividad práctica}}}}
\begin{enumerate}
\sphinxsetlistlabels{\arabic}{enumi}{enumii}{}{.}%
\item {} 
\sphinxAtStartPar
Lea el siguiente caso:
\begin{quote}

\sphinxAtStartPar
“Suponga que un extranjero visita \sphinxstylestrong{Puerto Boyacá} y quiere hacer turismo en el pueblo.
El extranjero se encuentra en las afueras del hospital \sphinxstylestrong{JOSE CAYETANO VASQUEZ} y quiere conocer cada una de las \sphinxstylestrong{iglesias católicas} en Puerto Boyacá.”
\begin{itemize}
\item {} 
\sphinxAtStartPar
Iglesia San José Obrero (Barrio la Paz).

\item {} 
\sphinxAtStartPar
Iglesia San Pedro Claver (Centro).

\item {} 
\sphinxAtStartPar
Iglesia Cristo Rey (Cra 4a, Calle 5).

\end{itemize}
\end{quote}


\begin{savenotes}\sphinxattablestart
\sphinxthistablewithglobalstyle
\centering
\begin{tabulary}{\linewidth}[t]{T}
\sphinxtoprule
\sphinxtableatstartofbodyhook

\begin{sphinxfigure-in-table}
\centering
\capstart
\noindent\sphinxincludegraphics[width=650\sphinxpxdimen]{{image114}.png}
\sphinxfigcaption{Mapa Puerto Boyacá}\label{\detokenize{1_guias/10programacion/G1023:id14}}\end{sphinxfigure-in-table}\relax
\\
\sphinxbottomrule
\end{tabulary}
\sphinxtableafterendhook\par
\sphinxattableend\end{savenotes}

\item {} 
\sphinxAtStartPar
En su \sphinxstylestrong{cuaderno}, cree un \sphinxstylestrong{diagrama de flujo} con 15 instrucciones que le permitan al extranjero llegar a cada una de las iglesias. Tenga en cuenta lo siguiente:
\begin{itemize}
\item {} 
\sphinxAtStartPar
5 instrucciones por iglesia (15 instrucciones).

\item {} 
\sphinxAtStartPar
2 direcciones por instrucción (30 direcciones).

\item {} 
\sphinxAtStartPar
Cree un inventario de las 30 direcciones por medio de variables así:

\begin{sphinxVerbatim}[commandchars=\\\{\}]
\PYG{n+nv}{d01}\PYG{+w}{ }\PYG{o}{=}\PYG{+w}{ }\PYG{l+s+s2}{\PYGZdq{}Hospital José Cayetano Vasquez, (Cra 5 con calle 29)\PYGZdq{}}
\PYG{n+nv}{d02}\PYG{+w}{ }\PYG{o}{=}\PYG{+w}{ }\PYG{l+s+s2}{\PYGZdq{}UPTC, (Cra 5 \PYGZsh{} 25 \PYGZhy{} 2)\PYGZdq{}}

\PYG{n+nv}{d10}\PYG{+w}{ }\PYG{o}{=}\PYG{+w}{ }\PYG{l+s+s2}{\PYGZdq{}Iglesia San José Obrero, (Calle 28 \PYGZsh{} 3A \PYGZhy{} 54)\PYGZdq{}}

\PYG{n+nv}{d20}\PYG{+w}{ }\PYG{o}{=}\PYG{+w}{ }\PYG{l+s+s2}{\PYGZdq{}Iglesia San Pedro Claver, (Calle 14 \PYGZsh{} 3 \PYGZhy{} 21)\PYGZdq{}}

\PYG{n+nv}{d30}\PYG{+w}{ }\PYG{o}{=}\PYG{+w}{ }\PYG{l+s+s2}{\PYGZdq{}Iglesia Cristo Rey, (Cra 4a con calle 5)\PYGZdq{}}
\end{sphinxVerbatim}

\item {} 
\sphinxAtStartPar
Cada instrucción(operación) debe describir como ir de una \sphinxstylestrong{dirección A} a una \sphinxstylestrong{dirección B}. Por ejemplo:

\end{itemize}


\begin{savenotes}\sphinxattablestart
\sphinxthistablewithglobalstyle
\centering
\begin{tabulary}{\linewidth}[t]{T}
\sphinxtoprule
\sphinxtableatstartofbodyhook

\begin{sphinxfigure-in-table}
\centering
\capstart
\noindent\sphinxincludegraphics[width=700\sphinxpxdimen]{{image124}.png}
\sphinxfigcaption{\sphinxstylestrong{Diagrama de flujo:} Orientaciones al extranjero}\label{\detokenize{1_guias/10programacion/G1023:id15}}\end{sphinxfigure-in-table}\relax
\\
\sphinxbottomrule
\end{tabulary}
\sphinxtableafterendhook\par
\sphinxattableend\end{savenotes}
\begin{itemize}
\item {} 
\sphinxAtStartPar
Usar una \sphinxstylestrong{página de cuaderno} por cada diagrama.

\item {} 
\sphinxAtStartPar
Para identificar las direcciones, se recomienda usar \sphinxhref{https://www.google.com/maps/@5.9736064,-74.5865216,7242m}{Google Maps}.

\item {} 
\sphinxAtStartPar
No olvide que el extranjero prefiere caminar (no moto, no cicla, no vehículo).

\item {} 
\sphinxAtStartPar
Antes de orientar al extranjero para ir a una nueva dirección, siempre verificar la ubicación actual para evitar que el extranjero se \sphinxstylestrong{teletransporte por error}.

\end{itemize}

\item {} 
\sphinxAtStartPar
En su cuaderno, adapte el algoritmo de la actividad anterior, de tal forma que contenga 3 condicionales.
\begin{quote}

\sphinxAtStartPar
Los condicionales permitiran al extranjero tomar caminos alternos en caso de la ocurrencia de algún evento inesperado (vías cerradas, semáforos en rojo, perros bravos, negocios cerrados, etc.).
\end{quote}

\sphinxAtStartPar
Por ejemplo:


\begin{savenotes}\sphinxattablestart
\sphinxthistablewithglobalstyle
\centering
\begin{tabulary}{\linewidth}[t]{T}
\sphinxtoprule
\sphinxtableatstartofbodyhook

\begin{sphinxfigure-in-table}
\centering
\capstart
\noindent\sphinxincludegraphics[width=700\sphinxpxdimen]{{image134}.png}
\sphinxfigcaption{\sphinxstylestrong{Diagrama de flujo:} Instrucciones seguras}\label{\detokenize{1_guias/10programacion/G1023:id16}}\end{sphinxfigure-in-table}\relax
\\
\sphinxbottomrule
\end{tabulary}
\sphinxtableafterendhook\par
\sphinxattableend\end{savenotes}
\begin{itemize}
\item {} 
\sphinxAtStartPar
Usar una \sphinxstylestrong{página de cuaderno} por cada diagrama.

\end{itemize}

\end{enumerate}


\bigskip\hrule\bigskip



\subsection{A04: Algoritmos (+3,0)}
\label{\detokenize{1_guias/10programacion/G1023:a04-algoritmos-3-0}}
\sphinxAtStartPar
\DUrole{sd-sphinx-override}{\DUrole{sd-badge}{\DUrole{sd-bg-danger}{\DUrole{sd-bg-text-danger}{Transferencia}}}}
\begin{enumerate}
\sphinxsetlistlabels{\arabic}{enumi}{enumii}{}{.}%
\item {} 
\sphinxAtStartPar
En su cuaderno desarrolle cada uno de los siguientes ejercicios, siguiendo la metodología para desarrollar un algoritmo.

\end{enumerate}
\subsubsection*{Adquirir un revista física de arte}

\sphinxAtStartPar
Desarrolle un algoritmo para que una persona adquiera una revista física de arte. Tenga en cuenta lo siguiente:
\begin{quote}

\sphinxAtStartPar
La persona inicialmente busca en INTERNET una revista física de arte.

\sphinxAtStartPar
Una vez encontrada la revista, revisar cuanto cuesta y ver si tiene el dinero suficiente para comprarla. Si no tiene el dinero suficiente volver a buscar la revista en INTERNET cuantas veces sea necesario. Si tiene el dinero suficiente alistar el dinero.

\sphinxAtStartPar
Luego la persona debe revisar en GOOGLE MAPS para localizar la dirección de la librería. Activar la función de GOOGLE MAPS y dirigirse caminando hacia la librería. Intentar entrar a la librería. Si la librería no está abierta: volver a buscar la revista en INTERNET e iniciar de nuevo. Si la librería está abierta, entrar a la librería.

\sphinxAtStartPar
Revisar si la revista está en STOCK. Si la revista no está en STOCK, volver a buscar la revista en INTERNET y empezar de nuevo. SI la revista si está en STOCK, comprar la revista física de arte.
\end{quote}
\subsubsection*{Entrar a una casa con llave}

\sphinxAtStartPar
Desarrolle un algoritmo para que una persona entre a una casa cuya cerradura está asegurada con llave. Tenga en cuenta lo siguiente:
\begin{quote}

\sphinxAtStartPar
La persona inicialmente buscará la llave en su bolsillo. La persona intentará abrir la cerradura con la llave de su bolsillo.

\sphinxAtStartPar
Si la cerradura no abre ahora buscará otra llave en un tapete exterior y volverá a intentar abrir la cerradura con dicha llave.

\sphinxAtStartPar
Si la cerradura no abre ahora buscará otra llave en la guantera de un carro y volverá a intentar abrir la cerradura con dicha llave.

\sphinxAtStartPar
Si la cerradura no abre, llamará a un cerrajero para que abra la puerta. Si el cerrajero no contesta, entonces seguir intentando hasta que conteste.

\sphinxAtStartPar
Si la cerradura si abre, intentar abrir la puerta. Si la puerta abre entrar a la casa y confirmar que se logró entrar. Si la cerradura no abre, confirmar que no se logró entrar.
\end{quote}
\subsubsection*{Saludar y despedirse de una persona}

\sphinxAtStartPar
Desarrolle un algoritmo para que el usuario salude y se despida de una persona. Tenga en cuenta lo siguiente:
\begin{quote}

\sphinxAtStartPar
La persona debe verificar inicialmente si ya saludo a la otra persona

\sphinxAtStartPar
Si la persona no ha saludado, saludar a la persona. Esperar a que la persona salude. Si la persona no saluda, seguir intentándolo hasta que salude.

\sphinxAtStartPar
Una vez que la persona salude, preguntar como esta. Si la persona responde NO esta bien, preguntar que le pasa, escuchar a la persona y decirle que el tiempo lo arregla todo. Volver a preguntar si la persona esta bien (cuantas veces sea necesario) hasta que la persona responda que esta bien.  Si la persona responder que SI esta bien entonces manifestar tranquilidad.

\sphinxAtStartPar
Volver a verificar si la persona ya saludo.

\sphinxAtStartPar
Si la persona ya saludo, entonces determinar un tiempo para poder irse.

\sphinxAtStartPar
El usuario debe establecer el tiempo para irse.

\sphinxAtStartPar
Verificar que el tiempo se cumplió. Mientras el tiempo no se cumpla, entonces hablar con la persona de cualquier tema mientras se cumple el tiempo. Ir disminuyendo el tiempo.

\sphinxAtStartPar
Una vez el tiempo esté cumplido, despedirse. Esperar a que la persona responda la despedida. Si la persona no contesta seguir despidiéndose hasta que la persona conteste.

\sphinxAtStartPar
Una vez se despida de la persona generar una mensaje de IRSE.
\end{quote}
\subsubsection*{Calmar a una persona}

\sphinxAtStartPar
Desarrolle un algoritmo para calmar una persona. Tenga en cuenta lo siguiente:
\begin{quote}

\sphinxAtStartPar
Identificar el estado de la persona.

\sphinxAtStartPar
Si la persona esta bien, entonces irse del lugar. Si la persona esta en mal estado realizar lo siguiente:

\sphinxAtStartPar
Identificar la edad de la persona.

\sphinxAtStartPar
Si la persona es un bebe entonces intentar calmarla con un tetero. Si la persona es un niño entonces calmarla dandole un dulce. Si la persona es un adolescente entonces intentar calmarla dandole una gaseosa. Si la persona es un adulto entonces intentar calmarla dandole una cerveza. Si la persona es un anciano entonces intentar calmarla dandole un café.

\sphinxAtStartPar
Luego volver a identificar el estado de la persona y si todavía no se ha calmado, seguir intentando hasta que se calme.
\end{quote}


\bigskip\hrule\bigskip


\sphinxstepscope


\section{Guía 1024: PYTHON (P3)}
\label{\detokenize{1_guias/10programacion/G1024:guia-1024-python-p3}}\label{\detokenize{1_guias/10programacion/G1024::doc}}

\begin{savenotes}\sphinxattablestart
\sphinxthistablewithglobalstyle
\centering
\begin{tabular}[t]{\X{35}{100}\X{65}{100}}
\sphinxtoprule
\sphinxtableatstartofbodyhook

\begin{savenotes}\sphinxattablestart
\sphinxthistablewithglobalstyle
\centering
\begin{tabulary}{\linewidth}[t]{T}
\sphinxtoprule
\sphinxtableatstartofbodyhook
\sphinxAtStartPar
\sphinxincludegraphics{{escudo}.png}
\\
\sphinxbottomrule
\end{tabulary}
\sphinxtableafterendhook\par
\sphinxattableend\end{savenotes}
&

\begin{savenotes}\sphinxattablestart
\sphinxthistablewithglobalstyle
\raggedright
\begin{tabular}[t]{*{1}{\X{1}{1}}}
\sphinxtoprule
\sphinxtableatstartofbodyhook
\sphinxAtStartPar
\sphinxstylestrong{INSTITUCIÓN EDUCATIVA}: John F. Kennedy
\begin{itemize}
\item {} 
\sphinxAtStartPar
\sphinxstylestrong{DOCENTE}: Julian Camilo Fonseca Romero

\item {} 
\sphinxAtStartPar
\sphinxstylestrong{ASIGNATURA}: Programación

\item {} 
\sphinxAtStartPar
\sphinxstylestrong{GRADO}: ONCE

\item {} 
\sphinxAtStartPar
\sphinxstylestrong{PERIODO}: 03

\item {} 
\sphinxAtStartPar
\sphinxstylestrong{TIEMPO}: 10 semanas (2 horas / semana)

\item {} 
\sphinxAtStartPar
\sphinxstylestrong{TIEMPO DE TRABAJO}: 20/20 horas

\end{itemize}
\\
\sphinxbottomrule
\end{tabular}
\sphinxtableafterendhook\par
\sphinxattableend\end{savenotes}
\\
\sphinxbottomrule
\end{tabular}
\sphinxtableafterendhook\par
\sphinxattableend\end{savenotes}

\begin{sphinxuseclass}{sd-tab-set}
\begin{sphinxuseclass}{sd-tab-item}\subsubsection*{Objetivo}

\begin{sphinxuseclass}{sd-tab-content}
\sphinxAtStartPar
\DUrole{sd-sphinx-override}{\DUrole{sd-badge}{\DUrole{sd-bg-light}{\DUrole{sd-bg-text-light}{Exploración}}}}
\DUrole{sd-sphinx-override}{\DUrole{sd-badge}{\DUrole{sd-bg-info}{\DUrole{sd-bg-text-info}{Estructuración}}}}
\DUrole{sd-sphinx-override}{\DUrole{sd-badge}{\DUrole{sd-bg-warning}{\DUrole{sd-bg-text-warning}{Actividad práctica}}}}
\DUrole{sd-sphinx-override}{\DUrole{sd-badge}{\DUrole{sd-bg-danger}{\DUrole{sd-bg-text-danger}{Transferencia}}}}
\begin{itemize}
\item {} 
\sphinxAtStartPar
\sphinxstylestrong{Temática}: Fundamentos del lenguaje de programación Python, historia, características y uso práctico en la consola interactiva.

\item {} 
\sphinxAtStartPar
\sphinxstylestrong{Objetivo de aprendizaje}: Comprender la importancia de Python y aprender a utilizar la consola interactiva para consultar documentación y explorar el lenguaje de forma autónoma.

\end{itemize}

\end{sphinxuseclass}
\end{sphinxuseclass}
\begin{sphinxuseclass}{sd-tab-item}\subsubsection*{Orientaciones curriculares}

\begin{sphinxuseclass}{sd-tab-content}\begin{itemize}
\item {} 
\sphinxAtStartPar
\sphinxstylestrong{Componente}: Uso y apropiación de la T\&I.

\item {} 
\sphinxAtStartPar
\sphinxstylestrong{Competencia}: Genero propuestas innovadoras para el uso y aprovechamiento de productos tecnológicos.

\item {} 
\sphinxAtStartPar
\sphinxstylestrong{Evidencia de aprendizaje}: Utilizo adecuadamente herramientas informáticas para la búsqueda, organización, procesamiento, sistematización, comunicación y difusión de ideas.

\end{itemize}

\end{sphinxuseclass}
\end{sphinxuseclass}
\begin{sphinxuseclass}{sd-tab-item}\subsubsection*{Criterios de evaluación}

\begin{sphinxuseclass}{sd-tab-content}\begin{itemize}
\item {} 
\sphinxAtStartPar
Actividades desarrolladas en su totalidad.

\item {} 
\sphinxAtStartPar
Participación en clase.

\item {} 
\sphinxAtStartPar
Explicación clara de conceptos.

\item {} 
\sphinxAtStartPar
Argumentación de ideas.

\item {} 
\sphinxAtStartPar
Cumplimiento de las reglas del aula.

\end{itemize}

\end{sphinxuseclass}
\end{sphinxuseclass}
\begin{sphinxuseclass}{sd-tab-item}\subsubsection*{Recursos (1)}

\begin{sphinxuseclass}{sd-tab-content}\begin{itemize}
\item {} 
\sphinxAtStartPar
\sphinxhref{https://www.online-python.com}{PYTHON ONLINE}. Ejercutar la consola interactiva PYTHON en el navegador.

\end{itemize}

\end{sphinxuseclass}
\end{sphinxuseclass}
\end{sphinxuseclass}

\subsection{Mensaje del autor}
\label{\detokenize{1_guias/10programacion/G1024:mensaje-del-autor}}\begin{quote}

\sphinxAtStartPar
“En esta guía didáctica, el estudiante se familiariza con la consola interactiva de Python, una herramienta fundamental para explorar y conocer las características básicas del lenguaje. Este primer acercamiento le permite al estudiante comprender cómo funciona Python y sentar las bases necesarias para un aprendizaje más profundo y estructurado en programación”.

\begin{flushright}
---Camilo Fonseca
\end{flushright}
\end{quote}


\bigskip\hrule\bigskip



\subsection{A00: Introducción (+0,1)}
\label{\detokenize{1_guias/10programacion/G1024:a00-introduccion-0-1}}
\sphinxAtStartPar
\DUrole{sd-sphinx-override}{\DUrole{sd-badge}{\DUrole{sd-bg-light}{\DUrole{sd-bg-text-light}{Exploración}}}}
\begin{enumerate}
\sphinxsetlistlabels{\arabic}{enumi}{enumii}{}{.}%
\item {} 
\sphinxAtStartPar
TRANSCRIBA en su cuaderno el objetivo, las orientaciones curriculares y los criterios de evaluación de la presente guía.

\item {} 
\sphinxAtStartPar
LEA el \sphinxstylestrong{mensaje del autor} y en su cuaderno responda las siguientes preguntas:
\begin{itemize}
\item {} 
\sphinxAtStartPar
¿Qué entiende por consola interactiva?

\item {} 
\sphinxAtStartPar
¿Qué es para usted la palabra interactivo?

\item {} 
\sphinxAtStartPar
Por medio de un lenguaje de programación ¿Puedo comunicarme con una máquina? ¿Por qué?

\item {} 
\sphinxAtStartPar
El proceso de comunicación, ¿Se limita a solo dar órdenes o instrucciones? ¿Por qué?

\end{itemize}

\end{enumerate}


\bigskip\hrule\bigskip



\subsection{A01: Lenguaje de programación PYTHON (+0,1)}
\label{\detokenize{1_guias/10programacion/G1024:a01-lenguaje-de-programacion-python-0-1}}
\sphinxAtStartPar
\DUrole{sd-sphinx-override}{\DUrole{sd-badge}{\DUrole{sd-bg-info}{\DUrole{sd-bg-text-info}{Estructuración}}}}
\begin{enumerate}
\sphinxsetlistlabels{\arabic}{enumi}{enumii}{}{.}%
\item {} 
\sphinxAtStartPar
Lea los siguientes conceptos:

\end{enumerate}


\bigskip\hrule\bigskip



\subsubsection{¿Qué es PYTHON?}
\label{\detokenize{1_guias/10programacion/G1024:que-es-python}}
\begin{sphinxuseclass}{sd-card}
\begin{sphinxuseclass}{sd-sphinx-override}
\begin{sphinxuseclass}{sd-w-25}
\begin{sphinxuseclass}{sd-mb-3}
\begin{sphinxuseclass}{sd-shadow-lg}
\begin{sphinxuseclass}{sd-card-hover}
\noindent\sphinxincludegraphics{{image0112}.png}

\begin{sphinxuseclass}{sd-card-body}
\begin{sphinxuseclass}{sd-card-title}
\begin{sphinxuseclass}{sd-font-weight-bold}PYTHON
\end{sphinxuseclass}
\end{sphinxuseclass}
\sphinxAtStartPar
Acceder a la web PYTHON

\end{sphinxuseclass}\sphinxhref{https://www.python.org/}{https://www.python.org/}
\end{sphinxuseclass}
\end{sphinxuseclass}
\end{sphinxuseclass}
\end{sphinxuseclass}
\end{sphinxuseclass}
\end{sphinxuseclass}
\sphinxAtStartPar
\sphinxstylestrong{Python} es un \sphinxstyleemphasis{lenguaje de programación} interpretado de alto nivel creado por \sphinxstylestrong{Guido van Rossum} y publicado por primera vez en \sphinxstylestrong{1991}. Está diseñado pensando en la legibilidad del código, y su sintaxis permite a los programadores expresar conceptos en menos líneas de \sphinxstylestrong{código}.

\begin{sphinxadmonition}{note}{¿Qué es un lenguaje de programación?}

\sphinxAtStartPar
Un \sphinxstylestrong{lenguaje de programación} es un conjunto de reglas y sintaxis que permite a los programadores dar \sphinxstyleemphasis{instrucciones} a una computadora para que realice tareas específicas. Es una forma de comunicarse con las máquinas.
\end{sphinxadmonition}


\bigskip\hrule\bigskip



\subsubsection{Principales características de Python}
\label{\detokenize{1_guias/10programacion/G1024:principales-caracteristicas-de-python}}
\sphinxAtStartPar
A continuación se muestran algunas de las características de Python que lo convierten en un lenguaje de programación tan versátil y ampliamente utilizado:
\begin{itemize}
\item {} 
\sphinxAtStartPar
\sphinxstylestrong{Legibilidad}. Python es conocido por su sintaxis clara y legible, que se parece en cierto modo a la inglesa.

\begin{sphinxVerbatim}[commandchars=\\\{\},numbers=left,firstnumber=1,stepnumber=1]
\PYG{k+kn}{import}\PYG{+w}{ }\PYG{n+nn}{time}

\PYG{n}{saludo} \PYG{o}{=} \PYG{l+s+s2}{\PYGZdq{}}\PYG{l+s+s2}{Buenos días}\PYG{l+s+s2}{\PYGZdq{}}

\PYG{k}{def}\PYG{+w}{ }\PYG{n+nf}{saludar}\PYG{p}{(}\PYG{p}{)}\PYG{p}{:}
    \PYG{n+nb}{print}\PYG{p}{(}\PYG{n}{saludo}\PYG{p}{)}

\PYG{n}{saludo}\PYG{p}{(}\PYG{p}{)}
\end{sphinxVerbatim}

\item {} 
\sphinxAtStartPar
\sphinxstylestrong{Fácil de aprender}. La legibilidad de Python hace que sea relativamente fácil para los principiantes entender lo que hace el código.

\item {} 
\sphinxAtStartPar
\sphinxstylestrong{Versatilidad}. Python no está limitado a un tipo de tarea: Usted puede utilizarlo en muchos campos. Si le interesa el desarrollo web, la automatización de tareas o la ciencia de datos, Python tiene las herramientas necesarias.

\item {} 
\sphinxAtStartPar
\sphinxstylestrong{Amplia compatibilidad con bibliotecas}. Incluye una gran biblioteca estándar con código pre\sphinxhyphen{}escrito para diversas tareas, lo que te ahorra tiempo y esfuerzo. Además, la vibrante comunidad de Python ha desarrollado miles de paquetes de terceros, que amplían aún más la funcionalidad de Python.

\item {} 
\sphinxAtStartPar
\sphinxstylestrong{Independencia de la plataforma}. Una de las grandes ventajas del lenguaje es que se puede escribir el código una vez y ejecutarlo en cualquier sistema operativo. Esta característica hace que Python sea una gran elección si trabaja en un equipo con diferentes sistemas operativos.

\item {} 
\sphinxAtStartPar
\sphinxstylestrong{Lenguaje interpretado}. Python es un lenguaje interpretado, lo que significa que el código se ejecuta línea a línea. Esto puede facilitar la depuración, ya que se puede probar pequeños fragmentos de código sin tener que compilar todo el programa.

\item {} 
\sphinxAtStartPar
\sphinxstylestrong{Código abierto y gratuito}. También es un lenguaje de código abierto, lo que significa que su código fuente está disponible libremente y puede distribuirse y modificarse. Esto ha llevado a una gran comunidad de desarrolladores a contribuir a su desarrollo y crear un vasto ecosistema de bibliotecas Python.

\begin{sphinxuseclass}{sd-card}
\begin{sphinxuseclass}{sd-sphinx-override}
\begin{sphinxuseclass}{sd-mb-3}
\begin{sphinxuseclass}{sd-shadow-lg}
\begin{sphinxuseclass}{sd-card-hover}
\begin{sphinxuseclass}{sd-card-body}
\begin{sphinxuseclass}{sd-card-title}
\begin{sphinxuseclass}{sd-font-weight-bold}PYTHON GIT
\end{sphinxuseclass}
\end{sphinxuseclass}
\sphinxAtStartPar
Acceder al código fuente PYTHON

\end{sphinxuseclass}\sphinxhref{https://github.com/python}{https://github.com/python}
\end{sphinxuseclass}
\end{sphinxuseclass}
\end{sphinxuseclass}
\end{sphinxuseclass}
\end{sphinxuseclass}
\item {} 
\sphinxAtStartPar
\sphinxstylestrong{Tipado dinámico:} Python está tipado dinámicamente, lo que significa que no es necesario declarar el tipo de datos de una variable. El intérprete de Python infiere el tipo, lo que hace que el código sea más flexible y sea más fácil trabajar con él.

\end{itemize}
\begin{quote}

\sphinxAtStartPar
PYTHON ES UN LENGUAJE DE PROGRAMACIÓN MULTI\sphinxhyphen{}PROPÓSITO.
\end{quote}


\bigskip\hrule\bigskip



\subsubsection{INDICE TIOBE}
\label{\detokenize{1_guias/10programacion/G1024:indice-tiobe}}
\sphinxAtStartPar
A lo largo de los años, Python se ha convertido en uno de los lenguajes de programación más populares debido a su sencillez, versatilidad y amplia gama de aplicaciones. Esto se comprueba en el \sphinxstylestrong{índice TIOBE}.

\begin{sphinxadmonition}{note}{¿Qué es el índice TIOBE?}

\begin{sphinxuseclass}{sd-card}
\begin{sphinxuseclass}{sd-sphinx-override}
\begin{sphinxuseclass}{sd-w-50}
\begin{sphinxuseclass}{sd-mb-3}
\begin{sphinxuseclass}{sd-shadow-lg}
\begin{sphinxuseclass}{sd-card-hover}
\noindent\sphinxincludegraphics{{logo-tiobe}.svg}

\begin{sphinxuseclass}{sd-card-img-overlay}
\begin{sphinxuseclass}{sd-card-body}
\end{sphinxuseclass}\sphinxhref{https://www.tiobe.com/}{https://www.tiobe.com/}
\end{sphinxuseclass}
\end{sphinxuseclass}
\end{sphinxuseclass}
\end{sphinxuseclass}
\end{sphinxuseclass}
\end{sphinxuseclass}
\end{sphinxuseclass}
\sphinxAtStartPar
El \sphinxstylestrong{índice de la comunidad de programación} (TIOBE) es un indicador de la popularidad de los lenguajes de programación creado por la compañía \sphinxstylestrong{TIOBE}. El índice se actualiza una vez al mes. Las calificaciones se basan en el número de ingenieros calificados en todo el mundo, cursos y proveedores de terceros.
\end{sphinxadmonition}

\begin{sphinxadmonition}{note}{¿Cuales son las fuentes del índice TIOBE?}

\sphinxAtStartPar
El índice TIOBE se basa en la consulta de sitios web populares como Google, Amazon, Wikipedia, Bing y más de 20 otros para calcular las calificaciones.

\sphinxAtStartPar
Es importante señalar que el índice TIOBE no se trata del mejor lenguaje de programación ni del lenguaje en el que se han escrito más líneas de código.
\end{sphinxadmonition}

\sphinxAtStartPar
En la siguiente gráfica se muestra la \sphinxstylestrong{evolución del uso de PYTHON desde el año 2002} hasta la actualidad:


\begin{savenotes}\sphinxattablestart
\sphinxthistablewithglobalstyle
\centering
\begin{tabulary}{\linewidth}[t]{T}
\sphinxtoprule
\sphinxtableatstartofbodyhook

\begin{sphinxfigure-in-table}
\centering
\capstart
\sphinxhref{https://www.tiobe.com/tiobe-index/python/}{\sphinxincludegraphics[width=700\sphinxpxdimen]{{image0210}.png}}
\sphinxfigcaption{Evolución uso de PYTHON 2002\sphinxhyphen{}2026}\label{\detokenize{1_guias/10programacion/G1024:id1}}\end{sphinxfigure-in-table}\relax
\\
\sphinxbottomrule
\end{tabulary}
\sphinxtableafterendhook\par
\sphinxattableend\end{savenotes}

\sphinxAtStartPar
La siguiente tabla muestra el ranking de lenguajes de programación más usados en \sphinxstylestrong{Diciembre de 2024 y Diciembre de 2025}:


\begin{savenotes}\sphinxattablestart
\sphinxthistablewithglobalstyle
\centering
\begin{tabulary}{\linewidth}[t]{T}
\sphinxtoprule
\sphinxtableatstartofbodyhook

\begin{sphinxfigure-in-table}
\centering
\capstart
\sphinxhref{https://www.tiobe.com/tiobe-index/}{\sphinxincludegraphics[width=500\sphinxpxdimen]{{image039}.png}}
\sphinxfigcaption{Ranking lenguajes de programación más usados}\label{\detokenize{1_guias/10programacion/G1024:id2}}\end{sphinxfigure-in-table}\relax
\\
\sphinxbottomrule
\end{tabulary}
\sphinxtableafterendhook\par
\sphinxattableend\end{savenotes}


\bigskip\hrule\bigskip



\subsubsection{Aplicaciones de PYTHON}
\label{\detokenize{1_guias/10programacion/G1024:aplicaciones-de-python}}
\begin{sphinxadmonition}{note}{Desarrollo de software}

\sphinxAtStartPar
Python se usa para la creación de scripts, automatización de tareas y pruebas que se ejecutan en \sphinxstylestrong{sistema operativo de computadores}. Por ejemplo:
\begin{itemize}
\item {} 
\sphinxAtStartPar
Scripts para leer y escribir archivos, mostrar mensajes en pantalla, mostrar interfaces gráficas con botones o formularios.

\item {} 
\sphinxAtStartPar
Automatización de tareas de escritorio para emular clics del mouse, entradas por teclado para tareas repetitivas, envío automático de correos electrónicos y recolección automática de datos de sitios web. (Herramienta \sphinxstylestrong{schedule})

\begin{sphinxuseclass}{sd-card}
\begin{sphinxuseclass}{sd-sphinx-override}
\begin{sphinxuseclass}{sd-w-25}
\begin{sphinxuseclass}{sd-mb-3}
\begin{sphinxuseclass}{sd-shadow-lg}
\begin{sphinxuseclass}{sd-card-hover}
\begin{sphinxuseclass}{sd-card-body}
\begin{sphinxuseclass}{sd-card-title}
\begin{sphinxuseclass}{sd-font-weight-bold}Schedule
\end{sphinxuseclass}
\end{sphinxuseclass}
\end{sphinxuseclass}
\noindent\sphinxincludegraphics{{image04}.svg}
\sphinxhref{https://pypi.org/project/schedule/}{https://pypi.org/project/schedule/}
\end{sphinxuseclass}
\end{sphinxuseclass}
\end{sphinxuseclass}
\end{sphinxuseclass}
\end{sphinxuseclass}
\end{sphinxuseclass}
\item {} 
\sphinxAtStartPar
Pruebas de depuración en funciones específicas de ejecución para encontrar comportamientos no deseados en la aplicación. (herramienta PyTest)

\begin{sphinxuseclass}{sd-card}
\begin{sphinxuseclass}{sd-sphinx-override}
\begin{sphinxuseclass}{sd-w-25}
\begin{sphinxuseclass}{sd-mb-3}
\begin{sphinxuseclass}{sd-shadow-lg}
\begin{sphinxuseclass}{sd-card-hover}
\begin{sphinxuseclass}{sd-card-body}
\begin{sphinxuseclass}{sd-card-title}
\begin{sphinxuseclass}{sd-font-weight-bold}Pytest
\end{sphinxuseclass}
\end{sphinxuseclass}
\end{sphinxuseclass}
\noindent\sphinxincludegraphics{{image057}.png}
\sphinxhref{https://docs.pytest.org/en/stable/}{https://docs.pytest.org/en/stable/}
\end{sphinxuseclass}
\end{sphinxuseclass}
\end{sphinxuseclass}
\end{sphinxuseclass}
\end{sphinxuseclass}
\end{sphinxuseclass}
\end{itemize}
\end{sphinxadmonition}

\begin{sphinxadmonition}{note}{Desarrollo web}

\sphinxAtStartPar
Con Python se puede desarrollar programas que se ejecutan en un \sphinxstylestrong{navegador web}. Por ejemplo:
\begin{itemize}
\item {} 
\sphinxAtStartPar
Con librerías como \sphinxstylestrong{Django} o \sphinxstylestrong{Flask}, se crean aplicaciones web que tengan sistemas de autenticación, paneles de administración e interacción con bases de datos.

\item {} 
\sphinxAtStartPar
\sphinxstylestrong{FASTAPI} permite la construcción sistemas que validen datos rápidamente entre el cliente y el servidor.

\end{itemize}

\begin{sphinxuseclass}{sd-container-fluid}
\begin{sphinxuseclass}{sd-sphinx-override}
\begin{sphinxuseclass}{sd-mb-4}
\begin{sphinxuseclass}{sd-row}
\begin{sphinxuseclass}{sd-row-cols-3}
\begin{sphinxuseclass}{sd-row-cols-xs-3}
\begin{sphinxuseclass}{sd-row-cols-sm-3}
\begin{sphinxuseclass}{sd-row-cols-md-3}
\begin{sphinxuseclass}{sd-row-cols-lg-3}
\begin{sphinxuseclass}{sd-col}
\begin{sphinxuseclass}{sd-d-flex-row}
\begin{sphinxuseclass}{sd-card}
\begin{sphinxuseclass}{sd-sphinx-override}
\begin{sphinxuseclass}{sd-w-100}
\begin{sphinxuseclass}{sd-shadow-sm}
\begin{sphinxuseclass}{sd-card-hover}
\begin{sphinxuseclass}{sd-card-body}
\begin{sphinxuseclass}{sd-card-title}
\begin{sphinxuseclass}{sd-font-weight-bold}Django
\end{sphinxuseclass}
\end{sphinxuseclass}
\end{sphinxuseclass}
\noindent\sphinxincludegraphics{{image067}.png}
\sphinxhref{https://www.djangoproject.com}{https://www.djangoproject.com}
\end{sphinxuseclass}
\end{sphinxuseclass}
\end{sphinxuseclass}
\end{sphinxuseclass}
\end{sphinxuseclass}
\end{sphinxuseclass}
\end{sphinxuseclass}
\begin{sphinxuseclass}{sd-col}
\begin{sphinxuseclass}{sd-d-flex-row}
\begin{sphinxuseclass}{sd-card}
\begin{sphinxuseclass}{sd-sphinx-override}
\begin{sphinxuseclass}{sd-w-100}
\begin{sphinxuseclass}{sd-shadow-sm}
\begin{sphinxuseclass}{sd-card-hover}
\begin{sphinxuseclass}{sd-card-body}
\begin{sphinxuseclass}{sd-card-title}
\begin{sphinxuseclass}{sd-font-weight-bold}Flask
\end{sphinxuseclass}
\end{sphinxuseclass}
\end{sphinxuseclass}
\noindent\sphinxincludegraphics{{image077}.png}
\sphinxhref{https://flask.palletsprojects.com/es/stable/}{https://flask.palletsprojects.com/es/stable/}
\end{sphinxuseclass}
\end{sphinxuseclass}
\end{sphinxuseclass}
\end{sphinxuseclass}
\end{sphinxuseclass}
\end{sphinxuseclass}
\end{sphinxuseclass}
\begin{sphinxuseclass}{sd-col}
\begin{sphinxuseclass}{sd-d-flex-row}
\begin{sphinxuseclass}{sd-card}
\begin{sphinxuseclass}{sd-sphinx-override}
\begin{sphinxuseclass}{sd-w-100}
\begin{sphinxuseclass}{sd-shadow-sm}
\begin{sphinxuseclass}{sd-card-hover}
\begin{sphinxuseclass}{sd-card-body}
\begin{sphinxuseclass}{sd-card-title}
\begin{sphinxuseclass}{sd-font-weight-bold}FastAPI
\end{sphinxuseclass}
\end{sphinxuseclass}
\end{sphinxuseclass}
\noindent\sphinxincludegraphics{{image087}.png}
\sphinxhref{https://fastapi.tiangolo.com}{https://fastapi.tiangolo.com}
\end{sphinxuseclass}
\end{sphinxuseclass}
\end{sphinxuseclass}
\end{sphinxuseclass}
\end{sphinxuseclass}
\end{sphinxuseclass}
\end{sphinxuseclass}
\end{sphinxuseclass}
\end{sphinxuseclass}
\end{sphinxuseclass}
\end{sphinxuseclass}
\end{sphinxuseclass}
\end{sphinxuseclass}
\end{sphinxuseclass}
\end{sphinxuseclass}
\end{sphinxuseclass}\end{sphinxadmonition}

\begin{sphinxadmonition}{note}{Desarrollo de juegos}

\sphinxAtStartPar
Python es usado para la creación de videojuegos que se pueden ejecutar en \sphinxstylestrong{sistema operativo de computadoras o consolas}. Por ejemplo:
\begin{itemize}
\item {} 
\sphinxAtStartPar
\sphinxstylestrong{Pygame} se usa para el desarrollo de juegos 2D. Ofrece módulos para manejar gráficos, sonidos y eventos lo que facilita la creación de juegos de manera sencilla.

\item {} 
\sphinxAtStartPar
Con \sphinxstylestrong{PANDA 3D} se desarrollan juegos en 3D. Funciona con Python y C++. Es ideal para crear entornos tridimensionales complejos e interactivos.

\end{itemize}

\begin{sphinxuseclass}{sd-container-fluid}
\begin{sphinxuseclass}{sd-sphinx-override}
\begin{sphinxuseclass}{sd-mb-4}
\begin{sphinxuseclass}{sd-row}
\begin{sphinxuseclass}{sd-row-cols-2}
\begin{sphinxuseclass}{sd-row-cols-xs-2}
\begin{sphinxuseclass}{sd-row-cols-sm-2}
\begin{sphinxuseclass}{sd-row-cols-md-2}
\begin{sphinxuseclass}{sd-row-cols-lg-2}
\begin{sphinxuseclass}{sd-col}
\begin{sphinxuseclass}{sd-d-flex-row}
\begin{sphinxuseclass}{sd-card}
\begin{sphinxuseclass}{sd-sphinx-override}
\begin{sphinxuseclass}{sd-w-100}
\begin{sphinxuseclass}{sd-shadow-sm}
\begin{sphinxuseclass}{sd-card-hover}
\begin{sphinxuseclass}{sd-card-body}
\begin{sphinxuseclass}{sd-card-title}
\begin{sphinxuseclass}{sd-font-weight-bold}PyGAME
\end{sphinxuseclass}
\end{sphinxuseclass}
\end{sphinxuseclass}
\noindent\sphinxincludegraphics{{image097}.png}
\sphinxhref{https://www.pygame.org/}{https://www.pygame.org/}
\end{sphinxuseclass}
\end{sphinxuseclass}
\end{sphinxuseclass}
\end{sphinxuseclass}
\end{sphinxuseclass}
\end{sphinxuseclass}
\end{sphinxuseclass}
\begin{sphinxuseclass}{sd-col}
\begin{sphinxuseclass}{sd-d-flex-row}
\begin{sphinxuseclass}{sd-card}
\begin{sphinxuseclass}{sd-sphinx-override}
\begin{sphinxuseclass}{sd-w-100}
\begin{sphinxuseclass}{sd-shadow-sm}
\begin{sphinxuseclass}{sd-card-hover}
\begin{sphinxuseclass}{sd-card-body}
\begin{sphinxuseclass}{sd-card-title}
\begin{sphinxuseclass}{sd-font-weight-bold}PANDA 3D
\end{sphinxuseclass}
\end{sphinxuseclass}
\end{sphinxuseclass}
\noindent\sphinxincludegraphics{{image106}.png}
\sphinxhref{https://www.panda3d.org/}{https://www.panda3d.org/}
\end{sphinxuseclass}
\end{sphinxuseclass}
\end{sphinxuseclass}
\end{sphinxuseclass}
\end{sphinxuseclass}
\end{sphinxuseclass}
\end{sphinxuseclass}
\end{sphinxuseclass}
\end{sphinxuseclass}
\end{sphinxuseclass}
\end{sphinxuseclass}
\end{sphinxuseclass}
\end{sphinxuseclass}
\end{sphinxuseclass}
\end{sphinxuseclass}
\end{sphinxuseclass}\end{sphinxadmonition}

\begin{sphinxadmonition}{note}{Inteligencia artificial}

\sphinxAtStartPar
Debido a su simplicidad y gran cantidad de herramientas disponibles, Python se usa ampliamente en el desarrollo de aplicaciones de Inteligencia Artificial y Machine Learning. Por ejemplo:
\begin{itemize}
\item {} 
\sphinxAtStartPar
Con la herramienta \sphinxstylestrong{TensorFlow} se pueden crear modelos de aprendizajes profundos. Los programadores pueden construir y entrenar redes neuronales complejas para aplicaciones como reconocimiento de voz o visión por computadora.

\begin{sphinxuseclass}{sd-card}
\begin{sphinxuseclass}{sd-sphinx-override}
\begin{sphinxuseclass}{sd-w-50}
\begin{sphinxuseclass}{sd-mb-3}
\begin{sphinxuseclass}{sd-shadow-lg}
\begin{sphinxuseclass}{sd-card-hover}
\begin{sphinxuseclass}{sd-card-body}
\begin{sphinxuseclass}{sd-card-title}
\begin{sphinxuseclass}{sd-font-weight-bold}TensorFlow
\end{sphinxuseclass}
\end{sphinxuseclass}
\end{sphinxuseclass}
\noindent\sphinxincludegraphics{{image11}.svg}
\sphinxhref{https://www.tensorflow.org/}{https://www.tensorflow.org/}
\end{sphinxuseclass}
\end{sphinxuseclass}
\end{sphinxuseclass}
\end{sphinxuseclass}
\end{sphinxuseclass}
\end{sphinxuseclass}
\end{itemize}
\end{sphinxadmonition}

\begin{sphinxadmonition}{note}{Ciencia de datos}

\sphinxAtStartPar
Python se usa ampliamente para la manipulación y análisis de grandes cantidades de datos. Por ejemplo:
\begin{itemize}
\item {} 
\sphinxAtStartPar
Con \sphinxstylestrong{PANDAS} se trabajan estructuras, limpieza, transformación y análisis de datos. Es ideal para manejar grandes conjuntos de datos de forma eficiente.

\begin{sphinxuseclass}{sd-card}
\begin{sphinxuseclass}{sd-sphinx-override}
\begin{sphinxuseclass}{sd-w-50}
\begin{sphinxuseclass}{sd-mb-3}
\begin{sphinxuseclass}{sd-shadow-sm}
\begin{sphinxuseclass}{sd-card-hover}
\begin{sphinxuseclass}{sd-card-body}
\begin{sphinxuseclass}{sd-card-title}
\begin{sphinxuseclass}{sd-font-weight-bold}PANDAS
\end{sphinxuseclass}
\end{sphinxuseclass}
\end{sphinxuseclass}
\noindent\sphinxincludegraphics{{image125}.png}
\sphinxhref{https://pandas.pydata.org}{https://pandas.pydata.org}
\end{sphinxuseclass}
\end{sphinxuseclass}
\end{sphinxuseclass}
\end{sphinxuseclass}
\end{sphinxuseclass}
\end{sphinxuseclass}
\item {} 
\sphinxAtStartPar
\sphinxstylestrong{NumPy} se usa para ejecutar cálculos numéricos. Soporta arreglos multidimensionales y funciones matemáticas de alto rendimiento. También se pueden trabajar sobre aplicaciones estadísticas que requieran analizar altos volúmenes de información.

\begin{sphinxuseclass}{sd-card}
\begin{sphinxuseclass}{sd-sphinx-override}
\begin{sphinxuseclass}{sd-w-25}
\begin{sphinxuseclass}{sd-mb-3}
\begin{sphinxuseclass}{sd-shadow-sm}
\begin{sphinxuseclass}{sd-card-hover}
\begin{sphinxuseclass}{sd-card-body}
\begin{sphinxuseclass}{sd-card-title}
\begin{sphinxuseclass}{sd-font-weight-bold}NumPy
\end{sphinxuseclass}
\end{sphinxuseclass}
\end{sphinxuseclass}
\noindent\sphinxincludegraphics{{image13}.svg}
\sphinxhref{https://numpy.org/}{https://numpy.org/}
\end{sphinxuseclass}
\end{sphinxuseclass}
\end{sphinxuseclass}
\end{sphinxuseclass}
\end{sphinxuseclass}
\end{sphinxuseclass}
\item {} 
\sphinxAtStartPar
\sphinxstylestrong{Matplotlib} es usado para visualización de datos con gráficos estadísticos estáticos, animados o interactivos.

\begin{sphinxuseclass}{sd-card}
\begin{sphinxuseclass}{sd-sphinx-override}
\begin{sphinxuseclass}{sd-w-50}
\begin{sphinxuseclass}{sd-mb-3}
\begin{sphinxuseclass}{sd-shadow-sm}
\begin{sphinxuseclass}{sd-card-hover}
\begin{sphinxuseclass}{sd-card-body}
\begin{sphinxuseclass}{sd-card-title}
\begin{sphinxuseclass}{sd-font-weight-bold}Matplotlib
\end{sphinxuseclass}
\end{sphinxuseclass}
\end{sphinxuseclass}
\noindent\sphinxincludegraphics{{image14}.svg}
\sphinxhref{https://matplotlib.org/}{https://matplotlib.org/}
\end{sphinxuseclass}
\end{sphinxuseclass}
\end{sphinxuseclass}
\end{sphinxuseclass}
\end{sphinxuseclass}
\end{sphinxuseclass}
\end{itemize}
\end{sphinxadmonition}


\bigskip\hrule\bigskip



\subsubsection{Consola Interactiva PYTHON}
\label{\detokenize{1_guias/10programacion/G1024:consola-interactiva-python}}
\sphinxAtStartPar
La consola de PYTHON es una herramienta que permite ejecutar algoritmos de forma interactiva.

\sphinxAtStartPar
Esta consola puede ejecutarse en cualquier sistema operativo que tenga instalado PYTHON.
\begin{itemize}
\item {} 
\sphinxAtStartPar
En WINDOWS se puede buscar el aplicativo: \sphinxstylestrong{CMD}

\item {} 
\sphinxAtStartPar
En LINUX se busca el aplicativo: \sphinxstylestrong{TERMINAL}

\end{itemize}

\sphinxAtStartPar
Una vez dentro del aplicativo (\sphinxstylestrong{CMD} o \sphinxstylestrong{TERMINAL}) se escribe el comando: \sphinxstylestrong{\sphinxguilabel{python3}}

\begin{sphinxVerbatim}[commandchars=\\\{\}]
C:\PYG{l+s+se}{\PYGZbs{}\PYGZgt{}}\PYG{+w}{ }python3
\end{sphinxVerbatim}

\sphinxAtStartPar
Una vez ejecutado este comando, el usuario estará dentro de la consola interactiva de PYTHON así como muestra la siguiente imagen:


\begin{savenotes}\sphinxattablestart
\sphinxthistablewithglobalstyle
\centering
\begin{tabulary}{\linewidth}[t]{T}
\sphinxtoprule
\sphinxtableatstartofbodyhook

\begin{sphinxfigure-in-table}
\centering
\capstart
\noindent\sphinxincludegraphics[width=700\sphinxpxdimen]{{image153}.png}
\sphinxfigcaption{Consola interactiva PYTHON}\label{\detokenize{1_guias/10programacion/G1024:id3}}\end{sphinxfigure-in-table}\relax
\\
\sphinxbottomrule
\end{tabulary}
\sphinxtableafterendhook\par
\sphinxattableend\end{savenotes}

\sphinxAtStartPar
La consola interactiva de Python inicialmente muestra la siguiente información:
\begin{itemize}
\item {} 
\sphinxAtStartPar
Versión de Python: \sphinxcode{\sphinxupquote{v3.13.1:06714517797}}

\item {} 
\sphinxAtStartPar
Fecha de compilación: \sphinxcode{\sphinxupquote{Dec 3 2024, 14:00:22}}

\item {} 
\sphinxAtStartPar
Herramienta de compilación de la consola PYTHON: \sphinxcode{\sphinxupquote{Clang 15.0.0}}.

\end{itemize}

\sphinxAtStartPar
Los comandos a introducir en esta consola se deben escribir después del prompt el cual son las tres signos mayor que (\sphinxcode{\sphinxupquote{>>>}}), los cuales se muestran al inicio.


\bigskip\hrule\bigskip



\subsection{A02: Taller (+0,8)}
\label{\detokenize{1_guias/10programacion/G1024:a02-taller-0-8}}
\sphinxAtStartPar
\DUrole{sd-sphinx-override}{\DUrole{sd-badge}{\DUrole{sd-bg-info}{\DUrole{sd-bg-text-info}{Estructuración}}}}
\begin{enumerate}
\sphinxsetlistlabels{\arabic}{enumi}{enumii}{}{.}%
\item {} 
\sphinxAtStartPar
En su cuaderno responda las siguientes preguntas:
\begin{itemize}
\item {} 
\sphinxAtStartPar
¿En que año fue creado PYTHON, quién lo desarrolló y por qué?

\end{itemize}

\item {} 
\sphinxAtStartPar
Cree un \sphinxstylestrong{diagrama de flujo} que explique las \sphinxstylestrong{principales características} de PYTHON. \sphinxstylestrong{No es necesario escribir la descripción de la característica}. Tenga en cuenta lo siguiente:
\begin{quote}

\sphinxAtStartPar
\sphinxstylestrong{INICIO} → PYTHON es legible.

\sphinxAtStartPar
Verificar si es fácil de aprender. Si “NO” es fácil de aprender entonces: “NO es versátil”. Si “SI” es fácil de aprender entonces tiene amplia compatibilidad con bibliotecas.

\sphinxAtStartPar
Al NO ser versátil entonces es dependiente de la plataforma.

\sphinxAtStartPar
Al ser dependiente de la plataforma, entonces será necesario volver a identificar si es fácil de aprender.

\sphinxAtStartPar
Al tener una amplia compatibilidad con bibliotecas, entonces es necesario verificar si es un lenguaje interpretado. Si “SI” es un lenguaje interpretado entonces es un lenguaje de tipado dinámico. Si “NO” es un lenguaje interpretado entonces será necesario verificar nuevamente Si es un lenguaje fácil de aprender.

\sphinxAtStartPar
Al ser un lenguaje de tipado dinámico entonces es necesario verificar que es de código abierto y gratuito. Si “SI” es de código abierto y gratuito, entonces mostrar un mensaje que diga: “EMPEZAR A PROGRAMAR” → \sphinxstylestrong{FIN}. Si “NO” es de código abierto, entonces  mostrar un mensaje que diga: “ESTE LENGUAJE NO ES PYTHON” → \sphinxstylestrong{FIN}.
\end{quote}

\item {} 
\sphinxAtStartPar
Responda si es \sphinxstylestrong{falso}, \sphinxstylestrong{verdadero} y \sphinxstylestrong{por qué}:

\sphinxAtStartPar
\sphinxstylestrong{a.} PYTHON solo se puede ejecutar en Windows. \_\_\_

\sphinxAtStartPar
\sphinxstylestrong{b.} La suscripción anual por el uso de PYTHON en proyectos privados cuesta \$1,000,000 COP. \_\_\_

\sphinxAtStartPar
\sphinxstylestrong{c.} El aplicativo web a desarrollar necesita cambiar cierta sintaxis del código fuente del lenguaje PYTHON. Lastima que no podamos cambiar nada o nos metemos en problemas legales. \_\_\_

\sphinxAtStartPar
\sphinxstylestrong{d.} Con PYTHON podemos trabajar en automatización, creación de videojuegos e inteligencia artificial. \_\_\_

\sphinxAtStartPar
\sphinxstylestrong{e.} Aprender PYTHON fue rápido. No sin antes aprender de algoritmia. Eso facilitó las cosas. \_\_\_

\item {} 
\sphinxAtStartPar
En su cuaderno, responda las siguientes preguntas:
\begin{itemize}
\item {} 
\sphinxAtStartPar
¿Qué es el índice TIOBE? y ¿Qué es TIOBE?

\item {} 
\sphinxAtStartPar
¿Cada cuanto se actualiza el índice TIOBE?

\item {} 
\sphinxAtStartPar
¿Cuáles son las fuentes del índice TIOBE?

\end{itemize}

\item {} 
\sphinxAtStartPar
Con la gráfica: \sphinxstylestrong{Evolución del uso de PYTHON desde el año 2002}, desarrolle la siguiente tabla en su cuaderno (identifique a diciembre del año):


\begin{savenotes}\sphinxattablestart
\sphinxthistablewithglobalstyle
\centering
\begin{tabulary}{\linewidth}[t]{TT}
\sphinxtoprule
\sphinxstyletheadfamily 
\sphinxAtStartPar
Año
&\sphinxstyletheadfamily 
\sphinxAtStartPar
Rating
\\
\sphinxmidrule
\sphinxtableatstartofbodyhook
\sphinxAtStartPar
2026
&
\sphinxAtStartPar

\\
\sphinxhline
\sphinxAtStartPar
2024
&
\sphinxAtStartPar
23,84\%
\\
\sphinxhline
\sphinxAtStartPar
2022
&
\sphinxAtStartPar

\\
\sphinxhline
\sphinxAtStartPar
2020
&
\sphinxAtStartPar

\\
\sphinxhline
\sphinxAtStartPar
2018
&
\sphinxAtStartPar

\\
\sphinxhline
\sphinxAtStartPar
2016
&
\sphinxAtStartPar

\\
\sphinxhline
\sphinxAtStartPar
2014
&
\sphinxAtStartPar

\\
\sphinxhline
\sphinxAtStartPar
2012
&
\sphinxAtStartPar

\\
\sphinxhline
\sphinxAtStartPar
2010
&
\sphinxAtStartPar

\\
\sphinxhline
\sphinxAtStartPar
2008
&
\sphinxAtStartPar

\\
\sphinxhline
\sphinxAtStartPar
2006
&
\sphinxAtStartPar

\\
\sphinxhline
\sphinxAtStartPar
2004
&
\sphinxAtStartPar

\\
\sphinxhline
\sphinxAtStartPar
2002
&
\sphinxAtStartPar

\\
\sphinxbottomrule
\end{tabulary}
\sphinxtableafterendhook\par
\sphinxattableend\end{savenotes}

\item {} 
\sphinxAtStartPar
Con la tabla: \sphinxstylestrong{Ranking de los lenguajes de programación más usados en Diciembre de 2024 y 2025}, cree un gráfica así:


\begin{savenotes}\sphinxattablestart
\sphinxthistablewithglobalstyle
\centering
\begin{tabulary}{\linewidth}[t]{T}
\sphinxtoprule
\sphinxtableatstartofbodyhook

\begin{sphinxfigure-in-table}
\centering
\capstart
\noindent\sphinxincludegraphics[width=240\sphinxpxdimen]{{image163}.png}
\sphinxfigcaption{EJEMPLO}\label{\detokenize{1_guias/10programacion/G1024:id4}}\end{sphinxfigure-in-table}\relax
\\
\sphinxbottomrule
\end{tabulary}
\sphinxtableafterendhook\par
\sphinxattableend\end{savenotes}
\begin{itemize}
\item {} 
\sphinxAtStartPar
Los lenguajes de programación con sus fechas van en el eje X (dos fechas por lenguaje Diciembre de 2024 y Diciembre 2025).

\item {} 
\sphinxAtStartPar
El Rating (\%) va en el eje Y.

\item {} 
\sphinxAtStartPar
Cruce con diagrama de barras.

\end{itemize}

\item {} 
\sphinxAtStartPar
Cree un \sphinxstylestrong{mapa conceptual} de herramientas usadas PYTHON. El mapa conceptual debe tener 3 niveles:
\begin{itemize}
\item {} 
\sphinxAtStartPar
1 NIVEL: Título Aplicaciones de PYTHON

\item {} 
\sphinxAtStartPar
2 NIVEL: 5 aplicaciones concretas

\item {} 
\sphinxAtStartPar
3 NIVEL: Herramientas por cada aplicación concreta.

\end{itemize}

\item {} 
\sphinxAtStartPar
En su cuaderno, COMPLETE la siguiente sopa de letras con 16 términos relacionados con el lenguaje de programación \sphinxstylestrong{PYTHON}. En la medida que encuentre un término, realice su descripción.


\begin{savenotes}\sphinxattablestart
\sphinxthistablewithglobalstyle
\centering
\begin{tabulary}{\linewidth}[t]{T}
\sphinxtoprule
\sphinxtableatstartofbodyhook

\begin{sphinxfigure-in-table}
\centering
\capstart
\noindent\sphinxincludegraphics[width=550\sphinxpxdimen]{{image173}.png}
\sphinxfigcaption{Términos PYTHON}\label{\detokenize{1_guias/10programacion/G1024:id5}}\end{sphinxfigure-in-table}\relax
\\
\sphinxbottomrule
\end{tabulary}
\sphinxtableafterendhook\par
\sphinxattableend\end{savenotes}

\end{enumerate}


\bigskip\hrule\bigskip



\subsection{A03: Consola interactiva Python (+0,5)}
\label{\detokenize{1_guias/10programacion/G1024:a03-consola-interactiva-python-0-5}}
\sphinxAtStartPar
\DUrole{sd-sphinx-override}{\DUrole{sd-badge}{\DUrole{sd-bg-warning}{\DUrole{sd-bg-text-warning}{Actividad práctica}}}}
\begin{quote}

\sphinxAtStartPar
En esta actividad, el estudiante aprende a usar la \sphinxstylestrong{consola interactiva} de PYTHON.
\end{quote}
\begin{enumerate}
\sphinxsetlistlabels{\arabic}{enumi}{enumii}{}{.}%
\item {} 
\sphinxAtStartPar
En su sistema operativo abra una terminal de sistema.
\begin{itemize}
\item {} 
\sphinxAtStartPar
Si es Windows use \sphinxstylestrong{CMD} (Command Prompt)

\item {} 
\sphinxAtStartPar
SI es LINUX use la \sphinxstylestrong{CLI} (Command Line Interface) de su distribución.

\end{itemize}

\item {} 
\sphinxAtStartPar
MAXIMIZAR la ventana de la terminal del sistema.

\item {} 
\sphinxAtStartPar
En la terminal del sistema, EJECUTE el comando \sphinxcode{\sphinxupquote{python3}}.

\begin{sphinxadmonition}{note}{¿Cómo ejecutar comandos?}

\sphinxAtStartPar
Dentro de la terminal, usted puede ejecutar comandos.

\sphinxAtStartPar
Para EJECUTAR un comando; es necesario escribir dicho comando y luego presionar la tecla \sphinxstylestrong{ENTER}.

\sphinxAtStartPar
Por ejemplo, AL EJECUTAR el comando \sphinxcode{\sphinxupquote{python3}} saldrá lo siguiente:

\begin{sphinxVerbatim}[commandchars=\\\{\}]
usuario@AC115\PYGZhy{}01:\PYGZti{}\PYGZdl{} python3
Python 3.13.0 (main, Oct 8 2024, 00:00:00) 
\end{sphinxVerbatim}
\end{sphinxadmonition}

\item {} 
\sphinxAtStartPar
AHORA usted puede ejecutar comandos en la consola interactiva de PYTHON.

\item {} 
\sphinxAtStartPar
IDENTIFIQUE el \sphinxstylestrong{prompt}.

\begin{sphinxadmonition}{note}{¿Qué es el prompt? (\sphinxstylestrong{>>>})}
\begin{itemize}
\item {} 
\sphinxAtStartPar
El prompt son tres caracteres de mayor que (>>>).

\item {} 
\sphinxAtStartPar
Los comandos a ejecutar (en la consola de PYTHON) se escriben después del prompt.

\item {} 
\sphinxAtStartPar
Se escribe el comando y luego se PULSA la tecla \sphinxstylestrong{ENTER}.

\end{itemize}
\end{sphinxadmonition}

\item {} 
\sphinxAtStartPar
Con la información mostrada en la consola, en su cuaderno complete la siguiente tabla:


\begin{savenotes}\sphinxattablestart
\sphinxthistablewithglobalstyle
\centering
\begin{tabulary}{\linewidth}[t]{TT}
\sphinxtoprule
\sphinxstyletheadfamily 
\sphinxAtStartPar
Consola interactiva de PYTHON
&\sphinxstyletheadfamily 
\sphinxAtStartPar
Descripción
\\
\sphinxmidrule
\sphinxtableatstartofbodyhook
\sphinxAtStartPar
Versión de PYTHON
&
\sphinxAtStartPar

\\
\sphinxhline
\sphinxAtStartPar
Fecha de compilación
&
\sphinxAtStartPar

\\
\sphinxhline
\sphinxAtStartPar
Herramienta de compilación
&
\sphinxAtStartPar

\\
\sphinxbottomrule
\end{tabulary}
\sphinxtableafterendhook\par
\sphinxattableend\end{savenotes}

\item {} 
\sphinxAtStartPar
En la consola interactiva de Python, a continuación escriba los siguientes comandos para obtener más información:
\begin{itemize}
\item {} 
\sphinxAtStartPar
EJECUTE el comando: \sphinxcode{\sphinxupquote{copyright}}. Con la información generada, en su cuaderno responda la siguiente pregunta:
\begin{itemize}
\item {} 
\sphinxAtStartPar
¿Cuáles son las cuatro organizaciones relacionadas con derechos reservados en el uso de PYTHON?

\end{itemize}

\item {} 
\sphinxAtStartPar
EJECUTE el comando: \sphinxcode{\sphinxupquote{credits}}. Traduzca al español la información generada y escríbala en su cuaderno.

\item {} 
\sphinxAtStartPar
EJECUTE el comando: \sphinxcode{\sphinxupquote{license()}}. Se genera una texto largo en la consola.

\begin{sphinxadmonition}{note}{Navegar en textos largos en consola}
\begin{itemize}
\item {} 
\sphinxAtStartPar
Cuando se generan textos largos en la consola, para avanzar PULSE la tecla \sphinxkeyboard{\sphinxupquote{ENTER}}.

\item {} 
\sphinxAtStartPar
Si se quiere salir del texto largo, PULSE la tecla \sphinxkeyboard{\sphinxupquote{Q}} y luego PULSE la tecla \sphinxkeyboard{\sphinxupquote{ENTER}}.

\end{itemize}
\end{sphinxadmonition}

\end{itemize}

\item {} 
\sphinxAtStartPar
NAVEGUE por todo el texto generado y en su cuaderno escriba solo los títulos (son 7 títulos). Por ejemplo:
\begin{itemize}
\item {} 
\sphinxAtStartPar
\sphinxstylestrong{A. HISTORY OF THE SOFTWARE}

\end{itemize}

\item {} 
\sphinxAtStartPar
Traduzca la información de solo el primer título \sphinxstylestrong{A. HISTORY OF THE SOFTWARE} (no escribir la traducción en su cuaderno!!!) y responda las siguientes preguntas en su cuaderno:
\begin{itemize}
\item {} 
\sphinxAtStartPar
¿En que año fue creado PYTHON?

\item {} 
\sphinxAtStartPar
¿Cuál es el nombre del creador de PYTHON?

\item {} 
\sphinxAtStartPar
¿En que lugar fue creado PYTHON?

\item {} 
\sphinxAtStartPar
¿Qué lenguaje de programación fue antecesor de PYTHON?

\item {} 
\sphinxAtStartPar
¿Qué eventos acontecieron en las siguientes fechas?:
\begin{itemize}
\item {} 
\sphinxAtStartPar
1995

\item {} 
\sphinxAtStartPar
Mayo del 2000

\item {} 
\sphinxAtStartPar
Octubre del 2000

\item {} 
\sphinxAtStartPar
2001

\end{itemize}

\item {} 
\sphinxAtStartPar
¿Qué significan las siguientes siglas?:
\begin{itemize}
\item {} 
\sphinxAtStartPar
CWI

\item {} 
\sphinxAtStartPar
CNRI

\item {} 
\sphinxAtStartPar
PSF

\end{itemize}

\end{itemize}

\item {} 
\sphinxAtStartPar
Complete la siguiente tabla relacionada con las versiones de PYTHON:


\begin{savenotes}\sphinxattablestart
\sphinxthistablewithglobalstyle
\centering
\begin{tabulary}{\linewidth}[t]{TTTTT}
\sphinxtoprule
\sphinxstyletheadfamily 
\sphinxAtStartPar
AÑO
&\sphinxstyletheadfamily 
\sphinxAtStartPar
Versión
&\sphinxstyletheadfamily 
\sphinxAtStartPar
Versión derivada
&\sphinxstyletheadfamily 
\sphinxAtStartPar
Propietario
&\sphinxstyletheadfamily 
\sphinxAtStartPar
Compatibilidad con GPL
\\
\sphinxmidrule
\sphinxtableatstartofbodyhook
\sphinxAtStartPar
1991\sphinxhyphen{}1995
&
\sphinxAtStartPar
0.9.0
&
\sphinxAtStartPar
N/A
&
\sphinxAtStartPar
CWI
&
\sphinxAtStartPar
SI
\\
\sphinxhline
\sphinxAtStartPar
1995\sphinxhyphen{}1999
&
\sphinxAtStartPar
1.3
&
\sphinxAtStartPar
1.2
&
\sphinxAtStartPar
CNRI
&
\sphinxAtStartPar
SI
\\
\sphinxhline
\sphinxAtStartPar
2000
&
\sphinxAtStartPar

&
\sphinxAtStartPar

&
\sphinxAtStartPar

&
\sphinxAtStartPar

\\
\sphinxbottomrule
\end{tabulary}
\sphinxtableafterendhook\par
\sphinxattableend\end{savenotes}

\item {} 
\sphinxAtStartPar
CIERRE la consola interactiva de PYTHON y la terminal de sistema.

\end{enumerate}


\bigskip\hrule\bigskip



\subsection{A04: Ayuda en la consola interactiva (+0,5)}
\label{\detokenize{1_guias/10programacion/G1024:a04-ayuda-en-la-consola-interactiva-0-5}}
\sphinxAtStartPar
\DUrole{sd-sphinx-override}{\DUrole{sd-badge}{\DUrole{sd-bg-warning}{\DUrole{sd-bg-text-warning}{Actividad práctica}}}}
\begin{enumerate}
\sphinxsetlistlabels{\arabic}{enumi}{enumii}{}{.}%
\item {} 
\sphinxAtStartPar
ABRA una \sphinxstylestrong{terminal de sistema} e ingrese a la \sphinxstylestrong{consola interactiva de PYTHON}.

\item {} 
\sphinxAtStartPar
EJECUTE el comando \sphinxcode{\sphinxupquote{help()}}. Aparecerá en pantalla lo siguiente:

\begin{sphinxVerbatim}[commandchars=\\\{\}]
\PYGZgt{}\PYGZgt{}\PYGZgt{} Welcome to Python 3.13\PYGZsq{}s help utility! If this is your first time using
Python, you should definitely check out the tutorial at
https://docs.python.org/3.13/tutorial/.

Enter the name of any module, keyword, or topic to get help on writing
Python programs and using Python modules.  To get a list of available
modules, keywords, symbols, or topics, enter \PYGZdq{}modules\PYGZdq{}, \PYGZdq{}keywords\PYGZdq{},
\PYGZdq{}symbols\PYGZdq{}, or \PYGZdq{}topics\PYGZdq{}.

Each module also comes with a one\PYGZhy{}line summary of what it does; to list
the modules whose name or summary contain a given string such as \PYGZdq{}spam\PYGZdq{},
enter \PYGZdq{}modules spam\PYGZdq{}.

To quit this help utility and return to the interpreter,
enter \PYGZdq{}q\PYGZdq{}, \PYGZdq{}quit\PYGZdq{} or \PYGZdq{}exit\PYGZdq{}.

help\PYGZgt{}
\end{sphinxVerbatim}

\begin{sphinxadmonition}{note}{Tenga en cuenta}

\sphinxAtStartPar
Tenga en cuenta que usted AHORA está en la ayuda de la consola interactiva de PYTHON. Por tanto el prompt ahora es:
\sphinxcode{\sphinxupquote{help>}}
\end{sphinxadmonition}

\item {} 
\sphinxAtStartPar
EJECUTE el comando: \sphinxcode{\sphinxupquote{topics}}. Cuente el número de temas que aparece en la tabla y escriba dicho número en su cuaderno. (Tenga en cuenta que algunos elementos de dicha tabla, están vacíos).

\item {} 
\sphinxAtStartPar
En los elementos de la tabla TOPICS, identifique: \sphinxstylestrong{BOOLEAN}, \sphinxstylestrong{FLOAT}, \sphinxstylestrong{INTEGER, NUMBERS} y \sphinxstylestrong{STRINGS}. (Guiarse por la primera letra ya que los elementos están organizados en orden alfabético).

\item {} 
\sphinxAtStartPar
EJECUTE el comando: \sphinxcode{\sphinxupquote{NUMBERS}}. Traduzca el contenido que aparece y escríbalo en su cuaderno.

\item {} 
\sphinxAtStartPar
EJECUTE el comando: \sphinxcode{\sphinxupquote{INTEGER}}. IDENTIFIQUE el tercer párrafo (\sphinxstyleemphasis{There is no limit…}), traducirlo y escribirlo en su cuaderno. IDENTIFIQUE también los diez ejemplos en el sexto párrafo (\sphinxstyleemphasis{Some examples of…}) y escribir dichos ejemplos en el cuaderno. (Para salir, PULSAR la tecla \sphinxstylestrong{Q}).

\item {} 
\sphinxAtStartPar
EJECUTE el comando: \sphinxcode{\sphinxupquote{FLOAT}}. IDENTIFIQUE los siete ejemplos del cuarto párrafo (\sphinxstyleemphasis{Some examples of…}) y escribirlos en su cuaderno.

\item {} 
\sphinxAtStartPar
EJECUTE el comando: \sphinxcode{\sphinxupquote{STRINGS}}. IDENTIFIQUE el cuarto párrafo (\sphinxstyleemphasis{In plain English…}), traducirlo y escribirlo en su cuaderno.

\item {} 
\sphinxAtStartPar
EJECUTE el comando: \sphinxcode{\sphinxupquote{BOOLEAN}}. IDENTIFIQUE el segundo párrafo (\sphinxstyleemphasis{In the context…}), traducirlo y escribirlo en su cuaderno.

\item {} 
\sphinxAtStartPar
EJECUTE el comando: \sphinxcode{\sphinxupquote{keywords}}. En su cuaderno ESCRIBA los elementos que aparecen en la tabla.

\item {} 
\sphinxAtStartPar
EJECUTE el comando: \sphinxcode{\sphinxupquote{symbols}}. En su cuaderno ESCRIBA los elementos que aparecen en la tabla.

\item {} 
\sphinxAtStartPar
CIERRE la consola interactiva de PYTHON y la terminal de sistema.

\end{enumerate}


\bigskip\hrule\bigskip



\subsection{A05: Codificación (+3,0)}
\label{\detokenize{1_guias/10programacion/G1024:a05-codificacion-3-0}}
\sphinxAtStartPar
\DUrole{sd-sphinx-override}{\DUrole{sd-badge}{\DUrole{sd-bg-danger}{\DUrole{sd-bg-text-danger}{Transferencia}}}}


\begin{savenotes}\sphinxattablestart
\sphinxthistablewithglobalstyle
\centering
\begin{tabulary}{\linewidth}[t]{TT}
\sphinxtoprule
\sphinxtableatstartofbodyhook
\sphinxAtStartPar
\sphinxstylestrong{TIEMPO}
&
\sphinxAtStartPar
5 minutos
\\
\sphinxhline
\sphinxAtStartPar
\sphinxstylestrong{EJERCICIOS}
&
\sphinxAtStartPar
3
\\
\sphinxhline
\sphinxAtStartPar
\sphinxstylestrong{OBJETIVO}
&
\sphinxAtStartPar
Comprobar conocimientos adquiridos.
\\
\sphinxhline
\sphinxAtStartPar
\sphinxstylestrong{ACTIVIDAD}
&
\sphinxAtStartPar
Ejecute los comandos relacionados con la consola interactiva de PYTHON.
\\
\sphinxbottomrule
\end{tabulary}
\sphinxtableafterendhook\par
\sphinxattableend\end{savenotes}

\begin{sphinxadmonition}{note}{Posibles ejercicios}
\begin{itemize}
\item {} 
\sphinxAtStartPar
Abra una terminal de sistema e ingrese a la consola de PYTHON.

\item {} 
\sphinxAtStartPar
IDENTIFIQUE el prompt de la consola de PYTHON.

\item {} 
\sphinxAtStartPar
INDIQUE la versión y la fecha de compilación de la consola de PYTHON.

\item {} 
\sphinxAtStartPar
ACCEDA a los derechos de autor(copyright) de la consola de PYTHON.

\item {} 
\sphinxAtStartPar
ACCEDA a los créditos(credits) de la consola de PYTHON.

\item {} 
\sphinxAtStartPar
ACCEDA al documento de la licencia de la consola de PYTHON y desplace la documentación hasta el final.

\item {} 
\sphinxAtStartPar
En la documentación de la licencia de la consola de PYTHON, IDENTIFIQUE la tabla relacionada con las versiones de PYTHON.

\item {} 
\sphinxAtStartPar
En la documentación de la licencia de la consola de PYTHON, IDENTIFIQUE las fechas históricas.

\item {} 
\sphinxAtStartPar
En la documentación de la licencia de la consola de PYTHON, IDENTIFIQUE el significado de las siglas: CWI, CNRI o PSF.

\item {} 
\sphinxAtStartPar
ACCEDA a la ayuda de la consola de PYTHON.

\item {} 
\sphinxAtStartPar
En la ayuda de la consola de PYTHON, ACCEDA a las diferentes temáticas(topics).

\item {} 
\sphinxAtStartPar
En la ayuda de la consola de PYTHON, ACCEDA a la temática relacionada con cualquiera de los tipos de variables en PYTHON.

\item {} 
\sphinxAtStartPar
En la ayuda de la consola de PYTHON, ACCEDA a los diferentes tipos de palabras reservadas(keywords).

\item {} 
\sphinxAtStartPar
En la ayuda de la consola de PYTHON, ACCEDA a los diferentes tipos de símbolos.

\end{itemize}
\end{sphinxadmonition}

\sphinxstepscope


\section{Guía 1025: Variables en Python (P4)}
\label{\detokenize{1_guias/10programacion/G1025:guia-1025-variables-en-python-p4}}\label{\detokenize{1_guias/10programacion/G1025::doc}}

\begin{savenotes}\sphinxattablestart
\sphinxthistablewithglobalstyle
\centering
\begin{tabular}[t]{\X{35}{100}\X{65}{100}}
\sphinxtoprule
\sphinxtableatstartofbodyhook

\begin{savenotes}\sphinxattablestart
\sphinxthistablewithglobalstyle
\centering
\begin{tabulary}{\linewidth}[t]{T}
\sphinxtoprule
\sphinxtableatstartofbodyhook
\sphinxAtStartPar
\sphinxincludegraphics{{escudo}.png}
\\
\sphinxbottomrule
\end{tabulary}
\sphinxtableafterendhook\par
\sphinxattableend\end{savenotes}
&

\begin{savenotes}\sphinxattablestart
\sphinxthistablewithglobalstyle
\raggedright
\begin{tabular}[t]{*{1}{\X{1}{1}}}
\sphinxtoprule
\sphinxtableatstartofbodyhook
\sphinxAtStartPar
\sphinxstylestrong{INSTITUCIÓN EDUCATIVA}: John F. Kennedy
\begin{itemize}
\item {} 
\sphinxAtStartPar
\sphinxstylestrong{DOCENTE}: Julian Camilo Fonseca Romero

\item {} 
\sphinxAtStartPar
\sphinxstylestrong{ASIGNATURA}: Programación

\item {} 
\sphinxAtStartPar
\sphinxstylestrong{GRADO}: ONCE

\item {} 
\sphinxAtStartPar
\sphinxstylestrong{PERIODO}: 04

\item {} 
\sphinxAtStartPar
\sphinxstylestrong{TIEMPO}: 10 semanas (2 horas / semana)

\item {} 
\sphinxAtStartPar
\sphinxstylestrong{TIEMPO DE TRABAJO}: 20/20 horas

\end{itemize}
\\
\sphinxbottomrule
\end{tabular}
\sphinxtableafterendhook\par
\sphinxattableend\end{savenotes}
\\
\sphinxbottomrule
\end{tabular}
\sphinxtableafterendhook\par
\sphinxattableend\end{savenotes}

\begin{sphinxuseclass}{sd-tab-set}
\begin{sphinxuseclass}{sd-tab-item}\subsubsection*{Objetivo}

\begin{sphinxuseclass}{sd-tab-content}
\sphinxAtStartPar
\DUrole{sd-sphinx-override}{\DUrole{sd-badge}{\DUrole{sd-bg-light}{\DUrole{sd-bg-text-light}{Exploración}}}}
\DUrole{sd-sphinx-override}{\DUrole{sd-badge}{\DUrole{sd-bg-info}{\DUrole{sd-bg-text-info}{Estructuración}}}}
\DUrole{sd-sphinx-override}{\DUrole{sd-badge}{\DUrole{sd-bg-warning}{\DUrole{sd-bg-text-warning}{Actividad práctica}}}}
\DUrole{sd-sphinx-override}{\DUrole{sd-badge}{\DUrole{sd-bg-danger}{\DUrole{sd-bg-text-danger}{Transferencia}}}}
\begin{itemize}
\item {} 
\sphinxAtStartPar
\sphinxstylestrong{Temática}: Variables y operaciones en PYTHON.

\item {} 
\sphinxAtStartPar
\sphinxstylestrong{Objetivo de aprendizaje}: Comprender, modelar y aplicar el concepto de variables y operadores aritméticos en el lenguaje de programación Python.

\end{itemize}

\end{sphinxuseclass}
\end{sphinxuseclass}
\begin{sphinxuseclass}{sd-tab-item}\subsubsection*{Orientaciones curriculares}

\begin{sphinxuseclass}{sd-tab-content}\begin{itemize}
\item {} 
\sphinxAtStartPar
\sphinxstylestrong{Componente}: Solución de problemas con T\&I.

\item {} 
\sphinxAtStartPar
\sphinxstylestrong{Competencia}: Propongo innovaciones tecnológicas e informáticas para la solución de problemas dando cumplimiento a restricciones, condiciones y especificaciones técnicas y contextuales.

\item {} 
\sphinxAtStartPar
\sphinxstylestrong{Evidencia de aprendizaje}: Identifico condiciones, especificaciones y restricciones de diseño, utilizadas en una solución tecnológica o del campo de la informática verificando su cumplimiento en diversos contextos.

\end{itemize}

\end{sphinxuseclass}
\end{sphinxuseclass}
\begin{sphinxuseclass}{sd-tab-item}\subsubsection*{Criterios de evaluación}

\begin{sphinxuseclass}{sd-tab-content}\begin{itemize}
\item {} 
\sphinxAtStartPar
Actividades desarrolladas en su totalidad.

\item {} 
\sphinxAtStartPar
Participación en clase.

\item {} 
\sphinxAtStartPar
Explicación clara de conceptos.

\item {} 
\sphinxAtStartPar
Argumentación de ideas.

\item {} 
\sphinxAtStartPar
Cumplimiento de las reglas del aula.

\end{itemize}

\end{sphinxuseclass}
\end{sphinxuseclass}
\begin{sphinxuseclass}{sd-tab-item}\subsubsection*{Recursos (2)}

\begin{sphinxuseclass}{sd-tab-content}\begin{itemize}
\item {} 
\sphinxAtStartPar
\sphinxhref{https://www.online-python.com}{PYTHON ONLINE}. Ejercutar la consola interactiva PYTHON en el navegador.

\item {} 
\sphinxAtStartPar
\sphinxhref{https://automatetheboringstuff.com}{Automate the Boring Stuff with Python}. Libro digital en inglés.

\end{itemize}

\end{sphinxuseclass}
\end{sphinxuseclass}
\end{sphinxuseclass}

\subsection{Mensaje del autor}
\label{\detokenize{1_guias/10programacion/G1025:mensaje-del-autor}}\begin{quote}

\sphinxAtStartPar
“Esta guía didáctica ha sido diseñada para facilitar la comprensión de las variables en Python, enfocándose en su definición, tipos y uso básico. El objetivo es que los estudiantes adquieran una base sólida para manejar datos y operaciones en programación, fomentando un aprendizaje claro y práctico que les permita avanzar con confianza en el mundo de la programación”.

\begin{flushright}
---Camilo Fonseca
\end{flushright}
\end{quote}


\bigskip\hrule\bigskip



\subsection{A00: Introducción (+0,1)}
\label{\detokenize{1_guias/10programacion/G1025:a00-introduccion-0-1}}
\sphinxAtStartPar
\DUrole{sd-sphinx-override}{\DUrole{sd-badge}{\DUrole{sd-bg-light}{\DUrole{sd-bg-text-light}{Exploración}}}}
\begin{enumerate}
\sphinxsetlistlabels{\arabic}{enumi}{enumii}{}{.}%
\item {} 
\sphinxAtStartPar
TRANSCRIBA en su cuaderno el objetivo, las orientaciones curriculares y los criterios de evaluación de la presente guía.

\item {} 
\sphinxAtStartPar
LEA el \sphinxstylestrong{mensaje del autor} y responda la siguiente pregunta en su cuaderno:
\begin{itemize}
\item {} 
\sphinxAtStartPar
¿Cómo los datos y las operaciones se relacionan con la programación?

\end{itemize}

\end{enumerate}


\bigskip\hrule\bigskip



\subsection{A01: Variables y operaciones (+0,1)}
\label{\detokenize{1_guias/10programacion/G1025:a01-variables-y-operaciones-0-1}}
\sphinxAtStartPar
\DUrole{sd-sphinx-override}{\DUrole{sd-badge}{\DUrole{sd-bg-info}{\DUrole{sd-bg-text-info}{Estructuración}}}}
\begin{enumerate}
\sphinxsetlistlabels{\arabic}{enumi}{enumii}{}{.}%
\item {} 
\sphinxAtStartPar
Lea los siguientes conceptos:

\end{enumerate}


\bigskip\hrule\bigskip



\subsubsection{¿Qué son las variables?}
\label{\detokenize{1_guias/10programacion/G1025:que-son-las-variables}}
\sphinxAtStartPar
%
\begin{footnote}[1]\sphinxAtStartFootnote
Sweigart, A. (2019). Automate the boring stuff with python, 2nd edition: Practical programming for total beginners. No Starch Press.
%
\end{footnote} Las \sphinxstylestrong{variables} almacenan información. Es por eso que a través de las variables es que se accede a la memoria RAM del computador. Tener en cuenta lo siguiente:
\begin{quote}

\sphinxAtStartPar
La diferencia entre una variable y otra radica en el \sphinxstylestrong{tipo de contenido} que esta almacena.
\end{quote}
\begin{quote}

\sphinxAtStartPar
El contenido de una variable puede \sphinxstylestrong{cambiar} cuantas veces sea necesario.
\end{quote}


\begin{savenotes}\sphinxattablestart
\sphinxthistablewithglobalstyle
\centering
\begin{tabulary}{\linewidth}[t]{T}
\sphinxtoprule
\sphinxtableatstartofbodyhook

\begin{sphinxfigure-in-table}
\centering

\noindent\sphinxincludegraphics[width=250\sphinxpxdimen]{{image0113}.png}
\end{sphinxfigure-in-table}\relax
\\
\sphinxbottomrule
\end{tabulary}
\sphinxtableafterendhook\par
\sphinxattableend\end{savenotes}

\sphinxAtStartPar
Una variable se compone de:
\begin{itemize}
\item {} 
\sphinxAtStartPar
El nombre de la variable.

\item {} 
\sphinxAtStartPar
Un signo igual (=) (llamado operador de asignación).

\item {} 
\sphinxAtStartPar
Un valor a almacenar.

\end{itemize}

\sphinxAtStartPar
En este ejemplo, la caja puede ser modelada así:

\begin{sphinxVerbatim}[commandchars=\\\{\},numbers=left,firstnumber=1,stepnumber=1]
\PYG{n}{spam} \PYG{o}{=} \PYG{l+m+mi}{42}
\end{sphinxVerbatim}

\sphinxAtStartPar
Igualmente las variables pueden cambiar su valor cuantas veces sea necesario. Por ejemplo la variable spam puede ser:

\begin{sphinxVerbatim}[commandchars=\\\{\},numbers=left,firstnumber=1,stepnumber=1]
\PYG{n}{spam} \PYG{o}{=} \PYG{l+m+mi}{42}
\PYG{n}{spam} \PYG{o}{=} \PYG{l+m+mf}{9.8}
\PYG{n}{spam} \PYG{o}{=} \PYG{l+s+s2}{\PYGZdq{}}\PYG{l+s+s2}{hola}\PYG{l+s+s2}{\PYGZdq{}}
\PYG{n}{spam} \PYG{o}{=} \PYG{k+kc}{True}
\end{sphinxVerbatim}


\bigskip\hrule\bigskip



\paragraph{¿Cómo se almacenan los datos en las variables?}
\label{\detokenize{1_guias/10programacion/G1025:como-se-almacenan-los-datos-en-las-variables}}
\sphinxAtStartPar
Los datos se llevan a las variables a través del signo (=).
\begin{quote}

\sphinxAtStartPar
En el caso de los \sphinxstyleemphasis{\sphinxstylestrong{algoritmos computacionales y programas}}, el signo (=) tiene una connotación diferente a la que se le da en matemáticas. El signo igual (=) significa que el computador va a realizar lo que está a la derecha del igual y lo va a almacenar en la variable que se encuentre a la izquierda del igual.
\end{quote}

\sphinxAtStartPar
Al momento de asignar valores correctamente a las variables, se debe tener en cuenta lo siguiente:
\begin{itemize}
\item {} 
\sphinxAtStartPar
Al lado izquierdo del igual (=) solo puede haber un nombre de variable.

\item {} 
\sphinxAtStartPar
Al lado derecho del igual (=) puede haber un valor constante, un valor variable o una expresión.

\item {} 
\sphinxAtStartPar
El computador siempre resuelve lo de la derecha del igual y su resultado lo almacena en la variable que esté a la izquierda del igual.

\end{itemize}
\begin{quote}

\sphinxAtStartPar
Cada vez que se asigna un nuevo valor a una variable, \sphinxstyleemphasis{\sphinxstylestrong{el valor anterior se pierde}}.
\end{quote}

\sphinxAtStartPar
Por ejemplo:


\begin{savenotes}\sphinxattablestart
\sphinxthistablewithglobalstyle
\centering
\begin{tabular}[t]{\X{35}{200}\X{165}{200}}
\sphinxtoprule
\sphinxstyletheadfamily 
\sphinxAtStartPar
Variables
&\sphinxstyletheadfamily 
\sphinxAtStartPar
Descripción
\\
\sphinxmidrule
\sphinxtableatstartofbodyhook
\sphinxAtStartPar
\sphinxcode{\sphinxupquote{A = 8}}
&
\sphinxAtStartPar
Le indica al computador que guarde el valor 8 en la variable de nombre \sphinxstylestrong{A}.
\\
\sphinxhline
\sphinxAtStartPar
\sphinxcode{\sphinxupquote{B = A}}
&
\sphinxAtStartPar
Le indica al computador que guarde en la variable \sphinxstylestrong{B} el contenido de la variable \sphinxstylestrong{A} que en la instrucción había sido “cargada” con 8, por lo tanto en la variable \sphinxstylestrong{B} queda el valor de 8 al igual que en la variable \sphinxstylestrong{B}.
\\
\sphinxhline
\sphinxAtStartPar
\sphinxcode{\sphinxupquote{C = A + B}}
&
\sphinxAtStartPar
Le indica al computador que en la variable \sphinxstylestrong{C} guarde el resultado de sumar el contenido de la variable \sphinxstylestrong{A} con el contenido de la variable \sphinxstylestrong{B}. Como la variable \sphinxstylestrong{A} tenía el contenido 8 y la variable \sphinxstylestrong{B} también tenía el contenido 8 entonces el computador sumara 8+8 y ese 16 de resultado lo almacenará en la variable \sphinxstylestrong{C}.
\\
\sphinxbottomrule
\end{tabular}
\sphinxtableafterendhook\par
\sphinxattableend\end{savenotes}


\bigskip\hrule\bigskip



\subsubsection{Tipos de variables}
\label{\detokenize{1_guias/10programacion/G1025:tipos-de-variables}}
\sphinxAtStartPar
Las variables son de 4 tipos:

\begin{sphinxadmonition}{note}{ENTERO (int)}

\sphinxAtStartPar
Un dato de tipo entero es un número \sphinxstylestrong{QUE NO TIENE PUNTO DECIMAL}, por lo tanto en sus operaciones jamás va a generar decimales. Una variable que se declare de tipo entero podrá almacenar solamente datos de tipo entero. Las variables enteras cumplen con las reglas de la Aritmética Entera.

\sphinxAtStartPar
Las variables de tipo entero pueden ser positivas o negativas.

\sphinxAtStartPar
Ejemplos de variables enteras son:

\begin{sphinxVerbatim}[commandchars=\\\{\},numbers=left,firstnumber=1,stepnumber=1]
\PYG{n}{a} \PYG{o}{=} \PYG{o}{\PYGZhy{}}\PYG{l+m+mi}{2}
\PYG{n}{b} \PYG{o}{=} \PYG{o}{\PYGZhy{}}\PYG{l+m+mi}{100}
\PYG{n}{e} \PYG{o}{=} \PYG{l+m+mi}{0}
\PYG{n}{numero} \PYG{o}{=} \PYG{l+m+mi}{20}
\PYG{n}{H} \PYG{o}{=} \PYG{l+m+mi}{34}
\PYG{n}{edad} \PYG{o}{=} \PYG{l+m+mi}{12}
\end{sphinxVerbatim}
\end{sphinxadmonition}

\begin{sphinxadmonition}{note}{DECIMAL (float)}

\sphinxAtStartPar
También llamadas variables de tipo flotante o simplemente float.

\sphinxAtStartPar
Una variable de tipo float ES UN NUMERO QUE TIENE PUNTO DECIMAL, por lo tanto en sus operaciones puede generar decimales.

\sphinxAtStartPar
Una variable que se declare de tipo float, solo podrá almacenar datos de tipo decimal.

\sphinxAtStartPar
Las variables de tipo float pueden ser positivas o negativas.

\sphinxAtStartPar
Ejemplos de variables float son:

\begin{sphinxVerbatim}[commandchars=\\\{\},numbers=left,firstnumber=1,stepnumber=1]
\PYG{n}{A} \PYG{o}{=} \PYG{o}{\PYGZhy{}}\PYG{l+m+mf}{1.25}
\PYG{n}{B} \PYG{o}{=} \PYG{o}{\PYGZhy{}}\PYG{l+m+mf}{1.0}
\PYG{n}{temperatura} \PYG{o}{=} \PYG{o}{\PYGZhy{}}\PYG{l+m+mf}{0.5}
\PYG{n}{velocidad} \PYG{o}{=} \PYG{l+m+mf}{0.0}
\PYG{n}{obs} \PYG{o}{=} \PYG{l+m+mf}{0.5}
\PYG{n}{sueldo} \PYG{o}{=} \PYG{l+m+mf}{1.2555555559}
\PYG{n}{PI} \PYG{o}{=} \PYG{l+m+mf}{3.1416}
\end{sphinxVerbatim}

\begin{sphinxadmonition}{note}{No confundir variables tipo \sphinxstylestrong{float} con \sphinxstylestrong{int}}

\sphinxAtStartPar
Se pueden dar variables así:

\begin{sphinxVerbatim}[commandchars=\\\{\},numbers=left,firstnumber=1,stepnumber=1]
\PYG{n}{A} \PYG{o}{=} \PYG{l+m+mi}{5}
\PYG{n}{B} \PYG{o}{=} \PYG{l+m+mf}{5.0}
\end{sphinxVerbatim}

\sphinxAtStartPar
Pese a que A y B contienen el mismo número 5; A y B son variables diferentes porque A contiene un entero y B contiene un float.
\end{sphinxadmonition}
\end{sphinxadmonition}

\begin{sphinxadmonition}{note}{CADENA (string)}

\sphinxAtStartPar
Cuando se tiene un conjunto de caracteres se dice técnicamente que se tiene una cadena, por lo tanto el nombre \sphinxcode{\sphinxupquote{“KENNEDY”}} es una cadena o string.

\sphinxAtStartPar
El contenido de una cadena se acostumbra encerrarlo entre comillas dobles.

\sphinxAtStartPar
Ejemplos de variables de tipo string:

\begin{sphinxVerbatim}[commandchars=\\\{\},numbers=left,firstnumber=1,stepnumber=1]
\PYG{n}{saludo} \PYG{o}{=} \PYG{l+s+s2}{\PYGZdq{}}\PYG{l+s+s2}{buenos días}\PYG{l+s+s2}{\PYGZdq{}}
\PYG{n}{nombre} \PYG{o}{=} \PYG{l+s+s2}{\PYGZdq{}}\PYG{l+s+s2}{Camilo}\PYG{l+s+s2}{\PYGZdq{}}
\PYG{n}{apellido} \PYG{o}{=} \PYG{l+s+s2}{\PYGZdq{}}\PYG{l+s+s2}{Fonseca}\PYG{l+s+s2}{\PYGZdq{}}
\PYG{n}{mensaje} \PYG{o}{=} \PYG{l+s+s2}{\PYGZdq{}}\PYG{l+s+s2}{Bienvenido a la sala 115}\PYG{l+s+s2}{\PYGZdq{}}
\PYG{n}{entrada} \PYG{o}{=} \PYG{l+s+s2}{\PYGZdq{}}\PYG{l+s+s2}{desea continuar con la operación? (si/no): }\PYG{l+s+s2}{\PYGZdq{}}
\end{sphinxVerbatim}

\begin{sphinxadmonition}{note}{No confundir variables tipo \sphinxstylestrong{string} con \sphinxstylestrong{int} o \sphinxstylestrong{float}}

\sphinxAtStartPar
Se pueden dar variables así:

\begin{sphinxVerbatim}[commandchars=\\\{\},numbers=left,firstnumber=1,stepnumber=1]
\PYG{n}{A} \PYG{o}{=} \PYG{l+s+s2}{\PYGZdq{}}\PYG{l+s+s2}{12}\PYG{l+s+s2}{\PYGZdq{}}
\PYG{n}{B} \PYG{o}{=} \PYG{l+m+mi}{12}
\end{sphinxVerbatim}

\sphinxAtStartPar
Pese a que A y B contienen el mismo número 12; A y B son variables diferentes porque A contiene un string y B contiene un entero.
\end{sphinxadmonition}
\end{sphinxadmonition}

\begin{sphinxadmonition}{note}{BOOLEAN}

\sphinxAtStartPar
Este tipo de variables solo pueden almacenar dos valores: \sphinxcode{\sphinxupquote{True}} o \sphinxcode{\sphinxupquote{False}}.

\sphinxAtStartPar
Ejemplos de variables tipo BOOLEANA:

\begin{sphinxVerbatim}[commandchars=\\\{\},numbers=left,firstnumber=1,stepnumber=1]
\PYG{n}{FACIL\PYGZus{}APRENDER} \PYG{o}{=} \PYG{k+kc}{True}
\PYG{n}{ES\PYGZus{}INDEPENDIENTE} \PYG{o}{=} \PYG{k+kc}{False}
\PYG{n}{ES\PYGZus{}GRATUITO} \PYG{o}{=} \PYG{k+kc}{False}
\end{sphinxVerbatim}

\begin{sphinxadmonition}{note}{Variables tipo boolean para tomar decisiones}

\sphinxAtStartPar
Al ser para toma de decisiones, entonces las variables booleanas se usan para:
\begin{itemize}
\item {} 
\sphinxAtStartPar
Condicionales

\item {} 
\sphinxAtStartPar
Bucles finitos o infinitos

\end{itemize}
\end{sphinxadmonition}
\end{sphinxadmonition}

\begin{sphinxadmonition}{note}{Ejemplos tipo de variables}


\begin{savenotes}\sphinxattablestart
\sphinxthistablewithglobalstyle
\centering
\begin{tabulary}{\linewidth}[t]{TT}
\sphinxtoprule
\sphinxstyletheadfamily 
\sphinxAtStartPar
Tipo de variables
&\sphinxstyletheadfamily 
\sphinxAtStartPar
Ejemplo
\\
\sphinxmidrule
\sphinxtableatstartofbodyhook
\sphinxAtStartPar
Enteros (int)
&
\sphinxAtStartPar
\sphinxhyphen{}2, \sphinxhyphen{}1, 0, 1, 2, 3, 4, 5
\\
\sphinxhline
\sphinxAtStartPar
Decimales (float)
&
\sphinxAtStartPar
\sphinxhyphen{}1.25, \sphinxhyphen{}1.0, \sphinxhyphen{}0.5, 0.0, 0.5, 1.0, 1.25
\\
\sphinxhline
\sphinxAtStartPar
Cadenas (strings)
&
\sphinxAtStartPar
“a”, “aa”, “aaa”, “Hello!”, “11 cats”
\\
\sphinxhline
\sphinxAtStartPar
Booleanos
&
\sphinxAtStartPar
True, False
\\
\sphinxbottomrule
\end{tabulary}
\sphinxtableafterendhook\par
\sphinxattableend\end{savenotes}
\end{sphinxadmonition}


\bigskip\hrule\bigskip



\subsubsection{Reglas para nombrar variables}
\label{\detokenize{1_guias/10programacion/G1025:reglas-para-nombrar-variables}}
\sphinxAtStartPar
Para nombrar correctamente una variable, se siguen las siguientes reglas:
\begin{itemize}
\item {} 
\sphinxAtStartPar
Solo puede ser una palabra sin espacios.

\item {} 
\sphinxAtStartPar
Solo pueden tener letras, números y barra al piso \sphinxcode{\sphinxupquote{\_}}.

\item {} 
\sphinxAtStartPar
Un nombre de variable NUNCA puede comenzar por un número.

\end{itemize}

\sphinxAtStartPar
A continuación se muestran ejemplos de nombramiento correcto e incorrecto de variables:


\begin{savenotes}\sphinxattablestart
\sphinxthistablewithglobalstyle
\centering
\begin{tabular}[t]{\X{30}{100}\X{70}{100}}
\sphinxtoprule
\sphinxstyletheadfamily 
\sphinxAtStartPar
Nombres de variables correctos
&\sphinxstyletheadfamily 
\sphinxAtStartPar
Nombres de variables incorrectos
\\
\sphinxmidrule
\sphinxtableatstartofbodyhook
\sphinxAtStartPar
\sphinxcode{\sphinxupquote{current\_balance}}
&
\sphinxAtStartPar
\sphinxcode{\sphinxupquote{current\sphinxhyphen{}balance}} (No se permiten guiones)
\\
\sphinxhline
\sphinxAtStartPar
\sphinxcode{\sphinxupquote{currentBalance}}
&
\sphinxAtStartPar
\sphinxcode{\sphinxupquote{current balance}} (No se permiten espacios)
\\
\sphinxhline
\sphinxAtStartPar
\sphinxcode{\sphinxupquote{account4}}
&
\sphinxAtStartPar
\sphinxcode{\sphinxupquote{4account}} (No se permite iniciar con un número)
\\
\sphinxhline
\sphinxAtStartPar
\sphinxcode{\sphinxupquote{\_42}}
&
\sphinxAtStartPar
\sphinxcode{\sphinxupquote{42}} (No puede ser un número)
\\
\sphinxhline
\sphinxAtStartPar
\sphinxcode{\sphinxupquote{TOTAL\_SUM}}
&
\sphinxAtStartPar
\sphinxcode{\sphinxupquote{TOTAL\_\$UM}} (caracteres especiales como \$ no se permiten)
\\
\sphinxhline
\sphinxAtStartPar
\sphinxcode{\sphinxupquote{hello}}
&
\sphinxAtStartPar
\sphinxcode{\sphinxupquote{'hello'}} (caracteres especiales como \sphinxcode{\sphinxupquote{'}} o \sphinxcode{\sphinxupquote{""}} no se permiten)
\\
\sphinxbottomrule
\end{tabular}
\sphinxtableafterendhook\par
\sphinxattableend\end{savenotes}


\bigskip\hrule\bigskip



\subsubsection{Instrucciones en PYTHON}
\label{\detokenize{1_guias/10programacion/G1025:instrucciones-en-python}}
\sphinxAtStartPar
La PC ejecuta instrucciones.
\begin{itemize}
\item {} 
\sphinxAtStartPar
Estas instrucciones son ejecutadas de forma serial. Es decir una a la vez.

\item {} 
\sphinxAtStartPar
Las instrucciones son dadas por el programador.

\end{itemize}

\sphinxAtStartPar
En la \sphinxstylestrong{consola interactiva de PYTHON}, cada instrucción se ejecuta cuando se presiona la tecla \sphinxkeyboard{\sphinxupquote{ENTER}}.

\begin{sphinxadmonition}{note}{EJEMPLO: Uso de la función \sphinxstyleliteralintitle{\sphinxupquote{print()}}}

\sphinxAtStartPar
La función \sphinxcode{\sphinxupquote{print()}} muestra en pantalla valores de variables. Los valores a mostrar deben ir dentro del paréntesis de la función \sphinxcode{\sphinxupquote{print()}}.  Por ejemplo:

\begin{sphinxVerbatim}[commandchars=\\\{\},numbers=left,firstnumber=1,stepnumber=1]
\PYG{g+gp}{\PYGZgt{}\PYGZgt{}\PYGZgt{} }\PYG{n+nb}{print}\PYG{p}{(}\PYG{l+s+s2}{\PYGZdq{}}\PYG{l+s+s2}{buenas tardes}\PYG{l+s+s2}{\PYGZdq{}}\PYG{p}{)}
\PYG{g+go}{buenas tardes}
\end{sphinxVerbatim}

\sphinxAtStartPar
En este ejemplo, se le ordena al PC que muestre en pantalla el mensaje \sphinxcode{\sphinxupquote{“buenas tardes”}}.

\sphinxAtStartPar
Después de escribir completamente la instrucción, al PULSAR \sphinxcode{\sphinxupquote{ENTER}} se ejecuta dicha instrucción.
\end{sphinxadmonition}

\sphinxAtStartPar
Se pueden ejecutar varias instrucciones en forma serial. Por ejemplo:

\begin{sphinxVerbatim}[commandchars=\\\{\},numbers=left,firstnumber=1,stepnumber=1]
\PYG{g+gp}{\PYGZgt{}\PYGZgt{}\PYGZgt{} }\PYG{n+nb}{print}\PYG{p}{(}\PYG{l+s+s2}{\PYGZdq{}}\PYG{l+s+s2}{instruccion01}\PYG{l+s+s2}{\PYGZdq{}}\PYG{p}{)}
\PYG{g+go}{instruccion01}
\PYG{g+gp}{\PYGZgt{}\PYGZgt{}\PYGZgt{} }\PYG{n+nb}{print}\PYG{p}{(}\PYG{l+s+s2}{\PYGZdq{}}\PYG{l+s+s2}{instruccion02}\PYG{l+s+s2}{\PYGZdq{}}\PYG{p}{)}
\PYG{g+go}{instruccion02}
\PYG{g+gp}{\PYGZgt{}\PYGZgt{}\PYGZgt{} }\PYG{n+nb}{print}\PYG{p}{(}\PYG{l+s+s2}{\PYGZdq{}}\PYG{l+s+s2}{instruccion03}\PYG{l+s+s2}{\PYGZdq{}}\PYG{p}{)}
\PYG{g+go}{instruccion03}
\PYG{g+gp}{\PYGZgt{}\PYGZgt{}\PYGZgt{} }\PYG{n+nb}{print}\PYG{p}{(}\PYG{l+s+s2}{\PYGZdq{}}\PYG{l+s+s2}{instruccion04}\PYG{l+s+s2}{\PYGZdq{}}\PYG{p}{)}
\PYG{g+go}{instruccion04}
\PYG{g+gp}{\PYGZgt{}\PYGZgt{}\PYGZgt{} }\PYG{n+nb}{print}\PYG{p}{(}\PYG{l+s+s2}{\PYGZdq{}}\PYG{l+s+s2}{instruccion05}\PYG{l+s+s2}{\PYGZdq{}}\PYG{p}{)}
\PYG{g+go}{instruccion05}
\end{sphinxVerbatim}


\subsubsection{Operadores en PYTHON}
\label{\detokenize{1_guias/10programacion/G1025:operadores-en-python}}
\sphinxAtStartPar
Los operadores son signos que nos permiten expresar relaciones entre variables y/o constantes, de las cuales normalmente se desprende un resultado. En algoritmos computacionales se utilizan los siguientes operadores:


\begin{savenotes}\sphinxattablestart
\sphinxthistablewithglobalstyle
\centering
\begin{tabulary}{\linewidth}[t]{TT}
\sphinxtoprule
\sphinxstyletheadfamily 
\sphinxAtStartPar
OPERADOR
&\sphinxstyletheadfamily 
\sphinxAtStartPar
Descripción
\\
\sphinxmidrule
\sphinxtableatstartofbodyhook
\sphinxAtStartPar
\sphinxcode{\sphinxupquote{+}}
&
\sphinxAtStartPar
Operador para expresar la suma
\\
\sphinxhline
\sphinxAtStartPar
\sphinxcode{\sphinxupquote{\sphinxhyphen{}}}
&
\sphinxAtStartPar
Operador para expresar la resta
\\
\sphinxhline
\sphinxAtStartPar
\sphinxcode{\sphinxupquote{*}}
&
\sphinxAtStartPar
Operador para expresar la multiplicación
\\
\sphinxhline
\sphinxAtStartPar
\sphinxcode{\sphinxupquote{**}}
&
\sphinxAtStartPar
Operador para exponente
\\
\sphinxhline
\sphinxAtStartPar
\sphinxcode{\sphinxupquote{/}}
&
\sphinxAtStartPar
Operador para expresar la división
\\
\sphinxhline
\sphinxAtStartPar
\sphinxcode{\sphinxupquote{//}}
&
\sphinxAtStartPar
Operador para expresar división entera
\\
\sphinxhline
\sphinxAtStartPar
\sphinxcode{\sphinxupquote{\%}}
&
\sphinxAtStartPar
Operador módulo (para determinar residuo de división)
\\
\sphinxbottomrule
\end{tabulary}
\sphinxtableafterendhook\par
\sphinxattableend\end{savenotes}
\begin{quote}

\sphinxAtStartPar
No olvidar que cada \sphinxstylestrong{nuevo valor} que se le asigne a una variable ya establecida; el \sphinxstylestrong{valor anterior} sera \sphinxstylestrong{eliminado}.
\end{quote}

\sphinxAtStartPar
A continuación se muestra un \sphinxstylestrong{ejemplo de uso de operadores}:
\begin{quote}

\sphinxAtStartPar
Resolver el siguiente conjunto de operaciones en su cuaderno:
\begin{equation*}
\begin{split} a = 10 \\ b = 15 \\ c = 20 \end{split}
\end{equation*}\begin{equation*}
\begin{split} a = a + b \\ b = b + 8  \\ c = c + a \end{split}
\end{equation*}\begin{equation*}
\begin{split} a = a + 5 \\ b = b + 3 \\ c = c + 2 \end{split}
\end{equation*}\begin{equation*}
\begin{split} a = a - b \\ b = a - b \\ c = a - b \end{split}
\end{equation*}\end{quote}

\sphinxAtStartPar
\sphinxstylestrong{Solución}:

\begin{sphinxVerbatim}[commandchars=\\\{\},numbers=left,firstnumber=1,stepnumber=1]
\PYG{n}{a} \PYG{o}{=} \PYG{l+m+mi}{10}
\PYG{n}{b} \PYG{o}{=} \PYG{l+m+mi}{15}
\PYG{n}{c} \PYG{o}{=} \PYG{l+m+mi}{20}

\PYG{n}{a} \PYG{o}{=} \PYG{n}{a} \PYG{o}{+} \PYG{n}{b}
\PYG{n}{a} \PYG{o}{=} \PYG{l+m+mi}{10} \PYG{o}{+} \PYG{l+m+mi}{15}
\PYG{n}{a} \PYG{o}{=} \PYG{l+m+mi}{25} \PYG{c+c1}{\PYGZsh{} ahora a vale 25, es decir a ya no vale 10}

\PYG{n}{b} \PYG{o}{=} \PYG{n}{b} \PYG{o}{+} \PYG{l+m+mi}{8}
\PYG{n}{b} \PYG{o}{=} \PYG{l+m+mi}{15} \PYG{o}{+} \PYG{l+m+mi}{8}
\PYG{n}{b} \PYG{o}{=} \PYG{l+m+mi}{23} \PYG{c+c1}{\PYGZsh{} ahora b vale 23, es decir b ya no vale 15}

\PYG{n}{c} \PYG{o}{=} \PYG{n}{c} \PYG{o}{+} \PYG{n}{a}
\PYG{n}{c} \PYG{o}{=} \PYG{l+m+mi}{20} \PYG{o}{+} \PYG{l+m+mi}{25}
\PYG{n}{c} \PYG{o}{=} \PYG{l+m+mi}{45} \PYG{c+c1}{\PYGZsh{} ahora c vale 45, es decir c ya no vale 20}

\PYG{n}{a} \PYG{o}{=} \PYG{n}{a} \PYG{o}{+} \PYG{l+m+mi}{5}
\PYG{n}{a} \PYG{o}{=} \PYG{l+m+mi}{25} \PYG{o}{+} \PYG{l+m+mi}{5}
\PYG{n}{a} \PYG{o}{=} \PYG{l+m+mi}{30}

\PYG{n}{b} \PYG{o}{=} \PYG{n}{b} \PYG{o}{+} \PYG{l+m+mi}{3}
\PYG{n}{b} \PYG{o}{=} \PYG{l+m+mi}{23} \PYG{o}{+} \PYG{l+m+mi}{3}
\PYG{n}{b} \PYG{o}{=} \PYG{l+m+mi}{26}

\PYG{n}{c} \PYG{o}{=} \PYG{n}{c} \PYG{o}{+} \PYG{l+m+mi}{2}
\PYG{n}{c} \PYG{o}{=} \PYG{l+m+mi}{45} \PYG{o}{+} \PYG{l+m+mi}{2}
\PYG{n}{c} \PYG{o}{=} \PYG{l+m+mi}{47}

\PYG{n}{a} \PYG{o}{=} \PYG{n}{a} \PYG{o}{\PYGZhy{}} \PYG{n}{b}
\PYG{n}{a} \PYG{o}{=} \PYG{l+m+mi}{30} \PYG{o}{\PYGZhy{}} \PYG{l+m+mi}{26}
\PYG{n}{a} \PYG{o}{=} \PYG{l+m+mi}{4}

\PYG{n}{b} \PYG{o}{=} \PYG{n}{a} \PYG{o}{\PYGZhy{}} \PYG{n}{b}
\PYG{n}{b} \PYG{o}{=} \PYG{l+m+mi}{4} \PYG{o}{\PYGZhy{}} \PYG{l+m+mi}{26}
\PYG{n}{b} \PYG{o}{=} \PYG{o}{\PYGZhy{}}\PYG{l+m+mi}{22}

\PYG{n}{c} \PYG{o}{=} \PYG{n}{a} \PYG{o}{\PYGZhy{}} \PYG{n}{b}
\PYG{n}{c} \PYG{o}{=} \PYG{l+m+mi}{4} \PYG{o}{\PYGZhy{}} \PYG{p}{(}\PYG{o}{\PYGZhy{}}\PYG{l+m+mi}{22}\PYG{p}{)} \PYG{c+c1}{\PYGZsh{} por ley de los signos, 4 se suma con el 22}
\PYG{n}{c} \PYG{o}{=} \PYG{l+m+mi}{26}

\PYG{c+c1}{\PYGZsh{} RESPUESTA:}
\PYG{n}{a} \PYG{o}{=} \PYG{l+m+mi}{4}
\PYG{n}{b} \PYG{o}{=} \PYG{o}{\PYGZhy{}}\PYG{l+m+mi}{22}
\PYG{n}{c} \PYG{o}{=} \PYG{l+m+mi}{26}
\end{sphinxVerbatim}


\bigskip\hrule\bigskip



\subsubsection{Linealización de operaciones}
\label{\detokenize{1_guias/10programacion/G1025:linealizacion-de-operaciones}}
\sphinxAtStartPar
Algo a tener en cuenta al momento de escribir una instrucción es que el computador solo entiende las operaciones en formato lineal. Esto quiere decir escritas en una sola línea. De tal manera que si se quiere escribir la ecuación:
\begin{equation*}
\begin{split}
X = {a + b \over c + d}
\end{split}
\end{equation*}
\sphinxAtStartPar
La anterior ecuación en una sola línea sería:
\begin{equation*}
\begin{split}
X = (a + b) / (c + d)
\end{split}
\end{equation*}
\sphinxAtStartPar
A continuación se muestra un \sphinxstylestrong{ejemplo de linealización} de una operación:
\begin{quote}

\sphinxAtStartPar
Linealizar la siguiente operación en su cuaderno:
\begin{equation}\label{equation:1_guias/10programacion/G1025:1_guias/10programacion/G1025:0}
\begin{split}X = \frac{a+\frac{b}{c}-d}{a+\dfrac{b}{c^{d}+\dfrac{d}{a-\frac{b}{c}*d}}}\end{split}
\end{equation}\end{quote}

\sphinxAtStartPar
\sphinxstylestrong{Solución}:

\begin{sphinxadmonition}{note}{PASO 01}
\begin{equation}\label{equation:1_guias/10programacion/G1025:1_guias/10programacion/G1025:1}
\begin{split}X = \frac{a+\colorbox{yellow}{(b/c)}-d}{a+\dfrac{b}{c^{d}+\dfrac{d}{{\large a-\colorbox{yellow}{(b/c)}*d}}}}\end{split}
\end{equation}
\sphinxAtStartPar
Primero se linealizan las divisiones \((b/c)\) tanto en el numerador como en el denominador.
\end{sphinxadmonition}

\begin{sphinxadmonition}{note}{PASO 02}
\begin{equation}\label{equation:1_guias/10programacion/G1025:1_guias/10programacion/G1025:2}
\begin{split}X = \frac{a+\colorbox{yellow}{(b/c)}-d}{a+\dfrac{b}{c^{d}+{\color{Green} (d/(a-\colorbox{yellow}{(b/c)}*d))}}}\end{split}
\end{equation}
\sphinxAtStartPar
Se linealiza la expresión: \((d/(a-(b/c) * d))\).
\end{sphinxadmonition}

\begin{sphinxadmonition}{note}{PASO 03}
\begin{equation}\label{equation:1_guias/10programacion/G1025:1_guias/10programacion/G1025:3}
\begin{split}X = \frac{a+\colorbox{yellow}{(b/c)}-d}{a+{\color{Red} ({b}/({c^{d}+{\color{Green} (d/(a-\colorbox{yellow}{(b/c)}*d))}}))}}\end{split}
\end{equation}
\sphinxAtStartPar
Se linealiza la expresión: \((b/(c^d+(d/(a-(b/c)*d))))\).
\end{sphinxadmonition}

\begin{sphinxadmonition}{note}{PASO 04}
\begin{equation}\label{equation:1_guias/10programacion/G1025:1_guias/10programacion/G1025:4}
\begin{split}X = (a+\colorbox{yellow}{(b/c)}-d)/(a+{\color{Red} ({b}/({c^{d}+{\color{Green} (d/(a-\colorbox{yellow}{(b/c)}*d))}}))})\end{split}
\end{equation}
\sphinxAtStartPar
Ahora se linealiza tanto numerador como denominador.
\end{sphinxadmonition}

\begin{sphinxadmonition}{note}{PASO 05}
\begin{equation}\label{equation:1_guias/10programacion/G1025:1_guias/10programacion/G1025:5}
\begin{split}X = (a+\colorbox{yellow}{(b/c)}-d)/(a+{\color{black} ({b}/(\colorbox{orange}{c**{d}}+{\color{Green} (d/(a-\colorbox{yellow}{(b/c)}*d))}))})\end{split}
\end{equation}
\sphinxAtStartPar
Finalmente la exponenciación en código python se realiza con (c**d).
\end{sphinxadmonition}


\bigskip\hrule\bigskip



\subsubsection{Precedencia (jerarquía de operadores)}
\label{\detokenize{1_guias/10programacion/G1025:precedencia-jerarquia-de-operadores}}
\sphinxAtStartPar
Algunas veces se tienen operaciones muy largas. Esto dificulta su solución. En estos casos se aplica la precedencia. La cual consiste en unas reglas que determinan el orden en el que una operación muy larga debe ser resuelta. Un ejemplo de una operación muy larga sería:
\begin{equation*}
\begin{split}
X = ( a + (b / c) – d ) / ( a + ( b / ( (c ** d) + (d / ( a – (b / c) * d ) ) ) ) )
\end{split}
\end{equation*}
\begin{sphinxadmonition}{note}{Reglas de presedencia}

\sphinxAtStartPar
Los operadores aplican expresiones aritméticas en una secuencia precisa. Esta secuencia es determinada por las siguientes reglas de precedencia:
\begin{enumerate}
\sphinxsetlistlabels{\arabic}{enumi}{enumii}{}{.}%
\item {} 
\sphinxAtStartPar
\sphinxstylestrong{PARÉNTESIS \sphinxguilabel{()}}: Los operadores en las expresiones contenidas dentro de pares de paréntesis se evalúan primero. Se dice que los paréntesis tienen el “nivel más alto de precedencia”. En casos de paréntesis anidados o incrustados, como:

\end{enumerate}
\begin{equation*}
\begin{split}
( a * ( b + c ) )
\end{split}
\end{equation*}
\sphinxAtStartPar
Los operadores en el par más interno de paréntesis se aplican primero.
En este caso primero se opera \sphinxstylestrong{a} y \sphinxstylestrong{b}. El resultado se multiplica por \sphinxstylestrong{a.}
\begin{enumerate}
\sphinxsetlistlabels{\arabic}{enumi}{enumii}{}{.}%
\setcounter{enumi}{1}
\item {} 
\sphinxAtStartPar
\sphinxstylestrong{MULTIPLICACIÓN} \sphinxguilabel{*} , \sphinxstylestrong{DIVISIÓN} \sphinxguilabel{/} , \sphinxstylestrong{MÓDULO} \sphinxguilabel{\%}: Las operaciones de multiplicación, división y módulo se aplican a continuación. Si una expresión contiene varias de esas operaciones, los operadores se aplican de izquierda a derecha. Se dice que los operadores de multiplicación, división y módulo tienen el mismo nivel de precedencia.

\item {} 
\sphinxAtStartPar
\sphinxstylestrong{SUMA} \sphinxguilabel{+}, \sphinxstylestrong{RESTA} \sphinxguilabel{\sphinxhyphen{}}: Las operaciones de suma y resta se aplican de último. Si una expresión contiene varias de esas operaciones, los operadores se aplican de izquierda a derecha. Los operadores de suma y resta tienen el mismo nivel de precedencia.

\end{enumerate}
\end{sphinxadmonition}

\sphinxAtStartPar
Las \sphinxstylestrong{reglas de precedencia de operadores} definen el orden en el que PYTHON aplica los operadores. Cuando se dice que ciertos operadores se aplican de izquierda a derecha, esto hace referencia a su asociatividad. Por ejemplo, los operadores de suma (+) en la expresión:


\begin{savenotes}\sphinxattablestart
\sphinxthistablewithglobalstyle
\centering
\begin{tabular}[t]{\X{25}{100}\X{75}{100}}
\sphinxtoprule
\sphinxtableatstartofbodyhook
\sphinxAtStartPar
\(X = a + b + c\)
&
\sphinxAtStartPar
Esta operación se asocian de izquierda a derecha, por lo que a + b se calcula primero, y después se agrega c a esa suma para determinar el valor de X.
\\
\sphinxbottomrule
\end{tabular}
\sphinxtableafterendhook\par
\sphinxattableend\end{savenotes}

\begin{sphinxadmonition}{note}{Tabla de precedencia}


\begin{savenotes}\sphinxattablestart
\sphinxthistablewithglobalstyle
\centering
\begin{tabulary}{\linewidth}[t]{TTT}
\sphinxtoprule
\sphinxstyletheadfamily 
\sphinxAtStartPar
OPERADOR
&\sphinxstyletheadfamily 
\sphinxAtStartPar
OPERACIONES
&\sphinxstyletheadfamily 
\sphinxAtStartPar
Orden de evaluación (Precedencia)
\\
\sphinxmidrule
\sphinxtableatstartofbodyhook
\sphinxAtStartPar
( )
&
\sphinxAtStartPar
Paréntesis
&
\sphinxAtStartPar
Se evalúa primero. Si los paréntesis son anidados como en la expresión: \((a * ( b + c / d + e ) )\) la expresión en el par más interno se evalúa primero
\\
\sphinxhline
\sphinxAtStartPar
\sphinxcode{\sphinxupquote{*}} \sphinxcode{\sphinxupquote{**}} \sphinxcode{\sphinxupquote{/}} \sphinxcode{\sphinxupquote{//}} \sphinxcode{\sphinxupquote{\%}}
&
\sphinxAtStartPar
Multiplicación, Exponente División, División entera Módulo
&
\sphinxAtStartPar
Se evalúan en segundo lugar. Si hay varios operadores de este tipo, se evalúan de izquierda a derecha.
\\
\sphinxhline
\sphinxAtStartPar
\sphinxcode{\sphinxupquote{+}} \sphinxcode{\sphinxupquote{\sphinxhyphen{}}}
&
\sphinxAtStartPar
Suma, Resta
&
\sphinxAtStartPar
Se evalúan al último. Si hay varios operadores de este tipo, se evalúan de izquierda a derecha.
\\
\sphinxbottomrule
\end{tabulary}
\sphinxtableafterendhook\par
\sphinxattableend\end{savenotes}
\end{sphinxadmonition}


\bigskip\hrule\bigskip



\subsection{A02: Taller (+0,8)}
\label{\detokenize{1_guias/10programacion/G1025:a02-taller-0-8}}
\sphinxAtStartPar
\DUrole{sd-sphinx-override}{\DUrole{sd-badge}{\DUrole{sd-bg-info}{\DUrole{sd-bg-text-info}{Estructuración}}}}
\begin{enumerate}
\sphinxsetlistlabels{\arabic}{enumi}{enumii}{}{.}%
\item {} 
\sphinxAtStartPar
Responda las siguientes preguntas en su cuaderno:
\begin{itemize}
\item {} 
\sphinxAtStartPar
¿Qué es una variable?

\item {} 
\sphinxAtStartPar
¿Por qué por medio de las variables se accede a la memoria RAM del computador?

\item {} 
\sphinxAtStartPar
¿Cómo se pueden diferenciar las variables?

\end{itemize}

\item {} 
\sphinxAtStartPar
En su cuaderno, realice tres dibujos que representen el nombramiento y la asignación de valores de diferentes variables. (\sphinxstylestrong{OBSERVE EL EJEMPLO DEL DIBUJO DE LA CAJA}). \sphinxstyleemphasis{En cada dibujo, identifique el \sphinxstylestrong{nombre} de la variable y el \sphinxstylestrong{valor} de dicha variable}.

\item {} 
\sphinxAtStartPar
En media hoja de cuaderno, responda la siguiente pregunta:
\begin{itemize}
\item {} 
\sphinxAtStartPar
¿Por que el signo (=) tiene un significado en algoritmos computacionales y otro significado diferente en matemáticas?

\end{itemize}

\item {} 
\sphinxAtStartPar
Describa cuatro aspectos que se deben tener en cuenta al momento de \sphinxstylestrong{asignar valores a variables}.

\item {} 
\sphinxAtStartPar
Describa un ejemplo en el que una \sphinxstylestrong{variable pueda contener otra variable}.

\item {} 
\sphinxAtStartPar
Describa los \sphinxstylestrong{cuatros tipos de variables} que existen e indique un ejemplo por cada una.

\item {} 
\sphinxAtStartPar
Dibuje la siguiente tabla e indique que \sphinxstylestrong{tipo de variable} es y explique por qué.

\end{enumerate}


\begin{savenotes}\sphinxattablestart
\sphinxthistablewithglobalstyle
\centering
\begin{tabulary}{\linewidth}[t]{TTT}
\sphinxtoprule
\sphinxstyletheadfamily 
\sphinxAtStartPar
Ejemplo
&\sphinxstyletheadfamily 
\sphinxAtStartPar
Tipo de variable
&\sphinxstyletheadfamily 
\sphinxAtStartPar
¿Por qué?
\\
\sphinxmidrule
\sphinxtableatstartofbodyhook
\sphinxAtStartPar
\sphinxcode{\sphinxupquote{A = True}}
&
\sphinxAtStartPar

&
\sphinxAtStartPar

\\
\sphinxhline
\sphinxAtStartPar
\sphinxcode{\sphinxupquote{b = 12.33}}
&
\sphinxAtStartPar

&
\sphinxAtStartPar

\\
\sphinxhline
\sphinxAtStartPar
\sphinxcode{\sphinxupquote{c = “12”}}
&
\sphinxAtStartPar

&
\sphinxAtStartPar

\\
\sphinxhline
\sphinxAtStartPar
\sphinxcode{\sphinxupquote{d = 12.0000099}}
&
\sphinxAtStartPar

&
\sphinxAtStartPar

\\
\sphinxhline
\sphinxAtStartPar
\sphinxcode{\sphinxupquote{e = A}}
&
\sphinxAtStartPar

&
\sphinxAtStartPar

\\
\sphinxhline
\sphinxAtStartPar
\sphinxcode{\sphinxupquote{f = “Este es el contenido”}}
&
\sphinxAtStartPar

&
\sphinxAtStartPar

\\
\sphinxhline
\sphinxAtStartPar
\sphinxcode{\sphinxupquote{g = “3.141516”}}
&
\sphinxAtStartPar

&
\sphinxAtStartPar

\\
\sphinxhline
\sphinxAtStartPar
\sphinxcode{\sphinxupquote{h = 19999999999991}}
&
\sphinxAtStartPar

&
\sphinxAtStartPar

\\
\sphinxhline
\sphinxAtStartPar
\sphinxcode{\sphinxupquote{i = False}}
&
\sphinxAtStartPar

&
\sphinxAtStartPar

\\
\sphinxhline
\sphinxAtStartPar
\sphinxcode{\sphinxupquote{j = h = d = b}}
&
\sphinxAtStartPar

&
\sphinxAtStartPar

\\
\sphinxbottomrule
\end{tabulary}
\sphinxtableafterendhook\par
\sphinxattableend\end{savenotes}
\begin{enumerate}
\sphinxsetlistlabels{\arabic}{enumi}{enumii}{}{.}%
\setcounter{enumi}{7}
\item {} 
\sphinxAtStartPar
Dibuje la siguiente tabla e identifique cuales son la \sphinxstylestrong{variables nombradas de forma correcta} y cuales son las \sphinxstylestrong{variables nombradas de forma incorrecta}.

\end{enumerate}


\begin{savenotes}\sphinxattablestart
\sphinxthistablewithglobalstyle
\centering
\begin{tabulary}{\linewidth}[t]{TTT}
\sphinxtoprule
\sphinxstyletheadfamily 
\sphinxAtStartPar
Variables
&\sphinxstyletheadfamily 
\sphinxAtStartPar
correcta o incorrecta
&\sphinxstyletheadfamily 
\sphinxAtStartPar
¿Por qué?
\\
\sphinxmidrule
\sphinxtableatstartofbodyhook
\sphinxAtStartPar
\sphinxcode{\sphinxupquote{balanceActual}}
&
\sphinxAtStartPar

&
\sphinxAtStartPar

\\
\sphinxhline
\sphinxAtStartPar
\sphinxcode{\sphinxupquote{balance\_actual}}
&
\sphinxAtStartPar

&
\sphinxAtStartPar

\\
\sphinxhline
\sphinxAtStartPar
\sphinxcode{\sphinxupquote{balance\sphinxhyphen{}futuro}}
&
\sphinxAtStartPar

&
\sphinxAtStartPar

\\
\sphinxhline
\sphinxAtStartPar
\sphinxcode{\sphinxupquote{“saludo}}
&
\sphinxAtStartPar

&
\sphinxAtStartPar

\\
\sphinxhline
\sphinxAtStartPar
\sphinxcode{\sphinxupquote{SALUDO}}
&
\sphinxAtStartPar

&
\sphinxAtStartPar

\\
\sphinxhline
\sphinxAtStartPar
\sphinxcode{\sphinxupquote{SUMA\_TOTAL}}
&
\sphinxAtStartPar

&
\sphinxAtStartPar

\\
\sphinxhline
\sphinxAtStartPar
\sphinxcode{\sphinxupquote{temperatura\_}}
&
\sphinxAtStartPar

&
\sphinxAtStartPar

\\
\sphinxhline
\sphinxAtStartPar
\sphinxcode{\sphinxupquote{32}}
&
\sphinxAtStartPar

&
\sphinxAtStartPar

\\
\sphinxhline
\sphinxAtStartPar
\sphinxcode{\sphinxupquote{\_a}}
&
\sphinxAtStartPar

&
\sphinxAtStartPar

\\
\sphinxhline
\sphinxAtStartPar
\sphinxcode{\sphinxupquote{4chan}}
&
\sphinxAtStartPar

&
\sphinxAtStartPar

\\
\sphinxhline
\sphinxAtStartPar
\sphinxcode{\sphinxupquote{\sphinxhyphen{}pool}}
&
\sphinxAtStartPar

&
\sphinxAtStartPar

\\
\sphinxbottomrule
\end{tabulary}
\sphinxtableafterendhook\par
\sphinxattableend\end{savenotes}


\bigskip\hrule\bigskip



\subsection{A03: Presentación (+0,25)}
\label{\detokenize{1_guias/10programacion/G1025:a03-presentacion-0-25}}
\sphinxAtStartPar
\DUrole{sd-sphinx-override}{\DUrole{sd-badge}{\DUrole{sd-bg-warning}{\DUrole{sd-bg-text-warning}{Actividad práctica}}}}
\begin{quote}

\sphinxAtStartPar
En la \sphinxstylestrong{consola interactiva de PYTHON}, realizar una presentación con sus datos.
\end{quote}
\begin{enumerate}
\sphinxsetlistlabels{\arabic}{enumi}{enumii}{}{.}%
\item {} 
\sphinxAtStartPar
Abra una \sphinxstylestrong{terminal de sistema} e ingrese a la \sphinxstylestrong{consola interactiva} de PYTHON.

\item {} 
\sphinxAtStartPar
CREAR diez variables con los siguientes nombres de variables y sus datos:
\begin{itemize}
\item {} 
\sphinxAtStartPar
nombre (\sphinxstylestrong{tipo STRING})

\item {} 
\sphinxAtStartPar
apellido (\sphinxstylestrong{tipo STRING})

\item {} 
\sphinxAtStartPar
edad (\sphinxstylestrong{tipo ENTERO})

\item {} 
\sphinxAtStartPar
día de nacimiento (\sphinxstylestrong{TIPO ENTERO})

\item {} 
\sphinxAtStartPar
mes de nacimiento (\sphinxstylestrong{TIPO ENTERO})

\item {} 
\sphinxAtStartPar
año de nacimiento (\sphinxstylestrong{TIPO ENTERO})

\item {} 
\sphinxAtStartPar
grado (\sphinxstylestrong{TIPO ENTERO})

\item {} 
\sphinxAtStartPar
curso (\sphinxstylestrong{TIPO STRING})

\item {} 
\sphinxAtStartPar
deporte (\sphinxstylestrong{TIPO STRING})

\item {} 
\sphinxAtStartPar
hobby (\sphinxstylestrong{TIPO STRING})

\end{itemize}

\begin{sphinxadmonition}{note}{EJEMPLO}

\begin{sphinxVerbatim}[commandchars=\\\{\},numbers=left,firstnumber=1,stepnumber=1]
\PYG{g+gp}{\PYGZgt{}\PYGZgt{}\PYGZgt{} }\PYG{n}{nombre} \PYG{o}{=} \PYG{l+s+s2}{\PYGZdq{}}\PYG{l+s+s2}{Camilo}\PYG{l+s+s2}{\PYGZdq{}}
\PYG{g+gp}{\PYGZgt{}\PYGZgt{}\PYGZgt{} }\PYG{n}{apellido} \PYG{o}{=} \PYG{l+s+s2}{\PYGZdq{}}\PYG{l+s+s2}{Fonseca}\PYG{l+s+s2}{\PYGZdq{}}
\PYG{g+gp}{\PYGZgt{}\PYGZgt{}\PYGZgt{} }\PYG{n}{edad} \PYG{o}{=} \PYG{l+m+mi}{35}
\PYG{g+gp}{\PYGZgt{}\PYGZgt{}\PYGZgt{} }\PYG{n}{dia\PYGZus{}nacimiento} \PYG{o}{=} \PYG{l+m+mi}{4}
\PYG{g+gp}{\PYGZgt{}\PYGZgt{}\PYGZgt{} }\PYG{n}{mes\PYGZus{}nacimiento} \PYG{o}{=} \PYG{l+m+mi}{12}
\PYG{g+gp}{\PYGZgt{}\PYGZgt{}\PYGZgt{} }\PYG{n}{año\PYGZus{}nacimiento} \PYG{o}{=} \PYG{l+m+mi}{1989}
\PYG{g+gp}{\PYGZgt{}\PYGZgt{}\PYGZgt{} }\PYG{n}{grado} \PYG{o}{=} \PYG{l+m+mi}{11}
\PYG{g+gp}{\PYGZgt{}\PYGZgt{}\PYGZgt{} }\PYG{n}{curso} \PYG{o}{=} \PYG{l+s+s2}{\PYGZdq{}}\PYG{l+s+s2}{A}\PYG{l+s+s2}{\PYGZdq{}}
\PYG{g+gp}{\PYGZgt{}\PYGZgt{}\PYGZgt{} }\PYG{n}{deporte} \PYG{o}{=} \PYG{l+s+s2}{\PYGZdq{}}\PYG{l+s+s2}{natacion}\PYG{l+s+s2}{\PYGZdq{}}
\PYG{g+gp}{\PYGZgt{}\PYGZgt{}\PYGZgt{} }\PYG{n}{hobby} \PYG{o}{=} \PYG{l+s+s2}{\PYGZdq{}}\PYG{l+s+s2}{videojuegos}\PYG{l+s+s2}{\PYGZdq{}}
\end{sphinxVerbatim}
\end{sphinxadmonition}

\item {} 
\sphinxAtStartPar
En consola muestre cada uno de los contenidos de la variables anteriores con la función \sphinxstylestrong{print()}.

\begin{sphinxadmonition}{note}{EJEMPLO}

\begin{sphinxVerbatim}[commandchars=\\\{\},numbers=left,firstnumber=1,stepnumber=1]
\PYG{g+gp}{\PYGZgt{}\PYGZgt{}\PYGZgt{} }\PYG{n+nb}{print}\PYG{p}{(}\PYG{n}{nombre}\PYG{p}{)}
\PYG{g+gp}{\PYGZgt{}\PYGZgt{}\PYGZgt{} }\PYG{n}{Camilo}
\PYG{g+gp}{\PYGZgt{}\PYGZgt{}\PYGZgt{} }\PYG{n+nb}{print}\PYG{p}{(}\PYG{n}{apellido}\PYG{p}{)}
\PYG{g+gp}{\PYGZgt{}\PYGZgt{}\PYGZgt{} }\PYG{n}{Fonseca}
\PYG{g+gp}{\PYGZgt{}\PYGZgt{}\PYGZgt{} }\PYG{n+nb}{print}\PYG{p}{(}\PYG{n}{edad}\PYG{p}{)}
\PYG{g+gp}{\PYGZgt{}\PYGZgt{}\PYGZgt{} }\PYG{l+m+mi}{35}
\PYG{g+gp}{\PYGZgt{}\PYGZgt{}\PYGZgt{} }\PYG{n+nb}{print}\PYG{p}{(}\PYG{n}{dia\PYGZus{}nacimiento}\PYG{p}{)}
\PYG{g+gp}{\PYGZgt{}\PYGZgt{}\PYGZgt{} }\PYG{l+m+mi}{4}
\PYG{g+gp}{\PYGZgt{}\PYGZgt{}\PYGZgt{} }\PYG{n+nb}{print}\PYG{p}{(}\PYG{n}{mes\PYGZus{}nacimiento}\PYG{p}{)}
\PYG{g+gp}{\PYGZgt{}\PYGZgt{}\PYGZgt{} }\PYG{l+m+mi}{12}
\PYG{g+gp}{\PYGZgt{}\PYGZgt{}\PYGZgt{} }\PYG{n+nb}{print}\PYG{p}{(}\PYG{n}{año\PYGZus{}nacimiento}\PYG{p}{)}
\PYG{g+gp}{\PYGZgt{}\PYGZgt{}\PYGZgt{} }\PYG{l+m+mi}{1989}
\PYG{g+gp}{\PYGZgt{}\PYGZgt{}\PYGZgt{} }\PYG{n+nb}{print}\PYG{p}{(}\PYG{n}{grado}\PYG{p}{)}
\PYG{g+gp}{\PYGZgt{}\PYGZgt{}\PYGZgt{} }\PYG{l+m+mi}{11}
\PYG{g+gp}{\PYGZgt{}\PYGZgt{}\PYGZgt{} }\PYG{n+nb}{print}\PYG{p}{(}\PYG{n}{curso}\PYG{p}{)}
\PYG{g+gp}{\PYGZgt{}\PYGZgt{}\PYGZgt{} }\PYG{n}{A}
\PYG{g+gp}{\PYGZgt{}\PYGZgt{}\PYGZgt{} }\PYG{n+nb}{print}\PYG{p}{(}\PYG{n}{deporte}\PYG{p}{)}
\PYG{g+gp}{\PYGZgt{}\PYGZgt{}\PYGZgt{} }\PYG{n}{natacion}
\PYG{g+gp}{\PYGZgt{}\PYGZgt{}\PYGZgt{} }\PYG{n+nb}{print}\PYG{p}{(}\PYG{n}{hobby}\PYG{p}{)}
\PYG{g+gp}{\PYGZgt{}\PYGZgt{}\PYGZgt{} }\PYG{n}{videojuegos}
\end{sphinxVerbatim}
\end{sphinxadmonition}

\item {} 
\sphinxAtStartPar
CREE una variable de nombre \sphinxstylestrong{mensaje00} y de tipo STRING para saludar.

\begin{sphinxadmonition}{note}{EJEMPLO}

\begin{sphinxVerbatim}[commandchars=\\\{\},numbers=left,firstnumber=1,stepnumber=1]
\PYG{g+gp}{\PYGZgt{}\PYGZgt{}\PYGZgt{} }\PYG{n}{mensaje00} \PYG{o}{=} \PYG{l+s+s2}{\PYGZdq{}}\PYG{l+s+s2}{Buenas tardes}\PYG{l+s+s2}{\PYGZdq{}}
\end{sphinxVerbatim}
\end{sphinxadmonition}

\item {} 
\sphinxAtStartPar
CREE una variable de nombre \sphinxstylestrong{mensaje01} y de tipo STRING para presentarse con su nombre y apellido.

\begin{sphinxadmonition}{note}{EJEMPLO}

\begin{sphinxVerbatim}[commandchars=\\\{\},numbers=left,firstnumber=1,stepnumber=1]
\PYG{g+gp}{\PYGZgt{}\PYGZgt{}\PYGZgt{} }\PYG{n}{mensaje01} \PYG{o}{=} \PYG{l+s+sa}{f}\PYG{l+s+s2}{\PYGZdq{}}\PYG{l+s+s2}{Mi nombre es: }\PYG{l+s+si}{\PYGZob{}}\PYG{n}{nombre}\PYG{l+s+si}{\PYGZcb{}}\PYG{l+s+s2}{ }\PYG{l+s+si}{\PYGZob{}}\PYG{n}{apellido}\PYG{l+s+si}{\PYGZcb{}}\PYG{l+s+s2}{\PYGZdq{}}
\end{sphinxVerbatim}
\end{sphinxadmonition}

\begin{sphinxadmonition}{note}{¿Cómo mostrar el valor de una variable en otra variable?}
\begin{itemize}
\item {} 
\sphinxAtStartPar
Es necesario agregar la letra \sphinxstylestrong{f}, antes de las primeras comillas dobles en el valor de la variable de tipo \sphinxstylestrong{STRING}.

\item {} 
\sphinxAtStartPar
Las variables a mostrar van dentro de llaves: \sphinxcode{\sphinxupquote{\{\}}}
Por ejemplo:

\end{itemize}

\begin{sphinxVerbatim}[commandchars=\\\{\},numbers=left,firstnumber=1,stepnumber=1]
\PYG{g+gp}{\PYGZgt{}\PYGZgt{}\PYGZgt{} }\PYG{n}{color} \PYG{o}{=} \PYG{l+s+s2}{\PYGZdq{}}\PYG{l+s+s2}{azul}\PYG{l+s+s2}{\PYGZdq{}}
\PYG{g+gp}{\PYGZgt{}\PYGZgt{}\PYGZgt{} }\PYG{n}{mensaje} \PYG{o}{=} \PYG{l+s+sa}{f}\PYG{l+s+s2}{\PYGZdq{}}\PYG{l+s+s2}{el cielo es de color }\PYG{l+s+si}{\PYGZob{}}\PYG{n}{color}\PYG{l+s+si}{\PYGZcb{}}\PYG{l+s+s2}{\PYGZdq{}}
\PYG{g+gp}{\PYGZgt{}\PYGZgt{}\PYGZgt{} }\PYG{n+nb}{print}\PYG{p}{(}\PYG{n}{mensaje}\PYG{p}{)}
\PYG{g+gp}{\PYGZgt{}\PYGZgt{}\PYGZgt{} }\PYG{n}{el} \PYG{n}{cielo} \PYG{n}{es} \PYG{n}{de} \PYG{n}{color} \PYG{n}{azul}
\end{sphinxVerbatim}
\end{sphinxadmonition}

\item {} 
\sphinxAtStartPar
CREE una variable de nombre \sphinxstylestrong{mensaje02} y de tipo STRING para mostrar su edad.

\begin{sphinxadmonition}{note}{EJEMPLO}

\begin{sphinxVerbatim}[commandchars=\\\{\},numbers=left,firstnumber=1,stepnumber=1]
\PYG{g+gp}{\PYGZgt{}\PYGZgt{}\PYGZgt{} }\PYG{n}{mensaje02} \PYG{o}{=} \PYG{l+s+sa}{f}\PYG{l+s+s2}{\PYGZdq{}}\PYG{l+s+s2}{Tengo }\PYG{l+s+si}{\PYGZob{}}\PYG{n}{edad}\PYG{l+s+si}{\PYGZcb{}}\PYG{l+s+s2}{ años de edad}\PYG{l+s+s2}{\PYGZdq{}}
\end{sphinxVerbatim}
\end{sphinxadmonition}

\item {} 
\sphinxAtStartPar
CREE una variable de nombre \sphinxstylestrong{mensaje03} y de tipo STRING para mostrar su fecha de nacimiento.

\begin{sphinxadmonition}{note}{EJEMPLO}

\begin{sphinxVerbatim}[commandchars=\\\{\},numbers=left,firstnumber=1,stepnumber=1]
\PYG{g+gp}{\PYGZgt{}\PYGZgt{}\PYGZgt{} }\PYG{n}{mensaje03} \PYG{o}{=} \PYG{l+s+sa}{f}\PYG{l+s+s2}{\PYGZdq{}}\PYG{l+s+s2}{Nací el día }\PYG{l+s+si}{\PYGZob{}}\PYG{n}{dia\PYGZus{}nacimiento}\PYG{l+s+si}{\PYGZcb{}}\PYG{l+s+s2}{ del mes }\PYG{l+s+si}{\PYGZob{}}\PYG{n}{mes\PYGZus{}nacimiento}\PYG{l+s+si}{\PYGZcb{}}\PYG{l+s+s2}{ del año }\PYG{l+s+si}{\PYGZob{}}\PYG{n}{año\PYGZus{}nacimiento}\PYG{l+s+si}{\PYGZcb{}}\PYG{l+s+s2}{\PYGZdq{}}
\end{sphinxVerbatim}
\end{sphinxadmonition}

\item {} 
\sphinxAtStartPar
CREE una variable de nombre \sphinxstylestrong{mensaje04} y de tipo STRING para mostrar el grado que cursa.

\begin{sphinxadmonition}{note}{EJEMPLO}

\begin{sphinxVerbatim}[commandchars=\\\{\},numbers=left,firstnumber=1,stepnumber=1]
\PYG{g+gp}{\PYGZgt{}\PYGZgt{}\PYGZgt{} }\PYG{n}{mensaje04} \PYG{o}{=} \PYG{l+s+sa}{f}\PYG{l+s+s2}{\PYGZdq{}}\PYG{l+s+s2}{Estoy en el grado }\PYG{l+s+si}{\PYGZob{}}\PYG{n}{grado}\PYG{l+s+si}{\PYGZcb{}}\PYG{l+s+si}{\PYGZob{}}\PYG{n}{curso}\PYG{l+s+si}{\PYGZcb{}}\PYG{l+s+s2}{\PYGZdq{}}
\end{sphinxVerbatim}
\end{sphinxadmonition}

\item {} 
\sphinxAtStartPar
CREE una variable de nombre \sphinxstylestrong{mensaje05} y de tipo STRING para mostrar el deporte que practica.

\begin{sphinxadmonition}{note}{EJEMPLO}

\begin{sphinxVerbatim}[commandchars=\\\{\},numbers=left,firstnumber=1,stepnumber=1]
\PYG{g+gp}{\PYGZgt{}\PYGZgt{}\PYGZgt{} }\PYG{n}{mensaje05} \PYG{o}{=} \PYG{l+s+sa}{f}\PYG{l+s+s2}{\PYGZdq{}}\PYG{l+s+s2}{Practico el deporte: }\PYG{l+s+si}{\PYGZob{}}\PYG{n}{deporte}\PYG{l+s+si}{\PYGZcb{}}\PYG{l+s+s2}{\PYGZdq{}}
\end{sphinxVerbatim}
\end{sphinxadmonition}

\item {} 
\sphinxAtStartPar
CREE una variable de nombre \sphinxstylestrong{mensaje06} y de tipo STRING para mostrar su hobby.

\begin{sphinxadmonition}{note}{EJEMPLO}

\begin{sphinxVerbatim}[commandchars=\\\{\},numbers=left,firstnumber=1,stepnumber=1]
\PYG{g+gp}{\PYGZgt{}\PYGZgt{}\PYGZgt{} }\PYG{n}{mensaje06} \PYG{o}{=} \PYG{l+s+sa}{f}\PYG{l+s+s2}{\PYGZdq{}}\PYG{l+s+s2}{En mi tiempo libre practico: }\PYG{l+s+si}{\PYGZob{}}\PYG{n}{hobby}\PYG{l+s+si}{\PYGZcb{}}\PYG{l+s+s2}{\PYGZdq{}}
\end{sphinxVerbatim}
\end{sphinxadmonition}

\item {} 
\sphinxAtStartPar
Finalmente muestre cada uno de los valores de las variables de los mensajes que acabó de crear. Para mostrarlos en pantalla use la función print().

\begin{sphinxadmonition}{note}{EJEMPLO}

\begin{sphinxVerbatim}[commandchars=\\\{\},numbers=left,firstnumber=1,stepnumber=1]
\PYG{g+gp}{\PYGZgt{}\PYGZgt{}\PYGZgt{} }\PYG{n+nb}{print}\PYG{p}{(}\PYG{n}{mensaje01}\PYG{p}{)}
\PYG{g+gp}{\PYGZgt{}\PYGZgt{}\PYGZgt{} }\PYG{n}{Mi} \PYG{n}{nombre} \PYG{n}{es}\PYG{p}{:} \PYG{n}{Camilo} \PYG{n}{Fonseca}
\PYG{g+gp}{\PYGZgt{}\PYGZgt{}\PYGZgt{} }\PYG{n+nb}{print}\PYG{p}{(}\PYG{n}{mensaje02}\PYG{p}{)}
\PYG{g+gp}{\PYGZgt{}\PYGZgt{}\PYGZgt{} }\PYG{n}{Tengo} \PYG{l+m+mi}{35} \PYG{n}{años} \PYG{n}{de} \PYG{n}{edad}
\PYG{g+gp}{\PYGZgt{}\PYGZgt{}\PYGZgt{} }\PYG{n+nb}{print}\PYG{p}{(}\PYG{n}{mensaje03}\PYG{p}{)}
\PYG{g+gp}{\PYGZgt{}\PYGZgt{}\PYGZgt{} }\PYG{n}{Naci} \PYG{n}{el} \PYG{n}{día} \PYG{l+m+mi}{4} \PYG{k}{del} \PYG{n}{mes} \PYG{l+m+mi}{12} \PYG{k}{del} \PYG{n}{año} \PYG{l+m+mi}{1989}
\PYG{g+gp}{\PYGZgt{}\PYGZgt{}\PYGZgt{} }\PYG{n+nb}{print}\PYG{p}{(}\PYG{n}{mensaje04}\PYG{p}{)}
\PYG{g+gp}{\PYGZgt{}\PYGZgt{}\PYGZgt{} }\PYG{n}{Estoy} \PYG{n}{en} \PYG{n}{el} \PYG{n}{grado} \PYG{l+m+mi}{11}\PYG{n}{A}
\PYG{g+gp}{\PYGZgt{}\PYGZgt{}\PYGZgt{} }\PYG{n+nb}{print}\PYG{p}{(}\PYG{n}{mensaje05}\PYG{p}{)}
\PYG{g+gp}{\PYGZgt{}\PYGZgt{}\PYGZgt{} }\PYG{n}{Practico} \PYG{n}{el} \PYG{n}{deporte}\PYG{p}{:} \PYG{n}{natación}
\PYG{g+gp}{\PYGZgt{}\PYGZgt{}\PYGZgt{} }\PYG{n+nb}{print}\PYG{p}{(}\PYG{n}{mensaje06}\PYG{p}{)}
\PYG{g+gp}{\PYGZgt{}\PYGZgt{}\PYGZgt{} }\PYG{n}{En} \PYG{n}{mi} \PYG{n}{tiempo} \PYG{n}{libre} \PYG{n}{practico}\PYG{p}{:} \PYG{n}{videojuegos}
\end{sphinxVerbatim}
\end{sphinxadmonition}

\end{enumerate}


\bigskip\hrule\bigskip



\subsection{A04: Orientaciones al extranjero (+0,25)}
\label{\detokenize{1_guias/10programacion/G1025:a04-orientaciones-al-extranjero-0-25}}
\sphinxAtStartPar
\DUrole{sd-sphinx-override}{\DUrole{sd-badge}{\DUrole{sd-bg-warning}{\DUrole{sd-bg-text-warning}{Actividad práctica}}}}
\begin{quote}

\sphinxAtStartPar
Con las \sphinxstylestrong{orientaciones} que previamente diseñó para al extranjero, usar la \sphinxstylestrong{consola interactiva de PYTHON} para mostrar dichas orientaciones.
\end{quote}
\begin{itemize}
\item {} 
\sphinxAtStartPar
Cree una serie de variables que contengan direcciones {[}\sphinxstylestrong{ejemplo: Cra 3A \# 7\sphinxhyphen{}01 (IE John F. Kennedy)}{]}. El nombre de cada variable debe ser \sphinxstylestrong{direccion01}, \sphinxstylestrong{direccion02} y así por cada dirección.

\item {} 
\sphinxAtStartPar
Cree una serie de variables que contengan cada una de las orientaciones. El nombre de cada variable debe ser \sphinxstylestrong{orientacion01}, \sphinxstylestrong{orientacion02} y así por cada orientación.

\item {} 
\sphinxAtStartPar
Las variables de \sphinxstylestrong{dirección} deben ser nombradas en una variable de \sphinxstylestrong{orientación}.

\item {} 
\sphinxAtStartPar
Al final muestre cada una de las orientaciones con la función \sphinxstylestrong{print()}

\end{itemize}

\begin{sphinxadmonition}{note}{EJEMPLO}

\begin{sphinxVerbatim}[commandchars=\\\{\},numbers=left,firstnumber=1,stepnumber=1]
\PYG{g+gp}{\PYGZgt{}\PYGZgt{}\PYGZgt{} }\PYG{n}{direccion01} \PYG{o}{=} \PYG{l+s+s2}{\PYGZdq{}}\PYG{l+s+s2}{Calle 28, Carrera 5 (Hospital José Cayetano Vasquez)}\PYG{l+s+s2}{\PYGZdq{}}
\PYG{g+gp}{\PYGZgt{}\PYGZgt{}\PYGZgt{} }\PYG{n}{direccion02} \PYG{o}{=} \PYG{l+s+s2}{\PYGZdq{}}\PYG{l+s+s2}{Calle 28}\PYG{l+s+s2}{\PYGZdq{}}
\PYG{g+gp}{\PYGZgt{}\PYGZgt{}\PYGZgt{} }\PYG{n}{orientacion01} \PYG{o}{=} \PYG{l+s+sa}{f}\PYG{l+s+s2}{\PYGZdq{}}\PYG{l+s+s2}{salga del hospital por la }\PYG{l+s+si}{\PYGZob{}}\PYG{n}{direccion01}\PYG{l+s+si}{\PYGZcb{}}\PYG{l+s+s2}{\PYGZdq{}}
\PYG{g+gp}{\PYGZgt{}\PYGZgt{}\PYGZgt{} }\PYG{n}{orientacion02} \PYG{o}{=} \PYG{l+s+sa}{f}\PYG{l+s+s2}{\PYGZdq{}}\PYG{l+s+s2}{camine 2 cuadras hacia el norte por }\PYG{l+s+si}{\PYGZob{}}\PYG{n}{direccion02}\PYG{l+s+si}{\PYGZcb{}}\PYG{l+s+s2}{\PYGZdq{}}
\PYG{g+gp}{\PYGZgt{}\PYGZgt{}\PYGZgt{} }\PYG{n+nb}{print}\PYG{p}{(}\PYG{n}{orientacion01}\PYG{p}{)}
\PYG{g+gp}{\PYGZgt{}\PYGZgt{}\PYGZgt{} }\PYG{n+nb}{print}\PYG{p}{(}\PYG{n}{orientacion02}\PYG{p}{)}
\end{sphinxVerbatim}
\end{sphinxadmonition}


\bigskip\hrule\bigskip



\subsection{A05: Operaciones en Python (+0,25)}
\label{\detokenize{1_guias/10programacion/G1025:a05-operaciones-en-python-0-25}}
\sphinxAtStartPar
\DUrole{sd-sphinx-override}{\DUrole{sd-badge}{\DUrole{sd-bg-warning}{\DUrole{sd-bg-text-warning}{Actividad práctica}}}}
\begin{quote}

\sphinxAtStartPar
Resuelva los siguientes ejercicios en su cuaderno (\sphinxstylestrong{EVIDENCIE EL PROCEDIMIENTO}) (\sphinxstylestrong{NO USAR CALCULADORA}).
\end{quote}
\begin{enumerate}
\sphinxsetlistlabels{\arabic}{enumi}{enumii}{}{.}%
\item {} 
\sphinxAtStartPar
Opere e identifique los valores finales para las variables a, b y c:

\end{enumerate}
\begin{quote}
\begin{equation*}
\begin{split}a &= 3
\\ b &= 2
\\ c &= 1
\\
\\ a &= a + 1
\\ b &= b + 2
\\ c &= c + 3\end{split}
\end{equation*}\end{quote}
\begin{enumerate}
\sphinxsetlistlabels{\arabic}{enumi}{enumii}{}{.}%
\setcounter{enumi}{1}
\item {} 
\sphinxAtStartPar
Opere e identifique los valores finales para las variables a, b y c:

\end{enumerate}
\begin{quote}
\begin{equation*}
\begin{split}a &= 10
\\ b &= 20
\\ c &= 5
\\
\\ a &= a + 3
\\ b &= b + 4 - a
\\ c &= a + b + c
\\
\\ a &= a + c
\\ b &= 4 - b
\\ c &= c + 3 - b + 2\end{split}
\end{equation*}\end{quote}
\begin{enumerate}
\sphinxsetlistlabels{\arabic}{enumi}{enumii}{}{.}%
\setcounter{enumi}{2}
\item {} 
\sphinxAtStartPar
Opere e identifique los valores finales para las variables a, b, c y d:

\end{enumerate}
\begin{quote}
\begin{equation*}
\begin{split}a &= 5
\\ b &= 18
\\ c &= 15
\\ d &= 25
\\
\\ a &= a + 10
\\ b &= b + 5 - c
\\ c &= c + 4 + b
\\ d &= d + b + a
\\
\\ a &= a + 1
\\ b &= b + c
\\ c &= b + c
\\ d &= b + b\end{split}
\end{equation*}\end{quote}
\begin{enumerate}
\sphinxsetlistlabels{\arabic}{enumi}{enumii}{}{.}%
\setcounter{enumi}{3}
\item {} 
\sphinxAtStartPar
Opere e identifique los valores finales para las variables a y b:

\end{enumerate}
\begin{quote}
\begin{equation*}
\begin{split}a &= 9
\\ b &= 6
\\
\\ a &= a + 4
\\ b &= b + 2
\\
\\ a &= a + 10
\\ b &= b - 25
\\
\\ a &= a - 20
\\ b &= b + 5
\\
\\ a &= a + 4
\\ b &= b + 2
\\
\\ a &= a + 10
\\ b &= b - 10\end{split}
\end{equation*}\end{quote}
\begin{enumerate}
\sphinxsetlistlabels{\arabic}{enumi}{enumii}{}{.}%
\setcounter{enumi}{4}
\item {} 
\sphinxAtStartPar
Opere e identifique los valores finales para las variables a, b, c y d:

\end{enumerate}
\begin{quote}
\begin{equation*}
\begin{split}a &= 18
\\ b &= 18
\\ c &= 18
\\ d &= 18
\\
\\ a &= a + b
\\ b &= a - b
\\ c &= a + b
\\ d &= a - b
\\
\\ a &= a - b
\\ b &= a + b
\\ c &= a - b
\\ d &= a + b\end{split}
\end{equation*}\end{quote}
\begin{enumerate}
\sphinxsetlistlabels{\arabic}{enumi}{enumii}{}{.}%
\setcounter{enumi}{5}
\item {} 
\sphinxAtStartPar
Opere e identifique los valores finales para las variables a y b:

\end{enumerate}
\begin{quote}
\begin{equation*}
\begin{split}a &= 10
\\ b &= 5
\\
\\ a &= a - 5
\\ b &= b + 6
\\
\\ a &= a + 18
\\ b &= b - 23
\\
\\ a &= a - 21
\\ b &= b - 5
\\
\\ a &= a - 4
\\ b &= b - 2
\\
\\ a &= a + 10
\\ b &= b + 10\end{split}
\end{equation*}\end{quote}
\begin{enumerate}
\sphinxsetlistlabels{\arabic}{enumi}{enumii}{}{.}%
\setcounter{enumi}{6}
\item {} 
\sphinxAtStartPar
Opere e identifique los valores finales para las variables a, b, c y d:

\end{enumerate}
\begin{quote}
\begin{equation*}
\begin{split}a &= 8
\\ b &= 7
\\ c &= 5
\\ d &= 8
\\
\\ a &= a + b - c + d
\\ b &= a + b - c + d
\\ c &= a + b - c + d
\\ d &= a + b - c + d
\\
\\ a &= a - b + c - d
\\ b &= a - b + c - d
\\ c &= a - b + c - d
\\ d &= a - b + c - d\end{split}
\end{equation*}\end{quote}
\begin{enumerate}
\sphinxsetlistlabels{\arabic}{enumi}{enumii}{}{.}%
\setcounter{enumi}{7}
\item {} 
\sphinxAtStartPar
Opere e identifique los valores finales para las variables a, b y c:

\end{enumerate}
\begin{quote}
\begin{equation*}
\begin{split}a &= 10
\\ b &= 5
\\ c &= 10
\\
\\ a &= a + b - 5
\\ b &= a + b - 5
\\ c &= a + b - 5
\\
\\ a &= a + 5 * b / 2
\\ b &= a + 5 * b / 2
\\ c &= a + 5 * b / 2\end{split}
\end{equation*}\end{quote}
\begin{enumerate}
\sphinxsetlistlabels{\arabic}{enumi}{enumii}{}{.}%
\setcounter{enumi}{8}
\item {} 
\sphinxAtStartPar
Opere e identifique los valores finales para las variables a, b y c:

\end{enumerate}
\begin{quote}
\begin{equation*}
\begin{split}a &= 1
\\ b &= 2
\\ c &= 3
\\
\\ a &= a + 2
\\ b &= a + 2 + b
\\ c &= a + 2 + c
\\
\\ a &= a / 2
\\ b &= b / 2
\\ c &= c / 2\end{split}
\end{equation*}\end{quote}
\begin{enumerate}
\sphinxsetlistlabels{\arabic}{enumi}{enumii}{}{.}%
\setcounter{enumi}{9}
\item {} 
\sphinxAtStartPar
Transcriba cada uno de los \sphinxstylestrong{ejercicios anteriores} en la \sphinxstylestrong{consola interactiva de python} y revise que sus \sphinxstylestrong{resultados en el cuaderno} \sphinxstyleemphasis{coincidan} con los \sphinxstylestrong{resultados en pantalla}.
\begin{itemize}
\item {} 
\sphinxAtStartPar
\sphinxstylestrong{EN CASO DE NO COINCIDIR}, verificar si el problema está en sus cálculos o en las variables establecidas en la \sphinxstylestrong{consola interactiva}.

\item {} 
\sphinxAtStartPar
Para mostrar los \sphinxstylestrong{resultados finales} de las variables en \sphinxstylestrong{pantalla}, usar la función \sphinxcode{\sphinxupquote{print()}}.

\end{itemize}

\end{enumerate}


\bigskip\hrule\bigskip



\subsection{A06: Linealización de Operaciones (+0,25)}
\label{\detokenize{1_guias/10programacion/G1025:a06-linealizacion-de-operaciones-0-25}}
\sphinxAtStartPar
\DUrole{sd-sphinx-override}{\DUrole{sd-badge}{\DUrole{sd-bg-warning}{\DUrole{sd-bg-text-warning}{Actividad práctica}}}}
\begin{enumerate}
\sphinxsetlistlabels{\arabic}{enumi}{enumii}{}{.}%
\item {} 
\sphinxAtStartPar
En su cuaderno responda:
\begin{itemize}
\item {} 
\sphinxAtStartPar
¿Cuales son las reglas de precedencia?

\end{itemize}

\item {} 
\sphinxAtStartPar
En su cuaderno dibuje una tabla y en ella escriba el orden de la precedencia y su descripción para los operadores: \sphinxcode{\sphinxupquote{()}}, \sphinxcode{\sphinxupquote{*}}, \sphinxcode{\sphinxupquote{**}}, \sphinxcode{\sphinxupquote{/}}, \sphinxcode{\sphinxupquote{//}}, \sphinxcode{\sphinxupquote{\%}}, \sphinxcode{\sphinxupquote{+}}, \sphinxcode{\sphinxupquote{\sphinxhyphen{}}}.

\item {} 
\sphinxAtStartPar
Linealizar las siguientes operaciones en su cuaderno (\sphinxstylestrong{EVIDENCIAR PROCEDIMIENTO}):

\sphinxAtStartPar
\sphinxstylestrong{3.1}
\begin{quote}
\begin{equation*}
\begin{split} X = { a + \frac{b}{c} \over {a \over b} + c } \end{split}
\end{equation*}\end{quote}


\bigskip\hrule\bigskip


\sphinxAtStartPar
\sphinxstylestrong{3.2}
\begin{quote}
\begin{equation*}
\begin{split} X = { a + b + { a \over b } \over c }  \end{split}
\end{equation*}\end{quote}


\bigskip\hrule\bigskip


\sphinxAtStartPar
\sphinxstylestrong{3.3}
\begin{quote}
\begin{equation*}
\begin{split} X = {\dfrac{a}{a + b} \over \dfrac{a}{a - b}}  \end{split}
\end{equation*}\end{quote}


\bigskip\hrule\bigskip


\sphinxAtStartPar
\sphinxstylestrong{3.4}
\begin{quote}
\begin{equation*}
\begin{split} X = {  a + { b \over a + b + { b \over c }  } \over a + { b \over c + a } } \end{split}
\end{equation*}\end{quote}


\bigskip\hrule\bigskip


\sphinxAtStartPar
\sphinxstylestrong{3.5}
\begin{quote}
\begin{equation*}
\begin{split} X = { a + b + c \over a + { b \over c } } \end{split}
\end{equation*}\end{quote}


\bigskip\hrule\bigskip


\sphinxAtStartPar
\sphinxstylestrong{3.6}
\begin{quote}
\begin{equation*}
\begin{split} X = { a + b + { c \over d * a } \over a + b * {c \over d} } \end{split}
\end{equation*}\end{quote}


\bigskip\hrule\bigskip


\sphinxAtStartPar
\sphinxstylestrong{3.7}
\begin{quote}
\begin{equation*}
\begin{split} X = { a + { b \over c } + d \over a } \end{split}
\end{equation*}\end{quote}


\bigskip\hrule\bigskip


\sphinxAtStartPar
\sphinxstylestrong{3.8}
\begin{quote}
\begin{equation*}
\begin{split} X = { { a \over b } + { b \over c } \over { a \over b } - { b \over c } } \end{split}
\end{equation*}\end{quote}


\bigskip\hrule\bigskip


\sphinxAtStartPar
\sphinxstylestrong{3.9}
\begin{quote}
\begin{equation*}
\begin{split} X = a + { a + { a + b \over c + d } \over a + { a \over b }  } \end{split}
\end{equation*}\end{quote}


\bigskip\hrule\bigskip


\sphinxAtStartPar
\sphinxstylestrong{3.10}
\begin{quote}
\begin{equation*}
\begin{split} X = a + b + { c \over d } + { { a \over b - c } \over { a \over b + c } } \end{split}
\end{equation*}\end{quote}

\end{enumerate}


\bigskip\hrule\bigskip



\subsection{A07: Modelo \% de notas (+3,0)}
\label{\detokenize{1_guias/10programacion/G1025:a07-modelo-de-notas-3-0}}
\sphinxAtStartPar
\DUrole{sd-sphinx-override}{\DUrole{sd-badge}{\DUrole{sd-bg-danger}{\DUrole{sd-bg-text-danger}{Transferencia}}}}


\begin{savenotes}\sphinxattablestart
\sphinxthistablewithglobalstyle
\centering
\begin{tabulary}{\linewidth}[t]{TT}
\sphinxtoprule
\sphinxtableatstartofbodyhook
\sphinxAtStartPar
\sphinxstylestrong{TIEMPO}
&
\sphinxAtStartPar
99 minutos
\\
\sphinxhline
\sphinxAtStartPar
\sphinxstylestrong{OBJETIVO}
&
\sphinxAtStartPar
Aplicar conocimientos adquiridos.
\\
\sphinxhline
\sphinxAtStartPar
\sphinxstylestrong{ACTIVIDAD}
&
\sphinxAtStartPar
Modelar con \sphinxstylestrong{variables y operaciones} la \sphinxstylestrong{calificación final} de cada asignatura del \sphinxstylestrong{periodo actual}. Posteriormente implementar y ejecutar el modelo en la \sphinxstylestrong{consola interactiva de PYTHON}.
\\
\sphinxbottomrule
\end{tabulary}
\sphinxtableafterendhook\par
\sphinxattableend\end{savenotes}

\begin{sphinxadmonition}{note}{MODELO DE NOTAS}
\begin{enumerate}
\sphinxsetlistlabels{\arabic}{enumi}{enumii}{}{.}%
\item {} 
\sphinxAtStartPar
En su \sphinxstylestrong{cuaderno} realice el \sphinxstylestrong{modelo de notas} del periodo actual para cada asignatura (si el estudiante ve 10 materias entonces serían 10 modelos). Cada modelo se compone de:
\begin{itemize}
\item {} 
\sphinxAtStartPar
Una \sphinxstylestrong{variable por cada nota} (si en una materia se sacan 5 notas, entonces serán 5 variables).

\item {} 
\sphinxAtStartPar
El nombre de las variables debe ser acorde al \sphinxstylestrong{contexto de la nota} y debe ser de \sphinxstylestrong{tipo decimal} (por ejemplo \sphinxcode{\sphinxupquote{notaQuizMatematicas = 4.5}}).

\item {} 
\sphinxAtStartPar
Una \sphinxstylestrong{ecuación con operaciones} donde se incluyan todas las \sphinxstylestrong{variables de las notas}. El valor de \( X \) será el valor de la nota final. Por ejemplo:
\begin{equation*}
\begin{split} X = \left( \left( \frac{notaGuia01+notaGuia02}{2} \right) * (0.50) \right) +  (notaINS*(0.40)) + (notaPerfil * (0.10)) \end{split}
\end{equation*}
\item {} 
\sphinxAtStartPar
La ecuación debe ser linealizada.

\item {} \begin{quote}

\sphinxAtStartPar
TENER EN CUENTA QUE CADA ASIGNATURA TIENE SUS PROPIOS PORCENTAJES Y NOTAS. Por tanto \sphinxstylestrong{cada modelo es diferente}.
\end{quote}

\end{itemize}

\item {} 
\sphinxAtStartPar
En la \sphinxstylestrong{consola interactiva de python}, aplique y ejecute el modelo diseñado en su cuaderno. Tenga en cuenta lo siguiente:
\begin{itemize}
\item {} 
\sphinxAtStartPar
Una consola interactiva por cada modelo.

\item {} 
\sphinxAtStartPar
Se deben crear variables para identificar el periodo actual, la materia y el nombre del estudiante.

\item {} 
\sphinxAtStartPar
Al final, el programa en PYTHON debe permitir mostrar los valores de las variables con la función \sphinxcode{\sphinxupquote{print()}}.

\end{itemize}

\end{enumerate}
\end{sphinxadmonition}


\bigskip\hrule\bigskip


\sphinxstepscope


\section{Guía 1029: Plan de refuerzo}
\label{\detokenize{1_guias/10programacion/G1029:guia-1029-plan-de-refuerzo}}\label{\detokenize{1_guias/10programacion/G1029::doc}}\begin{quote}\begin{description}
\sphinxlineitem{ASIGNATURA}
\sphinxAtStartPar
Programación

\sphinxlineitem{GRADO}
\sphinxAtStartPar
Décimo

\end{description}\end{quote}


\bigskip\hrule\bigskip



\subsection{PRIMER PERIODO}
\label{\detokenize{1_guias/10programacion/G1029:primer-periodo}}
\begin{sphinxadmonition}{note}{EXAMEN ESCRITO: Plan de refuerzo (10 minutos)}

\sphinxAtStartPar
Estudiar los conceptos a continuación. Se realiza examen de 5 preguntas. Si el estudiante no asiste en la fecha programada, este será evaluado la próxima vez que asista a clase.
\end{sphinxadmonition}


\subsubsection{Instrucción}
\label{\detokenize{1_guias/10programacion/G1029:instruccion}}
\sphinxAtStartPar
Una instrucción es una orden o indicación que se da para realizar una acción específica.


\subsubsection{Hardware}
\label{\detokenize{1_guias/10programacion/G1029:hardware}}
\sphinxAtStartPar
El hardware es el conjunto de componentes de un sistema informático o computacional que se puede ver o tocar. Por ejemplo:
\begin{itemize}
\item {} 
\sphinxAtStartPar
\sphinxstylestrong{Dispositivos de entrada:} Teclados, mouse, micrófonos.

\item {} 
\sphinxAtStartPar
\sphinxstylestrong{Dispositivos de salida:} Monitores, impresoras, bafles.

\item {} 
\sphinxAtStartPar
\sphinxstylestrong{Procesamiento:} CPU.

\item {} 
\sphinxAtStartPar
\sphinxstylestrong{Memoria:} memoria RAM, memoria CACHE.

\item {} 
\sphinxAtStartPar
\sphinxstylestrong{Almacenamiento:} Discos duros, memorias USB.

\end{itemize}


\subsubsection{Software}
\label{\detokenize{1_guias/10programacion/G1029:software}}
\sphinxAtStartPar
El software es el conjunto de programas, bases de datos, diferentes tipos de archivos y documentación que le dicen al hardware qué hacer. Es la parte lógica e intangible de un sistema informático. Ejemplos de software:
\begin{itemize}
\item {} 
\sphinxAtStartPar
Sistemas operativos

\item {} 
\sphinxAtStartPar
Aplicaciones

\item {} 
\sphinxAtStartPar
Drivers

\end{itemize}


\subsubsection{Información y datos}
\label{\detokenize{1_guias/10programacion/G1029:informacion-y-datos}}
\sphinxAtStartPar
Los datos son valores en bruto, hechos o símbolos que por sí solos no tienen significado. Por ejemplo:
\begin{itemize}
\item {} 
\sphinxAtStartPar
Números sueltos.

\item {} 
\sphinxAtStartPar
Caracteres sin contexto.

\item {} 
\sphinxAtStartPar
Valores sin procesar.

\end{itemize}

\sphinxAtStartPar
La información es el resultado de procesar, organizar e interpretar los datos de manera que tengan sentido y sean útiles. La información:
\begin{itemize}
\item {} 
\sphinxAtStartPar
Tiene significado y contexto.

\item {} 
\sphinxAtStartPar
Es comprensible para el usuario.

\item {} 
\sphinxAtStartPar
Ayuda en la toma de decisiones.

\end{itemize}


\subsubsection{Software Vs Programa}
\label{\detokenize{1_guias/10programacion/G1029:software-vs-programa}}

\paragraph{Software}
\label{\detokenize{1_guias/10programacion/G1029:id1}}
\sphinxAtStartPar
Un \sphinxstylestrong{software} es el conjunto de \sphinxstylestrong{programas}, bases de datos, archivos y documentación que permiten realizar diferentes tareas en un sistema computacional.


\paragraph{Programa}
\label{\detokenize{1_guias/10programacion/G1029:programa}}
\sphinxAtStartPar
Un \sphinxstylestrong{programa} es un \sphinxstylestrong{conjunto de instrucciones} que le dicen a una computadora COMO HACER una tarea específica. Las instrucciones están escritas en código y son diseñadas por programadores con conocimiento en un lenguaje de programación.


\bigskip\hrule\bigskip



\subsection{SEGUNDO PERIODO}
\label{\detokenize{1_guias/10programacion/G1029:segundo-periodo}}
\begin{sphinxadmonition}{note}{EXAMEN ESCRITO: Plan de refuerzo (10 minutos)}

\sphinxAtStartPar
Estudiar los conceptos a continuación. Se realiza examen de 5 preguntas. Si el estudiante no asiste en la fecha programada, este será evaluado la próxima vez que asista a clase.
\end{sphinxadmonition}

\sphinxstepscope


\chapter{Bibliografía}
\label{\detokenize{2_referencias/bibliografia:bibliografia}}\label{\detokenize{2_referencias/bibliografia::doc}}\begin{itemize}
\item {} 
\sphinxAtStartPar
Deitel, P. J., \& Deitel, H. (2014). C++: cómo programar (A. V. Romero Elizondo, Trans.). Pearson Educación.

\item {} 
\sphinxAtStartPar
Rohaut, S. Linux: preparación a la certificación LPIC 1 (exámenes LPI 101 y LPI 102), 3ª Edición. Barcelona: ENI ediciones, 2015.

\item {} 
\sphinxAtStartPar
Sweigart, A. (2019). Automate the boring stuff with python, 2nd edition: Practical programming for total beginners. No Starch Press.

\item {} 
\sphinxAtStartPar
Trejos Buriticá, O. I. (2017). Lógica de programación. Ediciones de la U.

\end{itemize}

\sphinxstepscope


\chapter{Webgrafía}
\label{\detokenize{2_referencias/webgrafia:webgrafia}}\label{\detokenize{2_referencias/webgrafia::doc}}\begin{itemize}
\item {} 
\sphinxAtStartPar
Antonio Fernández, B. (2025, May 9). Guía de iniciación 24.8 \sphinxhyphen{} Prefacio. Libreoffice.org. https://books.libreoffice.org/es/GS248/GS24800\sphinxhyphen{}Prefacio.html

\item {} 
\sphinxAtStartPar
Andrés, R. (2019, September 1). Raspberry Pi: la historia de un mini\sphinxhyphen{}PC que se ha convertido en un éxito de masas. Computer Hoy. https://computerhoy.20minutos.es/reportajes/tecnologia/historia\sphinxhyphen{}raspberry\sphinxhyphen{}pi\sphinxhyphen{}482477

\item {} 
\sphinxAtStartPar
Andreu, A. (2024, September 17). Qué es Mozilla Firefox, diferencias con otros navegadores y extensiones imprescindibles. Computer Hoy. https://computerhoy.20minutos.es/tecnologia/mozilla\sphinxhyphen{}firefox\sphinxhyphen{}diferencias\sphinxhyphen{}otros\sphinxhyphen{}navegadores\sphinxhyphen{}extensiones\sphinxhyphen{}imprescindibles\sphinxhyphen{}1404587

\item {} 
\sphinxAtStartPar
ARPANET. (2005, Feb 10). Sistemas.com. Retrieved April 25, 2026, from https://sistemas.com/arpanet.php

\item {} 
\sphinxAtStartPar
Audacity: una forma diferente de estar en la onda. (2023, November 20). INTEF. https://intef.es/observatorio\_tecno/audacity\sphinxhyphen{}una\sphinxhyphen{}forma\sphinxhyphen{}diferente\sphinxhyphen{}de\sphinxhyphen{}estar\sphinxhyphen{}en\sphinxhyphen{}la\sphinxhyphen{}onda/

\item {} 
\sphinxAtStartPar
Catucci, A. (2020, September 9). La historia de Apple: repaso por este caso de éxito tecnológico. Marketing Insider Review. https://marketinginsiderreview.com/historia\sphinxhyphen{}apple\sphinxhyphen{}exito\sphinxhyphen{}tecnologico/

\item {} 
\sphinxAtStartPar
Corrales, R. (2024, August 19). Microsoft, el gigante de la informática: historia, Windows, todos sus servicios y la IA de Copilot. Computer Hoy. https://computerhoy.20minutos.es/industria/microsoft\sphinxhyphen{}gigante\sphinxhyphen{}informatica\sphinxhyphen{}historia\sphinxhyphen{}windows\sphinxhyphen{}todos\sphinxhyphen{}servicios\sphinxhyphen{}ia\sphinxhyphen{}copilot\sphinxhyphen{}1399819

\item {} 
\sphinxAtStartPar
Delgado, J. M. (2024, September 15). Razones por las que GIMP es la mejor alternativa a Photoshop para editar tus fotos. Computer Hoy. https://computerhoy.20minutos.es/apps/razones\sphinxhyphen{}gimp\sphinxhyphen{}mejor\sphinxhyphen{}alternativa\sphinxhyphen{}photoshop\sphinxhyphen{}editar\sphinxhyphen{}fotos\sphinxhyphen{}1400576

\item {} 
\sphinxAtStartPar
Delgado, J. M. (2024, December 15). Funciones ocultas de Microsoft Office que deberías usar ahora mismo para mejorar tu productividad. Computer Hoy. https://computerhoy.20minutos.es/apps/funciones\sphinxhyphen{}ocultas\sphinxhyphen{}microsoft\sphinxhyphen{}office\sphinxhyphen{}deberias\sphinxhyphen{}usar\sphinxhyphen{}ahora\sphinxhyphen{}mismo\sphinxhyphen{}mejorar\sphinxhyphen{}productividad\sphinxhyphen{}1428291

\item {} 
\sphinxAtStartPar
Delgado, J. M. (2025, June 14). VLC sigue muy vivo: cosas sorprendentes que puedes hacer con el mejor reproductor de vídeo para PC. Computer Hoy. https://computerhoy.20minutos.es/tecnologia/vlc\sphinxhyphen{}sigue\sphinxhyphen{}muy\sphinxhyphen{}vivo\sphinxhyphen{}cosas\sphinxhyphen{}sorprendentes\sphinxhyphen{}puedes\sphinxhyphen{}hacer\sphinxhyphen{}mejor\sphinxhyphen{}reproductor\sphinxhyphen{}video\sphinxhyphen{}pc\sphinxhyphen{}1463759

\item {} 
\sphinxAtStartPar
de Luis, E. R. (2026, March 15). En 1990 un estudiante tuvo un problema mostrando imágenes en su Mac. La solución fue crear uno de los softwares más exitosos de la historia. Xataka.com; Xataka. https://www.xataka.com/aplicaciones/1987\sphinxhyphen{}tuvo\sphinxhyphen{}problema\sphinxhyphen{}mostrando\sphinxhyphen{}imagenes\sphinxhyphen{}su\sphinxhyphen{}mac\sphinxhyphen{}asi\sphinxhyphen{}que\sphinxhyphen{}creo\sphinxhyphen{}app\sphinxhyphen{}hoy\sphinxhyphen{}editor\sphinxhyphen{}imagenes\sphinxhyphen{}usado\sphinxhyphen{}historia

\item {} 
\sphinxAtStartPar
Eadicicco, L. (2025, May 11). La tecnologí­a que definió el internet moderno está cambiando y Silicon Valley finalmente lo admite. CNN en Español. https://cnnespanol.cnn.com/2025/05/11/ciencia/tecnologia\sphinxhyphen{}internet\sphinxhyphen{}moderno\sphinxhyphen{}silicon\sphinxhyphen{}valley\sphinxhyphen{}trax

\item {} 
\sphinxAtStartPar
Hill, S. (2024, March 24). Qué es Google One y cómo saber si lo necesitas. WIRED. https://es.wired.com/articulos/que\sphinxhyphen{}es\sphinxhyphen{}google\sphinxhyphen{}one\sphinxhyphen{}y\sphinxhyphen{}como\sphinxhyphen{}saber\sphinxhyphen{}si\sphinxhyphen{}lo\sphinxhyphen{}necesitas

\item {} 
\sphinxAtStartPar
Historia de Microsoft. (2012, November 10). Blog de tecnología; Sistemas Ibertrónica. https://ibertronica.es/blog/productos/microsoft\sphinxhyphen{}pasado\sphinxhyphen{}presente\sphinxhyphen{}y\sphinxhyphen{}futuro/

\item {} 
\sphinxAtStartPar
Huertos, A. A. (2019, May 28). La historia de los lenguajes de programación. Computer Hoy. https://computerhoy.20minutos.es/reportajes/tecnologia/historia\sphinxhyphen{}lenguajes\sphinxhyphen{}programacion\sphinxhyphen{}428041

\item {} 
\sphinxAtStartPar
Inkscape, tu ilustrador de cabecera. (2021, July 19). INTEF. https://intef.es/observatorio\_tecno/inkscape\sphinxhyphen{}tu\sphinxhyphen{}ilustrador\sphinxhyphen{}de\sphinxhyphen{}cabecera/

\item {} 
\sphinxAtStartPar
Lautrec, T. (n.d.). ¿Qué es AutoCAD y para qué sirve? Toulouse Lautrec. Retrieved May 7, 2026, from https://www.toulouselautrec.edu.pe/blogs/que\sphinxhyphen{}es\sphinxhyphen{}autocad

\item {} 
\sphinxAtStartPar
Maturana, J. (2023, June 2). 25 aniversario de Computer Hoy: La historia de Windows, cómo ha sido la evolución del servicio más popular de Microsoft. Computer Hoy. https://computerhoy.20minutos.es/windows/historia\sphinxhyphen{}windows\sphinxhyphen{}como\sphinxhyphen{}ha\sphinxhyphen{}sido\sphinxhyphen{}evolucion\sphinxhyphen{}servicio\sphinxhyphen{}popular\sphinxhyphen{}microsoft\sphinxhyphen{}1237508

\item {} 
\sphinxAtStartPar
Nava, J. (2026, February 19). ¿Qué es blender? sus usos y beneficios. IMSED Business \& Tech School. https://imsed.com/postgrados/tecnologia/que\sphinxhyphen{}es\sphinxhyphen{}blender\sphinxhyphen{}usos\sphinxhyphen{}beneficios/

\item {} 
\sphinxAtStartPar
(N.d.). Libreoffice.org. Retrieved May 17, 2026, from https://es.libreoffice.org/who\sphinxhyphen{}uses\sphinxhyphen{}libreoffice/

\item {} 
\sphinxAtStartPar
Onieva, D. (2020, February 17). LibreOffice, la suite gratuita que compite cara a cara con Office. SoftZone. https://www.softzone.es/programas/oficina/libreoffice/

\item {} 
\sphinxAtStartPar
Pastor, J. (2015, January 27). Altair 8800: todo comenzó hace ya 40 años. Xataka.com; Xataka. https://www.xataka.com/historia\sphinxhyphen{}tecnologica/altair\sphinxhyphen{}8800\sphinxhyphen{}todo\sphinxhyphen{}comenzo\sphinxhyphen{}hace\sphinxhyphen{}ya\sphinxhyphen{}40\sphinxhyphen{}anos

\item {} 
\sphinxAtStartPar
Pastor, J. (2018, April 28). Cuando Mosaic dominaba el mundo (de los navegadores). Xataka.com; Xataka. https://www.xataka.com/historia\sphinxhyphen{}tecnologica/cuando\sphinxhyphen{}mosaic\sphinxhyphen{}dominaba\sphinxhyphen{}el\sphinxhyphen{}mundo\sphinxhyphen{}de\sphinxhyphen{}los\sphinxhyphen{}navegadores

\item {} 
\sphinxAtStartPar
Rojas, J. (2022, May 17). El primer sitio web de la historia. Un hito que cambió vidas. Telefónica. https://www.telefonica.com/es/sala\sphinxhyphen{}comunicacion/blog/el\sphinxhyphen{}primer\sphinxhyphen{}sitio\sphinxhyphen{}web\sphinxhyphen{}de\sphinxhyphen{}la\sphinxhyphen{}historia\sphinxhyphen{}un\sphinxhyphen{}hito\sphinxhyphen{}que\sphinxhyphen{}cambio\sphinxhyphen{}vidas/

\item {} 
\sphinxAtStartPar
Superior, E. F. P. (2023, August 31). ¿Qué es la ofimática y para qué sirve?: herramientas. Esic.edu; ESIC. https://www.esic.edu/business/que\sphinxhyphen{}es\sphinxhyphen{}la\sphinxhyphen{}ofimatica\sphinxhyphen{}para\sphinxhyphen{}que\sphinxhyphen{}sirve\sphinxhyphen{}c

\item {} 
\sphinxAtStartPar
Textos: Jorge Eliécer Lozano / Ilustraciones: Gissell Andreína García Rodríguez. (n.d.). El software pirata representa un riesgo inminente. Edu.co. Retrieved May 7, 2026, from https://notired.univalle.edu.co/el\sphinxhyphen{}software\sphinxhyphen{}pirata\sphinxhyphen{}representa\sphinxhyphen{}un\sphinxhyphen{}riesgo\sphinxhyphen{}inminente/

\item {} 
\sphinxAtStartPar
The Conversation. (2024, January 22). La Mac cumple 40 años: debutó con este antológico aviso dirigido por Ridley Scott para mostrar cómo había otra forma de usar una computadora. LA NACION. https://www.lanacion.com.ar/tecnologia/la\sphinxhyphen{}mac\sphinxhyphen{}cumple\sphinxhyphen{}40\sphinxhyphen{}anos\sphinxhyphen{}debuto\sphinxhyphen{}con\sphinxhyphen{}este\sphinxhyphen{}antologico\sphinxhyphen{}aviso\sphinxhyphen{}dirigido\sphinxhyphen{}por\sphinxhyphen{}ridley\sphinxhyphen{}scott\sphinxhyphen{}para\sphinxhyphen{}mostrar\sphinxhyphen{}como\sphinxhyphen{}nid22012024/

\item {} 
\sphinxAtStartPar
Tipos de licencias de software: software libre, propietario y demás. (2022, July 19). Velneo.com. https://velneo.com/blog/licencias\sphinxhyphen{}software\sphinxhyphen{}libre\sphinxhyphen{}propietario\sphinxhyphen{}otros/

\item {} 
\sphinxAtStartPar
Universidad Nacional del Nordeste. (n.d.). Informática. Conceptos fundamentales. https://exa.unne.edu.ar/ingenieria/computacion/Tema1.pdf

\item {} 
\sphinxAtStartPar
Velásquez, W. (2022, April 1). 3 historias del software libre que deberías conocer. Wajari. https://wajari.com/blog/historias\sphinxhyphen{}software\sphinxhyphen{}libre\sphinxhyphen{}open\sphinxhyphen{}source/

\item {} 
\sphinxAtStartPar
(Visnuh), M. G. (2024, January 24). Con Microsoft Office 365 dominarás la productividad y con el precio de esta licencia en oferta no tendrás que empeñarte para ello. Xataka.com; Xataka. https://www.xataka.com/seleccion/microsoft\sphinxhyphen{}office\sphinxhyphen{}365\sphinxhyphen{}dominaras\sphinxhyphen{}productividad\sphinxhyphen{}precio\sphinxhyphen{}esta\sphinxhyphen{}licencia\sphinxhyphen{}oferta\sphinxhyphen{}no\sphinxhyphen{}tendras\sphinxhyphen{}que\sphinxhyphen{}empenarte\sphinxhyphen{}para\sphinxhyphen{}ello

\end{itemize}



\renewcommand{\indexname}{Índice}
\printindex
\end{document}